% =============================================================================
% THÈSE DE DOCTORAT - CRYPTOGRAPHIE INSTITUTIONNELLE
% =============================================================================
% Auteur  : Thierry Emmanuel Minka Mi Nguidjoï
% École   : ENSP / Université de Yaoundé I
% Version : 3.4 - Sommaires locaux avec minitoc
% Date    : 2026-09-20
% =============================================================================
\UseRawInputEncoding
%\RequirePackage[l2tabu]{nag}
\newcommand{\hmmax}{0}
\newcommand{\bmmax}{0}
\documentclass[a4paper,11pt,reqno,oneside]{memoir}
%\usepackage{fourier-orns}
\usepackage{style/these_ensp}

\addbibresource{annexes/bibliographie.bib}

\title{AUX FONDEMENTS DE LA CRYPTOGRAPHIE INSTITUTIONNELLE :\\
       Formalisation, Bornes d'Impossibilité, Conditions d'Optimalité\\
       et Vérification Formelle}
\author{MINKA MI NGUIDJOÏ Thierry Emmanuel}
\date{2025--2026}

% =============================================================================
\begin{document}

\ancrerTableaux
\activerBarrieresFlottants
% =============================================================================

% =============================================================================
% PRÉLIMINAIRES
% =============================================================================
\frontmatter
% -----------------------------------------------------------------------------
% COUVERTURES
% -----------------------------------------------------------------------------
\renewcommand{\enspDiplome}{Doctorat/Ph.D en G\'enie Informatique}
\renewcommand{\enspSpecialite}{Cryptographie et S\'ecurit\'e de l'Information}

\couvertureENSP%
  {AUX FONDEMENTS DE LA CRYPTOGRAPHIE INSTITUTIONNELLE :}%
  {Formalisation, Bornes d'Impossibilit\'e, Conditions d'Optimalit\'e
   et V\'erification Formelle}%
  {MINKA MI NGUIDJO\"I Thierry Emmanuel}%
  {17P318}%
  {Prof. Thomas \textbf{DJOTIO NDI\'E}, Ma\^itre de Conf\'erences, UYI}%
  {Dr. Flavien Serge \textbf{MANI ONANA}, Charg\'e de Cours, UYI}%
  {%
    \small
    \begin{tabular}{@{}lll@{}}
      \textbf{Pr\'esident :}    & Prof. [Prenom] \textbf{[NOM]}   & [Grade], [Institution]         \\[0.2cm]
      \textbf{Rapporteur 1 :}  & Prof. [Prenom] \textbf{[NOM]}    & [Grade], [Institution]  \\[0.2cm]
      \textbf{Rapporteur 2 :}  & Dr. [Prenom] \textbf{[NOM]}   & [Grade], [Institution]         \\[0.2cm]
      \textbf{Examinateur :} & Prof. [Prenom] \textbf{[NOM]}    & [Grade], [Institution]       \\[0.2cm]
      \textbf{Membre 1 :} & Prof.   [Prenom] \textbf{[NOM]}    & [Grade], [Institution]         \\[0.2cm]
      \textbf{Membre 2 :}    & Prof.    [Prenom] \textbf{[NOM]}  & [Grade], [Institution]       \\
    \end{tabular}%
  }%
  {2025-2026}%
  {jour mois 2026}%

\couvertureENSPblanche%
  {AUX FONDEMENTS DE LA CRYPTOGRAPHIE INSTITUTIONNELLE :}%
  {Formalisation, Bornes d'Impossibilit\'e, Conditions d'Optimalit\'e 
  et V\'erification Formelle}%
  {MINKA MI NGUIDJO\"I Thierry Emmanuel}%
  {2025-2026}%
  {jour mois 2026}%

% -----------------------------------------------------------------------------
% PAGES PRÉLIMINAIRES
% -----------------------------------------------------------------------------
\pagestyle{preliminaire}
% ============================================================================
% DEDICACES
% Fichier : preambule/dedicaces.tex
% ============================================================================

\chapter*{D\'EDICACE}
\markboth{D\'edicace}{}
\addcontentsline{toc}{chapter}{D\'edicace}
\thispagestyle{preliminaire}
\vspace*{\fill}
\begin{flushleft}
\Large
\textit{\`A \textbf{Minka Emmanuel} dit Herm\'etique.}\\
\vspace{0.3cm}
\textit{{\color{enspblue}\guillemotleft}~Grand-p\`ere, puisse ce travail continuer de faire vivre ton nom.~{\color{enspblue}\guillemotright}}
\end{flushleft}
\vspace*{\fill}

% ============================================================================
% REMERCIEMENTS (Acrostiche M-E-R-C-I)
% Fichier : preambule/remerciements.tex
% ============================================================================

\chapter*{REMERCIEMENTS}
\markboth{Remerciements}{}
\addcontentsline{toc}{chapter}{Remerciements}
\thispagestyle{preliminaire}
\begin{sloppypar}
\initial{M}es remerciements les plus sinc\`eres s'adressent premi\`erement \`a mes directeurs de th\`ese, le Professeur \textbf{Thomas DJOTIO NDI\'E} et le Docteur \textbf{Flavien Serge MANI ONANA}, pour leur encadrement rigoureux et leurs expertises combin\'ees. Je tiens \'egalement \`a exprimer ma profonde gratitude aux membres du jury, xxxxxxxxxxxxxxxxxxxxxxxxxxxxxxxxxxxxxxxxxxxxxxxxxxxxxxxxxxxxxxxxxxxxxxxxxxxx xxxxxxxxxxxxxxxxxxxxxxxxxxxxxxxxxxxxxxxxxxxxxxxx, qui ont accept\'e d'\'evaluer ce travail. Dans la même lanc\'ee, Professeurs \textbf{Georges KOUAMOU, Achille MELINGUI} et \textbf{Isidore NGONGO} de
l'\'Ecole Nationale Supérieure Polytechnique de Yaoundé I, merci pour vos remarques constructives lors du jury d'audition. Elles ont \'et\'e essentielles pour porter ce travail \`a ce niveau d'exigence.
\end{sloppypar}
\vspace{0.5cm}

\initial{E}nvers les \textbf{Services du Contr\^ole Sup\'erieur de l'\'Etat (CONSUPE)}, mon employeur actuel, j'exprime toute ma reconnaissance pour la bienveillance qui m'a permis de concilier responsabilit\'es administratives et exigences doctorales. Ma gratitude s'\'etend aux magistrats civils et militaires avec qui j'ai eu l'honneur d'\'echanger en ma qualit\'e d'Expert Judiciaire. Il en va de m\^eme pour mes \'echanges nourris avec certains responsables du Minist\`ere des Finances et de l'H\^opital G\'en\'eral de Garoua. Ces interactions ont ancr\'e ma r\'eflexion sur la cryptographie institutionnelle dans la r\'ealit\'e juridique et les enjeux de souverainet\'e num\'erique.
\vspace{0.5cm}
\initial{R}echercher l'excellence m'a conduit \`a m'impliquer activement dans la communaut\'e scientifique internationale, notamment au sein de l'\textbf{IACR}, de l'\textbf{IEEE} et de l'\textbf{ISACA} dont je suis membre. Mon appartenance \`a ces organisations refl\`ete ma volont\'e de lier la recherche acad\'emique aux standards professionnels de la s\'ecurit\'e de l'information. Je remercie \'egalement tous mes co-auteurs. Puisse notre collaboration continuer. \`A tous les contributeurs de la communaut\'e \textbf{Lean4} pour leur aide pr\'ecieuse dans la formalisation de mes preuves de s\'ecurit\'e, je dis simplement \textit{you rock!}
\vspace{0.5cm}
\initial{C}e travail est aussi le fruit de l'\'emulation v\'ecue au sein du \textbf{Laboratoire d'Ing\'enierie Math\'ematique et des Syst\`emes d'Information (LIMSI)}. Je pense aux encouragements du Professeur \textbf{Thomas BOUETOU}, aux discussions chaleureuses avec mes coll\`egues doctorants, aux critiques constructives des enseignants du D\'epartement de G\'enie Informatique pendant les diff\'erentes rencontres doctorales, mais également \`a mes \'etudiants de diverses promotions (\textbf{GI} et \textbf{HN}). Savoir que certains d'entre vous me consid\`erent comme un mod\`ele est un honneur qui m'engage quotidiennement.
\vspace{0.5cm}
\initial{I}ndispensable a \'et\'e le soutien de mon cercle le plus cher, ma famille. \`A mon \'epouse Elisabeth, pour sa patience~; \`a mes fils David et Axel~; \`a mes filles Sheila et Bissa\"i, \textbf{les MINKA} qui m'ont souvent vu inaccessible \`a cause de ce projet~; et \`a mes parents, \textbf{les Nguidjo\"i}, dont la fiert\'e est ma plus belle r\'ecompense. \textbf{Que ce travail ouvre grand les portes du doctorat pour notre lign\'ee}.
\vspace{0.4cm}
\begin{flushright}
\textit{\`A tous et \`a chacun, du plus profond de mon \^etre, \textbf{MERCI !}}
\end{flushright}
% ============================================================================
% GLOSSAIRE DES ACRONYMES ET NOTATIONS
% Fichier : preambule/glossaire.tex
% ============================================================================

\chapter*{GLOSSAIRE DES ACRONYMES ET DES NOTATIONS}
\markboth{Glossaire des acronymes et des notations}{}
\addcontentsline{toc}{chapter}{Glossaire des acronymes et des notations}
\thispagestyle{preliminaire}
\label{chap:glossaire}

\section*{ACRONYMES}
\subsection*{A-C}
\begin{sloppypar}
\begin{description}[style=nextline,leftmargin=2.4cm,font=\normalfont\bfseries]
\item[AML] Anti-Money Laundering (Lutte contre le blanchiment d'argent)
\item[Bâle III] Accords de Bâle III
\item[CEMAC] Communaut\'e \'Economique et Mon\'etaire de l'Afrique Centrale (\textit{Economic and Monetary Community of Central Africa})
\item[CIC] Calcul des Constructions Inductives (\textit{Calculus of Inductive Constructions}), syst\`eme de type sous-jacent \`a Lean4 et Coq
\item[CLO] Cryptographic Legal Opposability (Opposabilit\'e L\'egale Cryptographique), cadre th\'eorique de cette th\`ese
\item[CPTP] Completely Positive Trace-Preserving (Canal quantique compl\`etement positif pr\'eservant la trace)
\item[CRO] Confidentiality-Reliability-Opposability (Confidentialit\'e-Fiabilit\'e-Opposabilit\'e), trilemme fondamental
\item[$\mathsf{Conf}$] M\'etrique de confidentialit\'e (\textit{Confidentiality}), voir D\'efinition~\ref{def:confidentialite_clo}
\end{description}

\subsection*{D-F}

\begin{description}[style=nextline,leftmargin=2cm,font=\normalfont\bfseries]
\item[DDH] Decisional Diffie-Hellman (Hypoth\`ese de d\'ecision Diffie-Hellman)
\item[DLog] Discrete Logarithm (Logarithme discret)
\item[DPI] Data Processing Inequality (In\'egalit\'e du traitement des donn\'ees), $I(X;Z) \leq I(X;Y)$ pour $X \to Y \to Z$
\item[ECDSA] Elliptic Curve Digital Signature Algorithm (Algorithme de signature num\'erique sur courbe elliptique)
\item[eIDAS] Electronic IDentification, Authentication and trust Services (Identification \'electronique, authentification et services de confiance), R\`eglement EU 910/2014
\item[ENSP] \'Ecole Nationale Sup\'erieure Polytechnique (\textit{National Advanced School of Engineering})
\item[ESIGN] Electronic Signatures in Global and National Commerce Act, loi am\'ericaine de 2000
\item[FHE] Fully Homomorphic Encryption (Chiffrement totalement homomorphe)
\item[FIPS] Federal Information Processing Standards, standards f\'ed\'eraux am\'ericains (FIPS 203, 204, 205 pour PQC)
\item[FLP] Fischer-Lynch-Paterson (th\'eor\`eme d'impossibilit\'e du consensus asynchrone)
\end{description}

\subsection*{G-L}

\begin{description}[style=nextline,leftmargin=2cm,font=\normalfont\bfseries]
\item[GDPR] General Data Protection Regulation (R\`eglement g\'en\'eral sur la protection des donn\'ees), RGPD, R\`eglement EU 2016/679
\item[GMR] Goldwasser-Micali-Rackoff (auteurs originaux des preuves \`a divulgation nulle)
\item[Groth16] Sch\'ema zk-SNARK de Groth (2016), preuves succinctes de 128 octets (BN254, forme compress\'ee) ou 192 octets (BLS12-381)
\item[IACR] International Association for Cryptologic Research
\item[IND-CCA2] Indistinguishability under Adaptive Chosen Ciphertext Attack (Indistinguabilit\'e sous attaque adaptative \`a texte chiffr\'e choisi)
\item[IND-CPA] Indistinguishability under Chosen Plaintext Attack (Indistinguabilit\'e sous attaque \`a texte clair choisi)
\item[JL] Johnson-Lindenstrauss (transformation de compression dimensionnelle)
\item[KEM] Key Encapsulation Mechanism (M\'ecanisme d'encapsulation de cl\'e)
\item[KL] Kullback-Leibler (divergence de), $D_{\mathrm{KL}}(P\|Q) = \sum_x P(x)\log\frac{P(x)}{Q(x)}$
\item[LIMSI] Laboratoire d'Ing\'enierie Math\'ematique et Syst\`emes d'Information, laboratoire de th\`ese
\item[LWE] Learning With Errors (Apprentissage avec erreurs), probl\`eme cryptographique post-quantique
\end{description}

\subsection*{M-P}

\begin{description}[style=nextline,leftmargin=2cm,font=\normalfont\bfseries]
\item[MAP] Maximum A Posteriori (Estimateur du maximum a posteriori)
\item[Mathlib] Biblioth\`eque math\'ematique de Lean4
\item[ML-DSA] Module-Lattice-Based Digital Signature Algorithm, FIPS 204, bas\'e sur Dilithium
\item[ML-KEM] Module-Lattice-Based Key Encapsulation Mechanism, FIPS 203, bas\'e sur Kyber
\item[MPC] Multi-Party Computation (Calcul multi-parties)
\item[NIST] National Institute of Standards and Technology (Institut national des normes et de la technologie)
\item[NISQ] Noisy Intermediate-Scale Quantum (Appareils quantiques interm\'ediaires bruit\'es)
\item[NIZK] Non-Interactive Zero-Knowledge (Preuve à divulgation nulle non-interactive)
\item[OTP] One-Time Pad (Masque jetable), chiffrement parfait de Shannon
\item[$\mathsf{Opp}$] M\'etrique d'opposabilit\'e (\textit{Opposability}), voir D\'efinition~\ref{def:opposabilite_clo}
\item[Pedersen] Sch\'ema de commitment de Pedersen (1991), $g^m h^r$ parfaitement cachant
\item[PKE] Public-Key Encryption (Chiffrement \`a cl\'e publique)
\item[PKI] Public Key Infrastructure (Infrastructure \`a cl\'es publiques)
\item[Poseidon] Fonction de hachage optimis\'ee pour SNARKs
\item[POVM] Positive Operator-Valued Measure (Mesure \`a valeur op\'erateur positif), mesure quantique g\'en\'eralis\'ee
\item[PPT] Probabilistic Polynomial Time (Temps polynomial probabiliste)
\item[PQC] Post-Quantum Cryptography (Cryptographie post-quantique)
\item[PRG] Pseudorandom Generator (G\'en\'erateur pseudoal\'eatoire)
\end{description}

\subsection*{Q-S}

\begin{description}[style=nextline,leftmargin=2.5cm,font=\normalfont\bfseries]
\item[QPT] Quantum Polynomial Time (Temps polynomial quantique), adversaire quantique born\'e
\item[$\mathsf{Rel}$] M\'etrique de fiabilit\'e (\textit{Reliability}), voir D\'efinition~\ref{def:fiabilite_clo}
\item[RGPD] R\`eglement G\'en\'eral sur la Protection des Donn\'ees, voir GDPR
\item[RSA] Rivest-Shamir-Adleman (algorithme de chiffrement)
\item[SLH-DSA] Stateless Hash-Based Digital Signature Algorithm, FIPS 205, bas\'e sur SPHINCS+
\item[SNARK] Succinct Non-interactive ARgument of Knowledge (Argument succinct non-interactif de connaissance)
\item[SOTP] Semantic One-Time Pad (Masque jetable s\'emantique), famille de protocoles de cette th\`ese
\item[SOTP-Base] SOTP pour r\'egime de transition ($H^*(\lambda) \leq H(C) \leq \lambda/\log\lambda$)
\item[SOTP-High] SOTP pour r\'egime \`a contexte riche ($H(C) > \lambda/\log\lambda$)
\item[SOTP-Low] SOTP pour r\'egime \`a contexte rare ($H(C) < H^*(\lambda)$)
\item[STARK] Scalable Transparent ARgument of Knowledge (Argument transparent \'evolutif de connaissance)
\end{description}

\subsection*{T-Z}

\begin{description}[style=nextline,leftmargin=2cm,font=\normalfont\bfseries]
\item[TV] Total Variation (Distance de variation totale), $\mathsf{TV}(P,Q) = \frac{1}{2}\sum_x|P(x)-Q(x)|$
\item[UY1] Universit\'e de Yaound\'e I (\textit{University of Yaound\'e I})
\item[ZK] Zero-Knowledge (\`a divulgation nulle de connaissance)
\item[zk-SNARK] Zero-Knowledge Succinct Non-interactive ARgument of Knowledge (Argument succinct non-interactif de connaissance \`a divulgation nulle)
\end{description}

% ============================================================================
% SYMBOLES MATHÉMATIQUES
% ============================================================================

\newpage
\section*{SYMBOLES MATH\'EMATIQUES ET NOTATIONS}

\subsection*{Param\`etres institutionnels}

\begin{description}[style=nextline,leftmargin=1.6cm,font=\normalfont\bfseries]
\item[$\Delta$] Capacit\'e d'extraction institutionnelle (bits), $I(V;\Sigma|M,C) \leq \Delta$~; quelques dizaines de bits en calibration illustrative
\item[$\Delta_q$] Capacit\'e d'extraction institutionnelle adapt\'ee au quantum, $\Delta_q = \Delta + \delta_{\mathrm{enh}} - \delta_{\mathrm{degr}}$
\item[$\alpha$] Granularit\'e des verdicts institutionnels, $\log_2|\mathcal{V}_\lambda| \leq \alpha H(C)$
\item[$\beta$] Constante de r\'egularit\'e du protocole, $\mathrm{KL}(P_{\Sigma|m_0,c}\|P_{\Sigma|m_1,c}) \leq \beta \cdot I(M;\Sigma|C)$
\item[$\eta_q$] Coefficient de perte d'interpr\'etabilit\'e quantique, $I(V;\Lambda(\Sigma)|C) \geq I(V;\Sigma|C) - \eta_q H(C)$~; $\eta_q = 0$ pour une preuve qui est une cha\^{\i}ne de bits
\item[$H(C)$] Entropie du contexte institutionnel, $H(C) = H(L) + H(T|L) + H(R|L,T)$
\end{description}

\subsection*{Param\`etres cryptographiques}

\begin{description}[style=nextline,leftmargin=1.5cm,font=\normalfont\bfseries]
\item[$\lambda$] Param\`etre de s\'ecurit\'e cryptographique, typiquement $\lambda \in \{128, 192, 256\}$
\item[$\mathsf{negl}(\lambda)$] Fonction n\'egligeable en $\lambda$, $\forall c>0, \exists \lambda_0, \forall \lambda > \lambda_0: |f(\lambda)| < \lambda^{-c}$
\item[$c_0$] Constante de Pinsker, $c_0 = 1/(2\ln 2) \approx 0{,}721$ (classique) ou $c_0 \approx 0{,}3466$ (forme alternative)
\item[$H^*(\lambda)$] Crit\`ere de construction SOTP, $H^*(\lambda) = \sqrt{\lambda}$ (taille de l'encodage direct du contexte)
\item[$H^\sharp$] Seuil de la forme raffin\'ee, $H^\sharp = (\Delta + 1 + \kappa/\lambda)/\alpha$~; au-del\`a de ce seuil, la capacit\'e institutionnelle est contraignante (Remarque~\ref{rem:activite_delta})
\item[$H^\sharp_q$] Seuil quantique, $H^\sharp_q = (\Delta_q + \kappa/\lambda)/(\theta_q\alpha_q)$ (Remarque~\ref{rem:forme_raffinee_pq})
\end{description}

\subsection*{M\'etriques CRO}

\begin{description}[style=nextline,leftmargin=1.7cm,font=\normalfont\bfseries]
\item[$\mathsf{Conf}(\mathcal{P})$] Confidentialit\'e du protocole $\mathcal{P}$, $\mathsf{Conf} = 1 - \max_{m_0,m_1,c}\mathsf{TV}(P_{\Sigma|m_0,c}, P_{\Sigma|m_1,c})$
\item[$\Relrec(\mathcal{P})$] Fiabilit\'e de reconstruction, $\Relrec = 1 - \min_{\hat M}\Pr[\hat{M}(\Sigma,C) \neq M]$ (contr\^ol\'ee par Fano)
\item[$\Relver(\mathcal{P})$] Fiabilit\'e de v\'erification, compl\'etude et solidit\'e du v\'erificateur
\item[$\mathsf{Opp}(\mathcal{P})$] Opposabilit\'e du protocole $\mathcal{P}$, $\mathsf{Opp} = I(V;\Sigma|C)/H(V|C)$
\item[$\mathsf{Conf}^q$] Confidentialit\'e quantique, $1 - \max \tfrac12\|\rho_{m_0,c}-\rho_{m_1,c}\|_1$ (statistique)
\item[$\mathsf{Rel}^q$] Fiabilit\'e quantique, recouvrement optimal via POVM
\item[$\mathsf{Opp}^q$] Opposabilit\'e quantique, information extraite par un interpr\`ete admissible, rapport\'ee \`a $H(V|C)$
\end{description}

\subsection*{Variables al\'eatoires et espaces}

\begin{description}[style=nextline,leftmargin=1cm,font=\normalfont\bfseries]
\item[$M$] Variable al\'eatoire message, espace $\mathcal{M}_\lambda$ avec $|\mathcal{M}_\lambda| = 2^\lambda$
\item[$\Sigma$] Variable al\'eatoire preuve cryptographique, espace $\Sigma_\lambda$
\item[$C$] Variable al\'eatoire contexte institutionnel, espace $\mathcal{C}_\lambda$ avec distribution $P_C$
\item[$V$] Variable al\'eatoire verdict institutionnel, $V = \mathsf{Interpret}(\Sigma, C)$
\item[$\mathcal{Y}_\lambda$] Espace des verdicts, $|\mathcal{Y}_\lambda| \leq 2^{\alpha H(C)}$
\item[$\rho$] \'Etat quantique (matrice de densit\'e)
\item[$\Lambda$] Canal quantique (application CPTP)
\end{description}

\subsection*{Mesures d'information}

\begin{description}[style=nextline,leftmargin=2cm,font=\normalfont\bfseries]
\item[$H(X)$] Entropie de Shannon, $H(X) = -\sum_x P_X(x)\log_2 P_X(x)$
\item[$H(X|Y)$] Entropie conditionnelle, $H(X|Y) = H(X,Y) - H(Y)$
\item[$I(X;Y)$] Information mutuelle, $I(X;Y) = H(X) - H(X|Y) = H(Y) - H(Y|X)$
\item[$I(X;Y|Z)$] Information mutuelle conditionnelle, $I(X;Y|Z) = H(X|Z) - H(X|Y,Z)$
\item[$D_{\mathrm{KL}}(P\|Q)$] Divergence de Kullback-Leibler, $D_{\mathrm{KL}}(P\|Q) = \sum_x P(x)\log\frac{P(x)}{Q(x)}$
\item[$\mathsf{TV}(P,Q)$] Distance de variation totale, $\mathsf{TV}(P,Q) = \frac{1}{2}\sum_x|P(x)-Q(x)|$
\item[$S(\rho)$] Entropie de von Neumann, $S(\rho) = -\mathrm{Tr}(\rho\log\rho)$
\item[$\chi$] Information de Holevo, $\chi = S(\bar{\rho}) - \sum_x P_X(x)S(\rho_x)$
\item[$I_{\mathrm{acc}}$] Information accessible, supr\'emum sur toutes les POVMs
\end{description}

\subsection*{Notions quantiques}

\begin{description}[style=nextline,leftmargin=1.6cm,font=\normalfont\bfseries]
\item[$\|\rho - \sigma\|_1$] Distance de trace quantique, $\mathrm{Tr}\sqrt{(\rho-\sigma)^\dagger(\rho-\sigma)}$
\item[$F(\rho,\sigma)$] Fid\'elit\'e quantique, $F(\rho,\sigma) = \|\sqrt{\rho}\sqrt{\sigma}\|_1$
\item[$\|\Lambda\|_\diamond$] Norme diamant, mesure de distance entre canaux quantiques
\item[$\Lambda_{\mathrm{dep}}$] Canal d\'epolarisant, $\Lambda_{\mathrm{dep}}(\rho) = (1-p)\rho + \frac{p}{d}I_d$
\item[$\Lambda_{\mathrm{AD}}$] Canal d'amortissement d'amplitude (amplitude damping)
\item[$\Lambda_{\mathrm{deph}}$] Canal de d\'ephasing (d\'ecoh\'erence de phase)
\end{description}

\subsection*{Formalisation Lean4}

\begin{description}[style=nextline,leftmargin=4.3cm,font=\normalfont\bfseries]
\item[\texttt{ProtocolFamily}] Structure Lean4 repr\'esentant une famille de protocoles index\'ee par $\lambda$
\item[\texttt{ProtocolInstance}] Structure Lean4 pour une instance \`a $\lambda$ fix\'e
\item[\texttt{InstitutionalParams}] Structure Lean4 pour les param\`etres $(\Delta, \alpha, \beta, H)$
\item[\texttt{Negligible}] Pr\'edicat Lean4 pour les fonctions n\'egligeables
\item[\texttt{atTop}] Filtre Lean4 pour les comportements asymptotiques
%\item[Quasi-optimalit\'e] Hypoth\`ese $I(M;\Sigma|C) \leq 1$ pour $\lambda$ suffisamment grand
\end{description}

% ============================================================================
% CONCEPTS CLÉS
% ============================================================================

\newpage
\section*{CONCEPTS CL\'ES DE LA CRYPTOGRAPHIE INSTITUTIONNELLE}

\subsection*{Paradigme architectural}

\begin{description}[style=nextline,leftmargin=2.5cm,font=\normalfont\bfseries]
\item[Gold Layer] Couche de confidentialit\'e, masquage cryptographique via PRG ou ZK-proofs
\item[Iron Layer] Couche de fiabilit\'e, commitments assurant le binding
\item[Clay Layer] Couche d'opposabilit\'e, encodage pr\'eservant le contexte pour l'extraction institutionnelle
\end{description}

\subsection*{Th\'eor\`emes fondamentaux}

%\subsection*{Th\'eor\`emes fondamentaux}

\begin{description}[style=nextline,leftmargin=6cm,font=\normalfont\bfseries]
\item[Trilemme CRO] $\mathsf{Conf} \cdot \Relrec \cdot \mathsf{Opp} \leq \frac{1}{\gamma}\bigl(\frac{\Delta}{\lambda} + \frac{\kappa(\Delta+1)}{\lambda^2} + \frac{\kappa^2}{\lambda^3}\bigr)$, $\gamma = \alpha H(\mathcal{C}) - 1$~; forme raffin\'ee $\frac{1+\kappa/\lambda}{\lambda}\min\{1, \frac{\kappa/\lambda+\Delta}{\gamma}\}$
\item[Trilemme CRO-PQ] $\mathsf{Conf}^q \cdot \Relrec^q \cdot \mathsf{Opp}^q \leq \frac{1}{\gamma_q}\bigl(\frac{\Delta_q}{\lambda} + \dots\bigr)$~; majoration affaiblie $(\Delta_q + \eta_q H(C))/\lambda + O(\lambda^{-2})$
\item[Goulot d'\'etranglement contextuel] $I(V;\Sigma|C) \leq \min\{H(\Sigma|C), H(V|C)\} \leq \alpha H(\mathcal{C})$
\item[In\'egalit\'e de Fano] $H(X|\hat{X}) \leq h(p_e) + p_e\log(|\mathcal{X}|-1)$
\item[In\'egalit\'e de Pinsker] $\mathsf{TV}(P,Q)^2 \leq c_0 \cdot D_{\mathrm{KL}}(P\|Q)$
\item[In\'egalit\'e de Gentle Measurement] $\|\rho - \rho'\|_1 \leq 2\sqrt{\delta}$ pour mesure \`a probabilit\'e $\geq 1-\delta$
\item[Fuchs-van de Graaf] $1 - \frac{1}{2}\|\rho-\sigma\|_1 \leq F(\rho,\sigma) \leq \sqrt{1 - \frac{1}{4}\|\rho-\sigma\|_1^2}$
\end{description}

\subsection*{R\'egimes op\'erationnels}

\begin{description}[style=nextline,leftmargin=4cm,font=\normalfont\bfseries]
\item[Low-Capacity] $\Delta < 0{,}1\lambda$, protocoles transparents prioritaires
\item[Mid-Capacity] $0{,}1\lambda \leq \Delta \leq 0{,}5\lambda$, compromis \'equilibr\'es possibles
\item[High-Capacity] $\Delta > 0{,}5\lambda$, protocoles pr\'eservant la confidentialit\'e possibles
\item[R\'egime rare] $H(C) \leq H^\sharp$ (ou $H^\sharp_q$), l'in\'egalit\'e de Fano limite le produit
\item[R\'egime riche] $H(C) > H^\sharp$ (ou $H^\sharp_q$), la capacit\'e $\Delta$ (ou $\Delta_q$) limite le produit
\end{description}

\subsection*{Syst\`emes analys\'es}

\begin{description}[style=nextline,leftmargin=2.5cm,font=\normalfont\bfseries]
\item[Helios] Protocole de vote \`a v\'erification de bout en bout, $\mathsf{Conf} \approx 0{,}92$, $\Relver \approx 0{,}95$, $\mathsf{Opp} \approx 0{,}50$ (calibration)
\item[zk-SNARKs] Preuves succinctes de conformit\'e, $\mathsf{Conf} \approx 0{,}89$, $\Relver \approx 0{,}96$, $\mathsf{Opp} \approx 0{,}67$ (calibration)
\item[ECDSA] Signatures num\'eriques sur courbes elliptiques, $\mathsf{Conf} \approx 0{,}31$, $\Relver \approx 0{,}98$, $\mathsf{Opp} \approx 0{,}90$ (calibration)
\item[Dilithium] Signature PQC bas\'ee sur les r\'eseaux, standardis\'ee sous le nom ML-DSA (FIPS 204)
\item[Kyber] KEM PQC bas\'e sur les r\'eseaux, standardis\'e sous le nom ML-KEM (FIPS 203)~; hors du champ du Trilemme
\item[SPHINCS+] Signature PQC bas\'ee sur les hachages (FIPS 205)
\item[Falcon] Signature PQC bas\'ee sur NTRU (standardisation en cours)
\end{description}
\end{sloppypar}
% ============================================================================
% FIN GLOSSAIRE
% ============================================================================

% ============================================================================
% PREAMBULE
% Fichier : preambule/preambule.tex
% ============================================================================

\chapter*{PR\'EAMBULE}
\markboth{Pr\'eambule}{}
\addcontentsline{toc}{chapter}{Pr\'eambule}
\thispagestyle{preliminaire}

\begin{flushright}
  \begin{minipage}{0.72\textwidth}
    \itshape
    \guillemotleft~La cryptographie n'existe pas dans le vide. Elle est au service
    d'institutions humaines qui op\`erent sous des contraintes cognitives,
    organisationnelles et proc\'edurales.~\guillemotright
  \end{minipage}
\end{flushright}
\vspace{0.3cm}
\begin{flushright}
  \small\textit{- Principe fondateur de la Cryptographie
  Institutionnelle~\cite{MINKA2026CLO}}
\end{flushright}

\vspace{1.0cm}

Un arbitre institutionnel, qu'il soit juge, auditeur ou r\'egulateur, se trouve
confront\'e \`a une preuve fond\'ee sur des principes cryptographiques.
L'algorithme a effectu\'e toutes les v\'erifications n\'ecessaires et la
s\'ecurit\'e du sch\'ema n'est pas en cause. L'arbitre se trouve pourtant
dans l'incapacit\'e de rendre sa d\'ecision. Sa comp\'etence n'est pas en
cause. La preuve, bien que math\'ematiquement irr\'eprochable, reste opaque
du point de vue institutionnel. Elle atteste de quelque chose sans en pr\'eciser la nature,
dans les termes requis pour pouvoir statuer. Ce d\'ecalage entre la
rigueur du protocole et la clart\'e du verdict constitue le point de
d\'epart de cette th\`ese.

Le terme \guillemotleft~arbitre institutionnel~\guillemotright~fait r\'ef\'erence
\`a ce destinataire final, volontairement g\'en\'erique, qu'il s'agisse d'un
juge civil ou p\'enal, d'un auditeur des finances publiques, d'une commission
\'electorale ou d'un r\'egulateur sectoriel. Cet arbitre n'est ni cryptographe
ni sp\'ecialiste en informatique. Il est celui pour qui la preuve est
\'elabor\'ee. Sa conviction est la base de la d\'ecision qui s'impose. La
cryptographie a longtemps n\'eglig\'e cet aspect. Ce travail vise \`a en faire
une pr\'eoccupation centrale de la discipline.

Ce personnage est pr\'esent tout au long de la th\`ese. Au d\'ebut de chaque
chapitre, une note succincte permet de le repositionner dans le cadre de la
progression canonique
$\mathrm{F}{\to}\mathrm{B}{\to}\mathrm{E}{\to}\mathrm{S}{\to}\mathrm{C}{\to}\mathrm{X}$,
repr\'esentant respectivement formaliser, borner, \'etendre, saturer,
certifier, puis eXp\'erimenter. Ce fil directeur rappelle que la
preuve est con\c{c}ue pour un acteur humain dont la d\'ecision d\'epasse le
simple calcul.

L'arbitre n'est pas l'auteur de cette thèse. Le parcours professionnel décrit à la sous-section~\ref{subsubsec:experience_auteur}
a fourni le terrain d'observation, dont ce personnage est l'abstraction théorique.

\medskip
\noindent\textit{Cette sc\`ene ouvre la th\`ese.}

\vspace{0.8cm}

\begin{boiteArbitre}
L'arbitre franchit le seuil de la salle. Devant lui se pr\'esente une preuve
cryptographique attestant l'int\'egrit\'e de milliers de transactions. Un
algorithme en a valid\'e la v\'eracit\'e en quelques millisecondes. Aucun
adversaire op\'erant en temps polynomial ne saurait la contrefaire. Pourtant,
il lui revient de prononcer un verdict, ce qui suppose de saisir pr\'ecis\'ement ce
que la preuve \'etablit, au-del\`a du seul fait qu'une certification a eu
lieu. Or il en est incapable. Sur le plan
math\'ematique, la preuve est irr\'eprochable, tandis que, du point de vue
institutionnel, elle demeure opaque. Ce paradoxe, qu'il \'eprouve sans
parvenir \`a le formuler, constitue l'objet de cette th\`ese.
\end{boiteArbitre}
% ============================================================================
% R\'ESUM\'E FRAN\c{C}AIS
% Fichier : preambule/resume.tex
% ============================================================================

\chapter*{R\'ESUM\'E}
\markboth{R\'esum\'e}{}
\addcontentsline{toc}{chapter}{R\'esum\'e}
\thispagestyle{preliminaire}
\begin{sloppypar}
La cryptographie classique appr\'ecie la s\'ecurit\'e d'une preuve \`a
l'aune de sa robustesse face \`a des adversaires computationnellement
born\'es. \`A l'inverse, le paradigme institutionnel impose que la preuve
demeure intelligible pour un arbitre non sp\'ecialiste afin d'\'eclairer
une d\'ecision. Les d\'eploiements op\'erationnels r\'ev\`elent une
pathologie r\'ecurrente. Une preuve irr\'eprochable du point de vue
math\'ematique peut se r\'ev\'eler inutilisable par l'agent charg\'e de
trancher. Cette tension structurelle entre les deux paradigmes constitue
le paradoxe de l'authenticit\'e invisible, selon lequel plus une preuve
prot\`ege efficacement, plus elle devient opaque pour un tiers profane.

La th\`ese introduit une m\'etrique informationnelle in\'edite,
l'opposabilit\'e Opp, d\'efinie comme la proportion d'incertitude
institutionnelle qu'une preuve parvient \`a dissiper. Elle propose le
cadre CLO, qui agr\`ege confidentialit\'e (Conf), fiabilit\'e (Rel) et
opposabilit\'e (Opp) et l'articule \`a des param\`etres cognitifs
explicites, en particulier la capacit\'e d'extraction institutionnelle
$\Delta$ et la granularit\'e des verdicts $\alpha$. Elle \'etablit un
trilemme d'impossibilit\'e conditionnel. Sous quatre hypoth\`eses bornant
les capacit\'es de l'institution, le produit Conf$\cdot$Rel$\cdot$Opp
est major\'e par $\Delta/(\gamma\lambda)$, o\`u $\gamma = \alpha H(C) - 1$,
\`a un terme d'ordre $\lambda^{-2}$ et \`a un reste n\'egligeable pr\`es.
Une forme raffin\'ee identifie le r\'egime o\`u la capacit\'e
institutionnelle devient la contrainte liante. La ressource limitante
est alors la quantit\'e d'information qu'une institution peut
extraire humainement d'une preuve, bien plus que la taille des cl\'es.

Ce r\'esultat se prolonge au contexte post-quantique. Dans le sc\'enario
d'un support ou d'un v\'erificateur quantique, le bruit induit une perte
d'interpr\'etabilit\'e $\eta_q$, nulle pour les preuves classiques, dont
l'effet sur la borne cro\^it avec l'entropie contextuelle. La th\`ese
propose aussi la famille SOTP, param\'etr\'ee pour diff\'erents
r\'egimes de richesse contextuelle, dont elle \'etablit les propri\'et\'es
cryptographiques. Son optimalit\'e asymptotique est formul\'ee comme
conjecture. Le noyau alg\'ebrique du
trilemme est certifi\'e m\'ecaniquement dans Lean4, qui explicite
hypoth\`eses, constantes et domaine de validit\'e et fournit deux
t\'emoins formels garantissant la coh\'erence interne et la port\'ee du
r\'esultat.

Enfin, la th\`ese sp\'ecifie un programme exp\'erimental r\'efutable
comprenant six protocoles destin\'es \`a mesurer les param\`etres
institutionnels, avec estimateurs et seuils de r\'efutation, puis une
calibration illustrative de ces param\`etres sur des syst\`emes
d\'eploy\'es. Elle v\'erifie la plausibilit\'e des ordres de grandeur sans
constituer une validation empirique, le programme n'ayant pas \'et\'e ex\'ecut\'e en conditions
prospectives.

En \'erigeant l'opposabilit\'e au rang de propri\'et\'e cryptographique de
premi\`ere classe, ce travail constitue la cryptographie institutionnelle
comme domaine th\'eorique autonome. Il fournit des instruments d'aide au
choix des protocoles dans des secteurs sensibles tels que le vote
\'electronique, la conformit\'e r\'eglementaire et la transition
post-quantique et pr\'epare une normalisation future int\'egrant
l'interpr\'etabilit\'e parmi les crit\`eres de s\'ecurit\'e.

\vspace{0.5cm}
\noindent\textbf{Mots-cl\'es~:} cryptographie institutionnelle,
opposabilit\'e, Trilemme CRO, protocoles SOTP, v\'erification formelle,
programme exp\'erimental falsifiable, Lean4.
\end{sloppypar}
% ============================================================================
% ABSTRACT
% Fichier : preambule/abstract.tex
% ============================================================================

\chapter*{ABSTRACT}
\markboth{Abstract}{}
\addcontentsline{toc}{chapter}{Abstract}
\thispagestyle{preliminaire}

\begin{sloppypar}

Classical cryptography evaluates the security of a proof by the yardstick
of its robustness against computationally bounded adversaries. Conversely,
the institutional paradigm requires that the proof remain intelligible to
a non-specialist arbitrator in order to inform a decision. Operational
deployments reveal a recurrent pathology. A proof that is irreproachable
from the mathematical point of view may prove unusable to the agent tasked
with deciding. This structural tension between the two paradigms
constitutes the paradox of invisible authenticity, according to which the
more effectively a proof protects, the more opaque it becomes to a lay
third party.

The thesis introduces a novel informational metric, opposability Opp,
defined as the proportion of institutional uncertainty that a proof
succeeds in dissipating. It proposes the CLO framework, which aggregates
confidentiality (Conf), reliability (Rel) and opposability (Opp) and
articulates it with explicit cognitive parameters, in particular the
institutional extraction capacity $\Delta$ and the granularity of verdicts
$\alpha$. It establishes a conditional impossibility trilemma. Under four
hypotheses bounding the capacities of the institution, the product
Conf$\cdot$Rel$\cdot$Opp is bounded above by $\Delta/(\gamma\lambda)$,
where $\gamma = \alpha H(C) - 1$, up to a term of order $\lambda^{-2}$ and
a negligible remainder. A refined form identifies the regime in which institutional capacity
becomes the binding constraint. The limiting resource is then the quantity
of information that an institution can humanly extract from a proof, far
more than the size of the keys.

This result extends to the post-quantum setting. In the scenario of a
quantum support or verifier, the noise induces an interpretability loss $\eta_q$, which vanishes for
classical proofs and whose effect on the bound grows with the contextual
entropy. The thesis also proposes the SOTP family,
parameterised for different regimes of contextual richness, and
establishes its cryptographic properties. Its asymptotic optimality is
stated as a conjecture. The
algebraic core of the trilemma is mechanically certified in Lean4, which
makes explicit the hypotheses, constants and domain of validity and provides two formal witnesses guaranteeing the internal consistency and
the scope of the result.

Finally, the thesis specifies a refutable experimental programme
comprising six protocols intended to measure the institutional parameters,
with estimators and refutation thresholds, followed by an illustrative calibration of
these parameters on deployed systems. It checks the plausibility of the
orders of magnitude without constituting an empirical validation, since
the programme has not been executed under prospective conditions.

By raising opposability to the rank of a first-class cryptographic
property, this work constitutes institutional cryptography as an
autonomous theoretical domain. It provides instruments to guide the choice
of protocols in sensitive sectors such as electronic voting, regulatory
compliance and the post-quantum transition and prepares a future
standardisation integrating interpretability among the security criteria.

\vspace{0.5cm}
\noindent\textbf{Keywords:} Institutional cryptography, opposability,
CRO Trilemma, SOTP protocols, formal verification, falsifiable experimental
programme, Lean4.

\end{sloppypar}


% -----------------------------------------------------------------------------
% TABLE DES MATIÈRES GÉNÉRALE
% -----------------------------------------------------------------------------
\clearpage
\markboth{Table des mati\`eres}{}
\activerStylesChapitres
\tableofcontents
\thispagestyle{chapitre}


% -----------------------------------------------------------------------------
% LISTES
% -----------------------------------------------------------------------------
\clearpage
\markboth{Table des figures}{}
\listoffigures
\thispagestyle{chapitre}

\clearpage
\markboth{Liste des tableaux}{}
\listoftables
\thispagestyle{chapitre}

\clearpage
\markboth{Liste des listings}{}
\lstlistoflistings
\thispagestyle{chapitre}

\clearpage
\markboth{Liste des algorithmes}{}
\listofalgorithms
\thispagestyle{chapitre}
% =============================================================================
% CORPS DU DOCUMENT
% =============================================================================
\mainmatter
\setcounter{tocdepth}{3}
\setsecnumdepth{subsubsection}
\pagestyle{fancy}
%\input{annexes/Resume_etendu}

% -----------------------------------------------------------------------------
% INTRODUCTION GÉNÉRALE
% -----------------------------------------------------------------------------
\chapitreNonNumerote{INTRODUCTION G\'EN\'ERALE}{IG}
% ============================================================================
% INTRODUCTION G\'eN\'eRALE
% source: intro_conclu_gen/introduction_generale.tex
% Version : 2.1 — alignée sur la formalisation Lean4 révisée
% ============================================================================

%\chapter*{INTRODUCTION G\'EN\'ERALE}
%\addcontentsline{toc}{chapter}{Introduction g\'en\'erale}
\label{chap:introduction_generale}

% ----------------------------------------------------------------------------
\section{CONTEXTE}
\label{sec:ig_contexte}
% ----------------------------------------------------------------------------
\begin{sloppypar}

\subsection{La cryptographie classique}

La mission de la cryptographie s'est profond\'ement transform\'ee. Depuis les
travaux fondateurs de Diffie et Hellman~\cite{DiffieHellman76}, elle consistait
\`a prot\'eger des donn\'ees contre des adversaires computationnellement born\'es.
Les protocoles \'etaient alors jug\'es selon un crit\`ere unique, la
s\'ecurit\'e computationnelle~\cite{GoldwasserMicali84,Goldreich01}, que le
l\'egislateur a fini par prendre en compte. L'Electronic Signatures in Global and National Commerce (ESIGN) Act
am\'ericain (2000)~\cite{ESIGN00} puis le r\`eglement electronic IDentification, Authentication and trust Services (eIDAS) europ\'een
(2014)~\cite{eIDAS14} ont ainsi consacr\'e l'\'equivalence juridique des signatures
num\'eriques, en pariant que la rigueur math\'ematique suffisait \`a fonder la
confiance institutionnelle~\cite{Kerr04,Mason07}.

Ce pari se v\'erifie pour les signatures. Il ne r\'esiste plus en revanche d\`es
que le protocole prot\`ege activement l'information qu'un juge, un auditeur ou un
r\'egulateur doit interpr\'eter. Les preuves \`a divulgation nulle de connaissance~\cite{GMR85},
le chiffrement homomorphe~\cite{Gentry09,BrakerskiVaikuntanathan11} et les engagements avanc\'es~\cite{Pedersen91}
garantissent math\'ematiquement ce qu'ils cachent, sans garantir pour autant que
ce qu'ils r\'ev\`elent soit lisible par un non-expert~\cite{Bonneau12,Herley09}. De ce
d\'ecalage na\^it la tension centrale de cette th\`ese.

Nous avons nomm\'e cette tension \textbf{paradoxe de l'authenticit\'e invisible}~\cite{MINKA2026CLO}. Selon ce paradoxe, plus une preuve cryptographique est math\'ematiquement
impeccable, plus elle devient institutionnellement opaque, c'est-\`a-dire qu'elle perd son sens pour l'interpr\`ete humain non expert. Une meilleure documentation ne suffit pas \`a corriger cette opacit\'e~\cite{Whitten99,Adams1999}, qui est de nature structurelle~\cite{Barredo2020,Adadi2018} comme le montre la section~\ref{subsec:paradoxe_authenticite}.

\subsection{L'urgence post-quantique}

En ao\^ut 2024, le National Institute of Standards and Technology (NIST) a publi\'e les trois premi\`eres normes de cryptographie post-quantique~\cite{NISTPQC,Alagic2020,NISTPostQuantum2022}, \`a savoir FIPS~203 (ML-KEM, bas\'e sur Kyber~\cite{Kyber2022}), FIPS~204 (ML-DSA, bas\'e sur Dilithium~\cite{Dilithium2022}) et FIPS~205 (SLH-DSA, bas\'e sur SPHINCS+~\cite{SPHINCS2022}). Avec cette standardisation, la migration vers la cryptographie post-quantique a quitt\'e le terrain de la sp\'eculation th\'eorique pour devenir une r\'ealit\'e op\'erationnelle.

Cette transition d\'epasse le simple remplacement d'algorithmes, car elle rompt avec plusieurs paradigmes \'etablis. Les adversaires classiques sont limit\'es par la puissance de calcul, alors que les adversaires quantiques exploitent des requ\^etes en superposition~\cite{Boneh2011,Watrous2009} qui invalident les preuves de s\'ecurit\'e conventionnelles. La migration aggrave donc le paradoxe de l'authenticit\'e invisible, puisque la complexit\'e structurelle des sch\'emas post-quantiques (fond\'es sur les r\'eseaux ou sur les fonctions de hachage) entre en tension avec les besoins institutionnels. La mani\`ere dont ces effets quantiques modifient la valeur probante des preuves et leur compr\'ehension par les interpr\`etes institutionnels (juge, auditeur, r\'egulateur) demeure \`a ce jour largement inexplor\'ee.

\end{sloppypar}

% ----------------------------------------------------------------------------
\section{PROBL\'EMATIQUE CENTRALE}
\label{sec:ig_problematique}
% ----------------------------------------------------------------------------
\begin{sloppypar}

Le foss\'e entre s\'ecurit\'e cryptographique et utilit\'e institutionnelle ne
manque pas de t\'emoins. En 2009, la commission \'electorale de l'Universit\'e Catholique
de Louvain s'est trouv\'ee dans l'incapacit\'e d'exploiter le syst\`eme Helios sans recourir
\`a une expertise externe~\cite{Adida09UCL}. En novembre 2025, l'International Association
for Cryptologic Research (IACR) a annul\'e son \'election apr\`es la perte irr\'em\'ediable
d'une cl\'e partielle~\cite{IACR2025incident}.
De leur c\^ot\'e, les r\'egulateurs financiers destinataires de preuves zk-SNARK~\cite{Groth16,BenSasson2019,PLONK20} ne peuvent
pas cibler une transaction suspecte pr\'ecise, car l'opacit\'e qui prot\`ege les clients
fait aussi obstacle \`a l'investigation~\cite{Casey16,Garratt19}. Ces incidents, qui ne rel\`event pas de d\'efauts de mise en
\oe uvre, traduisent une incompatibilit\'e structurelle entre deux paradigmes
de v\'erification.

La cryptographie classique mesure l'\emph{avantage de l'adversaire}~\cite{GMR88,Shoup01},
d\'efini \`a partir d'un distingueur polynomial confront\'e \`a un d\'efi.
La v\'erification institutionnelle repose sur d'autres exigences, \`a savoir l'intelligibilit\'e
pour un juge, l'auditabilit\'e r\'etrospective par un contr\^oleur et la tra\c{c}abilit\'e par un r\'egulateur~\cite{Citron14,Casey11}.
Aucune m\'etrique ne rend compte de ces propri\'et\'es dans la litt\'erature cryptographique
classique~\cite{GMR85,Goldreich01}.

Ces constats conduisent \`a la probl\'ematique centrale de cette th\`ese, formul\'ee dans l'encadr\'e ci-dessous.

\begin{boiteProblematique}{Probl\'ematique centrale}
\textit{Comment concilier la s\'ecurit\'e cryptographique, qui requiert l'opacit\'e computationnelle, avec l'interpr\'etabilit\'e institutionnelle, qui requiert la transparence s\'emantique~? Existe-t-il des limites fondamentales aux compromis r\'ealisables entre confidentialit\'e, fiabilit\'e et opposabilit\'e~?}
\end{boiteProblematique}

Cette th\`ese formalise ces propri\'et\'es ($\Conf$~\ref{def:confidentialite_clo}, $\Rel$~\ref{def:fiabilite_clo}, $\Opp$~\ref{def:opposabilite_clo}) et traite cette probl\'ematique au moyen de \textbf{six questions/objectifs de recherche}, qui structurent nos contributions.

\medskip
\begin{boiteQuestion}{Obj1. Formalisation}
\textit{\'Etablir un cadre unifi\'e mesurant simultan\'ement la confidentialit\'e, la fiabilit\'e et l'opposabilit\'e d'un protocole cryptographique, en s'appuyant sur la th\'eorie de l'information.}
\end{boiteQuestion}

\medskip
\begin{boiteQuestion}{Obj2. Bornes d'impossibilit\'e classiques}
\textit{\'Etablir les limites fondamentales des compromis r\'ealisables, conditionnellement \`a la capacit\'e cognitive des institutions.}
\end{boiteQuestion}

\medskip
\begin{boiteQuestion}{Obj3. Extension post-quantique}
\textit{D\'eterminer comment les adversaires quantiques et les effets de canal quantique modifient ces bornes puis caract\'eriser les termes qui leur sont propres.}
\end{boiteQuestion}

\medskip
\begin{boiteQuestion}{Obj4. Constructions optimales}
\textit{Construire des protocoles qui atteignent asymptotiquement ces bornes afin de transformer une barri\`ere th\'eorique en objectif d'ing\'enierie.}
\end{boiteQuestion}

\medskip
\begin{boiteQuestion}{Obj5. Certification formelle}
\textit{Certifier m\'ecaniquement le noyau du th\'eor\`eme d'impossibilit\'e dans un assistant de preuve en explicitant les hypoth\`eses dont il d\'epend.}
\end{boiteQuestion}


\medskip
\begin{boiteQuestion}{Obj6. Exp\'erimentation}
\textit{Sp\'ecifier un dispositif exp\'erimental falsifiable permettant de mesurer les param\`etres institutionnels qui conditionnent ces bornes, puis confronter le cadre aux syst\`emes d\'eploy\'es.}
\end{boiteQuestion}

\medskip
La progression Obj1~$\to$~Obj6 organise la th\`ese en six \'etapes, qui
consistent \`a \textbf{formaliser, borner, \'etendre, saturer, certifier et
exp\'erimenter}. La probl\'ematique \'enonc\'ee plus haut demeure la question
unique de ce travail, \`a laquelle la r\'ealisation successive de ces six
objectifs apporte une r\'eponse.

\end{sloppypar}

% ----------------------------------------------------------------------------
\section{POSITIONNEMENT \'EPIST\'EMOLOGIQUE}
\label{sec:ig_positionnement}
% ----------------------------------------------------------------------------
\begin{sloppypar}

La \textbf{Cryptographie Institutionnelle} se situe \`a l'intersection de la th\'eorie
de l'information, de la cryptographie classique et post-quantique, du droit et des sciences cognitives~\cite{Simon82,Cowan2001,Ericsson95}.
Ce croisement est impos\'e par la nature du probl\`eme, car aucune de ces disciplines prise isol\'ement ne permet
de formaliser le compromis entre opacit\'e cryptographique et intelligibilit\'e
institutionnelle~\cite{Citron14,Casey10,Casey11}.

\begin{table}[htbp]
\centering
\caption{Disciplines fondatrices et leurs contributions \`a la Cryptographie Institutionnelle.}
\label{tab:disciplines_fondatrices}
{\footnotesize
\begin{tabular}{p{3.8cm}p{5.2cm}p{3.2cm}}
\toprule
\textbf{Discipline} & \textbf{Apport} & \textbf{Usage dans la th\`ese} \\
\midrule
Th\'eorie de l'information & Entropie, information mutuelle~\cite{CoverThomas06,Wyner75,WynerZiv76} & M\'etriques CLO \\
Th\'eorie de l'information quantique & von Neumann, Holevo~\cite{NielsenChuang2010,Wilde2017,Hiai1991} & Extension PQ \\
Cryptographie classique & Primitives, r\'eductions~\cite{Goldreich01,GoldwasserMicali84} & Cadre CRO \\
Cryptographie post-quantique & Lattices, NIST PQC~\cite{NISTPQC,Kyber2022,Dilithium2022,Falcon2022} & Trilemme PQ-CRO \\
Sciences cognitives & Rationalit\'e born\'ee~\cite{Simon82,Miller56,Kahneman11} & Param\`etre $\Delta$ \\
Droit & eIDAS, preuve juridique~\cite{eIDAS14,GDPR2016,ESIGN00,Daubert1993} & Contexte institutionnel \\
V\'erification formelle & Lean4, Mathlib~\cite{moura2015lean,Lean4,Mathlib,Wiedijk2007} & Certification \\
\bottomrule
\end{tabular}
}
\end{table}

Cette th\`ese r\'eunit la rigueur math\'ematique, la pertinence
institutionnelle et la certification formelle~\cite{Gonthier08,Leroy09,BhargavanEA2017}
dans une m\^eme d\'emarche, qui fournit les bases sur lesquelles
les parties suivantes construisent puis valident la th\'eorie.

\end{sloppypar}

% ----------------------------------------------------------------------------
\section{LECTEUR CIBLE ET PR\'EREQUIS}
\label{sec:ig_lecteur}
% ----------------------------------------------------------------------------
\begin{sloppypar}

Cette th\`ese s'adresse \`a plusieurs publics, dont les points d'entr\'ee diff\`erent.

\begin{itemize}
\item \textbf{Cryptographes.} Les sections~\ref{sec:cro_classique} et~\ref{sec:sotp_construction} contiennent des preuves originales utilisant des techniques standard (Fano~\cite{Fano61}, Pinsker~\cite{Pinsker1964}, r\'eductions cryptographiques~\cite{Shoup01,Damgard87}).

\item \textbf{Chercheurs en s\'ecurit\'e formelle.} La section~\ref{sec:lean4} fournit une formalisation compl\`ete avec axiomatisation explicite~\cite{BartheEA2013,Lochbihler2016,Affeldt2015}.

\item \textbf{Juristes et r\'egulateurs.} La section~\ref{sec:clo} montre l'applicabilit\'e aux syst\`emes de vote~\cite{Adida08,Benaloh94}, de conformit\'e~\cite{FATF12,AMLD4} et de preuve num\'erique~\cite{Mason07,Casey16}.

\item \textbf{Sp\'ecialistes post-quantique.} La section~\ref{sec:cro_pq} \'etend le cadre aux adversaires quantiques~\cite{Watrous2009,Unruh2016,Boneh2011} et aux standards NIST~\cite{NISTPQC,NISTPostQuantum2022}.
\end{itemize}

La lecture suppose une familiarit\'e avec les preuves \`a divulgation nulle~\cite{GMR85,BenSasson2019,Campanelli2019}, des notions de base de th\'eorie de l'information (entropie, information mutuelle)~\cite{CoverThomas06,Smith09} ainsi qu'une connaissance \'el\'ementaire de la th\'eorie de la complexit\'e. Le glossaire~\ref{chap:glossaire} plac\'e en d\'ebut de th\`ese d\'efinit l'ensemble des acronymes et notations utilis\'es.

\end{sloppypar}

\begin{boiteArbitre}
L'arbitre ne sait pas encore qu'il est pris dans un paradoxe~;
il sait seulement que quelque chose ne fonctionne pas entre la preuve qu'on
lui soumet et la d\'ecision qu'on lui demande de rendre. Ce chapitre lui donne
le nom du probl\`eme et la carte du voyage, avec six questions, six
contributions et un cycle. Il lui pr\'esente aussi les \'equipiers dont il aura
besoin, \`a savoir la th\'eorie de l'information, la cryptographie, les
sciences cognitives, le droit et la v\'erification formelle. Le voyage commence
\`a la section~\ref{sec:contexte}.
\end{boiteArbitre}

% ----------------------------------------------------------------------------
\section{CONTRIBUTIONS}
\label{sec:ig_contributions}
% ----------------------------------------------------------------------------
\begin{sloppypar}

\paragraph{C1. Nous d\'efinissons le cadre Cryptographic Legal Opposability (CLO)~\cite{MINKA2026CLO}~(section~\ref{sec:clo}).}
Publi\'e dans \emph{IEEE Access} (DOI: 10.1109/ACCESS.2026.3660105), cet article
introduit trois m\'etriques issues de la th\'eorie de l'information~\cite{Shannon48,CoverThomas06,Wyner75}, \`a savoir la confidentialit\'e
$\Conf$, la fiabilit\'e $\Rel$ et l'opposabilit\'e
\begin{equation}
  \Opp(\mathcal{P}) = \frac{I(V;\Sigma \mid C)}{H(V \mid C)},
\end{equation}
qui quantifie la fraction d'incertitude institutionnelle r\'esolue par la preuve
cryptographique. L'opposabilit\'e n'a pas d'\'equivalent dans la litt\'erature
ant\'erieure~\cite{Smith09,Alvim12}. Par le th\'eor\`eme du goulot d'\'etranglement contextuel,
$I(V;\Sigma|C) \leq \alpha H(\mathcal{C})$, l'opposabilit\'e accessible se trouve reli\'ee
\`a la granularit\'e des verdicts~\cite{ClarksonSchneider10}. La th\`ese s\'epare aussi
la fiabilit\'e de reconstruction, que contr\^ole l'in\'egalit\'e de Fano, de la fiabilit\'e
de v\'erification.

\paragraph{C2. Nous \'etablissons le Trilemme Confidentiality Reliability Opposability (CRO) conditionnel~\cite{MINKA2026CRO}~(section~\ref{sec:cro_classique}).}
Publi\'e dans \emph{IEEE Access} (DOI: 10.1109/ACCESS.2026.3669065), ce r\'esultat prouve une borne d'impossibilit\'e
\begin{equation}
  \Conf(\mathcal{P})\cdot\Rel(\mathcal{P})\cdot\Opp(\mathcal{P})
  \;\leq\;
  \frac{\Delta}{\gamma\lambda} + O(\lambda^{-2}),
  \qquad \gamma = \alpha H(\mathcal{C}) - 1,
\end{equation}
conditionnelle \`a la capacit\'e d'extraction institutionnelle $\Delta$, \`a la granularit\'e des verdicts $\alpha$, \`a la r\'egularit\'e du protocole $\beta$ ainsi qu'au r\'egime de confidentialit\'e $I(M;\Sigma|C) \leq \kappa/\lambda$. La forme publi\'ee dans~\cite{MINKA2026CRO}, $\Delta/\lambda + O(\lambda^{-2}) + \negl(\lambda)$, en est un corollaire. Une forme raffin\'ee identifie exactement le r\'egime o\`u la capacit\'e institutionnelle est la contrainte liante. L'hypoth\`ese d'extraction born\'ee re\c{c}oit aussi un fondement informationnel exact. Une calibration illustrative situe $\Delta$ dans quelques dizaines de bits, ordre de grandeur plausible au regard de la rationalit\'e born\'ee de Simon~\cite{Simon82,Simon1955} et des limites de m\'emoire de travail~\cite{Miller56,Cowan2001}. Cette impossibilit\'e reste
conditionnelle, \`a la mani\`ere de CAP, qui d\'epend des partitions r\'eseau~\cite{Brewer00,GilbertLynch02}, de FLP,
qui d\'epend de l'asynchronie~\cite{FLP85}, ou du th\'eor\`eme d'Arrow, qui suppose des pr\'ef\'erences rationnelles~\cite{Arrow50}.

\paragraph{C3. Nous \'etendons le Trilemme au contexte quantique~\cite{MINKA2026PQCRO}~(section~\ref{sec:cro_pq}).}
Publi\'e dans \emph{IEEE Access} (DOI: 10.1109/ACCESS.2026.3722742), ce r\'esultat \'etablit l'extension quantique du trilemme CRO
\begin{equation}
  \Conf^q(\mathcal{P})\cdot\Rel^q(\mathcal{P})\cdot\Opp^q(\mathcal{P})
  \;\leq\;
  \frac{\Delta_q + \eta_q \cdot H(C)}{\lambda} + O(\lambda^{-2}) + \negl(\lambda),
\end{equation}
o\`u $\Delta_q$ quantifie la capacit\'e d'extraction institutionnelle adapt\'ee au quantum et $\eta_q$ le coefficient de perte d'interpr\'etabilit\'e quantique. La preuve mobilise des outils de th\'eorie de l'information quantique~\cite{NielsenChuang2010,Wilde2017}, en particulier l'in\'egalit\'e de Pinsker quantique, le lemme de Gentle Measurement~\cite{Schumacher1997} et les bornes de Fano quantiques~\cite{Hiai1991}. La borne explicite conserve la forme classique, avec $\Delta_q$~; le terme $\eta_q H(C)$ n'appara\^it que dans la majoration affaiblie ci-dessus et s'annule pour les preuves qui sont des cha\^ines de bits. Une calibration illustrative sur les signatures standardis\'ees ML-DSA et SLH-DSA~\cite{Dilithium2022,SPHINCS2022} situe $\Delta_q$ dans quelques dizaines de bits~\cite{Mayers1997,Lo1997,Unruh2016}.

\paragraph{C4. Nous proposons la famille de constructions Semantic One-Time Pad (SOTP)~(section~\ref{sec:sotp_construction}).}
Soumis \`a \emph{IEEE TIT}, cet article pr\'esente trois constructions de r\'ef\'erence couvrant le spectre
d'entropie contextuelle $H(\mathcal{C})$~:
\begin{itemize}
    \item \textbf{SOTP-Low} pour les contextes rares, fond\'e sur un encodage d\'eterministe du contexte et un engagement de Pedersen~;
    \item \textbf{SOTP-Base} pour les contextes interm\'ediaires, fond\'e sur une compression Johnson-Lindenstrauss~\cite{JohnsonLindenstrauss84,DasguptaGupta03,Vempala05}~;
    \item \textbf{SOTP-High} pour les contextes riches, fond\'e sur les zk-SNARKs Groth16~\cite{Groth16}.
\end{itemize}
Leurs propri\'et\'es cryptographiques (confidentialit\'e, fiabilit\'e de v\'erification) sont prouv\'ees sous des hypoth\`eses standard~; l'existence de constructions approchant la borne du Trilemme est formul\'ee comme conjecture. Ces constructions transforment l'impossibilit\'e en guide de conception~\cite{Gennaro13,Parno13,BCTV14}.

\paragraph{C5. Nous certifions m\'ecaniquement le noyau alg\'ebrique du
Trilemme en Lean4 (section~\ref{sec:lean4}).}
Soumis \`a \emph{Journal of Automated Reasoning}, cet article pr\'esente
la premi\`ere formalisation Lean4 compl\`ete d'un th\'eor\`eme
d'impossibilit\'e en cryptographie institutionnelle~\cite{moura2015lean,Lean4,Mathlib}.
Enti\`erement v\'erifi\'ee sans \texttt{sorry}, la formalisation int\`egre
les quatre ingr\'edients informationnels (Fano, Pinsker, extraction
born\'ee, r\'egularit\'e) comme propri\'et\'es des instances de protocole
tout en explicitant les hypoth\`eses institutionnelles. Elle \'etablit
formellement la borne $\Delta/(\gamma\lambda) + O(\lambda^{-2})$ et sa forme
raffin\'ee puis fournit deux t\'emoins de coh\'erence et de port\'ee. Le d\'eveloppement
constitue \'egalement une biblioth\`eque Lean4 r\'eutilisable, publi\'ee avec une version fig\'ee de Lean et
de Mathlib~\cite{CROTrilemmaDepot}.

\medskip
\noindent\textbf{Apport m\'ethodologique.} La formalisation en Lean4 a
conduit \`a expliciter, dans les preuves du Trilemme CRO classique~\cite{MINKA2026CRO}
et de son extension post-quantique~\cite{MINKA2026PQCRO}, des hypoth\`eses qui y
restaient implicites, en l'occurrence la condition d'efficacit\'e du verdict et le r\'egime de
fuite sous lequel la borne s'applique. Les \'enonc\'es retenus dans cette th\`ese
sont donc conditionnels et explicites, assortis de leurs constantes et de leur
domaine de validit\'e~; ce que la certification a ainsi rendu visible est
d\'etaill\'e \`a la sous-section~\ref{subsec:lean4_erreurs}. Cette exp\'erience
conduit \`a une le\c{c}on m\'ethodologique sur le r\^ole de la v\'erification
formelle dans la production des r\'esultats interdisciplinaires, discut\'ee \`a la
sous-section~\ref{subsec:verification_corrige_lecon}.

\paragraph{C6. Nous sp\'ecifions un programme de validation falsifiable
pour les param\`etres institutionnels et quantiques
(chapitre~\ref{chap:experimentation}).}
Ce chapitre sp\'ecifie, jusqu'au niveau op\'erationnel, les protocoles
exp\'erimentaux permettant de mesurer les param\`etres $\Delta$,
$\alpha$, $\beta$, $\Delta_q$ et $\eta_q$ sur des acteurs
institutionnels r\'eels. La stratification des
\'echantillons, les estimateurs, les crit\`eres de fiabilit\'e
inter-\'evaluateurs, les seuils de falsification et le calendrier de
d\'eploiement y sont pr\'ecis\'es. Comme le Trilemme CRO est conditionnel \`a ces
param\`etres, leur mesure en fait un programme
scientifique falsifiable au sens de Popper~\cite{Popper1934}, ce qui conditionne
la port\'ee empirique du Trilemme.
Le volet~B du chapitre d\'erive aussi les valeurs d'opposabilit\'e
des trois syst\`emes de r\'ef\'erence (ECDSA, Helios, zk-SNARKs), lesquels ne sont
pas des instances du Trilemme. Il en tire une calibration illustrative des
param\`etres institutionnels, en attente de validation exp\'erimentale. Le
programme n'ayant pas encore \'et\'e d\'eploy\'e dans le cadre de cette
th\`ese, son statut est explicit\'e \`a la
sous-section~\ref{subsubsec:statut_deploiement}.
\medskip
Ces six contributions forment le cycle repr\'esent\'e ci-dessous.
\begin{equation}
\label{eq:progression_FBESCX}
  \underbrace{\text{CLO}}_{\text{Formalisation}}
  \;\xrightarrow{\quad}\;
  \underbrace{\text{CRO}}_{\text{Bornes classiques}}
  \;\xrightarrow{\quad}\;
  \underbrace{\text{PQ-CRO}}_{\text{Extension quantique}}
  \;\xrightarrow{\quad}\;
  \underbrace{\text{SOTP}}_{\text{Constructions de r\'ef\'erence}}
  \;\xrightarrow{\quad}\;
  \underbrace{\text{Lean4}}_{\text{Certification}}
  \;\xrightarrow{\quad}\;
  \underbrace{\text{Exp\'erimentation}}_{\text{Protocoles et validation}}
\end{equation}

\paragraph{Progression canonique.}
\label{par:progression_FBESCX}
La th\`ese suit un fil conducteur en six actes, \textbf{F}ormaliser, \textbf{B}orner, \textbf{E}tendre, \textbf{S}aturer, \textbf{C}ertifier et e\textbf{X}p\'erimenter.
Le chapitre~\ref{chap:contexte} pose les fondements en amont de~F. Le
chapitre~\ref{chap:cadre_formel} r\'ealise F$\to$B$\to$E, puis le
chapitre~\ref{chap:constructions} r\'ealise S$\to$C. Le
chapitre~\ref{chap:experimentation} r\'ealise~X, tandis que le
chapitre~\ref{chap:resultats_discussions} analyse les r\'esultats et ouvre les perspectives.

\end{sloppypar}

% ----------------------------------------------------------------------------
\section{PUBLICATIONS ET STATUT}
\label{sec:ig_publications}
% ----------------------------------------------------------------------------
\begin{sloppypar}

La th\`ese s'appuie sur cinq manuscrits, trois publi\'es et deux soumis, auxquels s'ajoute un programme exp\'erimental sp\'ecifi\'e. Le tableau~\ref{tab:contributions_these} indique, pour chacune des six contributions, l'\'etape du cycle canonique qu'elle r\'ealise, la section qui la d\'eveloppe et son statut \`a la date de soutenance. Les autres parties de la th\`ese y renvoient.

\begin{table}[htbp]
\centering
\caption{\'Etat des contributions de la th\`ese \`a la date de soutenance}
\label{tab:contributions_these}
{\footnotesize
\begin{tabular}{p{0.6cm}p{3.6cm}p{0.7cm}p{1.2cm}p{1.4cm}p{5.2cm}}
\toprule
\textbf{R\'ef.} & \textbf{Contribution} & \textbf{\'Etape} & \textbf{Section} & \textbf{Statut} & \textbf{Revue et r\'ef\'erence} \\
\midrule
{[P1]} & C1~: cadre CLO & \textbf{F} & \ref{sec:clo} & Publi\'e & \textit{IEEE Access}, DOI 10.1109/ACCESS.2026.3660105~\cite{MINKA2026CLO} \\
{[P2]} & C2~: Trilemme CRO conditionnel & \textbf{B} & \ref{sec:cro_classique} & Publi\'e & \textit{IEEE Access}, DOI 10.1109/ACCESS.2026.3669065~\cite{MINKA2026CRO} \\
{[P3]} & C3~: Trilemme CRO post-quantique & \textbf{E} & \ref{sec:cro_pq} & Publi\'e & \textit{IEEE Access}, DOI 10.1109/ACCESS.2026.3722742~\cite{MINKA2026PQCRO} \\
-- & C4~: famille SOTP & \textbf{S} & \ref{sec:sotp_construction} & Soumis & \textit{IEEE Transactions on Information Theory}~\cite{MINKA2026SOTP} \\
-- & C5~: formalisation Lean4 & \textbf{C} & \ref{sec:lean4} & Soumis & \textit{Journal of Automated Reasoning}~\cite{MINKA2026LEAN4} \\
-- & C6~: programme de validation & \textbf{X} & \ref{sec:protocoles_experimentaux} & Sp\'ecifi\'e & D\'ep\^ot de recherche~\cite{CROTrilemmaDepot} \\
\bottomrule
\end{tabular}
}
\end{table}


\end{sloppypar}

% ----------------------------------------------------------------------------
\section{ORGANISATION DE LA TH\`ESE}
\label{sec:ig_organisation}
% ----------------------------------------------------------------------------
\begin{sloppypar}

Outre la pr\'esente introduction et la conclusion g\'en\'erale, qui ne sont
pas num\'erot\'ees, le manuscrit comprend cinq chapitres.

\textbf{Chapitre~\ref{chap:contexte} -- Contexte, \'etat de l'art et probl\'ematique.}
Trois cas illustratifs (Helios, ECDSA, zk-SNARK) servent d'ancrage \`a la recherche
dans la section~\ref{sec:contexte}, qui en calcule analytiquement les profils CLO.
La section~\ref{sec:etat_art} rassemble ensuite les fondements tir\'es de la th\'eorie
de l'information classique et quantique, des primitives d'obscurcissement, de la
rationalit\'e born\'ee, des trilemmes d'impossibilit\'e et de la v\'erification formelle.
Le chapitre se termine par la section~\ref{sec:problematique}, o\`u sont formul\'es le
paradoxe de l'authenticit\'e invisible, la probl\'ematique centrale et les
six objectifs de recherche.

\textbf{Chapitre~\ref{chap:cadre_formel} -- Cadre formel et bornes d'impossibilit\'e.}
Les actes \textbf{F}, \textbf{B} et \textbf{E} y sont men\'es successivement.
Apr\`es la formalisation du cadre CLO \`a la section~\ref{sec:clo}
(Obj~\ref{obj:obj1_unification}), la section~\ref{sec:cro_classique}
\'etablit le Trilemme CRO conditionnel classique
(Obj~\ref{obj:obj2_impossibilite}), que la section~\ref{sec:cro_pq} \'etend ensuite au
r\'egime post-quantique (Obj~\ref{obj:obj3_postquantique}).

\textbf{Chapitre~\ref{chap:constructions} -- Constructions et certification formelle.}
Il couvre les actes \textbf{S} et \textbf{C}. Le Semantic One-Time Pad, prototype de
la famille, est construit \`a la section~\ref{sec:sotp_base}. La section~\ref{sec:sotp_construction}
g\'en\'eralise ce prototype en trois constructions de r\'ef\'erence dont elle conjecture
l'optimalit\'e universelle (Obj~\ref{obj:obj4_saturation}), avant que la section~\ref{sec:lean4}
ne soumette le noyau alg\'ebrique du trilemme \`a l'assistant de preuve Lean4
(Obj~\ref{obj:obj5_verification}).

\textbf{Chapitre~\ref{chap:experimentation} -- Protocoles de mesure et validation r\'etrospective.}
L'acte \textbf{X} (Obj~\ref{obj:obj6_experimentation}) fait l'objet de ce chapitre.
Six protocoles de mesure des param\`etres institutionnels sont sp\'ecifi\'es \`a la
section~\ref{sec:protocoles_experimentaux}. La
section~\ref{sec:derivation_opposabilite} d\'erive ensuite les valeurs d'opposabilit\'e
des syst\`emes de r\'ef\'erence, dont la
section~\ref{sec:validation_retrospective} tire une calibration
illustrative des param\`etres institutionnels en pr\'ecisant sa port\'ee et
ses limites.

\textbf{Chapitre~\ref{chap:resultats_discussions} -- R\'esultats, discussions et perspectives.}
La section~\ref{sec:resultats_theoriques} pr\'esente les validations
r\'etrospectives et le bilan formel du corpus. La
section~\ref{sec:discussions} interpr\`ete ces r\'esultats, les positionne
parmi les impossibilit\'es fondamentales et documente les limitations
reconnues. La section~\ref{sec:perspectives} ouvre le programme de recherche
en quatre axes.

\end{sloppypar}
 
\retablirNumerotation

% ==========================================================
% CHAPITRE 1 : CONTEXTE, ETAT DE L'ART ET PROBLEMATIQUE
% ==========================================================
\chapter{CONTEXTE, \'ETAT DE l'ART ET PROBL\'EMATIQUE}
        %{CONTEXTE, \'ETAT DE l'ART ET PROBL\'EMATIQUE}
\label{chap:contexte}
\debutChapitre%
   {Contexte, \'etat de l'art et probl\'ematique}%
  {%
    Ce chapitre pose les fondements th\'eoriques et empiriques de la
    recherche. Apr\`es avoir d\'ecrit le cadre institutionnel dans lequel
    s'inscrit ce travail, il parcourt l'\'etat de l'art interdisciplinaire
    sur lequel il s'appuie, qui couvre la th\'eorie de l'information classique
    et quantique, la cryptographie et les primitives d'obscurcissement, les
    trilemmes d'impossibilit\'e, les standards post-quantiques du NIST ainsi que
    les sciences cognitives. Il formule ensuite les six objectifs de recherche
    qui structurent la progression canonique de la th\`ese. Il met enfin en
    \'evidence le paradoxe de l'authenticit\'e invisible, selon lequel les
    protocoles cryptographiques les plus s\^urs, ceux qui maximisent la
    confidentialit\'e, induisent une opacit\'e institutionnelle s\'ev\`ere, alors
    que les protocoles transparents sacrifient la confidentialit\'e. Ce paradoxe,
    qui traverse l'ensemble de la contribution, motive l'introduction de la
    m\'etrique d'opposabilit\'e et structure la probl\'ematique centrale de la th\`ese.%
  }%
  {%
    cryptographie institutionnelle, opposabilit\'e juridique, paradoxe de
    l'authenticit\'e invisible, th\'eorie de l'information, rationalit\'e
    born\'ee, post-quantique%
  }%

% ============================================================================
% INTRODUCTION DU CHAPITRE 1 — CONTEXTE, \'ETAT DE L'ART ET PROBL\'eMATIQUE
% source : partie1/introduction_partie1.tex
% Version : 2.2 — alignée sur la formalisation Lean4 révisée
% ============================================================================

% ----------------------------------------------------------------------------
\section{INTRODUCTION}
\label{sec:introduction_contexte}
% ----------------------------------------------------------------------------

\begin{sloppypar}
En quatre d\'ecennies, la cryptographie prouvable a produit des garanties
d'une grande solidit\'e, qu'il s'agisse de la confidentialit\'e, mesur\'ee par
l'avantage d'un distingueur polynomial, ou de la fiabilit\'e, \'etablie par
r\'eduction \`a un probl\`eme computationnellement difficile. Ces r\'esultats,
exacts sur le plan math\'ematique, restent pourtant souvent inexploitables
lorsqu'une institution doit s'en saisir.

Un juge ne raisonne pas comme un distingueur polynomial, pas plus qu'un
auditeur financier ne v\'erifie un circuit arithm\'etique. Lorsque la preuve
cryptographique d'une transaction doit convaincre un r\'egulateur, ou qu'une
attestation zero-knowledge est vers\'ee au dossier d'une instruction p\'enale,
la question d\'ecisive devient celle de l'interpr\'etation de cette preuve
par un acteur institutionnel cognitivement born\'e. Cette th\`ese part de ce d\'eplacement de
perspective, dont le pr\'esent chapitre pose les fondements.

Trois syst\`emes cryptographiques d\'eploy\'es servent de construits
illustratifs, \`a savoir Helios pour le vote \'electronique, les zk-SNARKs
pour la conformit\'e financi\`ere et ECDSA pour les signatures
contractuelles. Leurs profils estim\'es laissent appara\^itre un motif que
l'\'etat de l'art interdisciplinaire permettra ensuite de nommer puis de
borner. La probl\'ematique est alors formul\'ee en six objectifs de
recherche qui structurent la progression canonique
$\mathbf{F}\to\mathbf{B}\to\mathbf{E}\to\mathbf{S}\to\mathbf{C}\to\mathbf{X}$
(Formaliser, Borner, \'Etendre, Saturer, Certifier, eXp\'erimenter) de
cette th\`ese. Apr\`es avoir situ\'e la recherche dans son cadre
institutionnel et professionnel (section~\ref{subsec:cadre_recherche}), nous
posons les fondements de la th\'eorie de l'information classique
(section~\ref{subsec:fondements_info}) avant de les \'etendre au cadre
quantique (section~\ref{subsec:fondements_info_quantique}). Les primitives
cryptographiques d'obscurcissement sont pr\'esent\'ees \`a la
section~\ref{subsec:primitives_crypto}, suivies d'une revue des trilemmes
d'impossibilit\'e fondateurs (section~\ref{subsec:trilemmes}). La
section~\ref{subsec:paradoxe_authenticite} formalise ensuite le paradoxe de
l'authenticit\'e invisible, dont la section~\ref{subsec:objectifs_recherche}
d\'erive les six objectifs de recherche et leurs contributions.
\end{sloppypar}
% ============================================================================
% CHAPITRE 1 — CONTEXTE, \'ETAT DE L'ART ET PROBL\'eMATIQUE
% Section 1.2 : Contexte
% source : partie1/contexte.tex
% ============================================================================

\section{CONTEXTE}
\label{sec:contexte}
\begin{sloppypar}
La cryptographie et le droit ont toujours entretenu une relation asym\'etrique~\cite{Mason07,Twining1999}.
La premi\`ere produit des garanties absolues, prouvables sous hypoth\`eses
math\'ematiques~\cite{Goldreich01,GMR88}, alors que le second exige des preuves que des
acteurs non-experts puissent comprendre et contester dans le cadre de
proc\'edures \'etablies~\cite{Daubert1993,Casey10}. Longtemps g\'erable, cette asym\'etrie s'est
creus\'ee avec les protocoles modernes au point de devenir une fracture structurelle~\cite{Citron14,Kerr04}.

Pour en d\'ecrire les contours, ce chapitre s'appuie sur trois syst\`emes
cryptographiques employ\'es comme construits illustratifs th\'eoriques, \`a savoir Helios pour le vote
\'electronique, les zk-SNARKs pour la conformit\'e financi\`ere et ECDSA pour
les signatures contractuelles. Leurs profils CLO ($\Conf$, $\Rel$, $\Opp$)
sont des estimations param\'etriques calcul\'ees analytiquement depuis leurs
sp\'ecifications publi\'ees, sous des hypoth\`eses institutionnelles d\'etaill\'ees
\`a la section~\ref{subsubsec:parametres_institutionnels}. Ils ne proviennent pas d'exp\'eriences
contr\^ol\'ees men\'ees avec des acteurs r\'eels. Ils servent \`a ancrer le paradoxe
dans des cas document\'es tout en motivant la formalisation des chapitres~\ref{chap:cadre_formel} \`a~\ref{chap:experimentation}.
\end{sloppypar}

% ----------------------------------------------------------------------------
\subsection{Cadre de recherche}
\label{subsec:cadre_recherche}
% ----------------------------------------------------------------------------
\begin{sloppypar}

\subsubsection{Le laboratoire LIMSI}
\label{subsubsec:limsi}

Les travaux pr\'esent\'es dans cette th\`ese ont \'et\'e conduits au
\textbf{Laboratoire d'Ing\'enierie Math\'ematique et Syst\`emes d'Information
(LIMSI)} de l'\'ecole Nationale Sup\'erieure Polytechnique de l'Universit\'e
de Yaound\'e~I, Cameroun. Le LIMSI accueille notamment des recherches situ\'ees \`a la
jonction de la cryptographie, des syst\`emes d'information, des m\'ethodes formelles et des math\'ematiques appliqu\'ees. Ce voisinage a permis de confronter les exigences de rigueur formelle aux difficult\'es que rencontre la pratique institutionnelle.

\subsubsection{L'exp\'erience professionnelle de l'auteur}
\label{subsubsec:experience_auteur}

Notre parcours professionnel associe plusieurs expertises compl\'ementaires, puisque nous exer\c{c}ons comme Auditeur des finances publiques au sein des \textbf{Services du Contr\^ole Sup\'erieur de l'\'Etat (CONSUPE)}, comme Expert judiciaire pr\`es les tribunaux du Cameroun ainsi que comme Expert certifi\'e en cybers\'ecurit\'e.

L'exercice de ces fonctions nous a souvent mis face \`a une situation d\'econcertante, dans laquelle des syst\`emes d'information robustes sur le plan cryptographique et s\^urs sur le plan math\'ematique ne parvenaient pas \`a convaincre les institutions qui devaient s'en servir~\cite{Bonneau16,Narayanan16}. En audit financier comme en expertise judiciaire, la validit\'e de la preuve technique suffisait rarement \`a fonder la confiance s\'emantique qu'exige une d\'ecision administrative ou juridictionnelle~\cite{Casey16,Grimmelmann15}.

\end{sloppypar}

% ----------------------------------------------------------------------------
\subsection{Gen\`ese de la question centrale de recherche}
\label{subsec:genese_question}
% ----------------------------------------------------------------------------
\begin{sloppypar}

La question centrale de cette th\`ese proc\`ede donc moins d'une r\'eflexion men\'ee \textit{in abstracto} que de l'observation r\'ep\'et\'ee d'un ph\'enom\`ene auquel nous avons donn\'e le nom de \textbf{paradoxe de l'authenticit\'e invisible}.

Le paradoxe tient \`a la conception m\^eme des protocoles cryptographiques actuels, qui ne prennent pas en compte la capacit\'e cognitive de leurs interpr\`etes~\cite{Simon82,Miller56,Herley09}, bien davantage qu'\`a un d\'efaut d'acculturation ou de formation des acteurs institutionnels.

Afin de d\'epasser le seul cadre national dans lequel ces constats avaient \'et\'e faits, nous avons retenu trois cas d'\'etude de port\'ee internationale (Helios pour le vote, les zk-SNARKs pour la conformit\'e, ECDSA pour la signature), dont les enseignements valent au-del\`a de leur contexte d'origine pour les questions de confiance num\'erique~\cite{DeFilippi16,Garay15}.

\end{sloppypar}

\begin{boiteArbitre}
L'arbitre a plusieurs dossiers sur sa table. Dans celui d'Helios figure une
\'election v\'erifiable de bout en bout que les garants \'electoraux n'ont pas su
d\'echiffrer sans assistance externe. Celui des zk-SNARKs contient une preuve de
quelques centaines d'octets qui atteste la conformit\'e de milliers de
transactions sans permettre au r\'egulateur de cibler la moindre op\'eration
suspecte. Le dossier d'ECDSA offre au contraire une transparence totale ainsi
qu'une tra\c{c}abilit\'e parfaite, obtenues au prix d'une confidentialit\'e nulle.
Chacun t\'emoigne d'une impuissance diff\'erente. Le pr\'esent chapitre les d\'ecrit
tandis que les chapitres suivants cherchent pourquoi elles convergent toutes
vers la m\^eme contrainte.
\end{boiteArbitre}

% ----------------------------------------------------------------------------
\subsection{Note de m\'ethodologie~: statut des exemples illustratifs}
\label{subsec:note_methodo}
% ----------------------------------------------------------------------------
\begin{sloppypar}
Cette th\`ese est une contribution th\'eorique. Les syst\`emes \'evoqu\'es au fil
du texte, Helios~\cite{Adida08,Adida09UCL}, ECDSA~\cite{Johnson01,NIST.FIPS.186-5},
zk-SNARKs~\cite{Groth16,Gennaro13,BCTV14} de conformit\'e, servent d'ancrages
illustratifs et de supports de validation th\'eorique, non de sujets d'\'etude
empirique. Les valeurs num\'eriques associ\'ees ($\Conf$, $\Rel$, $\Opp$) sont
des estimations param\'etriques d\'eriv\'ees analytiquement depuis les sp\'ecifications
publi\'ees de ces protocoles, sous des hypoth\`eses institutionnelles explicitement
formul\'ees (sous-section~\ref{subsubsec:hypotheses_capacite},
section~\ref{subsubsec:parametres_institutionnels}). Elles ne sont pas issues
d'exp\'eriences contr\^ol\'ees avec des acteurs institutionnels r\'eels. Ces
syst\`emes, dont la fiabilit\'e de v\'erification est proche de~$1$, ne sont pas
des instances du Trilemme CRO (Remarque~\ref{rem:portee_cas_etude}). Ils
servent \`a une \textbf{calibration illustrative} des ordres de grandeur des
param\`etres institutionnels, non \`a une validation, ni empirique au sens des
sciences exp\'erimentales~\cite{Tsybakov2009}, ni th\'eorique.
\end{sloppypar}

% ----------------------------------------------------------------------------
\subsection{Cas d'\'etude d\'eploy\'es}
\label{subsec:cas_etude}
% ----------------------------------------------------------------------------
\begin{sloppypar}
Les trois cas d'\'etude introduits \`a la section~\ref{subsec:genese_question} permettent de pr\'eciser le paradoxe. Leur choix tient \`a ce qu'ils
parcourent l'ensemble des compromis possibles entre confidentialit\'e et
opposabilit\'e, depuis les syst\`emes tr\`es confidentiels mais peu opposables
(Helios, zk-SNARKs) jusqu'\`a un syst\`eme tr\`es opposable dont la confidentialit\'e
est quasi nulle (ECDSA). Pris ensemble, ils font appara\^itre une r\'egularit\'e,
puisque aucun protocole d\'eploy\'e n'atteint un produit $\Conf \cdot \Rel \cdot \Opp$
sup\'erieur \`a $0{,}57$. Cette observation laisse supposer l'existence d'une
contrainte de fond.
\end{sloppypar}

% ----------------------------------------------------------------------------
\subsubsection{Vote \'electronique~: le cas Helios}
\label{subsubsec:cas_helios}
% ----------------------------------------------------------------------------
\begin{sloppypar}

\paragraph{Architecture cryptographique.}
\label{par:helios_architecture}

Le protocole Helios, con\c{c}u par Ben Adida en 2008~\cite{Adida08}, est
une r\'ef\'erence pour le vote \'electronique v\'erifiable de bout en bout~\cite{Benaloh94,Lindeman12}. Il
associe le chiffrement homomorphe ElGamal~\cite{ElGamal85}, gr\^ace auquel les
bulletins chiffr\'es peuvent \^etre agr\'eg\'es sans \^etre d\'echiffr\'es un \`a un, \`a des
preuves non-interactives zero-knowledge
(protocole de Chaum-Pedersen~\cite{ChaumPedersen92} rendu non interactif par Fiat-Shamir~\cite{FiatShamir87})
qui attestent la validit\'e de chaque bulletin sans en r\'ev\'eler le contenu.
Pour un bulletin binaire $m \in \{0,1\}$, le chiffr\'e est
$c = (g^r,\, h^r \cdot g^m)$ o\`u $h = g^x$ est la cl\'e publique
de l'\'election. Le d\'epouillement requiert la participation d'un seuil de garants ($t$
parmi $k$), la cl\'e de d\'echiffrement \'etant r\'epartie par le partage de secret de Shamir.

Les garanties qui en r\'esultent (confidentialit\'e individuelle, v\'erifiabilit\'e
universelle, r\'esistance partielle \`a la coercition~\cite{Benaloh94}) reposent
sur des fondements math\'ematiques solides~\cite{CramerDamgardSchoenmakers}. Rien
n'assure pour autant qu'elles soient accessibles aux institutions qui les
mettent en \oe uvre~\cite{Bonneau12,Whitten99}.

\paragraph{L'\'echec de Louvain (2009) et l'incident IACR (2025).}
\label{par:echec_louvain}

En mars 2009, l'Universit\'e Catholique de Louvain a d\'eploy\'e Helios pour
l'\'election de son Conseil \'etudiant~\cite{Adida09UCL}. Les garants
\'electoraux, des membres \'elus d\'epourvus de formation cryptographique, n'ont
pas pu conduire le d\'echiffrement distribu\'e sans assistance externe. La
v\'erifiabilit\'e universelle, pourtant garantie par construction, est ainsi
rest\'ee hors de port\'ee des observateurs auxquels elle \'etait destin\'ee.

En novembre 2025, l'IACR a d\^u annuler son \'election annuelle apr\`es qu'un
trustee eut perdu de fa\c{c}on irr\'em\'ediable sa cl\'e partielle~\cite{IACR2025incident}.
La configuration exigeait en effet les parts des trois garants (partage 3-sur-3).
Le scrutin a donc \'et\'e repris avec un seuil 2-sur-3. Les acteurs concern\'es
\textit{\'etaient} cette fois des experts en cryptographie, si bien que l'incident
met en \'evidence, ind\'ependamment de toute question de comp\'etence,
l'\'ecart entre la complexit\'e op\'erationnelle du protocole et les
contraintes concr\`etes d'un d\'eploiement institutionnel.

\paragraph{Profil CLO estim\'e.}
\label{par:helios_clo}

Sous les hypoth\`eses institutionnelles d\'etaill\'ees \`a la
section~\ref{sec:derivation_helios} et \`a la section~\ref{subsubsec:parametres_institutionnels},
les estimations param\'etriques analytiques retenues pour Helios sont
$\Conf \approx 0{,}92$ (confidentialit\'e op\'erationnelle du bulletin chiffr\'e ElGamal~\cite{TsiounisYung}, estim\'ee sur l'\'echelle CLO),
$\Relver \approx 0{,}95$ (compl\'etude et solidit\'e de la v\'erification des preuves de bien-formation),
$\Opp \approx 0{,}50$ (d\'ependance structurelle aux experts pour
l'interpr\'etation, avec capacit\'e d'extraction $\Delta \approx 56$ bits~\cite{MINKA2026CRO}),
ce qui donne le produit
\begin{equation}
\Conf \cdot \Rel \cdot \Opp \;\approx\; 0{,}437.
\label{eq:helios_cro}
\end{equation}
Cette valeur sert \`a la calibration illustrative des param\`etres
institutionnels (sous-section~\ref{subsec:analyse_comparative})~; le syst\`eme
n'est pas une instance du Trilemme CRO (Remarque~\ref{rem:portee_cas_etude}).

\end{sloppypar}

% ----------------------------------------------------------------------------
\subsubsection{Conformit\'e financi\`ere~: les zk-SNARKs}
\label{subsubsec:cas_zksnark}
% ----------------------------------------------------------------------------
\begin{sloppypar}

\paragraph{Succinctit\'e et opacit\'e.}
\label{par:zksnark_architecture}

Les arguments succincts non-interactifs zero-knowledge (zk-SNARKs)~\cite{Groth16,BCTV14,BenSasson2019}
ont nettement renforc\'e la confidentialit\'e des preuves. Une preuve $\pi$
de taille constante ($128$--$256$ octets) y certifie que
$\mathcal{C}(x, w) = 1$ pour un t\'emoin secret $w$, sa v\'erification ne
demandant que quelques millisecondes~\cite{Parno13,Wahby18}. Ces propri\'et\'es expliquent l'int\'er\^et
des zk-SNARKs pour la conformit\'e r\'eglementaire, puisqu'elles permettent de prouver
le respect des r\`egles AML~\cite{FATF12,AMLD4} ou de B\^ale~III~\cite{BaselIII}
sans divulguer les transactions sous-jacentes~\cite{Garratt19}.

Cette succinctit\'e a toutefois un co\^ut~\cite{Bitansky2017}, car une preuve qui ne r\'ev\`ele rien du t\'emoin
interdit aussi \`a un r\'egulateur d'auditer apr\`es coup une transaction suspecte.
Lors d'une crise syst\'emique comparable \`a celle des \textit{subprimes} de 2008~\cite{FCIC2011}, l'analyse de la contagion deviendrait impossible, l'opacit\'e
qui prot\`ege les clients faisant alors obstacle \`a l'investigation~\cite{Androulaki13,Meiklejohn13}. Le probl\`eme rel\`eve d'un vide de
contr\^ole~\cite{FinckMoscon19,DeFilippi16} bien plus que d'une d\'efaillance d'impl\'ementation.

\paragraph{Tensions avec le droit de la protection des donn\'ees.}
\label{par:zksnark_rgpd}

Le RGPD (R\`eglement G\'en\'eral pour la Protection des Donn\'ees \`a caract\`ere personnel)~\cite{GDPR2016} consacre un droit d'acc\`es (Article~15), un droit
\`a l'effacement (Article~17) ainsi qu'une obligation de rendre compte
(Article~5(2))~\cite{CCPA2018}. Chacune de ces exigences se heurte aux
propri\'et\'es zero-knowledge. Le droit d'acc\`es suppose en effet de pouvoir
d\'echiffrer, le droit \`a l'effacement s'oppose \`a l'immuabilit\'e de la
blockchain~\cite{Nakamoto08,Garay15} et l'obligation de justification pr\'esuppose
l'auditabilit\'e. La litt\'erature n'a pas encore r\'esolu ces
incompatibilit\'es~\cite{FinckMoscon19,Sergey18}.

\paragraph{Profil CLO estim\'e.}
\label{par:zksnark_clo}

Sous les m\^emes hypoth\`eses, les estimations param\'etriques analytiques donnent
$\Conf \approx 0{,}89$ (confidentialit\'e op\'erationnelle estim\'ee du dispositif de conformit\'e, au-del\`a de la seule preuve~\cite{Groth16}),
$\Relver \approx 0{,}96$ (compl\'etude et solidit\'e de la v\'erification),
$\Opp \approx 0{,}67$ (opposabilit\'e \'elev\'ee pour les institutions dot\'ees
d'une forte capacit\'e d'extraction $\Delta \approx 73$ bits~\cite{MINKA2026CRO}),
d'o\`u le produit
\begin{equation}
\Conf \cdot \Rel \cdot \Opp \;\approx\; 0{,}572.
\label{eq:zksnark_cro}
\end{equation}
Ce produit CRO relativement \'elev\'e s'explique par la forte capacit\'e
d'extraction des institutions financi\`eres sp\'ecialis\'ees ($\Delta \approx 73$ bits), qui disposent d'infrastructures
d'analyse, d'\'equipes d'experts et de cadres r\'eglementaires d\'etaill\'es~\cite{Ravi2020}.

\end{sloppypar}

% ----------------------------------------------------------------------------
\subsubsection{Signatures contractuelles~: ECDSA comme r\'ef\'erence d'opposabilit\'e}
\label{subsubsec:cas_ecdsa}
% ----------------------------------------------------------------------------
\begin{sloppypar}

\paragraph{Architecture et acceptation universelle.}
\label{par:ecdsa_architecture}

L'\emph{Elliptic Curve Digital Signature Algorithm} (ECDSA)~\cite{Johnson01}
se situe \`a l'autre extr\'emit\'e du spectre. Normalis\'e par le NIST dans FIPS~186-2 (2000) puis FIPS~186-5 (2023)~\cite{NIST.FIPS.186-5},
il est d\'eploy\'e dans SSL/TLS, dans Bitcoin~\cite{Nakamoto08,Narayanan16} ainsi que dans
les syst\`emes bancaires. Sa s\'ecurit\'e repose sur la difficult\'e du probl\`eme du logarithme discret sur
courbes elliptiques (ECDLP)~\cite{RFC6090,RFC6979}, la cl\'e publique $Q = d \cdot G$
\'etant obtenue \`a partir de la cl\'e secr\`ete $d \in [1, n-1]$. La signature d'un message $m$ produit
$(r, s)$ v\'erifiable par quiconque conna\^it $Q$ et $m$.

Cette transparence totale fonde l'opposabilit\'e d'ECDSA~\cite{ESIGN00,eIDAS14}, puisque la
v\'erification ne d\'epend d'aucun expert tiers, peut \^etre rejou\'ee \`a tout moment et
aboutit \`a une attribution univoque assortie d'une tra\c{c}abilit\'e compl\`ete~\cite{Sasson14}.
Elle a pour contrepartie directe l'absence de toute confidentialit\'e, le message
circulant en clair tandis que l'identit\'e est publique et la transaction trac\'ee~\cite{Androulaki13,ECDSA-Fails}.

\paragraph{Profil CLO estim\'e.}
\label{par:ecdsa_clo}

Les estimations param\'etriques analytiques sont
$\Conf \approx 0{,}31$ (confidentialit\'e op\'erationnelle estim\'ee du dispositif contractuel~; au sens statistique de la D\'efinition~\ref{def:confidentialite_clo}, un message publi\'e en clair a une confidentialit\'e nulle),
$\Relver \approx 0{,}98$ (infalsifiabilit\'e existentielle sous ECDLP~\cite{GoldwasserMicali84,GMR88}),
$\Opp \approx 0{,}90$ (v\'erification rejouable, non-r\'epudiation forte,
capacit\'e d'extraction $\Delta \approx 35$ bits suffisante~\cite{MINKA2026CRO}),
soit le produit
\begin{equation}
\Conf \cdot \Rel \cdot \Opp \;\approx\; 0{,}273.
\label{eq:ecdsa_cro}
\end{equation}
La faiblesse de ce produit ne signale aucune d\'efaillance de conception.
Elle d\'ecoule d'un choix architectural d\'elib\'er\'e, ECDSA privil\'egiant
l'opposabilit\'e au d\'etriment de la confidentialit\'e~\cite{BonehFranklin}.

\end{sloppypar}

% ----------------------------------------------------------------------------
\subsubsection{Analyse comparative et lacune th\'eorique}
\label{subsubsec:analyse_comparative}
% ----------------------------------------------------------------------------
\begin{sloppypar}
Le Tableau~\ref{tab:clo_resume_ch1} r\'esume les trois profils pr\'esent\'es plus haut
et appelle deux observations. En premier lieu, confidentialit\'e et opposabilit\'e
varient en sens inverse, les protocoles qui maximisent l'une minimisant
l'autre~\cite{ClarksonSchneider10}. Aucun des trois n'approche la r\'egion $(\Conf, \Rel, \Opp)
\approx (0{,}8, 0{,}8, 0{,}8)$, qui resterait pourtant en dessous de
$1$. Ensuite, les produits $\Conf \cdot \Rel \cdot \Opp$ varient de
$0{,}273$ \`a $0{,}572$, ce qui sugg\`ere l'existence d'une contrainte
qui d\'epend des param\`etres
institutionnels ($\Delta$, $\alpha$, $\beta$)~\cite{Simon82,Cowan2001}.

La cryptographie classique ne dispose d'aucun outil pour formaliser cette
contrainte~\cite{AbadiRogaway00}, car elle mesure l'avantage de l'adversaire~\cite{Goldreich01} sans rien dire de l'utilit\'e institutionnelle. La litt\'erature sur l'\emph{usable
security}~\cite{Bonneau12,Cranor05,Herley09,Dhamija06,Schechter07} \'etudie l'utilisabilit\'e, mais
qualitativement, sans m\'etriques quantifiables. L'informatique
juridique~\cite{Kerr04,Casey11,Mason07} analyse l'admissibilit\'e des preuves,
mais via la jurisprudence, sans formalisme math\'ematique.

Aucune de ces trois disciplines, qui s'ignorent largement, ne propose donc de
m\'etrique de l'opposabilit\'e~\cite{Smith09,Alvim12}. La section~\ref{subsec:trois_dimensions} s'attache \`a combler ce manque.
\end{sloppypar}

\begin{table}[htbp]
\centering
\caption{Profils CLO estimés~: trois systèmes représentatifs.
Toutes les valeurs sont des estimations paramétriques analytiques publiées
dans les travaux CLO~\cite{MINKA2026CLO} et CRO~\cite{MINKA2026CRO}, sur l'\'echelle
op\'erationnelle CLO (fiabilit\'e de v\'erification $\Relver$, Remarque~\ref{rem:rel_ver}).}
\label{tab:clo_resume_ch1}
\begin{tabular}{lcccc}
\toprule
\textbf{Système} & $\Conf$ & $\Relver$ & $\Opp$ & $\Conf \cdot \Relver \cdot \Opp$ \\
\midrule
Helios (vote)        & $0{,}92$ & $0{,}95$ & $0{,}50$ & $0{,}437$ \\
zk-SNARKs (finance) & $0{,}89$ & $0{,}96$ & $0{,}67$ & $0{,}572$ \\
ECDSA (contrats)    & $0{,}31$ & $0{,}98$ & $0{,}90$ & $0{,}273$ \\
\bottomrule
\end{tabular}
\end{table}

% ----------------------------------------------------------------------------
\subsection{Extension quantique~: les standards NIST PQC}
\label{subsec:extension_quantique_contexte}
% ----------------------------------------------------------------------------
\begin{sloppypar}

\subsubsection{La transition quantique}
\label{subsubsec:transition_quantique}

En ao\^ut 2024, le National Institute of Standards and Technology (NIST) a publi\'e
les trois premi\`eres normes de cryptographie post-quantique~\cite{NISTPQC,Alagic2020,NISTPostQuantum2022}, qui ouvrent une nouvelle
p\'eriode pour la discipline~\cite{Chen2022}. Il s'agit des normes suivantes.

\begin{itemize}
\item \textbf{FIPS~203 (ML-KEM)}, m\'ecanisme d'encapsulation de cl\'es fond\'e sur
les r\'eseaux euclidiens et d\'eriv\'e de Kyber~\cite{Kyber2022}, sert \`a l'\'etablissement de cl\'es
ainsi qu'au chiffrement hybride.

\item \textbf{FIPS~204 (ML-DSA)}, algorithme de signature num\'erique fond\'e sur
les r\'eseaux euclidiens et d\'eriv\'e de Dilithium~\cite{Dilithium2022}, est appel\'e \`a remplacer ECDSA et RSA
pour les signatures.

\item \textbf{FIPS~205 (SLH-DSA)}, algorithme de signature fond\'e sur les fonctions
de hachage et d\'eriv\'e de SPHINCS+~\cite{SPHINCS2022}, constitue une alternative sans trusted setup qui repose sur
des hypoth\`eses cryptographiques conservatrices.
\end{itemize}

Le standard FIPS~206 (FN-DSA), d\'eriv\'e de Falcon~\cite{Falcon2022}, en cours de finalisation, compl\'etera cet ensemble.

\subsubsection{Enjeux institutionnels de la migration PQC}
\label{subsubsec:enjeux_pqc}

La migration vers la cryptographie post-quantique fait na\^itre de nouvelles
tensions avec les exigences institutionnelles~\cite{Langlois2015,Boneh2011}, qui tiennent \`a la complexit\'e des preuves, \`a leur taille et \`a la nouveaut\'e des hypoth\`eses de s\'ecurit\'e.

\paragraph{Complexit\'e accrue des preuves.}
Les signatures Dilithium reposent sur des preuves de connaissance dans les r\'eseaux
euclidiens, dont la v\'erification fait intervenir des op\'erations matricielles et des
calculs de normes~\cite{Dilithium2022}. Un arbitre habitu\'e \`a ECDSA devra donc se familiariser avec les
modules sur $\mathbb{Z}_q$ ainsi qu'avec les probl\`emes LWE (Learning With Errors)~\cite{Regev2009}.

\paragraph{Taille des preuves.}
Une signature ML-DSA-65 (Dilithium-3) occupe 3\,309 octets et sa cl\'e publique 1\,952 octets, contre 64 octets pour une signature ECDSA-256.
De m\^eme, une cl\'e publique Kyber-1024 atteint 1\,568 octets, l\`a o\`u ECDSA se contente de 64 octets.
Cet accroissement complique l'examen manuel des preuves comme leur archivage~\cite{Ravi2020}.

\paragraph{Hypoth\`eses de s\'ecurit\'e non standard.}
LWE et les probl\`emes de r\'eseaux euclidiens constituent des hypoth\`eses r\'ecentes
(ann\'ees 2000), qui ne b\'en\'eficient pas de la confiance accumul\'ee au fil du temps par RSA (1977) ou
ECDLP (1985)~\cite{Kiltz2017}. Des institutions prudentes pourraient donc h\'esiter \`a reconna\^itre la m\^eme
force probante \`a des preuves fond\'ees sur ces hypoth\`eses nouvelles~\cite{Unruh2016}.

\subsubsection{Implications pour le Trilemme CRO}
\label{subsubsec:implications_cro}

La transition post-quantique conduit \`a se demander si les bornes du Trilemme CRO
\'etablies dans le contexte classique restent valides dans les trois situations
suivantes~\cite{NielsenChuang2010,Wilde2017}.

\begin{enumerate}
\item L'adversaire peut effectuer des requ\^etes quantiques en superposition~\cite{Watrous2009,Boneh2011}.
\item Les canaux de communication introduisent des d\'egradations quantiques~\cite{Holevo1973,Schumacher1997}.
\item Les param\`etres institutionnels ($\Delta$, $\alpha$, $\beta$) doivent
\^etre adapt\'es au contexte quantique~\cite{Mayers1997,Lo1997}.
\end{enumerate}

C'est pour y r\'epondre que la section~\ref{subsec:modele_formel_pq} d\'eveloppe
l'extension quantique du Trilemme CRO, en introduisant les param\`etres
$\Delta_q$ (capacit\'e d'extraction quantique) et $\eta_q$ (perte d'interpr\'etabilit\'e
quantique).

\end{sloppypar}

% ----------------------------------------------------------------------------
%\subsection{Synth\`ese}
%\label{sec:synthese_chap1}
% ----------------------------------------------------------------------------

\begin{boiteArbitre}
L'arbitre conna\^it ces dossiers pour les avoir tenus entre ses mains.
Helios, les zk-SNARKs et ECDSA reposent sur des architectures distinctes,
mais chacun a \'echou\'e \`a sa mani\`ere \`a le convaincre.

Ce qui l'a longtemps troubl\'e tient \`a ce que la d\'efaillance n'\'etait jamais
d'ordre technique. Les preuves \'etaient correctes et les algorithmes
irr\'eprochables, sans que cela lui permette pour autant de d\'ecider. Il lui
manquait moins la comp\'etence qu'un langage commun entre la rigueur
math\'ematique du protocole et les exigences proc\'edurales de sa fonction. Ce
chapitre a donn\'e \`a cette situation le nom de paradoxe de l'authenticit\'e
invisible.

Il en retient quatre enseignements. Une confidentialit\'e tr\`es \'elev\'ee
(sup\'erieure \`a $0{,}90$) produit une opacit\'e qui paralyse l'interpr\`ete
institutionnel, comme le montrent Helios et les zk-SNARKs. Les
produits $\Conf \cdot \Rel \cdot \Opp$ s'\'etagent de $0{,}273$ \`a $0{,}572$
selon les param\`etres institutionnels ($\Delta$, $\alpha$, $\beta$), ce qui
d\'esigne une contrainte qui d\'epend du contexte.
Les protocoles actuels ont aussi \'et\'e con\c{c}us sans tenir compte de
lui, puisque sa capacit\'e cognitive born\'ee ($\Delta$) n'appara\^it dans aucune
sp\'ecification. Enfin, la transition post-quantique risque d'aggraver cette
situation en y ajoutant une complexit\'e que les experts eux-m\^emes peinent \`a
ma\^itriser.

Il reste \`a passer du constat au th\'eor\`eme, ce que le chapitre suivant
pr\'epare en posant l'infrastructure th\'eorique n\'ecessaire.
\end{boiteArbitre}
% ============================================================================
% CHAPITRE 1 — CONTEXTE, \'ETAT DE L'ART ET PROBL\'EMATIQUE
% Section 1.3 : \'Etat de l'art
% source : partie1/etat_art.tex
% ============================================================================

\section{\'ETAT DE L'ART}
\label{sec:etat_art}
\begin{sloppypar}
La cryptographie institutionnelle ne peut \^etre formalis\'ee qu'en r\'eunissant six corpus th\'eoriques d'ordinaire s\'epar\'es, \`a savoir la th\'eorie de l'information classique~\cite{Shannon48,Shannon49,CoverThomas06}, la th\'eorie de l'information
quantique~\cite{NielsenChuang2010,Wilde2017}, les primitives cryptographiques d'obscurcissement~\cite{Pedersen91,Groth16}, les
d\'efinitions pr\'eliminaires du cadre CRO, les trilemmes d'impossibilit\'e~\cite{Arrow50,Brewer00,FLP85} ainsi que la v\'erification formelle~\cite{moura2015lean,Lean4,Coq,Isabelle}. Nous les parcourons ici dans le but de \emph{pr\'eparer des blocs directement utilisables}. Chaque section ne retient donc que les r\'esultats dont les sections~\ref{sec:clo}, \ref{sec:cro_classique},
\ref{sec:cro_pq}, \ref{sec:sotp_construction} et~\ref{sec:lean4}
auront besoin, en indiquant le r\^ole exact qu'ils jouent dans la d\'emonstration.

L'ordre suivi reprend la logique annonc\'ee au paragraphe~\ref{par:progression_FBESCX}, qui structure l'ensemble de la th\`ese. Les
sous-sections~\ref{subsec:fondements_info}, \ref{subsec:fondements_info_quantique}
et~\ref{subsec:primitives_crypto} rassemblent les outils n\'ecessaires pour
d\'efinir les m\'etriques (\textbf{F}). La sous-section~\ref{subsec:trilemmes}
examine ensuite les paradigmes d'impossibilit\'e conditionnelle dont s'inspireront
l'\'etablissement des bornes (\textbf{B}) puis leur extension
post-quantique (\textbf{E}). La sous-section~\ref{subsec:definitions_preliminaires_cro} fixe le
vocabulaire formel dans lequel seront \'enonc\'ees les constructions con\c{c}ues
pour approcher les bornes (\textbf{S}), avant que la sous-section~\ref{subsec:verification_formelle}
ne pr\'esente l'infrastructure de certification (\textbf{C}). Les fondements
cognitifs de la sous-section~\ref{subsec:fondements_cognitifs} motivent enfin
l'hypoth\`ese de travail sur la capacit\'e $\Delta$, que le programme
exp\'erimental devra mesurer (\textbf{X}).

Deux pr\'ecisions de m\'ethode doivent \^etre apport\'ees d\`es \`a pr\'esent. En premier lieu, les primitives de la
cryptographie institutionnelle sont celles de la cryptographie standard~\cite{Goldreich01,DiffieHellman76}.
Ce qui la distingue des sous-disciplines existantes tient \`a son \emph{objectif}, qui consiste \`a quantifier et \`a optimiser l'utilisabilit\'e
institutionnelle des garanties cryptographiques, en traitant notamment
l'\emph{opposabilit\'e} comme une propri\'et\'e de premi\`ere classe. En second lieu,
les trois m\'etriques ($\Conf$~\ref{def:confidentialite_clo}, $\Rel$~\ref{def:fiabilite_clo}, $\Opp$~\ref{def:opposabilite_clo}) seront d\'efinies ici dans
leur formulation issue de la th\'eorie de l'information, qui autorise le raisonnement
asymptotique. Leurs \'equivalents computationnels, \'etablis au moyen de jeux de
s\'ecurit\'e~\cite{GoldwasserMicali84,GMR88}, sont trait\'es \`a la section~\ref{sec:cro_classique}.
\end{sloppypar}

\begin{boiteArbitre}
L'arbitre ne dispose pas encore d'un langage pour décrire ce qu'il vit. Ce chapitre
lui en fournit les éléments, avec l'entropie pour mesurer ce qu'une preuve révèle,
l'inégalité de Fano pour borner ce qu'un interprète peut
récupérer d'un signal bruité, ou encore les trilemmes CAP et FLP qui lui montrent que son incapacité s'apparente à
d'autres impossibilités déjà connues. Il ignore encore que son cas admet une borne quantitative,
ce que le chapitre suivant lui apprendra.
\end{boiteArbitre}

% ----------------------------------------------------------------------------
\subsection{Fondements en th\'eorie de l'information classique}
\label{subsec:fondements_info}
% ----------------------------------------------------------------------------
\begin{sloppypar}
Nous avons retenu l'entropie, l'information
mutuelle et les in\'egalit\'es classiques de Shannon~\cite{Shannon48,Shannon49} afin
de quantifier confidentialit\'e et opposabilit\'e dans un m\^eme formalisme.
Les objets d\'efinis ci-apr\`es interviennent dans chacune des preuves des chapitres~\ref{chap:cadre_formel} et~\ref{chap:constructions}.

\subsubsection{Entropie et information mutuelle}
\label{subsubsec:entropie_info_mutuelle}

L'entropie de Shannon d'une variable al\'eatoire discr\`ete $X$ sur alphabet
$\mathcal{X}$ est d\'efinie par~\cite{Shannon48,CoverThomas06,Fano61}
\begin{equation}
H(X) = -\sum_{x \in \mathcal{X}} P(x)\log_2 P(x).
\label{eq:entropie_shannon}
\end{equation}
Elle mesure l'incertitude r\'esiduelle sur $X$ avant toute observation.
L'information mutuelle conditionnelle
\begin{equation}
I(X;Y|Z) = H(X|Z) - H(X|Y,Z)
\label{eq:mi_conditionnelle}
\end{equation}
d\'ecrit la r\'eduction d'incertitude sur $X$ apport\'ee par $Y$, une fois $Z$
connu. L'opposabilit\'e $\Opp$ sera d\'efinie \`a partir de cette quantit\'e, plus adapt\'ee que l'entropie brute~\cite{Wyner75,WynerZiv76}.

\subsubsection{Distances et divergences}
\label{subsubsec:distances_divergences}

La distance de variation totale entre deux distributions $P$ et $Q$, donn\'ee par
\begin{equation}
\TV(P,Q) = \tfrac{1}{2}\sum_x |P(x) - Q(x)|
\label{eq:tv_distance}
\end{equation}
quantifie la distinguabilit\'e au sens du pire adversaire d\'eterministe~\cite{CoverThomas06}. La
divergence de Kullback-Leibler, d\'efinie par
\begin{equation}
\KL(P \| Q) = \sum_x P(x)\log_2\frac{P(x)}{Q(x)}
\label{eq:kl_divergence}
\end{equation}
majore $\TV$ via l'in\'egalit\'e de Pinsker~\cite{Pinsker1964}, que nous utiliserons pour borner
la confidentialit\'e.

\subsubsection{L'in\'egalit\'e de Pinsker}
\label{subsubsec:pinsker_classique}

\begin{lemma}[In\'egalit\'e de Pinsker~\cite{Pinsker1964,CoverThomas06}]\label{lem:pinsker}
Pour toutes distributions $P$ et $Q$ sur un m\^eme alphabet fini, on a
\begin{equation}
\TV(P,Q)^2 \leq c_0 \cdot \KL(P \| Q),
\qquad c_0 = \frac{1}{2\ln 2} \approx 0{,}721.
\label{eq:pinsker}
\end{equation}
\end{lemma}

\noindent Dans la preuve du Trilemme CRO \`a la section~\ref{sec:cro_classique}, cette in\'egalit\'e relie la
confidentialit\'e $\Conf = 1 - \max_{m_0,m_1} \TV(P_{\Sigma|m_0},
P_{\Sigma|m_1})$ \`a l'information mutuelle $I(M;\Sigma|C)$ via
l'hypoth\`ese de $\beta$-r\'egularit\'e. Elle permet ainsi de convertir un
\'ecart entre distributions en une borne sur le produit $\Conf \cdot \Rel \cdot
\Opp$~\cite{Tsybakov2009}.

\subsubsection{Le lemme de Fano}
\label{subsubsec:fano_classique}

\begin{lemma}[Lemme de Fano~\cite{Fano61,CoverThomas06}]\label{lem:fano}
Soit $X$ une variable al\'eatoire sur un alphabet fini de cardinalit\'e $k$,
$\hat{X}$ un estimateur et $p_e = \Pr[X \neq \hat{X}]$. On a alors
\begin{equation}
H(X|Y) \leq h(p_e) + p_e \log_2(k-1),
\label{eq:fano_inequality}
\end{equation}
o\`u $h(p) = -p\log_2 p - (1-p)\log_2(1-p)$ est la fonction d'entropie
binaire.
\end{lemma}

Une fiabilit\'e \'elev\'ee (faible $p_e$) implique donc une information mutuelle \'elev\'ee.
La contrapos\'ee est tout aussi utile, puisque si $H(X|Y)$ reste grande, aucun
estimateur ne peut retrouver $X$ \`a partir de $Y$ avec pr\'ecision~\cite{HanVerdu}. Dans notre
cadre, il en d\'ecoule directement pour le Trilemme CRO l'in\'egalit\'e
\begin{equation}
1 - \Rel \;\geq\; \frac{H(M|C) - I(M;\Sigma|C) - 1}{\log_2 |\mathcal{M}|}.
\label{eq:rel_fano_bound}
\end{equation}

\subsubsection{L'in\'egalit\'e de traitement des donn\'ees}
\label{subsubsec:dpi}

\begin{lemma}[Data Processing Inequality~\cite{CoverThomas06,Lindblad}]\label{lem:dpi}
Pour toute cha\^ine de Markov $X \to Y \to Z$, on a
\begin{equation}
I(X;Z) \;\leq\; I(X;Y).
\label{eq:dpi}
\end{equation}
\end{lemma}

\noindent Ce r\'esultat, rarement mobilis\'e en
cryptographie, intervient directement dans notre analyse~\cite{Smith09,Alvim12}. La
transformation $M \to (\Sigma, C) \to V = \mathsf{Interpret}(\Sigma, C, K)$ forme
une cha\^ine de Markov d\`es que les ressources propres $K$ de l'institution sont
ind\'ependantes du message sachant la preuve et le contexte, que le verdict soit
d\'eterministe ou non. L'in\'egalit\'e \eqref{eq:dpi} garantit alors que
$I(M;V|C) \leq I(M;\Sigma|C)$, autrement dit ce que le verdict apprend du \emph{message} ne peut exc\'eder
ce que la preuve en contient. Cette in\'egalit\'e fournit l'un des deux termes de la
d\'ecomposition qui fonde, \`a la section~\ref{sec:cro_classique}, l'Hypoth\`ese
d'extraction born\'ee (Proposition~\ref{prop:fondement_extraction})~\cite{ClarksonSchneider10}.

\subsubsection{Information flow quantitatif}
\label{subsubsec:qif}

La litt\'erature sur le \emph{Quantitative Information Flow} (QIF)~\cite{Clarkson05,Smith09,Alvim12,Alvim18}
fournit un cadre formel pour mesurer les fuites d'information dans les syst\`emes
de s\'ecurit\'e. La mesure $g$-leakage~\cite{Alvim12} \'etend l'information mutuelle
en pond\'erant chaque fuite par son effet sur les gains de l'adversaire. Notre d\'efinition
de l'opposabilit\'e, con\c{c}ue comme un ratio de r\'eduction d'incertitude, s'en inspire directement.

\subsubsection{Capacit\'e de canal et codage source-canal}
\label{subsubsec:capacite_canal}

Le th\'eor\`eme de codage source-canal~\cite{Shannon48,CoverThomas06} \'etablit que
la transmission fiable d'une source $X$ sur un canal de capacit\'e $C$ est possible
si et seulement si $H(X) < C$. Transpos\'e au contexte institutionnel, ce r\'esultat conduit \`a
mod\'eliser l'interpr\`ete comme un canal de capacit\'e born\'ee $\Delta$, qui limite
la quantit\'e d'information qu'il peut extraire~\cite{WynerZiv76,Maurer1999}.
\end{sloppypar}

% ----------------------------------------------------------------------------
\subsection{Fondements en th\'eorie de l'information quantique}
\label{subsec:fondements_info_quantique}
% ----------------------------------------------------------------------------
\begin{sloppypar}
L'extension du Trilemme CRO
aux contextes quantiques n\'ecessite des outils de th\'eorie de l'information
quantique qui g\'en\'eralisent les notions classiques~\cite{NielsenChuang2010,Wilde2017,Holevo1973}. Cette section introduit
les concepts fondamentaux mobilis\'es \`a la section~\ref{sec:cro_pq}.

\subsubsection{\'Etats quantiques et entropie de von Neumann}
\label{subsubsec:von_neumann}

Un \'etat quantique est repr\'esent\'e par une matrice de densit\'e $\rho$, c'est-\`a-dire un
op\'erateur positif semi-d\'efini de trace unit\'e sur un espace de Hilbert $\mathcal{H}$~\cite{NielsenChuang2010}.
L'entropie de von Neumann, qui g\'en\'eralise l'entropie de Shannon, est d\'efinie par
\begin{equation}
S(\rho) = -\mathrm{Tr}(\rho \log \rho).
\label{eq:von_neumann}
\end{equation}
Pour un \'etat pur $|\psi\rangle$, l'entropie est nulle ($S(|\psi\rangle\langle\psi|) = 0$). Elle
vaut $S(\rho) = \log_2 d$ sur l'\'etat compl\`etement m\'elang\'e $\rho = I/d$. Lorsque
$\rho$ est diagonale dans la base computationnelle, on retrouve $S(\rho) = H(X)$, o\`u $X$
suit la distribution form\'ee par les probabilit\'es diagonales. La d\'efinition quantique
contient donc le cadre classique comme cas particulier~\cite{Wilde2017,Hiai1991}.

\subsubsection{Information mutuelle quantique}
\label{subsubsec:info_mutuelle_quantique}

Pour un \'etat bipartite $\rho_{AB}$, l'information mutuelle quantique s'\'ecrit
\begin{equation}
I(A;B)_\rho = S(\rho_A) + S(\rho_B) - S(\rho_{AB}).
\label{eq:qmi}
\end{equation}
L'information mutuelle quantique ob\'eit elle aussi \`a une in\'egalit\'e de traitement
des donn\'ees~\cite{Lindblad}, puisque l'application d'un canal CPTP $\Lambda$ au
syst\`eme $B$ ne peut que la r\'eduire,
\begin{equation}
I(A;B)_\rho \geq I(A; \Lambda(B))_\rho.
\label{eq:qdpi}
\end{equation}

\subsubsection{Information de Holevo}
\label{subsubsec:holevo}

L'information de Holevo borne l'information classique accessible depuis
un encodage quantique~\cite{Holevo1973,Schumacher1997}. Pour un ensemble $\{P_X(x), \rho_x\}$ d'\'etat moyen
$\bar{\rho} = \sum_x P_X(x)\rho_x$, elle vaut
\begin{equation}
\chi = S(\bar{\rho}) - \sum_x P_X(x) S(\rho_x).
\label{eq:holevo}
\end{equation}
Le th\'eor\`eme de Holevo \'etablit que l'information accessible $I_{\mathrm{acc}}$
par toute mesure POVM satisfait $I_{\mathrm{acc}} \leq \chi$. L'extension
post-quantique de l'opposabilit\'e s'appuie directement sur cette borne. Quels que
soient les outils quantiques dont il dispose, un interpr\`ete institutionnel n'extrait
en effet du \textit{proof state} quantique qu'au plus $\chi$ bits d'information~\cite{Winter}.

\subsubsection{Distance de trace et fid\'elit\'e}
\label{subsubsec:distance_trace_fidelite}

La distance de trace quantique, qui g\'en\'eralise la variation totale~\cite{NielsenChuang2010}, est d\'efinie par
\begin{equation}
\|\rho - \sigma\|_1 = \mathrm{Tr}\sqrt{(\rho-\sigma)^\dagger(\rho-\sigma)}.
\label{eq:trace_distance}
\end{equation}
Pour des \'etats classiques (diagonaux), $\|\rho - \sigma\|_1 = 2\TV(P,Q)$.
La fid\'elit\'e quantique~\cite{Jozsa}
\begin{equation}
F(\rho, \sigma) = \|\sqrt{\rho}\sqrt{\sigma}\|_1
\label{eq:fidelite}
\end{equation}
est reli\'ee \`a la distance de trace par les in\'egalit\'es de Fuchs-van de Graaf~\cite{FuchsGraaf}
\begin{equation}
1 - \tfrac{1}{2}\|\rho - \sigma\|_1 \leq F(\rho, \sigma) \leq \sqrt{1 - \tfrac{1}{4}\|\rho - \sigma\|_1^2}.
\label{eq:fuchs_vdg}
\end{equation}

\subsubsection{In\'egalit\'e de Pinsker quantique}
\label{subsubsec:pinsker_quantique}

\begin{lemma}[In\'egalit\'e de Pinsker quantique~\cite{Wilde2017,Hiai1991}]\label{lem:pinsker_quantique_etat_art}
Pour tous \'etats quantiques $\rho$ et $\sigma$, on a
\begin{equation}
\|\rho - \sigma\|_1^2 \leq 2 D(\rho\|\sigma),
\label{eq:qpinsker}
\end{equation}
o\`u $D(\rho\|\sigma) = \mathrm{Tr}(\rho\log\rho - \rho\log\sigma)$ est l'entropie
relative quantique.
\end{lemma}

Cette in\'egalit\'e sera utilis\'ee dans la section~\ref{sec:cro_pq} (Lemme~\ref{lem:pinsker_quantique}) pour borner la
confidentialit\'e quantique $\Conf^q$ en fonction de l'information mutuelle
quantique $I(M;\Sigma|C)$~\cite{Ohya1993}.

\subsubsection{Lemme de gentle measurement}
\label{subsubsec:gentle_measurement}

\begin{lemma}[Gentle measurement~\cite{Wilde2017,Winter}]\label{lem:mesure_douce_etat_art}
Soit $\rho$ un \'etat quantique et $\{M_m\}$ une POVM. Si l'issue $m$ se produit
avec probabilit\'e $p_m = \mathrm{Tr}(M_m\rho) \geq 1 - \delta$, alors l'\'etat
post-mesure satisfait
\begin{equation}
\|\rho - \rho'\|_1 \leq 2\sqrt{\delta}.
\label{eq:gentle}
\end{equation}
\end{lemma}

Ce lemme capture l'id\'ee qu'une mesure \guillemotleft~douce~\guillemotright (haute probabilit\'e) perturbe
peu l'\'etat quantique~\cite{OgawaNagaoka}. Dans le contexte institutionnel post-quantique, il borne
la perte d'interpr\'etabilit\'e $\eta_q$ induite par les effets de canal quantique.

\subsubsection{In\'egalit\'e de Fano quantique}
\label{subsubsec:fano_quantique}

\begin{lemma}[Lemme de Fano quantique~\cite{Wilde2017,Winter}]\label{lem:fano_quantique_etat_art}
Soient $X$ encod\'e dans l'\'etat quantique $\rho_X$ et $\hat{X}$ le r\'esultat d'une
mesure POVM avec probabilit\'e d'erreur $p_e$. Alors
\begin{equation}
S(X|\hat{X}) \leq h(p_e) + p_e \log_2(|\mathcal{X}| - 1).
\label{eq:qfano}
\end{equation}
\end{lemma}

Cette g\'en\'eralisation du lemme de Fano classique est utilis\'ee dans la section~\ref{sec:cro_pq} (Lemme~\ref{lem:fano_quantique}) pour borner la fiabilit\'e quantique $\Rel^q$~\cite{Hayashi}.

\subsubsection{Canaux quantiques (CPTP)}
\label{subsubsec:canaux_quantiques}

Un canal quantique est une application compl\`etement positive et pr\'eservant
la trace (CPTP) $\Lambda~: \mathcal{L}(\mathcal{H}_A) \to \mathcal{L}(\mathcal{H}_B)$~\cite{NielsenChuang2010}.
Par le th\'eor\`eme de dilatation de Stinespring~\cite{Stinespring}, tout canal admet une repr\'esentation
unitaire avec environnement auxiliaire de la forme
\begin{equation}
\Lambda(\rho) = \mathrm{Tr}_E\left[U (\rho \otimes |0\rangle\langle 0|_E) U^\dagger\right].
\label{eq:stinespring}
\end{equation}

Trois mod\`eles de canaux sont pertinents pour l'analyse post-quantique~\cite{Wilde2017}~:
\begin{itemize}
\item \textbf{Canal d\'epolarisant}, d\'efini par $\Lambda_{\mathrm{dep}}(\rho) = (1-p)\rho + \frac{p}{d}I_d$
\item \textbf{Canal d'amortissement d'amplitude} $\Lambda_{\mathrm{AD}}(\rho)$, qui mod\'elise la dissipation d'\'energie
\item \textbf{Canal de d\'ephasage}, d\'efini par $\Lambda_{\mathrm{deph}}(\rho) = (1-\lambda)\rho + \lambda Z\rho Z$
\end{itemize}

Ces canaux mod\'elisent les effets de bruit quantique qui d\'egradent
l'interpr\'etabilit\'e des preuves dans les syst\`emes hybrides ou quantiques~\cite{Zurek2003,Wootters1982}.

\subsubsection{Impossibilit\'es quantiques fondamentales}
\label{subsubsec:impossibilites_quantiques}

Deux r\'esultats d'impossibilit\'e quantique \'eclairent les contraintes du Trilemme PQ-CRO.

\textbf{No-cloning theorem.} Wootters et Zurek~\cite{Wootters1982} ont prouv\'e qu'un \'etat quantique inconnu ne peut \^etre copi\'e parfaitement. Cette contrainte interdit la duplication arbitraire des preuves quantiques.

\textbf{Impossibilit\'e du bit commitment quantique.} Mayers~\cite{Mayers1997} et Lo-Chau~\cite{Lo1997} ont montr\'e que, dans le mod\`ele quantique non-relativiste, aucun protocole de bit commitment ne garantit simultan\'ement le secret et la liaison. C'est pourquoi les engagements post-quantiques reposent sur des hypoth\`eses calculatoires~\cite{Unruh2016}.
\end{sloppypar}

\begin{boiteArbitre}
L'arbitre retient l'essentiel. Même équipé d'outils quantiques optimaux,
il ne tirera jamais d'un état de preuve quantique plus de $\chi$ bits.
Cette borne découle des lois de la physique, si bien qu'aucun progrès
des instruments ne la repoussera.
\end{boiteArbitre}

% ----------------------------------------------------------------------------
\subsection{Primitives cryptographiques d'obscurcissement}
\label{subsec:primitives_crypto}
% ----------------------------------------------------------------------------
\begin{sloppypar}
Les constructions SOTP du
chapitre~\ref{chap:constructions} font appel \`a trois primitives. Nous les pr\'esentons
bri\`evement en nous limitant aux propri\'et\'es que ces constructions exploitent,
sans viser l'exhaustivit\'e des sch\'emas.

\subsubsection{Engagements de Pedersen}
\label{subsubsec:pedersen}

Un engagement de Pedersen~\cite{Pedersen91} sur un groupe cyclique $G$ de
g\'en\'erateurs $g, h$ est le couple $\mathsf{Com}(m; r) = g^m h^r$, o\`u $m$
est le message secret et $r$ un al\'ea. Le sch\'ema est parfaitement occultant au sens
de la th\'eorie de l'information, puisque pour tout $m' \neq m$ il existe $r'$ tel que
$\mathsf{Com}(m'; r') = \mathsf{Com}(m; r)$. Sa propri\'et\'e de liaison n'est en revanche
que calculatoire, car produire $m' \neq m$ de m\^eme image revient \`a r\'esoudre le probl\`eme
du logarithme discret~\cite{Schnorr91}. La famille SOTP confie \`a ces engagements un double r\^ole.
Ils encodent le contexte institutionnel $C$ de mani\`ere v\'erifiable, au b\'en\'efice de $\Rel$,
tout en pr\'eservant la confidentialit\'e des donn\'ees sous-jacentes, au b\'en\'efice
de $\Conf$~\cite{ChaumPedersen92}.

\subsubsection{Arguments succincts zero-knowledge (zk-SNARKs)}
\label{subsubsec:zksnarks}

Pour un circuit $\mathcal{C}$, une entr\'ee publique $x$ et un t\'emoin secret $w$,
un zk-SNARK~\cite{Groth16,BCTV14,BenSasson2019} atteste que $\mathcal{C}(x, w)
= 1$ au moyen d'une preuve $\pi$ de taille constante ($\approx 128$--$256$ octets).
Le sch\'ema se d\'efinit par sa compl\'etude, qui garantit qu'un prouveur
honn\^ete convainc toujours le v\'erificateur, par sa solidit\'e (soundness) sous les
hypoth\`eses q-SDH et q-PKE~\cite{Gennaro13,Parno13} ainsi que par la propri\'et\'e de
zero-knowledge parfait, en vertu de laquelle le v\'erificateur
n'apprend rien sur $w$ au-del\`a de la validit\'e de l'assertion~\cite{GMR85}.
SOTP-High recourt \`a ces preuves pour maximiser $\Opp$ dans le r\'egime \`a haute
entropie contextuelle ($H(\mathcal{C}) > \lambda/\log\lambda$)~\cite{PLONK20,Wahby18,Campanelli2019}.

\subsubsection{Syst\`emes de preuves avanc\'es}
\label{subsubsec:systemes_preuves}

Plusieurs syst\`emes de preuves compl\`etent les zk-SNARKs dans la palette disponible.

\textbf{Bulletproofs.} B\"unz et al.~\cite{Bunz19} ont introduit des preuves \`a divulgation nulle sans trusted setup, avec des tailles logarithmiques en la taille du t\'emoin. Leur construction repose sur un argument de produit scalaire (\emph{inner-product argument}) et sur la difficult\'e du logarithme discret~; elle n'est donc pas post-quantique.

\textbf{STARKs.} Ben-Sasson et al.~\cite{BenSasson18} ont d\'evelopp\'e des preuves transparentes (sans trusted setup) et post-quantiques, bas\'ees sur des codes correcteurs d'erreurs et la technique FRI (Fast Reed-Solomon Interactive Oracle Proof of Proximity).

\textbf{Fiat-Shamir.} La transformation de Fiat et Shamir~\cite{FiatShamir87} convertit les preuves interactives en preuves non-interactives via l'application de hachage. Sa s\'ecurit\'e dans le mod\`ele de l'oracle al\'eatoire quantique a \'et\'e analys\'ee, pour les signatures, par Kiltz, Lyubashevsky et Schaffner~\cite{Kiltz2017}.

\subsubsection{La transformation de Johnson-Lindenstrauss}
\label{subsubsec:jl}

La transformation de Johnson-Lindenstrauss~\cite{JohnsonLindenstrauss84,DasguptaGupta03}
garantit qu'une famille de $n$ vecteurs de $\mathbb{R}^d$
peut \^etre projet\'ee dans $\mathbb{R}^k$ avec $k = O(\log n / \epsilon^2)$
en pr\'eservant toutes les distances euclidiennes \`a un facteur $1 \pm \epsilon$
pr\`es~\cite{Vempala05,Indyk98,Indyk06}. Dans SOTP-Base (r\'egime de transition), cette transformation compresse
l'espace contextuel, r\'eduisant la complexit\'e des engagements sans
sacrifier les distances s\'emantiques~\cite{DasguptaKumar14,Larsen17,Gionis99}.

\subsubsection{Fonctions de hachage et filtres de Bloom}
\label{subsubsec:hachage}

Les fonctions de hachage cryptographique~\cite{Damgard87} servent aux engagements et \`a la transformation Fiat-Shamir. Les filtres de Bloom~\cite{Bloom70} fournissent des structures de donn\'ees probabilistes pour la v\'erification d'appartenance avec contr\^ole du taux de faux positifs, pertinents pour les contextes institutionnels \`a grande cardinalit\'e.

\subsubsection{Chiffrement homomorphe}
\label{subsubsec:fhe}

Le chiffrement compl\`etement homomorphe (FHE), introduit par Gentry~\cite{Gentry09} et am\'elior\'e par Brakerski et Vaikuntanathan~\cite{BrakerskiVaikuntanathan11}, permet le calcul sur donn\'ees chiffr\'ees. Les sch\'emas bas\'es sur LWE~\cite{Regev2009} sont consid\'er\'es comme post-quantiques et utilis\'es dans les protocoles de calcul multipartite s\'ecuris\'e~\cite{Yao82,GMW87,Asharov12}.
\end{sloppypar}

% ----------------------------------------------------------------------------
\subsection{Fondements cognitifs~: rationalit\'e born\'ee et param\`etre~$\Delta$}
\label{subsec:fondements_cognitifs}
% ----------------------------------------------------------------------------
\begin{sloppypar}
Le choix de la capacit\'e
d'extraction institutionnelle $\Delta$ et de l'hypoth\`ese de r\'egularit\'e $\beta$
s'appuie sur des r\'esultats \'etablis en sciences cognitives et en th\'eorie
de l'information.

\subsubsection{Rationalit\'e born\'ee et capacit\'e de traitement}
\label{subsubsec:rationalite_bornee}

Selon Herbert Simon~\cite{Simon82,Simon1955}, les d\'ecideurs humains, arbitres compris,
ne recherchent pas l'optimum. Soumis \`a de s\'ev\`eres contraintes cognitives, ils \emph{satisficent},
c'est-\`a-dire qu'ils retiennent une solution jug\'ee suffisamment bonne en raisonnant sur des
mod\`eles mentaux simplifi\'es \`a $O(10)$--$O(100)$ \'el\'ements distinguables. George Miller~\cite{Miller56}
a estim\'e la capacit\'e de la m\'emoire de travail \`a $7 \pm 2$
\emph{chunks} d'information, chiffre que Cowan~\cite{Cowan2001} ram\`ene \`a environ
$4$ \'el\'ements. Pour une t\^ache d'interpr\'etation de preuve cryptographique, un ordre
de grandeur s'obtient en retenant $\approx 3$ bits par \'el\'ement et la borne haute de Miller,
ce qui donne

\begin{equation}
\Delta_{\text{individu}} \;\lesssim\; 7 \times 3 \;=\; 21 \;\text{bits}.
\label{eq:delta_individu}
\end{equation}
Ericsson et Kintsch~\cite{Ericsson95} ont montr\'e que les experts acc\`edent, via
la m\'emoire \`a long terme, \`a $10^3$--$10^4$ sch\'emas, ajoutant
$\log_2(10^4) \approx 13$ bits d'indexation. La d\'elib\'eration institutionnelle
collective amplifie encore cette capacit\'e~\cite{Argyris1978}. Faute de quantification
dans la litt\'erature, nous retenons pour l'institution l'ordre de grandeur
\begin{equation}
\Delta_{\text{institution}} \;\sim\; \text{quelques dizaines de bits}.
\label{eq:delta_institution}
\end{equation}
Ces estimations cumulent des grandeurs h\'et\'erog\`enes, \`a savoir la capacit\'e
de la m\'emoire de travail et la taille d'un r\'epertoire de sch\'emas. Elles ne valent
donc pas d\'erivation et servent seulement \`a motiver l'\emph{hypoth\`ese de travail}
d'une capacit\'e de l'ordre de $20$ \`a $80$ bits, qu'il reviendra au programme
exp\'erimental du chapitre~\ref{chap:experimentation} de mesurer.

\subsubsection{Biais cognitifs et d\'ecision}
\label{subsubsec:biais_cognitifs}

Kahneman~\cite{Kahneman11} a document\'e les biais syst\'ematiques affectant
la d\'ecision humaine, parmi lesquels l'aversion \`a la perte, l'ancrage et la disponibilit\'e. Ces biais
modulent la capacit\'e d'extraction effective, si bien qu'un juge peut disposer de
$\Delta$ bits de capacit\'e mais n'en utiliser qu'une fraction selon le contexte
\'emotionnel ou proc\'edural.

\subsubsection{Extension quantique~: $\Delta_q$ et $\eta_q$}
\label{subsubsec:extension_quantique_delta}

Dans le contexte post-quantique, la capacit\'e d'extraction institutionnelle
doit \^etre adapt\'ee. Nous la d\'efinissons par
\begin{equation}
\Delta_q = \Delta + \delta_{\mathrm{enh}} - \delta_{\mathrm{degr}},
\label{eq:delta_q}
\end{equation}
o\`u $\delta_{\mathrm{enh}}$ repr\'esente l'am\'elioration potentielle due aux
outils quantiques (analyse quantique assist\'ee) et $\delta_{\mathrm{degr}}$
repr\'esente la d\'egradation due \`a la complexit\'e accrue des preuves PQC~\cite{Chatterjee2021}.

Le coefficient de perte d'interpr\'etabilit\'e quantique $\eta_q$ mesure la
d\'egradation s\'emantique induite par un canal quantique $\Lambda$ sur le
support de la preuve. Il s'agit du plus petit r\'eel tel que
\begin{equation}
I(V;\Lambda(\Sigma)|C) \geq I(V;\Sigma|C) - \eta_q\, H(C)
\label{eq:eta_q}
\end{equation}
(D\'efinition~\ref{def:perte_interp_quantique}). Ces param\`etres sont analys\'es en d\'etail \`a la section~\ref{sec:cro_pq}.

\subsubsection{Justification issue de la th\'eorie de l'information pour $\Delta$}
\label{subsubsec:justification_delta}

L'interpr\`ete institutionnel (arbitre) poss\`ede une connaissance pr\'ealable $K$
non captur\'ee dans le contexte formel $C$. Comme il n'acc\`ede au message que par la
preuve, on a $K \perp M \mid (\Sigma, C)$. La r\`egle de cha\^inage appliqu\'ee de deux
mani\`eres \`a $I(V;M,\Sigma|C)$ donne alors
\begin{equation}
I(V;\Sigma|C) \;=\; I(V;M|C) \;+\; I(V;\Sigma|M,C),
\label{eq:chaining_knowledge}
\end{equation}
et l'in\'egalit\'e de traitement des donn\'ees sur $M \to (\Sigma,C) \to V$ donne
$I(V;M|C) \leq I(M;\Sigma|C)$. Si l'institution tire de la preuve au plus
$\Delta$ bits \emph{au-del\`a du message}, $I(V;\Sigma|M,C) \leq \Delta$, d'o\`u
\begin{equation}
I(V;\Sigma|C) \;\leq\; I(M;\Sigma|C) + \Delta.
\label{eq:extraction_bornee_etat_art}
\end{equation}
C'est l'Hypoth\`ese d'Extraction Born\'ee~\cite{Smith09}, formalis\'ee \`a la
section~\ref{sec:cro_classique} par l'\'equation~\eqref{eq:extraction_bornee} et
fond\'ee par la Proposition~\ref{prop:fondement_extraction}.
\end{sloppypar}

\begin{boiteArbitre}
L'arbitre comprend que $\Delta$ le d\'esigne, puisque ce param\`etre mesure ce qu'il tire
de la preuve au-del\`a du message. Il s'agit peut-\^etre de quelques dizaines de bits, qu'il faudra mesurer.
\end{boiteArbitre}

% ----------------------------------------------------------------------------
\subsection{D\'efinitions pr\'eliminaires du cadre CRO}
\label{subsec:definitions_preliminaires_cro}
% ----------------------------------------------------------------------------
\begin{sloppypar}
Cette section introduit les trois
m\'etriques dans leur formulation issue de la th\'eorie de l'information, formulation qui,
seule, autorise le raisonnement asymptotique du Trilemme CRO.

\subsubsection{Confidentialit\'e, fiabilit\'e, opposabilit\'e}
\label{subsubsec:definitions_preliminaires}

La confidentialit\'e mesure l'indistinguabilit\'e au pire cas~\cite{GoldwasserMicali84} selon la d\'efinition
\begin{equation}
\Conf(\mathcal{P}) = 1 - \max_{m_0,m_1,c} \TV(P_{\Sigma|M=m_0,C=c},
P_{\Sigma|M=m_1,C=c}).
\label{eq:conf_def}
\end{equation}
La fiabilit\'e, qui quantifie la r\'ecup\'eration correcte, vaut
\begin{equation}
\Rel(\mathcal{P}) = 1 - \Pr[\hat{M} \neq M],
\label{eq:rel_def}
\end{equation}
o\`u $\hat{M}$ est l'estimateur optimal. L'opposabilit\'e capture la
r\'eduction relative d'incertitude sur le verdict institutionnel~\cite{MINKA2026CLO}, d\'efinie par
\begin{equation}
\Opp(\mathcal{P}, C) = \frac{I(V;\Sigma|C)}{H(V|C)},
\label{eq:opp_def}
\end{equation}
o\`u $V = \mathsf{Interpret}(\Sigma, C)$ est le verdict produit par
l'interpr\`ete institutionnel.

\begin{remark}[Coexistence de deux notations de l'opposabilit\'e]\label{rem:deux_notations}
Deux notations de l'opposabilit\'e coexistent, renvoyant toutes les deux \`a la m\^eme r\'ealit\'e, mais chacune servant un but pr\'ecis.
La notation de l'\'equation~\eqref{eq:opp_def} sert \`a pr\'eciser le protocole sous revue et le contexte particulier de son examen ou ex\'ecution, c'est une notation purement p\'edagogique, utilis\'ee notamment dans les \'equations~\eqref{eq:opp_helios}, \eqref{eq:opp_zksnark} et \eqref{eq:opp_ecdsa}.
La notation de l'\'equation~\eqref{eq:opp_clo} est celle utilis\'ee dans les d\'emonstrations et autres d\'eveloppements, pour sa sym\'etrie avec celles des deux autres m\'etriques.
\end{remark}

\subsubsection{Relation avec les d\'efinitions cryptographiques standard}
\label{subsubsec:relation_crypto}

Les m\'etriques CRO s'articulent avec les d\'efinitions cryptographiques classiques.

\textbf{S\'emantique s\'ecurit\'e.} La confidentialit\'e $\Conf$ g\'en\'eralise la s\'ecurit\'e s\'emantique de Goldwasser et Micali~\cite{GoldwasserMicali84}, puisque $\Conf = 1$ implique l'indistinguabilit\'e s\'emantique parfaite.

\textbf{Infalsifiabilit\'e existentielle.} La fiabilit\'e de \emph{v\'erification} $\Relver$ (Remarque~\ref{rem:rel_ver}) est reli\'ee \`a la notion d'infalsifiabilit\'e existentielle sous attaque adaptative~\cite{GMR88}, car une signature infalsifiable v\'erifie $\Relver \to 1$. La fiabilit\'e de reconstruction $\Rel$ de l'\'equation~\eqref{eq:rel_def} mesure, elle, la capacit\'e \`a restituer le message.

\textbf{Non-r\'epudiation.} L'opposabilit\'e $\Opp$ n'a pas d'\'equivalent direct dans la litt\'erature cryptographique classique, bien qu'elle soit implicitement vis\'ee par les travaux sur la non-r\'epudiation~\cite{ZhouGollmann}.

\subsubsection{M\'etriques quantiques}
\label{subsubsec:metriques_quantiques}

Les d\'efinitions quantiques g\'en\'eralisent les m\'etriques classiques~\cite{Watrous2009,Boneh2011}.

\textbf{Confidentialit\'e quantique.} On d\'efinit $\Conf^q$ par
\begin{equation}
\Conf^q(\mathcal{P}) = 1 - \frac{1}{2}\max_{m_0,m_1,c}
\|\rho_\Sigma^{m_0,c} - \rho_\Sigma^{m_1,c}\|_1.
\label{eq:conf_q_def}
\end{equation}

\textbf{Fiabilit\'e quantique.} La fiabilit\'e $\Rel^q$ est d\'efinie via l'estimateur POVM optimal
maximisant la probabilit\'e de r\'ecup\'eration correcte~\cite{Wilde2017}.

\textbf{Opposabilit\'e quantique.} L'institution obtient le verdict $V$
\`a partir de l'issue classique d'une mesure sur l'\'etat de preuve et l'on pose
\begin{equation}
\Opp^q(\mathcal{P}) = \frac{I(V;\Sigma|C)}{H(V|C)}.
\label{eq:opp_q_def}
\end{equation}
La borne de Holevo joue ici un r\^ole d\'ecisif, puisque aucune mesure ne permet \`a un interpr\`ete
institutionnel d'extraire de l'\'etat de preuve plus que son information de
Holevo~\cite{Holevo1973}.
\end{sloppypar}

% ----------------------------------------------------------------------------
\subsection{Trilemmes et bornes d'impossibilit\'e}
\label{subsec:trilemmes}
% ----------------------------------------------------------------------------
\begin{sloppypar}
Trois r\'esultats
d'impossibilit\'e ont directement guid\'e la formulation du Trilemme CRO. Tous suivent
le m\^eme sch\'ema, qui consiste \`a identifier une contrainte externe puis \`a d\'emontrer
qu'elle borne irr\'em\'ediablement les compromis r\'ealisables.

\subsubsection{Le th\'eor\`eme d'Arrow}
\label{subsubsec:arrow}

Le th\'eor\`eme d'impossibilit\'e d'Arrow~\cite{Arrow50} \'etablit qu'aucune r\`egle
d'agr\'egation des pr\'ef\'erences individuelles ne peut satisfaire simultan\'ement
cinq conditions rationnelles (universalit\'e, unanimit\'e, indiff\'erence aux
alternatives non pertinentes, souverainet\'e, non-dictature). Ce r\'esultat
fondateur de la th\'eorie du choix social illustre comment des contraintes
apparemment mod\'er\'ees peuvent engendrer des impossibilit\'es profondes~\cite{Sen}.

\subsubsection{Le th\'eor\`eme CAP}
\label{subsubsec:cap}

Le th\'eor\`eme CAP de Brewer~\cite{Brewer00}, prouv\'e formellement par Gilbert
et Lynch~\cite{GilbertLynch02}, \'etablit qu'un syst\`eme distribu\'e ne peut
simultan\'ement garantir coh\'erence (Consistency), disponibilit\'e (Availability)
et tol\'erance aux partitions (Partition tolerance). L'impossibilit\'e est
conditionnelle, puisqu'elle ne s'applique qu'en pr\'esence de partitions r\'eseau~\cite{ChandraHadzilacos96}.

\subsubsection{L'impossibilit\'e FLP}
\label{subsubsec:flp}

Fischer, Lynch et Paterson~\cite{FLP85} ont prouv\'e qu'aucun algorithme
d\'eterministe ne peut r\'esoudre le consensus dans un syst\`eme asynchrone en
pr\'esence d'une seule d\'efaillance. L'impossibilit\'e est conditionn\'ee \`a l'asynchronie~\cite{BenOr83,DebratMerz2012}.

\subsubsection{Le trilemme blockchain}
\label{subsubsec:trilemme_blockchain}

Buterin~\cite{Buterin17} a formalis\'e que les syst\`emes blockchain ne peuvent
simultan\'ement atteindre d\'ecentralisation, s\'ecurit\'e et \'evolutivit\'e
(scalability). Ce trilemme, contrairement \`a CAP et FLP, n'est pas issu d'une preuve
formelle mais d'une heuristique d'ing\'enierie~\cite{Nakamoto08,Garay15,Narayanan16}.

\subsubsection{Impossibilit\'es en cryptographie}
\label{subsubsec:impossibilites_crypto}

Plusieurs r\'esultats d'impossibilit\'e sp\'ecifient les limites fondamentales en cryptographie.

\textbf{Obfuscation.} Barak et al.~\cite{Barak2001} ont prouv\'e l'impossibilit\'e
de l'obfuscation virtuelle black-box pour les programmes g\'en\'eraux.

\textbf{Privacy-Utility trade-off.} La litt\'erature sur la \textit{privacy}
diff\'erentielle~\cite{Dwork} \'etablit des bornes sur les compromis entre
confidentialit\'e et utilit\'e.

\medskip
Ces paradigmes partagent une m\^eme structure, qui consiste \`a identifier une
contrainte externe (partition, asynchronie, d\'ecentralisation) puis \`a
d\'emontrer qu'elle borne les compromis r\'ealisables. Le Trilemme CRO suit ce
sch\'ema avec $\Delta$ comme contrainte dominante~; son extension post-quantique
ajoute $\eta_q$ comme contrainte suppl\'ementaire li\'ee aux effets quantiques.
\end{sloppypar}

% ----------------------------------------------------------------------------
\subsection{V\'erification formelle}
\label{subsec:verification_formelle}
% ----------------------------------------------------------------------------
\begin{sloppypar}
La v\'erification formelle
modifie le statut \'epist\'emique d'une preuve. Une preuve
v\'erifi\'ee m\'ecaniquement garantit l'exactitude de chaque \'etape de calcul et r\'ev\`ele
les hypoth\`eses que la version papier laissait implicites~\cite{Wiedijk2007}.

\subsubsection{Panorama des assistants de preuve}
\label{subsubsec:assistants_preuve}

Coq~\cite{Coq} (Calcul des Constructions Inductives, preuve du Th\'eor\`eme
des Quatre Couleurs~\cite{Gonthier08}, compilateur CompCert~\cite{Leroy09}) et Isabelle/HOL~\cite{Isabelle}
sont matures et largement d\'eploy\'es. Lean4~\cite{Lean4,moura2015lean} combine fondements
logiques rigoureux et programmation fonctionnelle performante, avec la
biblioth\`eque Mathlib~\cite{Mathlib} pour l'appui math\'ematique.

\subsubsection{Formalisation de la cryptographie}
\label{subsubsec:formalisation_crypto}

Plusieurs projets ont formalis\'e des r\'esultats cryptographiques.

\textbf{EasyCrypt.} Barthe et al.~\cite{BartheEA2013} ont d\'evelopp\'e un assistant
de preuve sp\'ecialis\'e pour les preuves cryptographiques computationnelles,
utilisant le calcul de relations probabilistes.

\textbf{F* et HACL*.} Le projet Everest~\cite{BhargavanEA2017,FStar,HACL} vise
une impl\'ementation v\'erifi\'ee de HTTPS, avec la biblioth\`eque cryptographique
HACL* formellement v\'erifi\'ee.

\textbf{CryptoVerif.} Blanchet~\cite{CryptoVerif} a d\'evelopp\'e un outil automatique
pour prouver les propri\'et\'es de s\'ecurit\'e des protocoles cryptographiques.

\textbf{CryptHOL.} Lochbihler~\cite{Lochbihler2016} a formalis\'e les jeux
cryptographiques en Isabelle/HOL.

\textbf{Th\'eorie de l'information avec Coq/Isabelle.} Affeldt et al.~\cite{Affeldt2015}
et Hibon-Portier~\cite{HibonPortier20} ont formalis\'e la th\'eorie de l'information
dans Coq et Isabelle respectivement.

\subsubsection{Lean4 pour le trilemme CRO}
\label{subsubsec:lean4_trilemme}

Lean4 a \'et\'e retenu \`a l'issue d'une comparaison syst\'ematique sur quatre crit\`eres.

\textbf{Couverture math\'ematique.} Mathlib (v4.34.0) fournit l'analyse r\'eelle,
les probabilit\'es et la divergence de Kullback--Leibler
(\texttt{InformationTheory.KullbackLeibler})~\cite{Mathlib}~; les in\'egalit\'es
de Pinsker et de Fano n'y sont pas formalis\'ees et sont donc pos\'ees comme
propri\'et\'es des instances de protocole (section~\ref{sec:lean4}).

\textbf{Proximit\'e syntaxique.} La syntaxe de Lean4 reste lisible par un
math\'ematicien sans formation sp\'ecialis\'ee en v\'erification formelle~\cite{moura2015lean}.

\textbf{Performance de compilation.} Lean4 compile vers C natif et
la v\'erification du noyau alg\'ebrique du Trilemme CRO prend environ une
minute et demie une fois Mathlib t\'el\'echarg\'ee. Le code fourni en
Annexe~\ref{chap:annexe_lean4}, fig\'e sur Lean~v4.34.0 et Mathlib~v4.34.0 dans le
d\'ep\^ot compagnon~\cite{CROTrilemmaDepot}, peut aussi \^etre coll\'e dans
l'\'editeur en ligne \url{https://live.lean-lang.org/}, qui utilise la version
courante de Mathlib.

\textbf{Dynamisme de la communaut\'e.} Croissance forte depuis 2022 avec
contributions actives en cryptographie formelle.

\subsubsection{Apports \'epist\'emiques de la formalisation}
\label{subsubsec:apports_epistemiques}

Au-del\`a de la certification~\cite{AbadiRogaway00,BanaComon12}, la formalisation
Lean4 (section~\ref{sec:lean4}) a produit trois apports.

\begin{enumerate}
\item \textbf{Explicitation d'une hypoth\`ese.} La formalisation a impos\'e de poser
explicitement le r\'egime de pr\'eservation de la confidentialit\'e
$I(M;\Sigma|C) \leq \kappa/\lambda$ (Hypoth\`ese~\ref{ass:regime}), que la preuve
publi\'ee d\'eduisait d'une \'etape d'optimisation. La formulation alternative
retenue le pose comme hypoth\`ese.

\item \textbf{Formulations alternatives.} La formalisation a conduit \`a des
formulations plus pr\'ecises de trois \'enonc\'es publi\'es
(section~\ref{subsec:lean4_erreurs}). La forme raffin\'ee du trilemme, qui
d\'elimite le r\'egime o\`u la capacit\'e institutionnelle est contraignante, a
\'egalement \'et\'e certifi\'ee.

\item \textbf{S\'eparation explicite.} La formalisation s\'epare le noyau alg\'ebrique
v\'erifi\'e des propri\'et\'es pos\'ees (Fano, Pinsker, hypoth\`eses institutionnelles),
clarifiant le statut \'epist\'emique de chaque composant~\cite{Zellagui2018}.
\end{enumerate}
\end{sloppypar}

% ----------------------------------------------------------------------------
%\subsection{Synth\`ese}
%\label{sec:synthese_chap2}
% ----------------------------------------------------------------------------

\begin{boiteArbitre}
L'arbitre referme le dossier. Ce seul chapitre lui a fait traverser six corpus
(th\'eorie de l'information, m\'ecanique quantique, primitives cryptographiques,
sciences cognitives, trilemmes d'impossibilit\'e, v\'erification formelle),
ce qui lui a d'abord paru beaucoup. Il comprend maintenant que chaque pi\`ece \'etait n\'ecessaire.

Ce qui le frappe en premier tient \`a une absence. Aucun des corpus pass\'es en
revue ne nommait l'incapacit\'e qu'il \'eprouve chaque jour \`a extraire
un verdict d'une preuve pourtant irr\'eprochable. La quantit\'e $\Opp = I(V;\Sigma|C)/H(V|C)$
n'apparaissait nulle part. Ce travail l'introduit comme m\'etrique de premi\`ere
classe, ce que l'arbitre re\c{c}oit comme une reconnaissance.

Il note ensuite la coh\'erence de l'\'edifice. Les g\'en\'eralisations quantiques
(Holevo, Pinsker quantique, Fano quantique) s'encha\^inent avec la m\^eme rigueur
que leurs analogues classiques, si bien que le cadre s'\'etend au r\'egime quantique
sans se fracturer.

L'armature analytique repose sur trois in\'egalit\'es dont d\'epend toute la
m\'ecanique du Trilemme CRO. L'in\'egalit\'e de Pinsker relie la confidentialit\'e
\`a l'information mutuelle, celle de Fano borne la fiabilit\'e par la probabilit\'e
d'erreur, tandis que la DPI garantit que le verdict n'invente pas d'information.

Le r\^ole de Lean4 l'\'etonne davantage. Outre la certification, la formalisation
a contraint \`a expliciter une hypoth\`ese et \`a pr\'eciser plusieurs \'enonc\'es. L'arbitre
ne s'attendait pas \`a voir la rigueur formelle servir d'instrument de d\'ecouverte.

Il tourne la page. Six objectifs de recherche l'attendent (Obj~\ref{obj:obj1_unification},~Obj~\ref{obj:obj2_impossibilite},~Obj~\ref{obj:obj3_postquantique},~Obj~\ref{obj:obj4_saturation},~Obj~\ref{obj:obj5_verification} et~Obj~\ref{obj:obj6_experimentation}). Il veut savoir
s'ils tiennent ensemble.
\end{boiteArbitre}
% ============================================================================
% CHAPITRE 1 — CONTEXTE, \'ETAT DE L'ART ET PROBL\'EMATIQUE
% Section 1.4 : Probl\'ematique de Recherche
% source : partie1/problematique.tex
% ============================================================================

\section{PROBL\'EMATIQUE DE RECHERCHE}
\label{sec:problematique}
\begin{sloppypar}
La section~\ref{sec:contexte} a
document\'e trois \'echecs institutionnels r\'eels (Louvain~2009~\cite{Adida09UCL},
IACR~2025~\cite{IACR2025incident} et le vide de contr\^ole AML~\cite{FATF12,AMLD4}) ainsi qu'un succ\`es plus nuanc\'e, celui d'ECDSA~\cite{Johnson01,NIST.FIPS.186-5}.
La section~\ref{sec:etat_art} a ensuite rassembl\'e les outils n\'ecessaires
\`a leur quantification, \`a savoir l'entropie classique et quantique~\cite{Shannon48,NielsenChuang2010}, l'information mutuelle,
les in\'egalit\'es de Pinsker~\cite{Pinsker1964} et de Fano~\cite{Fano61} ainsi que les trilemmes d'impossibilit\'e~\cite{Arrow50,Brewer00,FLP85}. Il reste \`a convertir ces observations en un programme de recherche
formalis\'e, ce qui fait l'objet de la pr\'esente section.

Notre d\'emarche reprend la progression canonique d\'ecrite au paragraphe~\ref{par:progression_FBESCX}.
Nous formalisons d'abord, \`a la sous-section~\ref{subsec:paradoxe_authenticite},
l'incompatibilit\'e structurelle entre paradigme cryptographique et
paradigme institutionnel (\textbf{F}). La
sous-section~\ref{subsec:objectifs_recherche} \'etablit ensuite les six
objectifs de recherche (Obj~\ref{obj:obj1_unification}$\to$Obj~\ref{obj:obj2_impossibilite}$\to$Obj~\ref{obj:obj3_postquantique}$\to$Obj~\ref{obj:obj4_saturation}$\to$Obj~\ref{obj:obj5_verification}$\to$Obj~\ref{obj:obj6_experimentation}) ainsi que les contributions qui leur
correspondent. Il s'agit de borner les compromis classiques (\textbf{B}) puis
d'\'etendre ces bornes aux contextes post-quantiques (\textbf{E}). Viennent
ensuite l'approche des bornes par des constructions dont l'optimalit\'e
est examin\'ee (\textbf{S}) et la certification des r\'esultats par v\'erification formelle
(\textbf{C}), puis la sp\'ecification d'un programme exp\'erimental qui mesure les
param\`etres institutionnels (\textbf{X}). Enfin, la sous-section~\ref{subsec:objectifs_programme} examine
la coh\'erence interne du programme et en esquisse les retomb\'ees scientifiques,
pratiques et soci\'etales.
\end{sloppypar}

\begin{boiteArbitre}
Le foss\'e que l'arbitre constate entre validit\'e technique et conviction
institutionnelle trouve ici une formulation pr\'ecise, celle du paradoxe de
l'authenticit\'e invisible. Ce paradoxe na\^it de l'incompatibilit\'e structurelle entre
v\'erification algorithmique et v\'erification substantielle~\cite{MINKA2026CLO,eIDAS14}.
Il en d\'ecoule six questions, qui vont de la formalisation du paradoxe
jusqu'\`a son exp\'erimentation. L'arbitre se les posait d\'ej\`a sans parvenir
\`a les formuler~; loin de lui para\^itre acad\'emiques, elles r\'epondent \`a ce
qu'il attendait.
\end{boiteArbitre}

% ----------------------------------------------------------------------------
\subsection{Le paradoxe de l'authenticit\'e invisible~: formulation approfondie}
\label{subsec:paradoxe_authenticite}
% ----------------------------------------------------------------------------
\begin{sloppypar}
Deux s\'eries d'\'el\'ements convergent vers une m\^eme tension. La premi\`ere r\'eunit
les cas de la section~\ref{sec:contexte}, \`a savoir l'\'echec op\'erationnel de Helios \`a
Louvain~(2009)~\cite{Adida09UCL,Goodman15} et lors de l'\'election IACR~(2025)~\cite{IACR2025incident}, le vide de contr\^ole
des zk-SNARKs pour la conformit\'e AML/B\^ale~III~\cite{BaselIII,Garratt19}
ainsi que l'acceptation judiciaire universelle d'ECDSA~\cite{ESIGN00,eIDAS14}, obtenue au prix d'une
confidentialit\'e nulle. La seconde rassemble les fondements pos\'es \`a la section~\ref{sec:etat_art}, qu'il s'agisse de la th\'eorie de l'information de
Shannon~\cite{Shannon48,Shannon49,CoverThomas06}, de la th\'eorie quantique~\cite{Wilde2017,Holevo1973}, des primitives d'obscurcissement~\cite{Pedersen91,Groth16} ou des trilemmes
CAP/FLP/Blockchain~\cite{GilbertLynch02,Buterin17,Nakamoto08}. Nous formalisons maintenant cette tension comme probl\'ematique centrale.

\subsubsection{Le foss\'e cryptographie-institution~: analyse structurelle}
\label{subsubsec:fosse_crypto_institution}

Les syst\`emes cryptographiques con\c{c}us pour la s\'ecurit\'e math\'ematique
\'echouent dans les contextes institutionnels exigeant une interpr\'etation
humaine non-experte~\cite{Bonneau12,Whitten99,Herley09}. Cette d\'efaillance ne se corrige ni par une
meilleure impl\'ementation ni par une formation suppl\'ementaire des
utilisateurs, car elle proc\`ede d'une \textbf{incompatibilit\'e structurelle}
entre deux paradigmes de v\'erification.

\paragraph{Paradigme cryptographique~: v\'erification algorithmique.}
\label{par:paradigme_crypto}

La recherche cryptographique classique, depuis les travaux de
Goldwasser, Micali et Rackoff~\cite{GMR85}, mesure la s\'ecurit\'e via
des \textbf{jeux} opposant un challenger \`a un adversaire polynomial~\cite{Goldreich01,Shoup01}.
Trois propri\'et\'es en ressortent.

La confidentialit\'e est mesur\'ee par l'avantage du distingueur dans le
jeu IND-CPA~\cite{GoldwasserMicali84}, qui doit rester n\'egligeable
\begin{equation}
\label{eq:ind_cpa_advantage}
\mathsf{Adv}^{\text{IND-CPA}}_{\mathcal{A}}(\lambda) =
\left|\Pr[\mathcal{A} \text{ gagne}] - \tfrac{1}{2}\right| \leq \negl(\lambda).
\end{equation}
L'int\'egrit\'e est mesur\'ee par la r\'esistance aux collisions~\cite{Damgard87}, qui exige que, pour tout
adversaire polynomial $\mathcal{A}$,
\begin{equation}
\label{eq:collision_resistance}
\Pr[x \neq x' \land H(x) = H(x') : (x,x') \leftarrow \mathcal{A}(1^\lambda)]
\leq \negl(\lambda).
\end{equation}
L'authenticit\'e enfin est mesur\'ee par la r\'esistance \`a la \textit{forgery}
existentielle dans le jeu EUF-CMA~\cite{GMR88,Schnorr91}, c'est-\`a-dire par l'incapacit\'e d'un adversaire \`a produire une signature valide sur un message non autoris\'e
\begin{equation}
\label{eq:euf_cma}
\Pr[\mathsf{Verify}(vk, m^*, \sigma^*) = 1 \land m^* \notin \mathcal{Q}
: (m^*, \sigma^*) \leftarrow \mathcal{A}^{\mathcal{O}_{\text{sign}}}(vk)]
\leq \negl(\lambda).
\end{equation}

Ces trois d\'efinitions portent sur des propri\'et\'es \textbf{adversariales},
puisqu'elles quantifient la capacit\'e d'un attaquant \`a violer la s\'ecurit\'e.
Elles ne fournissent en revanche aucune m\'etrique des propri\'et\'es
\textbf{interpr\'etatives}, c'est-\`a-dire de la capacit\'e d'un utilisateur
l\'egitime \`a comprendre ce que le syst\`eme
certifie~\cite{Smith09,Alvim12,ClarksonSchneider10}.

\paragraph{Paradigme institutionnel~: v\'erification substantielle.}
\label{par:paradigme_institutionnel}

Les tribunaux, les organes de r\'egulation et les audits publics exigent quant
\`a eux une v\'erification \textbf{substantielle},
qualitativement distincte de la v\'erification algorithmique~\cite{Casey11,Casey16,Mason07}.
Celle-ci repose sur trois exigences.

La premi\`ere est l'intelligibilit\'e par des non-experts~\cite{Kerr04,Daubert1993}. Un juge, un
auditeur ou un r\'egulateur (un arbitre) doit pouvoir saisir \emph{ce qui est prouv\'e} sans recourir
syst\'ematiquement \`a des experts externes. Cette exigence entre en tension directe
avec le principe zero-knowledge, dont l'objet m\^eme est de minimiser
l'information r\'ev\'el\'ee~\cite{BenSasson2019,PLONK20}. La deuxi\`eme tient \`a l'auditabilit\'e r\'etrospective~\cite{Twining1999}.
Une d\'ecision institutionnelle doit rester r\'evisable plusieurs ann\'ees apr\`es
avoir \'et\'e prise, ce qui suppose des traces explicites du raisonnement suivi~;
or une preuve zero-knowledge ne laisse aucune trace exploitable
au-del\`a du r\'esultat binaire \guillemotleft~valide/invalide~\guillemotright. La troisi\`eme
concerne l'attribution des responsabilit\'es~\cite{Citron14}, puisqu'en cas de d\'efaillance ou de fraude
il doit \^etre possible d'identifier les acteurs responsables et de reconstituer
la cha\^ine causale.

\paragraph{Le paradoxe.}
\label{par:paradoxe_def}

Le \textbf{paradoxe de l'authenticit\'e invisible} na\^it de cet \'ecart. Il d\'esigne
la situation dans laquelle une preuve cryptographique peut \^etre \`a la fois math\'ematiquement
impeccable et institutionnellement opaque~\cite{Doshi-Velez17,Adadi2018,Barredo2020}. Il se
manifeste en pratique de trois fa\c{c}ons.

La premi\`ere est une v\'erification d\'epourvue de compr\'ehension. Un garant \'electoral peut ex\'ecuter
m\'ecaniquement un algorithme de v\'erification zero-knowledge et obtenir
un r\'esultat binaire sans comprendre \emph{pourquoi} la preuve est
valide ni \emph{ce qu'elle certifie exactement}. Une telle dissociation
entre v\'erification formelle et compr\'ehension substantielle affaiblit la
responsabilit\'e institutionnelle~\cite{Adams1999,Cranor05}. La deuxi\`eme manifestation tient \`a l'impossibilit\'e d'un audit a posteriori. Les
preuves zero-knowledge sont con\c{c}ues pour ne rien r\'ev\'eler
au-del\`a de la validit\'e de l'assertion~; cette propri\'et\'e, pr\'ecieuse pour
la confidentialit\'e, prive en revanche l'analyse r\'etrospective de sa mati\`ere~\cite{Androulaki13,Meiklejohn13}.
Lors d'une crise syst\'emique comparable \`a celle de 2008~\cite{FCIC2011}, la reconstitution des
m\'ecanismes de contagion devient ainsi structurellement impossible. La troisi\`eme manifestation est l'impossibilit\'e d'attribuer les responsabilit\'es.
Lorsqu'une preuve est rejet\'ee ou qu'une fraude est suspect\'ee, l'opacit\'e
cryptographique emp\^eche d'identifier les acteurs qui ont produit des preuves incorrectes ou
malveillantes~\cite{FinckMoscon19,DeFilippi16}. La figure suivante illustre ce compromis.

\begin{figure}[htbp]
\centering
\begin{tikzpicture}[
    scale=0.9,
    transform shape,
    box/.style={rectangle, draw=enspblue, fill=blue!8, thick,
                minimum width=3.5cm, minimum height=1.2cm,
                align=center, font=\small},
    arrow/.style={-{Stealth[length=3mm]}, thick, enspblue},
    lienTension/.style={-{Stealth[length=3mm]}, thick, enspred, dashed}
]

% Trois syst\`emes
\node[box] (helios) at (0,4)
    {\textbf{Helios}\\$\Conf \approx 0{,}95$ (\'elev\'e)\\$\Opp \approx 0{,}50$ (faible)};
\node[box] (zksnark) at (5,4)
    {\textbf{zk-SNARKs}\\$\Conf \approx 0{,}98$ (\'elev\'e)\\$\Opp \approx 0{,}67$ (mod\'er\'e)};
\node[box] (ecdsa) at (10,4)
    {\textbf{ECDSA}\\$\Conf \approx 0{,}00$ (faible)\\$\Opp \approx 0{,}90$ (\'elev\'e)};

% Propri\'et\'es
\node[box, fill=enspblue!12, draw=enspblue] (conf) at (0,0)
    {\textbf{Confidentialit\'e}\\Protection de la vie priv\'ee\\$\Conf \to 1$};
\node[box, fill=navyblue!12, draw=navyblue] (opp) at (10,0)
    {\textbf{Opposabilit\'e}\\Utilit\'e institutionnelle\\$\Opp \to 1$};

% Fl\`eches
\draw[arrow] (helios)  -- (conf)
    node[midway, left,  font=\footnotesize] {$\Conf \cdot \Opp \approx 0{,}48$};
\draw[arrow] (zksnark) -- (conf)
    node[midway, left,  font=\footnotesize, xshift=2cm]
        {$\Conf \cdot \Opp \approx 0{,}66$};
\draw[arrow] (ecdsa)   -- (opp)
    node[midway, right, font=\footnotesize] {$\Conf \cdot \Opp \approx 0$};

% Tension fondamentale
\draw[lienTension, <->] (conf) -- (opp)
    node[midway, above, font=\small, yshift=0.3cm]
        {\textbf{Tension fondamentale}};

% \'Etiquettes
\node[below=0.5cm of conf, font=\footnotesize, text=enspblue]
    {Confidentialit\'e forte};
\node[below=0.5cm of opp,  font=\footnotesize, text=navyblue]
    {Acc\`es institutionnel fort};

\end{tikzpicture}
\caption[Paradoxe de l'authenticit\'e invisible~: compromis confidentialit\'e--opposabilit\'e]{%
  Paradoxe de l'authenticit\'e invisible~: compromis entre confidentialit\'e et
  opposabilit\'e dans les syst\`emes cryptographiques d\'eploy\'es.%
}

\smallskip
\begin{minipage}{\linewidth}
  \small\textit{Note~:} Les syst\`emes optimis\'es pour la confidentialit\'e
  (Helios, zk-SNARKs) sacrifient l'interpr\'etabilit\'e institutionnelle,
  tandis que les syst\`emes transparents (ECDSA) sacrifient la confidentialit\'e.
\end{minipage}

\end{figure}

\begin{boiteArbitre}
L'arbitre comprend qu'une preuve peut \^etre \`a la fois correcte
et inutilisable. Il reconna\^it dans cette formulation une situation qu'il
rencontre chaque jour et qui trouve ici sa formalisation.
\end{boiteArbitre}

\subsubsection{Une nouvelle dimension du paradoxe}
\label{subsubsec:nouvelle_dimension}

La transition vers la cryptographie post-quantique, inaugur\'ee par les
standards NIST FIPS~203, 204 et 205 en ao\^ut~2024~\cite{NISTPQC,Alagic2020,NISTPostQuantum2022}, introduit une nouvelle dimension au paradoxe de l'authenticit\'e invisible.

\paragraph{Complexit\'e accrue des preuves PQC.}
\label{par:complexite_pqc}

Les signatures Dilithium (FIPS~204) reposent sur des preuves de connaissance
d\'efinies sur des r\'eseaux euclidiens~\cite{Dilithium2022,Langlois2015}, dont la v\'erification met en jeu des op\'erations matricielles
sur $\mathbb{Z}_q$ et des calculs de normes. Un juge habitu\'e \`a v\'erifier des signatures
ECDSA doit d\'esormais ma\^itriser les probl\`emes LWE (Learning With Errors)~\cite{Regev2009}
ainsi que les r\'eductions cryptographiques qui s'y rattachent. Le foss\'e entre validit\'e technique et interpr\'etabilit\'e
institutionnelle risque donc de se creuser~\cite{Ravi2020}.

\paragraph{Adversaires quantiques et canaux quantiques.}
\label{par:adversaires_quantiques}

Un adversaire quantique peut interroger les oracles cryptographiques par des \textbf{requ\^etes en
superposition}~\cite{Boneh2011,Watrous2009,Zhandry}, ce qui est susceptible de modifier
les bornes de s\'ecurit\'e classiques. Les canaux de
communication quantiques subissent aussi des d\'egradations intrins\`eques,
telles que la d\'ecoh\'erence ou le bruit de mesure, qui alt\`erent l'interpr\'etabilit\'e des
preuves transmises~\cite{Zurek2003,Wootters1982}. Ce sont ces effets quantiques~\cite{Mayers1997,Lo1997,Unruh2016} qui conduisent \`a introduire
les param\`etres $\Delta_q$ et $\eta_q$, formalis\'es \`a la section~\ref{sec:cro_pq}.

\subsubsection{Formulation math\'ematique~: implications pr\'eliminaires}
\label{subsubsec:formulation_mathematique}

Quatre conclusions de l'analyse pr\'ec\'edente structurent notre programme
de recherche.

Il faut d'abord faire de l'opposabilit\'e une propri\'et\'e de premi\`ere classe. Les
d\'efaillances observ\'ees (Louvain~2009, IACR~2025, vide de contr\^ole
AML) ne se r\'esorbent pas avec de meilleures interfaces ou des
formations suppl\'ementaires, car elles t\'emoignent de limites structurelles
qui appellent une refonte conceptuelle de la cryptographie institutionnelle~\cite{Bonneau16,Narayanan16}.

La capacit\'e cognitive humaine ($\Delta$) est la contrainte dominante.
Contrairement aux analyses cryptographiques classiques centr\'ees sur la
puissance de calcul de l'adversaire (mesur\'ee en temps polynomial), la
contrainte limitante en contexte institutionnel est la capacit\'e
d'extraction issue de la th\'eorie de l'information des arbitres. Estim\'ee
entre~20 et~80 bits selon la th\'eorie de la rationalit\'e
born\'ee~\cite{Simon82,Simon1955,Miller56,Cowan2001,Ericsson95}, elle est syst\'ematiquement inf\'erieure au param\`etre de s\'ecurit\'e cryptographique $\lambda \in [128,256]$ bits.

Les compromis observ\'es laissent aussi pr\'esager des bornes fondamentales.
Malgr\'e des d\'ecennies d'optimisation, aucun syst\`eme
d\'eploy\'e n'atteint un produit $\Conf \cdot \Rel \cdot \Opp > 0{,}6$.
Le fait que tous plafonnent au m\^eme niveau s'accorde avec l'existence de contraintes th\'eoriques
comparables au th\'eor\`eme~CAP~\cite{Brewer00,GilbertLynch02} ou au trilemme
blockchain~\cite{Buterin17,Garay15}.

L'extension post-quantique est n\'ecessaire mais non triviale. Les bornes
classiques ne s'appliquent pas directement aux contextes quantiques,
n\'ecessitant une reformulation avec des outils issus de la th\'eorie de l'information
quantique (entropie de von Neumann, information de Holevo, canaux CPTP)~\cite{NielsenChuang2010,Wilde2017,Hiai1991}.

\subsubsection{Formulation de la probl\'ematique}
\label{subsubsec:formulation_problematique}

Nous affinons ici la probl\'ematique annonc\'ee \`a la section~\ref{sec:ig_problematique}, qui s'\'enonce d\'esormais comme suit.

\begin{boiteProblematique}{Probl\'ematique centrale}
\textit{Quelles limites fondamentales contraignent les compromis entre confidentialit\'e, fiabilit\'e et opposabilit\'e dans les syst\`emes cryptographiques soumis \`a des exigences d'interpr\'etabilit\'e institutionnelle~?}
\end{boiteProblematique}
\end{sloppypar}

\begin{boiteArbitre}
L'arbitre comprend que cette question est la sienne.
Autrement dit, jusqu'o\`u peut-il exiger qu'on lui
explique ce qu'une preuve certifie~?
\end{boiteArbitre}

% ----------------------------------------------------------------------------
\subsection{Objectifs de recherche et contributions associ\'ees}
\label{subsec:objectifs_recherche}
% ----------------------------------------------------------------------------
\begin{sloppypar}
Les quatre implications pr\'eliminaires de la section~\ref{subsubsec:formulation_mathematique} engendrent six objectifs de recherche qui d\'ecoulent logiquement de la probl\'ematique. Elles sont pr\'esent\'ees ici dans l'ordre F$\to$B$\to$E$\to$S$\to$C.

\subsubsection{Unification th\'eorique, cadre CLO}
\label{subsubsec:obj1_clo}

\begin{objectif}[Unification Th\'eorique]\label{obj:obj1_unification}
\textit{\'Etablir un cadre unifi\'e qui formalise math\'ematiquement la confidentialit\'e, la fiabilit\'e et l'opposabilit\'e au moyen de mesures issues de la th\'eorie de l'information.}
\end{objectif}

\paragraph{Motivations.}
Sans formalisation commune, les observations qualitatives, \guillemotleft~Helios
est difficile \`a auditer~\guillemotright, restent des anecdotes. Leur transformation
en m\'etriques quantitatives ($\Opp = I(V;\Sigma|C)/H(V|C)$) est la
condition pr\'ealable \`a toute th\'eorie~\cite{Clarkson05,Smith09,Alvim12}.

\paragraph{Contribution~C1~: cadre CLO (Cryptographic Legal Opposability).}
Nous d\'efinissons trois mesures issues de la th\'eorie de l'information,
\`a savoir la confidentialit\'e $\Conf$, la fiabilit\'e $\Rel$ et l'opposabilit\'e $\Opp$ (voir
\'equations~\eqref{eq:conf_def}--\eqref{eq:opp_def} de la section~\ref{subsec:definitions_preliminaires_cro}). Le cadre \'etablit le Th\'eor\`eme du Goulot
d'\'etranglement Contextuel $I(V;\Sigma|C) \leq \alpha H(\mathcal{C})$,
fournit des r\`egles de composition (s\'equentielle et parall\`ele) et
propose la construction Semantic One-Time Pad (S-OTP). Les \'etudes de cas
illustratives portant sur Helios ($\Conf \cdot \Relver \cdot \Opp \approx
0{,}437$), zk-SNARKs ($\approx 0{,}572$) et ECDSA ($\approx 0{,}273$), sur
l'\'echelle op\'erationnelle CLO, font appara\^itre un spectre de compromis
qui motive la recherche de contraintes fondamentales~\cite{Adida08,Groth16,Johnson01}.

\textbf{Publication~:} Article publi\'e, IEEE Access
(DOI~: 10.1109/ACCESS.2026.3660105)~\cite{MINKA2026CLO}.

\subsubsection{Bornes d'impossibilit\'e classiques, Trilemme conditionnel CRO}
\label{subsubsec:obj2_cro}

\begin{objectif}[Bornes d'Impossibilit\'e Classiques]\label{obj:obj2_impossibilite}
\textit{\'Etablir les limites fondamentales, conditionn\'ees par la capacit\'e cognitive institutionnelle, des compromis r\'ealisables entre confidentialit\'e, fiabilit\'e et opposabilit\'e.}
\end{objectif}

\paragraph{Motivations.}
Le fait qu'aucune des estimations illustratives ne d\'epasse~$0{,}6$ pourrait
n'\^etre qu'un artefact de l'ing\'enierie disponible ou de l'estimation. Il reste donc \`a d\'eterminer si cette barri\`ere
est fondamentale. La r\'eponse requiert un th\'eor\`eme d'impossibilit\'e
rigoureusement prouv\'e, conditionnant l'impossibilit\'e \`a des param\`etres
institutionnels explicites, \`a l'image du th\'eor\`eme~CAP conditionn\'e
aux partitions r\'eseau~\cite{GilbertLynch02,ChandraHadzilacos96}.

\paragraph{Contribution~C2~: Trilemme CRO conditionnel.}
Nous \'etablissons une borne sup\'erieure sur le produit
$\Conf \cdot \Rel \cdot \Opp$ en fonction des param\`etres institutionnels
$\Delta$ (capacit\'e d'extraction), $\alpha$ (granularit\'e des verdicts),
$\beta$ (r\'egularit\'e du protocole)~\cite{Tsybakov2009}, sous la forme
\begin{equation}
\label{eq:cro_bound_intro}
\Conf \cdot \Rel \cdot \Opp \leq
\frac{\Delta}{\gamma\lambda} + O(\lambda^{-2}),
\qquad \gamma = \alpha H(\mathcal{C}) - 1,
\end{equation}
o\`u $\lambda$ est le param\`etre de s\'ecurit\'e. La forme publi\'ee
dans~\cite{MINKA2026CRO}, $\Conf \cdot \Rel \cdot \Opp \leq \Delta/\lambda +
O(\lambda^{-2}) + \negl(\lambda)$, en est un corollaire affaibli ($\gamma \geq 1$).
Une forme raffin\'ee identifie le r\'egime, $H(\mathcal{C}) > (\Delta + 1 +
\kappa/\lambda)/\alpha$, dans lequel la capacit\'e institutionnelle $\Delta$
devient la contrainte liante \`a la place de la seule in\'egalit\'e de Fano. Une
calibration illustrative situe $\Delta$ dans quelques dizaines de bits,
ordre de grandeur plausible au regard de la th\'eorie de la rationalit\'e
born\'ee~\cite{Simon82,Kahneman11}.

\textbf{Publication~:} Article publi\'e, IEEE Access
(DOI~: 10.1109/ACCESS.2026.3669065)~\cite{MINKA2026CRO}.

\subsubsection{Extension quantique, Trilemme conditionnel CRO}
\label{subsubsec:obj3_pqcro}

\begin{objectif}[Extension Post-Quantique]\label{obj:obj3_postquantique}
\textit{\'Etendre ces bornes aux adversaires quantiques et aux effets de canal quantique, puis caract\'eriser les termes propres au r\'egime post-quantique.}
\end{objectif}

\paragraph{Motivations.}
La migration vers les standards NIST PQC soul\`eve une question in\'edite, celle de savoir si
les bornes \'etablies dans le cadre classique restent valides face \`a
des adversaires capables de requ\^etes quantiques~\cite{Kiltz2017}. Elle conduit aussi \`a se demander si les canaux quantiques
introduisent une perte d'interpr\'etabilit\'e suppl\'ementaire~\cite{Schumacher1997}. Ces
questions sont urgentes car les institutions doivent d\`es maintenant
choisir des protocoles adapt\'es \`a l'\`ere post-quantique.

\paragraph{Contribution~C3~: Trilemme CRO post-quantique.}
Nous \'etendons le Trilemme CRO aux contextes quantiques~\cite{Holevo1973,Wilde2017} en \'etablissant la borne
\begin{equation}
\label{eq:pqcro_bound_intro}
\Conf^q \cdot \Rel^q \cdot \Opp^q \leq
\frac{\Delta_q + \eta_q \cdot H(C)}{\lambda} + O(\lambda^{-2}) + \negl(\lambda),
\end{equation}
o\`u $\Delta_q$ quantifie la capacit\'e d'extraction institutionnelle adapt\'ee
au quantum et $\eta_q$ le coefficient de perte d'interpr\'etabilit\'e quantique.
La preuve mobilise des outils de th\'eorie de l'information quantique,
\`a savoir l'in\'egalit\'e de Pinsker quantique, le lemme de Gentle Measurement~\cite{Wilde2017} et les bornes
de Fano quantiques. La borne explicite conserve la forme classique, avec
$\Delta_q$~; le terme $\eta_q H(C)$ n'appara\^it que dans la majoration
affaiblie ci-dessus. Une calibration illustrative sur des signatures
standardis\'ees (ML-DSA~\cite{Dilithium2022}, SLH-DSA~\cite{SPHINCS2022}) situe
$\Delta_q$ dans quelques dizaines de bits~\cite{Mayers1997,Lo1997,Unruh2016,Chatterjee2021}.

\textbf{Publication~:} Article publi\'e, IEEE Access, vol.~14, pp.~125130--125155
(DOI~: 10.1109/ACCESS.2026.3722742)~\cite{MINKA2026PQCRO}.

\subsubsection{Constructions optimales, famille SOTP}
\label{subsubsec:obj4_sotp}

\begin{objectif}[Constructions optimales]\label{obj:obj4_saturation}
\textit{Concevoir des protocoles qui atteignent asymptotiquement ces bornes puis \'etablir leur optimalit\'e, de sorte que le r\'esultat d'impossibilit\'e serve de guide \`a la conception.}
\end{objectif}

\paragraph{Motivations.}
Tant qu'une borne sup\'erieure n'est pas montr\'ee serr\'ee, elle garde le statut d'une
conjecture de m\'ethode. Sa saturation par des constructions explicites permet de faire du Trilemme~CRO
une caract\'erisation fondamentale des limites de la cryptographie
institutionnelle~\cite{Gennaro13,Parno13,BCTV14}.

\paragraph{Contribution~C4~: famille SOTP.}
Nous proposons trois constructions de r\'ef\'erence, chacune adapt\'ee \`a un
r\'egime de richesse contextuelle. Sous des conventions explicites, elles r\'ealisent
des compromis \'elev\'es sur l'\'echelle op\'erationnelle CLO
($\Conf$, $\Relver$, $\Opp$). SOTP-Low, destin\'ee aux contextes rares, repose sur un
encodage d\'eterministe du contexte. SOTP-Base exploite la compression
Johnson-Lindenstrauss~\cite{JohnsonLindenstrauss84,DasguptaGupta03,Vempala05},
tandis que SOTP-High, pr\'evue pour les contextes riches, s'appuie sur les zk-SNARKs de
Groth16~\cite{Groth16,BenSasson2019,PLONK20}. Les propri\'et\'es
cryptographiques de ces constructions sont d\'emontr\'ees sous les hypoth\`eses PRG et DLog.
L'existence de constructions approchant la borne du Trilemme, laquelle porte sur la fiabilit\'e
de reconstruction, demeure en revanche formul\'ee comme conjecture~\cite{Campanelli2019,Wahby18}.

\textbf{Publication~:} Soumis, IEEE TIT.

\subsubsection{V\'erification formelle, Lean4}
\label{subsubsec:obj5_lean4}

\begin{objectif}[V\'erification Formelle]\label{obj:obj5_verification}
\textit{Obtenir, dans un assistant de preuve, une certification m\'ecanique du noyau du th\'eor\`eme d'impossibilit\'e en explicitant les hypoth\`eses sur lesquelles il repose.}
\end{objectif}

\paragraph{Motivations.}
Les preuves du Trilemme~CRO encha\^inent, sur plusieurs dizaines de pages,
des in\'egalit\'es issues de la th\'eorie de l'information. Or l'exp\'erience montre
que des preuves de cette ampleur peuvent comporter des failles subtiles qui
\'echappent aux relectures humaines~\cite{Shoup01,Leroy09}.
Lorsque les r\'esultats engagent directement des enjeux soci\'etaux, comme la s\'ecurit\'e
\'electorale ou la r\'egulation financi\`ere, le niveau d'assurance requis d\'epasse
celui qu'offrent des preuves informelles~\cite{BartheEA2013,BhargavanEA2017}.

\paragraph{Contribution~C5~: formalisation Lean4.}
Le noyau alg\'ebrique du Trilemme~CRO est v\'erifi\'e m\'ecaniquement
dans l'assistant de preuve Lean4~\cite{Lean4,moura2015lean} (version~4.34.0,
biblioth\`eque Mathlib~\cite{Mathlib} v4.34.0). La formalisation apporte
quatre contributions. Elle fournit d'abord un mod\`ele formel de familles de protocoles,
$\mathtt{ProtocolFamily}$, qui capture la nature asymptotique du th\'eor\`eme.
Elle s\'epare ensuite explicitement le noyau v\'erifi\'e des in\'egalit\'es de Pinsker
et de Fano, pos\'ees comme propri\'et\'es des instances~\cite{AbadiRogaway00,BanaComon12}.
Elle a aussi conduit \`a expliciter l'Hypoth\`ese de r\'egime $I(M;\Sigma|C) \leq \kappa/\lambda$
et \`a pr\'eciser trois \'enonc\'es publi\'es. Elle certifie enfin la
forme raffin\'ee, qui d\'elimite le r\'egime o\`u la capacit\'e institutionnelle
est contraignante~\cite{Affeldt2015,HibonPortier20,Lochbihler2016}.

\textbf{Publication~:} Soumis, Journal of Automated Reasoning.
\end{sloppypar}


% ----------------------------------------------------------------------------
\subsubsection{Exp\'erimentation, protocoles de mesure et validation r\'etrospective}
\label{subsubsec:obj6_experimentation}
% ----------------------------------------------------------------------------
\begin{sloppypar}

\begin{objectif}[Exp\'erimentation]\label{obj:obj6_experimentation}
\textit{Sp\'ecifier un dispositif exp\'erimental falsifiable permettant de mesurer les param\`etres institutionnels qui conditionnent les bornes, puis confronter le cadre aux syst\`emes d\'eploy\'es.}
\end{objectif}

\paragraph{Motivations.}
Un th\'eor\`eme conditionnel ne vaut que par la mesurabilit\'e de ses
conditions. Les bornes \'etablies d\'ependent de six grandeurs,
$\Delta$, $\alpha$, $\beta$, $H(\mathcal{C})$, $\Delta_q$ et $\eta_q$, dont
aucune ne dispose \`a ce jour de proc\'edure de mesure normalis\'ee. Sans
d\'efinition op\'erationnelle, ces param\`etres demeurent des postulats et
les bornes restent muettes sur les syst\`emes r\'eels~\cite{Popper1934}. Un
r\'esultat n'acquiert sa port\'ee scientifique qu'\`a la condition d'\^etre
r\'efutable, ce qui oblige \`a pr\'eciser les mesures par lesquelles il pourrait l'\^etre.

\paragraph{Contribution~C6~: programme exp\'erimental et validation r\'etrospective.}
Nous sp\'ecifions six protocoles op\'erationnels consacr\'es respectivement
\`a la capacit\'e d'extraction institutionnelle, \`a l'entropie contextuelle,
\`a l'opposabilit\'e sur syst\`emes d\'eploy\'es, \`a l'estimation conjointe de la
capacit\'e et de la granularit\'e, \`a la capacit\'e d'extraction quantique et au
coefficient de perte d'interpr\'etabilit\'e quantique. Chacun est assorti d'une
d\'efinition op\'erationnelle, d'un estimateur explicite, d'un plan
d'\'echantillonnage et de crit\`eres de fiabilit\'e inter-\'evaluateurs. Nous d\'erivons
ensuite, par une premi\`ere voie, les valeurs d'opposabilit\'e de trois
syst\`emes de r\'ef\'erence, puis nous les confrontons aux param\`etres
institutionnels estim\'es par une seconde voie ind\'ependante. La
convergence des deux voies \'etablit la coh\'erence interne du cadre sans
en d\'emontrer la validit\'e empirique, puisque le programme n'a pas \'et\'e d\'eploy\'e dans le
cadre de cette th\`ese. Les seuils \`a partir desquels il invaliderait
le mod\`ele sont en revanche explicitement \'enonc\'es.

\textbf{Statut~:} programme sp\'ecifi\'e, d\'eploiement \`a venir.
\end{sloppypar}

% ----------------------------------------------------------------------------
\subsection{Objectifs scientifiques et programme}
\label{subsec:objectifs_programme}
% ----------------------------------------------------------------------------
\begin{sloppypar}

\subsubsection{Coh\'erence du programme~: le cycle \textbf{F}$\to$\textbf{B}$\to$\textbf{E}$\to$\textbf{S}$\to$\textbf{C}$\to$\textbf{X}}
\label{subsubsec:cycle_FBESCX}

Les six contributions sont interd\'ependantes et s'encha\^inent selon le
cycle logique
\begin{equation}
\label{eq:cycle_recherche}
\resizebox{0.95\textwidth}{!}{$ \underbrace{\text{Obj~\ref{obj:obj1_unification} (CLO)}}_{\mathbf{F}}
\xrightarrow{\text{d\'efinitions}}
\underbrace{\text{Obj~\ref{obj:obj2_impossibilite} (CRO)}}_{\mathbf{B}}
\xrightarrow{\text{borne classique}}
\underbrace{\text{Obj~\ref{obj:obj3_postquantique} (PQ-CRO)}}_{\mathbf{E}}
\xrightarrow{\text{extension}}
\underbrace{\text{Obj~\ref{obj:obj4_saturation} (SOTP)}}_{\mathbf{S}}
\xrightarrow{\text{saturation}}
\underbrace{\text{Obj~\ref{obj:obj5_verification} (Lean4)}}_{\mathbf{C}}
\xrightarrow{\text{certification}}
\underbrace{\text{Obj~\ref{obj:obj6_experimentation} (Exp\'erimentation)}}_{\mathbf{X}}
\xrightarrow{\text{mesure}}
\text{Obj~\ref{obj:obj2_impossibilite}}
 $}
\end{equation}
Chaque d\'ependance est fonctionnelle. Les d\'efinitions formelles de
$\Conf$, $\Rel$, $\Opp$ (Obj~\ref{obj:obj1_unification}) sont n\'ecessaires pour \'enoncer
le Trilemme~CRO (Obj~\ref{obj:obj2_impossibilite}), car sans op\'erationnalisation
de l'opposabilit\'e au sens de la th\'eorie de l'information, le trilemme resterait une
conjecture informelle~\cite{ClarksonSchneider10}. La borne d'impossibilit\'e classique (Obj~\ref{obj:obj2_impossibilite}) fournit
le point de d\'epart que l'extension post-quantique (Obj~\ref{obj:obj3_postquantique}) g\'en\'eralise avec
les param\`etres $\Delta_q$ et $\eta_q$~\cite{NielsenChuang2010,Wilde2017}. Les bornes classiques (Obj~\ref{obj:obj2_impossibilite})
et quantiques (Obj~\ref{obj:obj3_postquantique}) fournissent les cibles asymptotiques que les
constructions SOTP (Obj~\ref{obj:obj4_saturation}) visent. Sans ces bornes, SOTP ne serait qu'une
collection de protocoles ad~hoc~\cite{JohnsonLindenstrauss84,Groth16}. La
question de savoir si la borne~CRO est asymptotiquement serr\'ee reste
ouverte. Elle est formul\'ee comme conjecture au chapitre~\ref{chap:constructions}. La complexit\'e de la cha\^ine de
raisonnement justifie l'investissement dans la v\'erification formelle (Obj~\ref{obj:obj5_verification}),
car la certification m\'ecanis\'ee s'impose pour des th\'eor\`emes aux implications
soci\'etales directes~\cite{Gonthier08,Wiedijk2007}. Enfin, la formalisation Lean4~(Obj~\ref{obj:obj5_verification}) alimente
Obj~\ref{obj:obj1_unification} par r\'etroaction, puisqu'elle a impos\'e d'expliciter l'Hypoth\`ese de r\'egime
et de pr\'eciser trois \'enonc\'es publi\'es, ce qui a conduit \`a une forme
raffin\'ee du trilemme~\cite{Zellagui2018}.

\subsubsection{Correspondance objectifs-questions-contributions}
\label{subsubsec:correspondance_objectifs}

\begin{table}[htbp]
\centering
\caption{Correspondance entre objectifs scientifiques, questions de
  recherche et contributions}
\label{tab:correspondance_objectifs}
\small
\begin{tabular}{p{3.2cm}cp{7.5cm}}
\toprule
\textbf{Objectif} & \textbf{Q} &
\textbf{Contributions principales} \\
\midrule
Unification th\'eorique & Obj~\ref{obj:obj1_unification} &
  Cadre CLO~: $\Conf$, $\Rel$, $\Opp$~;
  Goulot contextuel~; S-OTP \\
\midrule
Bornes classiques & Obj~\ref{obj:obj2_impossibilite} &
  Trilemme CRO~: $\Delta, \alpha, \beta$~;
  Borne $\leq \Delta/(\gamma\lambda) + O(\lambda^{-2})$ \\
\midrule
Extension quantique & Obj~\ref{obj:obj3_postquantique} &
  Trilemme PQ-CRO~: $\Delta_q, \eta_q$~;
  Analyse NIST PQC \\
\midrule
Constructions de r\'ef\'erence & Obj~\ref{obj:obj4_saturation} &
  SOTP-Low/Base/High~;
  optimalit\'e (conjecture) \\
\midrule
V\'erification formelle & Obj~\ref{obj:obj5_verification} &
  Formalisation Lean4~;
  hypoth\`eses explicit\'ees~;
  forme raffin\'ee certifi\'ee \\
\midrule
Exp\'erimentation & Obj~\ref{obj:obj6_experimentation} &
  Six protocoles de mesure~; crit\`eres de r\'efutation~;
  calibration illustrative \\
\bottomrule
\end{tabular}
\end{table}

Le statut et la r\'ef\'erence de publication de chaque contribution figurent au
tableau~\ref{tab:contributions_these}.

\subsubsection{Impacts attendus}
\label{subsubsec:impacts_attendus}

\paragraph{Impacts scientifiques.}
\label{par:impact_scientifique}

La cryptographie institutionnelle se constitue en discipline
au croisement de la th\'eorie de l'information classique et
quantique, de la cryptographie appliqu\'ee, des sciences cognitives,
du droit compar\'e et de la v\'erification formelle~\cite{Shannon48,NielsenChuang2010,Simon82,Goldreich01,moura2015lean}. Chacune de ces composantes y est
n\'ecessaire sans qu'aucune suffise \`a elle seule.

Le cadre CLO-CRO-PQ-CRO-SOTP d\'egage une m\'ethodologie transf\'erable.
Celle-ci consiste \`a d\'efinir des m\'etriques issues de la th\'eorie de l'information pour les propri\'et\'es
d'int\'er\^et puis \`a \'etablir des bornes d'impossibilit\'e au moyen de Pinsker, de Fano et de la DPI.
Elle se poursuit par l'extension aux contextes quantiques, la construction de protocoles de r\'ef\'erence
visant les bornes et la v\'erification formelle des preuves complexes~\cite{Yao82,GMW87,Asharov12,Gentry09}. Une s\'equence de ce type
a d\'ej\`a \'et\'e explor\'ee pour le calcul multi-parties~; elle pourrait s'appliquer
au chiffrement homomorphe, \`a la blockchain ainsi qu'\`a d'autres domaines de
la cryptographie post-quantique~\cite{BrakerskiVaikuntanathan11,Nakamoto08,Zohar15}.

\paragraph{Impacts pratiques.}
\label{par:impact_pratique}

Le cadre trouve des applications imm\'ediates dans trois domaines.

Les autorit\'es \'electorales peuvent s'appuyer sur le cadre CLO pour choisir
des protocoles adapt\'es \`a leurs contraintes institutionnelles~\cite{Adida08,Lindeman12,Benaloh94}. L'incident
IACR~2025~\cite{IACR2025incident} montre qu'un protocole math\'ematiquement
sain peut \'echouer sur sa couche institutionnelle. La configuration
exigeait en effet les parts de cl\'e des trois garants (partage 3-sur-3)~; la perte
d'une seule part a donc impos\'e d'annuler le scrutin, avant qu'un seuil
2-sur-3 soit adopt\'e. Ce d\'efaut de conception op\'erationnelle, qui touche \`a la gestion des
cl\'es, rel\`eve de ce que le cadre CLO invite \`a expliciter sans constituer une mesure de $\Opp$.

Le R\`eglement G\'en\'eral sur la Protection des Donn\'ees~\cite{GDPR2016,CCPA2018} exige
simultan\'ement la protection des donn\'ees personnelles (Article~5) et la
responsabilit\'e d\'emontrable (Article~5(2)), exigences qui correspondent
directement \`a $\Conf$ et $\Opp$. Notre cadre permet de quantifier ces
compromis. Ainsi, atteindre $\Conf = 0{,}95$ contraint $\Opp$ selon la capacit\'e
$\Delta$, for\c{c}ant un choix explicite entre ces objectifs l\'egislatifs
potentiellement contradictoires~\cite{FinckMoscon19}. La loi camerounaise N\textsuperscript{o}2024-017 du 23 d\'ecembre 2024 relative \`a la protection des donn\'ees \`a caract\`ere personnel, qui s'inspire largement du RGPD, porte les m\^emes tensions.

La transition post-quantique b\'en\'eficie du Trilemme PQ-CRO qui permet
d'\'evaluer les standards NIST PQC sur leur ad\'equation aux capacit\'es
institutionnelles ($\Delta_q$) et aux pertes d'interpr\'etabilit\'e
quantiques ($\eta_q$), en plus de leur s\'ecurit\'e computationnelle~\cite{NISTPQC,Kyber2022,Dilithium2022,SPHINCS2022,Falcon2022}.

\paragraph{Impacts soci\'etaux et souverainet\'e num\'erique.}
\label{par:impact_societal}

Le paradoxe de l'authenticit\'e invisible pose un d\'efi de fond \`a
la d\'emocratie num\'erique, puisqu'il oblige \`a se demander comment des citoyens peuvent contr\^oler des
syst\`emes qu'ils ne comprennent pas~\cite{Citron14,Grimmelmann15}. Calibrer l'opposabilit\'e $\Opp$ selon la capacit\'e institutionnelle $\Delta$ fournit une r\'eponse
quantifiable universelle, selon laquelle les syst\`emes restent sous contr\^ole d\'emocratique
effectif \`a condition que $\Opp$ reste au-dessus du seuil impos\'e par $\Delta$.

Dans les pays du \textbf{Sud Global}, l'adoption de syst\`emes cryptographiques
con\c{c}us ailleurs sans \'evaluation critique de leur ad\'equation aux capacit\'es
institutionnelles locales constitue un enjeu de souverainet\'e
majeur~\cite{Borija2023Exploring,Agyekum2021Ghana}. L'\textbf{Afrique},
en particulier, fait face \`a des dynamiques de transfert technologique
asym\'etriques o\`u les protocoles d\'epassent souvent les capacit\'es
d'audit et de contr\^ole locaux.

Cette th\`ese fournit aux d\'ecideurs d'ici et d'ailleurs les outils pour
\'evaluer objectivement les technologies propos\'ees ($\Conf$, $\Rel$, $\Opp$),
identifier les d\'esalignements capacit\'e-protocole et exiger des
adaptations ou refuser des d\'eploiements inadapt\'es. La mod\'elisation de
 $\Delta$ comme param\`etre mesurable et am\'eliorable ouvre une approche
encadr\'ee pour le renforcement de capacit\'es. Des formations sp\'ecialis\'ees
pour magistrats, auditeurs et r\'egulateurs peuvent alors viser une cible quantifiable
(\guillemotleft~augmenter $\Delta$ de 20 bits~\guillemotright) plut\^ot qu'un objectif vague
(\guillemotleft~sensibiliser les juges~\guillemotright)~\cite{Argyris1978}.
\end{sloppypar}

% ----------------------------------------------------------------------------
%\subsection{Synth\`ese}
%\label{sec:synthese_chap3}
% ----------------------------------------------------------------------------

\begin{boiteArbitre}
L'arbitre pose son stylo. Ce chapitre a accompli une chose
inhabituelle \`a ses yeux, puisqu'il a converti en programme de recherche
formalis\'e ce qu'il ressent depuis des ann\'ees. Le foss\'e qu'il constatait entre validit\'e technique et
conviction institutionnelle porte d\'esormais un nom pr\'ecis, le paradoxe de
l'authenticit\'e invisible. Il dispose aussi d'une mesure, $\Opp = I(V;\Sigma|C)/H(V|C)$,
ainsi que d'une borne quantitative.

Il retient sans effort les six questions qui en d\'ecoulent, car elles suivent
l'ordre naturel de toute enqu\^ete s\'erieuse. Il s'agit d'abord de nommer les choses
(Obj~\ref{obj:obj1_unification}, le cadre CLO) puis d'en chercher les limites
(Obj~\ref{obj:obj2_impossibilite}, le Trilemme CRO). Ces limites doivent ensuite r\'esister \`a l'examen
quantique (Obj~\ref{obj:obj3_postquantique}) avant que l'on construise ce qui s'en approche
(Obj~\ref{obj:obj4_saturation}, la famille SOTP). Il reste \`a certifier que le raisonnement
tient (Obj~\ref{obj:obj5_verification}, Lean4) puis \`a dire comment le soumettre \`a la mesure
(Obj~\ref{obj:obj6_experimentation}, les protocoles exp\'erimentaux).

Il est surtout sensible \`a la coh\'erence du cycle, dans lequel chaque r\'eponse appelle
la suivante. Les d\'efinitions permettent d'\'enoncer les bornes
classiques, qui fournissent \`a leur tour leur cible aux constructions SOTP. Ces constructions
mettent \`a l'\'epreuve la finesse de la borne, qui reste une conjecture, tandis que la formalisation met au jour ce que
les preuves papier laissaient dans l'ombre, \`a savoir une hypoth\`ese implicite
et un \'ecart entre deux ordres de grandeur.

Reste \`a v\'erifier, au chapitre~\ref{chap:cadre_formel}, si les
th\'eor\`emes tiennent ce que les questions promettent.
\end{boiteArbitre}
% ============================================================================
% CONCLUSION DU CHAPITRE 1 
% source : partie1/conclusion_partie1.tex
% ============================================================================

\section{CONCLUSION DU CHAPITRE~\ref{chap:contexte}}
\label{sec:conclusion_contexte}

\begin{sloppypar}
Le chapitre~\ref{chap:contexte} a pos\'e progressivement les fondements sur lesquels
s'appuie l'ensemble de la th\`ese. Il en ressort trois acquis.

Le premier acquis, de nature illustrative, porte sur trois syst\`emes cryptographiques d\'eploy\'es,
\`a savoir Helios pour le vote \'electronique, les zk-SNARKs pour la conformit\'e
financi\`ere et ECDSA pour la signature contractuelle. Selon des
estimations param\'etriques sur l'\'echelle op\'erationnelle CLO, leur produit
$\Conf\cdot\Relver\cdot\Opp$ est compris entre $0{,}273$ et $0{,}572$. Ce motif
laisse entrevoir, sans la d\'emontrer, une tension structurelle entre deux paradigmes
de v\'erification.
Le paradigme cryptographique mesure la r\'esistance aux attaques, alors que le
paradigme institutionnel requiert l'intelligibilit\'e par des non-experts,
des traces exploitables pour un audit r\'etrospectif ainsi qu'une responsabilit\'e
attribuable en cas de d\'efaillance. Aucune m\'etrique formelle n'avait encore
\'etabli les limites de ce compromis.

Le deuxi\`eme acquis est conceptuel. L'opposabilit\'e
$\Opp = I(V;\Sigma|C)/H(V|C)$ entre comme troisi\`eme dimension
cryptographique aux c\^ot\'es de la confidentialit\'e et de la
fiabilit\'e. Deux arch\'etypes s'en d\'egagent, les syst\`emes \`a
haute confidentialit\'e (Helios, zk-SNARKs, produits $0{,}437$ et
$0{,}572$) s'opposant aux syst\`emes \`a haute opposabilit\'e (ECDSA,
produit $0{,}273$). L'\'equilibre hypoth\'etique
$\Conf \approx \Rel \approx \Opp \approx 0{,}8$ n'est approch\'e par aucune
des trois estimations. Cette d\'efinition transforme
l'interpr\'etabilit\'e institutionnelle d'une contrainte externe
subie en une propri\'et\'e formelle que l'on peut mesurer et optimiser.

Le troisi\`eme est programmatique. La cryptographie th\'eorique
formalise $\Conf$ et $\Rel$ via des jeux de s\'ecurit\'e mais ignore
$\Opp$~; la s\'ecurit\'e utilisable analyse l'utilisabilit\'e
qualitativement sans formalisation math\'ematique~; l'informatique
juridique raisonne sur l'admissibilit\'e de la preuve num\'erique
sans m\'etriques quantitatives. Les six objectifs de recherche
Obj~\ref{obj:obj1_unification} \`a Obj~\ref{obj:obj6_experimentation} d\'eclinent
le paradoxe de l'authenticit\'e invisible en probl\`emes formels
pr\'ecis, chacun associ\'e \`a une contribution et \`a une \'etape
de la progression canonique
\textbf{F}$\to$\textbf{B}$\to$\textbf{E}$\to$\textbf{S}$\to$\textbf{C}$\to$\textbf{X},
comblant ainsi le cloisonnement entre ces trois corpus.

Le chapitre~\ref{chap:cadre_formel} prend le relais. Il dispose d\'esormais
d'un vocabulaire fix\'e, d'une probl\'ematique formul\'ee avec pr\'ecision
et d'exemples calibr\'es. La mesure de $\Conf$, $\Rel$ et $\Opp$ remplace le d\'ebat qualitatif entre s\'ecurit\'e et transparence et transforme des choix de conception en
probl\`emes d'optimisation sous contraintes, susceptibles de r\'eponses
pr\'ecises. Il peut d\'esormais engager les six \'etapes du programme,
de la formalisation \`a l'exp\'erimentation.
\end{sloppypar}



% ==========================================================
% CHAPITRE 2 : CADRE FORMEL ET BORNES D'IMPOSSIBILITE
% ==========================================================
\chapter[CADRE FORMEL ET BORNES D'IMPOSSIBILIT\'E POUR L'OPPOSABILIT\'E L\'EGALE]%
        {CADRE FORMEL ET BORNES D'IMPOSSIBILIT\'E POUR L'OPPOSABILIT\'E L\'EGALE}
\label{chap:cadre_formel}
\debutChapitre%
  {Cadre formel et bornes d'impossibilit\'e pour l'opposabilit\'e l\'egale}%
  {%
    Ce chapitre pose le socle th\'eorique de la th\`ese en accomplissant les
    trois premiers actes du cycle canonique (formaliser, borner, \'etendre).
    Le cadre CLO y est d'abord d\'evelopp\'e afin de formaliser les trois
    dimensions concurrentes de la cryptographie institutionnelle que sont la
    confidentialit\'e, la fiabilit\'e et l'opposabilit\'e. Vient ensuite le
    Trilemme CRO conditionnel classique. Sous quatre hypoth\`eses explicites
    portant sur les limites de l'institution, il borne le produit des trois
    m\'etriques par une quantit\'e dont le terme dominant est le rapport entre
    la capacit\'e institutionnelle et le param\`etre de s\'ecurit\'e, pond\'er\'e par
    la richesse des verdicts. Sa forme raffin\'ee identifie le r\'egime
    o\`u cette capacit\'e est la contrainte liante. Le chapitre \'etend enfin ce
    r\'esultat au r\'egime post-quantique. La borne y garde la m\^eme structure,
    tandis qu'une perte d'interpr\'etabilit\'e $\eta_q$, nulle pour les preuves
    classiques, intervient dans le sc\'enario d'un support quantique.%
  }%
  {%
    cadre CLO, opposabilit\'e juridique, trilemme CRO, trilemme CRO
    post-quantique, bornes d'impossibilit\'e conditionnelles%
  }%

% ============================================================================
% CHAPITRE 2 : INTRODUCTION
% ============================================================================

% ----------------------------------------------------------------------------
\section{INTRODUCTION}
\label{sec:introduction_cadre_formel}
% ----------------------------------------------------------------------------

\begin{sloppypar}
La th\`ese progresse selon un cycle canonique en six actes,
\textbf{F}$\to$\textbf{B}$\to$\textbf{E}$\to$\textbf{S}$\to$\textbf{C}$\to$\textbf{X}.
Ce chapitre en r\'ealise les trois premiers. La section~\ref{sec:clo}
formalise le cadre CLO et d\'efinit les trois m\'etriques
($\Conf$, $\Rel$, $\Opp$) qui rendent le paradoxe tractable, en y
int\'egrant les jeux cryptographiques (IND-CONF, BIND-REL, EXT-OPP) qui
captent leur contenu op\'erationnel. La section~\ref{sec:cro_classique}
borne les compromis r\'ealisables dans le r\'egime classique par
le Trilemme CRO conditionnel. La section~\ref{sec:cro_pq} \'etend
ces bornes aux adversaires quantiques et aux canaux post-quantiques. La
section~\ref{sec:conclusion_cadre_formel} synth\'etise ces trois
contributions. Les trois actes suivants, saturer, certifier et
exp\'erimenter, font l'objet des chapitres~\ref{chap:constructions}
et~\ref{chap:experimentation}.

Ce chapitre \'etablit le socle th\'eorique de la th\`ese. Chaque section
correspond \`a une contribution publi\'ee et apporte une r\'eponse pr\'ecise
\`a l'un des objectifs de recherche formul\'es \`a la
sous-section~\ref{subsec:objectifs_recherche}.

% ----------------------------------------------------------------------------
\subsection{D\'emarche m\'ethodologique~: les actes
\textbf{F}, \textbf{B} et \textbf{E}}
\label{subsec:demarche_methodologique}
% ----------------------------------------------------------------------------

Avant d'entrer dans les sections techniques, nous exposons la logique de
chaque \'etape. Dans le cycle
\textbf{F}$\to$\textbf{B}$\to$\textbf{E}$\to$\textbf{S}$\to$\textbf{C}$\to$\textbf{X},
chaque \'etape r\'epond \`a une n\'ecessit\'e m\'ethodologique sp\'ecifique.

\paragraph{Formaliser (\textbf{F}).}
La cryptographie institutionnelle ne disposait d'aucun cadre m\'etrique
pour l'interpr\'etabilit\'e. Avant de borner, il fallait d\'efinir ce
qui doit \^etre born\'e. La m\'ethodologie adopt\'ee est
axiomatique-constructive. Elle part de la th\'eorie de l'information
(Shannon, Fano, Pinsker) pour d\'efinir des m\'etriques mesurables,
comparables et formellement manipulables. La formalisation pr\'ec\`ede
les bornes parce que sans d\'efinition pr\'ecise de l'opposabilit\'e,
toute borne serait circulaire.

\paragraph{Borner (\textbf{B}).}
Une fois le cadre CLO \'etabli, il reste \`a d\'eterminer les limites
fondamentales du r\'ealisable. La m\'ethodologie des
th\'eor\`emes d'impossibilit\'e conditionnelle, inspir\'ee des trilemmes
CAP et FLP, conditionne les bornes aux param\`etres cognitifs r\'eels des
institutions ($\Delta$, $\alpha$, $\beta$). Cette approche produit une
contrainte param\'etr\'ee par les acteurs institutionnels plut\^ot qu'un
obstacle th\'eorique absolu. Lorsque la granularit\'e des
verdicts d\'epasse la capacit\'e d'extraction institutionnelle $\Delta$, c'est
$\Delta$ qui plafonne les compromis r\'ealisables~; sinon, c'est l'in\'egalit\'e
de Fano.

\paragraph{\'Etendre (\textbf{E}).}
L'extension quantique introduit une perturbation qualitativement
diff\'erente, qui tient au bruit de canal quantique et \`a la perte
d'interpr\'etabilit\'e qu'il induit lorsque la preuve est port\'ee par un support quantique. Les outils de la th\'eorie
de l'information quantique, entropie de von Neumann, borne de Holevo (sous-sous-section~\ref{subsubsec:holevo}),
lemme de Gentle Measurement, permettent d'\'etendre les bornes \`a cet
environnement sans repartir de z\'ero. La borne explicite conserve sa forme,
avec une capacit\'e quantique $\Delta_q$~; le coefficient de perte
d'interpr\'etabilit\'e $\eta_q$ n'appara\^it que dans une majoration
affaiblie. Ce coefficient est nul pour les preuves qui sont des cha\^ines de bits.
\end{sloppypar}

% ============================================================================
% SECTION 2.2 : LE CADRE CLO — FORMALISATION DES MESURES
% Source : partie2/clo.tex (version condensée, niveau de détail maximal)
% Preuves complètes intégrées dans le corps (pas de renvoi aux annexes)
% ============================================================================

\section{LE CADRE CLO : FORMALISATION DES MESURES}
\label{sec:clo}
\begin{sloppypar}
Le chapitre~\ref{chap:contexte} a expos\'e le paradoxe de l'authenticit\'e
invisible~\cite{Goodman15,Casey16,Grimmelmann15}, selon lequel les protocoles qui maximisent la protection cryptographique
induisent une opacit\'e institutionnelle s\'ev\`ere, tandis que les protocoles
transparents sacrifient la confidentialit\'e~\cite{Bonneau16,Citron14}. La section~\ref{sec:etat_art}
a fourni les outils n\'ecessaires, th\'eorie de l'information de
Shannon~\cite{Shannon48,Shannon49,Wyner75}, in\'egalit\'es de Pinsker et
Fano~\cite{Fano61,Pinsker1964}, th\'eorie de l'information quantique
(von Neumann, Holevo~\cite{Holevo1973,NielsenChuang2010}), primitives d'obscurcissement (commitments,
zk-SNARKs~\cite{Groth16,BCTV14,Gennaro13}), trilemmes d'impossibilit\'e (CAP~\cite{Brewer00},
FLP~\cite{FLP85}), standards post-quantiques NIST~\cite{NISTPQC,Kyber2022,Dilithium2022}, sans encore les articuler 
en un cadre unifi\'e.

Nous d\'efinissons ici ce cadre. Le \textbf{cadre CLO (Cryptographic
Legal Opposability)}~\cite{MINKA2026CLO} constitue la premi\`ere formalisation
math\'ematique rigoureuse de la cryptographie institutionnelle. L'
Objectif~\ref{obj:obj1_unification}, comment \'etablir un cadre unifi\'e
formalisant confidentialit\'e, fiabilit\'e et opposabilit\'e via des mesures
issues de la th\'eorie de l'information permettant analyses quantitatives et comparaisons
objectives~? re\c{c}oit une r\'eponse en cinq \'etapes, qui portent successivement sur le
mod\`ele de protocole, les d\'efinitions des m\'etriques, les propri\'et\'es
fondamentales, le paradigme de conception et une construction prototype.

Notre apport conceptuel central est l'opposabilit\'e $\Opp$. La
litt\'erature cryptographique classique traite
l'interpr\'etabilit\'e institutionnelle comme contrainte externe~\cite{GMR85,Goldreich01,GoldwasserMicali84} ou lacune
\'educative des utilisateurs~\cite{Whitten99,Adams1999,Herley09}. Le cadre CLO en fait une dimension \`a
optimiser au m\^eme titre que $\Conf$ et $\Rel$.

Cette section correspond \`a l'\'etape \textbf{F (Formaliser)} de la progression \textbf{F}$\to$B$\to$E$\to$S$\to$C$\to$X et apporte la contribution \textbf{C1} de cette th\`ese.
\end{sloppypar}

\begin{boiteArbitre}
L'arbitre dispose d\'esormais d'une m\'etrique pour ce qu'il
ressentait. L'opposabilit\'e $\mathrm{Opp} = I(V;\Sigma|C)/H(V|C)$ mesure
exactement la fraction d'incertitude institutionnelle que la preuve
cryptographique r\'esout. Selon les estimations illustratives du
cadre, Helios lui laisse $50\%$ d'incertitude r\'esiduelle~; ECDSA lui en
laisse $10\%$. Il peut maintenant comparer et se demander ce qui, dans le
protocole ou dans l'institution, limite les compromis accessibles.
\end{boiteArbitre}

% -----------------------------------------------------------------------
\subsection{Mod\`ele de protocole CLO}
\label{subsec:modele_protocole_clo}
% -----------------------------------------------------------------------
\begin{sloppypar}

\subsubsection{D\'efinition formelle}
\label{subsubsec:definition_protocole}

Un protocole d'evidence cryptographique op\'erant en contexte institutionnel
doit satisfaire deux niveaux de v\'erification fondamentalement distincts,
\`a savoir la correction math\'ematique (v\'erification cryptographique) et l'extraction
de signification s\'emantique (interpr\'etation institutionnelle)~\cite{ESIGN00,eIDAS14,Kerr04}. Cette
dualit\'e, absente des d\'efinitions cryptographiques
classiques~\cite{GMR85,Goldreich01,GoldwasserMicali84} motive le mod\`ele formel qui suit.

\begin{definition}[Protocole CLO]\label{def:protocole_clo}
Un \textbf{protocole CLO} $\mathcal{P}$ consiste en quatre espaces et
quatre algorithmes param\'etr\'es par~$\lambda$.

\medskip\noindent\textit{Espaces.}
Le protocole manipule l'espace de messages $\mathcal{M}_\lambda$ (typiquement $\{0,1\}^\lambda$),
l'espace de preuves cryptographiques $\Sigma_\lambda$,
l'espace de contextes institutionnels $\mathcal{C}_\lambda$
et l'espace de verdicts $\mathcal{V}_\lambda$ (d\'ecisions institutionnelles).

\medskip\noindent\textit{Algorithmes.}
\begin{itemize}
\item $\mathsf{Setup}(1^\lambda) \to (pk, sk)$ : g\'en\'eration de cl\'es publique/secr\`ete
\item $\mathsf{Prove}(sk, m, c) \to \sigma$ : g\'en\'eration de preuve pour message $m$ et contexte $c$
\item $\mathsf{Verify}(pk, \sigma, c) \to \{0,1\}$ : v\'erification cryptographique de la preuve
\item $\mathsf{Interpret}(\sigma, c) \to v$ : \textbf{interpr\'etation institutionnelle} produisant verdict $v \in \mathcal{V}$
\end{itemize}

\medskip\noindent o\`u $c$ repr\'esente la distribution du contexte
institutionnel capturant l'environnement d\'ecisionnel (pr\'ec\'edents l\'egaux,
expertise disponible, proc\'edures \'etablies).
\end{definition}

\begin{remark}[Double couche de v\'erification]\label{rem:double_couche}
$\mathsf{Verify}$ r\'epond \`a la question \emph{\guillemotleft~La preuve est-elle math\'ematiquement
valide~?\guillemotright} $\mathsf{Interpret}$ r\'epond \`a la question \emph{\guillemotleft~Que signifie cette
preuve pour ma d\'ecision~?\guillemotright} La tension centrale de la cryptographie
institutionnelle \'emerge pr\'ecis\'ement de cette dualit\'e, puisqu'une preuve peut
\^etre cryptographiquement impeccable ($\mathsf{Verify}(\sigma) = 1$) mais
institutionnellement opaque, l'interpr\`ete \'etant incapable d'en extraire
$v$ de mani\`ere fiable.
\end{remark}

\begin{figure}[htbp]
\centering
\begin{tikzpicture}[
    scale=0.85,
    transform shape,
    box/.style={rectangle, draw=enspblue, fill=blue!8, thick,
                minimum width=2.8cm, minimum height=1cm,
                align=center, font=\small},
    cryptobox/.style={rectangle, draw=enspgreen, fill=enspgreen!8, thick,
                minimum width=2.8cm, minimum height=1cm,
                align=center, font=\small},
    instbox/.style={rectangle, draw=navyblue, fill=navyblue!8, thick,
                minimum width=2.8cm, minimum height=1cm,
                align=center, font=\small},
    dataarrow/.style={-{Stealth[length=3mm]}, thick, enspblue},
    interpretarrow/.style={-{Stealth[length=3mm]}, thick, navyblue, dashed}
]

% ---------------------------------------------------------------
% N\OE UDS
% ---------------------------------------------------------------

% Ligne 1 : entr\'ees
\node[box] (message) at (0, 6)   {\textbf{Message} $M$};
\node[box] (context) at (6, 6)   {\textbf{Contexte} $C$};

% Ligne 2 : Prouveur
\node[cryptobox, minimum width=4cm] (prover) at (3, 4)
    {\textbf{Prouveur}\\[2pt]
     $\mathsf{Prouver}(sk, m, c) \to \sigma$};

% Ligne 3 : Preuve
\node[cryptobox] (proof) at (3, 2)
    {\textbf{Preuve} $\Sigma$};

% Ligne 4 gauche : V\'erifier
\node[cryptobox] (verify) at (0, 0)
    {\textbf{V\'erifier}\\[2pt]
     $\mathsf{V\'erifier}(pk, \sigma, c)$};

% Ligne 5 gauche : r\'esultat binaire
\node[box] (vresult) at (0, -1.8)
    {$\{0,1\}$};

% Ligne 4 droite : Interpr\'eter
\node[instbox] (interpret) at (7, 0)
    {\textbf{Interpr\'eter}\\[2pt]
     $\mathsf{Interpr\'eter}(\sigma, c)$};

% Ligne 5 droite : Verdict
\node[instbox] (verdict) at (7, -1.8)
    {\textbf{Verdict} $V$};

% ---------------------------------------------------------------
% FL\`ECHES COUCHE CRYPTOGRAPHIQUE
% ---------------------------------------------------------------
\draw[dataarrow] (message.south) -- (prover.north west);
\draw[dataarrow] (context.south) -- (prover.north east);
\draw[dataarrow] (prover.south)  -- (proof.north);
\draw[dataarrow] (proof.south)   -- (verify.north);
\draw[dataarrow] (context.west)  -| (verify.north);
\draw[-{Stealth[length=3mm]}, thick]
                 (verify.south)  -- (vresult.north);

% ---------------------------------------------------------------
% FL\`ECHES COUCHE INSTITUTIONNELLE
% ---------------------------------------------------------------
\draw[interpretarrow] (proof.east)    -| (interpret.north);
\draw[interpretarrow] (context.east)  -| (interpret.north);
\draw[interpretarrow] (interpret.south) -- (verdict.north);

% ---------------------------------------------------------------
% \'ETIQUETTES FLOTTANTES (sans fit, pour \'eviter les chevauchements)
% ---------------------------------------------------------------
\node[font=\footnotesize, text=enspblue]  at (-1.2, 5.4)
    {Entr\'ee};
\node[font=\footnotesize, text=enspgreen] at (-1.4, 3.0)
    {Cryptographique};
\node[font=\footnotesize, text=enspgreen] at (1.5,  2.5)
    {Sortie};
\node[font=\footnotesize, text=navyblue]  at (8.6,  1.0)
    {Institutionnel};

% ---------------------------------------------------------------
% CADRES DE COUCHES --- dessin\'es MANUELLEMENT pour contr\^ole pr\'ecis
% ---------------------------------------------------------------

% Cadre couche cryptographique (englobe prover, proof, verify, vresult)
\draw[enspgreen, dashed, thick, rounded corners=4pt]
    (-1.6, -2.5) rectangle (5.0, 4.8);
\node[font=\small\bfseries, text=enspgreen] at (1.7, 4.5)
    {Couche cryptographique};

% Cadre couche institutionnelle (englobe interpret, verdict)
\draw[navyblue, dashed, thick, rounded corners=4pt]
    (5.4, -2.5) rectangle (9.0, 0.8);
\node[font=\small\bfseries, text=navyblue] at (7.2, 0.55)
    {Couche institutionnelle};

\end{tikzpicture}
\caption[Mod\`ele de Protocole CLO : architecture de v\'erification \`a deux couches]{%
  Mod\`ele de Protocole CLO : architecture de v\'erification \`a deux couches.%
}\label{fig:clo_model}

\smallskip
\begin{minipage}{\linewidth}
  \small\textit{Note :} La couche cryptographique (verte) assure la correction
  math\'ematique via les algorithmes $\mathsf{Prouver}$ et $\mathsf{V\'erifier}$.
  La couche institutionnelle (bleue marine) fournit l'interpr\'etation s\'emantique
  via la fonction $\mathsf{Interpr\'eter}$, qui projette les preuves cryptographiques
  sur des verdicts juridiquement significatifs en contexte institutionnel.
\end{minipage}

\end{figure}
\begin{boiteArbitre}
L'arbitre comprend que l'algorithme $\mathsf{Interpret}$ est le sien, car les trois
autres prot\`egent le message alors que celui-ci le sert, lui.
\end{boiteArbitre}

\subsubsection{Exemples d'instanciation}
\label{subsubsec:exemples_instanciation}

Les exemples qui suivent ont valeur illustrative~; les m\'etriques CLO
num\'eriques qui leur seront associ\'ees dans la suite sont des estimations
param\'etriques analytiques d\'eriv\'ees depuis les sp\'ecifications publi\'ees,
non des mesures empiriques contr\^ol\'ees.

\paragraph{Vote \'electronique (Helios)~\cite{Adida08,Adida09UCL}.}
Le message $m$ est le choix de vote $\{0,1\}$~; la preuve $\sigma$
consiste en un bulletin chiffr\'e ElGamal accompagn\'e d'une preuve NIZK~(Non-Interactive Zero-Knowledge)
de bien-formation~\cite{Chaum81,CramerDamgardSchoenmakers}. Le contexte $c$ recouvre le registre d'\'electeurs,
la structure du bulletin et les r\`egles \'electorales, tandis que le
verdict $v$ d\'esigne le d\'ecompte final ou la validation de l'\'election.
$\mathsf{Verify}$ contr\^ole la preuve NIZK et la signature du bulletin~;
$\mathsf{Interpret}$ consiste pour les garants \'electoraux \`a d\'echiffrer
l'agr\'egat homomorphe et en extraire le d\'ecompte~\cite{Goodman15}. Les m\'etriques CLO
estim\'ees (\'echelle op\'erationnelle CLO de~\cite{MINKA2026CLO}, fiabilit\'e
$\Relver$) sont $\Conf = 0{,}92$, $\Relver = 0{,}95$, $\Opp = 0{,}50$, avec
un produit $\Conf\cdot\Relver\cdot\Opp = 0{,}437$.

\paragraph{Conformit\'e financi\`ere (zk-SNARKs)~\cite{Groth16,BCTV14,Gennaro13,Parno13}.}
Ici $m$ encode transactions priv\'ees ou soldes~; $\sigma$ est une preuve
Groth16 de 128 octets certifiant la conformit\'e r\'eglementaire~\cite{BenSasson2019,PLONK20}. Le contexte
$c$ couvre les r\'egulations AML/B\^ale~III~\cite{BaselIII}, l'historique de
transactions et les politiques internes. $\mathsf{Verify}$ contr\^ole les
\'equations de \emph{pairing} de Groth16~\cite{Campanelli2019,Wahby18}~; $\mathsf{Interpret}$ demande au
r\'egulateur de corr\'eler le hachage de contexte avec ses bases de donn\'ees
pour rendre un verdict de conformit\'e~\cite{FATF15}. Les m\'etriques CLO estim\'ees sont
$\Conf = 0{,}89$, $\Relver = 0{,}96$, $\Opp = 0{,}67$, avec un produit
$\Conf\cdot\Relver\cdot\Opp = 0{,}572$.

\paragraph{Signatures contractuelles (ECDSA)~\cite{Johnson01,NIST.FIPS.186-5}.}
Le message $m$ est le document contractuel~; $\sigma = (r, s)$ est la
signature ECDSA. Le contexte $c$ comprend les parties contractantes, la
juridiction et les termes du contrat~; $v$ est la validit\'e du contrat et
l'identification de la partie responsable~\cite{Mason07,Kerr04}. $\mathsf{Verify}$ contr\^ole
l'\'equation $r = x_1 \bmod n$~; $\mathsf{Interpret}$ engage le juge dans
la v\'erification des identit\'es, l'examen des termes et le rendu du jugement.
Les m\'etriques CLO estim\'ees sont $\Conf = 0{,}31$, $\Relver = 0{,}98$, $\Opp = 0{,}90$,
avec un produit $\Conf\cdot\Relver\cdot\Opp = 0{,}273$.

Ces trois instanciations couvrent vote \'electronique, r\'egulation financi\`ere
et litige contractuel~\cite{Nakamoto08,Buterin17}, registre suffisant pour tester la g\'en\'eralit\'e du
mod\`ele. Le tableau~\ref{tab:clo_resume_ch1} r\'ecapitule ces m\'etriques. Ces
syst\`emes ne sont pas des instances du Trilemme CRO
(Remarque~\ref{rem:portee_cas_etude})~; ils servent \`a calibrer les ordres
de grandeur des param\`etres institutionnels.

% -----------------------------------------------------------------------
%\begin{table}[htbp]
%\centering
%\caption{M\'etriques CLO des syst\`emes d\'eploy\'es}
%\label{tab:clo_resume_ch2_clo}
%\begin{tabular}{@{}lcccc@{}}
%\toprule
%\textbf{Syst\`eme} & $\boldsymbol{\Conf}$ & $\boldsymbol{\Rel}$ & $\boldsymbol{\Opp}$ & $\boldsymbol{\Conf\cdot\Rel\cdot\Opp}$ \\
%\midrule
%Helios (vote \'electronique)~\cite{Adida08,Adida09UCL} & 0{,}92 & 0{,}95 & 0{,}50 & \textbf{0{,}437} \\
%zk-SNARKs (conformit\'e)~\cite{Groth16,BCTV14} & 0{,}89 & 0{,}96 & 0{,}67 & \textbf{0{,}572} \\
%ECDSA (contrats)~\cite{Johnson01,NIST.FIPS.186-5} & 0{,}31 & 0{,}98 & 0{,}90 & \textbf{0{,}273} \\
%\bottomrule
%\end{tabular}
%\begin{tablenotes}
%\item \footnotesize Source~: analyses param\'etriques d\'eriv\'ees des sp\'ecifications publi\'ees~\cite{Tsybakov2009}.
%\item \footnotesize Ces valeurs servent de validation th\'eorique, non de mesures empiriques.
%\end{tablenotes}
%\end{table}
\end{sloppypar}

% -----------------------------------------------------------------------
\subsection{Les trois dimensions : d\'efinitions formelles}
\label{subsec:trois_dimensions}
% -----------------------------------------------------------------------
\begin{sloppypar}
Nous d\'efinissons les trois m\'etriques fondamentales via la th\'eorie
de l'information (section~\ref{sec:etat_art},
sous-section~\ref{subsec:fondements_info})~\cite{Shannon48,CoverThomas06,HanVerdu}.

\subsubsection{Confidentialit\'e : indistinguabilit\'e au pire cas}
\label{subsubsec:confidentialite_clo}

La confidentialit\'e formalise l'impossibilit\'e de distinguer les preuves
g\'en\'er\'ees pour des messages distincts, g\'en\'eralisant l'indistinguabilit\'e
s\'emantique~\cite{GoldwasserMicali84}.

\begin{definition}[Confidentialit\'e]\label{def:confidentialite_clo}
Pour un protocole CLO $\mathcal{P}$ :
\begin{equation}
\Conf(\mathcal{P}) = 1 - \max_{\substack{m_0, m_1 \in \mathcal{M}_\lambda \\ c \in \mathcal{C}_\lambda}}
\TV\!\left(P_{\Sigma|M=m_0,C=c},\; P_{\Sigma|M=m_1,C=c}\right)
\label{eq:conf_clo}
\end{equation}
o\`u $\TV(P, Q) = \frac{1}{2} \sum_{\sigma \in \Sigma_\lambda} |P(\sigma) - Q(\sigma)|$
est la distance en variation totale.
\end{definition}

$\Conf = 1$ (confidentialit\'e parfaite au sens de Shannon~\cite{Shannon49,Shannon48})
signifie que les distributions de preuves sont identiques pour tous
messages et tout contexte~; m\^eme un adversaire computationnellement
illimit\'e ne peut distinguer~\cite{Wyner75,WynerZiv76}. \`a l'oppos\'e, $\Conf = 0$ trahit l'existence
de messages $m_0, m_1$ et d'un contexte $c$ pour lesquels les supports
sont disjoints, un observateur passif acquiert alors une certitude.
Le r\'egime pratique est interm\'ediaire.

Du point de vue des jeux de s\'ecurit\'e, notre d\'efinition majore l'avantage
de tout distingueur dans IND-CPA~\cite{GoldwasserMicali84,Bellare97}~:
\begin{equation}
\mathsf{Adv}^{\text{IND-CPA}}_{\mathcal{D}}(\lambda)
  \leq \max_{m_0,m_1,c} \TV\!\left(P_{\Sigma|m_0,c},\; P_{\Sigma|m_1,c}\right)
\label{eq:adv_ind_cpa}
\end{equation}
$\Conf$ fournit donc une borne informationnelle sur la s\'ecurit\'e cryptographique.

\subsubsection{Fiabilit\'e : reconstruction du message}
\label{subsubsec:fiabilite_clo}

\begin{definition}[Fiabilit\'e de reconstruction]\label{def:fiabilite_clo}
Soit $\hat{M} = \mathsf{Dec}(\Sigma, C)$ un estimateur de $M$ calculable \`a
partir de la preuve et du contexte. La fiabilit\'e est
\begin{equation}
\Rel(\mathcal{P}) \;:=\; \Relrec(\mathcal{P})
  = 1 - \min_{\mathsf{Dec}} \Pr\bigl[\mathsf{Dec}(\Sigma, C) \neq M\bigr],
\label{eq:rel_clo}
\end{equation}
la probabilit\'e \'etant prise sur $M$, $C$ et l'al\'ea du protocole
(\'equation~\eqref{eq:rel_def}).
\end{definition}

La fiabilit\'e mesure la capacit\'e de la preuve \`a restituer le message
original~\cite{Fano61,CoverThomas06}. C'est cette seule notion que
contr\^ole l'in\'egalit\'e de Fano. Un \emph{one-time pad} pur atteint
$\Conf = 1$ mais $\Rel = 2^{-\lambda}$ pour $M$ uniforme (pas de
r\'ecup\'eration sans cl\'e)~\cite{Shannon49}~; un engagement de
Pedersen~\cite{Pedersen91,ChaumPedersen92} est parfaitement masquant
($\Conf = 1$, donc $\Relrec = 2^{-\lambda}$) et calculatoirement liant sous
DLog~\cite{Damgard87}~; ECDSA~\cite{Johnson01,NIST.FIPS.186-5}, qui publie le message
en clair, a $\Relrec = 1$ et une confidentialit\'e nulle au sens de la
D\'efinition~\ref{def:confidentialite_clo}.

\begin{remark}[Fiabilit\'e de v\'erification]\label{rem:rel_ver}
La fiabilit\'e \emph{op\'erationnelle} de~\cite{MINKA2026CLO} est la
compl\'etude, assortie de la solidit\'e, de la v\'erification~:
\begin{equation}
\Relver(\mathcal{P}) = \min_{m,c}\Pr\bigl[\mathsf{Verify}(pk,\mathsf{Prove}(sk,m,c),c)=1\bigr].
\label{eq:rel_ver}
\end{equation}
Les deux notions ne co\"incident pas. Pour $M$ uniforme sur $\{0,1\}^\lambda$,
$\Conf = 1$ entra\^ine $\Sigma \perp M \mid C$, donc $\Relrec = 2^{-\lambda}$,
alors que $\Relver$ peut valoir~$1$ (engagement parfaitement masquant). Le
Trilemme (Th\'eor\`eme~\ref{thm:trilemme_cro_conditionnel}) porte sur
$\Relrec$, seule notion \`a laquelle s'applique l'in\'egalit\'e de Fano. Les
m\'etriques CLO des syst\`emes d\'eploy\'es et des constructions SOTP sont
exprim\'ees avec $\Relver$~; elles ne sont donc pas des instances du Trilemme
(Remarque~\ref{rem:portee_cas_etude}).
\end{remark}

\subsubsection{Opposabilit\'e : r\'esolution d'incertitude institutionnelle}
\label{subsubsec:opposabilite_clo}

\begin{definition}[Opposabilit\'e]\label{def:opposabilite_clo}
\begin{equation}
\Opp(\mathcal{P}) = \frac{I(V;\Sigma|C)}{H(V|C)}
\label{eq:opp_clo}
\end{equation}
\end{definition}

L'opposabilit\'e capture la fraction de l'incertitude du verdict
institutionnel $V$ que la preuve $\Sigma$ r\'esout, conditionnellement
au contexte $C$. Absente de la litt\'erature cryptographique
classique, cette m\'etrique quantifie
l'interpr\'etabilit\'e institutionnelle. Ainsi, $\Opp = 1$ signifie que la preuve
l\`eve toute ambigu\"it\'e~; $\Opp = 0$ que la preuve est institutionnellement
muette~\cite{Zohar15,Grimmelmann15}. L'information mutuelle conditionnelle $I(V;\Sigma|C)$ est born\'ee
par la taille de l'espace de verdicts, point que le th\'eor\`eme~\ref{thm:goulot_contexte}
exploitera. La d\'efinition n'est informative que si le verdict n'est pas une
fonction d\'eterministe de $(\Sigma,C)$, car dans ce cas $H(V|\Sigma,C)=0$ et
$\Opp=1$ pour tout protocole (Remarque~\ref{rem:verdict_deterministe}). Le
cadre retient donc des interpr\'etations admissibles au sens de la
D\'efinition~\ref{def:interp_admissible}, o\`u le verdict d\'epend aussi des
ressources propres de l'institution.

\end{sloppypar}

% -----------------------------------------------------------------------
\subsection{Propri\'et\'es fondamentales}
\label{subsec:proprietes_fondamentales}
% -----------------------------------------------------------------------
\begin{sloppypar}
Nous \'etablissons quatre propri\'et\'es structurant les m\'etriques CLO.
Les deux premi\`eres sont des bornes \'el\'ementaires~; les deux suivantes
r\'ev\`elent des tensions intrins\`eques.

\begin{theorem}[Bornes \'el\'ementaires des m\'etriques CLO]\label{thm:bornes_mesures}
Pour tout protocole CLO $\mathcal{P}$~:
\begin{equation}
  0 \leq \Conf(\mathcal{P}),\; \Rel(\mathcal{P}),\; \Opp(\mathcal{P}) \leq 1
  \label{eq:bornes_elementaires}
\end{equation}
\end{theorem}

\begin{Preuve}[Th\'eor\`eme~\ref{thm:bornes_mesures}~: Bornes \'el\'ementaires]\label{prv:bornes_mesures}

\begin{cmtEtape}{\'Etape unique. Montrer que chaque m\'etrique est dans $[0,1]$.}
Les trois m\'etriques sont d\'efinies comme des rapports ou des
compl\'ements de distances. On v\'erifie s\'epar\'ement que chacune
reste dans l'intervalle $[0,1]$.
\end{cmtEtape}

\textbf{Confidentialit\'e.}
Par d\'efinition, $\Conf(\mathcal{P}) = 1 - \max_{m_0,m_1}
\TV(P_{\Sigma|M=m_0}, P_{\Sigma|M=m_1})$.
\cmt{La confidentialit\'e est l'oppos\'e de la meilleure distinguabilit\'e.}

Puisque $\TV(P,Q) \in [0,1]$ pour toutes distributions $P$ et $Q$,
\cmt{La distance en variation totale est comprise entre 0 et 1.}

on a $0 \leq 1 - \max \TV(\cdot) \leq 1$,
\cmt{Si le max vaut 0, $\Conf = 1$. Si le max vaut 1, $\Conf = 0$.}

d'o\`u $\Conf \in [0,1]$.
\cmt{Conclusion pour la confidentialit\'e.}

\textbf{Fiabilit\'e.}
Par d\'efinition, $\Rel(\mathcal{P}) = 1 - \min_{\mathsf{Dec}} \Pr[\mathsf{Dec}(\Sigma,C) \neq M]$.
\cmt{La fiabilit\'e est la meilleure probabilit\'e de reconstruction du message.}

Une probabilit\'e d'erreur est dans $[0,1]$, d'o\`u $\Rel \in [0,1]$~; il en va
de m\^eme pour $\Relver$ (Remarque~\ref{rem:rel_ver}).
\cmt{Une probabilit\'e est toujours dans $[0,1]$.}

\textbf{Opposabilit\'e.}
Par d\'efinition, $\Opp(\mathcal{P}) = I(V;\Sigma|C)/H(V|C)$.
\cmt{Rapport de l'information extraite sur l'information disponible.}

Par les propri\'et\'es fondamentales de l'information mutuelle
conditionnelle~\cite{CoverThomas06},
\begin{equation}
0 \leq I(V;\Sigma|C) \leq H(V|C)
\end{equation}
\cmtFormule{L'information mutuelle est toujours born\'ee par les entropies marginales.}

ce qui donne $\Opp \in [0,1]$ en divisant par $H(V|C) > 0$.
\cmt{Le rapport de deux quantit\'es positives est dans $[0,1]$ si le num\'erateur est major\'e par le d\'enominateur.}
\end{Preuve}
\begin{lemma}[Tension confidentialit\'e--information mutuelle]\label{lem:compromis_conf_mi}
Pour les protocoles $\beta$-r\'eguliers (Hypoth\`ese~\ref{ass:regularite})~:
\begin{equation}
\Conf \geq 1 - \sqrt{c_0\,\beta \cdot I(M;\Sigma|C)}
\label{eq:compromis_conf_mi}
\end{equation}
o\`u $c_0 = 1/(2\ln 2) \approx 0{,}721$. Cette minoration est la forme
utilis\'ee dans toute la th\`ese~; elle est reprise, sous sa forme
d\'efinitive et avec ses hypoth\`eses compl\`etes, au
Lemme~\ref{lem:confidentiality_guarantee}.
\end{lemma}

\begin{Preuve}[Lemme~\ref{lem:compromis_conf_mi}~: Trade-off Conf-MI]\label{prv:compromis_conf_mi}

\begin{cmtEtape}{Objectif.}
\'Etablir que $\Conf \geq 1 - \sqrt{c_0 \beta I}$. Autrement dit, tant
que la preuve divulgue peu d'information sur le message, la
confidentialit\'e reste garantie~; c'est lorsque $I$ cro\^it que cette
garantie s'affaiblit. C'est la premi\`ere moiti\'e du m\'ecanisme du
trilemme.
\end{cmtEtape}

Par l'Hypoth\`ese~\ref{ass:regularite} ($\beta$-r\'egularit\'e) et l'in\'egalit\'e de
Pinsker (Lemme~\ref{lem:pinsker}), pour tout couple $(m_0, m_1)$ et
contexte $c$~:
\begin{align}
\TV^2(P_{\Sigma|m_0,c}, P_{\Sigma|m_1,c})
&\leq c_0 \KL(P_{\Sigma|m_0,c} \| P_{\Sigma|m_1,c})
\\
&\leq c_0 \beta \cdot I(M;\Sigma|C)
\end{align}
\cmtFormule{Application directe de Pinsker.}
\cmtFormule{Application de l'hypoth\`ese de $\beta$-r\'egularit\'e.}
\cmt{La distance TV est born\'ee par la divergence KL.}
\cmt{La divergence KL est born\'ee par $\beta \cdot I(M;\Sigma|C)$.}

En prenant le maximum sur tous les couples et contextes, on obtient
\begin{equation}
\max_{\substack{m_0, m_1 \in \mathcal{M}_\lambda \\ c \in \mathcal{C}_\lambda}}
\TV(P_{\Sigma|m_0,c}, P_{\Sigma|m_1,c})
\leq \sqrt{c_0\beta \cdot I(M;\Sigma|C)}
\end{equation}
\cmtFormule{Le pire cas de distinguabilit\'e est born\'e par $\sqrt{c_0 \beta I}$.}

Puisque $\Conf(\mathcal{P}) = 1 - \max \TV(\cdots)$ par d\'efinition
de $\Conf$, il vient
\begin{equation}
\Conf(\mathcal{P}) \geq 1 - \sqrt{c_0\,\beta \cdot I(M;\Sigma|C)}
\end{equation}
\cmtFormule{La garantie de confidentialit\'e d\'ecro\^it en racine de l'information mutuelle.}
\cmt{Aucune normalisation par $\lambda$ n'intervient, car l'argument de Pinsker porte sur des lois conditionnelles, non sur la taille de l'espace de messages.}
\end{Preuve}

\begin{remark}[Interpr\'etation du lemme]\label{rem:interpretation_lemme}
Ce lemme fournit une \emph{garantie} de confidentialit\'e. Plus
$I(M;\Sigma|C)$ est petit, plus le minorant de $\Conf$ est proche de~$1$,
selon une d\'ependance non lin\'eaire (en racine carr\'ee). La tension
$\Conf$--$\Rel$ elle-m\^eme vient de l'in\'egalit\'e de Fano, puisqu'une fiabilit\'e
de reconstruction \'elev\'ee exige que la preuve contienne de l'information
sur le message. Un minorant de $\Conf$ ne suffit pas, \`a lui seul, \`a
formaliser cette tension.
\end{remark}

\begin{lemma}[Lemme de Johnson-Lindenstrauss~\cite{JohnsonLindenstrauss84}]\label{lem:johnson_lindenstrauss}
Pour tout $0 < \varepsilon < 1$, tout entier $n > 0$ et tout ensemble 
 $S \subset \mathbb{R}^n$ avec $|S| = N$, il existe une application lin\'eaire 
 $\Phi: \mathbb{R}^n \to \mathbb{R}^d$ avec $d = O(\varepsilon^{-2} \log N)$ 
telle que pour tous $x, y \in S$ :
\begin{equation}
(1-\varepsilon)\|x-y\|^2 \leq \|\Phi x - \Phi y\|^2 \leq (1+\varepsilon)\|x-y\|^2
\label{eq:jl_preservation}
\end{equation}
avec haute probabilit\'e lorsque $\Phi$ est une matrice al\'eatoire avec des 
entr\'ees $\Phi_{ij} \sim \mathcal{N}(0, 1/d)$ ou sous-gaussiennes.
\end{lemma}

Ce lemme garantit la pr\'eservation des distances euclidiennes sous compression
al\'eatoire, propri\'et\'e sur laquelle repose la construction SOTP-Base
(section~\ref{sec:sotp_construction}, sous-section~\ref{subsec:sotp_base}),
puisque les distances entre repr\'esentations de contextes sont pr\'eserv\'ees \`a un
facteur $1 \pm \varepsilon$ pr\`es apr\`es compression. Que l'interpr\`ete
institutionnel puisse en extraire autant d'information d\'ecisionnelle
qu'avant compression est une hypoth\`ese de conception, non une cons\'equence
du lemme.

\begin{theorem}[Goulot d'\'etranglement contextuel]\label{thm:goulot_contexte}
Pour tout protocole CLO~:
\begin{equation}
I(V;\Sigma|C) \leq \min\bigl\{H(\Sigma|C),\; H(V|C)\bigr\}
  \leq \log_2|\mathcal{V}_\lambda| \leq \alpha\, H(\mathcal{C})
\label{eq:goulot_contextuel}
\end{equation}
\end{theorem}

\begin{Preuve}[Th\'eor\`eme~\ref{thm:goulot_contexte}~: Goulot d'\'etranglement]\label{prv:goulot_contexte}

\begin{cmtEtape}{Objectif.}
Montrer que $I(V;\Sigma|C) \leq \min\{H(\Sigma|C),\, H(V|C)\} \leq \alpha H(C)$.
Autrement dit, la preuve $\Sigma$ ne peut transmettre plus d'information
sur le verdict que ce qu'elle contient, ni plus que l'incertitude
initiale sur le verdict, laquelle est plafonn\'ee par la granularit\'e des
verdicts rapport\'ee \`a la richesse du contexte.
\end{cmtEtape}

\textbf{Borne 1 (contrainte par la capacit\'e de preuve).}
Par d\'efinition de l'information mutuelle conditionnelle, on a
\begin{equation}
I(V;\Sigma|C) = H(\Sigma|C) - H(\Sigma|V,\mathcal{C})
\end{equation}
\cmtFormule{L'information sur $V$ contenue dans $\Sigma$ est l'entropie de $\Sigma$ moins ce qui reste une fois $V$ connu.}

Puisque l'entropie conditionnelle est non-n\'egative~\cite{CoverThomas06},
$H(\Sigma|V,\mathcal{C}) \geq 0$, on obtient
\begin{equation}
I(V;\Sigma|C) \leq H(\Sigma|C)
\end{equation}
\cmtFormule{Enlever un terme positif donne une in\'egalit\'e.}
\cmt{On ne peut pas extraire plus d'information que la preuve n'en contient.}

\textbf{Borne 2 (contrainte par l'incertitude sur les verdicts).}
Par la propri\'et\'e sym\'etrique de l'information mutuelle, on a
\begin{equation}
I(V;\Sigma|C) = H(V|C) - H(V|\Sigma,\mathcal{C})
\end{equation}
\cmtFormule{Formulation sym\'etrique, o\`u l'information \'egale l'incertitude initiale moins l'incertitude r\'esiduelle.}

Puisque $H(V|\Sigma,\mathcal{C}) \geq 0$, il vient
\begin{equation}
I(V;\Sigma|C) \leq H(V|C)
\end{equation}
\cmtFormule{Le verdict ne peut \^etre r\'esolu au-del\`a de son incertitude initiale.}
\cmt{On ne peut pas r\'esoudre plus d'incertitude qu'il n'y en a.}

\textbf{Borne sur $H(V|C)$.}
L'entropie d'une variable est major\'ee par le logarithme de la taille de
son support~\cite{CoverThomas06}. Or la granularit\'e des verdicts
(D\'efinition~\ref{def:alpha_granularite},
\'equation~\eqref{eq:alpha_definition}) borne ce support, d'o\`u
\begin{equation}
H(V|C) \leq \log_2 |\mathcal{V}_\lambda| \leq \alpha\, H(\mathcal{C})
\end{equation}
\cmtFormule{L'incertitude sur le verdict ne peut d\'epasser le nombre de verdicts distincts.}
\cmt{C'est la granularit\'e $\alpha$, param\`etre institutionnel, qui relie le plafond au contexte.}

\textbf{Conclusion.}
La combinaison des trois bornes donne
\begin{equation}
I(V;\Sigma|C) \leq \min\{H(\Sigma|C),\; H(V|C)\} \leq \alpha\, H(C)
\end{equation}
\cmtFormule{Le plafond de l'extraction est fix\'e par la granularit\'e institutionnelle des verdicts.}
\end{Preuve}

\begin{remark}[Verdict d\'eterministe]\label{rem:verdict_deterministe}
Si $V$ est une fonction d\'eterministe de $(\Sigma,C)$, alors
$H(V|\Sigma,C)=0$ et $I(V;\Sigma|C)=H(V|C)$, de sorte que l'opposabilit\'e vaut $1$
d\`es que $H(V|C)>0$, quel que soit le protocole. La m\'etrique $\Opp$ ne
discrimine les protocoles que pour des interpr\'etations admissibles au sens
de la D\'efinition~\ref{def:interp_admissible}, o\`u le verdict d\'epend
aussi de ressources institutionnelles $K$ non r\'eductibles \`a $(\Sigma,C)$.
Ce fait est certifi\'e en Lean (\texttt{opp\_eq\_one\_of\_deterministic},
section~\ref{sec:lean4}).
\end{remark}

Ce th\'eor\`eme \'etablit qu'aucune sophistication cryptographique ne saurait
d\'epasser $\alpha H(\mathcal{C})$ bits d'information extractible par un interpr\`ete
institutionnel~\cite{Smith09,Alvim12}. La granularit\'e des verdicts rapport\'ee \`a
l'entropie contextuelle constitue un plafond sur
l'opposabilit\'e, plafond qui pr\'efigure le Trilemme~CRO classique de la section~\ref{sec:cro_classique} et son extension quantique de la section~\ref{sec:cro_pq}. Les protocoles dans des contextes
simples ($H(\mathcal{C}) \approx 20$ bits, comme ECDSA pour contrat
standardis\'e) peuvent atteindre haute opposabilit\'e sans effort~; les
contextes complexes ($H(\mathcal{C}) \approx 70$ bits, comme zk-SNARKs
pour conformit\'e AML) imposent un plafond sauf si la capacit\'e
institutionnelle $\Delta$ est exceptionnellement \'elev\'ee.

\begin{boiteArbitre}
L'arbitre comprend que le plafond tient moins \`a la cryptographie qu'\`a
la finesse de ses propres verdicts, rapport\'ee au contexte. Plus le contexte est simple, plus il peut
d\'ecider, sans que la preuve y puisse quoi que ce soit.
\end{boiteArbitre}

\begin{theorem}[R\`egles de composition]\label{thm:composition_sequentielle}
Soit $\mathcal{P}_1 \circ \mathcal{P}_2$ le protocole qui publie
$\Sigma_2 = \mathcal{P}_2(\Sigma_1)$, o\`u $\Sigma_1$ est produite par
$\mathcal{P}_1$ (post-traitement, \'eventuellement randomis\'e,
ind\'ependant de $M$ sachant $(\Sigma_1, C)$). Alors
\begin{align}
\Conf(\mathcal{P}_1 \circ \mathcal{P}_2) &\geq \Conf(\mathcal{P}_1),
  \label{eq:comp_seq_conf}\\
\Relrec(\mathcal{P}_1 \circ \mathcal{P}_2) &\leq \Relrec(\mathcal{P}_1).
  \label{eq:comp_seq_rel}
\end{align}
Soit $\mathcal{P}_1 \| \mathcal{P}_2$ le protocole qui publie
$(\Sigma_1,\Sigma_2)$ produits ind\'ependamment sachant $(M,C)$. Alors
\begin{align}
\Conf(\mathcal{P}_1 \| \mathcal{P}_2) &\leq \min\bigl\{\Conf(\mathcal{P}_1),\, \Conf(\mathcal{P}_2)\bigr\},
  \label{eq:comp_par_conf}\\
\Relrec(\mathcal{P}_1 \| \mathcal{P}_2) &\geq \max\bigl\{\Relrec(\mathcal{P}_1),\, \Relrec(\mathcal{P}_2)\bigr\}.
  \label{eq:comp_par_rel}
\end{align}
\end{theorem}

\begin{Preuve}[Th\'eor\`eme~\ref{thm:composition_sequentielle}~: R\`egles de composition]\label{prv:composition}

\begin{cmtEtape}{Objectif.}
\'Etablir le comportement de la confidentialit\'e et de la fiabilit\'e sous
post-traitement (composition s\'equentielle) et sous publication conjointe
(composition parall\`ele).
\end{cmtEtape}

\textbf{Composition s\'equentielle.}
$M \to (\Sigma_1,C) \to (\Sigma_2,C)$ est une cha\^ine de Markov.
L'in\'egalit\'e de traitement des donn\'ees pour la distance en variation
totale~\cite{CoverThomas06} donne, pour tous $m_0, m_1, c$~:
\begin{equation}
\TV(P_{\Sigma_2|m_0,c}, P_{\Sigma_2|m_1,c}) \leq \TV(P_{\Sigma_1|m_0,c}, P_{\Sigma_1|m_1,c}),
\end{equation}
\cmtFormule{Un traitement ne cr\'ee pas de distinguabilit\'e.}
d'o\`u~\eqref{eq:comp_seq_conf}. Tout d\'ecodeur de $M$ \`a partir de
$(\Sigma_2,C)$ en induit un \`a partir de $(\Sigma_1,C)$, obtenu en appliquant d'abord
$\mathcal{P}_2$ \`a $\Sigma_1$, puis ce d\'ecodeur. D'o\`u~\eqref{eq:comp_seq_rel}.
\cmt{Post-traiter une preuve ne peut que la rendre plus confidentielle et moins reconstructible.}

\textbf{Composition parall\`ele.}
La marginale de $(\Sigma_1,\Sigma_2)$ sur $\Sigma_i$ est un post-traitement
de la paire. Par la m\^eme in\'egalit\'e, la distance en variation totale
entre lois conjointes majore celle entre marginales, pour chaque $i$~;
d'o\`u~\eqref{eq:comp_par_conf}. Un d\'ecodeur qui ignore l'une des
composantes est admissible pour la paire, d'o\`u~\eqref{eq:comp_par_rel}.
\cmt{Publier deux preuves r\'ev\`ele au moins autant que chacune d'elles.}
\end{Preuve}

\begin{remark}[Opposabilit\'e compos\'ee]\label{rem:composition_opp}
Pour l'opposabilit\'e, la cha\^ine $V$--$(\Sigma_1,C)$--$(\Sigma_2,C)$ donne
$I(V;\Sigma_2|C) \leq I(V;\Sigma_1|C)$, donc
$\Opp(\mathcal{P}_1\circ\mathcal{P}_2) \leq \Opp(\mathcal{P}_1)$ \`a verdict
fix\'e. Aucune r\`egle g\'en\'erale ne vaut pour la composition parall\`ele,
car l'information mutuelle n'est ni sous- ni sur-additive en g\'en\'eral. Les
r\`egles de~\cite{MINKA2026CLO} ($\Opp \geq \min$ en s\'equentiel,
$\Opp \approx$ produit en parall\`ele) sont des heuristiques qui ne sont pas
reprises ici. La formulation initiale de ce th\'eor\`eme dans~\cite{MINKA2026CLO}
(minoration par le produit via l'in\'egalit\'e de Fr\'echet, $\Conf \geq \max$
en parall\`ele) est remplac\'ee ici par une formulation plus pr\'ecise,
l'in\'egalit\'e de Fr\'echet donnant $ab \geq a+b-1$.
\end{remark}

Ces r\`egles permettent d'analyser des protocoles composites. Ainsi, un syst\`eme
bancaire publiant conjointement une signature ECDSA (authentification) et une preuve zk-SNARK
(conformit\'e) h\'erite des bornes ci-dessus, puisque publier les deux preuves ne peut
qu'abaisser la confidentialit\'e par rapport \`a chacune d'elles.

\subsubsection{Jeux cryptographiques et m\'etriques CLO}
\label{subsubsec:jeux_crypto}

Les d\'efinitions informationnelles de $\Conf$, $\Rel$ et $\Opp$ sont
compl\'et\'ees par des jeux cryptographiques qui captent leur contenu
op\'erationnel en termes d'avantages d'adversaires PPT. Cette section
\'etablit le pont formel entre les deux formalismes.

\paragraph{Jeu IND-CONF : confidentialit\'e game-based.}
\label{subsubsubsec:jeu_ind_conf}

\begin{definition}[Jeu IND-CONF]\label{def:jeu_ind_conf}
Pour un protocole $\mathcal{P} = (\mathsf{Setup},\mathsf{Prove},
\mathsf{Verify})$ et adversaire $\mathcal{A}$~:
\begin{enumerate}
\item $\mathcal{A}$ re\c{c}oit $pp \leftarrow \mathsf{Setup}(1^\lambda)$
\item $\mathcal{A}$ choisit $m_0,m_1 \in \mathcal{M}_\lambda$
      ($|m_0|=|m_1|$) et $c \in \mathcal{C}_\lambda$
      \hfill\textit{(phase de d\'efi)}
      \cmt{L'adversaire choisit les deux messages \`a distinguer.}
\item Le challenger tire $b \xleftarrow{\$} \{0,1\}$ et calcule
      $\sigma \leftarrow \mathsf{Prove}(pp, m_b, c)$
      \cmt{Le challenger chiffre l'un des deux messages au hasard.}
\item $\mathcal{A}$ re\c{c}oit $\sigma$ et \'emet $b' \in \{0,1\}$
      \hfill\textit{(phase de r\'eponse)}
      \cmt{L'adversaire devine lequel des deux messages a \'et\'e chiffr\'e.}
\item L'exp\'erience produit 1 si $b' = b$, sinon 0
\end{enumerate}
\begin{equation}
\mathsf{Adv}^{\mathsf{IND\text{-}CONF}}_{\mathcal{A},\mathcal{P}}(\lambda)
= \left|\Pr\!\left[
\mathsf{Exp}^{\mathsf{IND\text{-}CONF}}_{\mathcal{A},\mathcal{P}}(\lambda)
= 1\right] - \frac{1}{2}\right|
\end{equation}
\cmtFormule{L'avantage est l'\'ecart \`a la devinette al\'eatoire (probabilit\'e 1/2).}
\cmt{Un avantage n\'egligeable signifie confidentialit\'e.}
\end{definition}

\begin{remark}[Pont IND-CONF $\leftrightarrow$ $\Conf$]
\label{rem:jeu_ind_conf_link}
Pour tout distingueur $\mathcal{A}$, non born\'e en calcul, qui choisit
$(m_0,m_1,c)$~:
\begin{equation}
\mathsf{Adv}^{\mathsf{IND\text{-}CONF}}_{\mathcal{A},\mathcal{P}}(\lambda)
\;\leq\; \frac{1}{2}\max_{m_0,m_1,c}\TV(P_{\Sigma|m_0,c},P_{\Sigma|m_1,c})
\;=\; \frac{1 - \Conf(\mathcal{P})}{2},
\end{equation}
\cmtFormule{L'avantage IND-CONF est born\'e par la distinguabilit\'e statistique.}
avec \'egalit\'e pour le test de Neyman--Pearson (Helstrom) sur la paire la
plus distinguable. $\Conf$ est donc la version \emph{statistique} de
IND-CONF, au sens o\`u $\Conf \geq 1 - \negl(\lambda)$ implique un avantage n\'egligeable
pour tout adversaire. La r\'eciproque est fausse, car pour des adversaires PPT
l'avantage peut \^etre n\'egligeable (s\'ecurit\'e calculatoire) alors que
$\Conf$ est petit (chiffrement \`a cl\'e publique, par exemple).
\cmt{La confidentialit\'e statistique implique la confidentialit\'e calculatoire, pas l'inverse.}
\end{remark}

\paragraph{Jeu BIND-REL : fiabilit\'e game-based.}
\label{subsubsubsec:jeu_bind_rel}

\begin{definition}[Jeu BIND-REL]\label{def:jeu_bind_rel}
Pour un protocole $\mathcal{P}$ avec $\mathsf{Open}$ et adversaire
$\mathcal{A}$~:
\begin{enumerate}
\item $\mathcal{A}$ re\c{c}oit $pp \leftarrow \mathsf{Setup}(1^\lambda)$
\item $\mathcal{A}$ produit $(c, \sigma, s_0, s_1)$
      \cmt{L'adversaire forge une preuve et deux t\'emoins.}
\item Calculer $m_i = \mathsf{Open}(pp, \sigma, s_i)$ pour $i \in \{0,1\}$
\item $\mathcal{A}$ \emph{gagne} si et seulement si
  $\mathsf{Verify}(pp,\sigma,c)=1$,
  $m_0 \neq m_1$,
  $m_0 \neq \bot$ et $m_1 \neq \bot$
  \cmt{Il gagne si la preuve passe ET les deux messages sont distincts et valides.}
\end{enumerate}
\begin{equation}
\mathsf{Win}^{\mathsf{BIND\text{-}REL}}_{\mathcal{A},\mathcal{P}}(\lambda)
= \Pr\!\left[\mathcal{A} \text{ gagne dans }
\mathsf{Exp}^{\mathsf{BIND\text{-}REL}}_{\mathcal{A},\mathcal{P}}(\lambda)
\right]
\end{equation}
\cmtFormule{Probabilit\'e que l'adversaire produise une collision d'engagement.}
\end{definition}

\begin{remark}[Pont BIND-REL $\leftrightarrow$ $\Relver$]
\label{rem:jeu_bind_rel_link}
BIND-REL contr\^ole la \emph{solidit\'e} de l'ouverture, donc la fiabilit\'e
de v\'erification (Remarque~\ref{rem:rel_ver}). Pour un protocole complet
dont la probabilit\'e de gain BIND-REL est au plus $\varepsilon_{\mathsf{bind}}$,
une preuve accept\'ee s'ouvre sur un unique message sauf avec probabilit\'e
$\varepsilon_{\mathsf{bind}}$, de sorte que
\begin{equation}
\Relver \geq 1 - \varepsilon_{\mathsf{bind}} - \negl(\lambda).
\end{equation}
\cmtFormule{Fiabilit\'e de v\'erification born\'ee inf\'erieurement par $1 - \varepsilon_{\mathsf{bind}}$.}
BIND-REL ne dit rien de $\Relrec$, puisqu'un engagement parfaitement masquant et
calculatoirement liant a $\Relrec = 2^{-\lambda}$.
\cmt{Lier une preuve \`a un message ne revient pas \`a le r\'ev\'eler.}
\end{remark}

\paragraph{Jeu EXT-OPP : opposabilit\'e game-based.}
\label{subsubsubsec:jeu_ext_opp}

\begin{definition}[Jeu EXT-OPP]\label{def:jeu_ext_opp}
Pour un protocole $\mathcal{P}$, un interpr\`ete $\mathsf{Interpret}$
de capacit\'e $\Delta$ et une distribution $\mathcal{D}_C$ fix\'ee \`a l'avance~:
\begin{enumerate}
\item Le challenger tire $c \leftarrow \mathcal{D}_C$,
      $m \leftarrow \mathcal{M}_\lambda$
      \cmt{Le contexte est tir\'e selon la distribution du domaine.}
\item Le challenger calcule $\sigma \leftarrow \mathsf{Prove}(pp,m,c)$
\item L'interpr\`ete re\c{c}oit $(\sigma,c)$ et produit
      $v \leftarrow \mathsf{Interpret}(pp,\sigma,c,\Delta)$
      \cmt{L'interpr\`ete rend un verdict \`a partir de la preuve et du contexte.}
\item Soit $v^*$ le verdict de r\'ef\'erence, r\'ealisation de la variable $V$
      (qui peut d\'ependre du message, du contexte et des ressources
      institutionnelles, D\'efinition~\ref{def:interp_admissible})
      \cmt{Le verdict de r\'ef\'erence est celui que l'institution rendrait en pleine connaissance de cause.}
\item L'interpr\`ete \emph{gagne} si $v = v^*$
\end{enumerate}
\begin{equation}
\mathsf{Err}^{\mathsf{EXT\text{-}OPP}}_{\mathsf{Interpret},\mathcal{P},\mathcal{D}_C}(\lambda)
= \Pr\!\left[
\mathsf{Interpret} \text{ perd dans }
\mathsf{Exp}^{\mathsf{EXT\text{-}OPP}}\right]
\end{equation}
\cmtFormule{Probabilit\'e que l'interpr\`ete rende un verdict incorrect.}
\end{definition}

\begin{remark}[Pont EXT-OPP $\leftrightarrow$ $\Opp$]
\label{rem:jeu_ext_opp_link}
Par l'in\'egalit\'e de Fano appliqu\'ee \`a l'estimation de $V$ \`a partir de
$(\Sigma,C)$, $H(V|\Sigma,C) \leq h(\varepsilon_{\mathsf{ext}}) +
\varepsilon_{\mathsf{ext}}\log_2(|\mathcal{V}_\lambda|-1)$, d'o\`u
\begin{equation}
\Opp = 1 - \frac{H(V|\Sigma,C)}{H(V|C)}
\geq 1 - \frac{h(\varepsilon_{\mathsf{ext}})
  + \varepsilon_{\mathsf{ext}}\log_2(|\mathcal{V}_\lambda|-1)}
  {H(V|C)}
\label{eq:jeu_fano_opp}
\end{equation}
\cmtFormule{L'opposabilit\'e est born\'ee inf\'erieurement par $1 - \varepsilon_{\mathsf{ext}}$.}
o\`u $\varepsilon_{\mathsf{ext}} = \mathsf{Err}^{\mathsf{EXT\text{-}OPP}}$
et $h(\cdot)$ est l'entropie binaire.
Un interpr\`ete avec faible erreur d'extraction implique une forte
opposabilit\'e.
\cmt{Moins l'interpr\`ete se trompe, plus l'opposabilit\'e est grande.}
\end{remark}

\end{sloppypar}

% -----------------------------------------------------------------------
\subsection{Extension aux m\'etriques quantiques}
\label{subsec:extension_postquantique}
% -----------------------------------------------------------------------
\begin{sloppypar}
La migration vers les standards post-quantiques NIST (FIPS~203, 204, 205)~\cite{NISTPQC,NISTPostQuantum2022}
requiert une extension du cadre CLO. Cette section esquisse les modifications
n\'ecessaires, d\'evelopp\'ees en d\'etail dans la section~\ref{sec:cro_pq}.

\subsubsection{Confidentialit\'e quantique}
\label{subsubsec:conf_quantique}

En pr\'esence d'adversaires quantiques, la d\'efinition de confidentialit\'e
doit consid\'erer les mesures POVM sur les \'etats quantiques de
preuve~\cite{Watrous2009,Boneh2011,Zhandry}. Comme dans le cas classique, la
m\'etrique retenue est \emph{statistique}, puisqu'elle quantifie la meilleure
distinction possible, sans restriction de calcul.

\begin{definition}[Confidentialit\'e quantique]\label{def:conf_quantique_clo}
Pour le protocole $\mathcal{P}$ avec param\`etre de s\'ecurit\'e $\lambda$, soit
$\rho_\Sigma^{m,c}$ l'\'etat quantique de preuve pour le message $m$ dans le
contexte $c$. La \emph{confidentialit\'e quantique} est
\begin{equation}
\Conf^q(\mathcal{P}) := 1 - \max_{m_0,m_1 \in \mathcal{M}_\lambda,\; c \in \mathcal{C}_\lambda}
\tfrac{1}{2}\bigl\|\rho_\Sigma^{m_0,c} - \rho_\Sigma^{m_1,c}\bigr\|_1 .
\label{eq:conf_q}
\end{equation}
Par le th\'eor\`eme de Helstrom, $\tfrac12\|\rho_0-\rho_1\|_1$ est l'avantage
maximal de toute mesure POVM pour distinguer $\rho_0$ de $\rho_1$.
\end{definition}

\begin{remark}[Compatibilit\'e avec la confidentialit\'e classique]
\label{rem:compatibilite_conf}
Pour les protocoles classiques o\`u $\Sigma$ prend des valeurs dans un alphabet
fini $\mathcal{S}$ (i.e., $\rho_\Sigma^{m,c} = \sum_{s} P_{\Sigma|M,C}(s|m,c)\,|s\rangle\langle s|$
est diagonal), la distance de trace se r\'eduit \`a la distance de variation totale~\cite{Levin,Pinsker1964}~:
\begin{equation}
\tfrac{1}{2}\|\rho_\Sigma^{m_0,c} - \rho_\Sigma^{m_1,c}\|_1
= \mathrm{TV}\!\bigl(P_{\Sigma|m_0,c},\, P_{\Sigma|m_1,c}\bigr)
\end{equation}
Donc $\Conf^q$ se r\'eduit \`a $\Conf$ du cadre classique~\cite{MINKA2026CLO}, ce qui fait de la
d\'efinition une g\'en\'eralisation stricte. Les notions calculatoires
(adversaires QPT) rel\`event des jeux de s\'ecurit\'e, comme IND-CONF dans le
cas classique (Remarque~\ref{rem:jeu_ind_conf_link}).
\end{remark}

\subsubsection{Perte d'interpr\'etabilit\'e quantique}
\label{subsubsec:perte_interpretabilite}

\begin{definition}[Perte d'interpr\'etabilit\'e quantique]\label{def:perte_interp_quantique}
Soit $\Lambda$ une application CPTP repr\'esentant les effets d'un canal
quantique sur le support de la preuve (stockage, transmission,
v\'erification). La perte d'interpr\'etabilit\'e quantique $\eta_q \in [0,1]$
est le plus petit r\'eel tel que
\begin{equation}
I(V;\Lambda(\Sigma)|C) \;\geq\; I(V;\Sigma|C) - \eta_q\, H(C).
\label{eq:def_eta_q}
\end{equation}
\end{definition}

L'opposabilit\'e apr\`es passage dans le canal est alors
\begin{equation}
\Opp^q(\mathcal{P}) = \frac{I(V;\Lambda(\Sigma)|C)}{H(V|C)} \;\geq\; \Opp(\mathcal{P}) - \eta_q\,\frac{H(C)}{H(V|C)},
\label{eq:opp_q}
\end{equation}
une minoration qui quantifie la d\'egradation s\'emantique dans les canaux
quantiques~\cite{Holevo1973,Wilde2017}. Par l'in\'egalit\'e de traitement des
donn\'ees, $I(V;\Lambda(\Sigma)|C) \leq I(V;\Sigma|C)$. Un canal ne peut donc que
r\'eduire l'opposabilit\'e, tandis que $\eta_q$ mesure l'ampleur maximale de cette
r\'eduction, rapport\'ee \`a la richesse du contexte.

\begin{remark}[Interpr\'etation physique]
\label{rem:interpretation_physique}
$\eta_q$ quantifie la quantit\'e de contenu s\'emantique perdue lorsque la
preuve cryptographique passe par des canaux quantiques~\cite{Zurek2003,Schumacher1997}. Pour les syst\`emes
dont la preuve est une cha\^ine de bits classique, y compris les signatures
post-quantiques standardis\'ees, $\Lambda$ est l'identit\'e et $\eta_q = 0$. Pour
un support ou un v\'erificateur quantique, $\eta_q > 0$ en raison du bruit
(d\'ecoh\'erence, erreurs de portes)~\cite{Holevo1973,Preskill2018}.
\end{remark}

Ces extensions seront formalis\'ees et exploit\'ees \`a la section~\ref{sec:cro_pq} pour \'etablir le Trilemme CRO conditionnel quantique.
\end{sloppypar}

% -----------------------------------------------------------------------
% Synth\`ese
% -----------------------------------------------------------------------
\begin{boiteArbitre}
L'arbitre referme le dossier CLO et ach\`eve ainsi l'\'etape \textbf{F}. Cette section lui a donn\'e un
langage pour mesurer ce qu'il vit, qu'aucun corpus pr\'ec\'edent ne lui avait fourni.

Il retient trois contributions, dans l'ordre o\`u elles se sont
construites. Le mod\`ele de protocole d'abord (deux couches et deux questions
distinctes, portant sur la validit\'e de la preuve puis sur ce qu'elle signifie pour lui).
Les trois m\'etriques ensuite ($\Conf$, $\Rel$, $\Opp$), dont la derni\`ere
le concerne directement~; il sait d\'esormais qu'elle
n'a de sens que si son propre jugement ne se r\'eduit pas \`a un calcul
m\'ecanique sur la preuve. Le goulot d'\'etranglement contextuel enfin,
qui lui a appris que le plafond est fix\'e par la finesse de ses verdicts.

La section~\ref{sec:cro_classique} l'attend. Il va falloir maintenant
savoir ce qui, sous des hypoth\`eses explicites sur l'institution, borne
le produit des trois m\'etriques.
\end{boiteArbitre}
% ============================================================================
% CHAPITRE: TRILEMME CRO CONDITIONNEL CLASSIQUE
% source  : partie2/cro_classique.tex
% Version : 2.2 — alignée sur la formalisation Lean4 (section~\ref{sec:lean4})
% ============================================================================

\section{LE TRILEMME CRO CONDITIONNEL CLASSIQUE}
\label{sec:cro_classique}
\begin{sloppypar}
La section~\ref{sec:clo} a fix\'e le langage, avec trois m\'etriques
d\'efinies via la th\'eorie de l'information~\cite{Shannon48,CoverThomas06} et un Th\'eor\`eme du Goulot d'\'etranglement
Contextuel montrant que $I(V;\Sigma|C) \leq \alpha H(\mathcal{C})$. Aucun r\'esultat n'interdit
pourtant encore d'atteindre simultan\'ement $\Conf = \Rel = \Opp = 0{,}9$. Il
reste donc \`a d\'eterminer o\`u s'arr\^ete le r\'ealisable.

Nous bornons ici la r\'eponse, sous forme du \textbf{Trilemme CRO
conditionnel classique}~\cite{MINKA2025CRO,MINKA2026CRO}. Ce r\'esultat, premi\`ere caract\'erisation quantitative
des limites fondamentales de la cryptographie institutionnelle dans le r\'egime
classique, borne sup\'erieurement le produit $\Conf \cdot \Rel \cdot \Opp$ en fonction de la capacit\'e
d'extraction institutionnelle~$\Delta$, de la granularit\'e des verdicts~$\alpha$ et
du param\`etre de s\'ecurit\'e~$\lambda$. Une forme raffin\'ee de la borne identifie
exactement le r\'egime, d\'etermin\'e par la granularit\'e des verdicts et la richesse
du contexte, dans lequel la contrainte liante tient \`a la capacit\'e d'extraction des interpr\`etes institutionnels, au-del\`a de la seule in\'egalit\'e de Fano.

Deux contributions additionnelles compl\`etent cette section. La classification en
r\'egimes op\'erationnels (sous-section~\ref{subsec:regimes_operationnels}) transforme la
borne abstraite en guide de conception concret selon le niveau d'entropie
contextuelle. Une calibration illustrative (sous-section~\ref{subsec:analyse_comparative})
situe les ordres de grandeur des param\`etres institutionnels \`a partir des trois syst\`emes d\'eploy\'es \'etudi\'es
dans le CLO~\cite{Adida08,Groth16,Johnson01}, qui ne sont pas des instances du th\'eor\`eme
(Remarque~\ref{rem:portee_cas_etude}), conform\'ement \`a la note de m\'ethodologie pr\'esent\'ee dans la sous-section~\ref{subsec:note_methodo}.

Cette section correspond \`a l'\'etape \textbf{B (Borner)} de la progression
F$\to$\textbf{B}$\to$E$\to$S$\to$C$\to$X et apporte la contribution \textbf{C2} de cette th\`ese.

La r\'eponse \`a l'Objectif~\ref{obj:obj2_impossibilite} est donc affirmative, puisque des
limites fondamentales existent dans le r\'egime classique. Leur extension
au contexte quantique sera \'etablie \`a la section suivante.
\end{sloppypar}
\begin{boiteArbitre}
L'arbitre apprend qu'il existe une borne, qui rel\`eve d'une limite fondamentale, ind\'ependante de la technique et repose sur des hypoth\`eses qu'il peut lire et discuter.
Quand ses verdicts sont plus fins que ce qu'il sait extraire, sa propre capacit\'e d'extraction~$\Delta$ le limite davantage que la puissance
de l'adversaire. Les protocoles qu'il examine ne sont donc pas mal con\c{c}us,
car ils se heurtent \`a une incompatibilit\'e structurelle d\'esormais \'enonc\'ee
et prouv\'ee.
\end{boiteArbitre}

% -----------------------------------------------------------------------
\subsection{Mod\`ele formel \'etendu avec param\`etres institutionnels}
\label{subsec:modele_formel_etendu}
% -----------------------------------------------------------------------
\begin{sloppypar}
Nous d\'efinissons trois param\`etres compl\'etant le mod\`ele CLO de la section~\ref{sec:clo}. Ces quantit\'es convertissent des contraintes
qualitatives, \guillemotleft~capacit\'e cognitive limit\'ee~\guillemotright, \guillemotleft~d\'ecisions nuanc\'ees possibles~\guillemotright,
\guillemotleft~protocoles bien con\c{c}us~\guillemotright, en grandeurs mesurables permettant des bornes
rigoureuses~\cite{ClarksonSchneider10,Smith09}.

\subsubsection{Param\`etres institutionnels}
\label{subsubsec:parametres_institutionnels}

\begin{definition}[Capacit\'e d'extraction institutionnelle]\label{def:delta_extraction}

$\Delta \geq 0$ (en bits) quantifie la connaissance pr\'ealable des interpr\`etes
institutionnels et leurs ressources cognitives au-del\`a du contexte formel~$C$~\cite{Simon82,Miller56}.
Ce param\`etre agr\`ege les pr\'ec\'edents l\'egaux et la jurisprudence applicable~\cite{Mason07,Kerr04},
l'expertise de domaine des d\'ecideurs~\cite{Ericsson95,Chi2006}, la documentation et les proc\'edures
\'etablies, les capacit\'es de d\'elib\'eration collaborative entre experts~\cite{Sunstein2005,Hong2016}, ainsi
que la m\'emoire institutionnelle et les bases de donn\'ees historiques~\cite{Argyris1978,Walsh1995}.

\end{definition}

\begin{definition}[Granularit\'e des verdicts]\label{def:alpha_granularite}

$\alpha > 0$ (sans dimension) mesure la finesse de distinction des verdicts
institutionnels relativement \`a la richesse contextuelle~:
\begin{equation}
\log_2 |\mathcal{V}_\lambda| \leq \alpha \cdot H(C)
\label{eq:alpha_definition}
\end{equation}
Un tribunal rendu apte \`a distinguer douze niveaux de fraude dans un contexte \`a
entropie $H(C) = 4$ bits correspond \`a $\alpha = \log_2 12/4 \approx 0{,}9$~\cite{Shannon48,CoverThomas06}.

\end{definition}

\begin{definition}[R\'egularit\'e du protocole]\label{def:beta_regularite}

$\beta \geq 1$ borne la divergence Kullback-Leibler entre distributions de
preuves pour messages distincts~\cite{Kullback1951,CoverThomas06}~:
\begin{equation}
\KL\!\left(P_{\Sigma|m_0,c} \,\|\, P_{\Sigma|m_1,c}\right) \leq \beta \cdot I(M;\Sigma|C)
\label{eq:beta_definition}
\end{equation}
Ce param\`etre \'ecarte les protocoles pathologiques dont la divergence KL exploserait
pour certaines paires de messages~\cite{Tsybakov2009}. Au sens statistique, un protocole
qui publie le message en clair (signature) a des lois de preuve \`a supports disjoints
et $\KL = \infty$, de sorte qu'il n'est pas $\beta$-r\'egulier. La $\beta$-r\'egularit\'e est une
hypoth\`ese sur la classe de protocoles vis\'ee (protocoles \`a fuite contr\^ol\'ee), non
une propri\'et\'e observ\'ee des protocoles d\'eploy\'es~\cite{Pedersen91,Groth16}.

\end{definition}
\end{sloppypar}

% -----------------------------------------------------------------------
\subsubsection{Hypoth\`eses sur les capacit\'es institutionnelles}
\label{subsubsec:hypotheses_capacite}
% -----------------------------------------------------------------------

\begin{sloppypar}
Nous posons quatre hypoth\`eses sur les capacit\'es
institutionnelles~\cite{Simon1955,Miller56,Ericsson95}. Chacune repose
sur des justifications multidisciplinaires ; aucune n'a \'et\'e valid\'ee
\`a ce stade par protocoles exp\'erimentaux contr\^ol\'es mesurant
directement $\Delta$, $\alpha$, $\beta$ aupr\`es d'acteurs
institutionnels r\'eels, cette limitation est reconnue explicitement
(sous-section~\ref{subsec:note_methodo}).

\begin{definition}[Interpr\'etation institutionnelle admissible]\label{def:interp_admissible}
Une interpr\'etation est admissible si le verdict s'\'ecrit
$V = \mathsf{Interpret}(\Sigma, C, K)$, o\`u $K$ repr\'esente les ressources
propres de l'institution (jurisprudence, expertise, d\'elib\'eration,
al\'eas de d\'ecision) et satisfait
\begin{equation}
K \perp M \mid (\Sigma, C).
\label{eq:interp_admissible}
\end{equation}
L'institution n'acc\`ede au message que par la preuve, de sorte que $M \to (\Sigma,C) \to V$
est une cha\^ine de Markov. Le verdict peut \^etre stochastique.
\end{definition}

\begin{assumption}[Extraction born\'ee]
\label{ass:extraction_bornee}
Pour tout protocole CLO :
\begin{equation}
I(V;\Sigma|C) \leq I(M;\Sigma|C) + \Delta
\label{eq:extraction_bornee}
\end{equation}
\end{assumption}

\begin{proposition}[Fondement informationnel de l'Hypoth\`ese~\ref{ass:extraction_bornee}]
\label{prop:fondement_extraction}
Pour toute interpr\'etation admissible,
\begin{equation}
I(V;\Sigma|C) = I(V;M|C) + I(V;\Sigma|M,C) \;\leq\; I(M;\Sigma|C) + I(V;\Sigma|M,C).
\label{eq:fondement_extraction}
\end{equation}
En cons\'equence, l'Hypoth\`ese~\ref{ass:extraction_bornee} est satisfaite
d\`es que
\begin{equation}
I(V;\Sigma|M,C) \leq \Delta,
\label{eq:extraction_bornee_forte}
\end{equation}
et $\Delta$ s'interpr\`ete exactement comme l'information que
l'institution tire de la preuve \emph{au-del\`a du message}.
\end{proposition}

\begin{Preuve}[Proposition~\ref{prop:fondement_extraction}]\label{prv:fondement_extraction}
\begin{cmtEtape}{Objectif.}
Justifier l'Hypoth\`ese~\ref{ass:extraction_bornee} \`a partir d'une
propri\'et\'e structurelle du processus d'interpr\'etation.
\end{cmtEtape}
Par la r\`egle de cha\^ine appliqu\'ee \`a $I(V;M,\Sigma|C)$ de deux
mani\`eres, on obtient
\begin{equation}
I(V;M|C) + I(V;\Sigma|M,C) = I(V;\Sigma|C) + I(V;M|\Sigma,C).
\end{equation}
\cmtFormule{Deux d\'ecompositions de la m\^eme information mutuelle.}
Sous~\eqref{eq:interp_admissible}, $V$ ne d\'epend de $M$ qu'\`a travers
$(\Sigma,C)$, donc $I(V;M|\Sigma,C)=0$, ce qui donne l'\'egalit\'e
de~\eqref{eq:fondement_extraction}. L'in\'egalit\'e de traitement des
donn\'ees sur $M \to (\Sigma,C) \to V$ donne
$I(V;M|C) \leq I(M;\Sigma|C)$.
\cmt{L'institution ne peut pas apprendre sur le message plus que la preuve n'en contient~; elle peut en revanche lire dans la preuve des traits qui ne sont pas le message.}
\end{Preuve}

\begin{assumption}[Verdict efficace]
\label{ass:verdicts_efficace}
La taille de l'espace de verdict satisfait $|\mathcal{V}_\lambda| \leq
2^{\alpha H(C)}$ pour une constante $\alpha > 0$, \emph{et} la fonction
de verdict utilise efficacement le contexte disponible :
\begin{equation}
H(V|C) \geq \alpha H(C) - 1 =: \gamma, \qquad \gamma \geq 1
\label{eq:verdicts_efficace}
\end{equation}
\end{assumption}

\begin{remark}[Statut de l'Hypoth\`ese~\ref{ass:verdicts_efficace}]
\label{rem:statut_verdict_efficace}
L'in\'egalit\'e $|\mathcal{V}_\lambda| \leq 2^{\alpha H(C)}$ est une
borne \emph{sup\'erieure} sur la taille de l'espace de verdict. Elle ne
permet pas, \`a elle seule, de d\'eduire la borne \emph{inf\'erieure}
$H(V|C) \geq \alpha H(C) - 1$. Cette borne inf\'erieure est la
propri\'et\'e d'\textbf{efficacit\'e} du verdict, qui doit \^etre
postul\'ee explicitement. Elle formalise l'id\'ee qu'une institution
exploitant effectivement le contexte disponible distingue au moins
$\gamma$ bits de verdicts. Pour les param\`etres institutionnels typiques
($H(C) \geq 10$ bits, $\alpha \approx 1$), on a $\gamma \geq 9$.
\end{remark}

\begin{assumption}[R\'egularit\'e]
\label{ass:regularite}
La divergence KL satisfait $\KL(P_{\Sigma|m_0,c} \| P_{\Sigma|m_1,c})
\leq \beta \cdot I(M;\Sigma|C)$.
\end{assumption}

\begin{assumption}[R\'egime de pr\'eservation de la confidentialit\'e]
\label{ass:regime}
Il existe une constante $\kappa \geq 0$, ind\'ependante de $\lambda$,
telle que :
\begin{equation}
I(M;\Sigma|C) \leq \frac{\kappa}{\lambda}
\label{eq:regime}
\end{equation}
\end{assumption}

\begin{remark}[Statut et r\^ole de l'Hypoth\`ese~\ref{ass:regime}]
\label{rem:statut_regime}
Dans la version initiale de la preuve, ce r\'egime \'etait pr\'esent\'e
comme une \emph{cons\'equence} de l'optimisation de $f(I)$ sur
$[0, \lambda]$. La formalisation a conduit \`a une r\'esolution alternative de cette
optimisation. La d\'eriv\'ee $f'(x) = 2x - \frac{5}{2}ax^{3/2}$ s'annule en
$\sqrt{x^*} = \frac{4}{5a}$ et ce point est un \emph{maximiseur} de $f$. Le r\'egime
$I = O(1/\lambda)$ ne peut donc pas \^etre d\'eduit de cette
optimisation.

L'Hypoth\`ese~\ref{ass:regime} est donc pos\'ee explicitement
comme condition du th\'eor\`eme. Elle est pr\'ecis\'ement le r\'egime
dans lequel les protocoles $\beta$-r\'eguliers garantissent la
confidentialit\'e, comme l'\'etablit le
Lemme~\ref{lem:confidentiality_guarantee}. Si
$I \leq \kappa/\lambda$, alors
$\Conf \geq 1 - \sqrt{c_0 \beta \kappa/\lambda} \to 1$ quand
$\lambda \to \infty$.
\end{remark}
\end{sloppypar}

% -----------------------------------------------------------------------
\begin{figure}[htbp]
\centering
\begin{tikzpicture}[scale=1.0, transform shape]
% --- Borne de Fano (1+kappa/lambda)/lambda ---
\draw[enspred, thick, dashed]
    plot[smooth] coordinates {
        (1.60, 6.348)
        (2.00, 5.062)
        (2.40, 4.210)
        (2.80, 3.603)
        (3.20, 3.149)
        (4.00, 2.516)
        (4.80, 2.094)
        (5.60, 1.794)
        (6.40, 1.569)
        (8.00, 1.254)
        (9.60, 1.044)
    }
    node[right, font=\small, text=enspred] {Fano seul};
% --- Borne raffin\'ee, Delta < gamma ---
\draw[enspblue, thick]
    plot[smooth] coordinates {
        (1.60, 1.676)
        (2.00, 1.336)
        (2.40, 1.110)
        (2.80, 0.950)
        (3.20, 0.830)
        (4.00, 0.663)
        (4.80, 0.552)
        (5.60, 0.472)
        (6.40, 0.413)
        (8.00, 0.330)
        (9.60, 0.275)
    }
    node[right, font=\small, text=enspblue, align=left] {Borne raffin\'ee\\$\Delta=5$, $\gamma=19$};
% --- Axes ---
\draw[-{Stealth[length=3mm]}, thick] (0,0) -- (10.5, 0)
    node[right, font=\small] {$\lambda$};
\draw[-{Stealth[length=3mm]}, thick] (0,0) -- (0, 7.0)
    node[above, font=\small] {Majorant de $\Conf \cdot \Rel \cdot \Opp$};
\foreach \y/\lab in {2.0/0{,}005, 4.0/0{,}010, 6.0/0{,}015} {
    \draw[thin, chaptercolor] (-0.1, \y) -- (0.1, \y);
    \node[left, font=\footnotesize] at (-0.1, \y) {$\lab$};
}
\node[left, font=\footnotesize] at (-0.1, 0) {$0$};
\foreach \x/\lab in {1.6/64, 3.2/128, 6.4/256, 9.6/384} {
    \draw[thin, chaptercolor] (\x, -0.1) -- (\x, 0.1);
    \node[below, font=\footnotesize] at (\x, -0.1) {$\lab$};
}
\end{tikzpicture}
\caption[Borne raffin\'ee du Trilemme CRO et borne de Fano]{%
  Majorants de $\Conf \cdot \Rel \cdot \Opp$ en fonction de $\lambda$
  ($\kappa = 1$)~: borne raffin\'ee~\eqref{eq:trilemme_cro_raffine} pour
  $\Delta = 5$, $\alpha = 1$, $H(\mathcal{C}) = 20$ bits ($\gamma = 19$, trait
  plein) et borne de Fano $(1+\kappa/\lambda)/\lambda$ (tirets). L'\'ecart entre
  les deux mesure l'effet de la capacit\'e institutionnelle lorsque
  $\Delta < \gamma$~; pour $\Delta \geq \gamma$, les deux courbes co\"incident.%
}\label{fig:cro_bound}
\end{figure}

% -----------------------------------------------------------------------
\subsection{\'Enonc\'e et preuve du Trilemme CRO classique}
\label{subsec:enonce_trilemme}
% -----------------------------------------------------------------------
\begin{sloppypar}
Nous \'etablissons le r\'esultat principal du r\'egime classique.

\begin{remark}[Version de la preuve]
\label{rem:version_preuve_classique}
La version de la preuve pr\'esent\'ee ici diff\`ere de celle de l'article
publi\'e~\cite{MINKA2026CRO}. Les formulations alternatives retenues, issues
de la certification Lean4, sont d\'etaill\'ees \`a la
section~\ref{subsec:lean4_erreurs}.
\end{remark}

\begin{lemma}[Garantie de confidentialit\'e pour les protocoles r\'eguliers]
\label{lem:confidentiality_guarantee}
Pour tout protocole $\beta$-r\'egulier (Hypoth\`ese~\ref{ass:regularite}) :
\begin{equation}
\Conf(\mathcal{P}) \geq 1 - \sqrt{c_0 \beta \cdot I(M;\Sigma|C)}
\label{eq:lemma4}
\end{equation}
o\`u $c_0 = 1/(2\ln 2) \approx 0{,}721$ est la constante de Pinsker.
\end{lemma}

\begin{proof}
Par le Lemme~\ref{lem:compromis_conf_mi} (Pinsker) et
l'Hypoth\`ese~\ref{ass:regularite}, pour tous $m_0, m_1 \in
\mathcal{M}_\lambda$ et $c \in \mathcal{C}_\lambda$, on a
\begin{equation}
\TV(P_{\Sigma|m_0,c}, P_{\Sigma|m_1,c}) \leq \sqrt{c_0 \KL(P_{\Sigma|m_0,c} \| P_{\Sigma|m_1,c})} \leq \sqrt{c_0 \beta \cdot I(M;\Sigma|C)}
\end{equation}
En prenant le maximum sur $(m_0, m_1, c)$ et en appliquant la
D\'efinition~\ref{def:confidentialite_clo}, on obtient~\eqref{eq:lemma4}.
\end{proof}

\begin{remark}[Interpr\'etation du Lemme~\ref{lem:confidentiality_guarantee}]
\label{rem:interp_lemme4}
Le Lemme~\ref{lem:confidentiality_guarantee} \'enonce que la faible
information mutuelle conditionnelle, combin\'ee \`a la
$\beta$-r\'egularit\'e, \emph{garantit} la confidentialit\'e. Si
$I(M;\Sigma|C) \leq \kappa/\lambda$, alors
$\Conf(\mathcal{P}) \geq 1 - \sqrt{c_0 \beta \kappa/\lambda} \to 1$
quand $\lambda \to \infty$.
\end{remark}

\begin{remark}[Constante $c_0$]
\label{rem:constante_c0}
Avec des divergences en bits, la constante de Pinsker serr\'ee est
$(\ln 2)/2 \approx 0{,}347$. La constante $c_0 = 1/(2\ln 2) \approx
0{,}721$ utilis\'ee dans ce chapitre, reprise de~\cite{MINKA2026CRO}, est plus
grande. Toute in\'egalit\'e \'enonc\'ee avec $c_0$ reste donc valide sans \^etre
serr\'ee. Le qualificatif \og serr\'ee\fg{} employ\'e dans l'article publi\'e
s'entend en ce sens.
\end{remark}

\begin{theorem}[Trilemme CRO conditionnel classique]
\label{thm:trilemme_cro_conditionnel}
Sous les Hypoth\`eses~\ref{ass:extraction_bornee},~\ref{ass:verdicts_efficace},
\ref{ass:regularite} et~\ref{ass:regime}, pour tout protocole CLO avec
$\lambda \geq 64$, $H(M) = \lambda$, $M \perp C$~:
\begin{equation}
\Conf \cdot \Rel \cdot \Opp \leq \frac{\Delta}{\lambda} + O(\lambda^{-2}) + \negl(\lambda)
\label{eq:trilemme_cro_conditionnel}
\end{equation}
Plus pr\'ecis\'ement, avec $\gamma = \alpha H(\mathcal{C}) - 1$~:
\begin{equation}
\Conf \cdot \Rel \cdot \Opp \leq \frac{1}{\gamma}\left(\frac{\Delta}{\lambda}
+ \frac{\kappa(\Delta+1)}{\lambda^2} + \frac{\kappa^2}{\lambda^3}\right)
\label{eq:trilemme_cro_explicite}
\end{equation}
et, sous la forme la plus pr\'ecise,
\begin{equation}
\Conf \cdot \Rel \cdot \Opp \leq \frac{1+\kappa/\lambda}{\lambda}\,
  \min\left\{1,\; \frac{\kappa/\lambda + \Delta}{\gamma}\right\}.
\label{eq:trilemme_cro_raffine}
\end{equation}
De plus, $\Conf(\mathcal{P}) \geq 1 - \sqrt{c_0 \beta \kappa/\lambda}$
par le Lemme~\ref{lem:confidentiality_guarantee}.
\end{theorem}

La preuve proc\`ede en cinq \'etapes. Les deux premi\`eres construisent les bornes
individuelles sur $\Rel$ et $\Conf$ via les in\'egalit\'es classiques de Fano et de
Pinsker~\cite{Fano61,CoverThomas06,Pinsker1964,Tsybakov2009}~; les \'etapes trois et quatre bornent $\Opp$ par les hypoth\`eses
institutionnelles~\cite{Simon82,Miller56}~; la derni\`ere combine le produit et applique
le r\'egime institutionnel.

\begin{Preuve}[Preuve d\'etaill\'ee]\label{prv:cro_classique}

\begin{cmtEtape}{Vue d'ensemble de la preuve.}
Posons $x = I(M;\Sigma|C)$, l'information que la preuve r\'ev\`ele sur le
message une fois le contexte connu. Toute la d\'emonstration consiste \`a
exprimer les trois m\'etriques en fonction de cette \emph{unique} quantit\'e,
puis \`a constater qu'elles ne peuvent cro\^itre ensemble. La fiabilit\'e
exige $x$ grand~; la confidentialit\'e exige $x$ petit~; l'opposabilit\'e,
elle, est plafonn\'ee par ce que l'institution parvient \`a extraire. Elle se
d\'eroule en cinq \'etapes, o\`u (1)~Fano borne $\Rel$ par $x$~; (2)~Pinsker et la r\'egularit\'e
garantissent $\Conf$ \`a partir de $x$~; (3)~l'extraction born\'ee borne
$\Opp$~; (4)~la borne sur les verdicts ach\`eve cette contrainte~;
(5)~on forme le produit et on applique le r\'egime.
\end{cmtEtape}

    \paragraph{\'Etape 1. Fano contraint la fiabilit\'e.}
    L'outil mobilis\'e ici est l'in\'egalit\'e de Fano, introduite \`a la
    sous-sous-section~\ref{subsubsec:fano_classique}. Elle \'enonce que, si l'on
    cherche \`a deviner une variable al\'eatoire \`a partir d'une observation,
    l'incertitude qui subsiste sur cette variable est born\'ee par la
    probabilit\'e de se tromper. Appliquons-la au message.

    Soit $\hat{M}$ l'estimateur optimal de $M$ construit \`a partir du couple
    $(\Sigma, C)$, c'est-\`a-dire la meilleure reconstruction possible du
    message \`a partir de la preuve et du contexte. Notons
    $p_e = 1 - \Rel$ sa probabilit\'e d'erreur. Fano donne alors
    \begin{equation}
    H(M|\Sigma,C) \leq h(p_e) + p_e\log_2(2^\lambda - 1)
    \label{eq:fano_applique}
    \end{equation}
    \cmtFormule{L'incertitude r\'esiduelle sur $M$ est born\'ee par la probabilit\'e d'erreur.}
    \cmt{$h(\cdot)$ d\'esigne l'entropie binaire, $h(p) = -p\log_2 p - (1-p)\log_2(1-p)$.}
    \cmt{Le second terme compte le nombre de messages restants \`a d\'epartager en cas d'erreur, soit $2^\lambda - 1$ messages.}

    L'entropie binaire ne d\'epasse jamais un bit, $h(p_e) \leq 1$, tandis que
    $\log_2(2^\lambda - 1) \leq \lambda$. En rempla\c{c}ant dans
    \eqref{eq:fano_applique}, on obtient
    \begin{equation}
    H(M|\Sigma,C) \leq 1 + (1 - \Rel)\,\lambda
    \end{equation}
    \cmt{Aucune restriction du type $\Rel \geq 1/2$ n'est n\'ecessaire ici, car la majoration $h(p_e) \leq 1$ vaut pour toute valeur de $p_e$.}

    Les messages \'etant uniformes sur $\mathcal{M}_\lambda$, leur entropie
    vaut $H(M) = \lambda$ (sous-sous-section~\ref{subsubsec:entropie_info_mutuelle}).
    L'information mutuelle \'etant la diff\'erence entre l'entropie et
    l'entropie conditionnelle, il vient
    \begin{equation}
    I(M;\Sigma,C) = H(M) - H(M|\Sigma,C) \geq \lambda - 1 - (1-\Rel)\lambda = \Rel \cdot \lambda - 1
    \label{eq:fano_mi}
    \end{equation}
    \cmtFormule{L'information extraite sur $M$ vaut au moins $\Rel\lambda - 1$.}

    Le protocole est con\c{c}u de sorte que le message soit ind\'ependant du
    contexte institutionnel, $M \perp C$~; autrement dit, savoir dans quelle juridiction on
    se trouve n'apprend rien sur le message transmis. Par cons\'equent
    $I(M;\Sigma,C) = I(M;\Sigma|C) = x$, d'o\`u
    \begin{equation}
    \boxed{\;\Rel \leq \frac{x + 1}{\lambda}\;}
    \label{eq:etape1_rel}
    \end{equation}
    \cmt{Lecture. Une preuve qui ne r\'ev\`ele presque rien sur le message ($x$ petit) ne permet pas de le reconstituer, donc la fiabilit\'e chute.}

    \paragraph{\'Etape 2. Pinsker et r\'egularit\'e garantissent la confidentialit\'e.}
    On mobilise cette fois l'in\'egalit\'e de Pinsker
    (sous-sous-section~\ref{subsubsec:pinsker_classique}), qui relie deux
    fa\c{c}ons de mesurer l'\'ecart entre deux lois. La premi\`ere est la distance
    en variation totale, qui gouverne la capacit\'e d'un adversaire \`a les
    distinguer~; la seconde est la divergence de Kullback-Leibler
    (sous-sous-section~\ref{subsubsec:distances_divergences}).

    Pour tout couple de messages $(m_0, m_1)$ et tout contexte $c$, Pinsker
    donne $\TV^2 \leq c_0 \KL$ avec $c_0 = 1/(2\ln 2) \approx 0{,}721$, tandis que
    l'Hypoth\`ese~\ref{ass:regularite} de $\beta$-r\'egularit\'e majore cette
    divergence par $\beta\,x$. En composant les deux, puis en passant au
    maximum sur les couples et les contextes, on obtient le
    Lemme~\ref{lem:confidentiality_guarantee}, soit
    \begin{equation}
    \boxed{\;\Conf \geq 1 - \sqrt{c_0 \beta x}\;}
    \label{eq:etape2_conf}
    \end{equation}
    \cmtFormule{Confidentialit\'e minor\'ee par une fonction d\'ecroissante de $x$.}
    \cmt{Il s'agit d'une \emph{garantie}, puisque tant que la preuve divulgue peu ($x \leq \kappa/\lambda$), la confidentialit\'e tend vers $1$.}
    \cmt{Cette \'etape ne sert pas \`a majorer le produit, mais \`a montrer que le r\'egime consid\'er\'e est bien celui o\`u la confidentialit\'e est pr\'eserv\'ee.}

    \paragraph{\'Etape 3. L'extraction born\'ee contraint l'opposabilit\'e.}
    Par d\'efinition, l'opposabilit\'e est la fraction de l'incertitude
    institutionnelle que la preuve r\'esout,
    $\Opp = I(V;\Sigma|C)/H(V|C)$. Son num\'erateur mesure ce que la preuve
    apporte au verdict, son d\'enominateur l'ampleur de la d\'ecision \`a
    prendre.

    L'Hypoth\`ese~\ref{ass:extraction_bornee} borne ce num\'erateur,
    puisque l'institution ne tire de la preuve que ce que celle-ci contient sur le
    message, $x$, augment\'e d'au plus $\Delta$ bits qu'elle extrait par ses
    propres moyens (pr\'ec\'edents, expertise, proc\'edures). Il vient
    \begin{equation}
    \Opp = \frac{I(V;\Sigma|C)}{H(V|C)} \leq \frac{x + \Delta}{H(V|C)}
    \label{eq:etape3_opp}
    \end{equation}
    \cmtFormule{Le num\'erateur est major\'e par $x + \Delta$.}
    \cmt{$\Delta$ est une capacit\'e cognitive, born\'ee par la rationalit\'e de l'interpr\`ete~; elle ne cro\^it pas avec le param\`etre de s\'ecurit\'e.}

    \paragraph{\'Etape 4. La borne sur les verdicts compl\`ete la contrainte sur $\Opp$.}
    Il reste \`a minorer le d\'enominateur. L'Hypoth\`ese~\ref{ass:verdicts_efficace}
    \'enonce qu'un verdict institutionnel exploite effectivement le contexte
    dont il dispose, ce qui s'\'ecrit
    \begin{equation}
    H(V|C) \geq \gamma := \alpha H(\mathcal{C}) - 1 \geq 1
    \label{eq:etape4_hvc}
    \end{equation}
    \cmtFormule{L'incertitude sur les verdicts est minor\'ee par $\gamma$.}
    \cmt{Un verdict d\'eg\'en\'er\'e, toujours identique, aurait $H(V|C) = 0$ et rendrait $\Opp$ ind\'efini, cas que l'hypoth\`ese \'ecarte.}

    En injectant \eqref{eq:etape4_hvc} dans \eqref{eq:etape3_opp}, on obtient
    \begin{equation}
    \boxed{\;\Opp \leq \frac{x + \Delta}{\gamma}\;}
    \label{eq:etape4_opp}
    \end{equation}

    \paragraph{\'Etape 5. Formation du produit et application du r\'egime.}
    Les trois m\'etriques sont d\'esormais exprim\'ees en fonction de $x$. En
    multipliant \eqref{eq:etape1_rel} et \eqref{eq:etape4_opp} puis en majorant
    trivialement $\Conf \leq 1$, on obtient
    \begin{equation}
    \Conf \cdot \Rel \cdot \Opp \leq \frac{x+1}{\lambda} \cdot \frac{x+\Delta}{\gamma}
    \label{eq:produit_brut}
    \end{equation}
    \cmt{On majore ici $\Conf$ par $1$ sans recourir \`a \eqref{eq:etape2_conf}, car dans le r\'egime consid\'er\'e la confidentialit\'e est proche de $1$ et la contrainte vient des deux autres facteurs.}

    Le membre de droite, produit de deux fonctions croissantes et
    positives de $x$, est donc croissant en $x$. Il suffit ainsi de
    l'\'evaluer \`a la borne sup\'erieure du r\'egime. Par
    l'Hypoth\`ese~\ref{ass:regime}, $x \leq \kappa/\lambda$, d'o\`u
    \begin{equation}
    \Conf \cdot \Rel \cdot \Opp \leq \frac{1}{\gamma}\left(\frac{\kappa/\lambda + 1}{\lambda}\right)\left(\frac{\kappa}{\lambda} + \Delta\right)
    \end{equation}
    \cmt{Substitution directe de $x = \kappa/\lambda$ dans \eqref{eq:produit_brut}.}

    Le d\'eveloppement du produit donne quatre termes, que l'on ordonne par
    puissance d\'ecroissante de $\lambda$~:
    \begin{align}
    \left(\frac{\kappa/\lambda + 1}{\lambda}\right)\left(\frac{\kappa}{\lambda} + \Delta\right)
    &= \frac{\kappa^2}{\lambda^3} + \frac{\kappa\Delta}{\lambda^2}
     + \frac{\kappa}{\lambda^2} + \frac{\Delta}{\lambda} \nonumber \\
    &= \frac{\Delta}{\lambda} + \frac{\kappa(\Delta + 1)}{\lambda^2} + \frac{\kappa^2}{\lambda^3}
    \end{align}
    \cmtFormule{Les deux termes crois\'es en $\kappa$ se regroupent en $\kappa(\Delta+1)/\lambda^2$.}

    On obtient ainsi la forme explicite annonc\'ee
    \begin{equation}
    \Conf \cdot \Rel \cdot \Opp \leq \frac{1}{\gamma}\left(\frac{\Delta}{\lambda} + \frac{\kappa(\Delta+1)}{\lambda^2} + \frac{\kappa^2}{\lambda^3}\right)
    \end{equation}

    Comme $\gamma \geq 1$, diviser par $\gamma$ ne fait que renforcer la
    majoration. Les deux derniers termes \'etant en $O(\lambda^{-2})$, il vient
    \begin{equation}
    \Conf \cdot \Rel \cdot \Opp \leq \frac{\Delta}{\lambda} + O(\lambda^{-2}) + \negl(\lambda)
    \end{equation}
    \cmtFormule{Le terme dominant est $\Delta/\lambda$.}
    \cmt{Lecture finale. Tant que $\Delta < \gamma$, la contrainte qui plafonne le produit est la capacit\'e institutionnelle~; au-del\`a, c'est la borne de Fano (Remarque~\ref{rem:activite_delta}).}

    \textbf{Forme raffin\'ee.} Comme $\Opp \leq 1$
    (Th\'eor\`eme~\ref{thm:bornes_mesures}), l'\'Etape~2 et l'Hypoth\`ese~\ref{ass:regime}
    donnent aussi $\Opp \leq \min\{1, (\kappa/\lambda+\Delta)/\gamma\}$~;
    combin\'ee \`a l'\'Etape~1 sous l'Hypoth\`ese~\ref{ass:regime},
    ($\Rel \leq (1+\kappa/\lambda)/\lambda$) ainsi qu'\`a $\Conf \leq 1$, elle
    donne~\eqref{eq:trilemme_cro_raffine}. Comme le minimum est major\'e par
    son second argument, la forme raffin\'ee implique la forme explicite.

\end{Preuve}

\begin{remark}[Quand la capacit\'e institutionnelle est-elle contraignante~?]
\label{rem:activite_delta}
Le facteur $\min\{1,(\kappa/\lambda+\Delta)/\gamma\}$ est strictement
inf\'erieur \`a~$1$ si et seulement si
\begin{equation}
\Delta < \gamma - \frac{\kappa}{\lambda}
\;\Longleftrightarrow\;
H(\mathcal{C}) > H^\sharp := \frac{\Delta + 1 + \kappa/\lambda}{\alpha}.
\label{eq:seuil_activite}
\end{equation}
Dans ce r\'egime (\emph{contexte riche relativement \`a la capacit\'e}), la
borne est port\'ee par le rapport $\Delta/\gamma$ et la capacit\'e
institutionnelle est la contrainte liante. Dans le r\'egime compl\'ementaire
($H(\mathcal{C}) \leq H^\sharp$), le produit est born\'e par la seule
in\'egalit\'e de Fano, $\Relrec \leq (1+\kappa/\lambda)/\lambda$~; la
forme~\eqref{eq:trilemme_cro_explicite} n'apporte alors pas d'information
suppl\'ementaire. Si le verdict \'etait une fonction d\'eterministe de
$(\Sigma,C)$, les Hypoth\`eses~\ref{ass:extraction_bornee}
et~\ref{ass:verdicts_efficace} imposeraient $\Delta \geq \gamma - I(M;\Sigma|C)$,
de sorte qu'on serait toujours dans le second r\'egime
(Remarque~\ref{rem:verdict_deterministe}). Ces faits sont certifi\'es
en Lean (\texttt{trilemma\_refined}, \texttt{refined\_bound\_is\_fano\_when\_delta\_ge\_gamma},
\texttt{deterministic\_forces\_capacity}, section~\ref{sec:lean4}).
\end{remark}

\begin{remark}[Port\'ee du Th\'eor\`eme~\ref{thm:trilemme_cro_conditionnel}]
\label{rem:scope_trilemme_classique}
L'Hypoth\`ese~\ref{ass:regime} est n\'ecessaire. Sans borne sur
$I(M;\Sigma|C)$, le produit $\Conf \cdot \Rel \cdot \Opp$ n'a pas besoin
de d\'ecro\^itre avec $\lambda$. En effet, le protocole qui r\'ev\`ele $M$ avec
probabilit\'e $1/2$ et produit sinon une preuve uniform\'ement
al\'eatoire est $\beta$-r\'egulier avec $\beta \leq 1{,}04$ pour tout
$\lambda \geq 64$, satisfait les
Hypoth\`eses~\ref{ass:extraction_bornee}--\ref{ass:regularite} et a
$\Conf = 1/2$, $\Rel = 1/2 + 2^{-\lambda-1}$ et, avec un contexte $C$
form\'e de deux bits uniformes ind\'ependants de $(M,\Sigma)$, $\alpha = 1$
(donc $\gamma = \alpha H(C) - 1 = 1$) et $V$ le premier bit de $\Sigma$
(donc $H(V|C) = 1 \geq \gamma$ et $\log_2|\mathcal{V}_\lambda| = 1 \leq \alpha H(C)$),
$\Opp = 1$, pour tout $\lambda$. L'Hypoth\`ese~\ref{ass:extraction_bornee} est
satisfaite avec $\Delta = 0$, puisque $I(V;\Sigma|C) = 1 \leq I(M;\Sigma|C)$. Ici
$I(M;\Sigma|C) = \lambda/2 - 1 + O(\lambda 2^{-\lambda})$ se situe hors
du r\'egime~\eqref{eq:regime}.
\end{remark}

\begin{remark}[Statut de v\'erification]
\label{rem:statut_verification_classique}
Le noyau alg\'ebrique de cette preuve est certifi\'e m\'ecaniquement en
Lean4 \`a la section~\ref{sec:lean4}. Les in\'egalit\'es de Fano et de
Pinsker, ainsi que les hypoth\`eses institutionnelles, y sont
formalis\'ees comme propri\'et\'es des instances de protocole. Le
d\'etail des formulations alternatives issues de la certification est pr\'esent\'e \`a la
section~\ref{subsec:lean4_erreurs}.
\end{remark}

\begin{corollary}[Forme simplifi\'ee]\label{cor:forme_simplifiee}

Pour $\lambda \geq 128$, $\beta = O(1)$, $\gamma = O(1)$~\cite{Goldreich01,KatzLindell14}~:
\begin{equation}
\Conf \cdot \Rel \cdot \Opp \lesssim \frac{\Delta}{\lambda}
\label{eq:trilemme_simplifie}
\end{equation}

\end{corollary}

Pour $\Delta = 5$ bits, $\alpha = 1$, $H(\mathcal{C}) = 20$ bits
($\gamma = 19$), $\kappa = 1$ et $\lambda = 128$, la
forme~\eqref{eq:trilemme_cro_raffine} donne
$\frac{1+1/128}{128}\cdot\frac{1/128 + 5}{19} \approx 2{,}1\cdot10^{-3}$,
contre $\approx 7{,}9\cdot10^{-3}$ pour la borne de Fano seule, de sorte que la capacit\'e
institutionnelle divise la borne par pr\`es de quatre. Pour $\Delta = 35$
bits dans le m\^eme contexte, $\Delta > \gamma$ et la borne retombe sur
celle de Fano ($\approx 7{,}9\cdot10^{-3}$)~; la forme publi\'ee
$\Delta/\lambda \approx 0{,}27$ est alors tr\`es l\^ache. Ces valeurs ont
statut d'illustration num\'erique~; elles sont calcul\'ees sous les hypoth\`eses
institutionnelles et ne constituent pas des mesures empiriques~\cite{Tsybakov2009}.

\textbf{Augmenter $\lambda$ n'am\'eliore pas le produit CRO si $\Delta$ reste fixe.}
\`A hypoth\`eses fix\'ees, la borne d\'ecro\^it en $1/\lambda$ parce que
l'Hypoth\`ese~\ref{ass:regime} force la fiabilit\'e de reconstruction \`a
d\'ecro\^itre~; la capacit\'e institutionnelle d\'etermine, dans le
r\'egime~\eqref{eq:seuil_activite}, la fraction de ce plafond que
l'institution peut exploiter. C'est l\`a que les limites de l'interpr\`ete, au-del\`a
de la seule puissance des adversaires, deviennent la contrainte liante.

\begin{boiteArbitre}
L'arbitre comprend que commander un algorithme plus fort ne changera rien,
puisque le goulot d'\'etranglement se situe en lui d\`es que ses verdicts
sont plus fins que ce qu'il sait extraire.
\end{boiteArbitre}

\begin{figure}[htbp]
\centering
\begin{tikzpicture}[
    scale=0.8,
    transform shape,
    every node/.style={font=\small}
]

\node[circle, draw=enspblue,  fill=enspblue!10,  minimum size=1.5cm, thick]
    (C) at ( 0, 5) {\textbf{Conf}};
\node[circle, draw=enspgreen, fill=enspgreen!10, minimum size=1.5cm, thick]
    (R) at (-4, 0) {\textbf{Rel}};
\node[circle, draw=navyblue,  fill=navyblue!10,  minimum size=1.5cm, thick]
    (O) at ( 4, 0) {\textbf{Opp}};

\draw[thick, enspblue!50]  (C) -- (R);
\draw[thick, enspgreen!50] (R) -- (O);
\draw[thick, navyblue!50]  (O) -- (C);

\node[draw=enspred, fill=enspred!8, rounded corners, thick,
      minimum width=3cm, minimum height=1.2cm] (bound) at (0, 2)
    {$\Conf \cdot \Rel \cdot \Opp \leq \dfrac{\Delta}{\lambda}$};

\draw[-{Stealth[length=2mm]}, thick, enspblue]  (C.south) -- (bound.north west);
\draw[-{Stealth[length=2mm]}, thick, enspgreen] (R.east)  -- (bound.south west);
\draw[-{Stealth[length=2mm]}, thick, navyblue]  (O.west)  -- (bound.south east);

\node[rotate=63,  font=\footnotesize, text=enspblue]  at (-2.3, 2.2)
    {Contrainte $\Delta$};
\node[rotate=-63, font=\footnotesize, text=navyblue]  at ( 2.3, 2.2)
    {Goulot $H(C)$};
\node[font=\footnotesize, text=enspgreen]              at (0,   -0.5)
    {Borne de Fano};

\node[above=0.3cm of C, font=\footnotesize, text=enspblue, align=center]
    {Protection\\de la vie priv\'ee};
\node[below left=0.3cm of R, font=\footnotesize, text=enspgreen, align=right]
    {Recouvrement\\du message};
\node[below right=0.3cm of O, font=\footnotesize, text=navyblue, align=left]
    {Interpr\'etation\\institutionnelle};

\end{tikzpicture}
\caption[Triangle du Trilemme CRO]{%
  Triangle du Trilemme CRO : trois propri\'et\'es de s\'ecurit\'e concurrentes.%
}\label{fig:trilemma_triangle}

\smallskip
\begin{minipage}{\linewidth}
  \small\textit{Note :} La borne centrale
  $\Conf \cdot \Rel \cdot \Opp \leq \Delta/\lambda$ limite l'optimisation simultan\'ee.
  Les \'etiquettes des ar\^etes indiquent le m\'ecanisme de contrainte principal
  reliant chaque paire de propri\'et\'es.
\end{minipage}

\end{figure}
\end{sloppypar}

% -----------------------------------------------------------------------
\subsection{Port\'ee du th\'eor\`eme et syst\`emes d\'eploy\'es}
\label{subsec:portee_trilemme}
% -----------------------------------------------------------------------
\begin{remark}[Port\'ee des \'etudes de cas]\label{rem:portee_cas_etude}
Sous l'Hypoth\`ese~\ref{ass:regime}, l'\'Etape~1 donne
$\Relrec \leq (1+\kappa/\lambda)/\lambda \approx 0{,}008$ \`a $\lambda = 128$.
Les syst\`emes d\'eploy\'es de r\'ef\'erence (Helios, zk-SNARKs, ECDSA),
dont la fiabilit\'e op\'erationnelle $\Relver$ est proche de~$1$, ne sont donc
\emph{pas} des instances du Th\'eor\`eme~\ref{thm:trilemme_cro_conditionnel},
car leurs m\'etriques CLO sont exprim\'ees sur l'\'echelle op\'erationnelle
de~\cite{MINKA2026CLO}, distincte des quantit\'es information-th\'eoriques des
D\'efinitions~\ref{def:confidentialite_clo}--\ref{def:opposabilite_clo}. Dans
la suite, ces syst\`emes servent \`a une \emph{calibration illustrative} des
ordres de grandeur des param\`etres institutionnels, non \`a une
validation de la borne. Le th\'eor\`eme s'applique aux protocoles \`a fuite
contr\^ol\'ee ($I(M;\Sigma|C) \leq \kappa/\lambda$), dont les engagements
parfaitement masquants et les preuves \`a divulgation nulle de connaissance
sont les exemples canoniques.
\end{remark}

% -----------------------------------------------------------------------
\subsection{R\'egimes op\'erationnels}
\label{subsec:regimes_operationnels}
% -----------------------------------------------------------------------
\begin{sloppypar}
Le Trilemme CRO classique distingue deux r\'egimes, s\'epar\'es par le seuil
$H^\sharp = (\Delta + 1 + \kappa/\lambda)/\alpha$ de la
Remarque~\ref{rem:activite_delta}. Ce seuil, d\'eriv\'e du th\'eor\`eme,
compare la richesse du contexte, pond\'er\'ee par la granularit\'e des
verdicts, \`a la capacit\'e d'extraction de l'institution.

\subsubsection{R\'egime \`a contexte rare~: $H(\mathcal{C}) \leq H^\sharp$}
\label{subsubsec:regime_faible_contexte}

La capacit\'e institutionnelle couvre la richesse du contexte~; le produit est
born\'e par l'in\'egalit\'e de Fano, $(1+\kappa/\lambda)/\lambda$. Avec les
param\`etres illustratifs de la calibration (Tableau~\ref{tab:validation_cro_classique}),
un contrat standardis\'e ($H(\mathcal{C}) \approx 20$ bits, $\alpha = 1$,
$\Delta \approx 35$ bits pour un juge g\'en\'eraliste~\cite{Mason07,Kerr04}, soit
$H^\sharp \approx 36$) et un scrutin de taille moyenne ($H(\mathcal{C}) \approx 40$
bits, $\alpha = 1{,}2$, $\Delta \approx 56$ bits pour un comit\'e de garants
sp\'ecialis\'es~\cite{Goodman15}, soit $H^\sharp \approx 48$) rel\`event de ce
r\'egime. La capacit\'e d'extraction n'y est pas la contrainte liante.

\subsubsection{R\'egime \`a contexte riche~: $H(\mathcal{C}) > H^\sharp$}
\label{subsubsec:regime_haute_contexte}
\label{subsubsec:regime_transition}

La granularit\'e des verdicts d\'epasse la capacit\'e d'extraction~; le produit
est born\'e par
\begin{equation}
\Conf \cdot \Rel \cdot \Opp \leq
\frac{1+\kappa/\lambda}{\lambda}\cdot\frac{\kappa/\lambda+\Delta}{\gamma},
\end{equation}
strictement plus petit que la borne de Fano. La conformit\'e financi\`ere via
zk-SNARKs~\cite{Groth16,BCTV14} en est l'exemple, avec $H(\mathcal{C}) \approx 70$
bits, $\alpha = 1{,}5$, $\Delta \approx 73$ bits pour un r\'egulateur
expert~\cite{BaselIII}, soit $H^\sharp \approx 49$ et $\gamma = 104$~; la borne
vaut environ $0{,}70$ fois la borne de Fano. C'est dans ce r\'egime que la
capacit\'e cognitive de l'institution, au-del\`a de la seule structure
cryptographique, plafonne le produit.

% -----------------------------------------------------------------------
\begin{table}[htbp]
\centering
\caption{Les deux r\'egimes op\'erationnels du Trilemme CRO classique}
\label{tab:regimes_operationnels}
\begin{tabular}{@{}lcc@{}}
\toprule
\textbf{Crit\`ere} & \textbf{Contexte rare} & \textbf{Contexte riche} \\
\midrule
Condition & $H(\mathcal{C}) \leq H^\sharp$ & $H(\mathcal{C}) > H^\sharp$ \\
Borne & $(1+\kappa/\lambda)/\lambda$ (Fano) & $\frac{1+\kappa/\lambda}{\lambda}\cdot\frac{\kappa/\lambda+\Delta}{\gamma}$ \\
Contrainte liante & fiabilit\'e de reconstruction & capacit\'e institutionnelle $\Delta$ \\
Construction de r\'ef\'erence & SOTP-Low / SOTP-Base & SOTP-High \\
Exemple (calibration) & ECDSA, Helios & zk-SNARKs \\
\bottomrule
\end{tabular}
\begin{tablenotes}
\item \footnotesize $H^\sharp = (\Delta + 1 + \kappa/\lambda)/\alpha$ (Remarque~\ref{rem:activite_delta}). Les exemples sont class\'es avec les param\`etres de calibration~; ils ne sont pas des instances du th\'eor\`eme (Remarque~\ref{rem:portee_cas_etude}).
\item \footnotesize Les constructions SOTP sont des constructions de r\'ef\'erence (chapitre~\ref{chap:constructions})~; leur optimalit\'e est une conjecture.
\end{tablenotes}
\end{table}
\end{sloppypar}

% -----------------------------------------------------------------------
\subsection{Calibration illustrative des param\`etres institutionnels}
\label{subsec:analyse_comparative}
% -----------------------------------------------------------------------
\begin{sloppypar}
Les syst\`emes d\'eploy\'es n'\'etant pas des instances du th\'eor\`eme
(Remarque~\ref{rem:portee_cas_etude}), leurs produits CLO ne peuvent ni
valider ni r\'efuter la borne. Ils permettent en revanche de calibrer, sur
l'\'echelle op\'erationnelle CLO, des ordres de grandeur plausibles de la
capacit\'e $\Delta$~\cite{Adida08,Groth16,Johnson01}.

% -----------------------------------------------------------------------
\begin{table}[htbp]
\centering
\caption{Calibration illustrative~: capacit\'es implicites sur l'\'echelle CLO}
\label{tab:validation_cro_classique}
\begin{tabular}{@{}lcccc@{}}
\toprule
\textbf{Syst\`eme} & $\boldsymbol{\Conf\cdot\Relver\cdot\Opp}$ & $\boldsymbol{\lambda}$ & $\boldsymbol{\Delta := \lambda \cdot \text{produit}}$ & $\boldsymbol{\Delta/\lambda}$ \\
\midrule
Helios~\cite{Adida08,Adida09UCL} & 0{,}437 & 128 & $\approx 56$ bits & 0{,}437 \\
zk-SNARKs~\cite{Groth16,BCTV14} & 0{,}572 & 128 & $\approx 73$ bits & 0{,}572 \\
ECDSA~\cite{Johnson01,NIST.FIPS.186-5} & 0{,}273 & 128 & $\approx 35$ bits & 0{,}273 \\
\bottomrule
\end{tabular}
\begin{tablenotes}
\item \footnotesize Les capacit\'es $\Delta$ sont \emph{d\'efinies} par $\Delta := \lambda \cdot (\Conf\cdot\Relver\cdot\Opp)$, c'est-\`a-dire par inversion de la forme publi\'ee $\Delta/\lambda$ sur l'\'echelle CLO~; les colonnes $\Delta/\lambda$ et produit co\"incident donc par construction.
\item \footnotesize Ces syst\`emes ne sont pas des instances du Th\'eor\`eme~\ref{thm:trilemme_cro_conditionnel} (Remarque~\ref{rem:portee_cas_etude}). Voir la note m\'ethodologique de la sous-section~\ref{subsec:note_methodo}.
\end{tablenotes}
\end{table}
% -----------------------------------------------------------------------

Chaque capacit\'e implicite peut \^etre rapproch\'ee d'une d\'ecomposition en
ressources cognitives et organisationnelles, d\'etaill\'ee au
chapitre~\ref{chap:experimentation}
(Tableau~\ref{tab:decomposition_delta})~:

\begin{itemize}
\item \textbf{Helios}~\cite{Adida08,Goodman15}, avec $\Delta \approx 56$ bits pour un comit\'e
\'electoral sp\'ecialis\'e (expertise, pr\'ec\'edents, proc\'edures, collaboration)~\cite{Ericsson95,Simon82}.

\item \textbf{zk-SNARKs}~\cite{Groth16,BCTV14}, avec $\Delta \approx 73$ bits pour une institution
financi\`ere experte (cadres r\'eglementaires, analyse de sch\'emas, \'equipes
forensiques, collaboration multi-agences)~\cite{BaselIII}.

\item \textbf{ECDSA}~\cite{Johnson01,Mason07}, avec $\Delta \approx 35$ bits pour un juge de
juridiction g\'en\'erale (pr\'ec\'edents, proc\'edures, expertise g\'en\'erale)~\cite{Kerr04,Casey11}.
\end{itemize}

Les capacit\'es implicites ainsi calibr\'ees ($35$--$73$ bits) tombent dans
l'intervalle de quelques dizaines de bits sugg\'er\'e par la litt\'erature sur la
rationalit\'e born\'ee~\cite{Miller56,Simon82,Ericsson95,Cowan2001}. Il s'agit
d'une v\'erification de plausibilit\'e des ordres de grandeur, distincte d'une validation du th\'eor\`eme. Cette calibration est \textbf{r\'etrospective et
illustrative}~; elle ne constitue pas une preuve empirique des hypoth\`eses
institutionnelles~\cite{Tsybakov2009,Kahneman11}.
\end{sloppypar}

% -----------------------------------------------------------------------
\subsection{Positionnement par rapport au r\'egime quantique}
\label{subsec:transition_postquantique}
% -----------------------------------------------------------------------
\begin{sloppypar}
Le Trilemme CRO classique \'etabli dans ce chapitre s'applique aux adversaires
computationnels classiques (PPT)~\cite{Goldreich01,KatzLindell14}. La migration vers les standards post-quantiques
NIST soul\`eve une question critique, celle de savoir si ces bornes tiennent face \`a des
adversaires quantiques~\cite{NISTPQC,Watrous2009,Boneh2011,Zhandry}.

La section~\ref{sec:cro_pq} \'etendra ce cadre en introduisant~:

\begin{itemize}
\item La capacit\'e d'extraction quantique $\Delta_q = \Delta + \delta_{\text{enh}} - \delta_{\text{degr}}$
combinaison de la capacit\'e classique et des modifications induites par les
techniques quantiques~\cite{NielsenChuang2010,Wilde2017}.

\item Le coefficient de perte d'interpr\'etabilit\'e quantique $\eta_q$ quantifiant
la d\'egradation s\'emantique dans les canaux quantiques (d\'ecoh\'erence, erreurs de portes)~\cite{Zurek2003,Wootters1982}.

\item La borne post-quantique~:
\begin{equation}
\Conf^q \cdot \Rel^q \cdot \Opp^q \leq \frac{\Delta_q + \eta_q \cdot H(C)}{\lambda} + O(\lambda^{-2})
\end{equation}
\end{itemize}

La borne explicite restera port\'ee par la capacit\'e quantique $\Delta_q$ et par
$\gamma_q$~; le terme $\eta_q \cdot H(C)/\lambda$ n'appara\^it que dans une
majoration affaiblie, commode pour les ordres de grandeur. Le bruit quantique
agit comme une r\'eduction de l'information interpr\'etable, dont une
formalisation par contraction est propos\'ee en perspective
(chapitre~\ref{chap:resultats_discussions}).
\end{sloppypar}

% -----------------------------------------------------------------------
\begin{boiteArbitre}
L'arbitre referme le dossier du Trilemme CRO classique et termine ainsi l'\'etape \textbf{B}. Cette section lui a
appris un r\'esultat contre-intuitif, puisque la limite se trouve en lui plut\^ot
que dans la cryptographie d\`es que ses verdicts sont fins.

Quatre contributions structurent ce qu'il retient. Le mod\`ele \'etendu
d'abord, avec trois param\`etres mesurables ($\Delta$, $\alpha$, $\beta$) qui
traduisent ses contraintes cognitives en grandeurs quantifiables. Les
hypoth\`eses ensuite, explicites et honn\^etes sur ce qu'elles ne prouvent pas
encore. Le th\'eor\`eme lui-m\^eme, $\Conf \cdot \Rel \cdot \Opp \leq
\Delta/(\gamma\lambda) + O(\lambda^{-2})$, puis sa forme raffin\'ee, qui lui dit
exactement quand sa capacit\'e d'extraction devient la contrainte liante. La
distinction des deux r\'egimes enfin, qui transforme ce r\'esultat n\'egatif en
guide de conception.

Il est surtout marqu\'e par la franchise sur la port\'ee. Helios, les
zk-SNARKs et ECDSA ne sont pas des instances du th\'eor\`eme et calibrent
seulement les ordres de grandeur~; les capacit\'es $\Delta$ qui en
d\'ecoulent ($35$ \`a $73$ bits) restent plausibles au regard de la
litt\'erature cognitive.

La borne est n\'ecessaire sans \^etre encore prouv\'ee serr\'ee de mani\`ere
g\'en\'erale. La section suivante l'\'etendra au r\'egime quantique.

Il veut voir si les constructions tiennent ce que le th\'eor\`eme promet.
\end{boiteArbitre}
% ============================================================================
% SECTION 2.4 : EXTENSION QUANTIQUE DU TRILEMME CRO CONDITIONNEL
% Source : partie2/cro_pq.tex
% Version : 2.2 — alignée sur la formalisation Lean4 (section~\ref{sec:lean4})
% ============================================================================

\section{EXTENSION QUANTIQUE DU TRILEMME CRO CONDITIONNEL}
\label{sec:cro_pq}
\begin{sloppypar}
La section~\ref{sec:cro_classique} a \'etabli le Trilemme CRO conditionnel
classique~\cite{MINKA2026CRO}, selon lequel, sous les hypoth\`eses de capacit\'e institutionnelle,
$\Conf \cdot \Rel \cdot \Opp \leq \Delta/(\gamma\lambda) + O(\lambda^{-2})$, avec une
forme raffin\'ee qui identifie le r\'egime o\`u la capacit\'e institutionnelle est la
contrainte liante. L'Objectif~\ref{obj:obj3_postquantique} reste ouvert, car il reste \`a savoir si ces bornes
tiennent face aux adversaires quantiques.

Nous \'etendons ici le cadre classique au r\'egime quantique.
Ce r\'esultat, \`a notre connaissance la premi\`ere caract\'erisation quantitative des limites
de la cryptographie institutionnelle sous adversaires quantiques~\cite{NISTPQC,Zhandry},
\'etablit une borne explicite de m\^eme forme que dans le cas classique, en
$\Delta_q$ et $\gamma_q$ (Th\'eor\`eme~\ref{thm:trilemme_pq}), dont d\'ecoule la
majoration affaiblie publi\'ee~:
\begin{equation}
\Conf^q \cdot \Rel^q \cdot \Opp^q \leq \frac{\Delta_q + \eta_q \cdot H(C)}{\lambda}
+ \mathcal{O}(\lambda^{-2})
\label{eq:cro_pq_borne_principale}
\end{equation}
o\`u $\Delta_q$ repr\'esente la capacit\'e d'extraction institutionnelle
consciente du quantique, $\eta_q$ le coefficient de perte d'interpr\'etabilit\'e
quantique et $H(C)$ l'entropie contextuelle. Le r\'esultat central tient \`a ce que la
structure du trilemme survit au passage quantique, puisque la borne explicite garde la m\^eme
forme que dans le cas classique. Le terme $\eta_q H(C)/\lambda$ n'appara\^it que
dans la majoration affaiblie et mesure l'effet additionnel du bruit sur
l'information interpr\'etable.

Deux contributions additionnelles compl\`etent cette section. Une calibration
illustrative (sous-section~\ref{subsec:analyse_coherence_pq}) situe les ordres de
grandeur des param\`etres quantiques \`a partir de six syst\`emes (trois classiques issus de la
section~\ref{sec:cro_classique} et trois sch\'emas standardis\'es par le NIST ou
associ\'es~\cite{Dilithium2022,Kyber2022,SPHINCS2022}), qui ne sont pas des
instances du th\'eor\`eme (Remarque~\ref{rem:portee_cas_etude}), conform\'ement
\`a la note de m\'ethodologie pr\'esent\'ee dans la sous-section~\ref{subsec:note_methodo}.
Les implications pour les standards NIST PQC
(sous-section~\ref{subsec:implications_standards}) fournissent des
recommandations concr\`etes pour la migration institutionnelle.

Cette section correspond \`a l'\'etape \textbf{E (\'Etendre)} de la progression
F$\to$B$\to$\textbf{E}$\to$S$\to$C$\to$X et apporte la contribution \textbf{C3} de
cette th\`ese.

La r\'eponse \`a l'Objectif~\ref{obj:obj3_postquantique} est affirmative puisque les
bornes CRO s'\'etendent au r\'egime quantique, param\'etr\'ees par la capacit\'e
$\Delta_q$ et la perte $\eta_q$ (D\'efinition~\ref{def:perte_interp_quantique}).
\end{sloppypar}

\begin{boiteArbitre}
L'arbitre de demain travaillera avec des signatures ML-DSA ($3\,309$
octets pour ML-DSA-65, l\`a o\`u ECDSA en produit~$64$) et avec des preuves sur r\'eseaux
euclidiens l\`a o\`u il v\'erifie un logarithme discret. Sa capacit\'e
$\Delta_q$ risque de diminuer pendant que la complexit\'e des preuves augmente.
Cette section d\'emontre que la borne du Trilemme CRO r\'esiste face aux
adversaires quantiques et que sa forme est inchang\'ee~; ce sont les
param\`etres institutionnels $\Delta_q$ et $\eta_q$ qui portent l'effet de la
transition. La contrainte cognitive persiste donc dans le monde quantique au lieu
de relever d'un artefact classique.
\end{boiteArbitre}

% -----------------------------------------------------------------------
\subsection{Pr\'eliminaires : th\'eorie de l'information quantique}
\label{subsec:preliminaires_quantiques}
% -----------------------------------------------------------------------
\begin{sloppypar}
Nous \'etendons les notations classiques de th\'eorie de l'information du
Trilemme CRO conditionnel aux param\`etres quantiques~\cite{MINKA2026CRO,Fano61,Pinsker1964}.
Cette extension suit Wilde~\cite{Wilde2017} et Nielsen \& Chuang~\cite{NielsenChuang2010},
assurant la compatibilit\'e arri\`ere avec le cadre classique~\cite{Shannon48,CoverThomas06}.

\subsubsection{Mesures d'information quantique}
\label{subsubsec:mesures_info_quantique}

\begin{definition}[Entropie de von Neumann~\cite{vonNeumann1932,NielsenChuang2010}]\label{def:entropie_conditionnelle_quantique}
Pour un \'etat quantique $\rho$~:
\begin{equation}
S(\rho) = -\operatorname{Tr}(\rho \log \rho)
\end{equation}
Pour un \'etat classique-quantique $\rho_{XE} = \sum_x P_X(x) |x\rangle\langle x|
\otimes \rho_E^x$, l'entropie conditionnelle est~:
\begin{equation}
S(X|E)_\rho = S(\rho_{XE}) - S(\rho_E)
\end{equation}
o\`u $\rho_E = \operatorname{Tr}_X(\rho_{XE})$.
\end{definition}

\begin{definition}[Information mutuelle quantique]\label{def:information_mutuelle_quantique}
Pour un \'etat biparti $\rho_{AB}$~:
\begin{equation}
I(A; B)_\rho = S(\rho_A) + S(\rho_B) - S(\rho_{AB})
\end{equation}
\end{definition}

\begin{definition}[Information quantique accessible]\label{def:info_accessible}
Pour un ensemble $\{P_X(x), \rho_x\}$, l'information accessible est~:
\begin{equation}
I_{\mathrm{acc}} = \sup_{\{\Pi_y\}} I(X; Y)
\end{equation}
o\`u le supremum est pris sur toutes les POVMs $\{\Pi_y\}$ sur le syst\`eme
quantique~\cite{NielsenChuang2010}. Selon la borne de Holevo~\cite{Holevo1973},
$I_{\mathrm{acc}} \leq \chi$.
\end{definition}

\begin{definition}[Entropie min-max liss\'ee]\label{def:entropie_min_max_lissee}
Pour $\varepsilon \in [0,1)$ et un \'etat $\rho_{AB}$, l'entropie min-max
liss\'ee est~:
\begin{equation}
H_{\min}^\varepsilon(A|B)_\rho = \max_{\bar{\rho}_{AB} \in \mathcal{B}^\varepsilon(\rho_{AB})}
H_{\min}(A|B)_{\bar{\rho}}
\end{equation}
o\`u $H_{\min}(A|B)_\rho = -\log_2 \min\{\mathrm{Tr}\,\sigma_B : \rho_{AB} \leq \mathbb{I}_A \otimes \sigma_B,\ \sigma_B \geq 0\}$
est l'entropie min conditionnelle et $\mathcal{B}^\varepsilon(\rho)$ la boule de
rayon $\varepsilon$ autour de $\rho$ pour la distance de purification.
\end{definition}

\begin{remark}[Relation avec l'entropie de Shannon]\label{rem:relation_shannon}
Lorsque $\rho = \sum_x P_X(x) |x\rangle\langle x|$ est diagonal dans la base
computationnelle (encodage classique)~:
\begin{equation}
S(\rho) = H(X)
\end{equation}
Ceci assure la compatibilit\'e arri\`ere avec le cadre CRO conditionnel pour
les protocoles classiques~\cite{CoverThomas06,HanVerdu}.
\end{remark}

\subsubsection{In\'egalit\'es quantiques cl\'es}
\label{subsubsec:inegalites_quantiques}

\begin{lemma}[In\'egalit\'e de Pinsker quantique]\label{lem:pinsker_quantique}
Pour tous \'etats quantiques $\rho, \sigma$~:
\begin{equation}
\|\rho - \sigma\|_1^2 \leq 2 \ln 2 \cdot D(\rho\|\sigma)
\end{equation}
\end{lemma}

\begin{Preuve}[Lemme~\ref{lem:pinsker_quantique}~: Pinsker quantique]\label{prv:pinsker_quantique}

\begin{cmtEtape}{Objectif.}
\'Etablir l'analogue quantique de l'in\'egalit\'e de Pinsker (sous-sous-section~\ref{subsubsec:pinsker_classique}) classique,
reliant la distance de trace \`a l'entropie relative quantique.
\end{cmtEtape}

Soit $\Pi$ le projecteur sur le sous-espace propre positif de
$\rho - \sigma$ et $\mathcal{M}$ la mesure \`a deux issues
$\{\Pi, \mathbb{I} - \Pi\}$ (mesure de Helstrom). Les lois de l'issue sont
$p = (\mathrm{Tr}\,\Pi\rho,\, 1 - \mathrm{Tr}\,\Pi\rho)$ et
$q = (\mathrm{Tr}\,\Pi\sigma,\, 1 - \mathrm{Tr}\,\Pi\sigma)$, de sorte que
\begin{equation}
\|p - q\|_1 = 2\,\mathrm{Tr}\,\Pi(\rho-\sigma) = \|\rho - \sigma\|_1 .
\end{equation}
\cmtFormule{La mesure de Helstrom pr\'eserve exactement la distance de trace.}

Par l'in\'egalit\'e de traitement de l'entropie relative quantique sous le
canal de mesure~\cite{Wilde2017}, $D(p\|q) \leq D(\rho\|\sigma)$.
\cmtFormule{Mesurer ne peut que r\'eduire la distinguabilit\'e.}

L'in\'egalit\'e de Pinsker classique, avec une entropie relative en bits,
donne $\|p - q\|_1^2 \leq 2\ln 2 \cdot D(p\|q)$. En combinant ces in\'egalit\'es, on obtient
\begin{equation}
\|\rho - \sigma\|_1^2 = \|p - q\|_1^2 \leq 2 \ln 2 \cdot D(p\|q) \leq 2\ln 2 \cdot D(\rho\|\sigma).
\end{equation}
\cmtFormule{C'est exactement l'in\'egalit\'e de Pinsker quantique.}
\cmt{Une diagonalisation simultan\'ee de $\rho$ et $\sigma$ n'est possible que s'ils commutent~; le passage par la mesure de Helstrom \'evite cette restriction.}
\end{Preuve}

\begin{lemma}[Mesure douce (gentle measurement)]\label{lem:mesure_douce}
Soit $\rho$ un \'etat quantique et $\{M_m\}$ une mesure POVM. Si le r\'esultat
$m$ se produit avec probabilit\'e $p_m = \operatorname{Tr}(M_m \rho) \geq 1 - \delta$
pour $\delta \in [0, 1)$, alors l'\'etat post-mesure
\begin{equation}
\rho' = \frac{\sqrt{M_m}\rho\sqrt{M_m}}{\operatorname{Tr}(M_m\rho)}
\end{equation}
satisfait
\begin{equation}
\|\rho - \rho'\|_1 \leq 3\sqrt{\delta}
\end{equation}
La constante optimale, $2\sqrt{\delta}$, est \'etablie
dans~\cite{Winter,OgawaNagaoka}~; la d\'emonstration donn\'ee ici,
volontairement \'el\'ementaire, fournit $3\sqrt{\delta}$, ce qui suffit
\`a tous les usages qui en sont faits dans cette th\`ese.
\end{lemma}

\begin{Preuve}[Lemme~\ref{lem:mesure_douce}~: Mesure douce]\label{prv:mesure_douce}

\begin{cmtEtape}{Objectif.}
Montrer qu'une mesure quasi-certaine ne perturbe presque pas l'\'etat.
\end{cmtEtape}

Par d\'efinition de $\rho'$~:
\begin{equation}
\rho' = \frac{\sqrt{M_m}\rho\sqrt{M_m}}{p_m}
\end{equation}
\cmtFormule{Formule de l'\'etat post-mesure.}

On a~:
\begin{equation}
\|\rho - \rho'\|_1 \leq \|\rho - \sqrt{M_m}\rho\sqrt{M_m}\|_1 + \|\sqrt{M_m}\rho\sqrt{M_m} - \rho'\|_1
\end{equation}
\cmtFormule{In\'egalit\'e triangulaire.}

Le second terme vaut
\begin{equation}
\|\sqrt{M_m}\rho\sqrt{M_m} - \rho'\|_1 = \left|1 - \frac{1}{p_m}\right| \cdot \|\sqrt{M_m}\rho\sqrt{M_m}\|_1 = \frac{1 - p_m}{p_m} \cdot p_m = 1 - p_m
\end{equation}
\cmtFormule{Calcul direct.}

Le premier terme est born\'e par l'in\'egalit\'e de Fuchs-van de Graaf, qui donne
\begin{equation}
\|\rho - \sqrt{M_m}\rho\sqrt{M_m}\|_1 \leq 2\sqrt{1 - \operatorname{Tr}(M_m\rho)} = 2\sqrt{\delta}
\end{equation}
\cmtFormule{Borne par la fid\'elit\'e.}

En combinant ces deux termes et en utilisant $1 - p_m \leq \delta$, on obtient
\begin{equation}
\|\rho - \rho'\|_1 \leq 2\sqrt{\delta} + \delta \leq 3\sqrt{\delta}
\end{equation}
\cmtFormule{La borne finale est $O(\sqrt{\delta})$.}
\end{Preuve}

\begin{lemma}[In\'egalit\'e de Fano quantique]\label{lem:fano_quantique}
Soit $X$ une variable al\'eatoire classique encod\'ee dans l'\'etat quantique
$\rho_X$ et $\hat{X}$ le r\'esultat d'une mesure POVM avec probabilit\'e
d'erreur $p_e = \Pr[\hat{X} \neq X]$. Alors
\begin{equation}
S(X|\hat{X}) \leq h_2(p_e) + p_e \log(|\mathcal{X}| - 1)
\end{equation}
o\`u $h_2(p) = -p\log p - (1-p)\log(1-p)$ est l'entropie binaire.
\end{lemma}

\begin{Preuve}[Lemme~\ref{lem:fano_quantique}~: Fano quantique]\label{prv:fano_quantique}

\begin{cmtEtape}{Objectif.}
\'Etablir l'analogue quantique de l'in\'egalit\'e de Fano (sous-sous-section~\ref{subsubsec:fano_classique}) classique,
reliant l'entropie conditionnelle \`a la probabilit\'e d'erreur de mesure.
\end{cmtEtape}

Soit $E$ l'\'ev\'enement $\{\hat{X} \neq X\}$. On d\'ecompose l'entropie
conditionnelle en observant que l'issue $\hat{X}$ de la mesure est une variable al\'eatoire classique~;
$(X,\hat{X})$ a donc une loi jointe classique et $S(X|\hat{X}) = H(X|\hat{X})$.
L'in\'egalit\'e de Fano classique s'applique alors directement et donne
\begin{equation}
S(X|\hat{X}) \leq h_2(p_e) + p_e \log(|\mathcal{X}| - 1)
\end{equation}
\cmtFormule{La m\^eme borne que Fano classique, appliqu\'ee \`a l'issue de la mesure.}
\cmt{Le contenu quantique du lemme tient au choix de la mesure, puisque la borne vaut pour toute POVM.}
\end{Preuve}

\begin{lemma}[Borne de Holevo-Helstrom]\label{lem:holevo_helstrom}
Pour deux \'etats quantiques $\rho_0, \rho_1$ avec probabilit\'es a priori
\'egales, la probabilit\'e d'erreur minimale de distinction est
\begin{equation}
p_{\mathrm{err}}^{\min} = \frac{1}{2}\left(1 - \frac{1}{2}\|\rho_0 - \rho_1\|_1\right)
\end{equation}
\end{lemma}

\subsubsection{Mod\`eles de canaux quantiques}
\label{subsubsec:modeles_canaux}

\begin{definition}[Canal d\'epolarisant]\label{def:canal_depolarisant}
Le canal d\'epolarisant avec param\`etre $p \in [0, 1]$ agit sur un \'etat
quantique de dimension $d$ comme
\begin{equation}
\Lambda_{\text{dep}}(\rho) = (1-p)\rho + \frac{p}{d} I_d
\end{equation}
Ce mod\`ele pr\'esente un bruit sym\'etrique dans les syst\`emes quantiques et
est le mod\`ele d'erreur dominant dans les dispositifs NISQ actuels~\cite{Preskill2018,Bharti2022}.
\end{definition}

\begin{definition}[Canal d'amortissement]\label{def:canal_amortissement}
Le canal d'amortissement d'amplitude avec param\`etre $g \in [0, 1]$ est d\'efini par
\begin{equation}
\Lambda_{\text{AD}}(\rho) = E_0 \rho E_0^\dagger + E_1 \rho E_1^\dagger
\end{equation}
o\`u $E_0 = \begin{pmatrix} 1 & 0 \\ 0 & \sqrt{1-g} \end{pmatrix}$ et
$E_1 = \begin{pmatrix} 0 & \sqrt{g} \\ 0 & 0 \end{pmatrix}$ (not\'e $g$ pour \'eviter
toute confusion avec $\gamma = \alpha H(C) - 1$).
\end{definition}

\begin{definition}[Canal de d\'ephasage]\label{def:canal_dephasage}
Le canal de d\'ephasage avec param\`etre $p \in [0, 1]$ (not\'e $p$ pour \'eviter
toute confusion avec le param\`etre de s\'ecurit\'e $\lambda$) est d\'efini par
\begin{equation}
\Lambda_{\text{ph}}(\rho) = (1-p)\rho + p\, Z \rho Z
\end{equation}
o\`u $Z$ est la matrice de Pauli $Z$.
\end{definition}

\begin{definition}[Canal d'effacement]\label{def:canal_effacement}
Le canal d'effacement avec param\`etre $e \in [0, 1]$ est d\'efini par
\begin{equation}
\Lambda_{\text{er}}(\rho) = (1-e)\rho + e |\bot\rangle\langle\bot|
\end{equation}
o\`u $|\bot\rangle$ est un \'etat orthogonal \`a l'espace original.
\end{definition}

\begin{definition}[Norme diamant]\label{def:norme_diamant}
La norme diamant d'une application lin\'eaire $\Phi$ entre espaces
d'op\'erateurs est
\begin{equation}
\|\Phi\|_\diamond = \sup_{\rho_{RA}} \|(\mathcal{I}_R \otimes \Phi)(\rho_{RA})\|_1
\end{equation}
o\`u le supremum est pris sur tous les \'etats $\rho_{RA}$ d'un syst\`eme
de r\'ef\'erence $R$ arbitrairement grand.
\end{definition}

% -----------------------------------------------------------------------
\begin{table}[htbp]
\centering
\caption{Comparaison des mod\`eles de canaux quantiques pour l'analyse CRO}
\label{tab:modeles_canaux_quantiques}
\begin{tabular}{@{}lccc@{}}
\toprule
\textbf{Type de canal} & \textbf{M\'ecanisme} & \textbf{Comportement $\eta_q$} & \textbf{Usage PQC} \\
\midrule
D\'epolarisant~\cite{NielsenChuang2010}      & Erreurs uniformes    & $\eta_q \propto p$           & R\'ef\'erence g\'en\'erale    \\
Amortissement~\cite{Breuer2002}     & Dissipation         & $\eta_q \propto \sqrt{g}$ & Stockage quantique   \\
D\'ephasage~\cite{Zurek2003}         & D\'ecoh\'erence      & $\eta_q \propto p$     & Perte de coh\'erence    \\
Effacement~\cite{Grassl1997}        & Perte compl\`ete    & $\eta_q = e$                 & Erreurs transmission \\
\bottomrule
\end{tabular}
\begin{tablenotes}
\item \footnotesize Les valeurs $\eta_q$ sont des ordres de grandeur indicatifs.
\end{tablenotes}
\end{table}
% -----------------------------------------------------------------------
\end{sloppypar}

% -----------------------------------------------------------------------
\subsection{Mod\`ele formel quantique}
\label{subsec:modele_formel_pq}
% -----------------------------------------------------------------------
\begin{sloppypar}
\textbf{Nous d\'efinissons} les extensions quantiques des trois m\'etriques
CLO, assurant la compatibilit\'e d\'efinitionnelle avec le cadre classique.

\subsubsection{Syntaxe des protocoles de preuve post-quantiques}
\label{subsubsec:syntaxe_pq}

\begin{definition}[Protocole de preuve quantique]\label{def:protocole_pq}
Un protocole de preuve post-quantique $\mathcal{P}$ avec param\`etre de
s\'ecurit\'e $\lambda$ comprend~\cite{Unruh2016,Chatterjee2021}~:
\begin{itemize}
\item \textbf{Espace de messages} $\mathcal{M}_\lambda$ avec $|\mathcal{M}_\lambda| = 2^\lambda$
\item \textbf{Espace de contexte} $\mathcal{C}_\lambda$ avec distribution $P_C$ et entropie $H(C)$
\item \textbf{Espace de preuves} $\Sigma_\lambda$ (peut \^etre classique ou quantique)
\item \textbf{Algorithme de preuve} $\mathcal{P}(\lambda, m, c) \to \sigma$
pour $m \in \mathcal{M}_\lambda, c \in \mathcal{C}_\lambda$
\item \textbf{Espace de verdicts} $\mathcal{Y}_\lambda$ de taille finie $k_\lambda$
\item \textbf{Algorithme de v\'erification quantique} $\mathcal{V}^q(\lambda, \sigma, c)$
pouvant utiliser des requ\^etes en superposition quantique~\cite{Watrous2009}
\end{itemize}
\end{definition}

\subsubsection{D\'efinitions CLO quantiques}
\label{subsubsec:definitions_clo_quantiques}

\begin{definition}[Confidentialit\'e quantique]\label{def:conf_quantique}
Pour le protocole $\mathcal{P}$ avec param\`etre de s\'ecurit\'e $\lambda$, soit
$\rho_\Sigma^{m,c}$ l'\'etat quantique de preuve pour le message $m$ dans le
contexte $c$. La confidentialit\'e quantique est la version statistique de la
D\'efinition~\ref{def:conf_quantique_clo}, soit
\begin{equation}
\Conf^q(\mathcal{P}) := 1 - \max_{m_0,m_1,c}\tfrac{1}{2}\bigl\|\rho_\Sigma^{m_0,c} - \rho_\Sigma^{m_1,c}\bigr\|_1 ,
\label{eq:conf_q_pq}
\end{equation}
c'est-\`a-dire un moins l'avantage maximal de toute mesure POVM (th\'eor\`eme de
Helstrom). La s\'ecurit\'e contre des adversaires QPT est une propri\'et\'e
calculatoire distincte, plus faible, qui rel\`eve des jeux de s\'ecurit\'e.
\end{definition}

\begin{definition}[Fiabilit\'e quantique]\label{def:fiabilite_quantique}
Soit $\hat{M}$ l'\textbf{estimateur quantique optimal de message}, d\'efini
comme la mesure POVM $\{\Pi_m\}$ sur l'espace de Hilbert de preuve
$\mathcal{H}_\Sigma$ qui maximise la probabilit\'e d'identifier correctement
le message.
La \textbf{fiabilit\'e quantique} est la fiabilit\'e de reconstruction, d\'efinie par
\begin{equation}
\Rel^q(\mathcal{P}) = 1 - \Pr[\hat{M} \neq M]
\label{eq:rel_q_pq}
\end{equation}
o\`u la probabilit\'e est sur la distribution jointe $(M, C, \Sigma)$ et les
r\'esultats de mesure de l'estimateur optimal.
\end{definition}

\begin{definition}[Opposabilit\'e quantique]\label{def:opposabilite_quantique}
Soit $V = \mathsf{Interpret}(Y, C, K)$ le verdict du v\'erificateur institutionnel
$\mathcal{J}$, o\`u $Y$ est l'issue classique de sa mesure sur la preuve (\'eventuellement
apr\`es passage dans un canal $\Lambda$) et $K$ ses ressources propres,
avec $K \perp M \mid (\Sigma, C)$ (D\'efinition~\ref{def:interp_admissible}). L'opposabilit\'e quantique vaut
\begin{equation}
\Opp^q(\mathcal{P}) = \frac{I(V; \Sigma|C)}{H(V|C)}
\label{eq:opp_q_pq}
\end{equation}
o\`u l'information mutuelle et l'entropie sont calcul\'ees sur la distribution
jointe $(M, C, \Sigma, V)$. Comme dans le cas classique, un verdict fonction
d\'eterministe de $(\Sigma,C)$ donnerait $\Opp^q = 1$
(Remarque~\ref{rem:verdict_deterministe}).
\end{definition}

\begin{remark}[Compatibilit\'e avec la confidentialit\'e classique]\label{rem:compatibilite_classique}
Pour les protocoles classiques o\`u $\Sigma$ prend des valeurs dans un alphabet
fini $\mathcal{S}$ (i.e., $\rho_\Sigma^{m,c} = \sum_{s} P_{\Sigma|M,C}(s|m,c)\,|s\rangle\langle s|$
est diagonal), la distance de trace se r\'eduit \`a la distance de variation
totale~\cite{Levin,Pinsker1964}~:
\begin{equation}
\tfrac{1}{2}\|\rho_\Sigma^{m_0,c} - \rho_\Sigma^{m_1,c}\|_1
= \mathrm{TV}\!\bigl(P_{\Sigma|m_0,c},\, P_{\Sigma|m_1,c}\bigr)
\end{equation}
Donc $\Conf^q$ se r\'eduit \`a $\Conf$ du cadre classique~\cite{MINKA2026CLO},
confirmant que la d\'efinition est une g\'en\'eralisation stricte.
\end{remark}

\subsubsection{Perte d'interpr\'etabilit\'e quantique}
\label{subsubsec:perte_interpretabilite_quantique}

\begin{definition}[Perte d'interpr\'etabilit\'e quantique]\label{def:perte_interp_quantique_pq}
La perte d'interpr\'etabilit\'e quantique $\eta_q \in [0,1]$ associ\'ee \`a un
canal $\Lambda$ agissant sur le support de la preuve est celle de la
D\'efinition~\ref{def:perte_interp_quantique}, c'est-\`a-dire le plus petit r\'eel tel que
\begin{equation}
I(V;\Lambda(\Sigma)|C) \;\geq\; I(V;\Sigma|C) - \eta_q\, H(C).
\label{eq:eta_q_definition}
\end{equation}
C'est l'unique d\'efinition de $\eta_q$ utilis\'ee dans cette th\`ese.
\end{definition}

\begin{remark}[Interpr\'etation physique]\label{rem:interpretation_physique_pq}
$\eta_q$ quantifie la quantit\'e de contenu s\'emantique perdue lorsque la
preuve cryptographique passe par des canaux
quantiques~\cite{Zurek2003,Schumacher1997}. Pour les preuves qui sont des
cha\^ines de bits classiques, y compris les signatures ML-DSA, SLH-DSA et Falcon,
$\Lambda$ est l'identit\'e et $\eta_q = 0$. Pour un support ou un v\'erificateur
quantique, sc\'enario encore hypoth\'etique, $\eta_q > 0$ en raison du bruit
quantique~\cite{Holevo1973,Preskill2018}.
\end{remark}

\subsubsection{Mod\`ele d'interaction institutionnelle quantique}
\label{subsubsec:interaction_institutionnelle}

\begin{definition}[V\'erificateur institutionnel quantique]\label{def:verificateur_institutionnel_quantique}
Un v\'erificateur institutionnel quantique $\mathcal{J}$ est un algorithme
QPT op\'erant sur un espace de Hilbert $\mathcal{H}_\mathcal{J}$ et
impl\'ementant une fonction de verdict~:
\begin{equation}
V^q: \Sigma_\lambda \times \mathcal{C}_\lambda \to \mathcal{Y}_\lambda
\end{equation}
Le v\'erificateur est caract\'eris\'e par sa capacit\'e d'extraction
$\Delta_q$ et sa granularit\'e $\alpha_q$.
\end{definition}

\begin{definition}[Protocole d'interaction quantique]\label{def:protocole_interaction_quantique}
Un protocole d'interaction quantique entre un prouveur $\mathcal{P}$ et
un v\'erificateur $\mathcal{J}$ est une suite d'\'echanges quantiques~:
\begin{equation}
(\rho_1, \rho_2, \ldots, \rho_r) = \mathsf{Interact}(\mathcal{P}, \mathcal{J}, m, c)
\end{equation}
o\`u chaque $\rho_i$ est un \'etat quantique et $r$ est le nombre de rounds.
\end{definition}

\begin{remark}[Capacit\'e institutionnelle quantique]\label{thm:capacite_institutionnelle_quantique}
L'in\'egalit\'e
$I(V^q; \Sigma|C) \leq I(M; \Sigma|C) + \Delta_q$
est le \emph{contenu} de l'Hypoth\`ese~\ref{ass:extraction_quantique}~;
elle ne d\'ecoule pas de l'in\'egalit\'e de traitement des donn\'ees, qui
donne seulement $I(M;V^q|C) \leq I(M;\Sigma|C)$ (une borne sur ce que le verdict
apprend du message, distinct de ce qu'il apprend de la preuve). Pour une interpr\'etation
admissible (D\'efinition~\ref{def:interp_admissible}, avec
$K \perp M \mid (\Sigma, C)$), la Proposition~\ref{prop:fondement_extraction}
s'applique \`a l'issue classique de la mesure de l'institution et fournit le
fondement $\Delta_q \geq I(V^q;\Sigma|M,C)$. La formulation initiale de ce
r\'esultat, \'enonc\'ee comme th\'eor\`eme avec l'information accessible $I_{\mathrm{acc}}$,
est remplac\'ee par cette formulation plus pr\'ecise. Un verdict \'egal \`a un bit
d'al\'ea de $\Sigma$ ind\'ependant de $M$ montre en effet que l'in\'egalit\'e de
traitement ne s'applique pas sous cette forme.
\end{remark}
\end{sloppypar}

% -----------------------------------------------------------------------
\subsection{Hypoth\`eses de capacit\'e institutionnelle quantique}
\label{subsec:hypotheses_quantiques}
% -----------------------------------------------------------------------

\begin{sloppypar}
Nous posons quatre hypoth\`eses sur les capacit\'es
institutionnelles~\cite{Simon1955,Miller56,Ericsson95,Kahneman11} en
prenant en compte les r\'ealit\'es du quantique. Chacune repose sur des
justifications multidisciplinaires. Aucune n'a toutefois \'et\'e valid\'ee \`a ce
stade par protocoles exp\'erimentaux contr\^ol\'es mesurant directement
$\Delta_q$, $\eta_q$, $\beta_q$, $\zeta_q$ aupr\`es d'acteurs institutionnels
r\'eels, limitation reconnue explicitement
(sous-section~\ref{subsec:note_methodo}).

\begin{remark}[Version de la preuve]
\label{rem:version_preuve_pq}
La version de la preuve pr\'esent\'ee ici diff\`ere de celle de l'article
publi\'e~\cite{MINKA2026PQCRO}. Les formulations alternatives retenues, issues
de la certification Lean4, sont d\'etaill\'ees \`a la
section~\ref{subsec:lean4_erreurs}.
\end{remark}

\begin{assumption}[Extraction institutionnelle born\'ee quantique]
\label{ass:extraction_quantique}
Pour tout interpr\`ete institutionnel avec fonction de verdict $V$
op\'erant dans un contexte quantique ou hybride :
\begin{equation}
I(V; \Sigma|C) \leq I(M; \Sigma|C) + \Delta_q
\label{eq:extraction_quantique}
\end{equation}
o\`u $\Delta_q \geq 0$ capture \`a la fois la capacit\'e institutionnelle
classique et les modifications induites par le quantique.
\end{assumption}

\begin{assumption}[Verdict efficace quantique]
\label{ass:verdict_quantique}
L'espace de verdict institutionnel satisfait $k_\lambda \leq
2^{\alpha_q H(C)}$ pour une constante $\alpha_q > 0$, \emph{et} la
fonction de verdict utilise efficacement le contexte disponible :
\begin{equation}
H(V|C) \geq \theta_q \cdot \alpha_q H(C) =: \gamma_q, \qquad \gamma_q \geq 1
\label{eq:verdict_quantique}
\end{equation}
pour une constante de non-d\'eg\'en\'erescence $\theta_q \in (0, 1]$.
\end{assumption}

\begin{remark}[Statut de l'Hypoth\`ese~\ref{ass:verdict_quantique}]
\label{rem:statut_verdict_quantique}
L'Hypoth\`ese~4 initiale \'enon\c{c}ait $H(V|C) \geq \theta_q \log_2
k_\lambda$, mais l'Hypoth\`ese~2 donne $\log_2 k_\lambda \leq
\alpha_q H(C)$, une borne \textbf{sup\'erieure}. La borne inf\'erieure ne
peut donc pas \^etre d\'eduite de l'Hypoth\`ese~2. Elle est maintenant
postul\'ee explicitement comme propri\'et\'e d'efficacit\'e.
\end{remark}

\begin{assumption}[R\'egularit\'e de protocole quantique]
\label{ass:regularite_quantique}
Un protocole $\mathcal{P}$ est $(\beta_q, \zeta_q)$-\emph{quantique-r\'egulier}
s'il existe des constantes $\beta_q \geq 1$ et $\zeta_q \geq 0$ telles que
pour tout contexte $c$ et toute paire de messages $(m_0, m_1)$ :
\begin{equation}
D(\rho_\Sigma^{m_0,c} \| \rho_\Sigma^{m_1,c}) \leq \beta_q \cdot \chi(\mathcal{E}_{m,c}) + \zeta_q \cdot D(\bar{\rho}_\Sigma^c \| \mathbb{I}/d_\Sigma)
\label{eq:regularite_quantique}
\end{equation}
La constante $\zeta_q$ est une constante de r\'egularit\'e~; elle est
distincte du coefficient de perte d'interpr\'etabilit\'e $\eta_q$
(D\'efinition~\ref{def:perte_interp_quantique_pq}).
\end{assumption}

\begin{assumption}[R\'egime de pr\'eservation de la confidentialit\'e quantique]
\label{ass:regime_quantique}
Il existe une constante $\kappa \geq 0$, ind\'ependante de $\lambda$,
telle que
\begin{equation}
\chi_{\max} = \max_c \chi(\mathcal{E}_{m,c}) \leq \frac{\kappa}{\lambda}
\label{eq:regime_quantique}
\end{equation}
\end{assumption}

\begin{remark}[Statut de l'Hypoth\`ese~\ref{ass:regime_quantique}]
\label{rem:statut_regime_quantique}
Puisque $I(M;\Sigma|C)$ est la moyenne de $\chi(\mathcal{E}_{m,c})$ sur
les contextes et que $I_{\mathrm{acc}} \leq I(M;\Sigma|C)$ par la borne de
Holevo, l'Hypoth\`ese~\ref{ass:regime_quantique} implique
$I_{\mathrm{acc}} \leq I(M;\Sigma|C) \leq \kappa/\lambda$. Par le
Lemme~\ref{lem:q_conf_mi_pq}, c'est le r\'egime dans lequel les
protocoles quantiques-r\'eguliers garantissent la confidentialit\'e, avec
$\Conf^q \geq 1 - \frac{1}{2}\sqrt{2\beta_q \kappa/\lambda + 2\zeta_q
D_{\max}}$, minorant informatif seulement si $\zeta_q D_{\max}$ est petit
(Lemme~\ref{lem:q_conf_mi_pq}).
\end{remark}
\end{sloppypar}


% -----------------------------------------------------------------------
\subsection{\'Enonc\'e et preuve du Trilemme CRO conditionnel quantique}
\label{subsec:enonce_preuve_pq}
% -----------------------------------------------------------------------
\begin{sloppypar}
Nous \'etablissons le r\'esultat principal du r\'egime post-quantique.

\subsubsection{Lemme pr\'eparatoire : confidentialit\'e quantique et information de Holevo}
\label{subsubsec:lemme_conf_mi_pq}

\begin{remark}[Lemme~6 de la version publi\'ee]
\label{rem:retrait_lemma6}
La formalisation a mis en \'evidence un cas limite pour l'in\'egalit\'e (68)
de la version publi\'ee, avec l'esp\'erance prise sur des tirages
\emph{ind\'ependants} de $m_0, m_1$. Pour deux messages encod\'es dans des \'etats orthogonaux
(classiques, parfaitement distinguables), $I_{\mathrm{acc}} = 1$ bit,
alors que le membre de droite de (68) vaut $1/2$. La pr\'esente formulation
se passe donc du Lemme~6 et de la borne sup\'erieure (72) du Lemme~7 qui en
d\'epend. Le d\'etail de ces formulations alternatives est pr\'esent\'e \`a la
section~\ref{subsec:lean4_erreurs}.
\end{remark}

\begin{lemma}[Garantie de confidentialit\'e quantique]
\label{lem:q_conf_mi_pq}
Pour tout protocole $(\beta_q, \zeta_q)$-quantique-r\'egulier
(Hypoth\`ese~\ref{ass:regularite_quantique}) :
\begin{equation}
\Conf^q(\mathcal{P}) \geq 1 - \frac{1}{2}\sqrt{2\beta_q \chi_{\max} + 2\zeta_q D_{\max}}
\label{eq:conf_lower_pq}
\end{equation}
o\`u $\chi_{\max} = \max_c \chi(\mathcal{E}_{m,c})$ et
$D_{\max} = \max_c D(\bar{\rho}_\Sigma^c \| \mathbb{I}/d_\Sigma)
= \max_c\bigl(\log_2 d_\Sigma - S(\bar{\rho}_\Sigma^c)\bigr)$.
\end{lemma}

\begin{proof}
Pour tous $m_0, m_1, c$, l'in\'egalit\'e de Pinsker quantique
(Lemme~\ref{lem:pinsker_quantique}) et
l'Hypoth\`ese~\ref{ass:regularite_quantique} donnent
\begin{equation}
\|\rho_\Sigma^{m_0,c} - \rho_\Sigma^{m_1,c}\|_1^2 \leq 2\beta_q \chi(\mathcal{E}_{m,c}) + 2\zeta_q D(\bar{\rho}_\Sigma^c \| \mathbb{I}/d_\Sigma) \leq 2\beta_q \chi_{\max} + 2\zeta_q D_{\max}
\end{equation}
En prenant le maximum sur $(m_0, m_1, c)$ dans la relation entre
$\Conf^q$ et la distance de trace
(D\'efinition~\ref{def:conf_quantique}), on obtient~\eqref{eq:conf_lower_pq}.
\end{proof}

\begin{remark}[Interpr\'etation du Lemme~\ref{lem:q_conf_mi_pq}]
\label{rem:interp_lemme7}
Seule la borne inf\'erieure~\eqref{eq:conf_lower_pq} est utilis\'ee dans
la preuve du Th\'eor\`eme~\ref{thm:trilemme_pq}. Ce minorant n'est
informatif que si $\zeta_q D_{\max}$ est petit. Pour une preuve classique de
$n$ bits encod\'ee dans $d_\Sigma = 2^n$ dimensions,
$D_{\max} = n - H(\Sigma|C=c)$, qui vaut plusieurs milliers de bits pour les
signatures NIST PQC (plus de $2{,}6\cdot10^4$ bits pour ML-DSA-65, dont la
signature compte $3\,309$ octets). Pour ces sch\'emas, la confidentialit\'e
s'\'etablit par le Lemme~\ref{lem:confidentiality_guarantee} appliqu\'e \`a l'issue classique. La version publi\'ee
posait $D_{\max} = \negl(\lambda)$ pour ces sch\'emas. La pr\'esente formulation
pr\'ecise ce point.
\end{remark}

\subsubsection{Preuve principale en trois \'etapes}
\label{subsubsec:preuve_trois_etapes}

\begin{theorem}[Trilemme CRO conditionnel quantique]
\label{thm:trilemme_pq}
Sous les Hypoth\`eses~\ref{ass:extraction_quantique},
\ref{ass:verdict_quantique}, \ref{ass:regularite_quantique} et
\ref{ass:regime_quantique}, pour tout protocole de preuve post-quantique
$\mathcal{P}$ avec param\`etre de s\'ecurit\'e $\lambda \geq 64$,
entropie de message $H(M) = \lambda$, $M$ ind\'ependant de $C$ et
param\`etres institutionnels quantiques $(\Delta_q, \alpha_q, \beta_q,
\eta_q)$ :
\begin{equation}
\Conf^q(\mathcal{P}) \cdot \Rel^q(\mathcal{P}) \cdot \Opp^q(\mathcal{P}) \leq \frac{1}{\gamma_q}\left(\frac{\Delta_q}{\lambda} + \frac{\kappa(\Delta_q+1)}{\lambda^2} + \frac{\kappa^2}{\lambda^3}\right) \leq \frac{\Delta_q + \eta_q H(C)}{\lambda} + O(\lambda^{-2}) + \negl(\lambda)
\label{eq:trilemme_pq}
\end{equation}
\end{theorem}

\begin{proof}
\begin{cmtEtape}{Vue d'ensemble.}
La structure est celle du cas classique, transpos\'ee aux outils de la
th\'eorie de l'information quantique. L'in\'egalit\'e de Fano y devient sa
version quantique (sous-sous-section~\ref{subsubsec:fano_quantique}) tandis que la
distinguabilit\'e de deux \'etats se mesure par la distance de trace
(sous-sous-section~\ref{subsubsec:distance_trace_fidelite}), analogue quantique de la distance en variation totale. Posons $x = I(M;\Sigma|C)$~; par
l'Hypoth\`ese~\ref{ass:regime_quantique}, $x \leq \kappa/\lambda$.
\end{cmtEtape}

\paragraph{\'Etape 1. Fano quantique contraint la fiabilit\'e.}
Soit $I_{\mathrm{acc}}$ l'information accessible, c'est-\`a-dire la
quantit\'e maximale d'information que l'institution peut extraire de la
preuve quantique par une mesure, quelle qu'elle soit
(sous-sous-section~\ref{subsubsec:holevo}). L'in\'egalit\'e de Fano
quantique (Lemme~\ref{lem:fano_quantique}) appliqu\'ee \`a l'estimateur
optimal de $M$ donne, exactement comme dans le cas classique,
\begin{equation}
I_{\mathrm{acc}} \geq \Rel^q \cdot \lambda - 1
\label{eq:fano_q_applique}
\end{equation}
\cmtFormule{Une fiabilit\'e \'elev\'ee exige une information accessible \'elev\'ee.}
\cmt{La majoration $h(p_e) \leq 1$ vaut ici aussi, de sorte qu'aucune restriction sur $\Rel^q$ n'est n\'ecessaire.}

La borne de Holevo~\cite{Holevo1973} \'enonce qu'aucune mesure ne peut
extraire plus d'information que ce que l'\'etat contient, soit
$I_{\mathrm{acc}} \leq I(M;\Sigma|C) = x$. C'est la diff\'erence
principale avec le cas classique, o\`u les deux quantit\'es co\"incident.
En cha\^inant avec \eqref{eq:fano_q_applique}, on obtient
\begin{equation}
\boxed{\;\Rel^q(\mathcal{P}) \leq \frac{x+1}{\lambda}\;}
\label{eq:relq_bound}
\end{equation}
\cmt{Le passage par $I_{\mathrm{acc}}$ laisse la borne intacte et la rend seulement conditionnelle \`a la mesurabilit\'e effective de l'information.}

\paragraph{\'Etape 2. Extraction quantique born\'ee et verdict efficace contraignent l'opposabilit\'e.}
Par la D\'efinition~\ref{def:opposabilite_quantique}, l'opposabilit\'e
quantique garde la m\^eme forme, $\Opp^q = I(V;\Sigma|C)/H(V|C)$.
L'Hypoth\`ese~\ref{ass:extraction_quantique} majore son num\'erateur, car
l'institution n'extrait de la preuve que $x$, augment\'e d'au plus
$\Delta_q$ bits issus de ses propres ressources, d'o\`u
\begin{equation}
I(V;\Sigma|C) \leq x + \Delta_q
\end{equation}
\cmtFormule{Capacit\'e d'extraction institutionnelle en r\'egime quantique.}
\cmt{$\Delta_q$ peut \^etre inf\'erieur \`a son homologue classique $\Delta$ si lire une preuve post-quantique co\^ute plus cher cognitivement~; c'est une hypoth\`ese \`a mesurer.}

L'Hypoth\`ese~\ref{ass:verdict_quantique} minore son d\'enominateur par
$\gamma_q = \theta_q\,\alpha_q H(C) \geq 1$, o\`u $\theta_q \in (0,1]$
mesure la non-d\'eg\'en\'erescence du verdict. En divisant, on obtient
\begin{equation}
\boxed{\;\Opp^q(\mathcal{P}) = \frac{I(V; \Sigma|C)}{H(V|C)} \leq \frac{x + \Delta_q}{\gamma_q}\;}
\label{eq:oppq_bound}
\end{equation}

\paragraph{\'Etape 3. Formation du produit et application du r\'egime.}
Les facteurs \'etant tous positifs et $\Conf^q \leq 1$ par d\'efinition
d'une probabilit\'e de distinction, le produit des bornes
\eqref{eq:relq_bound} et \eqref{eq:oppq_bound} donne
\begin{equation}
\Conf^q \cdot \Rel^q \cdot \Opp^q \leq \frac{x+1}{\lambda} \cdot \frac{x+\Delta_q}{\gamma_q}
\label{eq:produit_q_brut}
\end{equation}
\cmt{Comme dans le cas classique, seule la quantit\'e $x$ relie les trois m\'etriques.}

Ce membre de droite est croissant en $x$~; on l'\'evalue donc en
$x = \kappa/\lambda$, borne sup\'erieure du r\'egime fix\'e par
l'Hypoth\`ese~\ref{ass:regime_quantique}, puis on d\'eveloppe le produit
\begin{align}
\Conf^q \cdot \Rel^q \cdot \Opp^q
&\leq \frac{1}{\gamma_q}\left(\frac{1}{\lambda} + \frac{\kappa}{\lambda^2}\right)\left(\Delta_q + \frac{\kappa}{\lambda}\right) \nonumber \\
&= \frac{1}{\gamma_q}\left(\frac{\Delta_q}{\lambda} + \frac{\kappa\Delta_q}{\lambda^2} + \frac{\kappa}{\lambda^2} + \frac{\kappa^2}{\lambda^3}\right) \nonumber \\
&= \frac{1}{\gamma_q}\left(\frac{\Delta_q}{\lambda} + \frac{\kappa(\Delta_q+1)}{\lambda^2} + \frac{\kappa^2}{\lambda^3}\right)
\end{align}
\cmtFormule{Les deux termes en $\lambda^{-2}$ se regroupent, comme dans le cas classique.}

La derni\`ere in\'egalit\'e de~\eqref{eq:trilemme_pq} s'obtient en
observant que $\gamma_q \geq 1$ et que $\eta_q H(C) \geq 0$. La borne
\'enonc\'ee sous la forme $(\Delta_q + \eta_q H(C))/\lambda$ est donc plus
faible, mais elle fait appara\^itre explicitement la p\'enalit\'e
s\'emantique quantique.
\cmt{Lecture. Le terme $\eta_q H(C)$ n'appara\^it pas dans la borne explicite, qui ne d\'epend que de $\Delta_q$. Il figure dans la majoration affaiblie, o\`u il rend visible le co\^ut s\'emantique du canal quantique. L'effet du bruit sur les compromis r\'ealisables passe, lui, par deux autres voies, puisqu'il ampute la capacit\'e effective $\Delta_q$ et d\'egrade la garantie de confidentialit\'e du Lemme~\ref{lem:q_conf_mi_pq}.}
\end{proof}

\begin{remark}[Port\'ee du Th\'eor\`eme~\ref{thm:trilemme_pq}]
\label{rem:scope_trilemme_pq}
L'Hypoth\`ese~\ref{ass:regime_quantique} est n\'ecessaire. Pour
l'encodage classique diagonal ($\beta_q = \beta$, $\zeta_q = 0$,
$\eta_q = 0$), le protocole qui r\'ev\`ele $M$ avec probabilit\'e $1/2$ et
produit sinon une preuve uniform\'ement al\'eatoire satisfait les
Hypoth\`eses~\ref{ass:extraction_quantique}--\ref{ass:regularite_quantique}
avec $\beta_q \leq 1{,}04$ pour tout $\lambda \geq 64$ (contexte $C$ form\'e
de deux bits uniformes, $\alpha_q = 1$, $\theta_q = 1/2$, donc $\gamma_q = 1$~;
$V$ \'egal au premier bit de $\Sigma$, donc $H(V|C) = 1 \geq \gamma_q$). Ce protocole v\'erifie
$\Conf^q = 1/2$, $\Rel^q = 1/2 + 2^{-\lambda-1}$ et $\Opp^q = 1$, pour tout
$\lambda$ ; ici
$I(M;\Sigma|C) = \lambda/2 - 1 + O(\lambda 2^{-\lambda})$ se situe hors
du r\'egime~\eqref{eq:regime_quantique}. La borne ne tient donc pas
uniform\'ement pour $I = \Theta(\lambda)$.
\end{remark}

\begin{remark}[Statut de v\'erification]
\label{rem:statut_verification_pq}
Le noyau alg\'ebrique de cette preuve est analogue \`a celui du
Th\'eor\`eme~\ref{thm:trilemme_cro_conditionnel} classique, certifi\'e
m\'ecaniquement en Lean4 \`a la section~\ref{sec:lean4}. Le d\'etail des
formulations alternatives issues de la certification est pr\'esent\'e \`a la
section~\ref{subsec:lean4_erreurs}.
\end{remark}

\begin{remark}[Forme raffin\'ee quantique]\label{rem:forme_raffinee_pq}
Exactement comme dans le cas classique (Remarque~\ref{rem:activite_delta}),
$\Opp^q \leq 1$ donne la forme raffin\'ee
\begin{equation}
\Conf^q \cdot \Rel^q \cdot \Opp^q \leq \frac{1+\kappa/\lambda}{\lambda}\,
\min\left\{1,\; \frac{\kappa/\lambda + \Delta_q}{\gamma_q}\right\},
\label{eq:trilemme_pq_raffine}
\end{equation}
et la capacit\'e quantique est contraignante si et seulement si
$\Delta_q < \gamma_q - \kappa/\lambda$, soit
$H(C) > H^\sharp_q := (\Delta_q + \kappa/\lambda)/(\theta_q\alpha_q)$. Le
coefficient $\eta_q$ n'intervient pas dans cette forme.
\end{remark}

\begin{corollary}[Forme simplifi\'ee quantique]\label{cor:forme_simplifiee_pq}
Pour $\lambda \geq 128$, $\beta_q = O(1)$, $\gamma_q = O(1)$~:
\begin{equation}
\Conf^q \cdot \Rel^q \cdot \Opp^q \lesssim \frac{\Delta_q + \eta_q \cdot H(C)}{\lambda}
\label{eq:trilemme_pq_simplifie}
\end{equation}
\end{corollary}

\begin{boiteArbitre}
L'arbitre comprend que ce terme $\eta_q H(C)/\lambda$ rel\^ache la borne
au lieu de lui imposer un plancher. Si un jour ses preuves
transitent par des canaux quantiques, le bruit lui prendra une partie de ce
qu'elles disent. Pour les signatures qu'il recevra demain, qui sont
des cha\^ines de bits, $\eta_q$ reste en revanche nul.
\end{boiteArbitre}
\end{sloppypar}

% -----------------------------------------------------------------------
\subsection{R\'egimes op\'erationnels quantiques}
\label{subsec:regimes_quantiques}
% -----------------------------------------------------------------------
\begin{sloppypar}
Comme dans le cas classique, le Trilemme CRO conditionnel quantique
distingue deux r\'egimes, s\'epar\'es par le seuil
$H^\sharp_q = (\Delta_q + \kappa/\lambda)/(\theta_q\alpha_q)$ de la
Remarque~\ref{rem:forme_raffinee_pq}. Pour $H(C) \leq H^\sharp_q$, le produit
est born\'e par l'in\'egalit\'e de Fano, alors que pour $H(C) > H^\sharp_q$ la capacit\'e
quantique $\Delta_q$ est la contrainte liante. La transition post-quantique
d\'eplace une institution d'un r\'egime \`a l'autre lorsqu'elle modifie
$\Delta_q$ (formation, outils) ou la granularit\'e de ses verdicts.

\subsubsection{Profils institutionnels de migration (illustratifs)}
\label{subsubsec:regime_quantique_limite}
\label{subsubsec:regime_quantique_pret}
\label{subsubsec:regime_quantique_ameliore}

Trois profils de maturit\'e, qui rejoignent les recommandations des
r\'egulateurs europ\'eens~\cite{ENISA2021,ANSSI2022}, permettent de situer
les institutions. Ils d\'ecrivent des niveaux de capacit\'e. Le r\'egime de la borne d\'epend aussi de $H(C)$ et de la
granularit\'e des verdicts.

\begin{itemize}
\item \textbf{D\'ebut de migration.} La capacit\'e $\Delta_q$ est r\'eduite par le manque d'expertise
  PQC~; les protocoles doivent privil\'egier la transparence et
  l'explicabilit\'e.
\item \textbf{Institutions form\'ees.} La capacit\'e $\Delta_q$ est proche de la capacit\'e
  classique~; les constructions hybrides (classique + PQC) sont adapt\'ees.
\item \textbf{Institutions expertes.} La capacit\'e $\Delta_q$ est \'elev\'ee et l'outillage
  sp\'ecialis\'e~; les protocoles PQC purs sont envisageables.
\end{itemize}

% -----------------------------------------------------------------------
\begin{table}[htbp]
\centering
\caption{R\'egimes du Trilemme CRO quantique et profils institutionnels}
\label{tab:regimes_quantiques_comp}
\begin{tabular}{@{}lcc@{}}
\toprule
\textbf{Crit\`ere} & \textbf{Contexte rare} & \textbf{Contexte riche} \\
\midrule
Condition & $H(C) \leq H^\sharp_q$ & $H(C) > H^\sharp_q$ \\
Borne & $(1+\kappa/\lambda)/\lambda$ (Fano) & $\frac{1+\kappa/\lambda}{\lambda}\cdot\frac{\kappa/\lambda+\Delta_q}{\gamma_q}$ \\
Contrainte liante & fiabilit\'e de reconstruction & capacit\'e quantique $\Delta_q$ \\
Profil typique & institution form\'ee ou experte & d\'ebut de migration \\
\bottomrule
\end{tabular}
\begin{tablenotes}
\item \footnotesize $H^\sharp_q = (\Delta_q + \kappa/\lambda)/(\theta_q\alpha_q)$ (Remarque~\ref{rem:forme_raffinee_pq}). Le coefficient $\eta_q$ n'intervient pas dans la distinction des r\'egimes. Les niveaux de s\'ecurit\'e NIST (1 \`a 5) mesurent la r\'esistance cryptographique et sont sans rapport avec ces r\'egimes. Valeur illustrative (sous-section~\ref{subsec:note_methodo}).
\end{tablenotes}
\end{table}
\end{sloppypar}
% -----------------------------------------------------------------------

% -----------------------------------------------------------------------
\subsection{Calibration illustrative des param\`etres quantiques}
\label{subsec:analyse_coherence_pq}
% -----------------------------------------------------------------------
\begin{sloppypar}
Comme dans le cas classique, les syst\`emes d\'eploy\'es ne sont pas des
instances du th\'eor\`eme (Remarque~\ref{rem:portee_cas_etude}). Leurs
produits CLO permettent seulement de calibrer des ordres de grandeur
de $\Delta_q$~\cite{Adida08,Groth16,Johnson01,Dilithium2022,SPHINCS2022}.

\subsubsection{Syst\`emes classiques}

Pour les trois syst\`emes classiques, la preuve est une cha\^ine de bits et
$\eta_q = 0$ (Remarque~\ref{rem:interpretation_physique_pq}), de sorte que la calibration
quantique se r\'eduit \`a la calibration classique
(Tableau~\ref{tab:validation_cro_classique}), avec $\Delta_q = \Delta$.

% -----------------------------------------------------------------------
\begin{table}[htbp]
\centering
\caption{Calibration illustrative~: syst\`emes classiques en cadre quantique}
\label{tab:validation_classiques_pq}
\begin{tabular}{@{}lcccc@{}}
\toprule
\textbf{Syst\`eme} & $\boldsymbol{\eta_q}$ & $\boldsymbol{H(C)}$ \textbf{(domaine)} & $\boldsymbol{\Delta_q}$ \textbf{implicite} & $\boldsymbol{\Delta_q/\lambda}$ \\
\midrule
Helios~\cite{Adida08,Adida09UCL} & $0$ & 40 bits & $\approx 56$ bits & $0{,}437$ \\
zk-SNARKs~\cite{Groth16,BCTV14} & $0$ & 70 bits & $\approx 73$ bits & $0{,}570$ \\
ECDSA~\cite{Johnson01,NIST.FIPS.186-5} & $0$ & 20 bits & $\approx 35$ bits & $0{,}273$ \\
\bottomrule
\end{tabular}
\begin{tablenotes}
\item \footnotesize $\lambda = 128$. Valeurs de $H(C)$ tir\'ees du Tableau~\ref{tab:parametres_reference}. $\Delta_q := \lambda \cdot (\Conf\cdot\Relver\cdot\Opp) - \eta_q H(C)$, d\'efini par inversion de la forme publi\'ee (calibration circulaire par construction).
\end{tablenotes}
\end{table}
% -----------------------------------------------------------------------

\subsubsection{Sch\'emas post-quantiques standardis\'es}

% -----------------------------------------------------------------------
\begin{table}[htbp]
\centering
\caption{Calibration illustrative~: signatures post-quantiques}
\label{tab:validation_pqc}
\begin{tabular}{@{}lcccc@{}}
\toprule
\textbf{Sch\'ema} & \textbf{Type} & $\boldsymbol{\lambda}$ & $\boldsymbol{\eta_q}$ & $\boldsymbol{\Delta_q}$ \textbf{implicite} \\
\midrule
ML-DSA-65 (Dilithium-3)~\cite{Dilithium2022} & R\'eseau & 192 & $0$ & $\approx 54$ bits \\
ML-DSA-87 (Dilithium-5)~\cite{Dilithium2022} & R\'eseau & 256 & $0$ & $\approx 72$ bits \\
SLH-DSA (SPHINCS+)~\cite{SPHINCS2022} & Hachage & 256 & $0$ & $\approx 97$ bits \\
Falcon~\cite{Falcon2022} & R\'eseau & 256 & $0$ & $\approx 70$ bits \\
\bottomrule
\end{tabular}
\begin{tablenotes}
\item \footnotesize $\eta_q = 0$ car ces signatures sont des cha\^ines de bits classiques. Les valeurs $\eta_q \in [0{,}012 ; 0{,}020]$ de la formulation initiale correspondaient au sc\'enario hypoth\'etique d'un support quantique~; elles modifiaient $\Delta_q$ de moins d'un bit. ML-KEM (Kyber), m\'ecanisme d'encapsulation de cl\'e et non syst\`eme de preuve, ne figure pas dans le tableau.
\end{tablenotes}
\end{table}
% -----------------------------------------------------------------------

Les capacit\'es implicites ($54$ \`a $97$ bits) sont de l'ordre de quelques
dizaines de bits~; celle de SLH-DSA d\'epasse l'intervalle $[20, 80]$ bits
sugg\'er\'e par la litt\'erature sur la rationalit\'e
born\'ee~\cite{Miller56,Simon82,Ericsson95,Cowan2001}, ce qui signale soit une
sur\'evaluation du produit CLO retenu pour ce sch\'ema, soit une limite de
la calibration. Cette calibration est \textbf{r\'etrospective et
illustrative}~; elle ne constitue pas une preuve empirique des hypoth\`eses
institutionnelles quantiques.
\end{sloppypar}

% -----------------------------------------------------------------------
\subsection{Implications pour les standards NIST PQC}
\label{subsec:implications_standards}
% -----------------------------------------------------------------------
\begin{sloppypar}
Notre analyse r\'ev\`ele que les standards NIST PQC actuels se concentrent
exclusivement sur la s\'ecurit\'e computationnelle (IND-CCA2,
EUF-CMA)~\cite{Goldreich01,KatzLindell14} tout en n\'egligeant
l'interpr\'etabilit\'e institutionnelle. Nous fournissons trois recommandations
concr\`etes pour la migration institutionnelle.

\subsubsection{Tampon d'entropie pour la conformit\'e inter-juridictionnelle}

\textbf{Recommandation~1.} Inclure $\Delta H = \max_{J_1, J_2} |H_{J_1}(C) - H_{J_2}(C)|$
bits de redondance s\'emantique dans les formats de signature pour accommoder
les variations d'entropie contextuelle entre syst\`emes juridiques.

\subsubsection{M\'etadonn\'ees s\'emantiques pour applications \`a enjeux \'elev\'es}

\textbf{Recommandation~2.} Ajouter des champs de m\'etadonn\'ees
structur\'ees aux signatures PQC qui encodent les pr\'edicats l\'egaux, les
bornes temporelles et les exigences proc\'edurales pour l'interpr\'etation
institutionnelle.

\subsubsection{Constructions hybrides pour les p\'eriodes de transition}

\textbf{Recommandation~3.} Superposer des structures classiques riches en
s\'emantique sur des primitives r\'esistantes au quantique pendant la
p\'eriode de transition (2024--2035) o\`u les institutions manquent de
capacit\'e pour l'interpr\'etation PQC pure~:
\begin{equation}
\sigma_{\text{hybride}} = (\sigma_{\text{PQC}}, \sigma_{\text{explication}},
\pi_{\text{liaison}})
\end{equation}
o\`u $\sigma_{\text{explication}}$ fournit un contenu s\'emantique lisible par
l'homme et $\pi_{\text{liaison}}$ lie cryptographiquement les deux composants.
\end{sloppypar}

% -----------------------------------------------------------------------
\subsection{Positionnement par rapport au r\'egime classique}
\label{subsec:transition_classique_quantique}
% -----------------------------------------------------------------------
\begin{sloppypar}
Le Trilemme CRO conditionnel quantique g\'en\'eralise le Trilemme classique
de la section~\ref{sec:cro_classique}.
Le tableau~\ref{tab:comparaison_classique_quantique} r\'ecapitule les
diff\'erences structurelles.

\begin{table}[htbp]
\centering
\caption{Comparaison des Trilemmes CRO classique et quantique}
\label{tab:comparaison_classique_quantique}
\begin{tabular}{@{}lcc@{}}
\toprule
\textbf{Aspect} & \textbf{Classique} & \textbf{Quantique} \\
\midrule
Borne explicite & $\frac{1}{\gamma}(\frac{\Delta}{\lambda} + O(\lambda^{-2}))$ & $\frac{1}{\gamma_q}(\frac{\Delta_q}{\lambda} + O(\lambda^{-2}))$ \\
Majoration affaiblie & $\Delta/\lambda$ & $(\Delta_q + \eta_q H(C))/\lambda$ \\
Contrainte liante (contexte riche) & Capacit\'e $\Delta$ & Capacit\'e $\Delta_q$ \\
Terme $\eta_q$ & Absent & Majoration affaiblie seulement \\
Adversaire & Statistique (PPT pour les jeux) & Statistique (QPT pour les jeux) \\
D\'egradation & Calculatoire & Calculatoire + physique (si support quantique) \\
\bottomrule
\end{tabular}
\end{table}

La diff\'erence structurelle majeure est l'apparition du terme
$\eta_q \cdot H(C)$, sans analogue classique. Il ne figure pas dans la
borne explicite, qui conserve la forme du cas classique avec $\Delta_q$ \`a
la place de $\Delta$~; il appara\^it dans la majoration affaiblie et
mesure le co\^ut s\'emantique impos\'e par le canal quantique. Son effet
sur les compromis r\'ealisables est indirect mais r\'eel, puisque le bruit
diminue la capacit\'e d'extraction effective $\Delta_q$ et affaiblit la
garantie de confidentialit\'e, deux canaux par lesquels une institution
per\c{c}oit concr\`etement la d\'egradation.
\end{sloppypar}

\begin{boiteArbitre}
L'arbitre referme le dossier quantique au terme de l'\'etape \textbf{E}.
Ce qu'il en retient est \`a la fois rassurant et exigeant, puisque la structure
du trilemme ne change pas dans le monde quantique alors que sa capacit\'e
$\Delta_q$ d\'ecidera de sa place.

Il en retient quatre apports. Les outils (von Neumann, Holevo, Pinsker
quantique) g\'en\'eralisent proprement ce qu'il connaissait du cadre classique.
Les hypoth\`eses, honn\^etes sur ce qu'elles ne prouvent pas encore, restent
coh\'erentes avec la litt\'erature cognitive. Le th\'eor\`eme conserve la forme
du cas classique et s'accompagne de la majoration affaiblie
$(\Delta_q + \eta_q H(C))/\lambda$. Enfin, la distinction des deux r\'egimes,
compl\'et\'ee par des profils de migration, transforme ce r\'esultat en
guide.

Le terme $\eta_q H(C)/\lambda$, sans \'equivalent classique, retient
particuli\`erement son attention. Il laisse la borne explicite inchang\'ee et
chiffre ce que le bruit quantique co\^uterait \`a l'interpr\'etation si la
preuve transitait par un support quantique. Pour les signatures
post-quantiques qu'il recevra, qui sont des cha\^ines de bits, ce terme est nul.

La calibration reste illustrative, car les syst\`emes analys\'es ne sont pas
des instances du th\'eor\`eme. Les capacit\'es implicites, de quelques
dizaines de bits, sont plausibles sans \^etre mesur\'ees.

Les bornes \'etant \'etablies, il reste \`a savoir si des constructions les
approchent, question \`a laquelle s'attelle le chapitre suivant.
\end{boiteArbitre}
% ============================================================================
% CONCLUSION DU CHAPITRE 2 
% source  : partie2/conclusion_partie2.tex
% Version : 2.1 — alignée sur la formalisation Lean4 révisée
% ============================================================================

\section{CONCLUSION DU CHAPITRE~\ref{chap:cadre_formel}}
\label{sec:conclusion_cadre_formel}

\begin{sloppypar}
Le cadre CLO (section~\ref{sec:clo}) a fait de $\Conf$, $\Rel$ et $\Opp$
des grandeurs mesurables. Le
th\'eor\`eme du goulot d'\'etranglement contextuel \'etablit que
$I(V;\Sigma|C) \leq \alpha H(\mathcal{C})$, liant l'opposabilit\'e
accessible \`a la granularit\'e des verdicts et \`a la richesse du contexte
institutionnel. Le cadre distingue la fiabilit\'e de reconstruction
$\Relrec$, contr\^ol\'ee par Fano, de la fiabilit\'e de v\'erification
$\Relver$. Il exige aussi une interpr\'etation admissible pour que $\Opp$ discrimine
les protocoles. Les jeux cryptographiques (IND-CONF,
BIND-REL, EXT-OPP) compl\`etent cette formalisation en \'etablissant le pont
avec les r\'eductions de s\'ecurit\'e standard.

Le Trilemme CRO classique (section~\ref{sec:cro_classique}) a born\'e
ces m\'etriques sous quatre hypoth\`eses explicites sur les limites
institutionnelles, \`a savoir l'extraction born\'ee, le verdict efficace, la r\'egularit\'e
du protocole et le r\'egime de confidentialit\'e $I(M;\Sigma|C) \leq \kappa/\lambda$.
Le produit $\Conf\cdot\Rel\cdot\Opp$ ne peut
d\'epasser $\Delta/(\gamma\lambda) + O(\lambda^{-2})$, tandis que la forme raffin\'ee
identifie le r\'egime ($H(\mathcal{C}) > H^\sharp$) o\`u la capacit\'e
institutionnelle $\Delta$ est la contrainte liante~; l'Hypoth\`ese d'extraction
born\'ee y re\c{c}oit un fondement informationnel exact
(Proposition~\ref{prop:fondement_extraction}). Les syst\`emes d\'eploy\'es ne
sont pas des instances du th\'eor\`eme (Remarque~\ref{rem:portee_cas_etude}).
Ils servent seulement \`a calibrer des capacit\'es implicites,
$\Delta_{\text{Helios}} \approx 56$~bits,
$\Delta_{\text{zk-SNARKs}} \approx 73$~bits, $\Delta_{\text{ECDSA}}
\approx 35$~bits, plausibles au regard de la litt\'erature sur la
rationalit\'e born\'ee.

L'extension post-quantique (section~\ref{sec:cro_pq}) a montr\'e que
cette borne s'\'etend aux adversaires quantiques avec la m\^eme forme,
param\'etr\'ee par la capacit\'e $\Delta_q$. Le coefficient de perte
d'interpr\'etabilit\'e $\eta_q$ n'appara\^it que dans une majoration
affaiblie~; il est nul pour les preuves qui sont des cha\^ines de bits,
y compris les signatures post-quantiques standardis\'ees.

Ces trois acquis, des m\'etriques, une borne classique et son extension
quantique, d\'elimitent ce qui est interdit sans encore indiquer ce qui est
atteignable, question que traite le chapitre~\ref{chap:constructions}.
\end{sloppypar}


% ==========================================================
% CHAPITRE 3 : CONSTRUCTIONS OPTIMALES ET CERTIFICATION FORMELLE
% ==========================================================
\chapter[CONSTRUCTIONS ET CERTIFICATION FORMELLE]%
        {CONSTRUCTIONS ET CERTIFICATION FORMELLE}
\label{chap:constructions}
\debutChapitre%
  {Constructions et certification formelle}%
  {%
    Ce chapitre r\'ealise les actes saturer et certifier du cycle
    canonique. Il construit d'abord le Semantic One-Time Pad, prototype qui
    montre que le paradigme Gold-Iron-Clay est instanciable, puis la famille
    SOTP. Les trois variantes de cette famille, une par r\'egime op\'erationnel,
    ont des propri\'et\'es de confidentialit\'e et de fiabilit\'e de v\'erification
    prouv\'ees, alors que leur opposabilit\'e est guid\'ee par une heuristique de conception.
    Leur optimalit\'e universelle est formul\'ee comme conjecture. Il soumet enfin le noyau alg\'ebrique du
    trilemme \`a l'assistant de preuve Lean4, d\'emarche qui conduit \`a
    expliciter les hypoth\`eses de la preuve, \`a en donner les constantes, \`a en
    caract\'eriser le r\'egime de validit\'e et \`a pr\'eciser trois \'enonc\'es publi\'es.%
  }%
  {%
    famille SOTP, heuristique de conception, conjecture d'optimalit\'e,
    formalisation Lean4, certification m\'ecanique%
  }%

% ============================================================================
% CHAPITRE 3 : INTRODUCTION
% ============================================================================

% ----------------------------------------------------------------------------
\section{INTRODUCTION}
\label{sec:introduction_constructions}
% ----------------------------------------------------------------------------

\begin{sloppypar}
Une borne d'impossibilit\'e ne dit rien, \`a elle seule, de ce qu'il est
possible de construire. Le chapitre pr\'ec\'edent a \'etabli que le produit
$\Conf\cdot\Rel\cdot\Opp$ ne peut d\'epasser une quantit\'e gouvern\'ee par
la capacit\'e d'extraction institutionnelle. Deux questions restent
ouvertes. La premi\`ere porte sur le caract\`ere atteignable de cette borne~; la
seconde demande si la preuve qui l'\'etablit r\'esiste \`a
une v\'erification m\'ecanique.

Ce chapitre r\'ealise les actes \textbf{S} et \textbf{C} du cycle
canonique. La section~\ref{sec:sotp_base} construit le Semantic One-Time
Pad, prototype qui montre que le paradigme Gold-Iron-Clay est
instanciable. La section~\ref{sec:sotp_construction} g\'en\'eralise ce
prototype en une famille de trois protocoles, chacun taill\'e pour un
r\'egime d'entropie contextuelle, puis formule la conjecture de leur
optimalit\'e universelle. La section~\ref{sec:lean4} soumet le noyau alg\'ebrique du
trilemme \`a l'assistant de preuve Lean4. La
section~\ref{sec:conclusion_constructions} dresse le bilan de ces deux
actes. L'acte \textbf{X}, exp\'erimenter, fait l'objet du
chapitre~\ref{chap:experimentation}.

% ----------------------------------------------------------------------------
\subsection{D\'emarche m\'ethodologique~: les actes \textbf{S} et \textbf{C}}
\label{subsec:demarche_saturer_certifier}
% ----------------------------------------------------------------------------

\paragraph{Saturer (\textbf{S}).}
Une borne d'impossibilit\'e prend toute sa valeur op\'erationnelle
lorsqu'on sait l'approcher. La m\'ethodologie vis\'ee est celle de la
construction explicite, qui associe \`a chaque r\'egime d'entropie
contextuelle un protocole de la famille SOTP (SOTP-Low, SOTP-Base, SOTP-High)
propos\'e comme construction de r\'ef\'erence. Leurs propri\'et\'es
cryptographiques sont prouv\'ees~; la saturation de la borne, qui porte sur
la fiabilit\'e de reconstruction, reste une conjecture
(Remarque~\ref{rem:statut_sotp}).

\paragraph{Certifier (\textbf{C}).}
Un th\'eor\`eme qui croise trois disciplines peut receler des hypoth\`eses
implicites que l'intuition math\'ematique ne d\'etecte pas. La
m\'ethodologie Lean4 impose une explicitation compl\`ete, puisque chaque
hypoth\`ese non prouv\'ee dans le syst\`eme doit \^etre \'enonc\'ee
comme propri\'et\'e ou comme pr\'emisse. Cette exigence a conduit \`a
expliciter l'hypoth\`ese de r\'egime et \`a pr\'eciser trois \'enonc\'es
publi\'es.

\begin{remark}[Statut des r\'esultats de ce chapitre]\label{rem:statut_sotp}
Les constructions S-OTP et SOTP sont \'evalu\'ees sur l'\'echelle
op\'erationnelle CLO, qui combine confidentialit\'e statistique, fiabilit\'e de
v\'erification $\Relver$ (Remarque~\ref{rem:rel_ver}) et opposabilit\'e
sous des conventions d'interpr\'etation explicites. Ces m\'etriques ne sont
pas compar\'ees \`a la borne du Th\'eor\`eme~\ref{thm:trilemme_cro_conditionnel},
qui porte sur $\Relrec$ et \'enonce que, pour tout protocole \`a confidentialit\'e
statistique $1-\negl(\lambda)$, $\Relrec \leq (1+\kappa/\lambda)/\lambda$.
L'existence de protocoles qui approchent la borne raffin\'ee~\eqref{eq:trilemme_cro_raffine}
est formul\'ee comme conjecture (Conjecture~\ref{conj:optimalite_universelle}).
Les composantes de preuve qui ne d\'ependent que du contexte public $C$
n'apportent aucune information conditionnelle, puisque
$I(V; g(C) \mid C) = 0$ pour toute fonction $g$. La couche Clay est
donc analys\'ee comme un \emph{format de pr\'esentation} qui facilite
l'interpr\'etation (baisse du co\^ut cognitif) sans ajouter
d'information au sens de $\Opp$. Les expressions num\'eriques de
saturation (facteurs $0{,}75$, $0{,}5$, $0{,}45$) sont des heuristiques de
conception, sans statut de minoration prouv\'ee.
\end{remark}
\end{sloppypar}

% ============================================================================
% CHAPITRE 3 - LE SEMANTIC ONE-TIME PAD (bloc issu de clo.tex)
% ============================================================================

\section{LE SEMANTIC ONE-TIME PAD~: CONSTRUCTION DE BASE}
\label{sec:sotp_base}

% Bloc issu de clo.tex : paradigme Gold-Iron-Clay, S-OTP de base, metriques CLO atteintes
\subsection{Le paradigme Gold-Iron-Clay}
\label{subsec:gold_iron_clay}
% -----------------------------------------------------------------------
\begin{sloppypar}
Les trois m\'etriques \'etant d\'efinies, il faut une m\'ethodologie pour
concevoir des protocoles qui les \'equilibrent~\cite{Shoup01,Damgard87}. Le paradigme
\textbf{Gold-Iron-Clay} structure cette conception en trois couches
fonctionnelles. Sa d\'enomination, inspir\'ee de la m\'etallurgie m\'edi\'evale,
refl\`ete la relation entre mall\'eabilit\'e et r\'esistance.

\paragraph{Couche Gold (confidentialit\'e).}
Masquage cryptographique du message via one-time pads bas\'es sur PRG,
chiffrements s\'emantiquement s\^urs~\cite{GoldwasserMicali84,Bellare97}, ou preuves zero-knowledge~\cite{GMR85,BenOr88}. La couche
Gold r\'epond \`a la question \emph{\guillemotleft~Comment rendre $\sigma$ indiscernable entre deux
messages~?\guillemotright}

\paragraph{Couche Iron (fiabilit\'e).}
Commitments contraignants (\emph{binding}) permettant la r\'ecup\'eration
v\'erifiable, commitments de Pedersen~\cite{Pedersen91,ChaumPedersen92}, accumulateurs~\cite{CamenischLysyanskaya02}, signatures. La couche
Iron r\'epond \`a la question \emph{\guillemotleft~Comment garantir que $\sigma$
permet de retrouver $m$ avec probabilit\'e $1 - \negl(\lambda)$~?\guillemotright}

\paragraph{Couche Clay (opposabilit\'e).}
Encodage pr\'eservant le contexte institutionnel pour l'extraction
d'information par des interpr\`etes aux ressources born\'ees~\cite{Miller56,Simon82}, par encodage
d\'eterministe, compression Johnson-Lindenstrauss~\cite{JohnsonLindenstrauss84,DasguptaGupta03} ou hachage d\'eterministe. La couche
Clay r\'epond \`a la question \emph{\guillemotleft~Comment permettre \`a $\mathsf{Interpret}$ d'extraire
un maximum d'information sur $V$ depuis $(c, \sigma)$~?\guillemotright}

\begin{figure}[htbp]
\centering
\begin{tikzpicture}[
    scale=0.9,
    transform shape
]

% ---------------------------------------------------------------
% D\'EGRAD\'ES M\'ETALLIQUES
% ---------------------------------------------------------------

% Couche Or
\shade[left color=metalorfonce!60, right color=metalor!40,
       middle color=metalor!80]
    (-4, 4) rectangle (4, 5.5);
\draw[thick, metalorfonce] (-4, 4) rectangle (4, 5.5);

% Couche Fer
\shade[left color=metalferfonce!70, right color=metalfer!50,
       middle color=metalfer!90]
    (-4, 2) rectangle (4, 3.5);
\draw[thick, metalferfonce] (-4, 2) rectangle (4, 3.5);

% Couche Argile
\shade[left color=metalargilefonce!70, right color=metalargile!50,
       middle color=metalargile!90]
    (-4, 0) rectangle (4, 1.5);
\draw[thick, metalargilefonce] (-4, 0) rectangle (4, 1.5);

% ---------------------------------------------------------------
% \'ETIQUETTES DES COUCHES
% ---------------------------------------------------------------
\node[font=\bfseries\sffamily, text=metalorfonce] at (0, 4.75)
    {Couche Or \quad - \quad Confidentialit\'e};
\node[font=\bfseries\sffamily, text=white] at (0, 2.75)
    {Couche Fer \quad - \quad Fiabilit\'e};
\node[font=\bfseries\sffamily, text=white] at (0, 0.75)
    {Couche Argile \quad - \quad Opposabilit\'e};

% ---------------------------------------------------------------
% CONTENU DES COUCHES
% ---------------------------------------------------------------
\node[font=\small, align=center, text=metalorfonce] at (0, 4.25)
    {Masquage PRG, preuves ZK \quad $\Conf \to 1$};
\node[font=\small, align=center, text=white] at (0, 2.25)
    {Engagement liant \quad $\Relver \to 1$};
\node[font=\small, align=center, text=white] at (0, 0.25)
    {Encodage contextuel \quad $\Opp \to \max$};

% ---------------------------------------------------------------
% REFLETS M\'ETALLIQUES
% ---------------------------------------------------------------
\draw[white!70!metalor, very thick, opacity=0.4]
    (-3.6, 5.2) -- (3.6, 5.2);
\draw[white!50!metalfer, very thick, opacity=0.3]
    (-3.6, 3.2) -- (3.6, 3.2);
\draw[white!50!metalargile, very thick, opacity=0.3]
    (-3.6, 1.2) -- (3.6, 1.2);

% ---------------------------------------------------------------
% FL\`ECHES DE FLUX
% ---------------------------------------------------------------
\draw[-{Stealth[length=2.5mm]}, thick, metalorfonce]
    (-5.2, 4.75) -- (-4.2, 4.75)
    node[left, xshift=-1cm, font=\small] {$M$};
\draw[-{Stealth[length=2.5mm]}, thick, metalferfonce]
    (4.2, 2.75) -- (5.2, 2.75)
    node[right, font=\small] {$\Sigma$};
\draw[-{Stealth[length=2.5mm]}, thick, metalargilefonce]
    (4.2, 0.75) -- (5.2, 0.75)
    node[right, font=\small] {$V$};

% ---------------------------------------------------------------
% FL\`ECHES VERTICALES INTER-COUCHES
% ---------------------------------------------------------------
\draw[-{Stealth[length=2.5mm]}, thick, chaptercolor]
    (3.2, 3.9) -- (3.2, 3.6);
\draw[-{Stealth[length=2.5mm]}, thick, chaptercolor]
    (3.2, 1.9) -- (3.2, 1.6);

% ---------------------------------------------------------------
% ACCOLADES LAT\'ERALES (propri\'et\'es CRO \`a droite)
% ---------------------------------------------------------------
\draw[decorate, decoration={brace, amplitude=5pt, mirror},
      thick, enspblue]
    (4.1, 4.0) -- (4.1, 5.5)
    node[midway, right=6pt, font=\footnotesize, text=enspblue]
        {$\Conf$};
\draw[decorate, decoration={brace, amplitude=5pt, mirror},
      thick, enspgreen]
    (4.1, 2.0) -- (4.1, 3.5)
    node[midway, right=6pt, font=\footnotesize, text=enspgreen]
        {$\Relver$};
\draw[decorate, decoration={brace, amplitude=5pt, mirror},
      thick, navyblue]
    (4.1, 0.0) -- (4.1, 1.5)
    node[midway, right=6pt, font=\footnotesize, text=navyblue]
        {$\Opp$};

\end{tikzpicture}
\caption[Paradigme Gold-Iron-Clay (or, fer, argile)~: architecture \`a trois couches]{%
  Paradigme Gold-Iron-Clay (or, fer, argile)~: architecture \`a trois couches pour les protocoles
  cryptographiques institutionnels.%
}\label{fig:gold_iron_clay}

\smallskip
\begin{minipage}{\linewidth}
  \small\textit{Note :} La couche \textbf{Or} (dor\'ee, $\Conf$) assure le
  masquage cryptographique, la couche \textbf{Fer} (acier, $\Relver$) garantit
  le liage et la v\'erification fiable, tandis que la couche \textbf{Argile} (terre
  cuite, $\Opp$) pr\'esente le contexte sous une forme interpr\'etable~; elle
  n'accro\^it $\Opp$ que si elle transmet une information que l'institution
  ne d\'etient pas d\'ej\`a (Remarque~\ref{rem:statut_sotp}). Les accolades
  droites maintiennent la correspondance avec les propri\'et\'es du Trilemme CRO.
\end{minipage}

\end{figure}
\end{sloppypar}

% -----------------------------------------------------------------------
\subsection{Principes de conception \'equilibr\'ee}
\label{subsec:principes_conception}
% -----------------------------------------------------------------------
\begin{sloppypar}
Un protocole CLO optimal alloue son budget entropique selon les exigences
applicatives. Le budget total d\'efini par
\begin{equation}
H_{\text{total}} = H_{\text{Gold}} + H_{\text{Iron}} + H_{\text{Clay}} \approx 2\lambda + H(\mathcal{C})
\label{eq:entropy_budget_clo}
\end{equation}
vaut par exemple $296$ bits pour $\lambda = 128$ et un contexte de vote
Helios \`a $H(\mathcal{C}) = 40$ bits.

Lorsque l'application priorise la confidentialit\'e ($\Conf_{\min} = 0{,}9$),
l'allocation recommand\'ee affecte $60\%$ du budget \`a la couche Gold
($\approx 180$ bits), $30\%$ \`a Iron ($\approx 90$ bits) et les
$10\%$ restants \`a Clay (compression agressive du contexte). Lorsqu'au
contraire l'opposabilit\'e prime ($\Opp_{\min} = 0{,}8$), les proportions
s'inversent, avec $20\%$ pour Gold (masquage minimal), $30\%$ pour Iron
(binding standard) et $50\%$ pour Clay (pr\'eservation riche du contexte).

L'ind\'ependance modulaire des couches est n\'ecessaire, car Gold, Iron et Clay
doivent pouvoir \^etre substitu\'es sans refonte compl\`ete~\cite{canetti2001,Goldreich02}. Remplacer le PRG
par AES-GCM n'affecte pas Iron ni Clay~; remplacer Pedersen par un
accumulateur RSA laisse Gold et Clay intacts~; remplacer le hachage
contextuel par une compression JL n'impacte pas les deux premi\`eres
couches. Cette modularit\'e facilite l'adaptation aux \'evolutions
cryptographiques (post-quantique, homomorphe) et r\'eglementaires~\cite{NISTPQC,Kyber2022,Dilithium2022}.
\end{sloppypar}

% -----------------------------------------------------------------------
\subsection{Construction~: le Semantic One-Time Pad (S-OTP)}
\label{subsec:construction_sotp}
% -----------------------------------------------------------------------
\begin{sloppypar}
Nous montrons que le paradigme Gold-Iron-Clay est instanciable
dans la classe des protocoles bas\'es XOR, via la construction
\textbf{Semantic One-Time Pad (S-OTP)}, prototype de la famille SOTP
\'etendue \`a la section~\ref{sec:sotp_construction}.

\begin{theorem}[Propri\'et\'es de s\'ecurit\'e S-OTP]\label{thm:proprietes_sotp}
Sous hypoth\`eses PRG (g\'en\'erateur pseudo-al\'eatoire) et difficult\'e du
logarithme discret (DLog) dans $\mathbb{G}$~\cite{BlumMicali84,Yao82}, on a
\begin{align}
\Conf(\text{S-OTP}) &= 1 - \negl(\lambda) \\
\Relver(\text{S-OTP}) &= 1 - \negl(\lambda) \\
\Opp(\text{S-OTP}) &\leq \frac{H(h \mid C)}{H(V|C)}
\end{align}
o\`u $h$ est la composante auxiliaire de la preuve qui porte l'information
contextuelle.
\end{theorem}

\begin{remark}[Port\'ee du Th\'eor\`eme~\ref{thm:proprietes_sotp}]\label{rem:portee_sotp}
La fiabilit\'e \'etablie est la fiabilit\'e de v\'erification $\Relver$~; la
fiabilit\'e de reconstruction v\'erifie $\Relrec \leq (1+\kappa/\lambda)/\lambda$
d\`es que $I(M;\Sigma|C) \leq \kappa/\lambda$ (Remarque~\ref{rem:rel_ver}). Si $h$ est
une fonction du contexte public $C$, alors $H(h|C) = 0$ et $\Opp(\text{S-OTP}) = 0$.
L'opposabilit\'e de S-OTP suppose donc que $h$ transmette une information
contextuelle que l'institution ne d\'etient pas d\'ej\`a (contexte d'esp\`ece
priv\'e). L'expression $H(\mathcal{C})/\lambda$ de~\cite{MINKA2026CLO} correspond \`a la
convention $H(h|C) = H(\mathcal{C})$ et $H(V|C) \approx \lambda$.
\end{remark}

\begin{Preuve}[Th\'eor\`eme~\ref{thm:proprietes_sotp}~: Confidentialit\'e]\label{prv:sotp_conf}

\begin{cmtEtape}{Objectif.}
Montrer que $\Conf(\text{S-OTP}) = 1 - \negl(\lambda)$. Le masquage
one-time pad rend le message chiffr\'e indiscernable, propri\'et\'e dont
h\'erite l'engagement Pedersen.
\end{cmtEtape}

Par la construction du masquage one-time pad, $M^* = m \oplus \mathsf{mask}_m$
o\`u $\mathsf{mask}_m \xleftarrow{\$} \{0,1\}^{|m|}$. Pour tout couple
$(m_0, m_1) \in \mathcal{M}$, on a
\begin{equation}
\TV(P_{M^*|M=m_0}, P_{M^*|M=m_1}) = 0
\end{equation}
\cmtFormule{Le masque uniforme rend les deux distributions identiques.}
\cmt{Un XOR avec un al\'ea uniforme efface toute trace du message.}

car les deux distributions sont uniformes sur $\{0,1\}^{|m|}$.
\cmt{Quel que soit $m$, $m \oplus \mathsf{mask}$ est uniforme.}

L'engagement $c_\Sigma = \mathsf{Com}(M^*)$ h\'erite de cette
indiscernabilit\'e, car l'engagement de Pedersen est \emph{parfaitement} masquant
sans hypoth\`ese calculatoire~\cite{Pedersen91}. Il en r\'esulte que
\begin{equation}
\Conf(\text{S-OTP}) = 1 - \max_{m_0,m_1} \TV(P_{\Sigma|M=m_0}, P_{\Sigma|M=m_1})
= 1 - \negl(\lambda)
\end{equation}
\cmtFormule{La confidentialit\'e atteint son maximum \`a une quantit\'e n\'egligeable pr\`es.}
\cmt{Pedersen cache parfaitement $M^*$, inconditionnellement~; seul le caract\`ere liant repose sur DLog.}
\end{Preuve}

\begin{Preuve}[Th\'eor\`eme~\ref{thm:proprietes_sotp}~: Fiabilit\'e]\label{prv:sotp_rel}

\begin{cmtEtape}{Objectif.}
Montrer que $\Relver(\text{S-OTP}) = 1 - \negl(\lambda)$. Il faut
\'etablir deux propri\'et\'es, la compl\'etude (une preuve honn\^ete est accept\'ee) et
la soundness (une preuve forg\'ee est rejet\'ee).
\end{cmtEtape}

La compl\'etude suit directement de la compl\'etude du sch\'ema
d'engagement. Pour tout $(m, c)$, on a
\begin{equation}
\Pr[\mathsf{Verify}(pk, \mathsf{Prove}(sk,m,c), c) = 1]
= 1 - \delta_{\mathsf{Com}}
\end{equation}
\cmtFormule{Une preuve honn\^ete est accept\'ee sauf erreur du sch\'ema d'engagement.}
o\`u $\delta_{\mathsf{Com}} = \negl(\lambda)$ est l'erreur d'engagement.
\cmt{L'erreur d'engagement est n\'egligeable en $\lambda$.}

Pour \'etablir la \emph{soundness}, supposons qu'un adversaire $\mathcal{A}$ produise
$\sigma^* = (c_\Sigma^*, h^*, \ldots)$ avec $\mathsf{Verify}(\sigma^*, c) = 1$
sans passer par $\mathsf{Prove}$. Alors $\mathcal{A}$ aurait ouvert
l'engagement $c_\Sigma^*$ sur un contenu valide~; cela contredit la
propri\'et\'e de \emph{binding} calculatoire du sch\'ema de Pedersen sous
l'hypoth\`ese DLog (probabilit\'e $\leq \negl(\lambda)$)~\cite{Pedersen91}. Donc
$\Relver(\text{S-OTP}) = 1 - \negl(\lambda)$.
\cmt{Pour forger, il faudrait ouvrir un engagement sur deux valeurs distinctes.}
\cmt{Cela \'equivaut \`a r\'esoudre le logarithme discret, ce qui est difficile.}
\end{Preuve}

\begin{Preuve}[Th\'eor\`eme~\ref{thm:proprietes_sotp}~: Opposabilit\'e]\label{prv:sotp_opp}

\begin{cmtEtape}{Objectif.}
Montrer que $\Opp(\text{S-OTP}) \leq H(h|C)/H(V|C)$.
L'information pertinente pour le verdict ne peut transiter que par la
composante auxiliaire $h$.
\end{cmtEtape}

L'information auxiliaire $h$ encode $\ell_{\text{aux}} = H(\mathcal{C})$
bits de contexte. Par construction, on a
\begin{equation}
I(V;\Sigma|C) \leq I(V;h|C) = H(h|C) - H(h|V,C) \leq H(h|C)
\end{equation}
\cmtFormule{Seule la composante $h$ porte l'information utile au verdict.}
\cmt{Le masque $M^*$ est ind\'ependant des verdicts.}
\cmt{L'engagement $c_\Sigma$ est parfaitement cach\'e.}

En divisant par $H(V|C)$, on obtient
\begin{equation}
\Opp(\text{S-OTP}) = \frac{I(V;\Sigma|C)}{H(V|C)}
\leq \frac{H(h|C)}{H(V|C)}
\end{equation}
\cmtFormule{L'opposabilit\'e est born\'ee par ce que $h$ apprend \`a l'institution au-del\`a du contexte public.}
\cmt{Si $h$ est une fonction de $C$, alors $H(h|C) = 0$ (voir la Remarque~\ref{rem:portee_sotp}).}
\end{Preuve}

\begin{remark}[Optimalit\'e de S-OTP dans la classe OTP~: argument heuristique]\label{thm:optimalite_sotp}
Parmi les protocoles qui masquent le message par un one-time pad et
n'encodent le contexte que dans une composante auxiliaire $h$, S-OTP
semble r\'ealiser le meilleur compromis, puisque tout protocole de ce type
consacre $|m| = \lambda$ bits d'al\'ea au masquage pour la s\'ecurit\'e
parfaite et voit son opposabilit\'e born\'ee par $H(h|C)/H(V|C)$ (m\^eme argument
que ci-dessus). Une preuve d'optimalit\'e demanderait de d\'efinir
formellement la classe OTP et d'y d\'emontrer une borne sup\'erieure
uniforme~; nous la laissons ouverte. Toute construction atteignant une
opposabilit\'e sup\'erieure doit soit (i) abandonner la confidentialit\'e
information-th\'eorique au profit d'hypoth\`eses calculatoires, soit (ii) utiliser
des protocoles interactifs multi-tours, soit (iii) employer des primitives
diff\'erentes, au-del\`a du paradigme OTP.
\end{remark}

\subsection{M\'etriques CLO atteintes}
\label{subsec:metriques_sotp}

Pour les param\`etres typiques $\lambda = 128$, $\Delta = 50$ bits,
$H(\mathcal{C}) = 30$ bits (r\'egime contexte rare), on obtient
\begin{align}
\Conf(\text{S-OTP}) &\approx 0{,}999 \label{eq:sotp_conf_numerique} \\
\Relver(\text{S-OTP}) &\approx 0{,}999 \label{eq:sotp_rel_numerique} \\
\Opp(\text{S-OTP}) &\approx \min\!\left\{1,\; \frac{50}{30}\right\} = 1 \label{eq:sotp_opp_numerique}
\end{align}

Ces valeurs sont des estimations th\'eoriques calcul\'ees sous les hypoth\`eses
PRG et DLog~\cite{Goldreich01,KatzLindell14}, sans mesure empirique
contr\^ol\'ee. Le produit vaut alors
\begin{equation}
\Conf \cdot \Relver \cdot \Opp \approx 0{,}998
\label{eq:sotp_produit}
\end{equation}
Cette valeur porte sur la fiabilit\'e de v\'erification et sur une
opposabilit\'e calcul\'ee sous la convention $\Opp = \min\{1,\Delta/H(\mathcal{C})\}$,
qui suppose que $h$ transmet un contexte d'esp\`ece inconnu de l'institution~;
elle n'est pas comparable \`a la borne du
Th\'eor\`eme~\ref{thm:trilemme_cro_conditionnel}. En fiabilit\'e de
reconstruction, la confidentialit\'e statistique de S-OTP impose
$\Relrec \leq (1+\kappa/\lambda)/\lambda$ (Remarque~\ref{rem:rel_ver}),
tandis que~\cite{MINKA2026CLO} donne un produit de l'ordre de $H(\mathcal{C})/\lambda$.
Ce produit \'elev\'e n'est r\'ealisable qu'en r\'egime contexte rare ($H(\mathcal{C}) \leq \Delta$).
D\`es que le contexte s'enrichit ($H(\mathcal{C}) \gg \Delta$), l'opposabilit\'e
s'effondre \`a $\Opp \approx \Delta / H(\mathcal{C}) \ll 1$, ce qui motive
les variantes SOTP-Base et SOTP-High de la section~\ref{sec:sotp_construction}.

\begin{boiteArbitre}
L'arbitre comprend que $0{,}998$ correspond au meilleur cas. D\`es que le
contexte s'enrichit, l'opposabilit\'e s'effondre et d'autres constructions
deviennent n\'ecessaires.
\end{boiteArbitre}

La remarque~\ref{thm:optimalite_sotp} indique que S-OTP serait optimal
dans sa classe, sans le d\'emontrer. Cette optimalit\'e, m\^eme \'etablie,
resterait restreinte, car les variantes SOTP-Base et SOTP-High utilisent
des techniques au-del\`a des protocoles XOR (compression JL, zk-SNARKs)
et surpassent S-OTP dans les r\'egimes appropri\'es.
\end{sloppypar}

% -----------------------------------------------------------------------
% ============================================================================
% SECTION 2.5 : LA FAMILLE SOTP — CONSTRUCTIONS OPTIMALES
% Source : partie2/sotp_construction.tex
% Version : 2.2 — alignée sur la formalisation Lean4 (section~\ref{sec:lean4})
% ============================================================================

\section{LA FAMILLE SOTP~: UNE CONSTRUCTION PAR R\'EGIME}
\label{sec:sotp_construction}
\begin{sloppypar}
Le Trilemme CRO conditionnel classique~(section~\ref{sec:cro_classique}) \'etablit
que $\Conf \cdot \Rel \cdot \Opp \leq \Delta/(\gamma\lambda) +
O(\lambda^{-2})$, avec une forme raffin\'ee. Une borne d'impossibilit\'e sans
construction est incompl\`ete, car elle \'enonce une limite sans dire si
cette limite est approchable. La \textbf{famille SOTP} (Semantic One-Time
Pad) propose des constructions de r\'ef\'erence, une par r\'egime de richesse
contextuelle, dont les propri\'et\'es cryptographiques sont prouv\'ees et dont
l'optimalit\'e \`a facteur constant est formul\'ee comme conjecture
(Conjecture~\ref{conj:optimalite_universelle}). Cette mise en regard
borne/construction transforme un r\'esultat n\'egatif en guide de conception.

\medskip
\noindent\textbf{Statut de cette contribution.} La famille SOTP
constitue la contribution \textbf{C4} de cette th\`ese. Elle est
aussi importante que la borne elle-m\^eme, puisqu'elle
donne au Trilemme CRO sa port\'ee op\'erationnelle. Les trois variantes
(SOTP-Low pour le contexte rare, SOTP-Base pour la transition, SOTP-High
pour le contexte riche) couvrent l'ensemble du spectre
d'entropie contextuelle et fournissent chacune une construction
cryptographique concr\`ete, dont la saturation de la borne reste \`a d\'emontrer. Cette sym\'etrie
se retrouve dans les th\'eor\`emes d'impossibilit\'e f\'econds. CAP trouve
son sens dans les syst\`emes distribu\'es r\'eels et FLP dans les protocoles
partiels (Paxos, Raft). Le Trilemme CRO trouve le sien dans les trois
variantes SOTP.

Cette section correspond \`a l'\'etape \textbf{S (Saturer)} de la progression
F$\to$B$\to$E$\to$\textbf{S}$\to$C$\to$X et apporte la contribution \textbf{C4}
de cette th\`ese.

La philosophie architecturale
commune est d'abord pos\'ee en sous-section~\ref{subsec:philosophie_sotp}. Les trois
constructions sont d\'etaill\'ees ensuite par ordre croissant de complexit\'e
contextuelle, avec SOTP-Low pour contexte rare (sous-section~\ref{subsec:sotp_low}),
SOTP-Base pour le r\'egime de transition (sous-section~\ref{subsec:sotp_base}),
SOTP-High pour contexte riche (sous-section~\ref{subsec:sotp_high}). Les
crit\`eres de s\'election entre variantes pour guider la conception
d'architecture sont fournis en sous-section~\ref{subsec:synthese_selection}.
La section se conclut par la conjecture d'optimalit\'e universelle et
l'analyse de sensibilit\'e (sous-section~\ref{subsec:optimalite_sensibilite}).
\end{sloppypar}

\begin{boiteArbitre}
La borne existe sans \^etre pour autant inaccessible. L'arbitre dispose
d\'esormais d'un protocole con\c{c}u \emph{pour lui}, \`a savoir SOTP-Low quand son
contexte est rare, SOTP-Base quand il est interm\'ediaire, SOTP-High quand il
est riche. Chaque variante calibre ses couches Gold,
Iron, Clay sur ce que son $\Delta$ lui permet d'extraire,
sans lui en demander davantage. L'impossibilit\'e devient ainsi un guide de
conception. Savoir \`a quel point ces constructions approchent la borne reste
une question ouverte. La borne sert d\'esormais de carte au concepteur.
\end{boiteArbitre}

% -----------------------------------------------------------------------
\subsection{Philosophie de conception~: paradigme Gold-Iron-Clay}
\label{subsec:philosophie_sotp}
% -----------------------------------------------------------------------
\begin{sloppypar}
Nous d\'efinissons l'architecture commune \`a toutes les variantes SOTP.

Trois couches fonctionnelles structurent chaque protocole selon le paradigme
Gold-Iron-Clay~(sous-section~\ref{subsec:gold_iron_clay}).
La couche \textbf{Gold} accomplit le masquage cryptographique au moyen de one-time pads
pseudo-al\'eatoires dans les r\'egimes bas et interm\'ediaire~\cite{Shannon49,BlumMicali84}, preuves zero-knowledge dans
le r\'egime riche~\cite{GMR85,BenOr88}. La couche \textbf{Iron} assure le binding via des commitments
contraignants, Pedersen~\cite{Pedersen91,ChaumPedersen92} pour SOTP-Low/Base, Groth16~\cite{Groth16,BCTV14}
pour SOTP-High, permettant l'ouverture v\'erifiable du message. La couche
\textbf{Clay} enfin pr\'esente le contexte institutionnel sous une forme que des
interpr\`etes \`a capacit\'e born\'ee~$\Delta$ peuvent exploiter pour rendre leurs
verdicts~\cite{Simon82,Miller56}~; elle n'accro\^it $\Opp$ que si elle transmet une
information que l'institution ne d\'etient pas d\'ej\`a (Remarque~\ref{rem:statut_sotp}).

L'architecture elle-m\^eme \'etait d\'ej\`a esquiss\'ee par la
construction S-OTP de la section~\ref{subsec:construction_sotp}. L'apport propre \`a cette section tient \`a son adaptation
par r\'egime, puisque chaque variante SOTP sp\'ecialise la couche Clay selon l'entropie
contextuelle $H(\mathcal{C})$, substituant encodage d\'eterministe, compression
Johnson-Lindenstrauss~\cite{JohnsonLindenstrauss84,DasguptaGupta03,Vempala05}, ou hachage SNARK
Poseidon~\cite{Groth16,Poseidon2019} selon la richesse du contexte.

\subsubsection{Correspondance entre m\'etriques CLO et propri\'et\'es cryptographiques}
\label{subsubsec:interpretation_operationnelle}

Tandis que le cadre CLO d\'efinit $\Conf$, $\Rel$, $\Opp$ via la th\'eorie de
l'information~\cite{MINKA2026CLO,Shannon48,CoverThomas06}, les preuves de ce chapitre les
expriment en termes de propri\'et\'es de s\'ecurit\'e cryptographique standard admettant des
r\'eductions aux hypoth\`eses bien \'etablies~\cite{Goldreich01,KatzLindell14}.

La \textbf{confidentialit\'e} $\Conf$, notion statistique, a pour pendant calculatoire l'indistinguabilit\'e s\'emantique~\cite{GoldwasserMicali84,MicaliRackoff91}, selon laquelle, pour
toute paire de messages $(m_0, m_1)$ dans le m\^eme contexte $c$, nul adversaire polynomial
ne peut distinguer la preuve de $m_0$ de celle de $m_1$ avec avantage non n\'egligeable.
La \textbf{fiabilit\'e de v\'erification} $\Relver$ (Remarque~\ref{rem:rel_ver}) correspond au binding~\cite{Pedersen91,Damgard87}, propri\'et\'e selon laquelle aucun adversaire polynomial ne peut
produire deux ouvertures distinctes d'un m\^eme commitment sur des messages diff\'erents avec
probabilit\'e non n\'egligeable. Quant \`a l'\textbf{opposabilit\'e} $\Opp$, elle correspond \`a
la fraction d'information extractible sur le verdict, soit $I(V;\Sigma|C)/H(V|C)$, quantifi\'ee
via la capacit\'e d'extraction institutionnelle~$\Delta$ (D\'efinition~\ref{def:delta_extraction})~\cite{Simon82,Miller56,Ericsson95}.

Cette traduction rend le cycle des preuves explicite, car les garanties de confidentialit\'e
et de fiabilit\'e de v\'erification r\'esultent de r\'eductions formelles \`a des hypoth\`eses
cryptographiques standard (PRG~\cite{BlumMicali84,Yao82}, DLog~\cite{DiffieHellman76},
q-SDH/q-PKE~\cite{BonehBoyen08})~; l'opposabilit\'e, elle, d\'epend de l'interpr\'etation
institutionnelle et n'est pas minor\'ee par ces r\'eductions.

\subsubsection{Architecture de la couche Clay et budget entropique}
\label{subsubsec:couche_clay}

La couche Clay encode le contexte $c$ sous une forme maximisant l'information extractible
par l'interpr\`ete institutionnel tout en maintenant les garanties Gold et Iron. Chacun des
trois r\'egimes appelle une technique propre.

Pour le r\'egime rare ($H(\mathcal{C}) < H^*(\lambda)$), l'encodage d\'eterministe
$\phi:\mathcal{C} \to \{0,1\}^{H(\mathcal{C})}$ encode l'int\'egralit\'e du contexte sans
perte~\cite{Shannon48}. Pour la zone de transition, la compression Johnson-Lindenstrauss~\cite{JohnsonLindenstrauss84,Achlioptas2003} projette le
contexte depuis la dimension $d = \Theta(\Delta \cdot H(\mathcal{C}))$ vers $k =
O(\log n/\epsilon^2)$ dimensions, pr\'eservant les distances \`a $\epsilon$ pr\`es. Pour le
r\'egime riche enfin, le hachage Poseidon~\cite{Poseidon2019,Grassi2021} compresse un contexte arbitrairement complexe en
un condens\'e de taille fixe ($256$ bits), agr\'ement\'e de m\'etadonn\'ees structur\'ees pour les
recherches institutionnelles.

Le budget entropique total de chaque preuve SOTP s'\'ecrit
\begin{equation}
H_{\text{total}} = H_{\text{Gold}} + H_{\text{Iron}} + H_{\text{Clay}}
  \approx 2\lambda + H(\mathcal{C})
\label{eq:entropy_budget}
\end{equation}
Pour $\lambda = 128$ et $H(\mathcal{C}) = 40$ bits (vote type Helios~\cite{Adida08}), cela donne
$\approx 296$ bits. Ce budget contraint la couche Clay, car pr\'eserver trop de contexte
peut corr\'eler $\Sigma$ avec $m$, r\'eduisant $\Conf$~; compresser trop agressivement
r\'eduit $\Opp$ parce que l'interpr\`ete ne peut plus extraire les verdicts~\cite{Smith09,Alvim12}.

L'allocation \'equilibr\'ee est modulaire par conception~\cite{canetti2001,Goldreich02}. La couche Gold peut substituer
AES-GCM~\cite{AES01} \`a ChaCha20~\cite{ChaCha20} sans affecter Iron ou Clay~; la couche Iron peut passer d'un
commitment Pedersen \`a un accumulateur RSA~\cite{CamenischLysyanskaya02} sans toucher Gold ou Clay~; la couche Clay
peut changer de technique d'encodage ind\'ependamment. Cette modularit\'e facilite les
migrations quantiques et l'adaptation aux \'evolutions r\'eglementaires.
\end{sloppypar}

% -----------------------------------------------------------------------
\subsection{SOTP-Low~: r\'egime \`a contexte rare}
\label{subsec:sotp_low}
% -----------------------------------------------------------------------
\begin{sloppypar}
Nous proposons, pour le r\'egime \`a contexte rare, une construction
d\'eterministe simple dont nous \'etablissons les propri\'et\'es cryptographiques
(Remarque~\ref{rem:statut_sotp}).

\subsubsection{Motivation et champ d'application}
\label{subsubsec:motivation_low}

Le r\'egime cible est celui des contextes que l'on peut encoder int\'egralement dans
la preuve, $H(\mathcal{C}) < H^*(\lambda)$, avec $H^*(\lambda) = \sqrt{\lambda}$ comme
crit\`ere de construction (environ $11$ bits \`a $\lambda = 128$), distinct du seuil
$H^\sharp$ du th\'eor\`eme (Remarque~\ref{rem:activite_delta}). Les signatures contractuelles ECDSA~\cite{Johnson01,NIST.FIPS.186-5} en sont le cas d'\'ecole, puisqu'un contrat
standardis\'e et ses parties sont souvent param\'etr\'es par un contexte d'entropie
$H(\mathcal{C}) \approx 20$ bits~\cite{Mason07,Kerr04}. D\`es lors que $\Delta \gtrsim H(\mathcal{C})$, la
capacit\'e institutionnelle absorbe l'int\'egralit\'e du contexte de sorte que l'opposabilit\'e peut
approcher~1.

L'id\'ee consiste \`a encoder le contexte \emph{d\'eterministiquement} dans la couche Clay,
ce qui n'entra\^ine ni perte de compression ni circuit co\^uteux. La preuve se compose d'un masque PRG
(Gold), d'un commitment Pedersen (Iron) et d'un triplet $(\phi(c), \mathsf{Hash}(c), V(c))$
(Clay), cette derni\`ere composante r\'ev\`ele le contexte sans r\'ev\'eler le message, puisque
$M \perp C$ par construction.

\subsubsection{Construction SOTP-Low}
\label{subsubsec:construction_low}

\begin{algorithm}[H]
\caption{SOTP-Low}
\label{alg:sotp_low}
\begin{algorithmic}[1]
\State \textbf{Param\`etres} : $\lambda$, $H(\mathcal{C}) < H^*(\lambda)$, $(\Delta, \alpha, \beta)$
\State \textbf{$\mathsf{Setup}(1^\lambda)$} :
\State \quad G\'en\'erer groupe cyclique $(\mathbb{G}, q)$ avec $|q| \geq 2\lambda$ et DLog difficile
\State \quad \'echantillonner $g, h \xleftarrow{\$} \mathbb{G}$ tels que $\log_g h$ inconnu
\State \quad Choisir PRG $\mathsf{PRG}: \mathbb{Z}_q \to \{0,1\}^\lambda$
\State \quad Choisir $\mathsf{Hash}: \mathcal{C} \to \{0,1\}^\lambda$ r\'esistante aux collisions
\State \quad D\'efinir encodage d\'eterministe $\phi: \mathcal{C} \to \{0,1\}^{H(\mathcal{C})}$
\State \quad \textbf{Retourner} $pk = (\mathbb{G}, q, g, h, \mathsf{PRG}, \mathsf{Hash}, \phi)$
\State \textbf{$\mathsf{Prove}(pk, m, c)$} :
\State \quad \textcolor{blue}{// Gold}
\State \quad $s \xleftarrow{\$} \mathbb{Z}_q$ ;\; $\sigma_1 \leftarrow m \oplus \mathsf{PRG}(s)$
\State \quad \textcolor{blue}{// Iron}
\State \quad $\rho \xleftarrow{\$} \mathbb{Z}_q$ ;\; $\sigma_2 \leftarrow g^{s} \cdot h^\rho$~\cite{Pedersen91}
\State \quad \textcolor{blue}{// Clay}
\State \quad $\sigma_3 \leftarrow (\phi(c),\; \mathsf{Hash}(c),\; V(c))$
\State \quad \textbf{Retourner} $\sigma = (\sigma_1, \sigma_2, \sigma_3)$ ;\; t\'emoin secret $(s, \rho)$
\State \textbf{$\mathsf{Verify}(pk, \sigma, c)$} :
\State \quad V\'erifier $\mathsf{Hash}(c) \stackrel{?}{=} \sigma_3.\mathsf{hc}$ et
         $\phi(c) \stackrel{?}{=} \sigma_3.h_c$ et $\sigma_2 \in \mathbb{G}$
\State \textbf{$\mathsf{Open}(pk, \sigma, c, (m, s, \rho))$} :
\State \quad Accepter ssi $\sigma_2 = g^{s} h^{\rho}$ et $\sigma_1 = m \oplus \mathsf{PRG}(s)$
\end{algorithmic}
\end{algorithm}

L'algorithme op\`ere en trois phases distinctes qui correspondent exactement
aux trois couches du paradigme Gold-Iron-Clay.

La phase \textsf{Setup} installe les param\`etres cryptographiques partag\'es.
Elle g\'en\`ere un groupe cyclique $(\mathbb{G}, q)$ d'ordre premier $\lambda$
bits dans lequel le logarithme discret est difficile, ce qui fonde
la s\'ecurit\'e de la couche Iron. Elle tire al\'eatoirement deux g\'en\'erateurs
$g$ et $h$ dont le rapport discret $\log_g h$ est inconnu de tous les
participants~; cette inconnaissance rend l'engagement de Pedersen
statistiquement cachant. Elle fixe enfin trois primitives fonctionnelles, \`a savoir
un PRG pour le masquage de message, une fonction de hachage r\'esistante
aux collisions pour l'authentification du contexte et un encodage
d\'eterministe $\phi$ qui serialise le contexte institutionnel $c$ sur
exactement $H(\mathcal{C})$ bits. Ce dernier point est sp\'ecifique au
r\'egime \`a contexte rare. En effet, parce que $H(\mathcal{C}) < H^*(\lambda)$, la
totalit\'e du contexte tient dans cet encodage sans compression avec perte.

La phase \textsf{Prove} construit la preuve en trois \'etapes parall\`eles
ind\'ependantes. La couche Gold (ligne~11) tire une graine
secr\`ete $s$ uniform\'ement dans $\mathbb{Z}_q$ et produit le message
masqu\'e $\sigma_1 = m \oplus \mathsf{PRG}(s)$. Le message r\'eel est ainsi
int\'egralement dissimul\'e derri\`ere un pad pseudoal\'eatoire, rendant
$\sigma_1$ indistinguable d'une cha\^ine uniforme pour tout adversaire ne
connaissant pas $s$. La couche Iron (ligne~13) construit
un engagement de Pedersen sur la graine $s$, soit $\sigma_2 = g^{s} \cdot h^\rho$
o\`u $\rho$ est un scalaire al\'eatoire. Tirer la graine dans $\mathbb{Z}_q$ est
n\'ecessaire. Avec $s \in \{0,1\}^\lambda$ et un engagement sur $\mathsf{int}(s)
\bmod q$, deux graines $s$ et $s+q$ donneraient en effet le m\^eme engagement, ce qui ferait perdre le
caract\`ere liant. Cet
engagement lie le producteur de la preuve \`a la valeur de $s$ sans la
r\'ev\'eler~; il garantit que le message masqu\'e $\sigma_1$ ne peut \^etre
ouvert que sur un unique message $m$. La couche Clay
(ligne~15) produit le triplet $\sigma_3 = (\phi(c), \mathsf{Hash}(c), V(c))$.
L'encodage d\'eterministe $\phi(c)$ y transmet la structure s\'emantique du
contexte, le condens\'e $\mathsf{Hash}(c)$ en garantit l'int\'egrit\'e tandis que
$V(c)$ est le verdict institutionnel attendu. Cette couche porte
l'int\'egralit\'e de l'information contextuelle exploitable par l'arbitre.
La preuve compl\`ete $\sigma = (\sigma_1, \sigma_2, \sigma_3)$ est
retourn\'ee publiquement~; la paire $(s, \rho)$ reste secr\`ete chez le
producteur.

La phase \textsf{Verify} est d\'elib\'er\'ement l\'eg\`ere, puisqu'elle contr\^ole
uniquement que le condens\'e de hachage du contexte soumis correspond \`a
celui encod\'e dans $\sigma_3$, que l'encodage $\phi(c)$ co\"incide et que
$\sigma_2$ appartient bien au groupe $\mathbb{G}$. Elle ne v\'erifie pas la
coh\'erence entre $\sigma_1$ et $\sigma_2$, car ce lien n'est contr\^ol\'e que lors
de l'ouverture ($\mathsf{Open}$), qui requiert le t\'emoin $(m, s, \rho)$~;
sans ouverture, $\sigma_1$ et $\sigma_2$ ne sont pas li\'es, de sorte que
$\mathsf{Verify}$ seul accepte toute paire bien form\'ee. L'arbitre peut
donc v\'erifier la forme de la preuve sans acc\'eder au message $m$~; la
garantie de liaison porte sur l'ouverture. La couche Clay pr\'esente le
contexte sous une forme lisible, mais, \'etant une fonction du contexte public
$c$, elle n'apporte aucune information conditionnelle au sens de $\Opp$
(Remarque~\ref{rem:statut_sotp}).

\subsubsection{Propri\'et\'es de SOTP-Low}
\label{subsubsec:saturation_low}

\begin{proposition}[Propri\'et\'es cryptographiques de SOTP-Low]\label{thm:saturation_sotp_low}
Sous s\'ecurit\'e PRG~\cite{BlumMicali84,Yao82} et difficult\'e DLog~\cite{DiffieHellman76}~:
\begin{enumerate}
\item tout adversaire PPT a un avantage IND-CONF n\'egligeable contre SOTP-Low
      (confidentialit\'e calculatoire)~;
\item la compl\'etude est parfaite et, pour la proc\'edure d'ouverture,
      la probabilit\'e de gain BIND-REL est n\'egligeable, donc
      $\Relver(\mathrm{SOTP\text{-}Low}) = 1 - \negl(\lambda)$ au sens de la
      Remarque~\ref{rem:jeu_bind_rel_link}.
\end{enumerate}
\end{proposition}

\begin{remark}[Heuristique de conception de SOTP-Low (non prouv\'ee)]
\label{rem:facteur_limitant_low}
La formulation initiale de ce r\'esultat \'enon\c{c}ait la minoration
$\Conf \cdot \Rel \cdot \Opp \geq 0{,}75 \cdot \min\{\Delta, H(\mathcal{C})\}/\lambda$.
La formalisation a conduit \`a la remplacer, pour trois raisons. (i)~La composante $\sigma_3 = (\phi(c), \mathsf{Hash}(c), V(c))$
est une fonction du contexte public, donc $I(V;\sigma_3|C) = 0$~; (ii)~le passage
par $H(V|C) \approx \lambda$ contredit l'Hypoth\`ese~\ref{ass:verdicts_efficace}
($\log_2|\mathcal{V}_\lambda| \leq \alpha H(\mathcal{C})$)~; (iii)~une telle minoration
d\'epasserait le majorant du Th\'eor\`eme~\ref{thm:trilemme_cro_conditionnel} d\`es
que $\gamma > 4/3$ tout en portant sur $\Relver$ au lieu de $\Relrec$. L'expression
$0{,}75\,\Delta/\lambda$ est conserv\'ee comme \emph{heuristique de conception}, car
elle d\'ecrit la fraction de la capacit\'e d'extraction qu'un encodage entier de
$\lfloor H(\mathcal{C}) \rfloor$ bits permettrait d'exploiter si ce contexte
\'etait un contexte d'esp\`ece inconnu de l'institution ($1 - 1/\Delta \geq 0{,}75$
pour $\Delta \geq 4$). Sa justification rigoureuse, dans un mod\`ele \`a contexte
priv\'e, est une perspective (chapitre~\ref{chap:resultats_discussions}).
\end{remark}

\begin{Preuve}[Proposition~\ref{thm:saturation_sotp_low}~: propri\'et\'es cryptographiques de SOTP-Low]\label{prv:saturation_sotp_low}

\begin{cmtEtape}{Vue d'ensemble.}
SOTP-Low est constitu\'e de trois couches, Gold (masquage par PRG),
Iron (engagement de Pedersen sur la graine $s$) et Clay (encodage
d\'eterministe $\phi(c)$). La confidentialit\'e et la fiabilit\'e de
v\'erification sont \'etablies s\'epar\'ement.
\end{cmtEtape}

\noindent\textbf{Rappel de la construction.}
\begin{align}
\sigma_1 &= m \oplus \mathsf{PRG}(s),
\quad s \xleftarrow{\$} \mathbb{Z}_q
\tag{couche Gold}\\
\sigma_2 &= g^{s} \cdot h^\rho,
\quad \rho \xleftarrow{\$} \mathbb{Z}_q
\tag{couche Iron}\\
\sigma_3 &= \bigl(\phi(c),\; \mathsf{Hash}(c),\; V(c)\bigr)
\tag{couche Clay}
\end{align}
o\`u $(g, h)$ sont des g\'en\'erateurs du groupe $(\mathbb{G}, q)$ avec
$\log_g h$ inconnu. Le t\'emoin secret est $(s, \rho)$.

\medskip
\noindent\textbf{Partie I~. Confidentialit\'e calculatoire.}
La loi de $\sigma$ sachant $M = m_i$ ne d\'epend de $m_i$ que par
$\sigma_1$. En effet, $\sigma_2$ est parfaitement masquant (pour tout $s$, $g^s h^\rho$
est uniforme dans $\mathbb{G}$ lorsque $\rho$ est uniforme) tandis que $\sigma_3$ ne
d\'epend que de $c$, ind\'ependant de $m$ ($M \perp C$).
\cmt{Seule $\sigma_1$ porte de l'information sur $m$.}

Supposons qu'un distingueur PPT $\mathcal{D}$ distingue $m_0 \oplus
\mathsf{PRG}(s)$ de $m_1 \oplus \mathsf{PRG}(s)$ avec un avantage non
n\'egligeable $\mu(\lambda)$. L'adversaire $\mathcal{B}$ qui re\c{c}oit
$y$ (sortie du PRG ou cha\^ine uniforme), tire $b$, soumet $m_b \oplus y$ \`a
$\mathcal{D}$ et r\'epond \og PRG\fg{} si $\mathcal{D}$ retrouve $b$, distingue le
PRG avec un avantage $\mu(\lambda)/2$, puisque si $y$ est uniforme, $m_b \oplus y$ ne
contient aucune information sur $b$. Cela contredit la s\'ecurit\'e du
PRG~\cite{BlumMicali84}. L'avantage IND-CONF de tout adversaire PPT est donc
n\'egligeable.
\cmt{La confidentialit\'e est calculatoire. Au sens statistique de la D\'efinition~\ref{def:confidentialite_clo}, elle n'est garantie que si le PRG est remplac\'e par un masque uniforme.}

\medskip
\noindent\textbf{Partie II~. Fiabilit\'e de v\'erification.}

\textit{Compl\'etude.} Pour toute preuve honn\^ete, $\mathsf{Hash}(c) = \sigma_3.\mathsf{hc}$
et $\phi(c) = \sigma_3.h_c$ par construction, $\sigma_2 \in \mathbb{G}$ par
fermeture du groupe~; enfin $\mathsf{Open}$ accepte le t\'emoin honn\^ete $(m, s, \rho)$.
La compl\'etude est parfaite.
\cmtFormule{La v\'erification et l'ouverture acceptent toute preuve honn\^ete.}

\textit{Liaison (r\'eduction \`a DLog).} Supposons qu'un adversaire PPT
produise une preuve $\sigma$ et deux ouvertures accept\'ees
$(m_0, s_0, \rho_0)$ et $(m_1, s_1, \rho_1)$ avec $m_0 \neq m_1$. Comme
$\sigma_1 = m_0 \oplus \mathsf{PRG}(s_0) = m_1 \oplus \mathsf{PRG}(s_1)$ et
$m_0 \neq m_1$, on a $s_0 \neq s_1$ dans $\mathbb{Z}_q$. L'\'egalit\'e
$g^{s_0} h^{\rho_0} = g^{s_1} h^{\rho_1}$ impose alors $\rho_0 \neq \rho_1$ et
\begin{equation}
\log_g h = \frac{s_0 - s_1}{\rho_1 - \rho_0} \bmod q,
\end{equation}
\cmtFormule{Deux ouvertures distinctes r\'ev\`elent $\log_g h$.}
ce qui contredit la difficult\'e du logarithme discret~\cite{DiffieHellman76}.
La probabilit\'e de gain BIND-REL est donc n\'egligeable, d'o\`u
$\Relver(\mathrm{SOTP\text{-}Low}) = 1 - \negl(\lambda)$.
\cmt{C'est ici que le tirage de $s$ dans $\mathbb{Z}_q$ est indispensable.}
\end{Preuve}

\subsubsection{Support de validation~: signatures ECDSA}
\label{subsubsec:validation_ecdsa}

\`A titre d'illustration, sans valeur de validation (les m\'etriques CLO d'ECDSA sont
des estimations param\'etriques analytiques issues des sp\'ecifications
publi\'ees~\cite{Johnson01,NIST.FIPS.186-5} et ECDSA n'est pas une instance du
th\'eor\`eme, Remarque~\ref{rem:portee_cas_etude}), les param\`etres $\Delta = 35$ bits,
$\lambda = 128$ donnent pour l'heuristique de SOTP-Low la valeur
$0{,}75 \times 35/128 \approx 0{,}205$. Cette valeur n'est pas comparable au produit
CLO d'ECDSA ($0{,}273$).
\end{sloppypar}

% -----------------------------------------------------------------------
\subsection{SOTP-Base~: r\'egime de transition}
\label{subsec:sotp_base}
% -----------------------------------------------------------------------
\begin{sloppypar}
Pour les contextes interm\'ediaires, nous proposons une construction fond\'ee sur la
compression Johnson-Lindenstrauss~\cite{JohnsonLindenstrauss84}, qui r\'eduit la taille de la
repr\'esentation du contexte tout en pr\'eservant ses distances \`a $\varepsilon$ pr\`es.

\subsubsection{Motivation et champ d'application}
\label{subsubsec:motivation_base}

Helios~\cite{Adida08,Adida09UCL} est l'ach\'etype m\^eme de ce r\'egime, puisque le contexte \'electoral (registre des \'electeurs, r\`egles du
scrutin, proc\`es-verbal) repr\'esente $H(\mathcal{C}) \approx 40$--$50$ bits pour une
\'election de taille moyenne, la capacit\'e d'un comit\'e de garants sp\'ecialis\'es \'etant estim\'ee
\`a $\Delta \approx 56$ bits. L'encodage d\'eterministe de SOTP-Low ne passe pas \`a
l'\'echelle, car la preuve serait trop volumineuse. Les zk-SNARKs de SOTP-High
seraient \`a l'inverse surdimensionn\'es pour ce niveau d'entropie.

La compression Johnson-Lindenstrauss~\cite{JohnsonLindenstrauss84,DasguptaGupta03,Achlioptas2003} permet de
projeter le contexte depuis la dimension $d = \Theta(\Delta \cdot H(\mathcal{C}))$ vers
$k = O(\log n/\epsilon^2)$ dimensions via une matrice al\'eatoire $\Phi$, pr\'eservant les
distances \`a un facteur $1 \pm \epsilon$ pr\`es. Que l'information extractible soit
pr\'eserv\'ee dans la m\^eme proportion rel\`eve d'une hypoth\`ese de conception que le
lemme ne garantit pas.

\subsubsection{Construction SOTP-Base}
\label{subsubsec:construction_base}

\begin{algorithm}[H]
\caption{SOTP-Base~: construction avec compression JL}
\label{alg:sotp_base}
\begin{algorithmic}[1]
\State \textbf{Param\`etres} : $\lambda$, $H^*(\lambda) \leq H(\mathcal{C}) \leq \lambda/\log\lambda$, $\varepsilon = 0{,}1$
\State \textbf{$\mathsf{Setup}(1^\lambda)$} :
\State \quad G\'en\'erer $(\mathbb{G}, q, g, h)$, PRG, Hash comme SOTP-Low
\State \quad Fixer $d \leftarrow \lceil 4\Delta H(\mathcal{C})/\varepsilon^2 \rceil$
\State \quad Fixer $k \leftarrow \lceil 8\log(2^{H(\mathcal{C})})/\varepsilon^2 \rceil$
\State \quad G\'en\'erer matrice JL $\Phi \xleftarrow{\$} \mathcal{N}(0, 1/k)^{k \times d}$~\cite{DasguptaGupta03,Vempala05}
\State \quad \textbf{Retourner} $pk = (\mathbb{G}, q, g, h, \mathsf{PRG}, \mathsf{Hash}, \Phi, k, d)$
\State \textbf{$\mathsf{Prove}(pk, m, c)$} :
\State \quad Couches Gold et Iron identiques \`a SOTP-Low
\State \quad \textcolor{blue}{// Clay JL}
\State \quad $\vec{c} \leftarrow \mathsf{embed}(c) \in \mathbb{R}^d$
\State \quad $\sigma_3 \leftarrow \Phi \vec{c} \in \mathbb{R}^k$
\State \quad \textbf{Retourner} $\sigma = (\sigma_1, \sigma_2, \sigma_3)$
\State \textbf{$\mathsf{Verify}(pk, \sigma, c)$} :
\State \quad V\'erifier $\|\Phi \cdot \mathsf{embed}(c) - \sigma_3\| \leq \varepsilon$
\State \quad V\'erifier $\sigma_2 \in \mathbb{G}$ et coh\'erence Gold
\end{algorithmic}
\end{algorithm}

SOTP-Base partage avec SOTP-Low les couches Gold et Iron inchang\'ees.
La diff\'erence enti\`ere r\'eside dans la couche Clay, qui remplace
l'encodage d\'eterministe bijectif par une compression
Johnson-Lindenstrauss. Cette substitution est rendue n\'ecessaire par
le r\'egime contextuel. Ici $H(\mathcal{C}) \geq H^*(\lambda)$, ce
qui signifie que le contexte institutionnel est trop riche pour tenir
enti\`erement dans un encodage de taille fixe sans perte.

La phase \textsf{Setup} h\'erite du groupe, du PRG et du Hash de
SOTP-Low, puis calcule deux dimensions caract\'eristiques de la
compression. La dimension cible $d = \lceil 4\Delta
H(\mathcal{C})/\varepsilon^2 \rceil$ (ligne~4) calibre l'espace de
repr\'esentation interm\'ediaire sur le produit de la capacit\'e
d'extraction institutionnelle $\Delta$ et de l'entropie contextuelle, car
plus le contexte est riche et plus l'institution peut en extraire,
plus la repr\'esentation doit \^etre large pour pr\'eserver les
distances s\'emantiques pertinentes. La dimension de sortie
$k = \lceil 8\log(2^{H(\mathcal{C})})/\varepsilon^2 \rceil$
(ligne~5) d\'etermine la taille finale de la couche Clay transmise
dans la preuve~; elle est calibr\'ee pour que le lemme JL garantisse
la pr\'eservation des distances \`a $\varepsilon = 0{,}1$ pr\`es avec
haute probabilit\'e. La matrice $\Phi$ (ligne~6) est tir\'ee
al\'eatoirement selon une gaussienne centr\'ee de variance $1/k$. 
Cette normalisation assure que les normes euclidiennes sont pr\'eserv\'ees
en esp\'erance, propri\'et\'e sur laquelle repose la
garantie JL.

La phase \textsf{Prove} (lignes~8--13) reconstruit $\sigma_1$ et
$\sigma_2$ exactement comme SOTP-Low, puis traite le contexte
diff\'eremment. La ligne~11 plonge d'abord le contexte $c$ dans
$\mathbb{R}^d$ via la fonction $\mathsf{embed}$, qui traduit la
structure s\'emantique du contexte institutionnel en vecteur de r\'eels, de sorte que
deux contextes proches s\'emantiquement produisent des vecteurs proches
au sens de la distance euclidienne. La ligne~12 applique ensuite la
compression $\sigma_3 = \Phi\vec{c} \in \mathbb{R}^k$, par laquelle le vecteur de
$d$ dimensions est projet\'e vers un espace de $k \ll d$ dimensions.
Cette projection constitue la couche Clay de SOTP-Base. La
preuve transmise $\sigma_3$ est donc la projection compress\'ee du contexte,
distincte du contexte lui-m\^eme et de son encodage exact. Les distances
s\'emantiques entre contextes distincts y sont pr\'eserv\'ees \`a
$\varepsilon$ pr\`es par la garantie JL. La pr\'eservation de l'information
extractible reste une hypoth\`ese de conception.

La phase \textsf{Verify} (lignes~14--16) refl\`ete cette logique, puisqu'elle v\'erifie la proximit\'e $\|\Phi \cdot
\mathsf{embed}(c) - \sigma_3\| \leq \varepsilon$ au lieu de l'\'egalit\'e exacte
$\Phi \cdot \mathsf{embed}(c) = \sigma_3$. Ce crit\`ere souple
est in\'evitable, car la compression avec perte $\varepsilon = 0{,}1$
signifie que le v\'erificateur ne peut retrouver $\sigma_3$ qu'\`a
$10\%$ pr\`es depuis le contexte re\c{c}u. Le choix de $\varepsilon$
est ainsi un param\`etre de conception explicite qui gouverne le
compromis entre la fid\'elit\'e de la repr\'esentation contextuelle
et la dimension $k$ de la preuve transmise. R\'eduire $\varepsilon$
augmente la fid\'elit\'e mais exige un $k$ plus grand, alourdissant
la preuve. Avec $\varepsilon = 0{,}1$, $\Delta = 56$ et $H(\mathcal{C}) = 40$, la
dimension interm\'ediaire vaut $d = \lceil 4\Delta H(\mathcal{C})/\varepsilon^2 \rceil
= 896\,000$.

\subsubsection{Propri\'et\'es de SOTP-Base}
\label{subsubsec:saturation_base}

\begin{proposition}[Propri\'et\'es cryptographiques de SOTP-Base]\label{thm:saturation_sotp_base}
Sous s\'ecurit\'e PRG~\cite{BlumMicali84} et difficult\'e DLog~\cite{DiffieHellman76}, les
couches Gold et Iron de SOTP-Base \'etant identiques \`a celles de SOTP-Low,
SOTP-Base v\'erifie les deux propri\'et\'es de la
Proposition~\ref{thm:saturation_sotp_low}, \`a savoir la confidentialit\'e calculatoire et
$\Relver = 1 - \negl(\lambda)$ pour la proc\'edure d'ouverture. De plus, pour
tous contextes $c, c'$, avec probabilit\'e au moins $1 - \delta$ sur le tirage
de $\Phi$, $\|\Phi\,\mathsf{embed}(c) - \Phi\,\mathsf{embed}(c')\|^2 \in
(1 \pm \varepsilon)\|\mathsf{embed}(c) - \mathsf{embed}(c')\|^2$
(Lemme~\ref{lem:johnson_lindenstrauss}).
\end{proposition}

\begin{proof}
Les arguments de la Proposition~\ref{thm:saturation_sotp_low} s'appliquent sans
changement aux couches Gold et Iron~; la couche Clay n'intervient ni dans la
confidentialit\'e (elle ne d\'epend que de $c$, ind\'ependant de $m$) ni dans la
liaison. La derni\`ere propri\'et\'e est le lemme de Johnson--Lindenstrauss appliqu\'e
\`a l'ensemble fini des plongements de contextes.
\end{proof}

\begin{remark}[Heuristique de conception de SOTP-Base (non prouv\'ee)]
\label{rem:heuristique_sotp_base}
La formulation initiale de ce r\'esultat \'enon\c{c}ait
$\Conf \cdot \Rel \cdot \Opp \geq 0{,}5 \cdot \min\{\Delta/d,\, \beta\log_2 d/d\}$.
La formalisation a conduit \`a la remplacer pour les raisons expos\'ees \`a la
Remarque~\ref{rem:facteur_limitant_low}. En effet, $\sigma_3 = \Phi\,\mathsf{embed}(c)$ est
une fonction du contexte public, donc $I(V;\sigma_3|C) = 0$~; la pr\'eservation
des distances n'implique pas celle de l'information extractible~; enfin une
minoration ne peut se d\'eduire d'une majoration. L'expression
$0{,}5 \cdot \min\{\Delta/d, \Delta/(\beta\log_2 d)\}$ est conserv\'ee comme
heuristique de dimensionnement de la couche Clay.
\end{remark}

\subsubsection{Support de validation~: vote \'electronique Helios}
\label{subsubsec:validation_helios}

\`A titre d'illustration, l'heuristique de SOTP-Base pour les param\`etres
analytiques de Helios~\cite{Adida08,MINKA2026CRO} ($\Delta = 56$ bits, $\beta = 1{,}4$,
$d = 560$) donne $0{,}5 \times \min\{56/560,\, 56/(1{,}4 \log_2 560)\} \approx 0{,}050$.
Cette valeur n'est pas comparable au produit CLO de Helios ($0{,}437$), qui porte
sur $\Relver$ et n'est pas une instance du th\'eor\`eme
(Remarque~\ref{rem:portee_cas_etude}). La dimension $d = 560$ retenue ici est un
choix de dimensionnement, distinct de la r\`egle $d = \lceil 4\Delta H(\mathcal{C})/\varepsilon^2\rceil$.
\end{sloppypar}

% -----------------------------------------------------------------------
\subsection{SOTP-High~: r\'egime \`a contexte riche}
\label{subsec:sotp_high}
% -----------------------------------------------------------------------
\begin{sloppypar}
Pour les contextes riches, nous proposons une construction fond\'ee sur les zk-SNARKs
succincts Groth16, qui prot\`ege le t\'emoin (divulgation nulle de connaissance) au prix
d'une opposabilit\'e r\'eduite.

\subsubsection{Motivation et champ d'application}
\label{subsubsec:motivation_high}

La conformit\'e financi\`ere via zk-SNARKs~\cite{Groth16,Garratt15,BaselIII} illustre ce r\'egime, o\`u transactions complexes,
r\'eglementations multi-juridictionnelles et historiques longs produisent $H(\mathcal{C})
\approx 70$--$90$ bits. M\^eme la compression JL de SOTP-Base serait insuffisante, car la
dimension cible $k = \Omega(H(\mathcal{C}))$ resterait trop large pour absorption par
un interpr\`ete \`a capacit\'e $\Delta$ born\'ee~\cite{Simon82,Miller56}.

Les zk-SNARKs succincts~(Groth16~\cite{Groth16,BCTV14,Gennaro13}) l\`event cet obstacle en produisant une preuve
de taille constante ($128$ \`a $192$ octets selon la courbe) ind\'ependante de la complexit\'e du
circuit v\'erifi\'e. Le caract\`ere succinct brise le couplage lin\'eaire entre $H(\mathcal{C})$ et la
taille de preuve, permettant d'encoder des relations complexes sur un contexte de
plusieurs milliers de bits dans un condens\'e que l'interpr\`ete peut v\'erifier en quelques millisecondes.

Cette solution a une double contrepartie. Une \'etape de mise en place de confiance (\emph{trusted
setup}) est requise~\cite{Bowe2017}. L'opposabilit\'e est aussi r\'eduite, plafonn\'ee par la perte
entropique du hachage compressant le contexte en 256 bits.

\subsubsection{Construction SOTP-High}
\label{subsubsec:construction_high}

\begin{algorithm}[H]
\caption{SOTP-High~: construction avec zk-SNARKs Groth16}
\label{alg:sotp_high}
\begin{algorithmic}[1]
\State \textbf{Param\`etres} : $\lambda$, $H(\mathcal{C}) > \lambda/\log\lambda$, circuit $\mathcal{C}_{\text{circuit}}$
\State \textbf{$\mathsf{Setup}(1^\lambda, \mathcal{C}_{\text{circuit}})$} :
\State \quad \'echantillonner $\tau, \alpha_s, \beta_s, \gamma, \delta \xleftarrow{\$} \mathbb{F}_p^*$ (toxiques)
\State \quad Construire $pk_{\mathsf{prove}}$ et $vk_{\mathsf{verify}}$ (CRS Groth16)~\cite{Groth16,BCTV14}
\State \quad \textcolor{red}{// D\'etruire toxiques~: $\tau, \alpha_s, \beta_s, \gamma, \delta$}
\State \quad $(g, h) \xleftarrow{\$} \mathbb{G}^2$ pour couche Iron
\State \quad \textbf{Retourner} $pk = (pk_{\mathsf{prove}}, vk_{\mathsf{verify}}, g, h)$
\State \textbf{$\mathsf{Prove}(pk, m, c)$} :
\State \quad \textcolor{blue}{// Gold + Iron (Groth16 + Pedersen)}
\State \quad $\rho \xleftarrow{\$} \mathbb{Z}_q$ ;\; $c_m \leftarrow g^m \cdot h^\rho$~\cite{Pedersen91}
\State \quad $h_c \leftarrow \mathsf{Poseidon}(c)$~\cite{Poseidon2019,Grassi2021} ;\; $\pi \leftarrow \mathsf{Groth16.Prove}(pk_{\mathsf{prove}}, (h_c, c_m), (m, c, \rho))$
\State \quad \textit{relation~: $h_c = \mathsf{Poseidon}(c) \wedge c_m = g^m h^\rho \wedge \mathsf{Conforme}(m, c) = 1$}
\State \quad \textcolor{blue}{// Clay~: condens\'e hash + m\'etadonn\'ees}
\State \quad \textbf{Retourner} $\sigma = (h_c, \pi, c_m, \mathsf{metadata})$
\State \textbf{$\mathsf{Verify}(\sigma)$} :
\State \quad $b_1 \leftarrow \mathsf{Groth16.Verify}(vk_{\mathsf{verify}}, (h_c, c_m), \pi)$
\State \quad $b_2 \leftarrow [c_m \in \mathbb{G}]$
\State \quad \textbf{Retourner} $b_1 \wedge b_2$
\end{algorithmic}
\end{algorithm}

SOTP-High est la construction la plus complexe des trois. Elle est
con\c{c}ue pour le r\'egime o\`u le contexte institutionnel est
arbitrairement riche, $H(\mathcal{C}) > \lambda/\log\lambda$, ce qui
rend impossible toute repr\'esentation fid\`ele du contexte dans la
preuve sans r\'ev\'eler de l'information confidentielle. La solution
adopt\'ee s'\'ecarte donc de SOTP-Low et SOTP-Base. Au lieu de transmettre le contexte ou sa compression, on transmet une
\emph{preuve de connaissance} du contexte via un zk-SNARK Groth16.

La phase \textsf{Setup} (lignes~2--7) est la plus sensible des trois
constructions. Elle \'echantillonne cinq scalaires $\tau, \alpha_s,
\beta_s, \gamma, \delta$ dans le corps premier $\mathbb{F}_p^*$, qui
sont les param\`etres toxiques de l'\'etape de configuration
Groth16. Ils servent \`a construire la cl\'e de preuve $pk_{\mathsf{prove}}$
et la cl\'e de v\'erification $vk_{\mathsf{verify}}$ qui forment le
Common Reference String (CRS). La ligne~5, signal\'ee en rouge, est
un imp\'eratif de s\'ecurit\'e. Si l'un quelconque de ces cinq
scalaires subsistait apr\`es le \textsf{Setup}, son d\'etenteur pourrait
forger des preuves valides pour n'importe quelle assertion, brisant
enti\`erement la propri\'et\'e de soundness du syst\`eme. En pratique,
cette destruction est r\'ealis\'ee via une \'etape multi-partite
o\`u la compromission simultan\'ee de tous les participants serait
n\'ecessaire pour reconstituer les toxiques. Le Setup produit enfin
une paire $(g, h)$ pour la couche Iron de Pedersen, ind\'ependante
du CRS Groth16.

La phase \textsf{Prove} (lignes~8--14) fusionne les trois couches
en deux op\'erations principales. Les lignes~10 et~11 r\'ealisent les
trois couches. Le hash Poseidon $h_c =
\mathsf{Poseidon}(c)$ comprime le contexte arbitrairement riche $c$
vers un condens\'e de 256 bits qui constitue l'entr\'ee publique du
circuit~; l'engagement de Pedersen $c_m = g^m \cdot h^\rho$ joue le r\^ole de
la couche Iron, puisqu'il lie le prouveur au message $m$ sans le r\'ev\'eler. La
preuve $\pi$ atteste que le prouveur conna\^it $(m, c, \rho)$ tels que
$h_c = \mathsf{Poseidon}(c)$, $c_m = g^m h^\rho$ et que le pr\'edicat de
conformit\'e $\mathsf{Conforme}(m, c)$ est satisfait~; cette relation
lie l'engagement au contexte et \`a l'assertion prouv\'ee. Cette preuve ne
r\'ev\`ele rien sur $(m, c)$ au-del\`a de ce que les entr\'ees publiques
r\'ev\`elent d\'ej\`a~; la propri\'et\'e de divulgation nulle de connaissance de
Groth16 en garantit la confidentialit\'e. La preuve finale $\sigma = (h_c, \pi, c_m,
\mathsf{metadata})$ est ainsi constitu\'ee du condens\'e contextuel,
de la preuve zk-SNARK, de l'engagement sur le message et de
m\'etadonn\'ees permettant \`a l'arbitre de situer la preuve dans son
contexte institutionnel.

La phase \textsf{Verify} (lignes~15--18) est l\'eg\`ere
au regard de la richesse du contexte original. Elle r\'eduit \`a deux
v\'erifications bool\'eennes ind\'ependantes. La premi\`ere, $b_1$, v\'erifie que la
preuve Groth16 $\pi$ est valide pour le condens\'e public $h_c$ via
la cl\'e $vk_{\mathsf{verify}}$~; la seconde, $b_2$, v\'erifie que $c_m$ est un
point l\'egitime du groupe $\mathbb{G}$. La conjonction $b_1 \wedge
b_2$ suffit \`a \'etablir que le prouveur connaissait effectivement
un message $m$ et un contexte $c$ dont $\mathsf{Poseidon}(c) = h_c$,
sans que le v\'erificateur ait acc\`es \`a $m$ ou \`a
$c$. Le co\^ut du r\'egime \`a contexte riche se situe
l\`a. L'arbitre peut v\'erifier la coh\'erence de la
preuve, mais l'information contextuelle qu'il peut extraire de $h_c$
est born\'ee par la taille du condens\'e ($256$ bits). Elle est m\^eme nulle si $c$ est
d\'ej\`a connu de l'institution. Ce que la preuve apporte au verdict tient \`a
l'assertion certifi\'ee ($\mathsf{Conforme}(m,c) = 1$), non au condens\'e.

\subsubsection{Propri\'et\'es de SOTP-High}
\label{subsubsec:saturation_high}

\begin{proposition}[Propri\'et\'es cryptographiques de SOTP-High]\label{thm:saturation_sotp_high}
Sous les hypoth\`eses q-SDH et q-PKE (\emph{knowledge soundness} de
Groth16~\cite{BonehBoyen08,BCTV14}), la difficult\'e DLog et une c\'er\'emonie de mise en
place de confiance honn\^ete~\cite{Bowe2017}~:
\begin{enumerate}
\item la preuve $\pi$ est \`a divulgation nulle de connaissance parfaite et
      $c_m$ est parfaitement masquant, de sorte que $(\pi, c_m)$ ne r\'ev\`ele rien de $m$
      au-del\`a des entr\'ees publiques~;
\item la compl\'etude est parfaite et, par la \emph{knowledge soundness}, pour
      tout prouveur PPT qui fait accepter $\sigma$, il existe un
      \emph{extracteur} (qui rembobine le prouveur ou exploite la trappe de
      l'hypoth\`ese q-PKE, contrairement au v\'erificateur) produisant un t\'emoin
      $(m, c, \rho)$ satisfaisant la relation~; d'o\`u $\Relver = 1 - \negl(\lambda)$.
\end{enumerate}
\end{proposition}

\begin{remark}[Heuristique de conception de SOTP-High (non prouv\'ee)]
\label{rem:heuristique_sotp_high}
La formulation initiale de ce r\'esultat \'enon\c{c}ait
$\Conf \cdot \Rel \cdot \Opp \geq 0{,}45 \cdot \beta\log_2\lambda/\lambda$, avec une
constante $0{,}45$ \og calibr\'ee\fg{} sur les syst\`emes d\'eploy\'es. La formalisation a
conduit \`a la remplacer. Le condens\'e $h_c = \mathsf{Poseidon}(c)$ et les
m\'etadonn\'ees sont des fonctions du contexte public ($I(V;h_c|C) = 0$)~; trois
\'el\'ements de groupe repr\'esentent $\Theta(\lambda)$ bits et non $O(\log\lambda)$~;
le v\'erificateur ne peut extraire le t\'emoin, faute de quoi la preuve ne serait
pas \`a divulgation nulle~; enfin une constante ajust\'ee sur des donn\'ees n'a pas sa place
dans un \'enonc\'e de th\'eor\`eme. L'expression $0{,}45\,\beta\log_2\lambda/\lambda$ est
conserv\'ee comme heuristique de conception~; $\beta$ y d\'esigne la constante de
r\'egularit\'e du protocole (D\'efinition~\ref{def:beta_regularite}).
\end{remark}

\subsubsection{Support de validation~: conformit\'e zk-SNARKs}
\label{subsubsec:validation_snark}

Pour les param\`etres analytiques $\Delta = 73$ bits, $\alpha = 1{,}5$, $\beta = 1{,}5$,
$\lambda = 128$ (estimations param\'etriques sans mesure empirique contr\^ol\'ee),
l'heuristique de SOTP-High donne $0{,}45 \times 1{,}5 \times 7/128 \approx 0{,}037$.
Cette valeur n'est pas comparable au produit CLO des zk-SNARKs de conformit\'e
($\approx 0{,}572$), qui porte sur $\Relver$ et n'est pas une instance du
th\'eor\`eme (Remarque~\ref{rem:portee_cas_etude}).
\end{sloppypar}

% -----------------------------------------------------------------------
\subsection{Synth\`ese et s\'election de variante}
\label{subsec:synthese_selection}
% -----------------------------------------------------------------------
\begin{sloppypar}
\subsubsection{Comparaison des trois variantes}
\label{subsubsec:tableau_comparatif}

% -----------------------------------------------------------------------
\begin{table}[htbp]
\centering
\caption{Comparaison des variantes SOTP}
\label{tab:comparaison_sotp_construction}
\resizebox{\textwidth}{!}{%
\begin{tabular}{lp{4cm}p{4cm}p{4cm}}
\toprule
\textbf{Crit\`ere} & \textbf{SOTP-Low} & \textbf{SOTP-Base} & \textbf{SOTP-High} \\
\midrule
\multicolumn{4}{l}{\textit{\textbf{R\'egime op\'erationnel}}} \\
Entropie $H(\mathcal{C})$ & $< H^*(\lambda) \approx O(\sqrt{\lambda})$ & $H^*(\lambda) \leq H(\mathcal{C}) \leq \lambda/\log\lambda$ & $> \lambda/\log\lambda$ \\
Valeurs ($\lambda = 128$) & $< 11$ bits & $11$--$18$ bits & $> 18$ bits \\
\midrule
\multicolumn{4}{l}{\textit{\textbf{Techniques cryptographiques}}} \\
Couche Gold & PRG~\cite{BlumMicali84} & PRG~\cite{BlumMicali84} & ZK parfait (Groth16)~\cite{Groth16} \\
Couche Iron & Pedersen~\cite{Pedersen91} & Pedersen~\cite{Pedersen91} & knowledge soundness (q-SDH)~\cite{BonehBoyen08} \\
Couche Clay & $\phi:\mathcal{C}\to\{0,1\}^{H(\mathcal{C})}$ & JL $\Phi:\mathbb{R}^d\to\mathbb{R}^k$~\cite{JohnsonLindenstrauss84} & Poseidon + SNARK~\cite{Poseidon2019} \\
Primitives & SHA-256, Pedersen & Gaussienne, Pedersen & Groth16, Poseidon, BLS12-381~\cite{BLS12-381} \\
\midrule
\multicolumn{4}{l}{\textit{\textbf{M\'etriques CLO atteintes}}} \\
Confidentialit\'e & calculatoire (PRG) & calculatoire (PRG) & ZK parfait \\
$\Relver$ & $1 - \epsilon_{\text{DLog}}$ & $1 - \epsilon_{\text{DLog}}$ & $1 - \epsilon_{\text{q-SDH}}$ \\
$\Opp$ & d\'epend du contexte d'esp\`ece transmis & idem & idem \\
\midrule
\multicolumn{4}{l}{\textit{\textbf{Heuristique de conception (non prouv\'ee)}}} \\
Expression & $0{,}75 \cdot \Delta/\lambda$ & $0{,}5 \cdot \min\{\Delta/d, \Delta/(\beta\log_2 d)\}$ & $0{,}45 \cdot \beta\log_2\lambda/\lambda$ \\
Trusted setup & Non & Non & Oui~\cite{Bowe2017} \\
\midrule
\multicolumn{4}{l}{\textit{\textbf{Illustration (non comparable aux produits CLO)}}} \\
Syst\`eme de r\'ef\'erence & ECDSA~\cite{Johnson01} & Helios~\cite{Adida08} & zk-SNARKs conformit\'e~\cite{Groth16} \\
Valeur de l'heuristique & $\approx 0{,}205$ & $\approx 0{,}050$ & $\approx 0{,}037$ \\
Produit CLO (\'echelle $\Relver$) & $\approx 0{,}273$ & $\approx 0{,}437$ & $\approx 0{,}572$ \\
\bottomrule
\end{tabular}%
}
\end{table}
% -----------------------------------------------------------------------

L'association de chaque syst\`eme de r\'ef\'erence \`a une variante est
illustrative. Elle suit l'ordre relatif des entropies contextuelles de ces
syst\`emes (ECDSA, puis Helios, puis la conformit\'e financi\`ere). Les seuils
$H^*(\lambda)$ et $H^{**}(\lambda)$ de
l'Algorithme~\ref{alg:selection_variante_sotp} restent des crit\`eres de
dimensionnement \`a $\lambda$ fix\'e et ne servent pas \`a classer ces syst\`emes.

Les trois variantes forment un continuum adaptatif couvrant l'ensemble du spectre
d'entropie contextuelle. SOTP-Low privil\'egie la lisibilit\'e du contexte au prix d'une restriction sur
$H(\mathcal{C})$, id\'eal pour applications juridiques ou contractuelles~\cite{Mason07,Kerr04}. SOTP-Base
\'equilibre les trois m\'etriques via compression adaptative, polyvalent pour applications
gouvernementales et \'electorales~\cite{Adida08,Goodman15}. SOTP-High sacrifie l'opposabilit\'e pour la confidentialit\'e
parfaite et la scalabilit\'e, n\'ecessaire pour applications financi\`eres ou m\'edicales \`a
contexte tr\`es riche~\cite{Garratt15,BaselIII}.

\subsubsection{Crit\`ere de s\'election de variante}
\label{subsubsec:regle_selection}

\begin{algorithm}[H]
\caption{S\'election automatique de variante SOTP}
\label{alg:selection_variante_sotp}
\begin{algorithmic}[1]
\State \textbf{Entr\'ees} : $H(\mathcal{C})$, $\Delta$, $\lambda$, $\text{Priority}$,
  $\text{TrustedSetup\_OK}$
\State $H^*(\lambda) \leftarrow \sqrt{\lambda}$ ;\; $H^{**}(\lambda) \leftarrow \lambda/\log_2\lambda$
\If{$\text{TrustedSetup\_OK} = \texttt{False}$ \textbf{and} $H(\mathcal{C}) > H^{**}(\lambda)$}
  \State \textbf{ALERTE} : SOTP-High requis mais trusted setup indisponible~\cite{Bowe2017}
  \State \textbf{Recommander} : SOTP-Base avec perte d'opposabilit\'e, ou d\'eployer MPC
  \State \textbf{Retourner} SOTP-Base
\EndIf
\If{$H(\mathcal{C}) < H^*(\lambda)$}
  \State \textbf{Retourner} SOTP-Low
\ElsIf{$H(\mathcal{C}) \leq H^{**}(\lambda)$}
  \State \textbf{Retourner} SOTP-Base ;\; $d \leftarrow \lceil 4\Delta H(\mathcal{C})/\varepsilon^2 \rceil$
\Else
  \State \textbf{Retourner} SOTP-High ;\; compiler circuit $\mathcal{C}_{\text{circuit}}$
\EndIf
\end{algorithmic}
\end{algorithm}

Cet algorithme est un m\'eta-algorithme de d\'eploiement, distinct des
constructions cryptographiques. Il prend en entr\'ee les param\`etres
du contexte institutionnel et retourne la variante SOTP adapt\'ee, en
tenant compte des propri\'et\'es math\'ematiques du contexte
et des contraintes op\'erationnelles de l'environnement de
d\'eploiement.

La ligne~2 calcule les deux seuils qui partitionnent l'espace des
entropies contextuelles. $H^*(\lambda) = \sqrt{\lambda}$ est la
fronti\`ere inf\'erieure s\'eparant le r\'egime \`a contexte rare du r\'egime
de transition tandis que $H^{**}(\lambda) = \lambda/\log_2\lambda$ est la
fronti\`ere sup\'erieure s\'eparant le r\'egime de transition du r\'egime
\`a contexte riche. Ces deux seuils sont des fonctions croissantes de
$\lambda$, car plus le param\`etre de s\'ecurit\'e est \'elev\'e, plus les
fronti\`eres se d\'eplacent vers des entropies contextuelles plus grandes,
ce qui signifie qu'une m\^eme richesse contextuelle peut relever de
r\'egimes diff\'erents selon le niveau de s\'ecurit\'e cible.

Les lignes~3--7 constituent la garde de s\'ecurit\'e op\'erationnelle, qui
est la partie la plus importante de l'algorithme du point de vue du
d\'eploiement. SOTP-High requiert une \'etape de configuration
de confiance (\emph{trusted setup}) pour g\'en\'erer le CRS Groth16~;
si cette \'etape n'a pas \'et\'e conduite~(parce que
l'organisation n'a pas les ressources pour organiser un protocole
multi-parties, ou parce que le contexte r\'eglementaire ne l'autorise
pas) et que le contexte est n\'eanmoins riche ($H(\mathcal{C}) >
H^{**}(\lambda)$), l'algorithme \'emet une alerte explicite et ne
d\'egrade pas silencieusement la s\'ecurit\'e. Il recommande alors
soit d'accepter SOTP-Base avec la perte d'opposabilit\'e
associ\'ee au r\'egime pour lequel il n'est pas optimal, soit de d\'eployer
un protocole MPC pour r\'ealiser le trusted setup de mani\`ere
distribu\'ee. Ce traitement explicite de la d\'efaillance op\'erationnelle
distingue SOTP-High des deux autres variantes, qui n'ont aucune
d\'ependance envers une \'etape externe.

Les lignes~8--14 impl\'ementent la s\'election proprement dite par
comparaison en cascade. Si $H(\mathcal{C}) < H^*(\lambda)$, le contexte
est assez rare pour que l'encodage d\'eterministe de SOTP-Low soit
suffisant, sans qu'aucun param\`etre suppl\'ementaire ne soit n\'ecessaire. Si
$H(\mathcal{C})$ est compris entre les deux seuils, SOTP-Base est
s\'electionn\'e et la dimension de compression $d = \lceil 4\Delta
H(\mathcal{C})/\varepsilon^2 \rceil$ est calcul\'ee imm\'ediatement. Ce calcul est effectu\'e \`a la s\'election et non dans le \textsf{Setup}
pour permettre d'adapter $d$ \`a la valeur effective de $H(\mathcal{C})$
observ\'ee au moment du d\'eploiement. Si $H(\mathcal{C})$ exc\`ede
$H^{**}(\lambda)$ et que le trusted setup est disponible, SOTP-High
est s\'electionn\'e et le circuit arithm\'etique $\mathcal{C}_{\text{circuit}}$
est compil\'e. Cette compilation \'etant l'op\'eration la plus co\^uteuse
du processus, elle n'est d\'eclench\'ee qu'en
dernier recours.

La logique de s\'election refl\`ete ainsi une hi\'erarchie de pr\'ef\'erences
implicite. SOTP-Low est pr\'ef\'er\'e quand il est applicable (co\^ut
minimal, pas de d\'ependance externe), SOTP-Base est le choix par d\'efaut
en r\'egime interm\'ediaire (robuste, sans trusted setup), alors que SOTP-High
n'est d\'eploy\'e que lorsque la richesse contextuelle le rend
in\'evitable et que les conditions op\'erationnelles le permettent.

\subsubsection{M\'ethodologie de conception en cinq phases}
\label{subsubsec:methodologie_cinq_phases_sotp}

La s\'election algorithmique ci-dessus s'inscrit dans une m\'ethodologie plus large, organis\'ee en cinq phases.

\textbf{Phase 1. Analyse applicative.} Estimer $H(\mathcal{C})$ et identifier les
exigences minimales $\Conf_{\min}$, $\Rel_{\min}$, $\Opp_{\min}$.

\textbf{Phase 2. \'Evaluation institutionnelle.} Estimer $\Delta$ (capacit\'e de
d\'elib\'eration, expertise, m\'emoire institutionnelle) et $\alpha$ (granularit\'e des verdicts).
Ces estimations sont analytiques, car une mesure directe via protocoles cognitifs contr\^ol\'es
reste un travail futur.

\textbf{Phase 3. Test de faisabilit\'e.} V\'erifier l'in\'egalit\'e
\begin{equation}
\Conf_{\min} \cdot \mathsf{Rel}_{\mathrm{rec},\min} \cdot \Opp_{\min} \stackrel{?}{\leq}
\frac{1+\kappa/\lambda}{\lambda}\,\min\left\{1, \frac{\kappa/\lambda + \Delta}{\gamma}\right\}
\label{eq:test_faisabilite_sotp}
\end{equation}
pour les exigences exprim\'ees en fiabilit\'e de reconstruction
(\'equation~\eqref{eq:trilemme_cro_raffine}). Si l'in\'egalit\'e est viol\'ee, il est possible de
relaxer une exigence, d'augmenter $\Delta$ par investissement dans la
formation institutionnelle (effectif seulement si $\Delta < \gamma$) ou de r\'eviser la
granularit\'e des verdicts. Les exigences exprim\'ees en fiabilit\'e de v\'erification
$\Relver$ ne sont pas contraintes par ce test.

\textbf{Phase 4. S\'election et param\'etrage.} Appliquer l'Algorithme~\ref{alg:selection_variante_sotp}
et optimiser les primitives Gold/Iron/Clay dans le r\'egime identifi\'e.

\textbf{Phase 5. Preuves et validation.} \'Etablir les r\'eductions formelles aux
hypoth\`eses standard, puis mesurer $\Delta$ et l'opposabilit\'e obtenue selon les
protocoles du chapitre~\ref{chap:experimentation}.
\end{sloppypar}

% -----------------------------------------------------------------------
\subsection{Optimalit\'e universelle et analyse de sensibilit\'e}
\label{subsec:optimalite_sensibilite}
% -----------------------------------------------------------------------
\begin{sloppypar}
\subsubsection{Conjecture d'optimalit\'e universelle}
\label{subsubsec:optimalite_universelle}

\begin{conjecture}[Optimalit\'e universelle SOTP]\label{conj:optimalite_universelle}
\label{thm:optimalite_universelle}%
Sous les Hypoth\`eses~\ref{ass:extraction_bornee}--\ref{ass:regime}, pour chaque
r\'egime de richesse contextuelle, il existe une variante SOTP et une
constante $C_R \geq 1$ ind\'ependante de $\lambda$ telles que, pour tout
protocole CLO $\mathcal{P}'$ du m\^eme r\'egime,
\begin{equation}
\Conf(\mathcal{P}')\,\Relrec(\mathcal{P}')\,\Opp(\mathcal{P}')
\;\leq\; C_R \cdot \Conf(\text{SOTP})\,\Relrec(\text{SOTP})\,\Opp(\text{SOTP}).
\label{eq:optimalite_universelle}
\end{equation}
\end{conjecture}

\begin{remark}[Statut et argument heuristique]\label{rem:statut_optimalite}
La formulation initiale de cet \'enonc\'e, pr\'esent\'ee comme th\'eor\`eme sous la forme
$\text{produit}(\mathcal{P}') \leq c \cdot \text{produit}(\text{SOTP}) + o(1)$ avec
$c \in [0{,}45;\,0{,}75]$, a \'et\'e remplac\'ee par une formulation plus pr\'ecise, pour
quatre raisons. (i)~Avec $c < 1$ et $\mathcal{P}' = \text{SOTP}$, elle donnerait
$\text{produit} \leq 0{,}75 \times \text{produit}$, ce qui impose $C_R \geq 1$~;
(ii)~le terme additif $o(1)$ domine des produits tous en $O(1/\lambda)$~;
(iii)~son argument invoquait le trilemme sans l'Hypoth\`ese~\ref{ass:regime}~;
(iv)~elle s'appuyait sur les minorations des Remarques~\ref{rem:facteur_limitant_low},
\ref{rem:heuristique_sotp_base} et~\ref{rem:heuristique_sotp_high}, qui rel\`event
d'heuristiques de conception. La forme multiplicative ci-dessus, avec $C_R \geq 1$, est l'\'enonc\'e que
l'argument vise. Celui-ci proc\`ede en trois temps, qui consistent \`a majorer le produit de tout
$\mathcal{P}'$ par la forme raffin\'ee~\eqref{eq:trilemme_cro_raffine}~; minorer celui
de la variante SOTP du r\'egime~; former le ratio. La deuxi\`eme \'etape reste ouverte, car
elle requiert un mod\`ele dans lequel la couche Clay transmet un contexte d'esp\`ece
que l'institution ne d\'etient pas (perspective d\'evelopp\'ee au
chapitre~\ref{chap:resultats_discussions}). Les rapports indicatifs que donneraient
les heuristiques de conception, $1/0{,}75 \approx 1{,}33$ (SOTP-Low), $1/0{,}5 = 2$
(SOTP-Base) et $1/0{,}45 \approx 2{,}2$ (SOTP-High), fixent l'ordre de grandeur
de $C_R$ vis\'e.
\end{remark}

\begin{boiteArbitre}
L'arbitre comprend que la recherche du meilleur protocole peut
continuer. La borne lui dit ce qui est interdit~; savoir jusqu'o\`u les
constructions SOTP l'approchent reste une question ouverte, pos\'ee avec pr\'ecision.
\end{boiteArbitre}

\subsubsection{Analyse de sensibilit\'e}
\label{subsubsec:analyse_sensibilite}

L'analyse porte sur la borne raffin\'ee~\eqref{eq:trilemme_cro_raffine}, qui est le seul
\'enonc\'e prouv\'e, ainsi que sur les heuristiques de conception.

\textbf{Sensibilit\'e \`a $\Delta$.} Dans le r\'egime \`a contexte riche ($\Delta < \gamma$),
la borne est proportionnelle \`a $\kappa/\lambda + \Delta$. Augmenter $\Delta$ de $50$ \`a
$65$ bits (formation intensive des garants~\cite{Ericsson95}) rel\`eve la borne d'environ
$30\%$. Dans le r\'egime \`a contexte rare ($\Delta \geq \gamma$), la borne est celle de
Fano et ne d\'epend pas de $\Delta$.

\textbf{Sensibilit\'e \`a $\alpha$.} Dans le r\'egime \`a contexte riche, la borne varie en
$1/\gamma = 1/(\alpha H(\mathcal{C}) - 1)$, de sorte que des verdicts plus grossiers (passer de 12
niveaux \`a 6) rel\`event la borne et peuvent faire basculer l'institution dans le
r\'egime \`a contexte rare. Les institutions \`a verdicts tr\`es nuanc\'es (\'echelles de risque
fines, montants d'amendes gradu\'es) paient ainsi une p\'enalit\'e r\'eelle. L'optimum entre
pr\'ecision d\'ecisionnelle et compromis CRO d\'epend des priorit\'es applicatives.

\textbf{Sensibilit\'e \`a $\beta$.} La constante $\beta$ n'appara\^it pas dans la borne~;
elle intervient dans la garantie de confidentialit\'e du
Lemme~\ref{lem:confidentiality_guarantee}, o\`u un protocole plus r\'egulier (plus petit
$\beta$) resserre cette garantie. Dans l'heuristique de SOTP-High, qui est
proportionnelle \`a $\beta$, passer de $\beta = 1{,}5$ \`a $\beta = 1{,}2$ diminue
la valeur de $20\%$. $\beta$ est une propri\'et\'e des lois de preuve (divergence
KL) et ne d\'epend pas de la fonction de hachage utilis\'ee dans le circuit.
\end{sloppypar}

% -----------------------------------------------------------------------
\subsection{Extension quantique des constructions SOTP}
\label{subsec:extension_postquantique_sotp}
% -----------------------------------------------------------------------
\begin{sloppypar}
Les constructions SOTP pr\'esent\'ees dans cette section s'appliquent au r\'egime classique.
La section~\ref{sec:cro_pq} a \'etabli les bornes quantiques~; 
l'extension des constructions SOTP au r\'egime quantique requiert des adaptations 
propres \`a chaque variante.

\textbf{SOTP-Low post-quantique.} La couche Gold peut int\'egrer des PRG bas\'es sur 
les fonctions de hachage (dont l'algorithme de Grover~\cite{Grover96} r\'eduit la s\'ecurit\'e de $\lambda$
\`a $\lambda/2$ bits). L'engagement de Pedersen
reste parfaitement masquant, mais son caract\`ere liant repose sur DLog, que
l'algorithme de Shor r\'esout. Il doit donc \^etre remplac\'e par un engagement
post-quantique, par exemple \`a base de r\'eseaux sous SIS/LWE~\cite{Peikert2016}.

\textbf{SOTP-Base post-quantique.} La compression JL est ind\'ependante du mod\`ele 
adversaire~\cite{JohnsonLindenstrauss84,DasguptaGupta03}. Les adaptations concernent les couches Gold et Iron avec des primitives 
post-quantiques (engagement \`a base de r\'eseaux~\cite{Peikert2016}, signatures ML-DSA~\cite{Dilithium2022}).

\textbf{SOTP-High post-quantique.} Les zk-SNARKs Groth16 n\'ecessitent une attention 
particuli\`ere~\cite{Unruh2016,Chatterjee2021}. Les hypoth\`eses q-SDH et q-PKE peuvent \^etre affaiblies face \`a des 
adversaires quantiques (elles reposent sur des couplages, que l'algorithme de Shor
compromet). Les STARKs~\cite{BenSasson2019}, qui ne reposent que sur des fonctions de
hachage, constituent l'alternative post-quantique~; les Bulletproofs~\cite{Bulletproofs2018}
reposent sur DLog et ne sont pas post-quantiques.

Le coefficient de perte d'interpr\'etabilit\'e $\eta_q$
(D\'efinition~\ref{def:perte_interp_quantique}) s'applique \`a toutes les variantes
lorsque la preuve transite par un support quantique~:
\begin{equation}
\Opp^q \geq \Opp - \eta_q\,\frac{H(C)}{H(V|C)},
\end{equation}
et $\eta_q = 0$ lorsque la preuve reste une cha\^ine de bits classique.
\end{sloppypar}

% -----------------------------------------------------------------------
% Synth\`ese
% -----------------------------------------------------------------------
\begin{boiteArbitre}
L'arbitre referme le dossier SOTP. Il vient de terminer l'\'etape \textbf{S}. 
Cette section a modifi\'e sa mani\`ere de lire le Trilemme CRO. Une borne sup\'erieure
sans construction ressemble \`a une condamnation, alors qu'accompagn\'ee de trois constructions de
r\'ef\'erence elle devient une carte, dont une partie reste \`a lever.

Quatre contributions structurent ce qu'il retient. L'architecture commune
d'abord (Gold pour le masquage, Iron pour le liage, Clay pour lui), dont
la couche Clay s'adapte au r\'egime contextuel. Les trois constructions
ensuite, calibr\'ees chacune sur ce que son $\Delta$ lui permet d'extraire
sans lui en demander davantage, avec SOTP-Low quand le contexte est rare et
l'encodage d\'eterministe suffit, SOTP-Base quand la compression
Johnson-Lindenstrauss absorbe la transition, SOTP-High quand les preuves
succinctes Groth16 compriment un contexte arbitrairement riche en quelques
centaines d'octets. Viennent enfin la conjecture d'optimalit\'e universelle, qui
dit ce qu'il reste \`a d\'emontrer, puis la m\'ethodologie de s\'election, qui
transforme le r\'esultat en proc\'edure.

Il retient surtout la franchise du propos. Les propri\'et\'es cryptographiques des
constructions sont prouv\'ees~; les facteurs $0{,}75$, $0{,}5$ et $0{,}45$ sont des
heuristiques de conception que les syst\`emes qu'il conna\^it depuis le premier
chapitre ne servent pas \`a valider.

La certification formelle du trilemme, qui l'attend encore, fait l'objet de la section suivante.
\end{boiteArbitre}
% ============================================================================
% SECTION 2.6 : FORMALISATION DU TRILEMME CRO DANS LEAN4
% Source : partie2/cro_lean4.tex
% Version : 2.2 — alignée sur la formalisation Lean4 (section~\ref{sec:lean4})
% ============================================================================

\section{FORMALISATION DU TRILEMME CRO DANS LEAN4}
\label{sec:lean4}
\begin{sloppypar}
La section~\ref{sec:sotp_construction} a pr\'esent\'e les constructions de r\'ef\'erence
SOTP.
Il reste cependant \`a savoir si les preuves th\'eoriques \'etablies dans les
chapitres pr\'ec\'edents sont correctes. Cette question
\textbf{Obj~\ref{obj:obj5_verification}} motive l'\'etape ultime de notre progression
m\'ethodologique.

Nous v\'erifions formellement le Trilemme CRO conditionnel dans
l'assistant de preuve Lean4~\cite{Lean4,moura2015lean,Avigad2019}, constituant
l'\'etape \textbf{C (Certifier)} du processus F$\to$B$\to$E$\to$S$\to$\textbf{C}$\to$X.
Notre formalisation \'etablit la borne
$\gamma^{-1}(\Delta/\lambda + \kappa(\Delta+1)/\lambda^2 + \kappa^2/\lambda^3)$
et sa forme raffin\'ee. Elle fournit aussi deux t\'emoins de coh\'erence et de port\'ee, certifie
les cons\'equences d'un verdict d\'eterministe et constitue une biblioth\`eque Lean4
r\'eutilisable pour la cryptographie institutionnelle. Cette contribution \textbf{C5}
a aussi conduit \`a des formulations plus pr\'ecises de plusieurs \'etapes des
preuves publi\'ees, pr\'esent\'ees \`a la sous-section~\ref{subsec:lean4_erreurs}.

Les sous-sections suivantes pr\'ecisent d'abord ce que cette formalisation
couvre et ses limitations (section~\ref{subsec:portee_limitations}), puis les
conventions typographiques propres \`a cette section
(section~\ref{subsec:conventions_typographiques}). Elles rappellent ensuite les
\'el\'ements du Trilemme CRO s\'electionn\'es pour la formalisation, avec les points
critiques identifi\'es lors de l'analyse pr\'eliminaire
(section~\ref{subsec:elements_formaliser}), avant d'exposer les d\'efis techniques
de la formalisation (section~\ref{subsec:defis_formalisation}), le mod\`ele formel
et l'axiomatisation (section~\ref{subsec:modele_formel_lean4}) ainsi que la preuve
du c\oe ur alg\'ebrique (section~\ref{subsec:preuve_coeur_algebrique}). Les deux
derni\`eres \'enoncent le th\'eor\`eme principal v\'erifi\'e
(section~\ref{subsec:theoreme_principal_lean4}) puis appliquent la formalisation
aux familles de protocoles canoniques
(section~\ref{subsec:application_familles_canoniques}).

Le code source v\'erifi\'e est disponible en acc\`es ouvert (licence MIT,
\url{https://github.com/MINKAThierry/PhDThesis}~\cite{CROTrilemmaDepot}), avec une
version de Lean et de Mathlib fig\'ee (v4.34.0) et une int\'egration continue qui le
recompile \`a chaque modification.
\end{sloppypar}

\begin{boiteArbitre}
L'arbitre croyait que la preuve du Trilemme \'etait compl\`ete. La machine a
montr\'e que plusieurs \'etapes des preuves publi\'ees admettaient des formulations
plus pr\'ecises, retenues dans les versions pr\'esent\'ees ici. L'arbitre comprend
que la rigueur, comme l'opposabilit\'e, se v\'erifie, gr\^ace \`a la machine et \`a
la relecture qui se compl\`etent.
\end{boiteArbitre}

% -----------------------------------------------------------------------
\subsection{Port\'ee et limitations de la formalisation}
\label{subsec:portee_limitations}
% -----------------------------------------------------------------------
\begin{sloppypar}
Cette sous-section pr\'ecise ce que couvre la formalisation Lean4 pr\'esent\'ee dans
ce chapitre, ainsi que ses limitations explicites.

\subsubsection{Ce que la formalisation couvre}
\label{subsubsec:ce_que_couvre}

Notre formalisation couvre le c\oe ur alg\'ebrique du Trilemme CRO conditionnel.
Nous d\'efinissons une structure \texttt{ProtocolFamily} index\'ee par le
param\`etre de s\'ecurit\'e $\lambda$, avec des param\`etres institutionnels
fixes $(\Delta, \alpha, \beta, H(C))$. Nous
int\'egrons comme propri\'et\'es des instances de protocole les quatre pr\'emisses
du th\'eor\`eme, \`a savoir l'in\'egalit\'e de
Fano~\cite{Fano61,CoverThomas06}, l'in\'egalit\'e de Pinsker (sous-sous-section~\ref{subsubsec:pinsker_classique}) combin\'ee \`a la
r\'egularit\'e~\cite{Pinsker1964,Tsybakov2009}, l'hypoth\`ese d'extraction
institutionnelle~\cite{Simon82,Miller56} et la borne fiabilit\'e-information.
Nous prouvons formellement que, pour tout $I \le \kappa/\lambda$ (r\'egime de
pr\'eservation de la confidentialit\'e, Hypoth\`ese~\ref{ass:regime}), on a
\begin{equation}
\Conf \cdot \Rel \cdot \Opp \le \frac{1}{\gamma}\left(\frac{\Delta}{\lambda}
+ \frac{\kappa(\Delta+1)}{\lambda^2} + \frac{\kappa^2}{\lambda^3}\right)
\le \frac{\Delta}{\lambda}
+ \frac{\kappa(\Delta+1)}{\lambda^2} + \frac{\kappa^2}{\lambda^3},
\end{equation}
ainsi que la forme raffin\'ee~\eqref{eq:trilemme_cro_raffine}. Ces preuves
n'utilisent que l'alg\`ebre \'el\'ementaire et sont enti\`erement v\'erifi\'ees par
Lean4. Nous prouvons \'egalement que $(1/2)^\lambda$ est n\'egligeable, par un
passage \`a la limite (\texttt{Tendsto}) et un argument de filtre, ce qui comble
une lacune fr\'equente dans les expositions informelles.

\subsubsection{Ce que la formalisation ne couvre pas}
\label{subsubsec:ce_que_ne_couvre_pas}

Nous explicitons ci-dessous ce qui
n'est pas inclus dans cette version v\'erifi\'ee.

\textbf{Les axiomes (Fano, Pinsker) ne sont pas prouv\'es
ab initio.} Ils sont int\'egr\'es comme
propri\'et\'es des instances de protocole, laissant leur d\'erivation formelle \`a partir de la th\'eorie des
probabilit\'es comme travail futur.

\textbf{Les constructions SOTP ne sont pas encore couvertes par Lean4.}
La formalisation couvre le noyau alg\'ebrique du Trilemme, mais pas les
r\'eductions cryptographiques des constructions SOTP. Ces extensions
constituent une direction de recherche future.

\textbf{L'extension quantique du Trilemme CRO n'est pas couverte.}
La formalisation pr\'esent\'ee ici se limite au cadre classique. L'extension
aux syst\`emes post-quantiques n\'ecessite une
infrastructure th\'eorique distincte incluant les in\'egalit\'es de Pinsker
quantique, le lemme de mesure douce (sous-sous-section~\ref{subsubsec:gentle_measurement}) et les bornes de Fano
quantique. Cette extension est rel\'egu\'ee aux travaux futurs.
\end{sloppypar}

% -----------------------------------------------------------------------
\subsection{Conventions typographiques et linguistiques}
\label{subsec:conventions_typographiques}
% -----------------------------------------------------------------------
\begin{sloppypar}
Cette sous-section pr\'esente une sp\'ecificit\'e li\'ee \`a la nature de son
contenu. Le corps du texte, les d\'efinitions math\'ematiques, les
th\'eor\`emes et les explications utilisent la ponctuation et les conventions
typographiques fran\c{c}aises (espaces avant les deux-points, guillemets
fran\c{c}ais, etc.), en coh\'erence avec les autres chapitres de cette th\`ese.

En revanche, les listings de code Lean4 pr\'eservent la syntaxe anglaise de
ce langage de programmation pour garantir la validit\'e technique du code.
Les commentaires inclus dans ces listings sont r\'edig\'es en fran\c{c}ais pour
faciliter la compr\'ehension tout en maintenant la compatibilit\'e syntaxique.
Cette convention permet au code d'\^etre directement compilable par
l'assistant de preuve Lean4 tout en restant accessible pour un lecteur
francophone.
\end{sloppypar}

% -----------------------------------------------------------------------
\subsection{\'El\'ements du Trilemme CRO \`a formaliser}
\label{subsec:elements_formaliser}
% -----------------------------------------------------------------------
\begin{sloppypar}
Cette sous-section rappelle les \'el\'ements du Trilemme CRO conditionnel,
\'etabli \`a la section~\ref{sec:cro_classique}, qui seront formalis\'es dans
Lean4. Nous pr\'ecisons les simplifications op\'er\'ees par rapport \`a la
version th\'eorique, justifions le p\'erim\`etre retenu et identifions les
points critiques que la formalisation r\'ev\`ele.

\subsubsection{M\'etriques CLO \`a formaliser}
\label{subsubsec:metriques_formaliser}

Le cadre CLO d\'efinit trois
m\'etriques pour un protocole de preuve $\mathcal{P}$ avec espace de messages
$\mathcal{M}_\lambda$ ($|\mathcal{M}_\lambda| = 2^\lambda$), espace de
contexte $\mathcal{C}_\lambda$ et espace de preuves $\Sigma_\lambda$.

\begin{definition}[Confidentialit\'e]\label{def:conf_lean4}
\begin{equation}
\Conf(\mathcal{P}) = 1 - \max_{\substack{m_0,m_1\in\mathcal{M}_\lambda\\c\in\mathcal{C}_\lambda}}
   \mathrm{TV}(P_{\Sigma|M=m_0,C=c}, P_{\Sigma|M=m_1,C=c})
\end{equation}
o\`u $\mathrm{TV}$ est la distance de variation totale.
\end{definition}

\begin{definition}[Fiabilit\'e]\label{def:rel_lean4}
Soit $\hat{M}$ l'estimateur MAP de $M$ \`a partir de $(\Sigma, C)$. La fiabilit\'e
est alors d\'efinie par
\begin{equation}
\Rel(\mathcal{P}) = 1 - \Pr[\hat{M} \neq M]
\end{equation}
\end{definition}

\begin{definition}[Opposabilit\'e]\label{def:opp_lean4}
Soit $V: \Sigma \times \mathcal{C} \to \mathcal{Y}$ une fonction de
verdict. L'opposabilit\'e est alors d\'efinie par
\begin{equation}
\Opp(\mathcal{P}) = \frac{I(V;\Sigma|C)}{H(V|C)}
\end{equation}
o\`u $I$ et $H$ d\'esignent respectivement l'information mutuelle
conditionnelle et l'entropie conditionnelle.
\end{definition}

\subsubsection{Hypoth\`eses institutionnelles s\'electionn\'ees}
\label{subsubsec:hypotheses_formaliser}

Le trilemme repose sur quatre pr\'emisses, dont trois portent sur
l'institution interpr\`ete et le protocole, la quatri\`eme reliant fiabilit\'e et
information. Nous les formalisons telles quelles, sans simplification.
% Note de source : les labels LaTeX de ces quatre enonces ne correspondent pas
% a leur contenu ; seuls les titres font foi.

\begin{assumption}[Extraction born\'ee]\label{ass:fano_axiom}
Pour toute institution avec fonction de verdict $V$, on a
\begin{equation}
I(V;\Sigma|C) \le I(M;\Sigma|C) + \Delta
\end{equation}
o\`u $\Delta \ge 0$ quantifie la connaissance pr\'ealable ou la capacit\'e de
raisonnement au-del\`a du contexte formel $C$.
\end{assumption}

\begin{assumption}[Espace de verdict born\'e]\label{ass:pinsker_axiom}
La taille de l'espace de verdict satisfait
\begin{equation}
|\mathcal{Y}_\lambda| \le 2^{\alpha H(C)}
\end{equation}
pour une constante $\alpha > 0$ mesurant la granularit\'e de d\'ecision.
\end{assumption}

\begin{assumption}[R\'egularit\'e]\label{ass:extraction_axiom}
Un protocole $\mathcal{P}$ est $\beta$-r\'egulier si pour tout contexte $c$
et messages $m_0, m_1$, on a
\begin{equation}
D_{KL}(P_{\Sigma|m_0,c} \| P_{\Sigma|m_1,c}) \le \beta \cdot I(M;\Sigma|C)
\end{equation}
avec $\beta \ge 1$.
\end{assumption}

\begin{assumption}[Fiabilit\'e-information]\label{ass:reliability_axiom}
Pour toute instance de protocole, sans restriction sur $\Rel$, on a
\begin{equation}
\Rel \le \frac{I(M;\Sigma|C) + 1}{\lambda}
\end{equation}
La formulation initiale imposait la restriction $\Rel \ge 1/2$. La
formalisation a montr\'e qu'elle peut \^etre lev\'ee, car la majoration
$h(p_e) \le 1$ de l'entropie binaire vaut pour toute probabilit\'e d'erreur.
Le code Lean4 n'impose que $0 \le \Rel \le 1$.
\end{assumption}

\subsubsection{\'Enonc\'e du th\'eor\`eme original}
\label{subsubsec:enonce_theoreme_original}

\begin{theorem}[Trilemme CRO conditionnel]\label{thm:original_lean4}
Sous les Hypoth\`eses~\ref{ass:fano_axiom},~\ref{ass:pinsker_axiom} et
pour un protocole $\beta$-r\'egulier avec messages uniformes
($H(M)=\lambda$) et $\lambda \ge 64$, on a
\begin{equation}
\Conf(\mathcal{P}) \cdot \Rel(\mathcal{P}) \cdot \Opp(\mathcal{P})
   \le \frac{\Delta}{\lambda} + O(\lambda^{-2}) + \negl(\lambda)
\end{equation}
\end{theorem}

La preuve originale proc\`ede en
quatre \'etapes successives.
\begin{enumerate}
    \item L'in\'egalit\'e de Fano~\cite{Fano61,CoverThomas06} donne
          $I(M;\Sigma|C) \ge \Rel \cdot \lambda - 1$
    \item L'in\'egalit\'e de Pinsker et la
          r\'egularit\'e~\cite{Pinsker1964,Tsybakov2009} donnent
          $\Conf \ge 1 - \sqrt{c_0 \beta I}$
    \item Les hypoth\`eses d'extraction et de
          verdict~\cite{Simon82,Miller56} donnent
          $\Opp \le (I+\Delta)/\gamma$
    \item Le produit est combin\'e sous le r\'egime
          $I \le \kappa/\lambda$ (Hypoth\`ese~\ref{ass:regime}).
\end{enumerate}

\subsubsection{Points critiques et notre approche}
\label{subsubsec:points_critiques}

Notre analyse pr\'eliminaire a identifi\'e trois points qui demandent un
traitement rigoureux.

\paragraph{Statut de $I$ comme variable.}
L'information mutuelle $I(M;\Sigma|C)$ n'est pas une variable libre, puisqu'elle
est d\'etermin\'ee par le protocole. Le r\'egime de pr\'eservation de la
confidentialit\'e $I(\lambda) \le \kappa/\lambda$ (Hypoth\`ese~\ref{ass:regime})
est rendu explicite comme condition du th\'eor\`eme.

\paragraph{Nature asymptotique exacte.}
Nous rempla\c{c}ons le $O(\lambda^{-2})$ informel par une in\'egalit\'e
exacte faisant appara\^itre une constante $K$, une fonction
n\'egligeable $\varepsilon$ et un seuil $N$ explicites. Notre analyse
alg\'ebrique produit la constante $\kappa(\Delta+1) + \kappa^2$,
rendant explicite la d\'ependance en $\kappa$.

\paragraph{Constantes implicites.}
La notation $O(\lambda^{-2})$ dissimule la d\'ependance en
$1/(\alpha H+1)$ et $\beta$. Notre borne explicite
$\gamma^{-1}(\Delta/\lambda + \kappa(\Delta+1)/\lambda^2 + \kappa^2/\lambda^3)$
rend mesurables les param\`etres qui gouvernent la borne~\cite{Grossklags10,Herley09}.

Notre formalisation prouve la borne explicite
$O(\lambda^{-2})$ sous le r\'egime de confidentialit\'e
(Hypoth\`ese~\ref{ass:regime}), enti\`erement v\'erifi\'ee.

\subsubsection{Simplifications et justifications}
\label{subsubsec:simplifications}

\textbf{Simplification~1. Int\'egration des propri\'et\'es informationnelles
comme champs de structure.} Les th\'eor\`emes profonds (Fano, Pinsker) ne sont pas
d\'eclar\'es comme \texttt{axiom} globaux. Ils sont
int\'egr\'es comme propri\'et\'es des instances de protocole, rendant ces
propri\'et\'es falsifiables et v\'erifiables instance par instance.

\textbf{Simplification~2. R\'egime de confidentialit\'e explicite.} Le r\'egime
$I \le \kappa/\lambda$ est pos\'e explicitement comme condition du
th\'eor\`eme, garantissant la confidentialit\'e via le
Lemme~\ref{lem:confidentiality_guarantee}.
\end{sloppypar}

% -----------------------------------------------------------------------
\subsection{D\'efis de formalisation dans Lean4}
\label{subsec:defis_formalisation}
% -----------------------------------------------------------------------
\begin{sloppypar}
Le Trilemme CRO conditionnel pr\'esente plusieurs d\'efis techniques qui
rendent sa v\'erification informelle p\'erilleuse.
Cette sous-section analyse ces d\'efis et les solutions adopt\'ees.

\subsubsection{D\'efi~1~: mod\'eliser les asymptotiques}
\label{subsubsec:defi_asymptotiques}

Le th\'eor\`eme original m\'elange une in\'egalit\'e exacte (valide pour
chaque $\lambda$) avec une notation asymptotique ($O(\lambda^{-2})$,
$\negl(\lambda)$). Pour formaliser cela, nous
introduisons la notion de \textbf{famille de protocoles}, que d\'efinit la
structure suivante.

\begin{leancode}[caption={Famille de protocoles -- structure ProtocolFamily en Lean4}, label={lst:lean4_protocol_family}]
-- Famille de protocoles indexee par le parametre de securite
structure ProtocolFamily where
  params  : InstitutionalParams                -- Parametres institutionnels fixes
  inst    : (lam : ℕ) → ProtocolInstance lam   -- Instance pour chaque lambda
  hparams : ∀ lam, (inst lam).params = params  -- Parametres communs
\end{leancode}

Le th\'eor\`eme devient alors un \'enonc\'e existentiel sur une
constante $K$ et une fonction n\'egligeable $\varepsilon$, valable pour tout
$\lambda \geq 64$, comme l'exprime l'\'enonc\'e ci-dessous.

\begin{leancode}[caption={\'Enonc\'e asymptotique du th\'eor\`eme principal en Lean4}, label={lst:main_theorem_statement}]
-- Theoreme du Trilemme CRO conditionnel
theorem conditional_CRO_trilemma (F : ProtocolFamily) (κ : ℝ) (hκ : 0 ≤ κ)
    (hA : ∀ lam, InstitutionalAssumptions (F.inst lam))
    (hA4 : ∀ lam, (F.inst lam).I ≤ κ / lam) :
    ∃ K : ℝ, ∃ ε : ℕ → ℝ, Negligible ε ∧ ∀ lam : ℕ, 64 ≤ lam →
      (F.inst lam).Conf * (F.inst lam).Rel * (F.inst lam).Opp ≤
        F.params.Δ / lam + K / (lam : ℝ) ^ 2 + ε lam
\end{leancode}

\subsubsection{D\'efi~2~: int\'egrer les r\'esultats profonds}
\label{subsubsec:defi_axiomatiser}

Fano~\cite{Fano61} et Pinsker~\cite{Pinsker1964} sont des th\'eor\`emes
classiques de th\'eorie de l'information. Nous les int\'egrons comme
champs de la structure \texttt{ProtocolInstance}, de sorte que chaque
instance de protocole doit fournir la preuve que ces in\'egalit\'es
s'appliquent \`a ses param\`etres.

\begin{leancode}[caption={Propri\'et\'es informationnelles comme champs de structure}, label={lst:lean4_protocol_instance}]
structure ProtocolInstance (lam : ℕ) where
  params : InstitutionalParams
  Conf   : ℝ   -- confidentialite
  Rel    : ℝ   -- fiabilite
  Opp    : ℝ   -- opposabilite
  I      : ℝ   -- I(M;Σ|C)
  IV     : ℝ   -- I(V;Σ|C)
  HV     : ℝ   -- H(V|C)
  TVmax  : ℝ   -- max TV
  KLmax  : ℝ   -- max D_KL
  hConf  : Conf = 1 - TVmax
  hOpp   : Opp = IV / HV
  hTV    : 0 ≤ TVmax ∧ TVmax ≤ 1
  hRel   : 0 ≤ Rel ∧ Rel ≤ 1
  hI     : 0 ≤ I
  hIV    : 0 ≤ IV
  pinsker : TVmax ^ 2 ≤ c0 * KLmax             -- Lemme 1 (Pinsker)
  fano    : I ≥ Rel * lam - 1                   -- Lemme 2 (Fano), etape 1
\end{leancode}

\subsubsection{D\'efi~3~: r\'egime de confidentialit\'e explicite}
\label{subsubsec:defi_optimisation}

Le r\'egime $I = O(1/\lambda)$ est rendu explicite comme condition du
th\'eor\`eme (Hypoth\`ese~\ref{ass:regime}). La preuve du
Th\'eor\`eme~\ref{thm:trilemme_cro_conditionnel}
utilise directement cette hypoth\`ese sans n\'ecessiter d'optimisation
analytique sur $[0, \lambda]$.

\subsubsection{D\'efi~4~: prouver la n\'egligeabilit\'e}
\label{subsubsec:defi_negligeabilite}

La d\'efinition standard en
cryptographie, \`a savoir
\begin{equation}
f \text{ est n\'egligeable} \iff \forall c>0,\; \exists N,\; \forall n \ge N,\;
|f(n)| \le n^{-c}
\end{equation}
se formalise, avec le filtre \texttt{atTop} de Mathlib (\og pour tout $n$
suffisamment grand\fg~\cite{Mathlib,MathlibAnalysis}) de la mani\`ere suivante.

\begin{leancode}[caption={D\'efinition de n\'egligeabilit\'e en Lean4}, label={lst:lean4_negligible}]
-- Une fonction est negligeable si elle decroit plus vite que tout polynome
def Negligible (f : ℕ → ℝ) : Prop :=
  ∀ c : ℝ, 0 < c → ∀ᶠ n : ℕ in atTop, |f n| ≤ (n : ℝ) ^ (-c)
\end{leancode}

Prouver que $\varepsilon_0(\lambda) = (1/2)^\lambda$ est n\'egligeable
n\'ecessite de montrer que $2^n$ domine $n^c$. La preuve passe par la limite
$n^{\lceil c\rceil}/2^n \to 0$ (lemme \texttt{tendsto\_pow\_const\_div\_const\_pow\_of\_one\_lt}
de Mathlib) et un argument de filtre~; elle ne fournit pas de seuil explicite.
\end{sloppypar}

% -----------------------------------------------------------------------
\subsection{Mod\`ele formel et axiomatisation}
\label{subsec:modele_formel_lean4}
% -----------------------------------------------------------------------
\begin{sloppypar}
\subsubsection{Param\`etres institutionnels}
\label{subsubsec:params_institutionnels_lean4}

Les param\`etres institutionnels sont des constantes fixes pour une
institution donn\'ee, ind\'ependantes de $\lambda$.

\begin{leancode}[caption={Structure InstitutionalParams -- param\`etres institutionnels en Lean4}, label={lst:lean4_institutional_params}]
-- Parametres institutionnels : capacites cognitives et procedurales
structure InstitutionalParams where
  H : ℝ      -- H(C) : entropie contextuelle (bits)
  Δ : ℝ      -- Capacite d'extraction institutionnelle (bits)
  α : ℝ      -- Granularite de verdict (α > 0)
  β : ℝ      -- Constante de regularite du protocole (β ≥ 1)
  hH_nonneg : 0 ≤ H
  hΔ_nonneg : 0 ≤ Δ
  hα_pos    : 0 < α
  hβ_ge_one : 1 ≤ β
\end{leancode}

Cette structure capture les trois familles canoniques~\cite{Adida08,Groth16,Johnson01},
instanci\'ees ci-dessous.

\begin{leancode}[caption={Instances canoniques -- Helios, zk-SNARKs, audit traditionnel}, label={lst:lean4_instances_canoniques}]
noncomputable def HeliosParams : InstitutionalParams :=
  { H := 50, Δ := 56, α := 1.2, β := 1.4, ... }

noncomputable def zkSNARKParams : InstitutionalParams :=
  { H := 70, Δ := 73, α := 1.5, β := 1.5, ... }

noncomputable def TraditionalAuditParams : InstitutionalParams :=
  { H := 20, Δ := 35, α := 1.0, β := 1.0, ... }
\end{leancode}

\subsubsection{Instance de protocole}
\label{subsubsec:instance_protocole_lean4}

\begin{leancode}[caption={Structure ProtocolInstance pour un $\lambda$ donn\'e}, label={lst:lean4_protocol_instance_complete}]
structure ProtocolInstance (lam : ℕ) where
  params : InstitutionalParams
  Conf   : ℝ   -- confidentialite
  Rel    : ℝ   -- fiabilite
  Opp    : ℝ   -- opposabilite
  I      : ℝ   -- I(M;Σ|C)
  IV     : ℝ   -- I(V;Σ|C)
  HV     : ℝ   -- H(V|C)
  TVmax  : ℝ   -- max TV
  KLmax  : ℝ   -- max D_KL
  hConf  : Conf = 1 - TVmax
  hOpp   : Opp = IV / HV
  hTV    : 0 ≤ TVmax ∧ TVmax ≤ 1
  hRel   : 0 ≤ Rel ∧ Rel ≤ 1
  hI     : 0 ≤ I
  hIV    : 0 ≤ IV
  pinsker : TVmax ^ 2 ≤ c0 * KLmax
  fano    : I ≥ Rel * lam - 1
\end{leancode}

\subsubsection{Hypoth\`eses institutionnelles}
\label{subsubsec:axiomes_info_lean4}

\begin{leancode}[caption={Hypoth\`eses institutionnelles en Lean4}, label={lst:lean4_axiomes_extraction}]
structure InstitutionalAssumptions {lam : ℕ} (P : ProtocolInstance lam) : Prop where
  /-- Hypothese 1 (extraction bornee) : I(V;Σ|C) ≤ I(M;Σ|C) + Δ. -/
  A1 : P.IV ≤ P.I + P.params.Δ
  /-- Hypothese 2' (verdict efficace) : H(V|C) ≥ γ = αH(C) − 1 … -/
  A2 : gamma P.params ≤ P.HV
  /-- … avec γ ≥ 1. -/
  A2γ : 1 ≤ gamma P.params
  /-- Hypothese 3 (β-regularite) : D_KL ≤ β·I(M;Σ|C). -/
  A3 : P.KLmax ≤ P.params.β * P.I
\end{leancode}
\end{sloppypar}

% -----------------------------------------------------------------------
\subsection{Preuve du c\oe ur alg\'ebrique}
\label{subsec:preuve_coeur_algebrique}
% -----------------------------------------------------------------------
\begin{sloppypar}
\subsubsection{La constante de Pinsker}
\label{subsubsec:constante_pinsker}

\begin{leancode}[caption={Constante $c_0 = 1/(2\ln 2)$ -- d\'efinition et propri\'et\'es}, label={lst:lean4_c0}]
noncomputable def c0 : ℝ := 1 / (2 * Real.log 2)

lemma c0_pos : 0 < c0 := by
  unfold c0
  have h : 0 < Real.log 2 := Real.log_pos (by norm_num)
  positivity

lemma c0_gt_half : 1 / 2 < c0 := by
  unfold c0
  have h1 := Real.log_two_lt_d9
  have h2 : 0 < Real.log 2 := Real.log_pos (by norm_num)
  rw [lt_div_iff₀ (by positivity)]
  nlinarith
\end{leancode}

\subsubsection{N\'egligeabilit\'e de $\varepsilon_0$}
\label{subsubsec:negligeabilite_eps0}

\begin{lemma}[N\'egligeabilit\'e de $\varepsilon_0$~\cite{Goldreich01,KatzLindell14}]\label{lem:negligeabilite_eps0}
La fonction $\varepsilon_0(\lambda) = (1/2)^\lambda$ est
n\'egligeable.
\end{lemma}

\begin{leancode}[caption={Preuve de n\'egligeabilit\'e en Lean4}, label={lst:lean4_eps0}]
theorem negligible_eps0 : Negligible ε₀ := by
  intro c hc
  have hlim : Tendsto (fun n : ℕ => (n : ℝ) ^ (⌈c⌉₊) / (2 : ℝ) ^ n) atTop (nhds 0) :=
    tendsto_pow_const_div_const_pow_of_one_lt ⌈c⌉₊ (by norm_num)
  have hev : ∀ᶠ n : ℕ in atTop, (n : ℝ) ^ (⌈c⌉₊) / (2 : ℝ) ^ n < 1 :=
    hlim.eventually (gt_mem_nhds one_pos)
  filter_upwards [hev, eventually_ge_atTop 1] with n hn hn1
  -- (suite de la preuve : annexe B, listing complet)
\end{leancode}

\subsubsection{Lemme~4~: garantie de confidentialit\'e}
\label{subsubsec:lemme4_lean4}

\begin{leancode}[caption={Lemme~4 -- garantie de confidentialit\'e en Lean4}, label={lst:lean4_lemma4}]
theorem lemma4_confidentiality_guarantee {lam : ℕ} (P : ProtocolInstance lam)
    (hA3 : P.KLmax ≤ P.params.β * P.I) :
    1 - Real.sqrt (c0 * P.params.β * P.I) ≤ P.Conf := by
  have hTV0 := P.hTV.1
  have h1 : P.TVmax ≤ Real.sqrt (c0 * P.KLmax) := by
    have := Real.abs_le_sqrt P.pinsker
    rwa [abs_of_nonneg hTV0] at this
  have h2 : Real.sqrt (c0 * P.KLmax) ≤ Real.sqrt (c0 * P.params.β * P.I) := by
    apply Real.sqrt_le_sqrt
    rw [mul_assoc]
    exact mul_le_mul_of_nonneg_left hA3 c0_pos.le
  rw [P.hConf]
  linarith
\end{leancode}
\end{sloppypar}

% -----------------------------------------------------------------------
\subsection{Th\'eor\`eme principal}
\label{subsec:theoreme_principal_lean4}
% -----------------------------------------------------------------------

\subsubsection{\'Enonc\'e formel}
\label{subsubsec:enonce_formel_lean4}

\begin{sloppypar}
\begin{theorem}[Trilemme CRO conditionnel~: version v\'erifi\'ee]
\label{thm:conditional_CRO_trilemma}
Soit $F$ une famille de protocoles satisfaisant les conditions suivantes.
\begin{itemize}
    \item les Hypoth\`eses~\ref{ass:extraction_bornee}, \ref{ass:verdicts_efficace}
          et~\ref{ass:regularite} pour chaque instance (structure
          \texttt{InstitutionalAssumptions}, avec $\gamma = \alpha H(C) - 1 \geq 1$)~;
    \item $I_F(\lambda) \leq \kappa/\lambda$ pour un $\kappa \geq 0$
          ind\'ependant de $\lambda$ (r\'egime de pr\'eservation de la
          confidentialit\'e, Hypoth\`ese~\ref{ass:regime})~;
    \item des param\`etres institutionnels $(H(C),\Delta,\alpha,\beta)$ communs \`a
          toute la famille (structure \texttt{ProtocolFamily}).
\end{itemize}
Alors, l'in\'egalit\'e de Fano \'etant appliqu\'ee sous la forme
$I(M;\Sigma|C) \geq \Rel\cdot\lambda - 1$, valable pour tout $\Rel$, il existe une
constante $K$ et une fonction n\'egligeable $\varepsilon$ telles que, pour tout
$\lambda \geq 64$, on ait
\begin{equation}
\Conf_F(\lambda) \cdot \Rel_F(\lambda) \cdot \Opp_F(\lambda) \leq \frac{\Delta}{\lambda} + \frac{K}{\lambda^2} + \varepsilon(\lambda)
\label{eq:main_theorem_lean4}
\end{equation}
Instance par instance (\texttt{trilemma\_explicit}), on obtient la forme avec le facteur $1/\gamma$
\begin{equation}
\Conf_F(\lambda) \cdot \Rel_F(\lambda) \cdot \Opp_F(\lambda) \leq \frac{1}{\gamma}\left(\frac{\Delta}{\lambda} + \frac{\kappa(\Delta+1)}{\lambda^2} + \frac{\kappa^2}{\lambda^3}\right)
\label{eq:main_theorem_lean4_explicit}
\end{equation}
\end{theorem}
\end{sloppypar}

\subsubsection{T\'emoins de coh\'erence et de port\'ee}
\label{subsubsec:temoins_lean4}

\begin{sloppypar}
La formalisation Lean4 produit deux t\'emoins qui d\'emontrent la
non-vacuit\'e et la n\'ecessit\'e des hypoth\`eses du th\'eor\`eme.

\begin{theorem}[Coh\'erence du mod\`ele abstrait]
\label{thm:witness_consistency}
Il existe une instance du mod\`ele abstrait v\'erifiant simultan\'ement toutes
les hypoth\`eses, avec $\Conf = \Opp = 1$, $\Rel = 1/\lambda$, $I = 0$,
$\Delta = \gamma = 1$. Pour cette instance, on a
\begin{equation}
\Conf \cdot \Rel \cdot \Opp = \frac{1}{\lambda} = \frac{\Delta}{\gamma \lambda}
\end{equation}
Cette valeur \'egale exactement la borne du
Th\'eor\`eme~\ref{thm:conditional_CRO_trilemma}. Cela d\'emontre que les
hypoth\`eses sont non vides et que la borne est atteinte \emph{dans le
mod\`ele abstrait}.
\end{theorem}

\begin{remark}[Port\'ee du t\'emoin de coh\'erence]\label{rem:portee_temoin}
Une \texttt{ProtocolInstance} est un point du mod\`ele abstrait, c'est-\`a-dire un ensemble de r\'eels
qui satisfont les contraintes de Pinsker, de Fano (sous la forme rel\^ach\'ee
$I \geq \Rel\cdot\lambda - 1$) et les Hypoth\`eses~1 \`a~4. Le t\'emoin \'etablit la
non-vacuit\'e de ce mod\`ele. L'existence d'une loi de probabilit\'e qui le
r\'ealise est une autre question. Pour une loi r\'eelle avec $M$ uniforme sur $\lambda$ bits, $I = 0$
impose $\Relrec = 2^{-\lambda}$ au lieu de $1/\lambda$. Le t\'emoin atteint donc la borne
rel\^ach\'ee~; il ne montre pas qu'un protocole r\'eel l'atteint. Il a $\Delta = \gamma = 1$
et r\'ealise l'\'egalit\'e dans la forme rel\^ach\'ee de Fano. Il illustre ainsi le r\'egime
$\Delta \geq \gamma$ de la Remarque~\ref{rem:activite_delta}. Le t\'emoin de port\'ee
(Th\'eor\`eme~\ref{thm:scope_necessity}) est, lui, r\'ealis\'e par une loi explicite
(Remarque~\ref{rem:scope_trilemme_classique}).
\end{remark}

\begin{leancode}[caption={T\'emoin de coh\'erence~: instance atteignant la borne}, label={lst:lean4_witness_instance}]
theorem assumptions_consistent_and_bound_attained (lam : ℕ) (hlam : 1 ≤ lam) :
    InstitutionalAssumptions (witnessInstance lam hlam) ∧
    (witnessInstance lam hlam).I ≤ 0 / lam ∧
    (witnessInstance lam hlam).Conf * (witnessInstance lam hlam).Rel *
      (witnessInstance lam hlam).Opp =
      (witnessInstance lam hlam).params.Δ / (gamma (witnessInstance lam hlam).params * lam) := by
  refine ⟨⟨?_, ?_, ?_, ?_⟩, ?_, ?_⟩ <;>
    simp [witnessInstance, WitnessParams, gamma] <;> norm_num
\end{leancode}

\begin{theorem}[Port\'ee~: n\'ecessit\'e du r\'egime]
\label{thm:scope_necessity}
Il existe une instance de protocole v\'erifiant les Hypoth\`eses~1, 2$'$
et 3 avec $\Delta = 0$ mais hors du r\'egime de l'Hypoth\`ese~\ref{ass:regime}
($I = \lambda/2 - 1$). Pour cette instance, on a
\begin{equation}
\Conf \cdot \Rel \cdot \Opp = \frac{1}{4}
\end{equation}
pour tout $\lambda \geq 64$, ce produit restant constant. Cela d\'emontre que le r\'egime
$I \leq \kappa/\lambda$ est n\'ecessaire \`a la borne.
\end{theorem}

\begin{leancode}[caption={T\'emoin de port\'ee~: instance hors du r\'egime}, label={lst:lean4_scope_instance}]
theorem regime_assumption_necessary (lam : ℕ) (hlam : 64 ≤ lam) :
    InstitutionalAssumptions (scopeInstance lam hlam) ∧
    (scopeInstance lam hlam).params.Δ = 0 ∧
    (scopeInstance lam hlam).Conf * (scopeInstance lam hlam).Rel *
      (scopeInstance lam hlam).Opp = 1 / 4 ∧
    (scopeInstance lam hlam).I = lam / 2 - 1 := by
  have hl : (64 : ℝ) ≤ lam := by exact_mod_cast hlam
  refine ⟨⟨?_, ?_, ?_, ?_⟩, ?_, ?_, ?_⟩ <;>
    simp [scopeInstance, ScopeParams, gamma] <;> linarith
\end{leancode}

\begin{remark}[Statut de v\'erification]
\label{rem:statut_verification_lean4}
Le noyau alg\'ebrique de cette preuve est certifi\'e m\'ecaniquement en
Lean4. Les formulations alternatives issues de la certification sont
pr\'esent\'ees \`a la sous-section~\ref{subsec:lean4_erreurs}.
\end{remark}
\end{sloppypar}

% -----------------------------------------------------------------------
\subsection{Extensions~: verdict d\'eterministe et forme raffin\'ee}
\label{subsec:lean4_extensions}
% -----------------------------------------------------------------------
\begin{sloppypar}
Quatre \'enonc\'es compl\`etent la formalisation (version~2.2 du fichier,
listing~\ref{lst:lean4_extensions})~; ils certifient les deux faits
structurels \'etablis au chapitre~\ref{chap:cadre_formel}.
\begin{itemize}
\item \texttt{opp\_eq\_one\_of\_deterministic} \'etablit que, si $I(V;\Sigma|C) = H(V|C) > 0$
      (verdict d\'eterministe), alors $\Opp = 1$
      (Remarque~\ref{rem:verdict_deterministe})~;
\item \texttt{deterministic\_forces\_capacity} montre que, dans ce cas, les
      Hypoth\`eses~1 et~2$'$ imposent $\gamma - I(M;\Sigma|C) \leq \Delta$~;
\item \texttt{trilemma\_refined} certifie la forme raffin\'ee~\eqref{eq:trilemme_cro_raffine},
      sous la condition $I(V;\Sigma|C) \leq H(V|C)$, toujours vraie~;
\item \texttt{refined\_bound\_is\_fano\_when\_delta\_ge\_gamma} \'etablit que, si
      $\Delta \geq \gamma$, le facteur institutionnel vaut~$1$ et la borne se
      r\'eduit \`a celle de Fano (Remarque~\ref{rem:activite_delta}).
\end{itemize}
Chacun de ces \'enonc\'es ne d\'epend que des axiomes standard \texttt{propext},
\texttt{Classical.choice} et \texttt{Quot.sound}.

\begin{leancode}[caption={Extensions de la version 2.2}, label={lst:lean4_extensions}]
/-- **T1.** Verdict déterministe : I(V;Σ|C) = H(V|C) ⇒ Opp = 1. -/
theorem opp_eq_one_of_deterministic {lam : ℕ} (P : ProtocolInstance lam)
    (hdet : P.IV = P.HV) (hHV : 0 < P.HV) : P.Opp = 1 := by
  rw [P.hOpp, hdet]; exact div_self hHV.ne'

/-- **T1 bis.** Verdict déterministe + Hyp. 1 et 2′ ⇒ γ − I(M;Σ|C) ≤ Δ :
    la contrainte institutionnelle ne peut pas être active. -/
theorem deterministic_forces_capacity {lam : ℕ} (P : ProtocolInstance lam)
    (hA : InstitutionalAssumptions P) (hdet : P.IV = P.HV) :
    gamma P.params - P.I ≤ P.params.Δ := by
  have := hA.A1; have := hA.A2; linarith

/-- **T2 (forme raffinée).** Avec I(V;Σ|C) ≤ H(V|C) (toujours vrai) :
    Conf·Rel·Opp ≤ ((κ/λ+1)/λ) · min 1 ((κ/λ+Δ)/γ). -/
theorem trilemma_refined {lam : ℕ} (P : ProtocolInstance lam)
    (hA : InstitutionalAssumptions P) (κ : ℝ) (hκ : 0 ≤ κ)
    (hlam : 1 ≤ lam) (hA4 : P.I ≤ κ / lam) (hIVHV : P.IV ≤ P.HV) :
    P.Conf * P.Rel * P.Opp ≤
      ((κ / lam + 1) / lam) * min 1 ((κ / lam + P.params.Δ) / gamma P.params) := by
  have hl0 : (0 : ℝ) < lam := by exact_mod_cast hlam
  have hγ0 : 0 < gamma P.params := by linarith [hA.A2γ]
  have hHV : 0 < P.HV := by linarith [hA.A2, hA.A2γ]
  have hκl : 0 ≤ κ / lam := div_nonneg hκ hl0.le
  have hC1 : P.Conf ≤ 1 := by rw [P.hConf]; linarith [P.hTV.1]
  have hC0 : 0 ≤ P.Conf := by rw [P.hConf]; linarith [P.hTV.2]
  have hR0 := P.hRel.1
  have hR : P.Rel ≤ (κ / lam + 1) / lam := by
    rw [le_div_iff₀ hl0]; linarith [P.fano]
  have hO0 : 0 ≤ P.Opp := by rw [P.hOpp]; exact div_nonneg P.hIV hHV.le
  have hO1 : P.Opp ≤ 1 := by
    rw [P.hOpp, div_le_one hHV]; exact hIVHV
  have hO2 : P.Opp ≤ (κ / lam + P.params.Δ) / gamma P.params := by
    rw [P.hOpp, div_le_div_iff₀ hHV hγ0]
    have e1 := mul_le_mul_of_nonneg_right hA.A1 hγ0.le
    have e2 := mul_le_mul_of_nonneg_left hA.A2
      (by linarith [P.params.hΔ_nonneg] : (0 : ℝ) ≤ κ / lam + P.params.Δ)
    have e3 : P.I * gamma P.params ≤ κ / lam * gamma P.params :=
      mul_le_mul_of_nonneg_right hA4 hγ0.le
    nlinarith
  have hOm : P.Opp ≤ min 1 ((κ / lam + P.params.Δ) / gamma P.params) := le_min hO1 hO2
  calc P.Conf * P.Rel * P.Opp ≤ 1 * ((κ / lam + 1) / lam) * min 1 ((κ / lam + P.params.Δ) / gamma P.params) :=
        mul_le_mul (mul_le_mul hC1 hR hR0 zero_le_one) hOm hO0
          (by rw [one_mul]; exact div_nonneg (by linarith) hl0.le)
    _ = _ := by ring

/-- **T2 bis.** Si Δ ≥ γ, la borne raffinée coïncide avec la seule borne de Fano. -/
theorem refined_bound_is_fano_when_delta_ge_gamma {lam : ℕ} (P : ProtocolInstance lam)
    (hA : InstitutionalAssumptions P) (κ : ℝ) (hκ : 0 ≤ κ) (hlam : 1 ≤ lam)
    (hΔγ : gamma P.params ≤ P.params.Δ) :
    min 1 ((κ / lam + P.params.Δ) / gamma P.params) = 1 := by
  have hl0 : (0 : ℝ) < lam := by exact_mod_cast hlam
  have hγ0 : 0 < gamma P.params := by linarith [hA.A2γ]
  have hκl : 0 ≤ κ / lam := div_nonneg hκ hl0.le
  apply min_eq_left
  rw [le_div_iff₀ hγ0]; linarith
\end{leancode}
\end{sloppypar}

% -----------------------------------------------------------------------
\subsection{Ce que la formalisation a r\'ev\'el\'e}
\label{subsec:lean4_erreurs}
% -----------------------------------------------------------------------
\begin{sloppypar}
La formalisation Lean4 du noyau alg\'ebrique du Trilemme CRO a produit,
au-del\`a de la certification, des formulations alternatives de plusieurs
\'etapes des preuves publi\'ees~\cite{MINKA2026CRO,MINKA2026PQCRO}. Ces
formulations, plus pr\'ecises que les versions initiales, \'eclairent la
structure logique du r\'esultat et d\'elimitent son domaine de validit\'e.
Elles sont document\'ees ci-dessous, car elles constituent une contribution
m\'ethodologique \`a part enti\`ere.

\paragraph{Sens de l'in\'egalit\'e du Lemme~4.}
Le Lemme~4 publi\'e formulait sur $\Conf$ la borne \emph{sup\'erieure}
$\Conf \leq 1 - \sqrt{c_0 \beta I / \lambda}$. La formalisation a conduit
\`a une formulation alternative, dans laquelle l'\'enonc\'e prend la forme
d'une borne \emph{inf\'erieure} $\Conf \geq 1 - \sqrt{c_0 \beta I}$. Cette
formulation pr\'ecise la port\'ee de l'argument de Pinsker, qui ne requiert
pas de normalisation par $\lambda$. Le
Lemme~\ref{lem:confidentiality_guarantee} retient cette formulation
alternative.

\paragraph{Statut de l'Hypoth\`ese~2.}
L'Hypoth\`ese~2 publi\'ee \'enon\c{c}ait $|\mathcal{Y}_\lambda| \leq
2^{\alpha H(C)}$ (borne sup\'erieure) et en d\'eduisait
$H(V|C) \geq \alpha H(C) - 1$ (borne inf\'erieure). La formalisation a
montr\'e que ces deux \'enonc\'es sont logiquement ind\'ependants. La borne
inf\'erieure exprime une propri\'et\'e d'\emph{efficacit\'e} du verdict, que
la borne sup\'erieure ne suffit pas \`a \'etablir. La formulation alternative
retenue s\'epare explicitement les deux \'enonc\'es
(Hypoth\`ese~\ref{ass:verdicts_efficace}).

\paragraph{R\'esolution de l'optimisation.}
L'\'Etape~6 de la preuve publi\'ee r\'esolvait $f'(x) = 0$ avec
$f'(x) = 2x - \frac{5}{2}ax^{3/2}$, en retenant $\sqrt{x^*} =
\frac{5a}{4}$. La formalisation a conduit \`a une r\'esolution
alternative, $\sqrt{x^*} = \frac{4}{5a}$. Elle a montr\'e que ce point
est un \emph{maximiseur} de $f$. Le r\'egime $I = O(1/\lambda)$ ne
s'obtient donc pas par cette optimisation. La formulation alternative
pose ce r\'egime comme hypoth\`ese explicite (Hypoth\`ese~\ref{ass:regime}).

\paragraph{N\'ecessit\'e du r\'egime~: un cas limite.}
La formalisation a mis en \'evidence un cas limite qui \'etablit la
n\'ecessit\'e \emph{logique} du r\'egime de confidentialit\'e
(Hypoth\`ese~\ref{ass:regime}). Sans ce r\'egime, la borne du Trilemme
cesse d'\^etre valide, pour une raison structurelle que le cas ci-dessous
rend explicite.

Consid\'erons le protocole suivant, en encodage classique diagonal
($\Sigma$ prend ses valeurs dans un alphabet fini $\mathcal{S}$). Le
protocole r\'ev\`ele $M$ avec probabilit\'e $1/2$ et produit sinon une
preuve uniform\'ement al\'eatoire dans $\{0,1\}^\lambda$, ce qui s'\'ecrit
\begin{equation}
\Sigma =
\begin{cases}
M, & \text{avec probabilit\'e } 1/2, \\
U \xleftarrow{\$} \{0,1\}^\lambda, & \text{avec probabilit\'e } 1/2.
\end{cases}
\label{eq:counterexample_protocol}
\end{equation}
Soit $V$ le premier bit de $\Sigma$. Nous supposons de plus $M \perp C$.

Pour ce protocole, les trois m\'etriques valent
\begin{align}
\Conf &= 1/2, \\
\Rel &= 1/2 + 2^{-\lambda-1}, \\
\Opp &= 1,
\end{align}
et $I(M;\Sigma|C) = \lambda/2 - 1 + O(\lambda 2^{-\lambda})$. Le
protocole est $\beta$-r\'egulier avec $\beta \leq 1{,}04$ pour tout
$\lambda \geq 64$. Il satisfait donc les
Hypoth\`eses~\ref{ass:extraction_bornee}--\ref{ass:regularite}. Il
\emph{s'\'ecarte} cependant de l'Hypoth\`ese~\ref{ass:regime}, car
$I(M;\Sigma|C) = \Theta(\lambda)$ alors que le r\'egime exige $O(1/\lambda)$.

Pour ce protocole, le produit CRO vaut
\begin{equation}
\Conf \cdot \Rel \cdot \Opp \;\approx\; \frac{1}{2} \cdot \frac{1}{2} \cdot 1 \;=\; \frac{1}{4},
\end{equation}
ind\'ependant de $\lambda$. La borne du Trilemme
$\Conf \cdot \Rel \cdot \Opp \leq \Delta/\lambda + O(\lambda^{-2})$ ne
tient donc pas uniform\'ement pour $I = \Theta(\lambda)$, puisqu'elle
requiert le r\'egime de confidentialit\'e.

\textbf{Cons\'equence.} Ce cas limite \'etablit que
l'Hypoth\`ese~\ref{ass:regime} est une \emph{condition n\'ecessaire} de la
borne. Il justifie r\'etrospectivement la formulation alternative du
Th\'eor\`eme~\ref{thm:trilemme_cro_conditionnel}, dans laquelle cette
hypoth\`ese est explicite. Il a aussi motiv\'e la distinction formelle, dans
l'annexe~\ref{chap:annexe_lean4}, entre le \textbf{t\'emoin de coh\'erence}
(qui atteint la borne) et le \textbf{t\'emoin de port\'ee} (qui s'\'ecarte du
r\'egime et montre que la borne ne tient pas sans lui).

\paragraph{Propagation \`a l'extension quantique.}
Les formulations alternatives de la preuve classique ont conduit \`a
r\'eexaminer l'extension post-quantique du Trilemme~\cite{MINKA2026PQCRO},
qui h\'eritait de la m\^eme structure. Le Lemme~6 et la borne sup\'erieure
du Lemme~7 ont fait l'objet de formulations plus pr\'ecises. Pour le
Lemme~6, la formalisation a mis en \'evidence un cas limite form\'e de deux
\'etats orthogonaux. Les Hypoth\`eses~4 et~5 sont d\'esormais pos\'ees
explicitement au lieu d'\^etre d\'eduites. Les formulations correspondantes
sont document\'ees aux
sections~\ref{subsec:enonce_preuve_pq} et~\ref{subsec:hypotheses_quantiques}.

\paragraph{Comment la formalisation a produit ces formulations.}
Ces formulations alternatives sont issues de la formalisation. La relecture
des articles ne les avait pas fait appara\^itre. Elles sont n\'ees \`a deux
moments pr\'ecis. D'abord, l'\'enonc\'e publi\'e du Lemme~4 n'a pas pu \^etre
d\'emontr\'e \`a partir des hypoth\`eses disponibles, ce qui a conduit \`a
r\'eexaminer la preuve et \`a pr\'eciser le sens de l'in\'egalit\'e. Un \'echec
de preuve formelle ne tranche pas \`a lui seul~; l'examen qui l'a suivi a
\'etabli la formulation retenue. Ensuite, la formalisation du noyau
alg\'ebrique a montr\'e que l'optimisation publi\'ee ne pouvait pas \^etre
reproduite. Le cas limite ci-dessus a \'emerg\'e de l'effort pour pr\'eserver
la coh\'erence de la preuve formelle. Il est alors devenu n\'ecessaire
d'inscrire explicitement le r\'egime de confidentialit\'e dans l'\'enonc\'e du
th\'eor\`eme.

\paragraph{Statut des formulations alternatives.}
Les formulations alternatives des preuves classique et post-quantique sont
celles pr\'esent\'ees aux sections~\ref{sec:cro_classique}
et~\ref{sec:cro_pq} de cette th\`ese. Les versions publi\'ees initiales
restent accessibles pour r\'ef\'erence~\cite{MINKA2026CRO,MINKA2026PQCRO}.

\paragraph{La le\c{c}on m\'ethodologique.}
Les formulations alternatives et le cas limite qui \'etablit la n\'ecessit\'e
du r\'egime de confidentialit\'e conduisent \`a une le\c{c}on
m\'ethodologique sur le r\^ole de la formalisation dans le processus de
publication scientifique, d\'evelopp\'ee \`a la
section~\ref{subsec:verification_corrige_lecon}.
\end{sloppypar}

\begin{boiteArbitre}
L'arbitre referme le dossier. La machine a produit des formulations
alternatives que les auteurs ont retenues. La le\c{c}on qu'il en tire
d\'epasse la substitution d'une formule \`a une autre, car ni la relecture
humaine ni la machine ne suffisent seules. La machine v\'erifie les
d\'eductions tandis que la relecture v\'erifie que les d\'efinitions disent ce
qu'on croit. L'arbitre comprend que la formalisation est une n\'ecessit\'e
\'epist\'emologique pour tout r\'esultat interdisciplinaire dont aucune des
disciplines constituantes ne contr\^ole l'int\'egralit\'e.
\end{boiteArbitre}


% -----------------------------------------------------------------------
\subsection{Application aux familles de protocoles canoniques}
\label{subsec:application_familles_canoniques}
% -----------------------------------------------------------------------
\begin{sloppypar}
\subsubsection{R\'egimes de capacit\'e institutionnelle}
\label{subsubsec:regimes_capacite_lean4}

La formalisation contient aussi une classification indicative des
capacit\'es institutionnelles par le rapport $\Delta/\lambda$. Elle est distincte
des r\'egimes du th\'eor\`eme (seuil $H^\sharp$, Remarque~\ref{rem:activite_delta})
et des r\'egimes de construction SOTP (section~\ref{sec:sotp_construction})~; elle
sert seulement \`a illustrer l'usage de d\'efinitions calculables dans Lean. Les
param\`etres des familles canoniques cod\'ees en Lean sont ceux de la
calibration illustrative, \`a l'exception de $H(C)$, qui n'intervient pas dans
cette classification.

\begin{leancode}[caption={Classification des r\'egimes institutionnels en Lean4}, label={lst:lean4_capacity_regime}]
inductive InstitutionalRegime where
  | LowCapacity    -- Δ < 0.1 × λ : capacite faible
  | MidCapacity    -- 0.1 × λ ≤ Δ ≤ 0.5 × λ : capacite moyenne
  | HighCapacity   -- Δ > 0.5 × λ : capacite elevee

noncomputable def capacity_regime (Δ lam : ℝ) : InstitutionalRegime :=
  if Δ < 0.1 * lam then .LowCapacity
  else if Δ ≤ 0.5 * lam then .MidCapacity
  else .HighCapacity
\end{leancode}

\begin{theorem}[R\'egime Helios~\cite{Adida08}]\label{thm:regime_helios}
\`A $\lambda = 128$, Helios ($\Delta = 56$) est
dans le r\'egime de capacit\'e moyenne, car $56 \in [12{,}8,\, 64]$.
\end{theorem}

\begin{leancode}[caption={V\'erification du r\'egime Helios en Lean4}, label={lst:lean4_helios_regime}]
theorem Helios_regime_at_128 :
    capacity_regime HeliosParams.Δ 128 = .MidCapacity := by
  show capacity_regime 56 128 = _
  unfold capacity_regime
  norm_num
\end{leancode}

\begin{theorem}[R\'egime zk-SNARKs~\cite{Groth16,BCTV14}]\label{thm:regime_snarks}
\`A $\lambda = 256$, zk-SNARKs ($\Delta = 73$) est
dans le r\'egime de capacit\'e moyenne, car $73 \in [25{,}6,\, 128]$.
\end{theorem}

\begin{leancode}[caption={V\'erification du r\'egime zk-SNARKs en Lean4}, label={lst:lean4_zksnark_regime}]
theorem zkSNARK_regime_at_256 :
    capacity_regime zkSNARKParams.Δ 256 = .MidCapacity := by
  show capacity_regime 73 256 = _
  unfold capacity_regime
  norm_num
\end{leancode}

\begin{theorem}[R\'egime audit traditionnel~\cite{Mason07,Kerr04}]\label{thm:regime_audit}
\`A $\lambda = 128$, l'audit traditionnel ($\Delta = 35$) est dans le r\'egime de capacit\'e moyenne, car
$35 \in [12{,}8,\, 64]$.
\end{theorem}

\begin{leancode}[caption={V\'erification du r\'egime audit traditionnel en Lean4}, label={lst:lean4_audit_regime}]
theorem TraditionalAudit_regime_at_128 :
    capacity_regime TraditionalAuditParams.Δ 128 = .MidCapacity := by
  show capacity_regime 35 128 = _
  unfold capacity_regime
  norm_num
\end{leancode}

Ces v\'erifications sont enti\`erement automatis\'ees dans Lean4. Elles
v\'erifient des calculs sur des valeurs illustratives, sans valider le mod\`ele
de capacit\'e institutionnelle.
\end{sloppypar}

% -----------------------------------------------------------------------
% Bo\^ite finale
% -----------------------------------------------------------------------
\begin{boiteArbitre}
L'arbitre referme le dossier Lean4. Il a achev\'e l'\'etape \textbf{C}.
Cette section lui a appris que la preuve qu'il croyait compl\`ete gagnait
\`a \^etre pr\'ecis\'ee en plusieurs points.

Sept contributions structurent ce qu'il retient. Il y a d'abord le mod\`ele formel
(familles de protocoles index\'ees par $\lambda$, param\`etres institutionnels
fixes), puis l'int\'egration des propri\'et\'es informationnelles comme champs de
structure (Fano, Pinsker) et le noyau alg\'ebrique certifi\'e sans aucune
admission tacite. S'y ajoutent la preuve de n\'egligeabilit\'e, l\`a o\`u l'expos\'e
informel \'ecrivait simplement \og n\'egligeable\fg, le r\'egime de confidentialit\'e
rendu explicite comme condition du th\'eor\`eme ainsi que les constantes explicites
($\kappa(\Delta+1)$ et $\kappa^2$ au lieu d'un $O(\lambda^{-2})$ non chiffr\'e). La
forme raffin\'ee compl\`ete l'ensemble en lui indiquant exactement quand sa capacit\'e
d'extraction devient la contrainte liante.

L'encha\^inement des faits le marque particuli\`erement. Trois formulations
initiales (le sens d'une in\'egalit\'e, une borne inf\'erieure tir\'ee d'une borne
sup\'erieure, la r\'esolution d'une optimisation) ont \'et\'e pr\'ecis\'ees, ce que
la relecture des articles n'avait pas permis. La formalisation a oblig\'e \`a
justifier chaque \'etape. Le cas limite qui \'etablit la n\'ecessit\'e du
r\'egime est d'ailleurs n\'e de cette exigence.

La borne $\Delta/(\gamma\lambda) + O(\lambda^{-2})$ est maintenant prouv\'ee dans le
cadre du r\'egime, avec des termes correctifs chiffr\'es. Cette section en tire sa
port\'ee \'epist\'emologique, puisqu'elle montre que la rigueur se v\'erifie au lieu
de se supposer.
\end{boiteArbitre}
% ============================================================================
% CHAPITRE 3 : CONCLUSION
% ============================================================================

\section{CONCLUSION DU CHAPITRE~\ref{chap:constructions}}
\label{sec:conclusion_constructions}

\begin{sloppypar}
La famille SOTP (section~\ref{sec:sotp_construction}) propose trois
constructions de r\'ef\'erence, SOTP-Low, SOTP-Base et SOTP-High, adapt\'ees
chacune \`a un r\'egime de richesse contextuelle. Leurs propri\'et\'es
cryptographiques, \`a savoir la confidentialit\'e (calculatoire pour les
variantes \`a PRG, divulgation nulle de connaissance pour SOTP-High) et la
fiabilit\'e de v\'erification, sont prouv\'ees sous PRG, DLog et q-SDH/q-PKE. Les facteurs
$0{,}75$, $0{,}5$ et $0{,}45$ des formulations initiales sont conserv\'es comme
heuristiques de conception, tandis que l'optimalit\'e universelle
est formul\'ee comme conjecture (Conjecture~\ref{conj:optimalite_universelle}). Cette
r\'eserve a une raison structurelle, car une couche Clay qui ne d\'epend que du contexte public
n'apporte aucune information conditionnelle au sens de $\Opp$
(Remarque~\ref{rem:statut_sotp}).

La v\'erification Lean4 (section~\ref{sec:lean4}) a certifi\'e
le noyau alg\'ebrique. Les in\'egalit\'es de Pinsker et de Fano y sont
int\'egr\'ees comme propri\'et\'es des instances de protocole et les hypoth\`eses
institutionnelles y sont explicit\'ees~; la borne
$\Delta/(\gamma\lambda) + O(\lambda^{-2})$ et sa forme raffin\'ee sont v\'erifi\'ees
sans \texttt{sorry}. Deux t\'emoins formalisent la coh\'erence (une instance du
mod\`ele abstrait atteint la borne rel\^ach\'ee) et la port\'ee (une instance hors
r\'egime atteint $1/4$). Cette d\'emarche illustre la valeur
\'epist\'emique de la v\'erification formelle, qui certifie des \'enonc\'es et contraint aussi
\`a expliciter ce qui \'etait rest\'e implicite. Elle a conduit \`a des formulations
plus pr\'ecises de trois \'enonc\'es publi\'es (sous-section~\ref{subsec:lean4_erreurs}).

Quatre limitations demeurent. Les param\`etres
$\Delta$, $\alpha$, $\beta$ reposent sur des estimations analytiques. Leur mesure directe aupr\`es
d'acteurs institutionnels r\'eels rel\`eve des protocoles du
chapitre~\ref{chap:experimentation}. Les m\'etriques CLO
des syst\`emes de r\'ef\'erence ont \'et\'e fix\'ees par calibration param\'etrique
r\'etrospective, sans exp\'erimentation directe et ne valident pas les
heuristiques SOTP. La formalisation Lean4 couvre le noyau
alg\'ebrique. Les fondements probabilistes et les r\'eductions de s\'ecurit\'e
des constructions SOTP restent hors de son p\'erim\`etre. Le mod\`ele quantique utilise
principalement le canal d\'epolarisant, qui peut ne pas capturer toutes les
nuances des syst\`emes quantiques r\'eels.

La cryptographie institutionnelle dispose d\'esormais des m\'etriques, de la
loi d'impossibilit\'e, des constructions de r\'ef\'erence et de la certification
formelle que requiert tout sous-domaine th\'eorique. Ces r\'esultats doivent
encore \^etre confront\'es aux syst\`emes r\'eels, ce qui conduit \`a se demander
si les param\`etres institutionnels qui gouvernent ces bornes sont mesurables.
Cette question correspond \`a l'\'etape \textbf{X} de la progression et fait
l'objet du chapitre~\ref{chap:experimentation}.
\end{sloppypar}


% ==========================================================
% CHAPITRE 4 : PROTOCOLES DE MESURE ET VALIDATION RETROSPECTIVE
% ==========================================================
\chapter[PROTOCOLES DE MESURE ET VALIDATION R\'ETROSPECTIVE]%
        {PROTOCOLES DE MESURE ET VALIDATION R\'ETROSPECTIVE}
\label{chap:experimentation}
\debutChapitre%
  {Protocoles de mesure et validation r\'etrospective}%
  {%
    Les bornes \'etablies aux chapitres pr\'ec\'edents sont conditionnelles,
    de sorte que leur port\'ee pratique d\'epend enti\`erement de la mesurabilit\'e des
    param\`etres institutionnels. Ce chapitre r\'ealise l'acte exp\'erimenter du
    cycle canonique. Il sp\'ecifie d'abord le dispositif exp\'erimental complet
    permettant de mesurer ces param\`etres, sous la forme de six protocoles
    op\'erationnels portant sur la capacit\'e d'extraction institutionnelle,
    l'entropie contextuelle, l'opposabilit\'e sur syst\`emes d\'eploy\'es,
    l'estimation conjointe de la capacit\'e et de la granularit\'e, la capacit\'e
    d'extraction quantique et le coefficient de perte d'interpr\'etabilit\'e quantique. Il
    d\'erive ensuite les valeurs d'opposabilit\'e de trois syst\`emes
    d\'eploy\'es et les met en regard d'une d\'ecomposition des param\`etres
    institutionnels. Cette calibration illustrative v\'erifie la coh\'erence des ordres de
    grandeur, sans constituer une validation empirique, puisque ces
    syst\`emes ne sont pas des instances du th\'eor\`eme et que le programme
    exp\'erimental d\'ecrit n'a pas \'et\'e d\'eploy\'e dans le cadre de cette th\`ese.%
  }%
  {%
    protocole exp\'erimental, mesure de $\Delta$, entropie contextuelle,
    validation r\'etrospective, coh\'erence interne, reproductibilit\'e%
  }%

% ============================================================================
% CHAPITRE 3 — SECTION 3.1 : INTRODUCTION
% Fichier : partie3/introduction_validation.tex
% ============================================================================

\section{INTRODUCTION}
\label{sec:introduction_experimentation}
% ============================================================================
\begin{sloppypar}

Les Trilemmes CRO classique et post-quantique \'etablis au
chapitre~\ref{chap:cadre_formel} sont des r\'esultats conditionnels, car leur
port\'ee pratique d\'epend enti\`erement de la possibilit\'e de mesurer
empiriquement les param\`etres institutionnels qui les param\`etrent. Pour le
cadre classique (Contribution~C2), il s'agit de $\Delta$ (capacit\'e
d'extraction), $\alpha$ (granularit\'e des verdicts), $\beta$ (r\'egularit\'e)
et $H(\mathcal{C})$ (entropie contextuelle). Pour l'extension post-quantique
(Contribution~C3), s'y ajoutent $\Delta_q$ (capacit\'e quantique) et $\eta_q$
(coefficient de perte d'interpr\'etabilit\'e quantique). Sans proc\'edure de
mesure, ces grandeurs demeurent des postulats~; sans estimation, les bornes
demeurent muettes sur les syst\`emes r\'eels.

\subsection{Statut \'epist\'emique des deux volets}
\label{subsec:statut_epistemique}

Ce chapitre r\'epond \`a cette double exigence en deux temps de nature
distincte, qu'il importe de ne pas confondre.

\textbf{Premi\`erement}, les
sections~\ref{sec:protocoles_experimentaux}
\`a~\ref{subsec:calendrier_ressources} sp\'ecifient, jusqu'au niveau
op\'erationnel, le dispositif exp\'erimental permettant de mesurer chacun des
param\`etres, qu'il s'agisse de la stratification des \'echantillons, des estimateurs, des crit\`eres de
fiabilit\'e inter-\'evaluateurs ou des seuils de falsification. Ce programme n'a pas
\'et\'e d\'eploy\'e dans le cadre de cette th\`ese~; son statut exact est
explicit\'e \`a la sous-section~\ref{subsubsec:statut_deploiement}, qui referme
cette partie.

\textbf{Deuxi\`emement}, les
sections~\ref{sec:derivation_opposabilite}
\`a~\ref{sec:reproductibilite} proc\`edent \`a une analyse
\emph{r\'etrospective et analytique}. Elle d\'erive d'abord, pour trois
syst\`emes cryptographiques r\'eellement d\'eploy\'es, les valeurs
d'opposabilit\'e et les produits CRO qui servent de r\'ef\'erence \`a tout le
chapitre. Elle estime ensuite, par une voie enti\`erement distincte (triangulation de la litt\'erature cognitive, de l'analyse documentaire et des
sp\'ecifications techniques), les param\`etres institutionnels qui
param\`etrent les bornes avant de mettre les deux estimations en regard. Ces deux voies
ne sont pas ind\'ependantes, puisque les valeurs de $\Delta$ ont d'abord \'et\'e fix\'ees
par inversion de la borne. Les syst\`emes d\'eploy\'es ne sont pas non plus des
instances du th\'eor\`eme (Remarque~\ref{rem:portee_cas_etude}). Cette partie
constitue donc une \emph{calibration illustrative}. Elle v\'erifie que les ordres
de grandeur retenus sont mutuellement coh\'erents, sans constituer une validation
empirique.

Le chapitre sp\'ecifie d'abord les protocoles de mesure puis exploite les donn\'ees existantes.
Cet ordre m\'ethodologique r\'epond \`a une n\'ecessit\'e logique, car on ne peut extraire
d'information de mesures que si l'on sait d\'efinir ce qui doit \^etre mesur\'e.
La sp\'ecification des protocoles pr\'ec\`ede donc l'application aux
syst\`emes d\'eploy\'es. Les estimations r\'etrospectives restent ainsi ancr\'ees
dans une proc\'edure explicite et falsifiable.

\subsection{Param\`etres \`a mesurer et statut de leur validation}
\label{subsec:parametres_mesurer}

Le tableau~\ref{tab:parametres_mesure} r\'ecapitule, pour chaque contribution,
les param\`etres concern\'es, le protocole de mesure qui leur est d\'edi\'e et
l'\'etat de leur validation.

\begin{table}[htbp]
\centering
\caption{Param\`etres institutionnels~: protocole de mesure et statut de validation}
\label{tab:parametres_mesure}
\small
\begin{tabular}{@{}lllll@{}}
\toprule
\textbf{Contribution} & \textbf{Param\`etres} & \textbf{Protocole de mesure} &
\textbf{Statut de validation} & \textbf{Type} \\
\midrule
C2 (Trilemme classique)
  & $\Delta$, $\alpha$, $\beta$, $H(\mathcal{C})$
  & \S\ref{subsec:protocole_delta}--\S\ref{subsec:protocole_conjoint}
  & Calibration illustrative & R\'etrospectif \\
C3 (Trilemme post-quantique)
  & $\Delta_q$, $\eta_q$, $\beta_q$
  & \S\ref{subsec:protocole_delta_q}--\S\ref{subsec:protocole_combine}
  & Calibration illustrative & R\'etrospectif \\
C4 (Famille SOTP)
  & Opposabilit\'e r\'ealis\'ee
  & \S\ref{subsec:protocole_opp}
  & Calibration illustrative & R\'etrospectif \\
\bottomrule
\end{tabular}
\end{table}

La section~\ref{subsec:protocole_delta} d\'etaille la mesure de $\Delta$,
la section~\ref{subsec:protocole_HC} celle de $H(\mathcal{C})$ tandis que la
section~\ref{subsec:protocole_opp} porte sur la mesure de l'opposabilit\'e sur
syst\`emes d\'eploy\'es. Les protocoles post-quantiques
($\Delta_q$ en section~\ref{subsec:protocole_delta_q}, $\eta_q$ en
section~\ref{subsec:protocole_eta_q}) compl\`etent le dispositif, avant que
la section~\ref{subsec:protocole_combine} n'en articule la validation
conjointe.

La section~\ref{sec:derivation_opposabilite} \'etablit ensuite les
opposabilit\'es des trois syst\`emes classiques. La
section~\ref{sec:validation_retrospective} met ces d\'erivations en regard
des param\`etres institutionnels estim\'es par ailleurs. Cette mise en regard
v\'erifie la coh\'erence des ordres de grandeur. Le test du cadre rel\`eve du
programme exp\'erimental (crit\`eres de r\'efutation du
tableau~\ref{tab:criteres_refutation}).

\subsection{Table unique des param\`etres de r\'ef\'erence}
\label{subsec:parametres_reference}

Le tableau~\ref{tab:parametres_reference} fixe la seule source de r\'ef\'erence
pour $H(\mathcal{C})$, $\alpha$, $\Delta$ et $\Delta_q$. Toute valeur employ\'ee
ailleurs y renvoie. Il est reproduit dans le fichier
\texttt{06\_calcul/data/parametres\_reference.csv} du d\'ep\^ot
compagnon~\cite{CROTrilemmaDepot}.

\begin{table}[htbp]
\centering
\caption{Param\`etres institutionnels de r\'ef\'erence (valeurs illustratives)}
\label{tab:parametres_reference}
\small
\begin{tabular}{@{}lccccc@{}}
\toprule
\textbf{Syst\`eme} & $H(\mathcal{C})$ domaine & $H(\mathcal{C})$ esp\`ece & $\alpha$ & $\Delta$ & $\Delta_q$ / $\eta_q$ \\
\midrule
ECDSA (contrat) & 20 & 20 & 1{,}0 & 35 & --- \\
Helios (scrutin) & 40 & 8 & 1{,}2 & 56 & --- \\
zk-SNARKs (AML) & 70 & 11 & 1{,}5 & 73 & --- \\
ML-DSA-65 (contrat) & 20 & 10{,}5 & 1{,}0 & --- & 28 (hyp.) / $0$ \\
\bottomrule
\end{tabular}

\smallskip
\begin{minipage}{0.92\linewidth}\footnotesize
Entropies en bits. \og Domaine\fg{} d\'esigne la richesse du contexte institutionnel
g\'en\'eral (chapitre~\ref{chap:cadre_formel})~; \og esp\`ece\fg{} d\'esigne le contexte d'une
d\'ecision particuli\`ere (section~\ref{sec:derivation_opposabilite}). Pour
ML-DSA-65, $\eta_q = 0$ car la signature est une cha\^{\i}ne de bits classique
(les valeurs $\eta_q > 0$ du chapitre d\'ecrivent un sc\'enario hypoth\'etique)~;
$\Delta_q = 28$~bits est une hypoth\`ese de travail
(section~\ref{sec:etude_cas_dilithium}), \`a estimer par le protocole de la
section~\ref{subsec:protocole_delta_q}.
\end{minipage}
\end{table}

\end{sloppypar}
% ============================================================================
% CHAPITRE 3 — SECTION 3.2 : PROTOCOLES EXPÉRIMENTAUX
% Fichier : partie3/protocole_experimental.tex
% ============================================================================

\section{PROTOCOLES EXP\'ERIMENTAUX}
\label{sec:protocoles_experimentaux}
% ============================================================================
\begin{sloppypar}
Cette section sp\'ecifie les six protocoles de mesure destin\'es \`a estimer
empiriquement les param\`etres du cadre CRO, \`a savoir la capacit\'e
d'extraction institutionnelle, l'entropie contextuelle, l'opposabilit\'e,
ainsi que leurs extensions quantiques, en d\'etaillant pour chacun le design,
la proc\'edure, les m\'etriques et les pr\'edictions falsifiables, avant de
consolider le calendrier, les ressources et les crit\`eres de validation de
l'ensemble du programme.
\end{sloppypar}

\begin{boiteArbitre}
L'arbitre ouvre un dossier qui ne ressemble pas aux pr\'ec\'edents, puisqu'il
y trouve des protocoles sans r\'esultats, puisque rien n'a encore \'et\'e mesur\'e.
Il lit la mani\`ere dont on compte mesurer, qui est pr\'ecis\'ement ce qu'il
doit juger.
\end{boiteArbitre}

% -----------------------------------------------------------------------
\subsection{Protocole de mesure de $\Delta$~:
          capacit\'e d'extraction institutionnelle}
\label{subsec:protocole_delta}
% -----------------------------------------------------------------------
\begin{sloppypar}

\subsubsection{Objectif et d\'efinition op\'erationnelle}
\label{subsubsec:objectif_delta}

\noindent\textbf{Objectif.} Mesurer empiriquement la capacit\'e d'extraction
institutionnelle $\Delta$ (en bits), d\'efinie par l'hypoth\`ese d'extraction
born\'ee~\cite{MINKA2026CRO}~:
\begin{equation}
I(V;\Sigma|C) \leq I(M;\Sigma|C) + \Delta
\label{eq:extraction_bornee_op}
\end{equation}

\begin{definition}[Capacit\'e d'extraction institutionnelle $\Delta$]\label{def:capacite_extraction_op}
$\Delta$ est la quantit\'e d'information sur les verdicts qu'un acteur
institutionnel peut extraire \emph{au-del\`a} de l'information mutuelle brute
contenue dans la preuve cryptographique. En pratique, $\Delta$ capture la
connaissance pr\'ealable (pr\'ec\'edents l\'egaux, formation sp\'ecialis\'ee,
proc\'edures institutionnelles, capacit\'e de d\'elib\'eration collective) non
formalis\'ee dans le contexte $C$.
\end{definition}

\noindent\textbf{Estimateur.} Par la
Proposition~\ref{prop:fondement_extraction}, la quantit\'e contr\^ol\'ee par
l'Hypoth\`ese~\ref{ass:extraction_bornee} est l'information conditionnelle
$I(V;\Sigma|M,C)$. L'estimateur retenu est donc
\begin{equation}
\hat{\Delta} = \hat{I}(V;\Sigma \mid M, C)
\label{eq:estimateur_delta}
\end{equation}
o\`u $\hat{I}$ est un estimateur de l'information mutuelle conditionnelle
(plug-in avec correction de Miller--Madow, ou estimateur de type NSB pour
les petits \'echantillons), calcul\'e sur les triplets $(v_i, \sigma_i, m_i)$
recueillis \`a contexte $c$ fix\'e. Le message $m_i$ est connu de
l'exp\'erimentateur (v\'erit\'e terrain) mais non des participants. $\hat{\Delta}$ estime
la plus petite capacit\'e compatible avec l'Hypoth\`ese~\ref{ass:extraction_bornee}
sous la condition~\eqref{eq:extraction_bornee_forte}. L'estimateur est biais\'e
positivement en petit \'echantillon~; l'intervalle de confiance est obtenu par
bootstrap non param\'etrique \emph{sur les participants}.

\begin{remark}[Formulation initiale de l'estimateur]
La formulation initiale de ce protocole proposait
$\hat{\Delta} = I_{\mathrm{emp}}(V;\Sigma|C) - I_{\mathrm{vrai}}(M;\Sigma|C)$.
Cet estimateur prend des valeurs n\'egatives pour un protocole fiable
($I(M;\Sigma|C) \geq \Relrec\cdot\lambda - 1$ par Fano, alors que
$I(V;\Sigma|C) \leq \log_2|\mathcal{V}_\lambda|$) et ne fournit qu'une borne
inf\'erieure sur $\Delta$. La formulation retenue
est~\eqref{eq:estimateur_delta}.
\end{remark}

\subsubsection{Design exp\'erimental}
\label{subsubsec:design_delta}

\noindent\textbf{Participants.} $N \geq 100$ professionnels des institutions
cibles, stratifi\'es par niveau d'expertise et par type
d'institution~\cite{MINKA2026CRO} selon la r\'epartition du
tableau~\ref{tab:stratification_delta}.

\begin{table}[htbp]
\centering
\caption{Stratification de l'\'echantillon pour la mesure de $\Delta$}
\label{tab:stratification_delta}
\begin{tabular}{@{}llcc@{}}
\toprule
\textbf{Cat\'egorie} & \textbf{Sous-groupe} & \textbf{$n$} & \textbf{$\Delta$ pr\'edit} \\
\midrule
Juges & G\'en\'eralistes (juridictions civiles) & 20 & $\approx 25$ bits \\
      & Sp\'ecialis\'es (contrats, PKI, num\'erique) & 20 & $\approx 40$ bits \\
\midrule
R\'egulateurs & Juniors (< 5 ans d'exp\'erience) & 20 & $\approx 35$ bits \\
financiers  & Seniors ($\geq 5$ ans, superviseurs) & 20 & $\approx 65$ bits \\
\midrule
Auditeurs & Cour des comptes (secteur public) & 20 & $\approx 45$ bits \\
          & Cabinets priv\'es (Big 4 et \'equivalents) & 20 & $\approx 55$ bits \\
\midrule
\textbf{Total} & & \textbf{120} & \\
\bottomrule
\end{tabular}
\end{table}

\noindent\textbf{Mat\'eriel exp\'erimental.}
\begin{itemize}
\item 5 types de preuves cryptographiques (ECDSA, Helios avec ElGamal et ZK,
      Groth16 zk-SNARK, RSA-PSS et ElGamal classique)
\item 10 contextes standardis\'es par type de preuve (contractuels,
      \'electoraux, financiers), soit 50 paires contexte--preuve
\item Questionnaires valid\'es de 30 items par contexte
\end{itemize}

\subsubsection{Proc\'edure}
\label{subsubsec:procedure_delta}

La session exp\'erimentale pour chaque participant se d\'eroule en cinq
phases.

\noindent\textbf{Phase 0. Formation initiale (2 h).}
Introduction aux concepts cryptographiques de base sans r\'ev\'eler les
m\'ecanismes internes, afin d'homog\'en\'eiser le vocabulaire minimum.

\noindent\textbf{Phase 1. Session exp\'erimentale principale (4 h).}
Pour chacun des 10 contextes assign\'es (rotation \'equilibr\'ee)~:
\begin{enumerate}
\item[(a)] Pr\'esentation du contexte $c \in \mathcal{C}$
\item[(b)] Pr\'esentation de la preuve cryptographique $\sigma$
\item[(c)] Compl\'etion du questionnaire d'extraction de verdict
\item[(d)] Mesure du temps de r\'eponse et notation de confiance subjective
\end{enumerate}

\noindent\textbf{Phase 2. V\'erification test-retest (1 h).}
Re-pr\'esentation de 3 contextes d\'ej\`a vus apr\`es un d\'elai de 2 semaines.

\noindent\textbf{Phase 3. Calcul du gold standard.}
Deux cryptographes experts ind\'ependants rendent les verdicts \og corrects\fg{}
(accord inter-juge requis $\geq 90\%$).

\noindent\textbf{Phase 4. D\'ebriefing et restitution (30 min).}

\subsubsection{M\'etriques et estimateurs}
\label{subsubsec:metriques_delta}

Le protocole standardis\'e de~\cite{MINKA2026CRO} calcule $\hat\Delta$ selon
l'Algorithme~\ref{alg:estimation_delta}.

\begin{algorithm}[H]
\caption{Estimation empirique de $\Delta$}
\label{alg:estimation_delta}
\begin{algorithmic}[1]
\State \textbf{Entr\'ees~:} Institution $I$, types de preuves
       $\{\sigma_1, \ldots, \sigma_k\}$, contexte $C$
\State \textbf{Sortie~:} $\hat\Delta$, $\hat\alpha$
\State Pr\'esenter $\sigma_i$ avec $C$ aux acteurs institutionnels
\State Enregistrer verdicts $V$, scores de confiance, temps de d\'ecision
\State Associer \`a chaque r\'eponse le message $m$ connu de l'exp\'erimentateur
\State Tester l'admissibilit\'e $V \perp M \mid (\Sigma, C)$ (test de permutation)
\State Estimer $\hat\Delta = \hat I(V;\Sigma \mid M, C)$ (\'equation~\eqref{eq:estimateur_delta})
\State Estimer $\hat\alpha = \log_2|\mathcal{V}| / H(C)$
\State Valider par fiabilit\'e test-retest et coh\'erence inter-\'evaluateurs
\end{algorithmic}
\end{algorithm}

L'algorithme transpose sur le terrain la d\'efinition th\'eorique de la
capacit\'e d'extraction institutionnelle en encha\^inant quatre phases.

La phase d'\textbf{entr\'ees} (lignes~1--2) fixe le cadre de la mesure.
Elle re\c{c}oit une institution cible $I$, un ensemble de types de preuves
cryptographiques $\{\sigma_1, \ldots, \sigma_k\}$ et un contexte $C$, puis
d\'eclare comme sorties l'estimateur $\hat\Delta$ de la capacit\'e
d'extraction et l'estimateur $\hat\alpha$ de la granularit\'e des verdicts.
Aucun calcul n'intervient \`a ce stade, qui se borne \`a nommer les grandeurs
mesur\'ees. S\'eparer entr\'ees et sorties importe n\'eanmoins, car $\Delta$
n'est jamais observ\'e directement et se d\'eduit d'un contraste entre
deux quantit\'es informationnelles.

La phase de \textbf{collecte} (lignes~3--4) produit les donn\'ees brutes.
Elle pr\'esente chaque preuve $\sigma_i$ accompagn\'ee de son contexte $C$
aux acteurs institutionnels, puis enregistre les verdicts $V$ rendus, les
scores de confiance subjectifs et les temps de d\'ecision, trois s\'eries qui
doivent \^etre recueillies ensemble. Les verdicts seuls
suffiraient \`a calculer une distribution marginale, mais les scores de
confiance et les temps de d\'ecision permettent de d\'etecter les r\'eponses
prises au hasard, qui contamineraient l'estimation de l'information
mutuelle si elles \'etaient trait\'ees comme des jugements inform\'es.

La phase d'\textbf{appariement et de contr\^ole} (lignes~5--6) associe
chaque verdict au message effectivement engag\'e, que l'exp\'erimentateur
conna\^it et que les participants ignorent. Elle teste ensuite la condition
d'admissibilit\'e de la D\'efinition~\ref{def:interp_admissible}. Si $V$
d\'epend encore de $M$ une fois $(\Sigma,C)$ connus, les participants ont
acc\`es \`a une information hors preuve de sorte que l'\'etude sort du cadre.

La phase d'\textbf{estimation} (lignes~7--8) produit les deux param\`etres.
L'estimateur $\hat\Delta = \hat I(V;\Sigma|M,C)$ mesure ce que les acteurs
tirent de la preuve \emph{au-del\`a} de ce qu'elle dit du message, qu'il s'agisse
de la mise en forme, des m\'etadonn\'ees ou des indices que l'expertise institutionnelle
(pr\'ec\'edents l\'egaux, expertise de domaine, proc\'edures internes) rend
lisibles. Ce terme co\"{\i}ncide avec celui que la
Proposition~\ref{prop:fondement_extraction} isole dans la d\'ecomposition de
$I(V;\Sigma|C)$. L'estimateur $\hat\alpha$ est
calcul\'e ind\'ependamment comme le rapport
$\log_2|\mathcal{V}| / H(C)$ et mesure la finesse de l'espace de verdicts
relativement \`a la richesse contextuelle, conform\'ement \`a la
d\'efinition~\ref{def:alpha_granularite}. Les deux estimateurs sont donc
acquis par des voies distinctes, ce qui permet de les valider
s\'epar\'ement.

La phase de \textbf{validation} (ligne~9) contr\^ole la fiabilit\'e de la
mesure. Elle combine la stabilit\'e test-retest, qui
v\'erifie que les m\^emes acteurs produisent des r\'eponses coh\'erentes
lorsqu'ils sont r\'einterrog\'es apr\`es un d\'elai, avec la coh\'erence
inter-\'evaluateurs, qui v\'erifie que des acteurs du m\^eme profil
produisent des estimations convergentes. Sans cette \'etape, une valeur
\'elev\'ee de $\hat\Delta$ pourrait refl\'eter une simple variabilit\'e
al\'eatoire au lieu d'une capacit\'e d'extraction r\'eelle.

L'algorithme fournit ainsi une estimation de $\Delta$ dont la stabilit\'e est
contr\^ol\'ee, sans pr\'etendre \`a une valeur exacte. Le cadre CRO n'en demande
pas davantage, puisque $\Delta$ y intervient
comme constante du th\'eor\`eme et doit \^etre mesur\'e avec une incertitude
born\'ee pour que les pr\'edictions aient un sens op\'erationnel.

\subsubsection{Pr\'edictions v\'erifiables}
\label{subsubsec:predictions_delta}

Les pr\'edictions suivantes sont directement
falsifiables~\cite{MINKA2026CRO,MINKA2026SOTP}~:

\begin{enumerate}
\item $\hat\Delta_{\text{juge g\'en\'eral}} \approx 25$ bits
\item $\hat\Delta_{\text{juge PKI}} \approx 40$ bits
\item $\hat\Delta_{\text{r\'egulateur senior}} \approx 65$ bits
\item $\hat\Delta_{\text{ECDSA}} > \hat\Delta_{\text{Helios}}$ pour tout sous-groupe
\item Facteur d'amplification institutionnelle $1{,}8$ \`a $3{,}7$ fois par
      rapport \`a la capacit\'e individuelle Miller ($\approx 20$ bits)
\end{enumerate}
\end{sloppypar}

\begin{boiteArbitre}
L'arbitre note un \'el\'ement nouveau, \`a savoir des chiffres qui fixent des cibles,
avec vingt-cinq bits pour un juge g\'en\'eraliste et soixante-cinq pour un
r\'egulateur senior. Il y voit des paris, qui peuvent donc \^etre perdus.
\end{boiteArbitre}

% -----------------------------------------------------------------------
\subsection{Protocole de mesure de $H(\mathcal{C})$~:
          entropie contextuelle}
\label{subsec:protocole_HC}
% -----------------------------------------------------------------------
\begin{sloppypar}

\subsubsection{Objectif et enjeu}
\label{subsubsec:objectif_HC}

\noindent\textbf{Objectif.} Quantifier empiriquement l'entropie contextuelle
$H(\mathcal{C})$, c'est-\`a-dire la richesse informationnelle du contexte
institutionnel dans lequel s'op\`ere l'interpr\'etation des preuves
cryptographiques.

\noindent\textbf{Enjeu.} $H(\mathcal{C})$ est le param\`etre le plus difficile
\`a mesurer directement, car il caract\'erise la \emph{distribution}
des contextes possibles et non un contexte particulier. Cette distinction
d'\'echelle est explicit\'ee \`a la
section~\ref{sec:methodologie_calcul}. Les protocoles ci-dessous mesurent
l'entropie de la distribution des contextes d'un domaine institutionnel,
tandis que les d\'erivations de la section~\ref{sec:derivation_opposabilite}
mobilisent l'entropie du contexte d'esp\`ece. Les deux grandeurs portent le
m\^eme symbole et ne prennent pas les m\^emes valeurs.

\subsubsection{Approche 1~: d\'ecomposition structurelle}
\label{subsubsec:approche_structurelle}

\noindent\textbf{Principe.} Identifier les dimensions ind\'ependantes du
contexte institutionnel et sommer leurs entropies en supposant
l'ind\'ependance approximative~\cite{MINKA2026CLO}.

\noindent\textbf{Application au contexte contractuel.}
\begin{equation}
H(C) = H(C_{\text{parties}}) + H(C_{\text{clauses}}) + H(C_{\text{juridiction}}) + H(C_{\text{historique}})
\label{eq:decomposition_HC}
\end{equation}

\begin{table}[htbp]
\centering
\caption{Estimations structurelles de $H(\mathcal{C})$ par domaine institutionnel}
\label{tab:estimations_structurelles_HC}
\begin{tabular}{@{}llc@{}}
\toprule
\textbf{Domaine} & \textbf{Composantes dominantes} & $H(\mathcal{C})$ estim\'e \\
\midrule
Contentieux \'electoral & Registre, r\`egles, historique & 30--50 bits \\
Conformit\'e financi\`ere & Transactions, r\'eglementations, audit & 40--70 bits \\
Signature contractuelle & Parties, clauses, juridiction & 20--40 bits \\
\bottomrule
\end{tabular}
\end{table}

\subsubsection{Approche 2~: mod\'elisation probabiliste par MDS}
\label{subsubsec:approche_MDS}

\noindent\textbf{Design.} Pr\'esenter un ensemble de $n_{\text{ctx}} = 100$
contextes vari\'es \`a un panel de 5--10 experts en leur demandant une
\'evaluation de similarit\'e par paires.

\noindent\textbf{Proc\'edure d'analyse.}
\begin{enumerate}
\item Construction de la matrice de dissimilarit\'e $D \in \mathbb{R}^{100 \times 100}$
\item R\'eduction par MDS (multi-dimensional scaling)
\item Calcul du volume effectif $V_{\mathrm{eff}} = \det(\hat\Sigma)^{1/2}$
\item Estimation $H(\mathcal{C}) \approx \log_2 V_{\mathrm{eff}}$
\end{enumerate}

\subsubsection{Approche 3~: consensus d'experts par m\'ethode Delphi}
\label{subsubsec:approche_delphi}

\noindent\textbf{Principe.} La m\'ethode Delphi est une technique
d'\'elicitation it\'erative permettant d'obtenir un consensus
quantitatif~\cite{MINKA2026CRO}.

\noindent\textbf{Proc\'edure en quatre tours.} 20 experts estiment
$H(\mathcal{C})$ ind\'ependamment, puis r\'evisent apr\`es retour
anonymis\'e. Le processus s'arr\^ete lorsque l'IQR devient inf\'erieur \`a $5$ bits.
\end{sloppypar}

% -----------------------------------------------------------------------
\subsection{Protocole de mesure de $\Opp$~:
          opposabilit\'e sur syst\`emes d\'eploy\'es}
\label{subsec:protocole_opp}
% -----------------------------------------------------------------------
\begin{sloppypar}

\subsubsection{Objectif et syst\`eme cible}
\label{subsubsec:objectif_opp}

\noindent\textbf{Objectif.} Mesurer l'opposabilit\'e $\Opp$ empiriquement sur
un syst\`eme cryptographique r\'eellement d\'eploy\'e puis comparer la valeur
mesur\'ee \`a la pr\'ediction du mod\`ele SOTP~\cite{MINKA2026SOTP}.

\noindent\textbf{Syst\`eme cible.} Syst\`eme de vote Helios d\'eploy\'e pour
une \'election de d\'el\'egu\'es \'etudiants \`a l'ENSP (\'Ecole Nationale
Sup\'erieure Polytechnique, Universit\'e de Yaound\'e I).

\subsubsection{Design et participants}
\label{subsubsec:design_opp}

\noindent\textbf{Participants.} Commission \'electorale de $n = 10$ trustees
(\'etudiants doctorants et enseignants-chercheurs).

\noindent\textbf{Taille de l'\'election.} $\approx 500$ \'etudiants votants,
g\'en\'erant $\approx 500$ preuves cryptographiques.

\subsubsection{Proc\'edure}
\label{subsubsec:procedure_opp}

\noindent\textbf{Phase 1. Formation trustees (4 h).}
Formation structur\'ee sur Helios et les preuves z\'ero-connaissance.

\noindent\textbf{Phase 2. \'Election r\'eelle.}
D\'eploiement standard de Helios.

\noindent\textbf{Phase 3. Session d'audit post-\'election (3 h).}
Pour un \'echantillon al\'eatoire de 50 preuves (10\%)~:
\begin{enumerate}
\item Chaque trustee examine ind\'ependamment la preuve
\item Rend un verdict $v \in \mathcal{V}$
\item Note sa confiance sur 10
\end{enumerate}

\subsubsection{Calcul de l'opposabilit\'e empirique}
\label{subsubsec:calcul_opp_empirique}

L'opposabilit\'e est calcul\'ee selon~\cite{MINKA2026CLO}~:
\begin{equation}
\Opp_{\mathrm{emp}} = \frac{I_{\mathrm{emp}}(V;\Sigma|C)}{H_{\mathrm{emp}}(V|C)}
\label{eq:opp_empirique}
\end{equation}

\subsubsection{Validation crois\'ee avec la pr\'ediction SOTP}
\label{subsubsec:validation_croisee_sotp}

\noindent\textbf{Valeur de r\'ef\'erence.} Pour Helios, la calibration
illustrative de l'article CLO~\cite{MINKA2026CLO} donne, \`a l'\'echelle
op\'erationnelle (fiabilit\'e de v\'erification)~:
\begin{equation}
\Conf \cdot \Relver \cdot \Opp \approx 0{,}437
\label{eq:prediction_sotp_opp}
\end{equation}
La valeur cible pour l'opposabilit\'e seule est
$\Opp_{\mathrm{pr\'edit}} = 0{,}50$~; sa d\'erivation compl\`ete (espace de
verdicts \`a quatre d\'ecisions, signal expert binaire d'un bit) est
\'etablie \`a la section~\ref{sec:derivation_helios}. Le
pr\'esent protocole met \`a l'\'epreuve l'hypoth\`ese d'uniformit\'e des verdicts qui
sous-tend cette valeur. Il mesure la distribution empirique $H_{\mathrm{emp}}(V|C)$
au lieu de la postuler.

\noindent\textbf{Crit\`ere de validation.}
\begin{itemize}
\item Accord si $|\Opp_{\mathrm{emp}} - \Opp_{\mathrm{pr\'edit}}| \leq 0{,}05$
\item D\'esaccord (la valeur $0{,}50$ et l'hypoth\`ese d'uniformit\'e qui la fonde
      sont r\'efut\'ees) si $\Opp_{\mathrm{pr\'edit}}$ tombe hors de l'IC~95\,\% bootstrap
      de $\Opp_{\mathrm{emp}}$
\end{itemize}
\end{sloppypar}

% -----------------------------------------------------------------------
\subsection{Protocole standardis\'e d'estimation conjointe de $\Delta$ et $\alpha$}
\label{subsec:protocole_conjoint}
% -----------------------------------------------------------------------
\begin{sloppypar}
Pour une validation prospective compl\`ete du cadre CRO, les
param\`etres $(\Delta, \alpha, \beta)$ doivent \^etre mesur\'es
\emph{ind\'ependamment} et \emph{conjointement} avant que les pr\'edictions
CRO ne soient calcul\'ees.

\begin{algorithm}[H]
\caption{Protocole standardis\'e d'estimation conjointe de
         $(\hat\Delta, \hat\alpha)$~\cite{MINKA2026CRO}}
\label{alg:estimation_conjointe}
\begin{algorithmic}[1]
\State \textbf{Entr\'ees~:} Institution $I$, types de preuves, contexte $C$
\State \textbf{Sorties~:} $\hat\Delta$ (bits), $\hat\alpha$ (sans unit\'e)
\Require Accord institutionnel, \'ethique IRB, $n \geq 30$ participants
\State \textbf{\'Etape 1~:} Collecte de donn\'ees (verdicts, confiance, temps)
\State \textbf{\'Etape 2~:} Appariement des verdicts au message $m$ (v\'erit\'e terrain)
\State \textbf{\'Etape 3~:} Test d'admissibilit\'e $V \perp M \mid (\Sigma,C)$
\State \textbf{\'Etape 4~:} $\hat\Delta \leftarrow \hat I(V;\Sigma \mid M,C)$
\State \textbf{\'Etape 5~:} $\hat\alpha \leftarrow \log_2|\mathcal{V}| / H(C)$
\State \textbf{\'Etape 6~:} Validation (ICC $\geq 0{,}8$, $\kappa \geq 0{,}7$)
\end{algorithmic}
\end{algorithm}

L'algorithme formalise en six \'etapes le protocole de r\'ef\'erence pour
l'estimation conjointe des param\`etres institutionnels $(\Delta, \alpha)$,
en int\'egrant explicitement les pr\'ealables \'ethiques et
m\'ethodologiques que la version op\'erationnelle de la
section~\ref{subsec:protocole_delta} suppose implicites.

La phase de \textbf{param\'etrage} (lignes~1--3) fixe le cadre de
l'exp\'erience. Elle d\'eclare l'institution cible $I$, les types de
preuves \`a pr\'esenter et le contexte $C$, puis nomme les deux sorties,
$\hat\Delta$ en bits et $\hat\alpha$ sans unit\'e. Elle pose surtout trois
\emph{pr\'erequis} obligatoires, \`a savoir l'accord institutionnel pr\'ealable,
l'approbation d'un comit\'e d'\'ethique de la recherche et un effectif
minimal de $n \geq 30$ participants. Leur port\'ee d\'epasse la simple formalit\'e
proc\'edurale, puisqu'ils conditionnent la validit\'e statistique de
l'estimation. En de\c{c}\`a de trente participants, les estimateurs
d'information mutuelle deviennent instables~; l'approbation \'ethique
autorise quant \`a elle la collecte de donn\'ees sur des acteurs humains
identifiables. La ligne~3, marqu\'ee \texttt{Require}, signale que ces
conditions doivent \^etre satisfaites \emph{avant} toute ex\'ecution, car elles
sont des conditions d'entr\'ee de l'algorithme, distinctes de ses \'etapes.

L'\textbf{\'etape~1} (ligne~4) collecte les donn\'ees brutes. Elle
enregistre, pour chaque participant et chaque paire contexte-preuve, les
verdicts rendus, les scores de confiance subjectifs et les
temps de d\'ecision. Ces trois s\'eries sont reprises \`a l'identique de
l'algorithme~\ref{alg:estimation_delta}, dont le protocole standardis\'e
conserve la nature en se bornant \`a en fixer la s\'equence et les crit\`eres
de validit\'e. Les scores de confiance et les temps servent de variables
de contr\^ole pour d\'etecter les r\'eponses non inform\'ees.

Les \textbf{\'etapes~2 et 3} (lignes~5--6) apparient chaque verdict au
message engag\'e, connu du seul exp\'erimentateur, puis v\'erifient que les
participants n'ont pas acc\`es \`a $M$ hors de la preuve (condition
d'admissibilit\'e). Cette v\'erification conditionne l'interpr\'etation de
l'\'etape suivante.

Les \textbf{\'etapes~4 et 5} (lignes~7--8) produisent les deux
estimateurs. L'\'etape~4 applique l'estimateur~\eqref{eq:estimateur_delta}
$\hat\Delta \leftarrow \hat I(V;\Sigma|M,C)$.
L'\'etape~5 calcule $\hat\alpha$ par la formule
$\log_2|\mathcal{V}| / H(C)$, ind\'ependamment de $\hat\Delta$. Les deux calculs sont d\'elib\'er\'ement s\'epar\'es afin de valider
chaque estimateur selon ses propres crit\`eres, sans que l'erreur sur
l'un ne contamine l'autre~; un $\hat\Delta$ biais\'e n'entra\^ine donc pas
m\'ecaniquement un $\hat\alpha$ biais\'e.

L'\textbf{\'etape~6} (ligne~9) valide les estimations par deux
statistiques. Le coefficient de corr\'elation intra-classe
$\mathrm{ICC} \geq 0{,}8$ mesure la coh\'erence test-retest tandis que le
coefficient $\kappa$ de Cohen $\geq 0{,}7$ mesure l'accord
inter-\'evaluateurs au-del\`a du hasard. Ces deux seuils sont plus
exigeants que ceux de l'algorithme~\ref{alg:estimation_delta}, qui se
contentait d'une validation qualitative. Le protocole standardis\'e fixe
donc des crit\`eres num\'eriques pr\'ecis qui le rendent reproductible
par des \'equipes ind\'ependantes. Cette reproductibilit\'e justifie le qualificatif
de \emph{standardis\'e}, puisque deux
laboratoires appliquant les m\^emes \'etapes aux m\^emes entr\'ees doivent
obtenir des estimations compatibles.

Par rapport \`a l'algorithme~\ref{alg:estimation_delta}, cette version
explicite les pr\'erequis \'ethiques et statistiques, num\'erote les \'etapes
pour permettre de les citer individuellement et fixe des seuils quantitatifs
de validation. Ces ajouts font de la proc\'edure de mesure un protocole
standardis\'e, r\'eutilisable au-del\`a du cas d'esp\`ece.
\end{sloppypar}

% -----------------------------------------------------------------------
\subsection{Protocole de mesure de $\Delta_q$~:
          capacit\'e d'extraction quantique}
\label{subsec:protocole_delta_q}
% -----------------------------------------------------------------------
\begin{sloppypar}

\subsubsection{Objectif et d\'efinition op\'erationnelle}
\label{subsubsec:objectif_delta_q}

\noindent\textbf{Objectif.} Mesurer empiriquement la capacit\'e d'extraction
institutionnelle quantique $\Delta_q$ (en bits), d\'efinie
par~\cite{MINKA2026PQCRO}~:
\begin{equation}
I(V^q; \Sigma|C)_\rho \leq I(M;\Sigma|C)_\rho + \Delta_q
\label{eq:extraction_quantique_op}
\end{equation}

\begin{definition}[Capacit\'e d'extraction quantique $\Delta_q$]\label{def:capacite_extraction_quantique_op}
$\Delta_q$ est la quantit\'e d'information sur les verdicts qu'un acteur
institutionnel peut extraire dans un environnement post-quantique,
au-del\`a de l'information accessible de Holevo. Elle se d\'ecompose
comme $\Delta_q = \Delta + \delta_{\mathrm{enh}} - \delta_{\mathrm{degr}}$,
o\`u $\delta_{\mathrm{enh}}$ repr\'esente l'am\'elioration due aux
ressources quantiques et $\delta_{\mathrm{degr}}$ la d\'egradation due au
bruit quantique.
\end{definition}

\noindent\textbf{Estimateur.}
\begin{equation}
\hat{\Delta}_q = \hat I(V^q;\Sigma \mid M, C)
\label{eq:estimateur_delta_q}
\end{equation}
calcul\'e comme~\eqref{eq:estimateur_delta} \`a partir des r\'eponses des
participants, le message $m$ \'etant connu de l'exp\'erimentateur. Pour des
signatures post-quantiques qui sont des cha\^{\i}nes de bits classiques, cet
estimateur est identique \`a l'estimateur classique~; il ne diff\`ere que dans le
sc\'enario hypoth\'etique d'un support ou d'un v\'erificateur quantique.

\subsubsection{Design exp\'erimental}
\label{subsubsec:design_delta_q}

\noindent\textbf{Participants.} $N \geq 80$ professionnels form\'es aux
syst\`emes post-quantiques, stratifi\'es par type d'institution selon la
r\'epartition du tableau~\ref{tab:stratification_delta_q}.

\begin{table}[htbp]
\centering
\caption{Stratification de l'\'echantillon pour la mesure de $\Delta_q$}
\label{tab:stratification_delta_q}
\begin{tabular}{@{}llcc@{}}
\toprule
\textbf{Cat\'egorie} & \textbf{Sous-groupe} & \textbf{$n$} & \textbf{$\Delta_q$ pr\'edit} \\
\midrule
R\'egulateurs PQC & Juniors (formation initiale) & 20 & $\approx 28$ bits \\
                  & Seniors (formation avanc\'ee) & 20 & $\approx 45$ bits \\
\midrule
Auditeurs crypto & Classiques (transition PQC) & 20 & $\approx 35$ bits \\
                 & Experts PQC certifi\'es & 20 & $\approx 67$ bits \\
\midrule
\textbf{Total} & & \textbf{80} & \\
\bottomrule
\end{tabular}
\end{table}

\noindent\textbf{Mat\'eriel exp\'erimental.}
\begin{itemize}
\item 2 sch\'emas de signature post-quantiques, ML-DSA-65 (FIPS~204,
      ex-Dilithium-3) et SLH-DSA (FIPS~205, ex-SPHINCS+)~; ML-KEM (FIPS~203,
      ex-Kyber) est un m\'ecanisme d'encapsulation de cl\'e et ne produit pas de
      preuve, il est donc exclu
\item Simulateur d'environnement quantique avec bruit de d\'epolarisation
      param\'etrable ($p \in [0{,}001;\,0{,}1]$)
\item Interface de v\'erification post-quantique avec m\'etriques de performance
\item Questionnaires adapt\'es au contexte NIST PQC (FIPS 203/204/205)
\end{itemize}

\subsubsection{Proc\'edure}
\label{subsubsec:procedure_delta_q}

La session exp\'erimentale se d\'eroule en six phases.

\noindent\textbf{Phase 0. Formation PQC (3 h).}
Introduction aux standards NIST PQC portant sur la structure des cl\'es, le processus de
signature et les diff\'erences avec RSA/ECDSA. Formation sur les implications
institutionnelles de la transition post-quantique.

\noindent\textbf{Phase 1. Session exp\'erimentale principale (4 h).}
Pour chaque type de preuve post-quantique~:
\begin{enumerate}
\item[(a)] Pr\'esentation du contexte (transaction, document, conformit\'e)
\item[(b)] Pr\'esentation de la signature ML-DSA ou SLH-DSA
\item[(c)] Extraction de verdict avec et sans aide \`a l'interpr\'etation
\item[(d)] Mesure des temps de d\'ecision et scores de confiance
\end{enumerate}

\noindent\textbf{Phase 2. Simulation quantique (2 h).}
Exposition \`a des sc\'enarios simul\'es avec bruit quantique variable~:
\begin{itemize}
\item Condition A, sans bruit (r\'ef\'erence classique)
\item Condition B, bruit faible ($p = 0{,}001$, typique des syst\`emes corrig\'es)
\item Condition C, bruit mod\'er\'e ($p = 0{,}01$, typique NISQ)
\item Condition D, bruit \'elev\'e ($p = 0{,}1$, test de robustesse)
\end{itemize}

\noindent\textbf{Phase 3. Test-retest (1 h).}
Re-pr\'esentation de sc\'enarios apr\`es 2 semaines.

\noindent\textbf{Phase 4. Gold standard PQC.}
Experts en cryptographie post-quantique \'etablissent les verdicts corrects.

\noindent\textbf{Phase 5. D\'ebriefing (30 min).}

\subsubsection{M\'etriques et estimateurs}
\label{subsubsec:metriques_delta_q}

Le protocole calcule $\hat\Delta_q$ et $\hat\eta_q$ selon
l'Algorithme~\ref{alg:estimation_delta_q}.

\begin{algorithm}[H]
\caption{Estimation empirique de $\Delta_q$}
\label{alg:estimation_delta_q}
\begin{algorithmic}[1]
\State \textbf{Entr\'ees~:} Institution $I$, preuves PQC $\{\sigma_1, \sigma_2\}$,
       param\`etre de bruit $p$
\State \textbf{Sortie~:} $\hat\Delta_q$, $\hat\eta_q$
\State Pr\'esenter $\sigma_i$ avec contexte $C$ dans environnement simul\'e
\State Enregistrer verdicts $V^q$, confiance, temps, erreurs d'interpr\'etation
\State Apparier chaque verdict au message $m$ (v\'erit\'e terrain)
\State Tester l'admissibilit\'e $V^q \perp M \mid (\Sigma, C)$
\State Estimer $\hat\Delta_q = \hat I(V^q;\Sigma \mid M, C)$
\State Estimer $\hat\eta_q$ par r\'egression sur les niveaux de bruit
\State Valider par test-retest et accord inter-\'evaluateurs
\end{algorithmic}
\end{algorithm}

L'algorithme transpose au r\'egime post-quantique la proc\'edure
d'estimation de la capacit\'e d'extraction institutionnelle, en y
int\'egrant deux diff\'erences structurelles. L'environnement de
pr\'esentation y est simul\'e avec un param\`etre de bruit quantique
contr\^ol\'e~; la sortie comprend aussi, en plus de $\hat\Delta_q$,
l'estimateur $\hat\eta_q$ du coefficient de perte
d'interpr\'etabilit\'e.

La phase d'\textbf{entr\'ees} (lignes~1--2) fixe le cadre. Elle re\c{c}oit
l'institution cible $I$, un ensemble de deux preuves post-quantiques
$\{\sigma_1, \sigma_2\}$ correspondant aux deux standards de signature NIST
PQC retenus ainsi qu'un param\`etre de bruit $p$ qui contr\^ole l'intensit\'e
de la d\'egradation quantique simul\'ee. Elle d\'eclare comme sorties
l'estimateur $\hat\Delta_q$ de la capacit\'e d'extraction quantique et
l'estimateur $\hat\eta_q$ du coefficient de perte s\'emantique. Cette double
sortie distingue l'algorithme de son homologue classique, puisque
l'algorithme~\ref{alg:estimation_delta} ne mesurait qu'une capacit\'e alors que la
version quantique mesure \'egalement la d\'egradation introduite par le
canal.

La phase de \textbf{collecte} (lignes~3--4) produit les donn\'ees brutes
dans un environnement simul\'e. Chaque preuve $\sigma_i$ est pr\'esent\'ee
avec son contexte $C$, mais contrairement \`a la version classique,
l'environnement de pr\'esentation est param\'etr\'e par $p$, de sorte que les
acteurs institutionnels voient les preuves \`a travers un canal quantique
dont le bruit est contr\^ol\'e. Aux verdicts $V^q$, aux scores de confiance et
aux temps de d\'ecision s'ajoute une quatri\`eme s\'erie, celle des erreurs
d'interpr\'etation. Elle permet de saisir ce que les acteurs extraient
\emph{incorrectement} sous l'effet du bruit, faute de quoi la d\'egradation
s\'emantique resterait invisible dans les donn\'ees.

La phase d'\textbf{appariement et de contr\^ole} (lignes~5--6) reprend la
structure classique, puisque chaque verdict est associ\'e au message engag\'e
avant que l'admissibilit\'e de l'interpr\`ete soit test\'ee. Dans le sc\'enario d'un
support quantique, les quantit\'es sont index\'ees par l'\'etat $\rho$~;
pour des signatures qui sont des cha\^{\i}nes de bits, elles co\"{\i}ncident
avec leurs analogues classiques.

La phase d'\textbf{estimation} (lignes~7--8) produit les deux param\`etres
par des voies distinctes. L'estimateur $\hat\Delta_q = \hat I(V^q;\Sigma|M,C)$ est le
m\^eme que dans le r\'egime classique et quantifie ce que les acteurs
institutionnels tirent de la preuve au-del\`a de ce qu'elle dit du message. L'estimateur $\hat\eta_q$ provient en revanche
d'une r\'egression sur les niveaux de bruit. En r\'ep\'etant
la mesure pour plusieurs valeurs de $p$, on observe comment
l'information extraite d\'ecro\^it avec le bruit~; la pente de cette
d\'ecroissance fournit $\hat\eta_q$. Cette s\'eparation des deux
estimations s'impose, car $\Delta_q$ et $\eta_q$ ont des
interpr\'etations distinctes (l'une mesure une capacit\'e, l'autre une
perte) et doivent \^etre estim\'es ind\'ependamment pour que chacun reste
interpr\'etable dans la borne post-quantique.

La phase de \textbf{validation} (ligne~9) reprend les crit\`eres de
stabilit\'e test-retest et de coh\'erence inter-\'evaluateurs de la
version classique. Ces crit\`eres sont n\'eanmoins plus difficiles \`a
satisfaire en r\'egime quantique, car le bruit introduit une variabilit\'e
suppl\'ementaire dans les r\'eponses. Une faible stabilit\'e test-retest
peut alors refl\'eter soit une incoh\'erence des acteurs, soit une
instabilit\'e intrins\`eque du canal quantique~; distinguer les deux
n\'ecessite de r\'ep\'eter la mesure \`a bruit constant avant de faire
varier $p$.

Par rapport \`a son homologue classique, l'algorithme op\`ere dans un
environnement simul\'e param\'etr\'e, enregistre les erreurs d'interpr\'etation,
calcule des informations mutuelles index\'ees par $\rho$ et produit par
r\'egression un second estimateur $\hat\eta_q$. Ces diff\'erences pr\'eservent
la logique du contraste informationnel, qu'elles \'etendent au r\'egime
quantique en y ajoutant la mesure explicite de la d\'egradation
s\'emantique.

\subsubsection{Pr\'edictions v\'erifiables pour $\Delta_q$}
\label{subsubsec:predictions_delta_q}

\begin{enumerate}
\item $\hat\Delta_q < \hat\Delta$ pour les acteurs non form\'es (p\'enalit\'e d'adaptation)
\item $\hat\Delta_q \to \hat\Delta$ pour les experts PQC (adaptation r\'eussie)
\item $\hat\Delta_q \in [28, 67]$ bits pour les institutions bien pr\'epar\'ees
\item D\'ecroissance de $\Delta_q$ avec le niveau de bruit quantique
\item Corr\'elation positive entre formation PQC et $\Delta_q$
\end{enumerate}
\end{sloppypar}

% -----------------------------------------------------------------------
\subsection{Protocole de mesure de $\eta_q$~:
          coefficient de perte d'interpr\'etabilit\'e quantique}
\label{subsec:protocole_eta_q}
% -----------------------------------------------------------------------
\begin{sloppypar}

\subsubsection{Objectif et d\'efinition op\'erationnelle}
\label{subsubsec:objectif_eta_q}

\noindent\textbf{Objectif.} Mesurer empiriquement le coefficient de perte
d'interpr\'etabilit\'e quantique $\eta_q$, d\'efini
par~\cite{MINKA2026PQCRO}~:
\begin{equation}
I(V^q; \Lambda(\rho_\Sigma)|C)_\rho \;\geq\; I(V^q; \rho_\Sigma|C)_\rho - \eta_q \cdot H(C)
\label{eq:perte_interpretabilite_op}
\end{equation}
($\eta_q$ est le plus petit coefficient pour lequel cette minoration est
vraie~; dans l'autre sens, l'in\'egalit\'e serait satisfaite pour $\eta_q = 0$
par l'in\'egalit\'e de traitement des donn\'ees, ce qui viderait la d\'efinition de son contenu).

\noindent\textbf{Port\'ee.} Les signatures ML-DSA, SLH-DSA et Falcon produisent
des cha\^{\i}nes de bits classiques, pour lesquelles $\eta_q = 0$. Le protocole
ne s'applique qu'\`a un sc\'enario de stockage ou de v\'erification quantique,
qui reste hypoth\'etique \`a ce jour.

\noindent\textbf{D\'efinition op\'erationnelle.} $\eta_q$ quantifie la perte
d'information s\'emantique lorsqu'une preuve cryptographique traverse un canal
quantique (stockage, transmission, v\'erification avec ressources quantiques).

\noindent\textbf{Estimateur.}
\begin{equation}
\hat\eta_q = \frac{I(V^q; \rho_\Sigma|C)_\rho - I(V^q; \Lambda(\rho_\Sigma)|C)_\rho}{H(C)}
\label{eq:estimateur_eta_q}
\end{equation}

\subsubsection{Design exp\'erimental}
\label{subsubsec:design_eta_q}

\noindent\textbf{Principe.} Comparer l'extraction d'information depuis des
preuves avec et sans d\'egradation quantique simul\'ee.

\noindent\textbf{Mat\'eriel.}
\begin{itemize}
\item Simulateur de canaux quantiques (d\'epolarisation, d\'ephasage,
      amortissement d'amplitude)
\item Param\`etres de bruit calibr\'es sur les sp\'ecifications de mat\'eriel
      quantique actuel (IBM Quantum, IonQ, etc.)
\item Protocoles PQC avec et sans assistance quantique
\end{itemize}

\begin{table}[htbp]
\centering
\caption{Param\`etres de canaux quantiques pour la mesure de $\eta_q$}
\label{tab:val_canaux_quantiques}
\begin{tabular}{@{}lccc@{}}
\toprule
\textbf{Canal} & \textbf{Param\`etre} & \textbf{Plage test\'ee} & \textbf{$\eta_q$ attendu} \\
\midrule
D\'epolarisation & $p_{\mathrm{dep}}$ & $[0{,}001;\,0{,}1]$ & $0{,}02$--$0{,}20$ \\
D\'ephasage & $p_{\mathrm{deph}}$ & $[0{,}001;\,0{,}05]$ & $0{,}01$--$0{,}10$ \\
Amortissement & $\gamma_{\mathrm{AD}}$ & $[0{,}001;\,0{,}02]$ & $0{,}005$--$0{,}05$ \\
Combin\'e (NISQ typique) & --- & --- & $0{,}05$--$0{,}15$ \\
\bottomrule
\end{tabular}
\end{table}

\subsubsection{Proc\'edure}
\label{subsubsec:procedure_eta_q}

\noindent\textbf{Phase 1. Calibration.}
\'Etablir la ligne de base sans bruit quantique en mesurant
$I(V^q; \rho_\Sigma|C)_\rho$ pour chaque type de preuve.

\noindent\textbf{Phase 2. D\'egradation contr\^ol\'ee.}
Appliquer des canaux quantiques avec param\`etres croissants~:
\begin{equation}
\Lambda_{\mathrm{dep}}(\rho) = (1-p)\rho + \frac{p}{d}I_d
\label{eq:canal_depolarisant_op}
\end{equation}
Mesurer $I(V^q; \Lambda(\rho_\Sigma)|C)_\rho$ pour chaque niveau.

\noindent\textbf{Phase 3. Analyse de r\'egression.}
Ajuster le mod\`ele~:
\begin{equation}
\Delta I = I(V^q; \rho_\Sigma|C) - I(V^q; \Lambda(\rho_\Sigma)|C) = \eta_q \cdot H(C)
\end{equation}
pour extraire $\hat\eta_q$.

\subsubsection{M\'etriques compl\'ementaires}
\label{subsubsec:metriques_eta_q}

\begin{itemize}
\item \textbf{Fid\'elit\'e moyenne~:} $F_{\mathrm{avg}}(\Lambda) = \mathbb{E}[F(\Lambda(\rho), \rho)]$
\item \textbf{Distance de trace~:} $\|\Lambda(\rho_\Sigma) - \rho_\Sigma\|_1$
\item \textbf{Perte d'information par bit contextuel~:} $\Delta I / H(C)$
\end{itemize}

\subsubsection{Pr\'edictions v\'erifiables pour $\eta_q$}
\label{subsubsec:predictions_eta_q}

\begin{enumerate}
\item $\hat\eta_q \in [0{,}05;\,0{,}15]$ pour un support simul\'e aux taux d'erreur
      du mat\'eriel actuel (valeurs a priori, \`a pr\'e-enregistrer)
\item Ordre de grandeur de $\hat\eta_q\,H(C)$ compatible avec la borne de
      continuit\'e de l'information mutuelle conditionnelle (Alicki--Fannes--Winter),
      exprim\'ee en fonction de la distance de trace
      $\tfrac12\|\Lambda(\rho_\Sigma)-\rho_\Sigma\|_1 \leq \sqrt{1-F}$ (heuristique)
\item Relation croissante, approximativement lin\'eaire pour bruit faible, entre
      $p_{\mathrm{dep}}$ et $\hat\eta_q$
\item $\hat\eta_q = 0$, aux fluctuations d'\'echantillonnage pr\`es, lorsque la
      preuve est lue comme cha\^{\i}ne de bits (contr\^ole n\'egatif)
\end{enumerate}
\end{sloppypar}

% -----------------------------------------------------------------------
\subsection{Protocole combin\'e de validation quantique}
\label{subsec:protocole_combine}
% -----------------------------------------------------------------------
\begin{sloppypar}

\subsubsection{Objectif}
\label{subsubsec:objectif_combine}

Valider conjointement les param\`etres post-quantiques $(\Delta_q, \eta_q)$ et
comparer les pr\'edictions du Trilemme CRO post-quantique avec les mesures
empiriques sur les standards NIST PQC.

\subsubsection{Design int\'egr\'e}
\label{subsubsec:design_combine}

\begin{algorithm}[H]
\caption{Protocole combin\'e post-quantique}
\label{alg:protocole_combine}
\begin{algorithmic}[1]
\State \textbf{Entr\'ees~:} Institution $I$, standards PQC, canaux quantiques
\State \textbf{Sorties~:} $\hat\Delta_q$, $\hat\eta_q$, $\hat\beta_q$
\State \textbf{Partie A. Mesure de $\Delta_q$}
\State Mesurer extraction sur ML-DSA et SLH-DSA sans bruit
\State Calculer $\hat\Delta_q$ par comparaison avec gold standard PQC
\State \textbf{Partie B. Mesure de $\eta_q$}
\State Appliquer canaux quantiques \`a intensit\'e croissante
\State Mesurer d\'ecroissance de l'information mutuelle
\State Calculer $\hat\eta_q$ par r\'egression
\State \textbf{Partie C. Validation de la borne CRO$^q$}
\State V\'erifier les hypoth\`eses du th\'eor\`eme (dont $H(V|C) \geq \gamma$)
\State Si elles tiennent, confronter $\Conf^q \cdot \Relrec^q \cdot \Opp^q$ \`a la borne~\eqref{eq:trilemme_pq_raffine}
\State \textbf{Partie D. Validation crois\'ee}
\State Comparer pr\'ediction vs mesure pour les deux standards PQC
\State Calculer erreur relative et intervalles de confiance
\end{algorithmic}
\end{algorithm}

L'algorithme structure la validation post-quantique en quatre parties
ind\'ependantes, chacune mesurant ou v\'erifiant une composante distincte
de la borne CRO$^q$. L\`a o\`u les algorithmes pr\'ec\'edents isolaient
un param\`etre, celui-ci les articule puisqu'il produit les trois estimateurs
$(\hat\Delta_q, \hat\eta_q, \hat\beta_q)$, puis les confronte \`a la
borne th\'eorique et aux pr\'edictions sur les standards NIST PQC.

La phase d'\textbf{entr\'ees et sorties} (lignes~1--2) fixe le cadre.
Elle re\c{c}oit l'institution cible $I$, les standards PQC \`a tester et
les canaux quantiques disponibles. Elle d\'eclare trois sorties, \`a savoir
$\hat\Delta_q$ (capacit\'e d'extraction quantique), $\hat\eta_q$ (perte
d'interpr\'etabilit\'e) et $\hat\beta_q$ (r\'egularit\'e du protocole en
r\'egime quantique). Cette troisi\`eme sortie est nouvelle par rapport
aux algorithmes pr\'ec\'edents, car elle correspond au param\`etre $\beta$
de l'hypoth\`ese de r\'egularit\'e, transpos\'e au r\'egime quantique. Sa
mesure est n\'ecessaire pour v\'erifier que la borne post-quantique
s'applique bien aux protocoles test\'es.

La \textbf{partie~A} (lignes~3--5) mesure $\Delta_q$ en l'absence de
bruit. Elle \'evalue l'extraction institutionnelle sur les deux
sch\'emas de signature PQC (ML-DSA, SLH-DSA) dans leur
environnement nominal, sans d\'egradation quantique ajout\'ee. Le terme
$\hat\Delta_q$ est ensuite obtenu par comparaison avec un
\emph{gold standard} PQC, constitu\'e des verdicts de r\'ef\'erence \'etablis par
des experts en cryptographie post-quantique. Cette comparaison isole la
capacit\'e d'extraction proprement institutionnelle de toute
d\'egradation li\'ee au canal. La partie~A produit donc la valeur de
$\Delta_q$ qui serait observ\'ee sur un syst\`eme PQC id\'eal, sans
bruit.

La \textbf{partie~B} (lignes~6--9) mesure $\eta_q$ par d\'egradation
contr\^ol\'ee. Elle applique des canaux quantiques \`a intensit\'e
croissante, typiquement un canal d\'epolarisant dont le param\`etre
$p$ varie, puis observe la d\'ecroissance de l'information mutuelle
extraite. Cette d\'ecroissance, non lin\'eaire en $p$,
d\'epend de la structure du protocole et de la nature du canal. La
r\'egression sur les niveaux de bruit fournit $\hat\eta_q$ comme pente
de la relation entre perte d'information et bruit. La partie~B isole
donc l'effet du canal, ind\'ependamment de la capacit\'e
institutionnelle mesur\'ee en partie~A.

La \textbf{partie~C} (lignes~10--12) confronte les mesures \`a la borne
CRO$^q$. Elle v\'erifie d'abord que le cas \'etudi\'e satisfait les hypoth\`eses
du th\'eor\`eme, faute de quoi aucune comparaison n'est probante
(Remarque~\ref{rem:portee_cas_etude}). Si elles tiennent, le produit
$\Conf^q \cdot \Relrec^q \cdot \Opp^q$ est compar\'e \`a la forme raffin\'ee de la
borne (\'equation~\eqref{eq:trilemme_pq_raffine}). Un d\'epassement, hypoth\`eses
v\'erifi\'ees, r\'efute l'une d'elles selon les crit\`eres du
tableau~\ref{tab:criteres_refutation}~; il ne peut pas \^etre absorb\'e par un
r\'eajustement a posteriori de $\hat\Delta_q$.

La \textbf{partie~D} (lignes~13--15) proc\`ede \`a la validation
crois\'ee. Elle compare les pr\'edictions th\'eoriques aux mesures
empiriques pour chacun des deux standards PQC, puis calcule l'erreur
relative et les intervalles de confiance. Au-del\`a d'un contr\^ole de
coh\'erence, cette comparaison permet d'identifier les
standards pour lesquels le mod\`ele pr\'edit correctement les
compromis ainsi que ceux pour lesquels des ph\'enom\`enes non captur\'es par
le mod\`ele interviennent. Les intervalles de confiance quantifient
l'incertitude r\'esiduelle, qui est plus \'elev\'ee en r\'egime
post-quantique en raison de la variabilit\'e suppl\'ementaire introduite
par $\eta_q$.

L'algorithme se distingue des pr\'ec\'edents sur trois points. Il
articule quatre parties ind\'ependantes au lieu de mesurer un seul
param\`etre. Il produit un troisi\`eme estimateur $\hat\beta_q$. Enfin, il valide explicitement la borne
CRO$^q$ au lieu de se contenter de mesurer ses composantes. Cette articulation
fait du protocole combin\'e l'outil de validation complet du r\'egime
post-quantique.

\subsubsection{Crit\`eres de validation quantique}
\label{subsubsec:criteres_validation_quantique}

Les crit\`eres de r\'efutation de l'ensemble du programme, classique et
post-quantique, sont fix\'es avant toute collecte dans le
tableau~\ref{tab:criteres_refutation}. Les valeurs seuils sont indicatives et
seront arr\^et\'ees lors du pr\'e-enregistrement~; leur forme est
l'engagement m\'ethodologique.

\begin{table}[htbp]
\centering
\caption{Crit\`eres de r\'efutation pr\'e-enregistr\'es}
\label{tab:criteres_refutation}
\small
\begin{tabular}{@{}p{3.4cm}p{9.8cm}@{}}
\toprule
\textbf{Param\`etre} & \textbf{L'hypoth\`ese est r\'efut\'ee si\ldots} \\
\midrule
$\Delta$ (Hyp.~\ref{ass:extraction_bornee}) & la borne inf\'erieure de l'IC~95\,\% de $\hat{I}(V;\Sigma|M,C)$ d\'epasse la capacit\'e fix\'ee \emph{avant} l'exp\'erience pour la population \'etudi\'ee. \\
Admissibilit\'e (D\'ef.~\ref{def:interp_admissible}) & un test de permutation rejette $V \perp M \mid (\Sigma,C)$ au seuil $0{,}01$. \\
$\gamma$ (Hyp.~\ref{ass:verdicts_efficace}) & la borne sup\'erieure de l'IC~95\,\% de $\hat{H}(V|C)$ est inf\'erieure \`a $\alpha H(\mathcal{C}) - 1$. \\
$\eta_q$ (sc\'enario quantique) & l'IC~95\,\% de $\hat{\eta}_q$ exclut l'intervalle pr\'edit avant l'exp\'erience. \\
Borne (hypoth\`eses v\'erifi\'ees) & le produit mesur\'e d\'epasse la forme raffin\'ee de la borne avec une confiance de 95\,\%. \\
\bottomrule
\end{tabular}
\end{table}

\end{sloppypar}

\begin{boiteArbitre}
L'arbitre fait le point. Trois param\`etres classiques, trois param\`etres
quantiques, soit six protocoles qui visent tous \`a transformer des hypoth\`eses
en mesures. Il retient surtout leur sym\'etrie, puisque chaque quantit\'e
admise plus haut dispose d\'esormais d'une proc\'edure pour \^etre test\'ee. La
th\'eorie, d\'ej\`a coh\'erente, devient ainsi v\'erifiable.
\end{boiteArbitre}

% -----------------------------------------------------------------------
\subsection{Calendrier, ressources et statut de d\'eploiement}
\label{subsec:calendrier_ressources}
% -----------------------------------------------------------------------
\begin{sloppypar}

\subsubsection{Tableau de synth\`ese}
\label{subsubsec:tableau_synthese}

\begin{table}[htbp]
\centering
\caption{Programme exp\'erimental complet~: calendrier et ressources}
\label{tab:programme_experimental}
\begin{threeparttable}
\begin{tabular}{@{}lllcc@{}}
\toprule
\textbf{Protocole} & \textbf{Section} & \textbf{Dur\'ee} &
\textbf{Participants} & \textbf{Co\^ut estim\'e\tnote{$\star$}} \\
\midrule
\multicolumn{5}{c}{\textit{Param\`etres classiques (Contribution C2)}} \\
\midrule
Mesure de $\Delta$ & \S\ref{subsec:protocole_delta} & 6 mois & 120 institutionnels & 15\,000\,000 FCFA \\
Mesure de $H(\mathcal{C})$ & \S\ref{subsec:protocole_HC} & 3 mois & 20 experts & 5\,000\,000 FCFA \\
Mesure de $\Opp$ sur Helios & \S\ref{subsec:protocole_opp} & 4 mois & 10 trustees & 8\,000\,000 FCFA \\
\midrule
\multicolumn{5}{c}{\textit{Param\`etres quantiques (Contribution C3)}} \\
\midrule
Mesure de $\Delta_q$ & \S\ref{subsec:protocole_delta_q} & 6 mois & 80 professionnels PQC & 12\,000\,000 FCFA \\
Mesure de $\eta_q$ & \S\ref{subsec:protocole_eta_q} & 4 mois & 40 experts + simulation & 8\,000\,000 FCFA \\
Validation combin\'ee PQC & \S\ref{subsec:protocole_combine} & 3 mois & 40 validateurs & 6\,000\,000 FCFA \\
\midrule
\textbf{Total} & & \textbf{26 mois} & \textbf{310} &
\textbf{54\,000\,000 FCFA} \\
\bottomrule
\end{tabular}
\begin{tablenotes}
\footnotesize
\item[$\star$] Les co\^uts indiqu\'es constituent des ordres de grandeur
indicatifs destin\'es \`a fournir une vision globale du programme~;
ils feront l'objet d'une \'evaluation budg\'etaire affin\'ee lors de la
mise en \oe uvre effective de chaque protocole.
\end{tablenotes}
\end{threeparttable}
\end{table}

\subsubsection{Calendrier d\'etaill\'e}
\label{subsubsec:calendrier_detaille}

\begin{table}[htbp]
\centering
\caption{Calendrier du programme exp\'erimental}
\label{tab:calendrier_experimental}
\begin{tabular}{@{}lcccccc@{}}
\toprule
\textbf{Activit\'e} & \textbf{M1--M6} & \textbf{M7--M12} & \textbf{M13--M18} &
\textbf{M19--M24} & \textbf{M25--M30} & \textbf{M31--M36} \\
\midrule
Mesure $\Delta$ (\S\ref{subsec:protocole_delta}) & $\times\times$ & $\times\times$ & & & & \\
Mesure $H(\mathcal{C})$ (\S\ref{subsec:protocole_HC}) & $\times$ & $\times$ & & & & \\
Mesure $\Opp$ Helios (\S\ref{subsec:protocole_opp}) & & $\times\times$ & $\times$ & & & \\
Mesure $\Delta_q$ (\S\ref{subsec:protocole_delta_q}) & & & $\times\times$ & $\times\times$ & & \\
Mesure $\eta_q$ (\S\ref{subsec:protocole_eta_q}) & & & & $\times\times$ & $\times$ & \\
Validation PQC (\S\ref{subsec:protocole_combine}) & & & & & $\times\times$ & $\times$ \\
Analyse longitudinale & & & & & & $\times\times$ \\
\bottomrule
\end{tabular}
\end{table}

\subsubsection{Statut de d\'eploiement du programme}
\label{subsubsec:statut_deploiement}

Le programme d\'ecrit aux sections~\ref{subsec:protocole_delta}
\`a~\ref{subsec:calendrier_ressources} constitue une sp\'ecification
m\'ethodologique. Aucun de ses six
protocoles n'a \'et\'e d\'eploy\'e dans le cadre de cette th\`ese, qu'il s'agisse de la
campagne de mesure de $\Delta$ aupr\`es des 120 acteurs institutionnels
(section~\ref{subsec:protocole_delta}), de l'audit post-\'electoral
Helios \`a l'ENSP (section~\ref{subsec:protocole_opp}) ou des sessions
post-quantiques sur ML-DSA et SLH-DSA
(sections~\ref{subsec:protocole_delta_q}
et~\ref{subsec:protocole_eta_q}).

Ce choix d\'ecoule du positionnement \'epist\'emologique retenu au
chapitre~\ref{chap:cadre_formel}, selon lequel ce travail est d'abord th\'eorique. La
certification m\'ecanique en Lean4 garantit la correction de la d\'erivation \`a
partir des hypoth\`eses. La v\'erit\'e des hypoth\`eses elles-m\^emes rel\`eve de la
validation comportementale, que seules des mesures peuvent \'etablir. Trois pr\'ealables
non techniques conditionnent aussi le d\'eploiement, \`a savoir l'accord des
juridictions et r\'egulateurs approch\'es, l'approbation d'un comit\'e
d'\'ethique de la recherche ainsi que le financement de l'ordre de grandeur
report\'e au tableau~\ref{tab:programme_experimental}, pour vingt-six mois de
travail cumul\'es r\'epartis sur un calendrier de trente-six mois
(tableau~\ref{tab:calendrier_experimental}).

Il en r\'esulte, pour la lecture des sections suivantes, que les
d\'erivations et les estimations qui y figurent sont de
nature r\'etrospective et analytique. Cela vaut pour les valeurs
d'opposabilit\'e d\'eriv\'ees \`a la
section~\ref{sec:derivation_opposabilite} comme pour les param\`etres
institutionnels estim\'es \`a la
section~\ref{sec:validation_retrospective}, puisque les unes et les autres
reposent sur les sp\'ecifications publi\'ees des protocoles et sur des
contextes institutionnels mod\'elis\'es, sans mesure aupr\`es de juges ou
de r\'egulateurs r\'eels. Elles \'etablissent que les param\`etres
institutionnels admettent des valeurs plausibles et mutuellement
coh\'erentes. Seule une mesure ind\'ependante dira si ces valeurs sont les
bonnes. L'audit
post-\'electoral Helios, le plus l\'eger en ressources et le seul
enti\`erement internalisable \`a l'ENSP, constitue le point d'entr\'ee
naturel du programme~; il est repris \`a ce titre dans les perspectives
du chapitre~\ref{chap:resultats_discussions}.

\end{sloppypar}

\begin{boiteArbitre}
L'arbitre referme le dossier des protocoles. Il n'a rien ex\'ecut\'e, rien
mesur\'e, rien valid\'e. Il tient une promesse de mesure rigoureuse, chiffr\'ee et falsifiable, qui reste une promesse.

Il a compris que ces protocoles \'etablissent le caract\`ere \emph{testable}
du cadre, ce qui est moins que sa v\'erit\'e sans \^etre peu. Une th\'eorie qu'on
ne peut pas tester reste sans prise sur le r\'eel, m\^eme si elle est belle.

Il rel\`eve la transparence du dispositif, dont les pr\'edictions
sont num\'eriques, les crit\`eres d'invalidation explicites et les
limites nomm\'ees. L'audit Helios lui-m\^eme est pr\'esent\'e comme le
point d'entr\'ee le plus modeste, qui est aussi le seul internalisable \`a l'ENSP.

Il ne sait pas encore si les chiffres tiennent, mais il sait d\'esormais
\emph{comment} il le saura, ce qui correspond exactement \`a ce qu'il attendait d'un
protocole exp\'erimental.

Les d\'erivations peuvent commencer.
\end{boiteArbitre}
% Section du chapitre 4 (ex-annexe C)
% Annexe issue de l'ancien chapitre 3 (restructuration en 3 chapitres)
% ============================================================================
% CHAPITRE 3 — SECTION 3.3 : DÉRIVATION DES VALEURS D'OPPOSABILITÉ
% Fichier : partie3/derivation_opposabilite.tex
% ============================================================================

\section{D\'ERIVATION DES VALEURS D'OPPOSABILIT\'E DES SYST\`EMES DE R\'EF\'ERENCE}
\label{sec:derivation_opposabilite}
% ============================================================================
\begin{sloppypar}

Cette section s'ouvre sur la grandeur dont tout le reste d\'epend. Les
sections~\ref{sec:validation_retrospective} \`a~\ref{sec:reproductibilite}
manipulent trois produits CLO de r\'ef\'erence ($0{,}273$ pour ECDSA,
$0{,}437$ pour Helios, $0{,}572$ pour les zk-SNARK), calcul\'es avec la
fiabilit\'e de v\'erification $\Relver$. Encore
faut-il savoir d'o\`u ils viennent. La pr\'esente section les d\'erive,
syst\`eme par syst\`eme, \`a partir de la d\'efinition m\^eme de l'opposabilit\'e
\'etablie au chapitre~\ref{chap:cadre_formel}.

Helios, les zk-SNARK de type Groth16 et ECDSA, trois syst\`emes cryptographiques
r\'eellement d\'eploy\'es, servent de r\'ef\'erence. Ils ont \'et\'e
retenus parce qu'ils occupent trois r\'egimes distincts du compromis
confidentialit\'e--opposabilit\'e et que leurs sp\'ecifications
publi\'ees autorisent une analyse dimension par dimension. Les valeurs
obtenues sont des estimations \emph{analytiques}, \'etablies \`a partir des
sp\'ecifications de protocole et de contextes institutionnels mod\'elis\'es. Elles ne
reposent sur aucune mesure aupr\`es de juges ou de r\'egulateurs r\'eels~\cite{MINKA2026CLO}
et leur statut est celui \'etabli \`a la
section~\ref{subsubsec:statut_deploiement}.

\begin{boiteArbitre}
L'arbitre remonte \`a la source. Trois chiffres circulent depuis le d\'ebut de
ce chapitre ($0{,}273$, $0{,}437$, $0{,}572$) et il veut savoir d'o\`u ils
sortent. On lui annonce d'embl\'ee qu'il s'agit de d\'erivations, c'est-\`a-dire
d'estimations analytiques, sans relev\'e aupr\`es de juges r\'eels. Il note cette
distinction, qui tiendra pour tout ce qui suit.
\end{boiteArbitre}

% -----------------------------------------------------------------------
\subsection{M\'ethodologie de calcul en trois \'etapes}
\label{sec:methodologie_calcul}
% -----------------------------------------------------------------------

L'opposabilit\'e a \'et\'e d\'efinie \`a la section~\ref{sec:clo} comme la
fraction d'incertitude institutionnelle r\'esolue par la preuve~:
\begin{equation}
\Opp(\mathcal{P}, C) = \frac{I(V;\Sigma|C)}{H(V|C)}
\label{eq:opp_definition_calcul}
\end{equation}
o\`u $V = \mathsf{Interpret}(\Sigma, C)$ d\'esigne le verdict institutionnel,
$\Sigma$ la preuve cryptographique et $C$ le contexte. Le th\'eor\`eme du goulot
d'\'etranglement contextuel, \'etabli au m\^eme endroit, borne le num\'erateur
par $\min\{H(\Sigma|C), H(V|C)\}$ et garantit donc $\Opp \leq 1$, contrainte que chacune des trois d\'erivations ci-dessous doit respecter.
Chaque estimation suit trois \'etapes, appliqu\'ees identiquement aux trois
syst\`emes.

\noindent\textbf{\'Etape 1. Mod\'elisation de l'espace de verdicts.}
\'Enum\'erer les d\'ecisions institutionnelles possibles \`a partir des cadres
l\'egaux et r\'eglementaires pertinents, puis estimer $H(V|C)$ en supposant une
distribution approximativement uniforme sur ces verdicts. Cette hypoth\`ese
d'uniformit\'e est une ligne de base conservatrice en l'absence de donn\'ees sur
la jurisprudence r\'eelle~; ses limites sont discut\'ees \`a la
section~\ref{sec:limites_derivations}.

\noindent\textbf{\'Etape 2. Estimation de l'extraction d'information.}
Analyser ce qu'un interpr\`ete institutionnel, dot\'e d'une connaissance
contextuelle $H(\mathcal{C})$, peut effectivement tirer de la preuve $\Sigma$
compte tenu de son expertise et de ses moyens de v\'erification. L'information
extractible est quantifi\'ee en bits, d\'ecompos\'ee dimension par dimension puis
born\'ee par l'hypoth\`ese d'extraction born\'ee
$I(V;\Sigma|C) \leq I(M;\Sigma|C) + \Delta$ \'etablie au
chapitre~\ref{chap:cadre_formel}.

\noindent\textbf{\'Etape 3. Calcul de l'opposabilit\'e.}
Appliquer la d\'efinition~\eqref{eq:opp_definition_calcul}, puis composer avec
les m\'etriques $\Conf$ et $\Relver$ pour obtenir le produit CRO.

\noindent\textbf{Deux \'echelles d'entropie contextuelle.} Une pr\'ecaution de
notation s'impose avant de calculer. Les valeurs employ\'ees dans les trois
d\'erivations qui suivent ($8$, $11$ et $20$ bits) d\'esignent l'entropie
du \emph{contexte d'esp\`ece}, c'est-\`a-dire la richesse informationnelle effectivement
mobilis\'ee par l'institution consid\'er\'ee face \`a un litige donn\'e. Elles ne
doivent pas \^etre confondues avec les plages du
tableau~\ref{tab:estimations_structurelles_HC}, de $20$ \`a $70$ bits, qui
portent sur l'entropie de la \emph{distribution des contextes} d'un domaine
institutionnel entier et servent au dimensionnement des protocoles de la
section~\ref{sec:protocoles_experimentaux}. Le rapport entre les deux
\'echelles est celui d'un cas particulier \`a la population dont il est tir\'e.
La distinction est sans effet sur les valeurs d'opposabilit\'e calcul\'ees
ci-dessous, o\`u $H(\mathcal{C})$ n'intervient pas dans le quotient, mais elle
conditionne l'interpr\'etation des \'ecarts entre syst\`emes.

% -----------------------------------------------------------------------
\subsection{Helios~: vote \'electronique en contentieux \'electoral}
\label{sec:derivation_helios}
% -----------------------------------------------------------------------

\noindent\textbf{Protocole.} Helios~\cite{Adida08} est un syst\`eme de vote
ouvert reposant sur le chiffrement ElGamal et sur des preuves
z\'ero-connaissance construites sur des groupes de logarithme discret de
$2048$ bits.

\noindent\textbf{Contexte institutionnel mod\'elis\'e.} Tribunal de juridiction
g\'en\'erale statuant sur un litige \'electoral, avec
$H(\mathcal{C}) \approx 8$ bits, soit environ $16$ cat\'egories de preuves
distinguables multipli\'ees par $16$ contextes proc\'eduraux distincts.

\noindent\textbf{Formulation CLO.}
\begin{itemize}
\item Messages~: votes individuels $v \in \{1,\ldots,K\}$
\item Preuves~: $\sigma = (\mathsf{Enc}(v), \pi_{\text{valid}}, \pi_{\text{tally}})$,
      o\`u $\mathsf{Enc}$ d\'esigne un chiffrement ElGamal et $\pi$ des preuves
      z\'ero-connaissance de type Fiat-Shamir
\item Verdicts~: $\mathcal{V} = \{$\'election valide, \'election invalide,
      recomptage ordonn\'e, investigation demand\'ee$\}$
\end{itemize}

\noindent\textbf{\'Etape 1. Estimation de $H(V|C)$.} L'espace de verdicts
comporte quatre d\'ecisions distinctes, d'o\`u
\begin{equation}
H(V|C) = \log_2|\mathcal{V}| = \log_2(4) = 2 \text{ bits}.
\label{eq:hvc_helios}
\end{equation}
Cette t\'etralogie d\'ecisionnelle (accepter les r\'esultats, les annuler,
ordonner un recomptage, exiger une investigation) correspond aux standards
des juridictions \'electorales document\'ees~\cite{MINKA2026CLO}.

\noindent\textbf{\'Etape 2. Estimation de $I(V;\Sigma|C)$.} Un juge de
juridiction g\'en\'erale ne peut v\'erifier ind\'ependamment ni les calculs de
logarithme discret ElGamal sur $2048$ bits, ni les transformations
Fiat-Shamir, ni la correction du d\'ecompte homomorphe. En pratique, la preuve
lui parvient comme un \emph{signal expert binaire} transmis par
t\'emoignage, sous la forme \guillemotleft~l'expert cryptographique confirme la
validit\'e~\guillemotright~contre \guillemotleft~l'expert identifie une faille
cryptographique~\guillemotright. Cette cat\'egorisation binaire transmet
\begin{equation}
I(V;\Sigma|C) \approx 1 \text{ bit}.
\label{eq:ivc_helios}
\end{equation}
Cette pauvret\'e informationnelle d\'ecoule de la conception m\^eme du protocole,
puisque la propri\'et\'e z\'ero-connaissance masque le contenu des votes tout en
attestant la correction du calcul. Le juge ne peut
distinguer ni les nuances proc\'edurales, ni les irr\'egularit\'es partielles, ni
les anomalies de comportement des trustees \`a partir des preuves seules.

\noindent\textbf{\'Etape 3. Calcul de l'opposabilit\'e.}
\begin{equation}
\Opp(\text{Helios},\, \mathcal{C}_{\text{tribunal}})
= \frac{1 \text{ bit}}{2 \text{ bits}} = 0{,}50
\label{eq:opp_helios}
\end{equation}
Les preuves cryptographiques r\'esolvent exactement la moiti\'e de
l'incertitude institutionnelle. Le bit r\'esiduel rel\`eve du raisonnement
juridique sur la conformit\'e proc\'edurale, l'interpr\'etation de l'intention
des \'electeurs et le droit \'electoral de la juridiction. La r\'ealit\'e
institutionnelle de cet \'ecart est document\'ee par les incidents IACR~2025 et
UCL~2009, o\`u des trustees pourtant experts en cryptographie ont rencontr\'e
des difficult\'es op\'erationnelles avec la gestion des cl\'es
distribu\'ees~\cite{MINKA2026CLO}. Le protocole de mesure de la
section~\ref{subsec:protocole_opp} confronterait cette valeur de $0{,}50$ \`a
$\Opp_{\mathrm{emp}}$ sur une \'election r\'eellement d\'eploy\'ee.

\noindent\textbf{Produit CRO.}
\begin{equation}
\Conf \cdot \Relver \cdot \Opp
\approx 0{,}92 \times 0{,}95 \times 0{,}50 \approx 0{,}437
\label{eq:produit_helios}
\end{equation}
La confidentialit\'e $\Conf \approx 0{,}92$ d\'ecoule de la s\'ecurit\'e
s\'emantique d'ElGamal sous l'hypoth\`ese DDH~; la fiabilit\'e
$\Relver \approx 0{,}95$, de la soundness des preuves ZK
($< 2^{-80}$)~\cite{MINKA2026CLO}.

% -----------------------------------------------------------------------
\subsection{zk-SNARK Groth16~: conformit\'e r\'eglementaire}
\label{sec:derivation_zksnark}
% -----------------------------------------------------------------------

\noindent\textbf{Protocole.} Preuves z\'ero-connaissance pour la
v\'erification de conformit\'e pr\'eservant la vie priv\'ee. Une institution
financi\`ere y prouve que l'ensemble de ses transactions satisfait les r\`egles de
lutte contre le blanchiment sans en r\'ev\'eler le d\'etail. Construction
Groth16 sur courbe BN254, preuves de $128$ octets en forme compress\'ee~\cite{Groth16}.

\noindent\textbf{Contexte institutionnel mod\'elis\'e.} Audit r\'eglementaire
par une autorit\'e financi\`ere sp\'ecialis\'ee, avec
$H(\mathcal{C}) \approx 11$ bits, soit environ $64$ cat\'egories de preuves
documentaires multipli\'ees par $32$ cadres proc\'eduraux.

\noindent\textbf{Formulation CLO.}
\begin{itemize}
  \item Messages~: bases de transactions $\{(t_i, a_i, v_i)\}$
        (transactions, comptes, montants)
  \item Preuves~: zk-SNARK $\pi$ (construction Groth16 sur BN254)
  \item Verdicts~:
    \[
      \mathcal{V} = \left\{
      \begin{array}{l}
        \text{Conforme},\ \text{Non-conforme},\
        \text{Investigation requise},\\
        \text{\'Evaluation de p\'enalit\'e},\
        \text{Action corrective},\ \text{Escalade},\\
        \text{Approbation conditionnelle},\
        \text{Audit forensique complet}
      \end{array}
      \right\}
    \]
\end{itemize}

\noindent\textbf{\'Etape 1. Estimation de $H(V|C)$.} Huit d\'ecisions
distinctes, d'o\`u
\begin{equation}
H(V|C) = \log_2(8) = 3 \text{ bits}.
\label{eq:hvc_zksnark}
\end{equation}
Les d\'ecisions r\'eglementaires financi\`eres sont structurellement plus
nuanc\'ees que les verdicts \'electoraux, car un r\'egulateur peut approuver sous
conditions, exiger des mesures correctives sans sanction imm\'ediate ou
escalader vers l'application de la loi. Ces huit cat\'egories correspondent aux
pratiques r\'eglementaires document\'ees.

\noindent\textbf{\'Etape 2. Estimation de $I(V;\Sigma|C)$.} Les op\'erations de
couplage elliptique $e: \mathbb{G}_1 \times \mathbb{G}_2 \to \mathbb{G}_T$, la
validation des param\`etres du trusted setup et la v\'erification des
contraintes R1CS demeurent inaccessibles \`a un r\'egulateur sans formation
sp\'ecialis\'ee. Deux bits restent n\'eanmoins extractibles~:
\begin{enumerate}
\item \textbf{V\'erification computationnelle} ($1$ bit). La preuve passe ou
      \'echoue les contr\^oles de couplage, ce qui fournit une information binaire accessible via
      un v\'erificateur certifi\'e sans comprendre la cryptographie sous-jacente.
\item \textbf{Cat\'egorisation de conformit\'e} ($1$ bit). Il s'agit d'une classification
      grossi\`ere lisible dans la structure de la preuve (\guillemotleft~satisfait toutes les contraintes~\guillemotright,
      \guillemotleft~conformit\'e limite~\guillemotright,
      \guillemotleft~violation claire~\guillemotright), accessible via les
      m\'etadonn\'ees du circuit.
\end{enumerate}
\begin{equation}
I(V;\Sigma|C) \approx 2 \text{ bits}.
\label{eq:ivc_zksnark}
\end{equation}
L'avantage sur Helios ($2$ contre $1$ bit) tient au contexte plus riche
($11$ contre $8$ bits), qui autorise une cat\'egorisation plus fine. L'opacit\'e
cryptographique reste la m\^eme, puisque la garantie z\'ero-connaissance de
Groth16 ne peut \^etre contourn\'ee par l'expertise.

\noindent\textbf{\'Etape 3. Calcul de l'opposabilit\'e.}
\begin{equation}
\Opp(\text{zk-SNARK},\, \mathcal{C}_{\text{r\'egulateur}})
= \frac{2 \text{ bits}}{3 \text{ bits}} \approx 0{,}67
\label{eq:opp_zksnark}
\end{equation}
Le tiers restant exige un jugement r\'eglementaire sur la mat\'erialit\'e, le
risque syst\'emique et la priorit\'e d'application, dimensions inaccessibles
\`a la v\'erification cryptographique.

\noindent\textbf{Produit CRO.}
\begin{equation}
\Conf \cdot \Relver \cdot \Opp
\approx 0{,}89 \times 0{,}96 \times 0{,}67 \approx 0{,}572
\label{eq:produit_zksnark}
\end{equation}
La confidentialit\'e $\Conf \approx 0{,}89$ refl\`ete le z\'ero-connaissance
computationnel sous hypoth\`eses de couplage~; la fiabilit\'e
$\Relver \approx 0{,}96$ est l\'eg\`erement inf\'erieure \`a l'unit\'e en raison des
exigences du trusted setup~\cite{MINKA2026CLO}.

% -----------------------------------------------------------------------
\subsection{ECDSA~: litige contractuel avec infrastructure PKI}
\label{sec:derivation_ecdsa}
% -----------------------------------------------------------------------

\noindent\textbf{Protocole.} Authentification par ECDSA sur courbe NIST
P-256~\cite{Johnson01}.

\noindent\textbf{Contexte institutionnel mod\'elis\'e.} Tribunal civil statuant
sur un litige contractuel, avec infrastructure PKI \'etablie et
$H(\mathcal{C}) \approx 20$ bits, d\'ecompos\'es en compr\'ehension PKI
($\approx 2^{10}$ types de certificats, $10$ bits), pr\'ec\'edents
jurisprudentiels ($\approx 2^{8}$ cat\'egories accumul\'ees en d\'ecennies de
droit des signatures num\'eriques, $8$ bits) et cadres l\'egaux (ESIGN, eIDAS et
r\'eglementations sp\'ecifiques, $2$ bits)~\cite{MINKA2026CLO}.

\noindent\textbf{Formulation CLO.}
\begin{itemize}
\item Messages~: documents contractuels, publics par destination
\item Preuves~: signatures ECDSA $(r, s)$ avec \'equation de v\'erification
      publique $sG = R + eP_{\text{pub}}$ sur NIST P-256
\item Verdicts~: espace d\'ecisionnel multidimensionnel couvrant validit\'e,
      force d'authentification, poids juridique, admissibilit\'e proc\'edurale,
      validit\'e temporelle, confirmation d'autorit\'e et pertinence contextuelle
\end{itemize}

\noindent\textbf{\'Etape 1. Estimation de $H(V|C)$.} L'espace de verdicts est
ici multidimensionnel par nature, le tribunal devant \'evaluer simultan\'ement
plusieurs facettes ind\'ependantes de la preuve num\'erique~:
\begin{equation}
H(V|C) \approx 20 \text{ bits}.
\label{eq:hvc_ecdsa}
\end{equation}
Cette valeur co\"incide avec $H(\mathcal{C}) = 20$ bits, ce qui traduit un
contexte institutionnel calibr\'e pour r\'esoudre exactement le niveau
d'incertitude d\'ecisionnelle qu'il affronte.

\noindent\textbf{\'Etape 2. Estimation de $I(V;\Sigma|C)$.} Contrairement aux
preuves z\'ero-connaissance, ECDSA est \emph{transparent par conception}, car
l'\'equation de v\'erification expose les relations math\'ematiques explicites
entre cl\'e publique, hachage du message et composants de la signature. Le
tribunal en extrait environ $18$ bits r\'epartis sur cinq
dimensions~\cite{MINKA2026CLO}~:
\begin{enumerate}
\item \textbf{V\'erification de propri\'et\'e de la cl\'e publique} ($3$ bits), qui recouvre la
      validation de la cha\^{\i}ne de certificats, \'evaluation de la confiance
      en l'\'emetteur, v\'erification des restrictions d'usage de la cl\'e.
\item \textbf{Liaison message--signature} ($4$ bits), qui recouvre l'int\'egrit\'e du hachage,
      la conformit\'e du format de signature, la validation des param\`etres
      algorithmiques et les propri\'et\'es de non-r\'epudiation.
\item \textbf{Validit\'e temporelle} ($3$ bits), qui recouvre l'horodatage, la p\'eriode de
      validit\'e du certificat et le statut de r\'evocation (CRL/OCSP).
\item \textbf{\'Evaluation de l'autorit\'e l\'egale} ($3$ bits), qui recouvre l'autorisation du
      signataire, les conditions de formation du contrat et la conformit\'e
      juridictionnelle.
\item \textbf{Pertinence contextuelle} ($5$ bits), qui recouvre l'alignement avec le type de
      transaction, les exigences r\'eglementaires, les standards industriels, le profil de
      risque et l'ad\'equation au contexte litigieux.
\end{enumerate}
\begin{equation}
I(V;\Sigma|C) \approx 3 + 4 + 3 + 3 + 5 = 18 \text{ bits}.
\label{eq:ivc_ecdsa}
\end{equation}
ECDSA ne cherche pas \`a masquer l'information sur le signataire et le
message, puisqu'un contrat a pr\'ecis\'ement vocation \`a \^etre interpr\'et\'e par des
tiers. La transparence cryptographique, qui p\'enalise lourdement la
confidentialit\'e, est exactement ce qui rend le syst\`eme institutionnellement
puissant. Il s'agit donc d'un compromis d\'elib\'er\'e de conception.

\noindent\textbf{\'Etape 3. Calcul de l'opposabilit\'e.}
\begin{equation}
\Opp(\text{ECDSA},\, \mathcal{C}_{\text{tribunal+PKI}})
= \frac{18 \text{ bits}}{20 \text{ bits}} = 0{,}90
\label{eq:opp_ecdsa}
\end{equation}
Cette interpr\'etabilit\'e de $90\%$ montre que la s\'ecurit\'e
cryptographique peut coexister avec l'utilit\'e institutionnelle.
Les dix pour cent r\'esiduels refl\`etent la discr\'etion judiciaire
irr\'eductible sur l'interpr\'etation contractuelle, l'\'evaluation de la bonne
foi et les consid\'erations \'equitables, qui transcendent la v\'erification
cryptographique~\cite{MINKA2026CLO}.

\noindent\textbf{Produit CRO.}
\begin{equation}
\Conf \cdot \Relver \cdot \Opp
\approx 0{,}31 \times 0{,}98 \times 0{,}90 \approx 0{,}273
\label{eq:produit_ecdsa}
\end{equation}
La faible confidentialit\'e ($\Conf \approx 0{,}31$) est la cons\'equence
directe de la transparence, puisque les documents contractuels sont publics et que seules
les cl\'es priv\'ees demeurent cach\'ees. La tr\`es haute fiabilit\'e
($\Relver \approx 0{,}98$) tient \`a la v\'erification d\'eterministe, la
contrefa\c{c}on supposant la r\'esolution du probl\`eme ECDLP.

% -----------------------------------------------------------------------
\subsection{\'Etude de cas~: ML-DSA-65 (Dilithium-3) en contexte judiciaire}
\label{sec:etude_cas_dilithium}
% -----------------------------------------------------------------------

\noindent\textbf{Protocole.} Signature post-quantique ML-DSA-65 (FIPS~204,
ex-Dilithium-3, cat\'egorie de s\'ecurit\'e NIST~3) sur r\'eseaux euclidiens module-LWE~\cite{Dilithium2022}.

\noindent\textbf{Contexte institutionnel mod\'elis\'e.} Tribunal civil statuant
sur un litige contractuel avec infrastructure PQC en transition, avec
$H(\mathcal{C}) \approx 10{,}5$ bits \`a l'\'echelle de l'esp\`ece (la table de
r\'ef\'erence retient $20$~bits \`a l'\'echelle du domaine,
tableau~\ref{tab:parametres_reference}).

\noindent\textbf{Formulation CLO.}
\begin{itemize}
\item Messages~: documents contractuels
\item Preuves~: signatures Dilithium-3 avec v\'erification sur r\'eseaux
\item Verdicts~: espace d\'ecisionnel similaire \`a ECDSA, mais avec
      incertitude accrue sur l'interpr\'etation des preuves PQC
\end{itemize}

\noindent\textbf{\'Etape 1. Estimation de $H(V|C)$.} L'espace de verdicts
comporte les m\^emes dimensions qu'ECDSA, mais la granularit\'e est
l\'eg\`erement r\'eduite par l'inadaptation institutionnelle~:
\begin{equation}
H(V|C) \approx 10{,}5 \text{ bits}.
\label{eq:hvc_dilithium}
\end{equation}

\noindent\textbf{\'Etape 2. Estimation de $I(V;\Sigma|C)$.} La v\'erification
Dilithium expose des relations alg\'ebriques sur r\'eseaux, moins famili\`eres
aux acteurs institutionnels que les relations ECDSA. L'extraction est born\'ee
par
\begin{equation}
I(V;\Sigma|C) \approx 8 \text{ bits}.
\label{eq:ivc_dilithium}
\end{equation}

\noindent\textbf{\'Etape 3. Calcul de l'opposabilit\'e.}
\begin{equation}
\Opp^q(\text{Dilithium-3},\, \mathcal{C}_{\text{tribunal PQC}})
= \frac{8 \text{ bits}}{10{,}5 \text{ bits}} \approx 0{,}76
\label{eq:opp_q_dilithium}
\end{equation}

\noindent\textbf{Param\`etres post-quantiques.} Pour ML-DSA-65,
$\lambda = 192$. La signature est une cha\^{\i}ne de bits classique (3\,309 octets),
donc $\eta_q = 0$ (Remarque~\ref{rem:forme_raffinee_pq}). La capacit\'e d'extraction
est fix\'ee, \`a titre illustratif, \`a partir de celle d'ECDSA dans le m\^eme usage
(tableau~\ref{tab:parametres_reference})~:
\begin{align}
\Delta_q &= \Delta + \delta_{\text{enh}} - \delta_{\text{degr}}
\approx 35 - 7 = 28 \text{ bits (hypoth\`ese de travail)}
\label{eq:delta_q_dilithium}\\
\eta_q &= 0
\label{eq:eta_q_dilithium}
\end{align}
La confidentialit\'e et la fiabilit\'e se lisent comme pour ECDSA. Le contrat
sign\'e \'etant public, on a $\Conf^q = 0$ au sens de la
D\'efinition~\ref{def:confidentialite_clo}~; la fiabilit\'e de v\'erification est
quant \`a elle proche de~$1$ (infalsifiabilit\'e sous MLWE/MSIS)~:
\begin{align}
\Conf^q &= 0
\label{eq:conf_q_dilithium}\\
\Relver^q &\approx 0{,}98
\label{eq:rel_q_dilithium}
\end{align}
La formulation initiale de cette \'etude retenait $\Conf^q \approx 0{,}95$, valeur
difficile \`a concilier avec $\Conf \approx 0{,}31$ attribu\'ee \`a ECDSA dans le m\^eme usage.

\noindent\textbf{Lecture.} Avec $\Conf^q = 0$ au sens statistique, le produit
CLO$^q$ est nul et la borne est trivialement respect\'ee. La formulation initiale
calculait $0{,}95 \times 0{,}98 \times 0{,}76 \approx 0{,}71$ et le comparait \`a
$(\Delta_q + \eta_q H(C))/\lambda \approx 0{,}147$~:
\begin{equation}
0{,}71 > 0{,}147.
\label{eq:produit_dilithium}
\end{equation}
Ce d\'epassement s'explique comme pour les syst\`emes classiques. Dans ce sc\'enario,
la signature ML-DSA n'est pas une instance du
Th\'eor\`eme~\ref{thm:trilemme_cro_conditionnel} (message public avec $\Relver$ proche
de~$1$ au lieu de $\Relrec$). La pr\'esente formulation pr\'ecise aussi
l'interpr\'etation initiale, selon laquelle \og la borne ne s'applique qu'\`a la
composante cryptographique\fg{}, puisque le th\'eor\`eme borne le produit entier.
L'\'etude de cas illustre la mani\`ere dont les param\`etres $\Delta_q$ et $\eta_q$
s'estimeraient, sans constituer une v\'erification de la borne.

\begin{boiteArbitre}
L'arbitre note l'\'ecart, puis sa cause. Une signature sur un contrat public
ne cache rien. Sa confidentialit\'e \'etant nulle, la comparer \`a la borne n'a
pas de sens. L'\'etude de cas lui montre comment on estimerait $\Delta_q$,
sans pr\'etendre v\'erifier la borne.
\end{boiteArbitre}

% -----------------------------------------------------------------------
\subsection{Synth\`ese des d\'erivations}
\label{sec:synthese_derivations}
% -----------------------------------------------------------------------

\begin{table}[htbp]
\centering
\caption{Composantes du calcul d'opposabilit\'e pour les trois syst\`emes de
r\'ef\'erence. Toutes les valeurs sont des estimations analytiques \'etablies
\`a partir des sp\'ecifications de protocole et de contextes institutionnels
mod\'elis\'es~\cite{MINKA2026CLO}.}
\label{tab:composantes_opposabilite}
\small
\begin{tabular}{@{}lcccccc@{}}
\toprule
\textbf{Syst\`eme} & $H(\mathcal{C})$ & $H(V|C)$ & $I(V;\Sigma|C)$ & $\Opp$ &
$\Conf$ & \textbf{Produit} \\
\midrule
Helios (tribunal g\'en\'eral) & 8 bits & 2 bits & 1 bit & 0{,}50 & 0{,}92 & 0{,}437 \\
zk-SNARK (r\'egulateur financier) & 11 bits & 3 bits & 2 bits & 0{,}67 & 0{,}89 & 0{,}572 \\
ECDSA (tribunal PKI) & 20 bits & 20 bits & 18 bits & 0{,}90 & 0{,}31 & 0{,}273 \\
\bottomrule
\end{tabular}
\end{table}

\noindent\textbf{Deux \'echelles d'entropie contextuelle.} Les valeurs de
$H(\mathcal{C})$ de ce tableau d\'ecrivent le contexte d'une \emph{d\'ecision
particuli\`ere} (\'echelle de l'esp\`ece)~; celles du chapitre~\ref{chap:cadre_formel}
d\'ecrivent la richesse du contexte institutionnel \emph{g\'en\'eral} (\'echelle du
domaine). Les deux sont r\'eunies dans le tableau~\ref{tab:parametres_reference}.

\noindent\textbf{Mise en regard de deux estimations.} Les produits de la
derni\`ere colonne ont \'et\'e obtenus par la voie \emph{descendante}
(espace de verdicts, information extractible, composition avec $\Conf$ et
$\Relver$). La section~\ref{sec:estimation_params_classiques} propose une
d\'ecomposition additive des capacit\'es $\Delta$ ($35$, $56$ et $73$
bits), dont les ordres de grandeur sont compatibles avec les rapports
$\Delta/\lambda$ correspondants~:
\begin{equation}
\frac{35}{128} \approx 0{,}273, \qquad
\frac{56}{128} \approx 0{,}437, \qquad
\frac{73}{128} \approx 0{,}570.
\label{eq:concordance_voies}
\end{equation}
Cette mise en regard ne vaut pas validation, pour deux raisons. (i)~Les valeurs de $\Delta$ ont
\'et\'e initialement fix\'ees par inversion $\Delta := \lambda\cdot(\Conf\cdot\Relver\cdot\Opp)$
(tableau~\ref{tab:validation_cro_classique}), de sorte que les deux voies ne sont
pas ind\'ependantes. (ii)~Ces syst\`emes ne satisfont pas les hypoth\`eses du
th\'eor\`eme (Remarque~\ref{rem:portee_cas_etude}). En particulier,
$H(V|C) = 2 < \gamma = 8{,}6$ pour Helios et $3 < 15{,}5$ pour les zk-SNARKs, avec
les valeurs de $H(\mathcal{C})$ de cette section. Lus comme des instances, les
produits d\'epasseraient $\Delta/(\gamma\lambda)$ d'un facteur de l'ordre de $8{,}6$,
$15{,}5$ et $19$. La mise en regard montre seulement que les ordres de grandeur
retenus sont mutuellement coh\'erents, ce qui est le r\^ole d'une calibration
illustrative.

% -----------------------------------------------------------------------
\subsection{Limites propres \`a ces d\'erivations}
\label{sec:limites_derivations}
% -----------------------------------------------------------------------

Deux limitations affectent sp\'ecifiquement la m\'ethode de calcul expos\'ee
ci-dessus, en sus de la r\'eserve g\'en\'erale \'etablie \`a la
section~\ref{subsubsec:statut_deploiement}.

\noindent\textbf{Uniformit\'e des verdicts.} L'\'etape~1 suppose une
distribution approximativement uniforme sur l'espace des verdicts, faute de
donn\'ees sur la jurisprudence r\'eelle. La distribution effective des
d\'ecisions judiciaires dans des litiges impliquant des preuves
cryptographiques pourrait \^etre fortement asym\'etrique, comme le sugg\`ere par exemple l'acceptation
quasi syst\'ematique d'une signature ECDSA bien \'etablie. Une
telle asym\'etrie abaisserait $H(V|C)$ et rel\`everait m\'ecaniquement $\Opp$. Les
valeurs report\'ees au tableau~\ref{tab:composantes_opposabilite} en
d\'ependent, dans un sens qui d\'epend aussi de l'effet de l'asym\'etrie sur
$I(V;\Sigma|C)$. Le protocole de la
section~\ref{subsec:protocole_opp} mesure pr\'ecis\'ement la distribution
empirique $H_{\mathrm{emp}}(V|C)$~; lever cette hypoth\`ese fait partie de ses
objectifs.

\noindent\textbf{Caract\`ere statique du mod\`ele.} L'analyse traite
$H(\mathcal{C})$ et $\Delta$ comme des constantes. En r\'ealit\'e, les
institutions apprennent et s'adaptent, car l'accumulation de pr\'ec\'edents
augmente $\Delta$ tandis que le changement technologique modifie $H(\mathcal{C})$. Un
mod\`ele dynamique reste \`a construire~; sa calibration rel\`everait de l'analyse
longitudinale pr\'evue au calendrier (tableau~\ref{tab:calendrier_experimental}),
sur trois \`a cinq ans.

\begin{boiteArbitre}
L'arbitre referme la d\'erivation. Il a vu trois syst\`emes occupant chacun un r\'egime
propre, avec des chiffres d\'efendus dimension par dimension. Il a vu aussi, \`a la fin,
deux limites nomm\'ees sans d\'etour, \`a savoir l'uniformit\'e des verdicts dont
d\'ependent les valeurs ainsi que le caract\`ere statique d'un mod\`ele qui ignore
l'apprentissage des institutions.

Il retient que rien n'est cach\'e, puisque les hypoth\`eses sont pos\'ees, les limites
dites et les protocoles susceptibles de les lever d\'esign\'es. Il y voit une estimation
transparente, dont on conna\^it la port\'ee comme les limites.

Les trois chiffres ont d\'esormais une origine explicite. Ce sont des
calibrations que la suite du chapitre manipule comme telles.
\end{boiteArbitre}
\end{sloppypar}
% Section du chapitre 4 (ex-annexe D)
% Annexe issue de l'ancien chapitre 3 (restructuration en 3 chapitres)
% ============================================================================
% CHAPITRE 3 — SECTION 3.4 : VALIDATION RÉTROSPECTIVE
% Fichier : partie3/donnees_validation.tex
% ============================================================================

\section{VALIDATION R\'ETROSPECTIVE}
\label{sec:validation_retrospective}
% ============================================================================
\begin{sloppypar}

La section pr\'ec\'edente a sp\'ecifi\'e \emph{comment} mesurer les
param\`etres institutionnels. La pr\'esente section \'etablit \emph{quelles
valeurs} ces param\`etres prennent, par estimation r\'etrospective sur les
trois syst\`emes de r\'ef\'erence classiques et les standards
post-quantiques. Aucune mesure nouvelle n'est introduite, de sorte que les chiffres qui
suivent sont des calibrations illustratives, d\'eriv\'ees des sp\'ecifications
publi\'ees et mises en regard des produits CLO de la
section~\ref{sec:derivation_opposabilite}.
\end{sloppypar}

\begin{boiteArbitre}
L'arbitre change de posture, car apr\`es un plan, il assiste \`a une
confrontation. Ce qui avait \'et\'e admis plus haut ($\Delta \in [35, 73]$
bits, $\alpha \in [0{,}8, 1{,}5]$, $\beta = O(1)$) revient maintenant sous
forme de valeurs calibr\'ees, syst\`eme par syst\`eme. Il s'agit de savoir
si les ordres de grandeur retenus se tiennent entre eux, sans v\'erifier le
th\'eor\`eme, dont ces syst\`emes ne sont pas des instances.
\end{boiteArbitre}

% -----------------------------------------------------------------------
\subsection{Estimation d\'etaill\'ee des param\`etres institutionnels classiques}
\label{sec:estimation_params_classiques}
% -----------------------------------------------------------------------
\begin{sloppypar}

La section~\ref{sec:derivation_opposabilite} a d\'eriv\'e les produits CLO des
trois syst\`emes de r\'ef\'erence par la voie descendante, en partant des
espaces de verdicts. La pr\'esente section emprunte la voie oppos\'ee et
d\'ecompose les param\`etres institutionnels $(\Delta, \alpha, \beta)$ qui
param\`etrent les bornes du chapitre~\ref{chap:cadre_formel}. Les deux voies ne
sont pas ind\'ependantes (section~\ref{sec:synthese_derivations})~; leur mise en
regard est une calibration.

\subsubsection{Composantes g\'en\'eriques de la capacit\'e d'extraction}
\label{subsec:composantes_generiques}

La capacit\'e d'extraction institutionnelle $\Delta$ agr\`ege six familles de
ressources, dont le tableau~\ref{tab:composantes_delta} donne les plages
observ\'ees dans la litt\'erature cognitive et organisationnelle.

\begin{table}[htbp]
\centering
\caption{Familles de composantes de la capacit\'e d'extraction institutionnelle
         $\Delta$ et plages associ\'ees}
\label{tab:composantes_delta}
\begin{tabular}{@{}lcc@{}}
\toprule
\textbf{Composante} & \textbf{Plage} & \textbf{Base} \\
\midrule
Pr\'ec\'edents l\'egaux et jurisprudence & 10--15 bits & Cat\'egories de cas document\'es \\
Formation d'expert et connaissances & 10--20 bits & Domaine de sp\'ecialisation \\
R\`egles proc\'edurales et m\'emoire institutionnelle & 5--10 bits & Protocoles \'ecrits \\
D\'elib\'eration collaborative & 0--15 bits & Amplification collective \\
Capacit\'e temporelle de d\'elib\'eration & 5--10 bits & Temps allou\'e \\
Ressources externes et consultation & 5--15 bits & Experts, r\'ef\'erences \\
\bottomrule
\end{tabular}
\end{table}

\noindent\textbf{Lecture du tableau.} Ces plages ne sont pas cumulatives, car une
institution donn\'ee ne mobilise qu'un sous-ensemble de ces familles. Certaines se
recouvrent d'ailleurs partiellement, puisque la consultation d'experts externes
et la formation interne alimentent la m\^eme ressource cognitive. La
sommation na\"{\i}ve des bornes sup\'erieures produirait $85$ bits, valeur
sup\'erieure \`a toutes les estimations retenues dans cette th\`ese. Le
tableau~\ref{tab:decomposition_delta} proc\`ede donc syst\`eme par syst\`eme,
en ne retenant que les composantes effectivement mobilis\'ees dans chaque
contexte institutionnel.

\subsubsection{D\'ecomposition de $\Delta$ par syst\`eme}
\label{subsec:decomposition_delta}

La capacit\'e d'extraction institutionnelle $\Delta$ est d\'ecompos\'ee, \`a titre
illustratif, en composantes cognitives et organisationnelles. Les totaux
($15+10+10+0 = 35$, $15+20+15+6 = 56$, $25+20+15+13 = 73$) sont ceux de la table
de r\'ef\'erence (tableau~\ref{tab:parametres_reference})~; la d\'ecomposition du
tableau~\ref{tab:validation_cro_classique} du chapitre~\ref{chap:cadre_formel} y renvoie.

\begin{table}[htbp]
\centering
\caption{D\'ecomposition d\'etaill\'ee de $\Delta$ (bits) par syst\`eme classique}
\label{tab:decomposition_delta}
\begin{tabular}{@{}lcccl@{}}
\toprule
\textbf{Composante de $\Delta$} & \textbf{ECDSA} & \textbf{Helios} &
\textbf{zk-SNARK} & \textbf{Justification} \\
\midrule
Pr\'ec\'edent l\'egal / jurisprudence
  & 15 & 15 & 25
  & Doctrine \'etablie \\
Expertise technique / formation
  & 10 & 20 & 20
  & Formation sp\'ecialis\'ee \\
Proc\'edures / cadres institutionnels
  & 10 & 15 & 15
  & Standards, protocoles \\
D\'elib\'eration collaborative
  & 0 & 6 & 13
  & Comit\'es multi-parties \\
\midrule
\textbf{Total $\Delta$} & \textbf{35} & \textbf{56} & \textbf{73} & \\
\midrule
\textbf{R\'egime (Rem.~\ref{rem:activite_delta})} & rare & rare & riche & \\
\bottomrule
\end{tabular}
\end{table}

\subsubsection{Coh\'erence avec la rationalit\'e born\'ee}
\label{subsec:coherence_rationalite}

Les valeurs $\Delta \in [35, 73]$ bits correspondent \`a des facteurs
d'amplification de $1{,}8$ \`a $3{,}7$ par rapport \`a une capacit\'e
individuelle de r\'ef\'erence fix\'ee \`a $\approx 20$ bits. Ce rep\`ere est une
analogie inspir\'ee des travaux sur les limites de la m\'emoire de
travail~\cite{Miller56}. Miller ne l'a pas mesur\'e~; la d\'ecomposition
ci-dessous est donc heuristique.

\begin{equation}
\text{Amplification} = 1{,}5_{\text{doc}} \times 1{,}5_{\text{collab}} \times 1{,}3_{\text{sp\'ec}} \approx 2{,}9\times
\end{equation}

\subsubsection{Param\`etres $\alpha$ et $\beta$}
\label{subsec:parametres_alpha_beta}

\begin{table}[htbp]
\centering
\caption{Param\`etres $\alpha$ et $\beta$ par syst\`eme classique}
\label{tab:parametres_alpha_beta}
\begin{tabular}{@{}lcccl@{}}
\toprule
\textbf{Syst\`eme} & $\alpha$ & $\beta$ & \textbf{Variante SOTP} &
\textbf{Justification} \\
\midrule
ECDSA & 1{,}0 & 1{,}0 & SOTP-Low
  & Validit\'e binaire, nonce uniforme \\
Helios & 1{,}2 & 1{,}4 & SOTP-Base
  & Audit d\'etaill\'e, irr\'egularit\'e ZK \\
zk-SNARKs & 1{,}5 & 1{,}5 & SOTP-High
  & Conformit\'e gradu\'ee, circuit optimis\'e \\
\bottomrule
\end{tabular}
\end{table}
\end{sloppypar}

% -----------------------------------------------------------------------
\subsection{Estimation des param\`etres quantiques}
\label{sec:estimation_params_quantiques}
% -----------------------------------------------------------------------
\begin{sloppypar}

\subsubsection{D\'ecomposition de $\Delta_q$ par syst\`eme PQC}
\label{subsec:decomposition_delta_q}

La capacit\'e d'extraction quantique $\Delta_q$ int\`egre les modifications
induites par la transition quantique (section~\ref{sec:cro_pq}).

\begin{table}[htbp]
\centering
\caption{D\'ecomposition d\'etaill\'ee de $\Delta_q$ (bits) par syst\`eme PQC}
\label{tab:decomposition_delta_q}
\footnotesize
\begin{tabular}{@{}lcccl@{}}
\toprule
\textbf{Composante} & \textbf{ML-DSA-65} & \textbf{ML-KEM-1024}$^\dagger$ &
\textbf{SLH-DSA} & \textbf{Justification} \\
\midrule
Capacit\'e classique $\Delta$ & 35 & 45 & 30 & Contexte institutionnel \\
P\'enalit\'e d'adaptation $(-\delta_{\text{degr}})$ & $-7$ & $-6$ & $-5$ & Nouveau paradigme \\
Am\'elioration quantique $(+\delta_{\text{enh}})$ & $+0$ & $+0$ & $+0$ & Pas encore exploitable \\
\midrule
\textbf{Total $\Delta_q$} & \textbf{28} & \textbf{39} & \textbf{25} & \\
\midrule
\textbf{Plage NIST PQC} & \multicolumn{3}{c}{$\Delta_q \in [25, 67]$ bits} & Hypoth\`ese \\
\bottomrule
\end{tabular}

\smallskip
{\footnotesize $^\dagger$ML-KEM est un m\'ecanisme d'encapsulation de cl\'e, qui ne produit
pas de preuve. Il est donc hors du champ du Trilemme et n'est conserv\'e que pour la
tra\c{c}abilit\'e avec les valeurs publi\'ees (fichier \texttt{06\_calcul/data/calibration\_pq.csv}).\par}
\end{table}

\noindent\textbf{Justification de la p\'enalit\'e d'adaptation.} La transition
vers les standards NIST PQC impose une courbe d'apprentissage aux institutions.
La p\'enalit\'e $\delta_{\text{degr}}$ refl\`ete~:
\begin{itemize}
\item Nouveau format de cl\'es (structur\'e diff\'eremment de RSA/ECDSA)
\item Nouveaux algorithmes de v\'erification (lattice-based vs discrete log)
\item Manque de pr\'ec\'edents l\'egaux pour les signatures post-quantiques
\item Documentation technique encore en maturation
\end{itemize}

\subsubsection{Estimation de $\eta_q$ par niveau de bruit quantique}
\label{subsec:estimation_eta_q}

Pour les signatures post-quantiques, qui sont des cha\^{\i}nes de bits
classiques, $\eta_q = 0$. Le tableau~\ref{tab:estimation_eta_q} d\'ecrit le
\emph{sc\'enario hypoth\'etique} d'un support ou d'un v\'erificateur quantique~;
les valeurs sont des ordres de grandeur a priori.

\begin{table}[htbp]
\centering
\caption{Estimation de $\eta_q$ par type de canal quantique}
\label{tab:estimation_eta_q}
\begin{tabular}{@{}lcccc@{}}
\toprule
\textbf{Canal / Environnement} & \textbf{Param\`etre} & \textbf{Fid\'elit\'e $F$} &
$\eta_q$ \textbf{estim\'e} & \textbf{Source} \\
\midrule
Syst\`eme corrig\'e (id\'eal) & $p \leq 10^{-4}$ & $\geq 0{,}9999$ & $0{,}005$--$0{,}01$ & Th\'eorique \\
NISQ actuel (IBM, IonQ) & $p \approx 10^{-3}$ & $\approx 0{,}99$ & $0{,}05$--$0{,}10$ & Mesures 2024--2025 \\
NISQ bruit\'e & $p \approx 10^{-2}$ & $\approx 0{,}95$ & $0{,}10$--$0{,}20$ & Sc\'enario pessimiste \\
Simulation hybride & Variable & Variable & $0{,}02$--$0{,}15$ & Param\'etrable \\
\midrule
\textbf{Valeur de sc\'enario} & \multicolumn{3}{c}{$\eta_q \approx 0{,}08$ (m\'ediane NISQ)} & \cite{MINKA2026PQCRO} \\
\bottomrule
\end{tabular}
\end{table}

\noindent\textbf{Relation heuristique.}
\begin{equation}
\eta_q \lesssim 2\sqrt{1 - F_{\mathrm{avg}}}
\label{eq:relation_eta_q_fidelite}
\end{equation}
Cette relation n'a pas valeur de th\'eor\`eme. L'in\'egalit\'e de Fuchs--van de Graaf
borne la distance de trace par $\sqrt{1-F}$ tandis que la continuit\'e de l'information
mutuelle conditionnelle (Alicki--Fannes--Winter) convertit cette distance en perte
d'information avec un facteur qui d\'epend de la dimension du syst\`eme au lieu de
$H(C)$. Elle fixe un ordre de grandeur pour le sc\'enario.

\subsubsection{Param\`etres $\alpha_q$ et $\beta_q$ pour PQC}
\label{subsec:parametres_alpha_q_beta_q}

\begin{table}[htbp]
\centering
\caption{Param\`etres post-quantiques par syst\`eme NIST PQC ($H(\mathcal{C})$ \`a l'\'echelle de l'esp\`ece)}
\label{tab:parametres_alpha_q_beta_q}
\begin{tabular}{@{}lcccl@{}}
\toprule
\textbf{Syst\`eme} & $\alpha_q$ & $\beta_q$ & $H(\mathcal{C})$ &
\textbf{Justification} \\
\midrule
ML-DSA-65 & 1{,}0 & 1{,}2 & 10{,}5 bits & Signature (FIPS~204) \\
ML-KEM-1024$^\dagger$ & 1{,}1 & 1{,}3 & 12{,}0 bits & KEM, hors champ \\
SLH-DSA & 0{,}9 & 1{,}1 & 9{,}0 bits & Signature (FIPS~205) \\
\bottomrule
\end{tabular}
\end{table}
\end{sloppypar}

% -----------------------------------------------------------------------
\subsection{Comparaison des produits CLO classiques et quantiques}
\label{sec:comparaison_produits_cro}
% -----------------------------------------------------------------------
\begin{sloppypar}

Le tableau~\ref{tab:comparaison_produits_cro} rassemble les produits CLO
illustratifs des syst\`emes \'etudi\'es. Pour les trois syst\`emes classiques,
la colonne \guillemotleft~Produit CLO~\guillemotright{} reprend les valeurs
d\'eriv\'ees \`a la section~\ref{sec:derivation_opposabilite} (avec $\Relver$).
Pour les standards post-quantiques, o\`u aucune d\'erivation descendante n'est
disponible faute de contexte institutionnel document\'e, les produits sont
obtenus par transposition du rapport $\Delta_q/\lambda$ major\'e d'une correction
d'adaptation, si bien qu'ils constituent des hypoth\`eses de travail.

% -----------------------------------------------------------------------
\subsection{Tableau comparatif principal}
\label{sec:tableau_comparatif_principal}
% -----------------------------------------------------------------------

\begin{table}[htbp]
\centering
\caption{Produits CLO illustratifs, classiques et post-quantiques ($\lambda = 128$)}
\label{tab:comparaison_produits_cro}
\begin{tabular}{@{}llcccc@{}}
\toprule
\textbf{Type} & \textbf{Syst\`eme} & $\Delta/\Delta_q$ & $\eta_q$ (sc\'enario) &
\textbf{Produit CLO} & \textbf{\'Ecart vs classique} \\
\midrule
\multicolumn{6}{c}{\textit{Syst\`emes classiques (calibration)}} \\
\midrule
Classique & ECDSA & 35 & --- & 0{,}273 & --- \\
Classique & Helios & 56 & --- & 0{,}437 & --- \\
Classique & zk-SNARKs & 73 & --- & 0{,}572 & --- \\
\midrule
\multicolumn{6}{c}{\textit{Syst\`emes post-quantiques (hypoth\`ese de travail)}} \\
\midrule
PQC & ML-DSA-65 & 28 & 0{,}08 & 0{,}35 & --- \\
PQC & ML-KEM-1024$^\dagger$ & 39 & 0{,}08 & 0{,}42 & --- \\
PQC & SLH-DSA & 25 & 0{,}08 & 0{,}30 & --- \\
\midrule
\multicolumn{6}{c}{\textit{Comparaison dans des usages voisins}} \\
\midrule
Classique $\to$ PQC & ECDSA $\to$ ML-DSA & 35 $\to$ 28 & 0 $\to$ 0{,}08 &
  $0{,}273 \to 0{,}35$ & $+28\%$ \\
Classique $\to$ PQC & RSA $\to$ SLH-DSA & 40 $\to$ 25 & 0 $\to$ 0{,}08 &
  $0{,}35 \to 0{,}30$ & $-14\%$ \\
\bottomrule
\end{tabular}

\smallskip
{\footnotesize $^\dagger$Hors champ (m\'ecanisme d'encapsulation de cl\'e).\par}
\end{table}

\noindent\textbf{Observation.} La p\'enalit\'e d'adaptation ($\Delta_q < \Delta$)
abaisse la borne. Le terme $\eta_q H(C)$, qui n'intervient que dans le sc\'enario
d'un support quantique, la rel\`eve l\'eg\`erement. Le produit
ML-DSA augmente pourtant par rapport \`a ECDSA alors que $\Delta_q$ diminue, ce qui tient
au fait que la transposition utilis\'ee pour ces produits n'est pas d\'eriv\'ee de la borne.

\subsubsection{Analyse de la contribution du terme $\eta_q \cdot H(C)$}

La majoration affaiblie $(\Delta_q + \eta_q \cdot H(C))/\lambda$ agr\`ege deux
termes d'ordres de grandeur tr\`es diff\'erents. Le
tableau~\ref{tab:decomposition_borne_pq} en donne la d\'ecomposition et la
confronte au produit illustratif du tableau~\ref{tab:comparaison_produits_cro}.

\begin{table}[htbp]
\centering
\caption{D\'ecomposition de la majoration CRO$^q$ et confrontation au produit
         illustratif ($\lambda = 128$, sc\'enario $\eta_q > 0$)}
\label{tab:decomposition_borne_pq}
\small
\begin{tabular}{@{}lcccccc@{}}
\toprule
\textbf{Syst\`eme} & $\Delta_q/\lambda$ & $\eta_q \cdot H(C)/\lambda$ &
\textbf{Borne} & \textbf{Part $\eta_q$} & \textbf{Produit} & \textbf{D\'epass\'ee~?} \\
\midrule
ML-DSA-65 & $0{,}219$ & $0{,}006$ & $0{,}225$ & $2{,}7\%$ & $0{,}35$ & oui \\
ML-KEM-1024$^\dagger$ & $0{,}305$ & $0{,}007$ & $0{,}312$ & $2{,}2\%$ & $0{,}42$ & oui \\
SLH-DSA & $0{,}195$ & $0{,}006$ & $0{,}201$ & $3{,}0\%$ & $0{,}30$ & oui \\
\midrule
\multicolumn{7}{l}{\textbf{Sc\'enario bruit\'e} ($\eta_q = 0{,}15$)} \\
\midrule
ML-DSA-65 & $0{,}219$ & $0{,}012$ & $0{,}231$ & $5{,}2\%$ & $0{,}35$ & oui \\
ML-KEM-1024$^\dagger$ & $0{,}305$ & $0{,}014$ & $0{,}319$ & $4{,}4\%$ & $0{,}42$ & oui \\
\bottomrule
\end{tabular}
\end{table}

\noindent\textbf{Interpr\'etation du poids quantique.} Dans le sc\'enario
d'un support quantique ($\eta_q \approx 0{,}08$), la contribution du terme
quantique reste modeste, avec $2$ \`a $3\,\%$ de la majoration pour cette valeur contre $5\,\%$ pour
$\eta_q = 0{,}15$. Pour les signatures r\'eelles, $\eta_q = 0$ et la contribution
est nulle. La capacit\'e d'extraction reste le terme dominant.

\noindent\textbf{D\'epassement de la borne.} Les trois produits illustratifs
d\'epassent la majoration ($0{,}35 > 0{,}225$~; $0{,}42 > 0{,}312$~;
$0{,}30 > 0{,}201$). Ces syst\`emes ne sont
pas des instances du th\'eor\`eme (message public, $\Relver$ au lieu de $\Relrec$,
Remarque~\ref{rem:portee_cas_etude}) et les produits eux-m\^emes restent des
hypoth\`eses de travail. Le d\'epassement ne r\'efute donc pas le Trilemme. Il
emp\^eche d'affirmer que la borne est respect\'ee et d\'esigne $\Delta_q$ comme
premi\`ere cible du protocole de la section~\ref{subsec:protocole_delta_q}.

\end{sloppypar}

\begin{boiteArbitre}
L'arbitre s'arr\^ete sur un chiffre qui ne colle pas. Les produits sont
au-dessus de la borne, de dix \`a treize centi\`emes. Loin de le cacher, le texte
le nomme en donnant sa raison, puisque ces syst\`emes ne sont pas des instances
du th\'eor\`eme. L'arbitre note que la borne n'est ni confirm\'ee ni r\'efut\'ee
ici et qu'il reviendra \`a la mesure de $\Delta_q$ de trancher.
\end{boiteArbitre}

% -----------------------------------------------------------------------
\subsection{Propagation d'incertitude}
\label{sec:intervalles_bootstrap}
% -----------------------------------------------------------------------
\begin{sloppypar}

\subsubsection{M\'ethodologie}
\label{subsec:methodologie_bootstrap}

L'incertitude portant sur les param\`etres institutionnels est propag\'ee
jusqu'au rapport $\Delta/\lambda$ en affectant $\Delta$ d'un coefficient de
variation de $20\,\%$ pour les syst\`emes classiques et de $25\,\%$ pour les
syst\`emes post-quantiques, o\`u la p\'enalit\'e d'adaptation ajoute une source
d'incertitude. Ces coefficients sont des hypoth\`eses de travail~; ils ne sont
pas tir\'es d'une mesure.

\subsubsection{R\'esultats classiques}
\label{subsec:resultats_classiques_bootstrap}

Un coefficient de variation de $20\,\%$ sur $\Delta$ se propage int\'egralement
\`a $\Delta/\lambda$, dont l'intervalle \`a $95\,\%$ vaut environ $[0{,}17;\,0{,}38]$
(ECDSA), $[0{,}27;\,0{,}61]$ (Helios) et $[0{,}35;\,0{,}79]$ (zk-SNARKs). Cette
largeur traduit l'incertitude sur les param\`etres institutionnels~; elle
rappelle qu'en l'absence de donn\'ees, les valeurs de ce chapitre sont des
ordres de grandeur.

\begin{remark}[Intervalles de la formulation initiale]
\label{rem:ic_anterieurs}
La formulation initiale de cette section donnait des intervalles de demi-largeur
d'environ $1{,}8\,\%$ ($[0{,}268;\,0{,}278]$ pour ECDSA, par exemple), centr\'es sur
les valeurs calibr\'ees. Elle concluait que les \og mesures CLO\fg{} y tombaient.
Ces intervalles ne sont pas compatibles avec le coefficient de variation annonc\'e
(un tirage de $200\,000$ \'echantillons le montre, voir le script
\texttt{06\_calcul/python/verifications.py} du d\'ep\^ot)~; la comparaison portait
sur des valeurs identiques. La pr\'esente formulation ne les reprend pas.
\end{remark}

\subsubsection{R\'esultats quantiques}
\label{subsec:resultats_quantiques_bootstrap}

Avec un coefficient de variation de $25\,\%$ sur $\Delta_q$, le rapport
$\Delta_q/\lambda$ ($\lambda = 128$) a pour intervalle \`a $95\,\%$ environ
$[0{,}11;\,0{,}33]$ pour ML-DSA ($\Delta_q = 28$) et $[0{,}10;\,0{,}29]$ pour SLH-DSA
($\Delta_q = 25$). Le terme $\eta_q H(C)/\lambda$ est nul pour ces signatures,
qui sont des cha\^{\i}nes de bits classiques~; dans le sc\'enario hypoth\'etique
d'un support quantique ($\eta_q \approx 0{,}08$), il vaut moins d'un centi\`eme
et ne modifie pas ces intervalles. Les intervalles post-quantiques de
la formulation initiale ne sont pas repris, pour les m\^emes raisons.
\end{sloppypar}

% -----------------------------------------------------------------------
\subsection{Coh\'erence avec nos travaux ant\'erieurs}
\label{sec:validation_croisee_benchmarks}
% -----------------------------------------------------------------------
\begin{sloppypar}

\subsubsection{Rep\`eres classiques}
\label{subsec:benchmarks_classiques}

Le tableau~\ref{tab:benchmarks_classiques} situe les trois valeurs de $\Delta$
retenues par rapport aux rep\`eres propos\'es dans nos travaux ant\'erieurs et
dans la litt\'erature sur la rationalit\'e born\'ee. Hormis la r\'ef\'erence basse,
ces rep\`eres sont des auto-citations qui \'etablissent la coh\'erence avec nos
publications sans apporter de confirmation ind\'ependante.

\begin{table}[htbp]
\centering
\caption{Rep\`eres pour $\Delta$~: coh\'erence avec nos travaux ant\'erieurs}
\label{tab:benchmarks_classiques}
\begin{tabular}{@{}lcll@{}}
\toprule
\textbf{Institution} & $\Delta$ \textbf{estim\'e} &
\textbf{Source} & \textbf{Statut} \\
\midrule
Tribunaux petites cr\'eances & $\approx 25$ bits & \cite{Simon82} & Ordre de grandeur (analogie) \\
Autorit\'es fiscales & $\approx 65$ bits & \cite{MINKA2026CRO} & Travaux ant\'erieurs \\
R\'egulateurs financiers & $\approx 75$ bits & \cite{MINKA2026CRO} & Travaux ant\'erieurs \\
\midrule
ECDSA (tribunaux contrats) & 35 bits & \cite{MINKA2026SOTP} & Calibration \\
Helios (commissions \'elect.) & 56 bits & \cite{MINKA2026SOTP} & Calibration \\
zk-SNARKs (r\'egulateurs) & 73 bits & \cite{MINKA2026SOTP} & Calibration \\
\bottomrule
\end{tabular}
\end{table}

\subsubsection{Rep\`eres quantiques projet\'es}
\label{subsec:benchmarks_quantiques}

Aucun rep\`ere historique n'existe pour $\Delta_q$, car la transition PQC est en
cours et les institutions n'ont pas encore accumul\'e de pratique. Le
tableau~\ref{tab:benchmarks_quantiques} propose une grille de projection
index\'ee sur le niveau de pr\'eparation institutionnelle, qui rel\`eve de
l'hypoth\`ese.

\begin{table}[htbp]
\centering
\caption{Grille projet\'ee pour $\Delta_q$ selon le niveau de pr\'eparation PQC (hypoth\`ese)}
\label{tab:benchmarks_quantiques}
\begin{tabular}{@{}lcll@{}}
\toprule
\textbf{Niveau de pr\'eparation} & $\Delta_q$ \textbf{projet\'e} &
\textbf{Contexte} & \textbf{Horizon} \\
\midrule
Non pr\'epar\'e & $15$--$25$ bits & Transition d\'ebutante & 2025--2026 \\
Partiellement pr\'epar\'e & $25$--$45$ bits & Formation en cours & 2026--2027 \\
Bien pr\'epar\'e & $45$--$67$ bits & Standards int\'egr\'es & 2027--2028 \\
Expert PQC & $60$--$80$ bits & Anticipation avanc\'ee & 2028+ \\
\midrule
\multicolumn{4}{l}{\textit{Valeurs de calibration retenues}} \\
\midrule
ML-DSA-65 (non pr\'epar\'e) & 28 bits & Signature & 2025 \\
SLH-DSA (non pr\'epar\'e) & 25 bits & Signature (hachage) & 2025 \\
\bottomrule
\end{tabular}
\end{table}
\end{sloppypar}

% -----------------------------------------------------------------------
\subsection{\'Elicitation de plausibilit\'e}
\label{sec:evaluation_experts}
% -----------------------------------------------------------------------
\begin{sloppypar}

Les estimations ont fait l'objet d'\'echanges informels avec des interlocuteurs
de profils juridique, technique et r\'eglementaire. Faute d'un protocole
document\'e (fonctions des personnes consult\'ees, date, grille d'\'evaluation,
consentement), ces \'echanges ne sont pas pr\'esent\'es comme un r\'esultat de
cette th\`ese. Les scores de la formulation initiale ne sont pas repris. Une
\'elicitation structur\'ee, de type Delphi, pourrait \^etre int\'egr\'ee au
programme de la section~\ref{sec:protocoles_experimentaux}.
\end{sloppypar}

% -----------------------------------------------------------------------
\subsection{D\'ependance en $\lambda$}
\label{sec:stabilite_lambda}
% -----------------------------------------------------------------------
\begin{sloppypar}
\label{subsec:stabilite_classiques}\label{subsec:stabilite_quantiques}

\`A $\Delta$ fix\'e, la borne d\'ecro\^it en $1/\lambda$. Le rapport $\Delta/\lambda$
vaut ainsi $0{,}273$ pour ECDSA \`a $\lambda = 128$ et $0{,}137$ \`a $\lambda = 256$. Il s'agit
d'une pr\'ediction du mod\`ele, qui reste \`a tester. La formulation initiale
pr\'esentait des produits \og stables\fg{} en $\lambda$ ($0{,}273 \to 0{,}270$). La
forme de la borne conduit \`a ne pas reprendre ces tableaux. Que le produit CLO d'un
syst\`eme d\'eploy\'e varie peu avec $\lambda$ est plausible, mais cela illustre
justement que ces syst\`emes ne sont pas des instances du th\'eor\`eme
(Remarque~\ref{rem:portee_cas_etude}).
\end{sloppypar}

% -----------------------------------------------------------------------
\subsection{Test de significativit\'e}
\label{sec:test_significativite}
% -----------------------------------------------------------------------
\begin{sloppypar}
\label{subsec:test_classiques}\label{subsec:test_quantiques}

La formulation initiale rapportait une corr\'elation $r = 0{,}9997$ entre les deux
voies d'estimation, avec $p < 10^{-6}$ par permutation. Ce test n'est pas repris,
pour deux raisons. Avec trois points, un test de permutation ne compte que $3! = 6$
permutations, de sorte que la plus petite $p$-valeur atteignable est $1/6$. Surtout, la
corr\'elation mesurait une identit\'e, puisque $\Delta$ avait \'et\'e fix\'e par
$\Delta := \lambda \cdot (\Conf \cdot \Relver \cdot \Opp)$.

Pour les syst\`emes post-quantiques, la confrontation des produits CLO$^q$
illustratifs \`a la borne figure au tableau~\ref{tab:decomposition_borne_pq}, qui montre que
les trois produits la d\'epassent. Ce d\'epassement ne r\'efute pas le
th\'eor\`eme, dont ces syst\`emes ne sont pas des instances. Il emp\^eche en
revanche d'affirmer que la borne est \og respect\'ee\fg.
\end{sloppypar}

% -----------------------------------------------------------------------
\subsection{Donn\'ees brutes et format CSV}
\label{sec:donnees_brutes}
% -----------------------------------------------------------------------
\begin{sloppypar}

\subsubsection{Structure du jeu de donn\'ees}
\label{subsec:structure_csv}

Les valeurs de ce chapitre sont consign\'ees dans deux fichiers du d\'ep\^ot
compagnon. Le premier, \texttt{06\_calcul/data/parametres\_reference.csv}, reproduit le
tableau~\ref{tab:parametres_reference} et les triplets CLO des syst\`emes
classiques~; le second, \texttt{06\_calcul/data/calibration\_pq.csv}, reproduit le
tableau~\ref{tab:decomposition_borne_pq}. Leur sch\'ema complet figure dans
\texttt{06\_calcul/data/SCHEMA\_fr.md}~; le tableau~\ref{tab:structure_csv} en donne les
colonnes principales. Toutes les valeurs sont des param\`etres de calibration
illustratifs, sans mesure.

\begin{table}[htbp]
\centering
\caption{Colonnes principales des fichiers de donn\'ees}
\label{tab:structure_csv}
\begin{tabular}{@{}lll@{}}
\toprule
\textbf{Colonne} & \textbf{Type} & \textbf{Description} \\
\midrule
\texttt{system} & cha\^{\i}ne & Identifiant syst\`eme \\
\texttt{H\_C\_domaine}, \texttt{H\_C\_espece} & r\'eel & Entropie contextuelle (deux \'echelles) \\
\texttt{Delta}, \texttt{Delta\_q} & r\'eel & Capacit\'e d'extraction (calibration) \\
\texttt{alpha}, \texttt{beta} & r\'eel & Granularit\'e, r\'egularit\'e \\
\texttt{eta\_q} & r\'eel & Perte d'interpr\'etabilit\'e (sc\'enario~; $0$ pour une cha\^{\i}ne de bits) \\
\texttt{lambda} & entier & Param\`etre de s\'ecurit\'e \\
\texttt{product\_clo} & r\'eel & Produit CLO (calibration illustrative, $\Relver$) \\
\texttt{borne\_these} & r\'eel & Majoration affaiblie $(\Delta_q + \eta_q H(C))/\lambda$ \\
\bottomrule
\end{tabular}
\end{table}

\label{subsec:echantillon_classiques}\label{subsec:echantillon_quantiques}
Les colonnes \texttt{product\_pred}, \texttt{product\_meas} et \texttt{borne} de la
formulation initiale deviennent \texttt{product\_clo} et \texttt{borne\_these}, car
le jeu de donn\'ees ne contient aucune mesure et la \og pr\'ediction\fg{} co\"{\i}ncidait
avec le produit CLO.
\end{sloppypar}

% -----------------------------------------------------------------------
\subsection{Synth\`ese}
\label{sec:synthese_validation}
% -----------------------------------------------------------------------
\begin{sloppypar}

Le tableau~\ref{tab:synthese_validation} r\'ecapitule ce que cette section
\'etablit ainsi que ce qu'elle n'\'etablit pas.

\begin{table}[htbp]
\centering
\caption{Synth\`ese de la calibration r\'etrospective}
\label{tab:synthese_validation}
\small
\begin{tabular}{@{}llccl@{}}
\toprule
\textbf{Crit\`ere} & \textbf{M\'etrique} & \textbf{Classique} & \textbf{PQC} &
\textbf{Statut} \\
\midrule
Plausibilit\'e des $\Delta$ & Intervalle cognitif $[20,80]$ & 3/3 & 2/2 & Ordre de grandeur \\
Hypoth\`ese~2$'$ & $H(V|C) \geq \gamma$ & 1/3 (ECDSA) & --- & Hors champ sinon \\
Produit $\leq$ borne & --- & non applicable & d\'epass\'ee & Syst\`emes hors champ \\
Incertitude & IC~95\,\% de $\Delta/\lambda$ & $\pm 39\,\%$ & $\pm 49\,\%$ & CV suppos\'e \\
\bottomrule
\end{tabular}
\end{table}

\noindent\textbf{Conclusion.}
\begin{itemize}
\item \textbf{Classique (C2/C4).} Les ordres de grandeur retenus sont
  mutuellement coh\'erents~; aucune validation empirique n'est revendiqu\'ee.
\item \textbf{Post-quantique (C3).} Les produits illustratifs d\'epassent la
  borne~; les param\`etres $\Delta_q$ sont des postulats \`a mesurer en priorit\'e.
\item \textbf{R\'eserve.} La validation prospective demeure n\'ecessaire. Les
  protocoles de la section~\ref{sec:protocoles_experimentaux} fournissent les
  mesures de $\Delta$, $\Delta_q$ et $\eta_q$ qui restent \`a produire, avec les
  crit\`eres de r\'efutation du tableau~\ref{tab:criteres_refutation}~; leur
  statut de d\'eploiement est \'etabli \`a la
  section~\ref{subsubsec:statut_deploiement}.
\end{itemize}
\end{sloppypar}

% -----------------------------------------------------------------------
\subsection{Reproductibilit\'e}
\label{sec:reproductibilite}
% -----------------------------------------------------------------------
\begin{sloppypar}

\subsubsection{D\'ep\^ot}
\label{subsec:depot}

\begin{center}
\url{https://github.com/MINKAThierry/PhDThesis}\\
(version 1.0.0, archiv\'ee sous le DOI \texttt{10.5281/zenodo.XXXXXXX})
\end{center}

Le d\'ep\^ot~\cite{CROTrilemmaDepot} s'organise en huit r\'epertoires, d\'ecrits au
tableau~\ref{tab:depot_structure}. Chacun contient un fichier \texttt{README} bilingue
(fran\c{c}ais et anglais) qui indique le statut de son contenu, publi\'e, pr\'eprint ou
proposition de th\`ese.

\begin{table}[htbp]
\centering
\caption{Organisation du d\'ep\^ot}
\label{tab:depot_structure}
\small
\begin{tabularx}{\textwidth}{@{}l X l@{}}
\toprule
\textbf{R\'epertoire} & \textbf{Contenu} & \textbf{Partie de la th\`ese} \\
\midrule
00\_these & Sources \LaTeX{} et PDF du manuscrit & Ensemble \\
01\_publie & Les trois articles publi\'es dans IEEE Access (P1, P2, P3) & Chap.~1 \`a 3, annexe~\ref{chap:annexe_publications} \\
02\_preprints & Fiches des neuf pr\'eprints IACR ePrint cit\'es & Chap.~1, annexe~\ref{chap:annexe_publications} \\
03\_resultats\_these & \'Enonc\'es propres \`a la th\`ese et formulations pr\'ecis\'ees par rapport aux articles & Chap.~2 \`a 4 \\
04\_sotp & Sp\'ecification de la famille SOTP, impl\'ementations de r\'ef\'erence et preuves game-based & Chap.~\ref{chap:constructions}, annexe~\ref{chap:annexe_gamebased} \\
05\_formalisation\_lean & Fichier ConditionalCRO.lean (version 2.2) et configuration Lean/Mathlib & Section~\ref{sec:lean4}, annexe~\ref{chap:annexe_lean4} \\
06\_calcul & Scripts, donn\'ees de calibration, sorties de r\'ef\'erence et tests & Section~\ref{sec:validation_retrospective} \\
07\_protocoles\_terrain & Instruments des protocoles P1 \`a P7, non d\'eploy\'es & Section~\ref{sec:protocoles_experimentaux} \\
\bottomrule
\end{tabularx}
\end{table}

Le r\'epertoire \texttt{06\_calcul/} r\'eunit les
scripts ci-dessous (\texttt{python/}), les donn\'ees de calibration (\texttt{data/}),
les sorties de r\'ef\'erence avec leurs empreintes SHA-256 (\texttt{results/}) et une
batterie de tests qui recalcule chaque valeur num\'erique cit\'ee dans ce chapitre
(\texttt{tests/}). La commande \texttt{./reproduire.sh}, lanc\'ee dans
\texttt{06\_calcul/}, r\'eg\'en\`ere l'ensemble.

\subsubsection{Environnement}
\label{subsec:environnement}

\begin{table}[htbp]
\centering
\caption{Environnement de calcul pour reproductibilit\'e}
\label{tab:environnement_calcul}
\begin{tabular}{@{}ll@{}}
\toprule
\textbf{Composante} & \textbf{Version} \\
\midrule
Python & 3.10 \`a 3.12 (r\'ef\'erence~: 3.11) \\
NumPy & $\geq 1.24$ \\
mpmath & $\geq 1.3$ \\
Lean~4 / Mathlib & v4.34.0 / v4.34.0 \\
\bottomrule
\end{tabular}
\end{table}

\subsubsection{Scripts}
\label{subsec:scripts_validation}

Les commandes suivantes s'ex\'ecutent depuis \texttt{06\_calcul/python/}.

\noindent\textbf{Calibration classique~:}
\begin{verbatim}
$ python validate_sotp.py --system all --type classical --lambda 128
\end{verbatim}

\noindent\textbf{Calibration post-quantique~:}
\begin{verbatim}
$ python validate_sotp.py --system all --type postquantum --lambda 128
\end{verbatim}

\noindent\textbf{Analyse de sensibilit\'e $\eta_q$ (sc\'enario hypoth\'etique)~:}
\begin{verbatim}
$ python sensitivity_eta_q.py --range 0.01:0.20 --steps 20
\end{verbatim}
\end{sloppypar}

\begin{boiteArbitre}
L'arbitre pose la plume. Il a vu des ordres de grandeur coh\'erents entre eux,
des intervalles dont la largeur est assum\'ee sans fard ainsi que des d\'epassements de borne
nomm\'es comme tels. Il a vu aussi des tests \'ecart\'es parce qu'ils
n'apportaient aucune information.

Il retient surtout la derni\`ere ligne, selon laquelle la validation prospective reste
\`a faire. Elle est nomm\'ee comme telle, avec des crit\`eres de r\'efutation
fix\'es \`a l'avance. L'arbitre ne confond pas calibration et preuve~; le dossier non
plus.
\end{boiteArbitre}

% ============================================================================
% CHAPITRE 3 — SECTION 3.5 : CONCLUSION
% Fichier : partie3/conclusion_validation.tex
% Aligné sur la structure visée : 3.5 Conclusion du Chapitre 3
% ============================================================================

\section{CONCLUSION DU CHAPITRE~\ref{chap:experimentation}}
\label{sec:conclusion_experimentation}
% ============================================================================
\begin{sloppypar}

Ce chapitre, qui correspond \`a l'\'etape \textbf{X} de la progression
F$\to$B$\to$E$\to$S$\to$C$\to$\textbf{X}, a trait\'e une question que les bornes du
chapitre~\ref{chap:cadre_formel} laissaient enti\`ere, celle de la valeur
d'in\'egalit\'es dont les param\`etres ne sont pas mesur\'es. Il en r\'esulte trois acquis
assortis d'une r\'eserve.

Le premier acquis est m\'ethodologique. Les six param\`etres
institutionnels deviennent des grandeurs mesurables. Chacun dispose
d\'esormais d'une d\'efinition op\'erationnelle, d'un estimateur explicite,
d'un plan d'\'echantillonnage chiffr\'e et de crit\`eres de fiabilit\'e
emprunt\'es aux standards des sciences du comportement, tels que l'accord
inter-\'evaluateurs, la fid\'elit\'e test-retest et les seuils ICC et $\kappa$. La
cryptographie institutionnelle dispose ainsi de son protocole de mesure avant
d'avoir ses mesures. Cet ordre s'impose puisqu'un r\'esultat n'est
falsifiable que si la proc\'edure de falsification est sp\'ecifi\'ee \`a
l'avance.

Le deuxi\`eme acquis est une calibration. La d\'erivation descendante
(section~\ref{sec:derivation_opposabilite}) part des espaces de verdicts, tandis
que la d\'ecomposition des capacit\'es d'extraction
(section~\ref{sec:estimation_params_classiques}) part des ressources des
acteurs. Les deux voies aboutissent \`a des ordres de grandeur compatibles, r\'eunis dans une
table unique de r\'ef\'erence (tableau~\ref{tab:parametres_reference}). Cette
compatibilit\'e ne vaut pas validation. Les deux voies ne sont pas
ind\'ependantes, puisque les valeurs de $\Delta$ ont d'abord \'et\'e fix\'ees par
inversion de la borne. Les syst\`emes d\'eploy\'es ne satisfont pas non plus toutes les
hypoth\`eses du th\'eor\`eme (Remarque~\ref{rem:portee_cas_etude}). Cette compatibilit\'e
montre seulement que les grandeurs retenues ne se contredisent pas.

Le troisi\`eme acquis, qui est n\'egatif, se r\'ev\`ele le plus instructif. Sur les
standards post-quantiques, les produits CLO$^q$ illustratifs d\'epassent la borne
(tableau~\ref{tab:decomposition_borne_pq})~; il en va de m\^eme de l'\'etude de cas
ML-DSA. Ces d\'epassements ne r\'efutent pas l'extension quantique, dont les
syst\`emes consid\'er\'es ne sont pas des instances~; ils montrent que les
param\`etres $\Delta_q$ retenus sont des postulats qui constituent la premi\`ere
cible du protocole de la section~\ref{subsec:protocole_delta_q}.

La r\'eserve, enfin, a d\'ej\`a \'et\'e \'enonc\'ee \`a la
section~\ref{subsubsec:statut_deploiement} et ne sera pas att\'enu\'ee ici.
Le programme exp\'erimental sp\'ecifi\'e n'a pas \'et\'e d\'eploy\'e.
Tant qu'il ne l'aura pas \'et\'e, les Trilemmes CRO conservent le
statut que leur \'enonc\'e leur conf\`ere, celui de r\'esultats math\'ematiquement
\'etablis sous hypoth\`eses, dont les param\`etres institutionnels restent \`a
mesurer. Le
chapitre~\ref{chap:resultats_discussions} interpr\`ete ces bornes et en tire
les implications~; il le fait dans les limites que le pr\'esent
chapitre~\ref{chap:experimentation} vient de tracer.

\end{sloppypar}

% ==========================================================
% CHAPITRE 5 : RESULTATS, DISCUSSIONS ET PERSPECTIVES
% ==========================================================

\chapter{R\'ESULTATS, DISCUSSIONS ET PERSPECTIVES}%
        %{R\'ESULTATS, DISCUSSIONS ET\\ PERSPECTIVES}
\label{chap:resultats_discussions}
\debutChapitre%
  {R\'esultats, discussions et perspectives}%
  {%
    Ce chapitre pr\'esente le bilan et l'interpr\'etation des
    contributions. Il r\'esume d'abord la calibration illustrative sur des
    syst\`emes d\'eploy\'es, classiques et post-quantiques, qui ne sont pas des
    instances du th\'eor\`eme. Il dresse ensuite le bilan formel du corpus,
    dans lequel les r\'esultats sont class\'es selon les actes du cycle canonique et selon
    leur statut (prouv\'e sous hypoth\`eses, heuristique ou conjectur\'e,
    calibr\'e). Le chapitre analyse aussi les implications th\'eoriques et
    pratiques de ces r\'esultats. Il les positionne parmi les impossibilit\'es
    fondamentales, d\'egage quatre sch\'emas structurels du compromis, \'etablit
    des seuils d'opposabilit\'e par r\'egime juridictionnel, examine les
    implications pour les standards post-quantiques et formule une hypoth\`ese de travail sur l'audit des syst\`emes de r\'emun\'eration
    de l'\'Etat. Il documente enfin
    les limitations et propose un programme de recherche futur
    structur\'e en quatre axes.%
  }%
  {%
    calibration r\'etrospective, coh\'erence interne, implications
    r\'eglementaires, questions ouvertes, standards post-quantiques,
    perspectives de recherche%
  }%

% ============================================================================
% CHAPITRE 4 — SECTION 4.1 : INTRODUCTION
% Fichier : partie4/introduction_partie3.tex

% ============================================================================

\section{INTRODUCTION}
\label{sec:introduction_resultats_discussions}
% ============================================================================
\begin{sloppypar}

Le cadre th\'eorique est maintenant pos\'e. Le chapitre~\ref{chap:contexte}
avait mis en \'evidence le paradoxe de l'authenticit\'e invisible,
formalis\'e six objectifs de recherche et construit l'outillage
th\'eorique n\'ecessaire. Les chapitres~\ref{chap:cadre_formel} et~\ref{chap:constructions} y avaient
r\'epondu par les \'etapes F, B, E, S et C de la progression
F$\to$B$\to$E$\to$S$\to$C$\to$X, en mobilisant le cadre CLO pour mesurer les compromis, le Trilemme CRO classique pour
en \'etablir les contraintes fondamentales, l'extension post-quantique
pour g\'en\'eraliser aux syst\`emes r\'esistants au quantum, la famille
SOTP comme ensemble de constructions de r\'ef\'erence dont l'optimalit\'e reste
conjecturale (Conjecture~\ref{conj:optimalite_universelle}) ainsi que la
formalisation Lean4 pour certifier le noyau alg\'ebrique. Le
chapitre~\ref{chap:experimentation} a sp\'ecifi\'e l'\'etape X.

La question qui se pose d\'esormais est diff\'erente. Les outils
th\'eoriques pr\'edisent-ils correctement le comportement des syst\`emes
r\'eels~? Les bornes d'impossibilit\'e correspondent-elles \`a des
limites observables~? Les hypoth\`eses institutionnelles sur lesquelles
repose le Trilemme, en particulier les valeurs de $\Delta$ et $\Delta_q$,
sont-elles coh\'erentes avec ce que l'on sait des capacit\'es cognitives
r\'eelles des institutions et des effets quantiques~? Ce chapitre
pr\'esente les \'el\'ements de r\'eponse disponibles, discute leurs
implications et trace les directions restant ouvertes.

Quatre axes structurent ce chapitre. Le premier est la calibration
illustrative du cadre sur des syst\`emes d\'eploy\'es, qui ne sont pas des
instances du th\'eor\`eme (Remarque~\ref{rem:portee_cas_etude}). Les ordres
de grandeur retenus y sont mutuellement coh\'erents, sans que cette
coh\'erence constitue une validation, tandis que les produits post-quantiques
illustratifs d\'epassent la borne. Le deuxi\`eme est la plausibilit\'e des
capacit\'es institutionnelles retenues, $\Delta \in [35, 73]$~bits et
$\Delta_q \in [25, 67]$~bits, au regard de la litt\'erature sur la
rationalit\'e born\'ee. Le troisi\`eme est le bilan de la formalisation
Lean4, qui comprend la compilation r\'eussie, l'hypoth\`ese de r\'egime \'etablie
comme n\'ecessaire, la comparaison chiffr\'ee entre la formulation initiale de
la borne et sa formulation certifi\'ee ainsi que la port\'ee exacte du t\'emoin de coh\'erence. Le quatri\`eme est l'expos\'e
des limitations, puisque la validation reste r\'etrospective et analytique en
attendant le programme exp\'erimental prospectif.

Deux types de validation doivent \^etre distingu\'es d\`es l'introduction. La coh\'erence interne, qui v\'erifie que les ordres de
grandeur retenus pour les syst\`emes existants ne se contredisent pas, est
ce que ce chapitre \'etablit, sans tester le cadre. La validation
exp\'erimentale prospective, qui mesurerait directement $\Delta$, $\alpha$,
$\beta$, $\Delta_q$ et $\eta_q$ aupr\`es d'acteurs institutionnels r\'eels
dans des conditions contr\^ol\'ees, constitue le programme de recherche
futur esquiss\'e \`a la section~\ref{sec:perspectives}.

La section~\ref{sec:resultats_theoriques} confronte les pr\'edictions du
cadre aux estimations analytiques et pr\'esente les r\'esultats formels de
la formalisation Lean4. La section~\ref{sec:discussions} interpr\`ete
ces r\'esultats, les positionne parmi les impossibilit\'es fondamentales
connues, examine leurs implications pour la conception et la
r\'eglementation, puis dresse le bilan des limitations et
probl\`emes ouverts. La section~\ref{sec:perspectives} trace enfin le
programme de recherche du sous-domaine naissant selon quatre axes
structur\'es.

\end{sloppypar}
% ============================================================================
% CHAPITRE 4 — SECTION 4.2 : RÉSULTATS THÉORIQUES ET FORMELS
% Fichier : partie4/resultats.tex
% Version : 2.2 — alignée sur la formalisation Lean4 révisée
% ============================================================================

\section{R\'ESULTATS TH\'EORIQUES ET FORMELS}
\label{sec:resultats_theoriques}
% ============================================================================
\begin{sloppypar}
Deux questions distinctes structurent cette section. La premi\`ere, quantitative, consiste \`a savoir si les ordres de grandeur retenus pour des syst\`emes d\'eploy\'es,
classiques et post-quantiques, sont coh\'erents entre eux et avec le cadre
CLO-CRO-PQ-SOTP. La seconde, formelle, porte sur les r\'esultats math\'ematiques
effectivement \'etablis et sur la hi\'erarchie de certitude qui les ordonne.
Ce sont deux tests ind\'ependants, auxquels correspondent deux sous-sections
distinctes.
\end{sloppypar}

\begin{boiteArbitre}
L'arbitre arrive \`a un moment charni\`ere. Jusqu'ici, on lui a donn\'e un
cadre, des hypoth\`eses et des d\'erivations~; on lui demande maintenant de
v\'erifier. Deux questions ind\'ependantes lui sont soumises. La
premi\`ere porte sur la coh\'erence mutuelle des chiffres retenus pour les syst\`emes r\'eels,
la seconde sur ce qui a \'et\'e
math\'ematiquement \'etabli et sur le degr\'e de certitude qui s'y attache. Il note qu'on
attend de lui qu'il tienne les deux ensemble, sans choisir entre elles.
\end{boiteArbitre}

% -----------------------------------------------------------------------
\subsection{Calibration r\'etrospective}
\label{subsec:validation_theorique}
% -----------------------------------------------------------------------
\begin{sloppypar}
Le cadre CLO-CRO-PQ-SOTP a \'et\'e mis en regard des estimations analytiques
disponibles pour des syst\`emes cryptographiques repr\'esentant des contextes
institutionnels distincts, \`a savoir ECDSA~\cite{Johnson01,NIST.FIPS.186-5} pour les
signatures contractuelles, Helios pour le vote \'electronique
d\'emocratique~\cite{Adida08,Adida09UCL}, les zk-SNARKs Groth16 pour la
conformit\'e financi\`ere~\cite{Groth16,Garratt15,BCTV14}, ainsi que
ML-DSA~\cite{Dilithium2022,NIST.FIPS.204} et SLH-DSA~\cite{SPHINCS2022,NIST.FIPS.205}
pour les signatures post-quantiques
NIST~\cite{NIST.PQ.Standardization,Chen2016,Alagic2022}~; ML-KEM~\cite{Kyber2022,NIST.FIPS.203},
m\'ecanisme d'encapsulation de cl\'e, n'est mentionn\'e que pour la tra\c{c}abilit\'e.
Les m\'etriques $\Conf$, $\Relver$, $\Opp$ utilis\'ees comme r\'ef\'erence ont \'et\'e
calcul\'ees par analyse param\'etrique \`a partir de sp\'ecifications publi\'ees
sans mesure aupr\`es d'acteurs institutionnels r\'eels
(sous-section~\ref{subsubsec:statut_deploiement}). Ces syst\`emes ne sont pas des
instances du Trilemme (Remarque~\ref{rem:portee_cas_etude}), de sorte que ce qui suit
rel\`eve d'une calibration illustrative.

\subsubsection{M\'ethodologie}
\label{subsubsec:methodologie_validation}

\paragraph{Param\`etres institutionnels.}
Les param\`etres ($\Delta$, $\alpha$, $\beta$, $H(\mathcal{C})$) et, pour les
syst\`emes post-quantiques, ($\Delta_q$, $\eta_q$) sont ceux de la table unique
de r\'ef\'erence (tableau~\ref{tab:parametres_reference}). $\Delta$ est rendu
plausible par la litt\'erature sur la m\'emoire de travail et la rationalit\'e
born\'ee~\cite{Miller56,Cowan2001,Simon82,Simon1955} et sur l'expertise~\cite{Ericsson95,Chi2006}~;
$\Delta_q$ est d\'ecompos\'e en $\Delta_q = \Delta + \delta_{\text{enh}} - \delta_{\text{degr}}$.
Les valeurs ont \'et\'e fix\'ees par calibration r\'etrospective
(chapitre~\ref{chap:constructions})~; elles ne constituent pas des pr\'edictions
ind\'ependantes.

\paragraph{Forme de la borne.}
Pour les syst\`emes post-quantiques, la forme de r\'ef\'erence est la forme
raffin\'ee de la Remarque~\ref{rem:forme_raffinee_pq}
(\'equation~\eqref{eq:trilemme_pq_raffine}), qui ne s'applique qu'aux protocoles
satisfaisant les hypoth\`eses du th\'eor\`eme.

\label{eq:ecart_relatif}

\subsubsection{Syst\`emes classiques}
\label{subsubsec:resultats_classiques}

Le tableau~\ref{tab:coherence_classiques} rapproche, pour les trois syst\`emes
classiques, les valeurs calcul\'ees par les heuristiques SOTP du
chapitre~\ref{chap:constructions} et les produits CLO de la
section~\ref{sec:derivation_opposabilite}.

\begin{table}[htbp]
\centering
\caption{Calibration illustrative~: heuristiques SOTP et produits CLO
(syst\`emes classiques, $\lambda = 128$)}
\label{tab:coherence_classiques}
\small
\begin{tabular}{lcccc}
\toprule
\textbf{Syst\`eme} & \textbf{Variante} & \textbf{$H(\mathcal{C})$ domaine} &
\textbf{Valeur SOTP (ch.~\ref{chap:constructions})} & \textbf{Produit CLO ($\Relver$)} \\
\midrule
ECDSA signatures & SOTP-Low & 20 & $0{,}205$ & $0{,}273$ \\
Helios Vote & SOTP-Base & 40 & $0{,}050$ & $0{,}437$ \\
zk-SNARKs conformit\'e & SOTP-High & 70 & $0{,}037$ & $0{,}572$ \\
\bottomrule
\end{tabular}
\end{table}

Les deux colonnes ne mesurent pas la m\^eme chose. Les valeurs SOTP sont des
illustrations de l'heuristique de conception de chaque variante~; les produits
CLO sont calcul\'es avec la fiabilit\'e de v\'erification, \`a l'\'echelle
op\'erationnelle. La formulation initiale de ce tableau affichait des
\og pr\'edictions\fg{} ($0{,}205$--$0{,}273$, $0{,}430$, $0{,}568$) et des \'ecarts de
$0$ \`a $1{,}6\,\%$, obtenus par ajustement des param\`etres sur les produits CLO.
La pr\'esente formulation ne les reprend pas.

\subsubsection{Syst\`emes post-quantiques}
\label{subsubsec:resultats_postquantiques}

Le tableau~\ref{tab:coherence_pqc} rappelle les param\`etres post-quantiques
retenus (tableau~\ref{tab:parametres_reference} et
section~\ref{sec:estimation_params_quantiques}).

\begin{table}[htbp]
\centering
\caption{Param\`etres post-quantiques retenus ($\lambda = 128$, calibration)}
\label{tab:coherence_pqc}
\small
\begin{tabular}{lccccc}
\toprule
\textbf{Sch\'ema} & $\boldsymbol{\eta_q}$ \textbf{(r\'eel / sc\'enario)} & $\boldsymbol{\Delta_q}$ &
\textbf{Majoration} & \textbf{Produit illustratif} & \textbf{D\'epass\'ee~?} \\
\midrule
ML-DSA-65 & $0$ / $0{,}08$ & 28 & $0{,}225$ & $0{,}35$ & oui \\
SLH-DSA & $0$ / $0{,}08$ & 25 & $0{,}201$ & $0{,}30$ & oui \\
\bottomrule
\end{tabular}
\end{table}

La formulation initiale de ce tableau retenait un autre jeu de valeurs
($\Delta_q \in [54, 97]$~bits, $\lambda \in \{192, 256\}$) et des r\'egimes
\og Q-Ready\fg{} et \og Q-Enhanced\fg{} non d\'efinis, distincts de ceux du
chapitre~\ref{chap:experimentation}. La table unique de r\'ef\'erence les remplace.
\label{eq:decomposition_delta_q_validation}

\subsubsection{Analyse comparative}
\label{subsubsec:analyse_comparative_resultats}

Le tableau~\ref{tab:analyse_cro_comparative} d\'etaille les composantes
individuelles $\Conf$, $\Relver$, $\Opp$.

\begin{table}[htbp]
\centering
\caption{Composantes CLO (calibration illustrative)}
\label{tab:analyse_cro_comparative}
\small
\begin{tabular}{lcccccc}
\toprule
\textbf{Syst\`eme} & $\Conf$ & $\Relver$ & $\Opp$ & \textbf{Produit CLO} &
$H(\mathcal{C})$ esp\`ece & $\Delta$ / $\Delta_q$ \\
 & & & & & \textbf{(bits)} & \textbf{(bits)} \\
\midrule
\multicolumn{7}{c}{\textit{Syst\`emes classiques}} \\
\midrule
ECDSA Signatures & $0{,}31$ & $0{,}98$ & $0{,}90$ & $\mathbf{0{,}273}$ & 20 & 35 \\
Helios Vote & $0{,}92$ & $0{,}95$ & $0{,}50$ & $\mathbf{0{,}437}$ & 8 & 56 \\
zk-SNARKs Conformit\'e & $0{,}89$ & $0{,}96$ & $0{,}67$ & $\mathbf{0{,}572}$ & 11 & 73 \\
\midrule
\multicolumn{7}{c}{\textit{Signatures post-quantiques (hypoth\`ese de travail)}} \\
\midrule
ML-DSA-65 & n.d. & n.d. & n.d. & $\mathbf{0{,}35}$ & 10{,}5 & 28 \\
SLH-DSA & n.d. & n.d. & n.d. & $\mathbf{0{,}30}$ & 9{,}0 & 25 \\
\bottomrule
\end{tabular}

\vspace{0.3em}
\footnotesize
Les valeurs classiques sont d\'eriv\'ees \`a la
section~\ref{sec:derivation_opposabilite}~; les $H(\mathcal{C})$ report\'es
sont ceux du contexte d'esp\`ece (voir la mise en garde de la
section~\ref{sec:methodologie_calcul}). Pour les standards
post-quantiques, la d\'ecomposition n'est pas d\'eriv\'ee faute de contexte
institutionnel document\'e.
\end{table}

Trois constats ressortent de ce tableau.

\textbf{Constat 1. Anti-corr\'elation $\Conf$--$\Opp$.} ECDSA
($\Conf = 0{,}31$, $\Opp = 0{,}90$) s'oppose \`a Helios ($\Conf = 0{,}92$,
$\Opp = 0{,}50$) et aux zk-SNARK ($\Conf = 0{,}89$, $\Opp = 0{,}67$), ce qui signifie que la
confidentialit\'e et l'opposabilit\'e varient en sens inverse sur ces trois
syst\`emes. Ce motif est compatible avec l'intuition du Trilemme, sans que trois points
calibr\'es suffisent \`a l'\'etablir. Il est analys\'e \`a la
section~\ref{subsec:quatre_patterns}.

\textbf{Constat 2. Pas de plafond commun d\'eductible.} Les produits CLO
restent inf\'erieurs \`a $0{,}6$, mais ils sont calcul\'es avec $\Relver$~; ils ne se
comparent pas \`a la borne du Trilemme, qui porte sur $\Relrec$.

\textbf{Constat 3. Les produits post-quantiques illustratifs d\'epassent la
majoration.} Ils ne sont pas pour autant sup\'erieurs aux produits classiques
($0{,}30$--$0{,}35$ contre $0{,}273$--$0{,}572$). L'analyse figure \`a la
section~\ref{sec:comparaison_produits_cro}.

\subsubsection{Port\'ee de la mise en regard}
\label{subsubsec:implications_concordance}

La mise en regard ne valide pas le cadre. Elle a une implication
conditionnelle, puisque \emph{si} les param\`etres $\Delta$, $\alpha$, $\beta$,
$\Delta_q$, $\eta_q$ d'un syst\`eme satisfaisant les hypoth\`eses du th\'eor\`eme
sont mesur\'es a priori, la forme raffin\'ee de la borne permet d'anticiper le
compromis avant d\'eploiement. Que des syst\`emes con\c{c}us ind\'ependamment
pr\'esentent un motif compatible avec le mod\`ele est un indice, que seule la
validation exp\'erimentale prospective pourra confirmer ou infirmer.

\subsubsection{Pr\'edictions consolid\'ees}
\label{subsubsec:predictions_consolidees}

Les valeurs attendues du programme exp\'erimental sont report\'ees au
tableau~\ref{tab:predictions_consolidees}. Elles sont fix\'ees avant toute
mesure~; les crit\`eres de r\'efutation associ\'es sont ceux du
tableau~\ref{tab:criteres_refutation}.

\begin{table}[htbp]
\centering
\caption{Valeurs attendues du programme exp\'erimental (\`a pr\'e-enregistrer)}
\label{tab:predictions_consolidees}
\small
\begin{tabular}{@{}p{5cm}p{3.5cm}p{2.5cm}@{}}
\toprule
\textbf{Grandeur} & \textbf{Valeur attendue} & \textbf{Contribution} \\
\midrule
\multicolumn{3}{c}{\textit{Classique (C2)}} \\
\midrule
$\hat\Delta = \hat I(V;\Sigma|M,C)$, juges g\'en\'eraux & $\approx 25$ bits & C2 \\
$\hat\Delta$, r\'egulateurs seniors & $\approx 65$ bits & C2 \\
$\Opp_{\mathrm{emp}}(\text{Helios})$ & $0{,}50 \pm 0{,}05$ & C2, C4 \\
\midrule
\multicolumn{3}{c}{\textit{Post-quantique (C3)}} \\
\midrule
$\hat\Delta_q$ (r\'egulateurs form\'es) & $[25, 67]$ bits & C3 \\
$\hat\eta_q$ (signature lue comme cha\^{\i}ne de bits) & $0$ & C3 \\
$\hat\eta_q$ (sc\'enario de support quantique simul\'e) & $[0{,}05;\,0{,}15]$ & C3 \\
\bottomrule
\end{tabular}
\end{table}

\subsubsection{Crit\`eres de r\'efutation}
\label{subsubsec:criteres_reussite}

Un th\'eor\`eme ne se \og valide\fg{} pas \`a la majorit\'e, car une seule violation
de la borne, sur un protocole qui satisfait les hypoth\`eses et avec des
param\`etres mesur\'es, r\'efute l'une des hypoth\`eses. Les crit\`eres retenus
sont donc ceux du tableau~\ref{tab:criteres_refutation}, qui portent sur
chaque hypoth\`ese (extraction born\'ee, admissibilit\'e de l'interpr\`ete,
richesse des verdicts, perte $\eta_q$) et sur la borne elle-m\^eme. Les crit\`eres
majoritaires de la formulation initiale (\og la borne n'est pas viol\'ee pour au
moins 2 des 3 standards\fg{}, \og 4 des 5 pr\'edictions \`a $\pm 20\,\%$\fg{}) ne
sont pas repris.
Si les capacit\'es mesur\'ees s'\'ecartent syst\'ematiquement des plages
attendues, le mod\`ele d'extraction born\'ee devra \^etre r\'evis\'e~; ce
r\'esultat serait lui-m\^eme instructif pour la cryptographie institutionnelle.
\end{sloppypar}

\begin{boiteArbitre}
L'arbitre referme la partie quantitative. Il y a
trouv\'e des ordres de grandeur coh\'erents entre eux, des d\'epassements de
borne expliqu\'es par le fait que les syst\`emes \'etudi\'es sortent du champ du
th\'eor\`eme ainsi que des crit\`eres de r\'efutation fix\'es \`a l'avance.

Sans constituer une preuve ni m\^eme un signe fort, ce dossier est d\'esormais pr\^et \`a
\^etre test\'e.

La seconde question l'attend maintenant, \`a savoir ce qui a \'et\'e
formellement \'etabli.
\end{boiteArbitre}

% -----------------------------------------------------------------------
\subsection{R\'esultats formels}
\label{subsec:resultats_formels}
% -----------------------------------------------------------------------
\begin{sloppypar}
Au-del\`a de la calibration, cette section dresse le bilan du
corpus th\'eorique constitu\'e aux sections~\ref{sec:clo}--\ref{sec:lean4}.
Les r\'esultats sont pr\'esent\'es selon leur \'etape dans la progression
canonique F$\to$B$\to$E$\to$S$\to$C$\to$X,
puis selon leur niveau de certitude.

\subsubsection{\'Etape F~: les th\'eor\`emes du cadre CLO}
\label{subsubsec:resultats_etape_F}

Nous rappelons les d\'efinitions \'etablies par la
section~\ref{sec:clo} et les th\'eor\`emes qui en d\'ecoulent directement.

Le \textbf{th\'eor\`eme des bornes \'el\'ementaires}~(th\'eor\`eme~\ref{thm:bornes_mesures})
\'etablit que $\Conf, \Relrec, \Opp \in [0,1]$ pour tout protocole CLO, par
construction des d\'efinitions issues de la th\'eorie de l'information.

Le \textbf{th\'eor\`eme du goulot d'\'etranglement contextuel}~(th\'eor\`eme~\ref{thm:goulot_contexte})
\'etablit $I(V;\Sigma|C) \leq \min\{H(\Sigma|C),H(V|C)\} \leq \alpha H(\mathcal{C})$, ce qui signifie que
l'information extractible sur le verdict est born\'ee par la granularit\'e des
verdicts rapport\'ee au contexte, ind\'ependamment de la puissance de calcul
de l'interpr\`ete. Si le verdict est une fonction d\'eterministe de $(\Sigma,C)$,
alors $\Opp = 1$ (Remarque~\ref{rem:verdict_deterministe}).

Le \textbf{th\'eor\`eme de composition s\'equentielle}~(th\'eor\`eme~\ref{thm:composition_sequentielle})
\'etablit les r\`egles de composition. Le post-traitement pr\'eserve la
confidentialit\'e et ne peut accro\^itre la fiabilit\'e de reconstruction,
$\Conf(\mathcal{P}_1 \circ \mathcal{P}_2) \geq \Conf(\mathcal{P}_1)$ et
$\Relrec(\mathcal{P}_1 \circ \mathcal{P}_2) \leq \Relrec(\mathcal{P}_1)$~; la publication
conjointe fait l'inverse,
$\Conf(\mathcal{P}_1 \| \mathcal{P}_2) \leq \min\{\Conf(\mathcal{P}_1),\Conf(\mathcal{P}_2)\}$ et
$\Relrec(\mathcal{P}_1 \| \mathcal{P}_2) \geq \max\{\Relrec(\mathcal{P}_1),\Relrec(\mathcal{P}_2)\}$.
Cette formulation pr\'ecise les formulations initiales de ce r\'esultat.

\subsubsection{\'Etape B~: les th\'eor\`emes du Trilemme CRO classique}
\label{subsubsec:resultats_etape_B}

Nous rappelons les bornes \'etablies par la
section~\ref{sec:cro_classique}.

Le \textbf{Trilemme CRO conditionnel}
(Th\'eor\`eme~\ref{thm:trilemme_cro_conditionnel}) est le r\'esultat
central. Dans sa forme explicite, avec $\gamma = \alpha H(C) - 1 \geq 1$ :
\begin{equation}
\Conf \cdot \Relrec \cdot \Opp \leq \frac{1}{\gamma}\left(\frac{\Delta}{\lambda} + \frac{\kappa(\Delta+1)}{\lambda^2} + \frac{\kappa^2}{\lambda^3}\right)
\label{eq:trilemme_cro_resultat_corrige}
\end{equation}
et, sous la d\'efinition admissible de l'interpr\`ete
(D\'efinition~\ref{def:interp_admissible}), la forme raffin\'ee
\eqref{eq:trilemme_cro_raffine}, qui rend compte du r\'egime \og rare\fg{} o\`u Fano
limite l'opposabilit\'e et du r\'egime \og riche\fg{} o\`u $\Delta$ la limite
(Remarque~\ref{rem:activite_delta}).
Comme $\gamma \geq 1$, on obtient la forme simplifi\'ee
\begin{equation}
\Conf \cdot \Relrec \cdot \Opp \leq \frac{\Delta}{\lambda} + \frac{\kappa(\Delta+1)}{\lambda^2} + \frac{\kappa^2}{\lambda^3}
\label{eq:trilemme_cro_simplifie_corrige}
\end{equation}
o\`u $\kappa \geq 0$ est la constante du r\'egime de pr\'eservation de la
confidentialit\'e $I(M;\Sigma|C) \leq \kappa/\lambda$
(Hypoth\`ese~\ref{ass:regime}).

\begin{remark}[Port\'ee de la borne]
\label{rem:portee_borne_corrigee}
La borne~\eqref{eq:trilemme_cro_resultat_corrige} est valide
sous l'Hypoth\`ese~\ref{ass:regime} (r\'egime de pr\'eservation
de la confidentialit\'e). Elle est structurellement identique \`a celle
v\'erifi\'ee m\'ecaniquement dans l'annexe~\ref{chap:annexe_lean4}. La borne
publi\'ee $\Delta/\lambda + O(\lambda^{-2}) + \negl(\lambda)$ s'obtient
comme cas particulier.
\end{remark}

\subsubsection{\'Etape E~: les th\'eor\`emes du Trilemme CRO post-quantique}
\label{subsubsec:resultats_etape_E}

Nous rappelons l'extension quantique \'etablie \`a la
section~\ref{sec:cro_pq}.

Le \textbf{Trilemme CRO post-quantique}
(Th\'eor\`eme~\ref{thm:trilemme_pq}) g\'en\'eralise la borne classique
\begin{equation}
\Conf^q \cdot \Relrec^q \cdot \Opp^q \leq \frac{1}{\gamma_q}\left(\frac{\Delta_q}{\lambda} + \frac{\kappa(\Delta_q+1)}{\lambda^2} + \frac{\kappa^2}{\lambda^3}\right) \leq \frac{\Delta_q + \eta_q H(C)}{\lambda} + O(\lambda^{-2}) + \negl(\lambda)
\label{eq:trilemme_pq_resultat_corrige}
\end{equation}
o\`u $\Delta_q = \Delta + \delta_{\text{enh}} - \delta_{\text{degr}}$
avec $\delta_{\text{enh}} = 0$ pour la technologie actuelle (voir
section~\ref{sec:cro_pq}) tandis que $\eta_q = 0$ pour toute preuve qui est une
cha\^{\i}ne de bits classique.

\begin{remark}[Port\'ee de la borne post-quantique]
\label{rem:portee_borne_pq_corrigee}
Le terme $\eta_q H(C)/\lambda$ n'appara\^it pas dans la borne explicite.
Il est d'ordre $\lambda^{-1}$, de sorte qu'aucun terme en
$O(\lambda^{-2})$ ne l'absorbe~; il est simplement absent de cette borne. Il figure uniquement dans
la majoration affaiblie, obtenue en ajoutant une quantit\'e positive,
o\`u il rend visible le co\^ut s\'emantique du canal. La borne principale
est donc structurellement identique \`a la borne classique, avec
$\Delta_q$ \`a la place de $\Delta$, sous
l'Hypoth\`ese~\ref{ass:regime_quantique} (r\'egime de pr\'eservation de
la confidentialit\'e quantique).
\end{remark}

\subsubsection{\'Etape S~: les constructions SOTP}
\label{subsubsec:resultats_etape_S}

Nous \'etablissons les propri\'et\'es cryptographiques des trois
constructions, \`a savoir la confidentialit\'e calculatoire sous
PRG~\cite{GoldreichKrawczyk99,BlumMicali84} et masquage parfait
de l'engagement ainsi que la fiabilit\'e de v\'erification sous DLog~\cite{Pedersen91} et
q-PKE~\cite{Groth16,Boneh2008} (Propositions du chapitre~\ref{chap:constructions}
et annexe~\ref{chap:annexe_gamebased}). Les minorations d'opposabilit\'e et les
fractions de la borne atteintes (\og $75\,\%$, $50\,\%$, $45\,\%$\fg{} dans la
formulation initiale) rel\`event d'heuristiques de conception, car elles supposent une extraction du verdict que la preuve seule ne
garantit pas (Remarques~\ref{rem:facteur_limitant_low},
\ref{rem:heuristique_sotp_base} et~\ref{rem:heuristique_sotp_high}).
\label{eq:saturation_resultat}

L'optimalit\'e universelle de la famille SOTP est formul\'ee comme conjecture
(Conjecture~\ref{conj:optimalite_universelle}). L'argument initial reposait sur
ces heuristiques et ne distinguait pas fiabilit\'e de reconstruction et fiabilit\'e
de v\'erification (Remarque~\ref{rem:statut_optimalite}).
\label{eq:optimalite_resultat}

\subsubsection{\'Etape C~: la formalisation Lean4}
\label{subsubsec:resultats_etape_C}

Nous v\'erifions m\'ecaniquement le noyau alg\'ebrique du Trilemme
CRO~(section~\ref{sec:lean4}).

La formalisation Lean4 couvre l'int\'egralit\'e du noyau alg\'ebrique du
Trilemme, organis\'e en neuf sections, suivies d'un bloc de v\'erification des axiomes~:
\begin{itemize}
    \item La constante de Pinsker $c_0 = 1/(2\ln 2)$ et ses propri\'et\'es
    \item La n\'egligeabilit\'e de $2^{-\lambda}$ via les outils de Mathlib
    \item Les param\`etres institutionnels et $\gamma = \alpha H(C) - 1$
    \item La structure \texttt{ProtocolInstance} avec les propri\'et\'es
          informationnelles (Pinsker, Fano) int\'egr\'ees
    \item Les hypoth\`eses institutionnelles 1, 2$'$ et 3
    \item Le Lemme 4 (garantie de confidentialit\'e) et le
          Th\'eor\`eme 1 en cinq \'etapes
    \item Le th\'eor\`eme principal asymptotique et les t\'emoins de
          coh\'erence et de port\'ee
    \item La classification par r\'egimes institutionnels
    \item Les extensions (section~\ref{subsec:lean4_extensions}), \`a savoir le verdict
          d\'eterministe ($\Opp = 1$), la forme raffin\'ee de la borne et sa
          co\"{\i}ncidence avec Fano lorsque $\Delta \geq \gamma$
\end{itemize}

La formalisation \'etablit la borne
\begin{equation}
\Conf \cdot \Relrec \cdot \Opp \leq \frac{\Delta}{\lambda}
+ \frac{\kappa(\Delta+1)}{\lambda^2} + \frac{\kappa^2}{\lambda^3} + \varepsilon(\lambda)
\end{equation}
enti\`erement v\'erifi\'ee sans \texttt{sorry} ni axiome d\'eclar\'e (seuls les
axiomes standard \texttt{propext}, \texttt{Classical.choice} et
\texttt{Quot.sound} apparaissent). Fano et Pinsker, absents de Mathlib, sont
pos\'es comme propri\'et\'es des instances de protocole~; l'extraction born\'ee et
le r\'egime de confidentialit\'e sont des hypoth\`eses du th\'eor\`eme. La
certification porte donc sur la d\'erivation~; les pr\'emisses restent des hypoth\`eses.

Le p\'erim\`etre couvert garantit la correction du noyau alg\'ebrique
central. Les constructions SOTP compl\`etes et l'extension post-quantique
constituent des directions de recherche future. Les formulations
alternatives issues de la certification sont pr\'esent\'ees \`a la
sous-section~\ref{subsec:lean4_erreurs}.

\subsubsection{Coh\'erence et port\'ee du th\'eor\`eme v\'erifi\'e}
\label{subsubsec:coherence_portee}

La formalisation Lean4 produit deux t\'emoins qui d\'emontrent la
non-vacuit\'e et la n\'ecessit\'e des hypoth\`eses du th\'eor\`eme.

\paragraph{T\'emoin de coh\'erence}

Une instance du mod\`ele abstrait v\'erifie simultan\'ement toutes les
hypoth\`eses avec $\Conf = \Opp = 1$, $\Rel = 1/\lambda$, $I = 0$. Pour cette
instance~:
\begin{equation}
\Conf \cdot \Rel \cdot \Opp = \frac{1}{\lambda}
= \frac{\Delta}{\gamma \lambda}
\end{equation}
Cela d\'emontre que les hypoth\`eses du mod\`ele formel sont compatibles et que
la borne y est atteinte. Ce t\'emoin n'est pas r\'ealisable par un protocole, car
avec $I = 0$ Fano impose $\Relrec = 2^{-\lambda}$ au lieu de $1/\lambda$. Le mod\`ele
formel utilise la forme rel\^ach\'ee de Fano, $I \geq \Rel\cdot\lambda - 1$, qui
autorise ce point (Remarque~\ref{rem:portee_temoin}).

\paragraph{T\'emoin de port\'ee}

Une instance v\'erifie les Hypoth\`eses 1, 2$'$ et 3 avec $\Delta = 0$
mais s'\'ecarte du r\'egime de confidentialit\'e ($I = \lambda/2 - 1$)~; elle
pr\'esente :
\begin{equation}
\Conf \cdot \Rel \cdot \Opp = \frac{1}{4}
\end{equation}
constant pour tout $\lambda \geq 64$. Cela d\'emontre que le r\'egime
$I \leq \kappa/\lambda$ est n\'ecessaire \`a la borne. Le cas limite
complet est document\'e \`a la sous-section~\ref{subsec:lean4_erreurs}.

\paragraph{Hi\'erarchie de certitude}

Le tableau~\ref{tab:hierarchie_certitude} synth\'etise les r\'esultats
formels class\'es selon leur niveau de certitude.

\begin{table}[htbp]
\centering
\caption{Bilan des r\'esultats formels et couverture Lean4}
\label{tab:hierarchie_certitude}
\small
\begin{tabular}{p{3.2cm}p{3cm}cp{2.8cm}c}
\toprule
\textbf{Cat\'egorie} & \textbf{Niveau de certitude} & \textbf{Nb} &
\textbf{Statut de v\'erification} & \textbf{Couvert. Lean4} \\
\midrule
Trilemme (forme raffin\'ee), Lemme~\ref{lem:confidentiality_guarantee}, Proposition~\ref{prop:fondement_extraction} & Prouv\'es sous Hypoth\`eses~\ref{ass:extraction_bornee}--\ref{ass:regime} & 3 & Preuves papier & \textbf{Oui} (noyau, annexe~\ref{chap:annexe_lean4}) \\
Goulot, composition, bornes \'el\'ementaires (Sec.~\ref{sec:clo}) & Prouv\'es (formulations pr\'ecis\'ees) & 3 & Preuves papier & Non (verdict d\'eterministe~: oui) \\
Trilemme PQ (Sec.~\ref{sec:cro_pq}) & Sous Hypoth\`eses~\ref{ass:extraction_quantique}--\ref{ass:regime_quantique} & 1 & Preuve papier & Non \\
Propri\'et\'es SOTP ($\Conf$, $\Relver$) & Prouv\'ees sous PRG et DLog (q-PKE pour High) & 3 & R\'eductions (annexe~\ref{chap:annexe_gamebased}) & Non \\
Minorations d'opposabilit\'e SOTP, optimalit\'e & Heuristiques / conjecture & --- & Non prouv\'ees & Non \\
Valeurs num\'eriques des syst\`emes d\'eploy\'es & Calibration illustrative & --- & Non mesur\'ees & --- \\
\bottomrule
\end{tabular}
\end{table}

\begin{remark}[P\'erim\`etre de la couverture Lean4]
\label{rem:perimetre_lean4}
La formalisation Lean4 (annexe~\ref{chap:annexe_lean4}) couvre le
noyau alg\'ebrique du Trilemme CRO classique, y compris les
t\'emoins de coh\'erence et de port\'ee. Les constructions SOTP et
l'extension post-quantique constituent des directions de recherche
futures. La colonne
\og Couverture Lean4 \fg{} rend ce p\'erim\`etre explicite.
\end{remark}

Cette hi\'erarchie, prouv\'e sous hypoth\`eses / heuristique ou conjectur\'e /
calibr\'e, vise la transparence \'epist\'emique que recommandent les
promoteurs de la preuve formelle~\cite{Wiedijk2008,Avigad2014}, puisque chaque
\'enonc\'e du manuscrit y est rattach\'e \`a un statut explicite.

\subsubsection{Analyse de la contribution du terme $\eta_q \cdot H(C)$}
\label{subsubsec:analyse_contribution_eta_q}

La borne post-quantique $(\Delta_q + \eta_q \cdot H(C))/\lambda$ agr\`ege
deux termes d'ordres de grandeur tr\`es diff\'erents. Le
tableau~\ref{tab:decomposition_borne_pq} en donne la d\'ecomposition
compl\`ete et la confronte, pour chaque standard NIST~PQC, au produit
illustratif.

Pour les signatures post-quantiques r\'eelles, $\eta_q = 0$ et ce terme
dispara\^it. Dans le sc\'enario hypoth\'etique d'un support quantique
($\eta_q \approx 0{,}08$), il ne repr\'esente que $2$ \`a $5\,\%$ de la
majoration, de sorte que la capacit\'e d'extraction reste le terme dominant. Les produits
illustratifs d\'epassent la majoration~; ces syst\`emes ne sont pas des
instances du th\'eor\`eme (section~\ref{sec:comparaison_produits_cro}).

\begin{remark}[Clarification sur les valeurs num\'eriques des tableaux]
\label{rem:clarification_section_VI_B}
Comme le pr\'ecise la Remarque~\ref{rem:portee_cas_etude}, les produits CLO des syst\`emes
d\'eploy\'es sont calcul\'es avec $\Relver$. Les capacit\'es $\Delta$ qui en sont
d\'eduites rel\`event d'une calibration illustrative~; le
Th\'eor\`eme~\ref{thm:trilemme_cro_conditionnel} ne les implique pas.
\end{remark}
\end{sloppypar}

\begin{boiteArbitre}
L'arbitre retient trois choses. Les chiffres qu'on lui montre sont des
calibrations, ce qu'on prend soin de lui dire. Sa capacit\'e $\Delta$ d\'epasse
d\'esormais le statut de m\'etaphore puisqu'elle dispose d'un estimateur, $\hat I(V;\Sigma|M,C)$,
ainsi que de crit\`eres de r\'efutation. Enfin, la rigueur formelle a certifi\'e la
d\'erivation de la borne
$\gamma^{-1}(\Delta/\lambda + \kappa(\Delta+1)/\lambda^2 + \kappa^2/\lambda^3)$ en
disant exactement ce qu'elle suppose.
Cette explicitation des hypoth\`eses rel\`eve \`a ses yeux de la d\'emarche scientifique elle-m\^eme.
\end{boiteArbitre}
% ============================================================================
% CHAPITRE 4 — SECTION 4.3 : DISCUSSIONS
% Fichier : partie4/discussions.tex
% Version : 2.2 — alignée sur la formalisation Lean4 révisée
% ============================================================================

\section{DISCUSSIONS}
\label{sec:discussions}
% ============================================================================
\begin{sloppypar}
La
section~\ref{sec:resultats_theoriques} a mis le cadre CLO-CRO-PQ-SOTP en
regard des estimations analytiques disponibles pour cinq syst\`emes d\'eploy\'es
(trois classiques et deux signatures post-quantiques), dans une calibration
illustrative. Elle a dress\'e le bilan d'un corpus th\'eorique articul\'e selon la
progression F$\to$B$\to$E$\to$S$\to$C$\to$X.
\'Etablir ces faits ne suffit toutefois pas \`a en fixer la port\'ee. La pr\'esente section
examine ce qu'ils signifient au-del\`a de la cryptographie institutionnelle consid\'er\'ee en vase
clos, c'est-\`a-dire pour l'avancement de la cryptographie comme discipline, pour les
praticiens qui con\c{c}oivent des syst\`emes et pour les probl\`emes ouverts
qui restent non r\'esolus.

L'analyse proc\`ede en onze temps. Elle examine d'abord ce que la
formalisation a r\'ev\'el\'e sur le Trilemme~CRO
(sous-section~\ref{subsec:ce_que_formalisation_revele}), puis positionne nos
r\'esultats parmi les impossibilit\'es fondamentales en pr\'esentant la
classification \'epist\'emologique des contributions
(sous-section~\ref{subsec:positionnement_impossibilites}). Elle d\'etaille ensuite l'apport
sp\'ecifique de l'extension post-quantique
(sous-section~\ref{subsec:apport_extension_quantique}), d\'egage les quatre sch\'emas
structurels du compromis CLO (sous-section~\ref{subsec:quatre_patterns}) et mesure la sensibilit\'e de
l'opposabilit\'e \`a la capacit\'e d'extraction
(sous-section~\ref{subsec:analyse_sensibilite}). Les implications pour
la conception et la r\'eglementation (sous-section~\ref{subsec:implications_conception})
pr\'ec\`edent la situation de la contribution
vis-\`a-vis des travaux existants (sous-section~\ref{subsec:positionnement_travaux}) ainsi qu'un bilan
des limitations et probl\`emes ouverts
(sous-section~\ref{subsec:limitations_problemes_ouverts}). Les trois derni\`eres sous-sections montrent comment la
famille SOTP transforme l'impossibilit\'e en guide
(sous-section~\ref{subsec:portee_constructive}), tirent la le\c{c}on
m\'ethodologique de la certification
(sous-section~\ref{subsec:verification_corrige_lecon}) et clarifient le
statut \'epist\'emique du programme exp\'erimental
(sous-section~\ref{subsec:statut_protocoles_experimentaux}).
\end{sloppypar}

\begin{boiteArbitre}
L'arbitre dispose maintenant des faits, \`a savoir des bornes prouv\'ees sous hypoth\`eses,
des constructions dont les propri\'et\'es cryptographiques sont \'etablies mais
dont l'opposabilit\'e reste heuristique ainsi qu'une formalisation qui a mis au jour
des lacunes du raisonnement initial. Ce
chapitre lui demande ce que tout cela \emph{signifie} pour la discipline
naissante, pour les praticiens qui choisissent des protocoles et pour les
questions qui restent ouvertes. Une impossibilit\'e prouv\'ee constitue \`a ses yeux une
forme nette de connaissance, puisqu'elle lui apprend ce qu'il ne pourra jamais obtenir
et pour quelle raison.
\end{boiteArbitre}

% -----------------------------------------------------------------------
\subsection{Ce que la formalisation a r\'ev\'el\'e}
\label{subsec:ce_que_formalisation_revele}
% -----------------------------------------------------------------------
\begin{sloppypar}
La v\'erification m\'ecanique du Trilemme~CRO dans Lean4 a produit,
au-del\`a de la certification du c\oe ur alg\'ebrique, cinq r\'ev\'elations
\'epist\'emologiques que le raisonnement informel n'avait pas permis
d'identifier.

\subsubsection{Hypoth\`ese explicite des propri\'et\'es d'instance}
\label{subsubsec:hypothese_quasi_optimalite}

La preuve originale~\cite{MINKA2026CRO} utilise les propri\'et\'es
informationnelles (Fano, Pinsker) sans les expliciter comme conditions
v\'erifiables du th\'eor\`eme. La formalisation Lean4 les rend
explicites en les int\'egrant comme champs de la structure
\texttt{ProtocolInstance}~:
\begin{equation}
\texttt{pinsker} : \TV^2 \leq c_0 \cdot \KL, \qquad
\texttt{fano} : I \geq \Relrec \cdot \lambda - 1.
\label{eq:hypothese_quasi_optimalite}
\end{equation}
Cette explicitation s'aligne avec la rationalit\'e
born\'ee~\cite{Simon82,Miller56,Ericsson95}, puisque les institutions ne peuvent
extraire qu'une information limit\'ee d'une preuve cryptographique, ce qui
force les protocoles \`a op\'erer dans des r\'egimes de faible information
relativement \`a $\lambda$. La formalisation Lean4 transforme des
intuitions diffuses en conditions v\'erifiables.

\subsubsection{Structure en trois \'etapes de la preuve formalis\'ee}
\label{subsubsec:non_derivabilite_borne}

La formalisation \'etablit la borne
$\gamma^{-1}(\Delta/\lambda + \kappa(\Delta+1)/\lambda^2 + \kappa^2/\lambda^3)$
en trois \'etapes principales. (1)~Fano borne $\Relrec$~; (2)~l'extraction
born\'ee et le verdict efficace bornent $\Opp$~; (3)~on forme le produit
puis on applique le r\'egime de confidentialit\'e
$I(M;\Sigma|C) \leq \kappa/\lambda$. Cette structure, plus simple que
l'optimisation sur $[0,\lambda]$, pr\'eserve la borne dominante
$\Delta/\lambda$ et la propri\'et\'e qualitative centrale du trilemme.
La formalisation compl\`ete de l'optimisation analytique sur
$[0,\lambda]$ n\'ecessiterait des outils d'analyse r\'eelle non encore
disponibles dans Mathlib~\cite{MathlibAnalysis,Tsybakov2009,Scerri14}.

\subsubsection{N\'egligeabilit\'e~: une preuve par limite}
\label{subsubsec:necessite_analyse_cas}

L'\'enonc\'e informel \og $(1/2)^\lambda$ est n\'egligeable\fg{} est tenu
pour \'evident dans la litt\'erature
cryptographique~\cite{Goldreich01,KatzLindell14}. Dans la formalisation, la
n\'egligeabilit\'e est d\'efinie par un filtre (\og pour tout $c$, \`a partir d'un
certain rang, $\varepsilon(\lambda)\,\lambda^c \leq 1$\fg)~; la preuve passe par
la limite $\lambda^c/2^\lambda \to 0$ fournie par Mathlib. Elle n'est pas
constructive, puisqu'elle ne produit pas de seuil explicite. La formulation initiale
de ce paragraphe d\'ecrivait une analyse de cas ($c \le 4$ / $c > 4$) et un seuil
$N(c) = \lceil 2c \rceil + 4$~; la pr\'esente formulation suit le code compil\'e. Le ph\'enom\`ene relev\'e par Scerri~\cite{Scerri14} demeure, car
formaliser un fait asymptotique \og \'evident\fg{} oblige \`a choisir une
d\'efinition pr\'ecise.

\subsubsection{Transparence des constantes implicites}
\label{subsubsec:transparence_constantes}

La notation $O(\lambda^{-2})$ dissimule une d\'ependance sur
$1/\gamma$ o\`u $\gamma = \alpha H(C) - 1$. Notre borne explicite
$\gamma^{-1}(\Delta/\lambda + \kappa(\Delta+1)/\lambda^2 +
\kappa^2/\lambda^3)$ expose cette d\'ependance, puisque la constante
d\'epend \`a la fois de $\Delta$ (capacit\'e d'extraction
institutionnelle) et de $\gamma$ (granularit\'e des verdicts
efficace). Pour la cryptographie
institutionnelle, o\`u les param\`etres doivent \^etre empiriquement
mesurables, cette transparence constitue une
exigence op\'erationnelle.

\subsubsection{Les param\`etres institutionnels sont des constantes du th\'eor\`eme}
\label{subsubsec:parametres_constantes}

La formalisation \'etablit que $\Delta$, $\alpha$, $\beta$, $H(C)$ ont le
statut de constantes
du th\'eor\`eme, distinct de celui de param\`etres libres, de sorte que leur mesure empirique devient un
pr\'erequis pour son application. Un syst\`eme d\'eployant un protocole
sans mesurer pr\'ealablement sa capacit\'e d'extraction $\Delta$ ne peut
pas \'evaluer si la borne th\'eorique est satisfaite~\cite{Grossklags10,Herley09}.

\subsubsection{Coh\'erence et port\'ee du th\'eor\`eme principal}
\label{subsubsec:consequences_empiriques_borne}

La formalisation Lean4 produit deux t\'emoins qui d\'emontrent la
non-vacuit\'e et la n\'ecessit\'e des hypoth\`eses. Le \textbf{t\'emoin
de coh\'erence} est une instance qui v\'erifie toutes les hypoth\`eses
avec $\Conf = \Opp = 1$, $\Rel = 1/\lambda$, $I = 0$~; pour cette
instance, $\Conf \cdot \Rel \cdot \Opp = 1/\lambda$, ce qui \'egale
exactement la borne du th\'eor\`eme. Il s'agit d'un point du mod\`ele abstrait,
qu'aucun protocole ne r\'ealise (avec $I = 0$, $\Relrec = 2^{-\lambda}$~;
Remarque~\ref{rem:portee_temoin}). Le \textbf{t\'emoin de port\'ee}
est une instance qui v\'erifie les Hypoth\`eses 1, 2$'$ et 3 avec
$\Delta = 0$ mais s'\'ecarte du r\'egime de confidentialit\'e
($I = \lambda/2 - 1$)~; elle pr\'esente $\Conf \cdot \Rel \cdot \Opp = 1/4$
constant. Ces deux t\'emoins \'etablissent que les hypoth\`eses sont
satisfiables et que le r\'egime est n\'ecessaire \`a la borne.
\end{sloppypar}

\begin{boiteArbitre}
L'arbitre s'arr\^ete sur ce que la machine a trouv\'e, soit cinq points que le
raisonnement \`a la main n'avait pas vus. Des propri\'et\'es
informationnelles y deviennent des conditions v\'erifiables, une
structure de preuve pr\'eserve la borne dominante et un
fait \og \'evident \fg{} exige une d\'efinition pr\'ecise. S'y ajoutent des constantes que
la notation $O(\cdot)$ cachait ainsi que des param\`etres institutionnels qui
se r\'ev\`elent \^etre des constantes du th\'eor\`eme, distinctes de variables
libres.

Il note que la v\'erification formelle a fait davantage que confirmer le
r\'esultat, puisqu'elle l'a approfondi, ce qu'il tient pour un gain r\'eel.
\end{boiteArbitre}

% -----------------------------------------------------------------------
\subsection{Positionnement parmi les impossibilit\'es fondamentales}
\label{subsec:positionnement_impossibilites}
% -----------------------------------------------------------------------
\begin{sloppypar}

\subsubsection{La cryptographie institutionnelle comme domaine propre}
\label{subsubsec:crypto_institutionnelle_domaine}

La m\'etrique $\Opp = I(V;\Sigma|C)/H(V|C)$ est absente de tout cadre
cryptographique ant\'erieur. Le calcul multi-parties
s\'ecuris\'e~\cite{Yao82,GMW87,Goldreich1987,Evans2018} optimise le
nombre de rounds et la confidentialit\'e du calcul distribu\'e~; le
chiffrement homomorphe~\cite{Gentry09,Vaikuntanathan2012,Brakerski2014}
optimise la profondeur de circuits sur donn\'ees chiffr\'ees~; les
syst\`emes zero-knowledge~\cite{GMR85,Groth16,BCTV14} minimisent la
fuite d'information tout en prouvant l'appartenance. Aucune de ces
sp\'ecialisations ne pose la question de savoir \emph{dans quelle mesure un
acteur institutionnel born\'e cognitivement peut extraire une
information utile de cette preuve pour rendre un verdict
inform\'e}~\cite{Simon82,Miller56,Kahneman11}.

Les travaux fondateurs de Goldwasser-Micali-Rackoff~\cite{GMR85} ont
axiomatis\'e la puissance computationnelle de
l'adversaire~\cite{Goldreich01,Bellare97}~; nous axiomatisons ici la
capacit\'e interpr\'etative de l'interpr\`ete. Il s'agit du m\^eme objet, un acteur
born\'e, consid\'er\'e selon une sym\'etrie invers\'ee.

\begin{table}[htbp]
\centering
\caption{Positionnement de la cryptographie institutionnelle parmi les
sp\'ecialisations cryptographiques}
\label{tab:comparaison_domaines_crypto}
\begin{tabular}{@{}p{3.5cm}p{3.5cm}p{4.5cm}@{}}
\toprule
\textbf{Domaine} & \textbf{Optimise pour} & \textbf{M\'etrique centrale} \\
\midrule
Cryptographie classique   & R\'esistance adversaires
  & Avantage distingueur~\cite{Goldreich01} \\
MPC                       & Calcul distribu\'e priv\'e
  & Rounds, parties~\cite{Yao82,Evans2018} \\
FHE                       & Calcul sur donn\'ees chiffr\'ees
  & Profondeur circuits~\cite{Gentry09,Brakerski2014} \\
ZKP                       & Preuves sans r\'ev\'elation
  & Taille preuve, compl\'etude~\cite{GMR85,Groth16} \\
\textbf{Crypto institutionnelle}   & \textbf{Interpr\'etabilit\'e l\'egale}
  & \textbf{Opposabilit\'e $\Opp$}~\cite{MINKA2026CLO,MINKA2026CRO} \\
\textbf{Crypto PQ institutionnelle} & \textbf{R\'esilience quantique}
  & \textbf{$\Delta_q$, $\eta_q$}~\cite{MINKA2026PQCRO} \\
\bottomrule
\end{tabular}
\end{table}

\subsubsection{Le Trilemme CRO parmi les impossibilit\'es fondamentales}
\label{subsubsec:trilemme_parmi_impossibilites}

Le Trilemme CRO conditionnel joue dans la cryptographie institutionnelle
le m\^eme r\^ole structurant que les th\'eor\`emes d'impossibilit\'e
dans d'autres branches de l'informatique
th\'eorique~\cite{Abadi2012,Attiya2004}. Le th\'eor\`eme
FLP~\cite{FLP85} \'etablit qu'aucun protocole d\'eterministe
ne garantit le consensus en pr\'esence d'une seule d\'efaillance
asynchrone~; le th\'eor\`eme CAP~\cite{Brewer00,GilbertLynch02} \'etablit
qu'un syst\`eme distribu\'e ne peut simultan\'ement garantir coh\'erence
stricte, disponibilit\'e totale et tol\'erance aux partitions~; le
Trilemme CRO \'etablit, sous ses hypoth\`eses, que le produit $\Conf\cdot\Relrec\cdot\Opp$ est
born\'e par
$\gamma^{-1}(\Delta/\lambda + \kappa(\Delta+1)/\lambda^2 + \kappa^2/\lambda^3)$
(classique) ou par son analogue quantique
$\gamma_q^{-1}(\Delta_q/\lambda + \kappa(\Delta_q+1)/\lambda^2 + \kappa^2/\lambda^3)$
(\cite{MINKA2026CRO} et section~\ref{sec:cro_pq}).

\begin{table}[htbp]
\centering
\caption{Comparaison des th\'eor\`emes d'impossibilit\'e fondamentaux}
\label{tab:comparaison_impossibilites}
\begin{tabular}{@{}p{2.8cm}p{4cm}p{4.5cm}@{}}
\toprule
\textbf{Th\'eor\`eme} & \textbf{Contrainte fondamentale}
  & \textbf{Hypoth\`eses} \\
\midrule
FLP & Consensus impossible (asynchrone + 1 d\'efaillance)
  & Mod\`ele asynchrone~\cite{FLP85,Attiya2004} \\
CAP & C\,+\,A\,+\,P impossibles simultan\'ement
  & Syst\`eme distribu\'e~\cite{Brewer00,GilbertLynch02} \\
Blockchain Trilemma & D\'ecentralisation + S\'ecurit\'e + Scalabilit\'e
  & Blockchain~\cite{Buterin2014,Abadi2012} \\
\textbf{CRO classique} & $\Conf\cdot\Relrec\cdot\Opp \le \gamma^{-1}(\Delta/\lambda + \kappa(\Delta+1)/\lambda^2 + \kappa^2/\lambda^3)$
  & \textbf{$\Delta,\alpha,\beta$ institutionnels}~\cite{MINKA2026CRO} \\
\textbf{CRO post-quantique} &
  $\Conf^q\!\cdot\!\Relrec^q\!\cdot\!\Opp^q \le
   \gamma_q^{-1}(\Delta_q/\lambda + \kappa(\Delta_q+1)/\lambda^2 + \kappa^2/\lambda^3)$
  & \textbf{$\Delta_q$, $\eta_q$ quantiques}~\cite{MINKA2026PQCRO} \\
\bottomrule
\end{tabular}
\end{table}

Une distinction s'impose toutefois, car FLP et CAP ont \'et\'e prouv\'es dans des
mod\`eles de calcul formels bien
\'etablis alors que le Trilemme CRO est conditionnel
aux hypoth\`eses sur $\Delta$, $\alpha$, $\beta$, $\Delta_q$,
$\eta_q$, param\`etres dont les valeurs n'ont pas encore \'et\'e
valid\'ees par des protocoles exp\'erimentaux contr\^ol\'es sur des
acteurs institutionnels r\'eels. Le Trilemme
CRO est un th\'eor\`eme correct sous ses hypoth\`eses~; sa pertinence
empirique reste \`a \'etablir.

Sous cette r\'eserve, le d\'eplacement de perspective est net. Que
$\Delta \in [35, 73]$~bits s'av\`ere plus restrictif que
$\lambda \in [128, 256]$~bits signifie que la contrainte dominante sur
les syst\`emes cryptographiques institutionnels tient \`a la capacit\'e
interpr\'etative de l'institution, probl\`eme largement ignor\'e
jusqu'ici. La puissance de calcul de l'adversaire, que la cryptographie sait
traiter, passe au second plan.

\subsubsection{Classification \'epist\'emologique des r\'esultats}
\label{subsubsec:classification_epistemologique}

Les r\'esultats de cette th\`ese, qui ne sont pas de m\^eme nature, doivent
\^etre distingu\'es selon leur niveau de
certitude.

\begin{table}[htbp]
\centering
\caption{Hi\'erarchie de certitude des r\'esultats de la th\`ese}
\label{tab:classification_epistemologique}
\begin{tabular}{@{}p{3.8cm}p{3.2cm}p{4.5cm}@{}}
\toprule
\textbf{R\'esultat} & \textbf{Niveau de certitude} & \textbf{Statut} \\
\midrule
R\'esultats CLO (bornes, goulot, composition) & Prouv\'es (formulations pr\'ecis\'ees)
  & Preuves papier~(section~\ref{sec:clo}) \\
Trilemme classique (forme raffin\'ee) & Sous hypoth\`eses institutionnelles
  & Noyau v\'erifi\'e Lean4~(\cite{MINKA2026CRO} et section~\ref{sec:lean4}) \\
Trilemme PQ & Sous hypoth\`eses quantiques
  & Preuve papier~(\cite{MINKA2026PQCRO} et section~\ref{sec:cro_pq}) \\
Propri\'et\'es SOTP ($\Conf$, $\Relver$) & Sous PRG, DLog, q-PKE
  & R\'eductions~(annexe~\ref{chap:annexe_gamebased}) \\
Opposabilit\'e SOTP, optimalit\'e & Heuristique / conjecture
  & Non prouv\'ees~(Conjecture~\ref{conj:optimalite_universelle}) \\
Formalisation Lean4 & M\'ecaniquement certifi\'ee
  & D\'erivation v\'erifi\'ee, pr\'emisses pos\'ees~(section~\ref{sec:lean4}) \\
\midrule
\multicolumn{3}{l}{\textit{Falsifiables exp\'erimentalement~:}} \\
Valeurs de $\Delta$, $\Delta_q$, $\eta_q$ & Empirique
  & \`A valider par \'etudes contr\^ol\'ees~(chapitre~\ref{chap:experimentation}) \\
\bottomrule
\end{tabular}
\end{table}

Distinguer explicitement ces niveaux est une contribution
m\'ethodologique propre~\cite{Nosek2015,IOC2019}, puisqu'elle prot\`ege contre
la sur\'evaluation des r\'esultats conditionnels et ouvre pr\'ecis\'ement
le programme exp\'erimental futur. La fronti\`ere entre \og th\'eor\`eme
correct sous hypoth\`eses \fg{} et \og fait empirique \'etabli \fg{} doit
\^etre maintenue avec rigueur~\cite{Lakatos1978}.
\end{sloppypar}

% -----------------------------------------------------------------------
\subsection{L'apport de l'extension quantique}
\label{subsec:apport_extension_quantique}
% -----------------------------------------------------------------------
\begin{sloppypar}
L'extension quantique (section~\ref{sec:cro_pq}) introduit deux
innovations th\'eoriques sans analogue dans le r\'egime classique.

\subsubsection{P\'enalit\'e s\'emantique quantique}
\label{subsubsec:penalite_semantique}

Le terme $\eta_q \cdot H(C)/\lambda$ repr\'esente une contrainte
fondamentalement nouvelle, puisqu'il chiffre la part de la preuve qui, sous
bruit de canal quantique, ne peut plus \^etre interpr\'et\'ee. Il
n'intervient pas dans la borne explicite. Il est aussi nul pour toute preuve
qui est une cha\^{\i}ne de bits classique, y compris les signatures ML-DSA et
SLH-DSA. Il ne devient pertinent que dans le sc\'enario, hypoth\'etique \`a ce
jour, d'un support ou d'un v\'erificateur quantique, o\`u il fournit une
pr\'ediction mesurable.
Cette p\'enalit\'e \'emerge de la th\'eorie de l'information quantique
(borne de Holevo~\cite{Holevo1973}, lemme de la mesure
douce~\cite{Winter1999,Ozawa2000}).

\subsubsection{D\'ecomposition de la capacit\'e quantique}
\label{subsubsec:decomposition_capacite_quantique}

La formule $\Delta_q = \Delta + \delta_{\text{enh}} - \delta_{\text{degr}}$
d\'ecompose explicitement les effets quantiques sur la capacit\'e
institutionnelle~:
\begin{equation}
\Delta_q = \Delta + \delta_{\text{enh}} - \delta_{\text{degr}}
\label{eq:decomposition_delta_q_discussion}
\end{equation}
\begin{itemize}
  \item $\delta_{\text{enh}}$ d\'esigne l'am\'elioration par outils quantiques
    (recherche Grover~\cite{Grover96}, apprentissage machine
    quantique~\cite{Biamonte2017})~;
  \item $\delta_{\text{degr}}$ d\'esigne la d\'egradation par bruit
    quantique~\cite{Preskill1998,Devoret2013} et inadaptation
    institutionnelle.
\end{itemize}
Cette d\'ecomposition permet de mod\'eliser l'\'evolution temporelle de
$\Delta_q$ \`a mesure que les institutions s'adaptent aux syst\`emes
post-quantiques et que les technologies quantiques
m\^urissent. Les valeurs retenues, $\Delta_q \in [25, 67]$~bits, sont
des hypoth\`eses de travail que le protocole de la
section~\ref{subsec:protocole_delta_q} permettra de mesurer.
\end{sloppypar}

% -----------------------------------------------------------------------
\subsection{Quatre sch\'emas fondamentaux du compromis CLO}
\label{subsec:quatre_patterns}
% -----------------------------------------------------------------------
\begin{sloppypar}

Les valeurs d\'eriv\'ees au chapitre~\ref{chap:experimentation}
(tableau~\ref{tab:composantes_opposabilite}) r\'ev\`elent quatre sch\'emas
structurels r\'ecurrents.

\medskip
\noindent\textbf{Sch\'ema 1. Tension vie priv\'ee--opposabilit\'e.}
Les syst\`emes maximisant la confidentialit\'e (Helios $\Conf = 0{,}92$,
zk-SNARK $\Conf = 0{,}89$) atteignent une opposabilit\'e mod\'er\'ee
($\Opp \approx 0{,}50$--$0{,}67$), tandis que le syst\`eme \`a haute
opposabilit\'e (ECDSA $\Opp = 0{,}90$) sacrifie la confidentialit\'e
($\Conf = 0{,}31$). Ce sch\'ema mat\'erialise la tension d\'ecrite par le
paradoxe de l'authenticit\'e invisible~\cite{MINKA2026CLO}.

\noindent\textit{Formalisation.} Ce sch\'ema est captur\'e par le Lemme
Conf-MI (sous-section~\ref{subsec:proprietes_fondamentales}), selon lequel plus
$I(M;\Sigma|C)$ est grand, plus $\Conf$ est faible. ECDSA op\`ere \`a une
extr\'emit\'e de ce compromis ($I$ grand, $\Conf$ faible). Helios est \`a
l'autre extr\'emit\'e pour la confidentialit\'e, mais sa fiabilit\'e
$\Rel = 0{,}95$ est une fiabilit\'e de \emph{v\'erification}, car sous $\Relrec$ Fano
imposerait $I(M;\Sigma|C) \geq 0{,}95\lambda - 1 \approx 121$~bits, incompatible
avec une confidentialit\'e \'elev\'ee. Le sch\'ema illustre donc la tension, sans
en \^etre une instance.

\medskip
\noindent\textbf{Sch\'ema 2. D\'ependance au contexte institutionnel.}
L'opposabilit\'e varie significativement avec l'expertise de
l'institution. Les r\'egulateurs sp\'ecialis\'es
($H(\mathcal{C}) = 11$ bits) extraient plus d'information que les
tribunaux g\'en\'eraux ($H(\mathcal{C}) = 8$ bits), tandis que les
tribunaux dot\'es d'une expertise PKI ($H(\mathcal{C}) = 20$ bits)
atteignent une interpr\'etabilit\'e substantiellement plus \'elev\'ee.
$H(\mathcal{C})$ est donc un param\`etre de conception critique, sur
lequel le concepteur peut agir. Les valeurs cit\'ees ici sont celles du contexte
d'esp\`ece~; elles ne se confondent pas avec les entropies de domaine du
tableau~\ref{tab:estimations_structurelles_HC}, distinction \'etablie
\`a la section~\ref{sec:methodologie_calcul}.

\noindent\textit{Implication pratique.} D\'eployer un zk-SNARK devant un
tribunal g\'en\'eral ($H(\mathcal{C}) \approx 8$ bits) au lieu d'un
r\'egulateur sp\'ecialis\'e ($H(\mathcal{C}) \approx 11$ bits)
r\'eduirait $I(V;\Sigma|C)$ de $2$ bits \`a environ $1$ bit ainsi que
l'opposabilit\'e de $0{,}67$ \`a environ $0{,}33$.

\medskip
\noindent\textbf{Sch\'ema 3. Plafond d'extraction des syst\`emes
z\'ero-connaissance.} M\^eme avec une entropie contextuelle \'elev\'ee,
les syst\`emes cryptographiquement opaques limitent fondamentalement
l'extraction institutionnelle \`a un ou deux bits. Ce plafond manifeste
l'hypoth\`ese d'extraction born\'ee. En effet, la
propri\'et\'e z\'ero-connaissance garantit $I(M;\Sigma|C) \approx 0$, de
sorte que l'hypoth\`ese d'extraction born\'ee donne
$I(V;\Sigma|C) \leq 0 + \Delta$~; seule la connaissance pr\'ealable
$\Delta$ permet alors une extraction non nulle.

\medskip
\noindent\textbf{Sch\'ema 4. La transparence active l'interpr\'etabilit\'e.}
La structure de v\'erification transparente d'ECDSA permet une
\'evaluation institutionnelle multidimensionnelle, $18$ bits extraits
sur $20$ disponibles, ce qui montre qu'une structure appropri\'ee peut \^etre
r\'ev\'el\'ee sans compromettre la solidit\'e cryptographique, tout en activant
l'utilit\'e institutionnelle.
\end{sloppypar}

% -----------------------------------------------------------------------
\subsection{Analyse de sensibilit\'e \`a la capacit\'e d'extraction}
\label{subsec:analyse_sensibilite}
% -----------------------------------------------------------------------
\begin{sloppypar}

\noindent\textbf{Impact de $\Delta$ sur l'extraction.} Le param\`etre
$\Delta$ majore l'information que l'interpr\`ete tire de la preuve
\emph{au-del\`a} de ce qu'elle dit du message. Son impact sur
$I(V;\Sigma|C)$, donc sur $\Opp$, d\'epend du r\'egime du protocole. Les
valeurs du tableau~\ref{tab:sensibilite_opp_delta} ne d\'erivent d'aucun mod\`ele
de r\'eponse. Elles illustrent des tendances plausibles, que le protocole de
la section~\ref{subsec:protocole_delta} permettrait de mesurer.

\begin{table}[htbp]
\centering
\caption{Sensibilit\'e de $\Opp$ \`a la capacit\'e d'extraction $\Delta$ (valeurs illustratives, sans mod\`ele de r\'eponse)}
\label{tab:sensibilite_opp_delta}
\small
\begin{tabular}{@{}llcccc@{}}
\toprule
\textbf{Syst\`eme} & \textbf{Sc\'enario} & $\Delta$ (bits) & $I(V;\Sigma|C)$ & $H(V|C)$ & $\Opp$ \\
\midrule
ECDSA & Standard (juge PKI) & 35 & 18 & 20 & 0,90 \\
ECDSA & Expert \'etendu & 50 & 19 & 20 & 0,95 \\
ECDSA & Novice & 15 & 10 & 20 & 0,50 \\
\midrule
Helios & Standard (juge g\'en\'eral) & 56 & 1 & 2 & 0,50 \\
Helios & Tribunal sp\'ecialis\'e & 70 & 1,4 & 2 & 0,70 \\
Helios & Novice & 30 & 0,6 & 2 & 0,30 \\
\midrule
zk-SNARK & Standard (r\'egulateur financier) & 73 & 2 & 3 & 0,67 \\
zk-SNARK & Expert cryptographe & 90 & 2,5 & 3 & 0,83 \\
zk-SNARK & R\'egulateur novice & 40 & 1,2 & 3 & 0,40 \\
\bottomrule
\end{tabular}
\end{table}

\noindent\textbf{Observation 1. Impact fort sur les syst\`emes
transparents.} Pour ECDSA, passer de $\Delta = 15$ (juge non initi\'e)
\`a $\Delta = 50$ (expert \'etendu) fait varier $\Opp$ de $0{,}50$ \`a
$0{,}95$, soit quarante-cinq points. La transparence de la structure de
preuve permet \`a l'expertise suppl\'ementaire de se traduire directement
en capacit\'e d\'ecisionnelle.

\noindent\textbf{Observation 2. Impact limit\'e sur les syst\`emes
z\'ero-connaissance.} Pour Helios, m\^eme lorsque $\Delta$ passe de
$30$ \`a $70$ bits, $\Opp$ ne varie que de $0{,}30$ \`a $0{,}70$. La
propri\'et\'e z\'ero-connaissance cr\'ee un plafond sur
l'extraction, ind\'ependamment de la capacit\'e institutionnelle, de sorte que
l'opacit\'e cryptographique ne se compense que partiellement par
l'expertise pr\'ealable. Ce plafond est compatible avec le goulot
contextuel, qui majore $I(V;\Sigma|C)$ par $H(V|C) = 2$~bits.

\noindent\textbf{Observation 3. H\'et\'erog\'en\'eit\'e des r\'egimes.}
Les trois syst\`emes appartiennent \`a des r\'egimes fondamentalement
diff\'erents selon le rapport $I(V;\Sigma|C)/H(V|C)$ avant l'apport de
$\Delta$~:
\begin{itemize}
\item ECDSA rel\`eve d'un r\'egime riche en information ($I \approx H(V|C)$ m\^eme sans $\Delta$)~;
\item Helios rel\`eve d'un r\'egime pauvre en information ($I \approx H(V|C)/2$ m\^eme avec $\Delta$ maximal)~;
\item zk-SNARK rel\`eve d'un r\'egime interm\'ediaire, sensible \`a $\Delta$ mais born\'e.
\end{itemize}
\end{sloppypar}

% -----------------------------------------------------------------------
\subsection{Implications pour la conception et la r\'eglementation}
\label{subsec:implications_conception}
% -----------------------------------------------------------------------
\begin{sloppypar}
Le Trilemme CRO conditionnel r\'eoriente la conception de protocoles
cryptographiques, qui passe de l'optimisation abstraite (\og quelles propri\'et\'es
de s\'ecurit\'e sont atteignables~? \fg) \`a l'ing\'enierie contrainte
par la capacit\'e (\og quelles propri\'et\'es cette institution peut-elle
effectivement utiliser, \'etant donn\'e sa capacit\'e mesurable
$\Delta$~? \fg).

\subsubsection{Choix de la variante SOTP}
\label{subsubsec:tableau_comparatif_sotp}\label{subsubsec:algo_selection_variante}

Les trois variantes SOTP sont compar\'ees au tableau~\ref{tab:comparaison_sotp_construction}
et l'Algorithme~\ref{alg:selection_variante_sotp} en formalise le choix
(chapitre~\ref{chap:constructions}).

L'algorithme part des deux seuils qui d\'elimitent les r\'egimes d'entropie
contextuelle, $H^*(\lambda) = \sqrt{\lambda}$ et $H^{**}(\lambda) =
\lambda/\log_2\lambda$. Il traite d'abord le cas bloquant. Si le contexte
est assez riche pour exiger SOTP-High mais que l'institution ne dispose pas
d'un trusted setup fiable, aucune recommandation optimale n'existe. La
proc\'edure signale l'impasse sans la masquer~; elle rabat alors
sur SOTP-Base, au prix d'une perte d'opposabilit\'e explicite. Dans les
autres cas, elle compare $H(\mathcal{C})$ aux deux seuils et retourne la
variante du r\'egime correspondant, soit SOTP-Low en contexte rare, SOTP-Base
en r\'egime de transition (o\`u elle dimensionne au passage la projection
al\'eatoire) ou SOTP-High en contexte riche, o\`u le circuit de preuve doit
\^etre compil\'e. Le co\^ut de la d\'ecision se r\'eduit \`a deux
comparaisons~; l'essentiel du travail reste l'estimation pr\'ealable de
$H(\mathcal{C})$ et de $\Delta$.

\subsubsection{M\'ethodologie de conception en six phases}
\label{subsubsec:methodologie_cinq_phases}

La progression F$\to$B$\to$E$\to$S$\to$C$\to$X fournit un cadre de d\'ecision explicite l\`a o\`u la
pratique actuelle consiste souvent \`a d\'eployer le protocole
cryptographiquement le plus puissant disponible, ind\'ependamment de
la capacit\'e institutionnelle de l'environnement cible.

\begin{table}[htbp]
\centering
\caption{M\'ethodologie de conception
F$\to$B$\to$E$\to$S$\to$C$\to$X}
\label{tab:methodologie_FBESCX}
\begin{tabular}{@{}p{1.8cm}p{2.8cm}p{6.5cm}@{}}
\toprule
\textbf{Phase} & \textbf{Action} & \textbf{Outils} \\
\midrule
\textbf{F} & Formaliser le contexte
  & Estimation $H(\mathcal{C})$~\cite{Shannon48,CoverThomas06} \\
\textbf{B} & Borner la capacit\'e
  & Triangulation cognitive ($\Delta$)~\cite{Miller56,Simon82,Ericsson95,Chi2006} \\
\textbf{E} & \'Etendre au r\'egime~PQ
  & \'Evaluer $\Delta_q$, $\eta_q$~(section~\ref{sec:cro_pq}) \\
\textbf{S} & Saturer~: retenir la variante SOTP du r\'egime
  & Variante SOTP (Low/Base/High), heuristique~(section~\ref{sec:sotp_construction}) \\
\textbf{C} & Certifier la faisabilit\'e
  & Exigences comparées \`a la borne raffin\'ee~\cite{MINKA2026CLO,GDPR2016,eIDAS14} \\
\textbf{X} & eXp\'erimenter la mesure
  & Protocoles exp\'erimentaux~(chapitre~\ref{chap:experimentation}) \\
\bottomrule
\end{tabular}
\end{table}

\noindent\textbf{Phase C. Test de faisabilit\'e.} V\'erifier~:
\begin{equation}
\Conf_{\min} \cdot \mathsf{Rel}_{\mathrm{rec},\min} \cdot \Opp_{\min} \stackrel{?}{\leq}
\frac{1+\kappa/\lambda}{\lambda}\cdot\min\Bigl\{1,\;\frac{\kappa/\lambda+\Delta}{\gamma}\Bigr\}
\label{eq:test_faisabilite}
\end{equation}
(forme raffin\'ee~\eqref{eq:trilemme_cro_raffine}, applicable aux protocoles qui
satisfont les hypoth\`eses du th\'eor\`eme).
Si l'in\'egalit\'e est viol\'ee, trois options s'offrent au concepteur, \`a savoir relaxer une exigence,
augmenter $\Delta$ par investissement dans la formation institutionnelle
ou augmenter $\lambda$ (impact marginal si $\Delta$ reste fixe).

\subsubsection{Strat\'egie de conception par r\'egime}
\label{subsubsec:strategie_par_regime}

\begin{definition}[R\'egime de basse capacit\'e]\label{def:regime_low}
Une institution op\`ere en r\'egime de \textbf{basse capacit\'e} si
$\Delta < 0{,}1\lambda$. Elle ne dispose pas de l'expertise n\'ecessaire
pour interpr\'eter des preuves cryptographiques
complexes~\cite{Herley09}. \textit{Recommandation.} Des m\'ecanismes
transparents (ECDSA avec pistes d'audit claires). La cryptographie
pr\'eservant la vie priv\'ee est inadapt\'ee \`a ce niveau.
\end{definition}

\begin{definition}[R\'egime de capacit\'e moyenne]\label{def:regime_mid}
Une institution op\`ere en r\'egime de \textbf{capacit\'e moyenne} si
$0{,}1\lambda \le \Delta \le 0{,}5\lambda$. Elle peut atteindre des
compromis CRO \'equilibr\'es~\cite{Grossklags10}.
\textit{Recommandation.} Des protocoles comme Helios avec des aides
\`a l'interpr\'etation structur\'ees.
\end{definition}

\begin{definition}[R\'egime de haute capacit\'e]\label{def:regime_high}
Une institution op\`ere en r\'egime de \textbf{haute capacit\'e} si
$\Delta > 0{,}5\lambda$. Elle peut soutenir des preuves cryptographiques
avanc\'ees~\cite{Grossklags10,Simon82}. \textit{Recommandation.}
Des zk-SNARKs, avec un investissement substantiel en formation et infrastructure
d'expertise.
\end{definition}

L'impl\'ementation se d\'eroule en trois \'etapes. \textbf{\'Etape~1.
Profilage institutionnel.} \'Evaluer $\Delta$ par inventaire de la
disponibilit\'e d'experts, tests de charge cognitive sur des preuves
repr\'esentatives, analyse de d\'ecisions historiques, audit de la
qualit\'e de documentation~\cite{Miller56,Ericsson95,Simon82}.
\textbf{\'Etape~2. S\'election protocolaire.} Choisir les protocoles
correspondant au r\'egime mesur\'e~; ne jamais d\'eployer un protocole
dont le r\'egime cible d\'epasse la capacit\'e actuelle sans planification
de transition. \textbf{\'Etape~3. Renforcement de capacit\'e.}
Investir dans la formation, les outils d'aide \`a la d\'ecision et les
trajectoires de migration progressives.

\subsubsection{Trois recommandations de conception}
\label{subsubsec:trois_recommandations}

L\`a o\`u la strat\'egie par r\'egime ci-dessus indique \emph{quel}
protocole retenir pour une capacit\'e donn\'ee, l'analyse de sensibilit\'e
de la sous-section~\ref{subsec:analyse_sensibilite} indique \emph{o\`u
investir} pour faire \'evoluer cette capacit\'e. Elle se traduit en trois
recommandations directement actionnables.

\noindent\textit{Recommandation 1 (ad\'equation contexte--protocole).} Le
choix du protocole doit \^etre guid\'e par le $H(\mathcal{C})$ de
l'institution cible. Un tribunal g\'en\'eral ($\Delta \approx 35$ bits) ne
tire aucun b\'en\'efice de zk-SNARK avanc\'es ($\Opp \approx 0{,}40$)
l\`a o\`u ECDSA offrirait $\Opp \approx 0{,}90$.

\noindent\textit{Recommandation 2 (investissement en formation).}
L'investissement en formation institutionnelle ($\Delta$ croissant) a des
rendements d\'ecroissants pour les syst\`emes z\'ero-connaissance mais
croissants pour les syst\`emes transparents. \`A budget de formation
\'egal, le gain d'opposabilit\'e est plus \'elev\'e sur ECDSA que sur Helios.

\noindent\textit{Recommandation 3 (compromis explicite).} Tout syst\`eme
visant $\Opp > 0{,}7$ dans un contexte \`a faible $H(\mathcal{C})$ doit
accepter une r\'eduction de $\Conf$. La fronti\`ere de Pareto
$(\Conf, \Opp)$ est born\'ee par le Trilemme CRO
conditionnel~\cite{MINKA2026CRO}, pour les protocoles qui satisfont ses
hypoth\`eses. Les valeurs chiffr\'ees de ces recommandations sont
illustratives (tableau~\ref{tab:sensibilite_opp_delta}).

\subsubsection{Seuils d'opposabilit\'e et admissibilit\'e juridique}
\label{subsubsec:seuils_opposabilite}

La m\'etrique $\Opp$ permet de proposer des seuils minimaux
diff\'erenci\'es par contexte juridictionnel~\cite{Kerr04,Casey11,Mason07}.

\begin{table}[htbp]
\centering
\caption{Seuils d'opposabilit\'e sugg\'er\'es par r\'egime juridictionnel}
\label{tab:seuils_opposabilite}
\begin{tabular}{@{}p{3.5cm}p{3.5cm}c@{}}
\toprule
\textbf{R\'egime} & \textbf{Standard de preuve} & \textbf{$\Opp_{\min}$} \\
\midrule
P\'enal & Certitude raisonn\'ee~\cite{Blackstone1769,Wigmore1983}
  & $\ge 0{,}70$ \\
Civil & Balance des probabilit\'es~\cite{Federal1988}
  & $\ge 0{,}50$ \\
Arbitrage commercial & Pr\'epond\'erance~\cite{Born2014}
  & $\ge 0{,}40$ \\
Conformit\'e r\'eglementaire & Audit interne~\cite{COSO2013}
  & $\ge 0{,}30$ \\
\bottomrule
\end{tabular}
\end{table}

Ces seuils reposent sur des estimations analytiques~; la variabilit\'e
intra-institutionnelle de $\Delta$ introduit une
incertitude~\cite{Kahneman11}~; ils constituent un guide
indicatif~\cite{Casey11}.

\noindent\textbf{Implications pour l'audit des syst\`emes de
r\'emun\'eration de l'\'Etat.} L'analyse des syst\`emes de gestion de la
solde et des pensions \`a l'\'echelle nationale, particuli\`erement dans
le contexte du Cameroun, pose un d\'efi pour l'audit
informatique en environnement institutionnel contraint. L'extraction et
l'analyse exhaustive des donn\'ees de paiement d\'epassent la v\'erification
comptable, car elles sont susceptibles de d\'evoiler
des informations sensibles pour la s\'ecurit\'e de l'\'Etat. Le cadre
CLO-CRO sugg\`ere une \emph{hypoth\`ese de travail} selon laquelle, dans les contextes o\`u
$\Delta \ll \gamma$, la forme raffin\'ee de la borne pr\'edit un plafond
d'opposabilit\'e bas pour les protocoles tr\`es confidentiels, de sorte que
l'auditeur pourrait devoir arbitrer entre la force probante de sa preuve et
la protection des donn\'ees. Cette hypoth\`ese ne d\'ecoule pas du cadre
sans mesure de $\Delta$ et de $\gamma$ dans ce contexte~; elle est \`a
\'evaluer empiriquement.

\subsubsection{Implications pour les standards NIST PQC}
\label{subsubsec:implications_nist_pqc}

L'analyse des standards NIST~PQC r\'ev\`ele des implications directes
pour le d\'eploiement~\cite{NIST.FIPS.203,NIST.FIPS.204,NIST.FIPS.205,Alagic2022}.

\begin{table}[htbp]
\centering
\caption{Implications CRO pour les standards NIST~PQC}
\label{tab:implications_nist_pqc}
\begin{tabular}{@{}p{3.2cm}p{2.5cm}p{2.8cm}p{2.8cm}@{}}
\toprule
\textbf{Standard} & \textbf{R\'egime} & \textbf{Recommandation}
  & \textbf{Pr\'ecaution} \\
\midrule
FIPS~203 (ML-KEM) & Hors champ (KEM) & --- & Pas de preuve \`a interpr\'eter~\cite{Kyber2022,Peikert2016} \\
FIPS~204 (ML-DSA) & Signature, $\eta_q = 0$ & Institutions form\'ees
  & Formation requise~\cite{Dilithium2022} \\
FIPS~205 (SLH-DSA) & Signature, $\eta_q = 0$ & Haute s\'ecurit\'e
  & Taille de signature \'elev\'ee~\cite{SPHINCS2022,Bernstein2015} \\
\bottomrule
\end{tabular}
\end{table}

Pour les institutions \`a faible capacit\'e ($\Delta_q < 30$~bits), les
standards post-quantiques n\'ecessitent des investissements significatifs
en formation et outils d'interpr\'etation avant tout d\'eploiement
effectif.

\subsubsection{Opportunit\'es de standardisation}
\label{subsubsec:opportunites_standardisation}

Trois possibilit\'es de standardisation sont identifi\'ees.
Pour \textbf{NIST~PQC}, l'int\'egration de l'opposabilit\'e comme crit\`ere
additionnel permettrait une s\'election guid\'ee par les contraintes de
d\'eploiement r\'eelles. Pour \textbf{ISO/IEC~27xxx}, une norme sur les
m\'etriques CLO fournirait un protocole standardis\'e de calcul de
$H(\mathcal{C})$ et $\Delta$. Pour la \textbf{CEMAC}, une directive sur la
cryptographie institutionnelle pour la monnaie mobile permettrait \`a six
pays d'harmoniser leurs exigences r\'eglementaires.
\end{sloppypar}

% -----------------------------------------------------------------------
\subsection{Positionnement vis-\`a-vis des travaux existants}
\label{subsec:positionnement_travaux}
% -----------------------------------------------------------------------
\begin{sloppypar}

\subsubsection{Cadres cryptographiques~: ZKP et flux d'information quantitatif}
\label{subsubsec:cadres_crypto_zkp}

Les syst\`emes zero-knowledge atteignent une efficacit\'e remarquable
mais cr\'eent une vuln\'erabilit\'e pr\'ecise dans les contextes
institutionnels, puisqu'une preuve zk-SNARK offre au r\'egulateur une
v\'erification binaire sans explication de la logique
sous-jacente~\cite{Groth16,BCTV14,PLONK20}. Le paradigme
\og ne rien r\'ev\'eler \fg{} optimise pour la confidentialit\'e mais
n\'eglige intrins\`equement l'opposabilit\'e~\cite{Bonneau16,Grimmelmann15}.
CLO formalise cette tension via $\Opp$~\cite{MINKA2026CLO}.

Le flux d'information quantitatif (QIF)~\cite{Smith09,Alvim12,Clarkson05}
mesure la fuite d'information pour la minimiser, alors que CLO cherche \`a
quantifier combien d'information est rendue interpr\'etable par
l'institution pour l'optimiser~; les deux objectifs sont sym\'etriques mais
oppos\'es. Les \'etudes cognitives calibrent $\Delta$ et
$\alpha$~\cite{Miller56,Simon82,Kahneman11}~; les pr\'edictions CLO
guident la conception des \'etudes.

\subsubsection{S\'ecurit\'e utilisable}
\label{subsubsec:securite_utilisable}

La recherche en s\'ecurit\'e utilisable a \'etabli des r\'esultats
empiriquement riches (folk models~\cite{Wash2010}, rationalit\'e de
l'ins\'ecurit\'e~\cite{Herley09}). Elle ne formalise toutefois pas
math\'ematiquement ces limites, n'en d\'erive pas de bornes d'impossibilit\'e
et ne produit pas de constructions optimales
garanties~\cite{Cranor05,Bonneau12}. La m\'etrique $\Opp$ exprime
formellement ce que les \'etudes d'utilisateurs observent
qualitativement~\cite{MINKA2026CLO}, de sorte que les deux approches se compl\`etent.

\subsubsection{Recherche juridique et responsabilit\'e algorithmique}
\label{subsubsec:recherche_juridique}

Citron et Pasquale~\cite{Citron2008,Citron14} ont argument\'e que les
syst\`emes automatis\'es doivent fournir une explicabilit\'e pour
satisfaire les exigences de due process technologique. Ces travaux
identifient pr\'ecis\'ement les probl\`emes que CLO cherche \`a
formaliser, mais s'arr\^etent \`a l'analyse
normative~\cite{Casey11,Kerr04}. CLO leur fournit un outil
quantitatif, en compl\'ement de l'analyse normative. La valeur de cette
instrumentation ne sera \'etablie qu'une fois que des \'etudes
contr\^ol\'ees auront confirm\'e que les valeurs analytiques de $\Opp$
pr\'edisent effectivement les jugements d'admissibilit\'e r\'eels.
\end{sloppypar}

% -----------------------------------------------------------------------
\subsection{Limitations et probl\`emes ouverts}
\label{subsec:limitations_problemes_ouverts}
% -----------------------------------------------------------------------
\begin{sloppypar}

\subsubsection{Limitations techniques}
\label{subsubsec:limitations_techniques}

\textbf{Les axiomes de Fano et Pinsker ne sont pas d\'eriv\'es ab initio.}
Notre formalisation les int\`egre comme propri\'et\'es des instances de
protocole sans les d\'eriver \`a partir de la th\'eorie des
probabilit\'es. Une formalisation compl\`ete n\'ecessiterait l'int\'egration
des biblioth\`eques de mesure-th\'eorie de Mathlib.

\textbf{Les constructions SOTP ne sont pas encore couvertes par Lean4.}
La formalisation couvre le noyau alg\'ebrique du Trilemme, mais pas les
r\'eductions cryptographiques des constructions SOTP ni l'extension
post-quantique. Ces extensions constituent une direction de recherche
future.

\subsubsection{Limitations m\'ethodologiques (L1--L6)}
\label{subsubsec:limitations_methodologiques}

Six limitations d\'elimitent pr\'ecis\'ement le p\'erim\`etre
d'applicabilit\'e des r\'esultats.

\paragraph{L1. Nature analytique des estimations CLO.}
Les m\'etriques $\Conf$, $\Relver$, $\Opp$ des syst\`emes
analys\'es ont \'et\'e calcul\'ees par analyse param\'etrique \`a partir
de sp\'ecifications publi\'ees, sans mesure aupr\`es d'acteurs
institutionnels r\'eels~\cite{MINKA2026CLO}. Des \'etudes contr\^ol\'ees
avec juges, auditeurs et garants en situation r\'eelle sont n\'ecessaires
pour \'etablir la correspondance entre m\'etriques analytiques et
comportement observ\'e~\cite{Casey11,Goodman15,Campbell1976,Saltelli2008}. De
plus, ces syst\`emes ne sont pas des instances du Trilemme
(Remarque~\ref{rem:portee_cas_etude}), de sorte que leurs valeurs calibrent le cadre sans le
tester.

\paragraph{L2. Quantification non standardis\'ee de $H(\mathcal{C})$.}
La m\'ethode de calcul de l'entropie contextuelle n'est pas
standardis\'ee. Deux chercheurs analysant le m\^eme syst\`eme pourraient
d\'eriver des valeurs diff\'erentes selon les hypoth\`eses retenues sur
la granularit\'e du contexte. La voie de r\'esolution passe par un
consensus communautaire.

\paragraph{L3. Statut de publication des security games.}
Les bornes CLO publi\'ees dans IEEE Access~\cite{MINKA2026CLO,MINKA2026CRO,MINKA2026PQCRO}
sont \'etablies dans le formalisme information-th\'eorique. Les jeux
cryptographiques (IND-CONF, BIND-REL, EXT-OPP) et les r\'eductions
game-based des constructions SOTP sont d\'evelopp\'es dans les
pr\'epublications~\cite{MINKA2025CRO,MINKA2025Q2CSI,MINKA2025ZKNR_UC}
mais n'ont pas encore fait l'objet d'une publication dans une revue
\`a comit\'e de lecture.

\paragraph{L4. \'Equivalence formelle incompl\`ete.}
Les ponts \'etablis entre les deux formalismes montrent que la
r\'esistance aux jeux \emph{implique} des bornes inf\'erieures sur les
m\'etriques CLO. La r\'eciproque, selon laquelle une forte opposabilit\'e
information-th\'eorique impliquerait une r\'esistance aux jeux,
n'est pas \'etablie. Cette asym\'etrie est une limitation
structurelle, puisque les garanties information-th\'eoriques sont
\emph{n\'ecessaires} mais pas \emph{suffisantes} pour la s\'ecurit\'e
computationnelle.

\paragraph{L5. Mod\`eles quantiques simplifi\'es.}
L'extension quantique utilise principalement le canal d\'epolarisant
comme mod\`ele de bruit. Des ph\'enom\`enes plus fins (amortissement
d'amplitude~\cite{Devoret2013}, d\'ephasage~\cite{Breuer2002}, effets de
m\'emoire quantique~\cite{Rivas2012}) peuvent ne pas y \^etre
captur\'es. Des mod\`eles de canaux sp\'ecifiques pourraient donner des
bornes plus fines.

\paragraph{L6. Port\'ee limit\'ee aux preuves non interactives.}
L'analyse couvre principalement les syst\`emes de preuves non
interactives~\cite{Groth16,BCTV14}. Les protocoles \`a rounds
multiples~\cite{Goldreich1987}, les architectures
multi-prouveurs~\cite{CramerDamgardSchoenmakers} et les verdicts
distribu\'es~\cite{Evans2018} peuvent pr\'esenter des caract\'eristiques
de compromis substantiellement diff\'erentes.

\subsubsection{Limitations conceptuelles}
\label{subsubsec:limitations_conceptuelles}

\textbf{Mod\`ele de contexte statique.} Nous traitons le contexte $C$
comme fixe. Or les institutions apprennent et s'adaptent, de sorte que $H(C)$ et
$\Delta$ \'evoluent par accumulation de pr\'ec\'edents, changement
technologique et mises \`a jour r\'eglementaires~\cite{Simon82}.
\textbf{Fonctions de verdict d\'eterministes.} Nous mod\'elisons
$V(\Sigma, C)$ comme d\'eterministe dans les exemples (auquel cas $\Opp = 1$,
Remarque~\ref{rem:verdict_deterministe}), alors que le th\'eor\`eme autorise un
verdict qui d\'epend aussi de ressources institutionnelles. La prise de d\'ecision r\'eelle
implique des processus stochastiques, des d\'esaccords multi-personnes et
des biais cognitifs~\cite{Casey10,Kahneman11}. \textbf{Hypoth\`ese de
$\beta$-r\'egularit\'e.} Le Th\'eor\`eme~\ref{thm:conditional_CRO_trilemma}
requiert $\beta = O(1)$~; des constructions pathologiques pourraient la
violer. \textbf{Port\'ee du t\'emoin.} Le t\'emoin de coh\'erence Lean4
atteint la borne dans le mod\`ele abstrait (forme rel\^ach\'ee de Fano), mais il
n'est pas r\'ealisable par un protocole, de sorte que la question d'un protocole concret
qui approche la borne reste ouverte.
\end{sloppypar}

\begin{boiteArbitre}
L'arbitre dresse l'inventaire de ce qui manque. Certains axiomes sont admis sans
\^etre d\'eriv\'es, tandis que la couverture Lean4 n'inclut encore ni SOTP ni
l'extension post-quantique. Six limitations m\'ethodologiques sont
nomm\'ees sans d\'etour, parmi lesquelles les estimations analytiques, l'entropie non
standardis\'ee, le r\'ealisme computationnel, le bruit quantique simplifi\'e
et la restriction aux preuves non interactives. Quatre limitations conceptuelles touchent
aussi le mod\`ele lui-m\^eme, \`a savoir le contexte statique, les verdicts d\'eterministes,
la r\'egularit\'e suppos\'ee et la port\'ee du t\'emoin.

Il note que rien n'est dissimul\'e, puisque chaque limite est pos\'ee et que chaque
cons\'equence en est tir\'ee. Il y voit un dossier en cours d'instruction,
l'instruction exigeant de savoir exactement
o\`u sont les trous.
\end{boiteArbitre}

% -----------------------------------------------------------------------
\subsection{La port\'ee constructive~: de l'impossibilit\'e au guide}
\label{subsec:portee_constructive}
% -----------------------------------------------------------------------
\begin{sloppypar}

Les th\'eor\`emes d'impossibilit\'e f\'econds sont souvent accompagn\'es de
constructions qui en montrent la port\'ee. CAP trouve son sens dans les
syst\`emes distribu\'es r\'eels (Spanner, CockroachDB)~\cite{Brewer00,GilbertLynch02}.
FLP trouve le sien dans les protocoles partiels (Paxos,
Raft)~\cite{FLP85}. Le Trilemme CRO se prolonge dans la
\textbf{famille SOTP}, qui r\'eunit trois constructions cryptographiques, une par r\'egime
d'entropie contextuelle, dont les propri\'et\'es de confidentialit\'e et de
fiabilit\'e de v\'erification sont prouv\'ees mais dont l'opposabilit\'e est
guid\'ee par une heuristique de conception.

\subsubsection{Ce que la construction apporte et ce qu'elle n'\'etablit pas}

\textbf{La borne devient utilisable pour la conception.} Un ing\'enieur qui
conna\^it la forme raffin\'ee de la borne sait quelles exigences sont
incompatibles~; avec SOTP, il dispose aussi d'un point de d\'epart par
r\'egime (contexte rare $\to$ SOTP-Low, transition $\to$ SOTP-Base, contexte
riche $\to$ SOTP-High).

\textbf{La finesse de la borne reste ouverte.} Les fractions de la borne
atteintes ($0{,}45$--$0{,}75$) rel\`event d'heuristiques et le t\'emoin de
coh\'erence Lean n'est pas r\'ealisable. Qu'un protocole concret approche la borne \`a facteur
constant est une question ouverte (Conjecture~\ref{conj:optimalite_universelle}).

\textbf{La borne devient un guide.} Elle dit
quel compromis est exclu~; les constructions indiquent comment s'en
approcher, sans garantie d'optimalit\'e.

\subsubsection{Port\'ee de la contribution SOTP}

\textbf{SOTP-Low} couvre le r\'egime des contextes rares
($H(C) < \sqrt{\lambda}$). Elle utilise un encodage d\'eterministe du contexte,
avec confidentialit\'e calculatoire sous PRG~\cite{BlumMicali84} et fiabilit\'e
de v\'erification sous DLog~\cite{DiffieHellman76}.

\textbf{SOTP-Base} couvre le r\'egime de transition. Elle utilise la projection
Johnson-Lindenstrauss~\cite{JohnsonLindenstrauss84,DasguptaGupta03}, dimensionn\'ee
pour pr\'eserver les distances entre contextes \`a $\varepsilon$ pr\`es.

\textbf{SOTP-High} couvre le r\'egime des contextes riches
($H(C) > \lambda/\log\lambda$). Elle utilise des zk-SNARKs
Groth16~\cite{Groth16,BCTV14} et un hachage Poseidon~\cite{Poseidon2019} pour
engager un contexte arbitrairement riche dans un condens\'e de 256 bits.

L'optimalit\'e universelle de la famille, qui interdirait toute am\'elioration
super-constante du produit par rapport \`a la meilleure variante, est une
conjecture (Conjecture~\ref{conj:optimalite_universelle}).

\begin{boiteArbitre}
L'arbitre comprend que la borne seule ne suffit pas, puisqu'une limite
d\'epourvue de construction ressemble \`a une porte ferm\'ee dont on ignorerait le plan. La famille SOTP
esquisse ce plan en indiquant par o\`u commencer dans chaque r\'egime, sans
pr\'etendre avoir atteint la limite. L'impossibilit\'e devient ainsi un guide.
\end{boiteArbitre}

\end{sloppypar}

% -----------------------------------------------------------------------
\subsection{Ce que la formalisation enseigne~: le\c{c}on m\'ethodologique}
\label{subsec:verification_corrige_lecon}
% -----------------------------------------------------------------------
\begin{sloppypar}

Les formulations alternatives issues de la formalisation Lean4, y compris le
cas limite qui \'etablit la n\'ecessit\'e du r\'egime de confidentialit\'e, sont
pr\'esent\'ees \`a la sous-section~\ref{subsec:lean4_erreurs}. La pr\'esente
sous-section se concentre sur la le\c{c}on m\'ethodologique qui en d\'ecoule.

\paragraph{La rigueur se v\'erifie.}
Les formulations pr\'ecis\'ees touchaient \`a la structure des preuves~: le sens
d'une in\'egalit\'e, une borne inf\'erieure tir\'ee d'une borne sup\'erieure et la
r\'esolution d'une optimisation. Ces points \'echappent facilement \`a la relecture
parce qu'ils produisent des r\'esultats \emph{plausibles}. La v\'erification
m\'ecanique a mis en \'evidence les questions de sens d'in\'egalit\'e et de
quantification. Elle ne dispense pas de l'examen des d\'efinitions, qui a conduit
\`a d'autres pr\'ecisions (distinction entre fiabilit\'e de reconstruction et
fiabilit\'e de v\'erification, d\'efinition de $\eta_q$), hors de port\'ee d'une
formalisation qui prend ces d\'efinitions pour donn\'ees.

\paragraph{La formalisation est un instrument de d\'ecouverte.}
Au-del\`a de la certification de r\'esultats d\'ej\`a prouv\'es, la formalisation
r\'ev\`ele les hypoth\`eses implicites et les points structurels \`a pr\'eciser. Ce
constat prolonge les observations de Scerri~\cite{Scerri14} et de
Wiedijk~\cite{Wiedijk2008}, selon lesquelles revisiter les mod\`eles de preuve
r\'ev\`ele des probl\`emes fondationnels que ni la th\'eorie pure ni
l'exp\'erimentation n'identifient.

\paragraph{La formalisation gagne \`a pr\'ec\'eder la publication.}
Dans notre cas, la formalisation Lean4 a \'et\'e entreprise \emph{apr\`es} la
publication. Elle a conduit \`a des formulations plus pr\'ecises qui auraient pu
figurer d\`es la soumission. Pour les r\'esultats interdisciplinaires combinant
th\'eorie de l'information, cryptographie et sciences cognitives, la
formalisation m\'ecanique a sa place dans le processus de preuve lui-m\^eme, en
amont de la publication.

\paragraph{Cons\'equences pour la pratique.}
Cette exp\'erience conduit \`a trois recommandations pour la communaut\'e.
\textbf{Pour les auteurs}, il s'agit de formaliser le noyau alg\'ebrique avant
soumission, m\^eme partiellement~; une formalisation partielle fait d\'ej\`a
appara\^itre les points structurels \`a pr\'eciser. \textbf{Pour les revues}, il
convient d'inclure syst\'ematiquement un relecteur \og v\'erification formelle \fg{}
pour les articles dont le noyau math\'ematique est formalisable~; la charge est
faible, le b\'en\'efice \'elev\'e. \textbf{Pour la formation}, il importe d'enseigner
la formalisation comme partie int\'egrante du m\'etier de chercheur. Il s'agit
d'une comp\'etence transf\'erable qui b\'en\'eficie \`a toutes les disciplines
math\'ematis\'ees.

\begin{boiteArbitre}
L'arbitre referme le dossier. Des formulations plus pr\'ecises sont issues de la
machine ou d'une relecture attentive des d\'efinitions, puis retenues par les
auteurs. Il en retient que ni la relecture seule ni la machine seule ne
suffisent, car l'une produit des preuves plausibles tandis que l'autre certifie
ce qu'on lui donne. L'arbitre comprend que la formalisation est une n\'ecessit\'e
\'epist\'emologique pour tout r\'esultat interdisciplinaire dont aucune des
disciplines constituantes ne contr\^ole l'int\'egralit\'e.
\end{boiteArbitre}

\end{sloppypar}

% -----------------------------------------------------------------------
\subsection{Le statut des protocoles exp\'erimentaux~: un programme
falsifiable}
\label{subsec:statut_protocoles_experimentaux}
% -----------------------------------------------------------------------
\begin{sloppypar}

Le chapitre~\ref{chap:experimentation} occupe une position
inhabituelle dans l'\'economie de cette th\`ese, puisqu'il sp\'ecifie des
protocoles exp\'erimentaux sans les avoir d\'eploy\'es. Cette position
m\'erite d'\^etre clarifi\'ee, car elle d\'etermine la port\'ee
\'epist\'emique de l'ensemble du travail.

\subsubsection{Statut \'epist\'emique}

Le chapitre~\ref{chap:experimentation} se situe entre r\'esultats
empiriques (aucun protocole n'a \'et\'e d\'eploy\'e) et th\'eorie pure (les
protocoles sont sp\'ecifi\'es jusqu'\`a l'op\'erationnel). Il rel\`eve de la
conception m\'ethodologique, puisqu'il sp\'ecifie \emph{comment} valider empiriquement
le cadre, sans encore l'avoir fait. Cette sp\'ecification constitue une
contribution en soi, pour trois raisons.

\textbf{Elle rend le cadre falsifiable.} Un cadre th\'eorique
conditionnel \`a des param\`etres non mesurables est non falsifiable au
sens de Popper~\cite{Popper1934}. En sp\'ecifiant les protocoles de
mesure, les seuils de falsification et les pr\'edictions testables, le
chapitre~\ref{chap:experimentation} transforme le Trilemme CRO en
programme scientifique falsifiable. Cette transformation conditionne sa
port\'ee empirique.

\textbf{Elle guide la communaut\'e.} Les protocoles exp\'erimentaux sont
sp\'ecifi\'es \`a un niveau op\'erationnel rare dans la litt\'erature
th\'eorique, qu'il s'agisse de la stratification des \'echantillons, des estimateurs avec biais
et variance, des crit\`eres de fiabilit\'e inter-\'evaluateurs ou du calendrier de
d\'eploiement. Cette sp\'ecification permet \`a d'autres chercheurs de
reproduire, de critiquer ou d'\'etendre le programme.

\textbf{Elle documente les limitations.} Le volet~B du chapitre d\'erive
les valeurs d'opposabilit\'e des trois syst\`emes de r\'ef\'erence et
\'etablit la coh\'erence interne du cadre. Il fournit une premi\`ere
\textbf{calibration analytique} en attente de validation
exp\'erimentale tout en documentant explicitement ce qui reste \`a mesurer.

\subsubsection{Parall\`eles avec d'autres disciplines}

Ce statut se retrouve dans les disciplines o\`u la th\'eorie pr\'ec\`ede
l'exp\'erience.

\textbf{Physique th\'eorique.} Une th\'eorie qui pr\'edit des faits
falsifiables et sp\'ecifie comment les tester est une contribution,
m\^eme si les tests n'ont pas encore \'et\'e men\'es. La pr\'ediction du
boson de Higgs a pr\'ec\'ed\'e sa d\'etection de plusieurs d\'ecennies.

\textbf{Biologie th\'eorique.} Les mod\`eles de plissement
prot\'eique~\cite{Anfinsen1973} ont pr\'ec\'ed\'e les m\'ethodes de
validation exp\'erimentale (AlphaFold~\cite{Jumper2021}).

\textbf{\'Economie comportementale.} La th\'eorie des
perspectives~\cite{Kahneman1979} a sp\'ecifi\'e des protocoles
exp\'erimentaux avant leur d\'eploiement.

Dans tous ces cas, la \textbf{sp\'ecification du programme de
validation} est reconnue comme une contribution scientifique. Il
en va de m\^eme pour le chapitre~\ref{chap:experimentation}.

\subsubsection{Cons\'equences pour la lecture de la th\`ese}

Cette clarification engage trois mani\`eres de lire ce qui pr\'ec\`ede.

\textbf{Les bornes sont conditionnelles et leurs conditions sont
\'enonc\'ees.} Le Trilemme CRO, classique comme post-quantique, ne dit rien
en dehors du r\'egime o\`u il est \'etabli, lequel suppose l'extraction born\'ee, un verdict
efficace, la r\'egularit\'e du protocole et le r\'egime de fuite
$I(M;\Sigma|C) \leq \kappa/\lambda$. Ces conditions d\'elimitent le domaine de validit\'e du
r\'esultat, de sorte qu'un protocole plac\'e hors de ce domaine \'echappe \`a la
borne. Lire le trilemme comme une interdiction absolue serait donc le
mal lire. Sa validit\'e math\'ematique est \'etablie~; sa port\'ee
empirique d\'epend de grandeurs qui restent \`a mesurer.

\textbf{Le mod\`ele est r\'efutable et les seuils de r\'efutation sont
publics} (tableau~\ref{tab:criteres_refutation}). Si les mesures de $\Delta$
ou de $\Delta_q$ s'\'ecartent syst\'ematiquement des plages attendues, ou si le
produit observ\'e d\'epasse la borne sur un syst\`eme qui satisfait les
hypoth\`eses du th\'eor\`eme, le mod\`ele
d'extraction born\'ee devra \^etre r\'evis\'e, sans r\'einterpr\'etation des
mesures. Le chapitre~\ref{chap:experimentation} fixe ces seuils
avant toute campagne de mesure, ce qui interdit de les ajuster apr\`es
coup.

\textbf{La coh\'erence \'etablie est interne.} Les ordres
de grandeur retenus pour les syst\`emes de r\'ef\'erence sont mutuellement
coh\'erents et compatibles avec la litt\'erature cognitive~; les deux voies
d'estimation ne sont pas ind\'ependantes. Cela ne constitue pas une
validation exp\'erimentale, que cette th\`ese ne revendique d'ailleurs pas.

\textbf{La port\'ee du travail est donc quadruple.} Elle est
th\'eorique par les m\'etriques et les bornes (C1--C3),
constructive par les protocoles qui les illustrent (C4),
certifi\'ee par la v\'erification m\'ecanique de leur noyau (C5) et
exp\'erimentale par le programme falsifiable qui les rend
testables (C6). Cette quadruple port\'ee structure la conclusion
g\'en\'erale.

\begin{boiteArbitre}
L'arbitre comprend que sp\'ecifier comment tester une th\'eorie est
aussi important que la th\'eorie elle-m\^eme. Une th\'eorie non testable
est une sp\'eculation. Une th\'eorie testable, m\^eme non encore
test\'ee, est une science. Le
chapitre~\ref{chap:experimentation} est ce qui fait passer le
Trilemme CRO de la premi\`ere cat\'egorie \`a la seconde.
\end{boiteArbitre}

\end{sloppypar}

% -----------------------------------------------------------------------
\subsection{Le\c{c}ons m\'ethodologiques}
\label{subsec:lecons_methodologiques}
% -----------------------------------------------------------------------
\begin{sloppypar}
Les difficult\'es rencontr\'ees produisent des enseignements
transf\'erables \`a la communaut\'e travaillant \`a l'intersection de
la cryptographie, de la th\'eorie de l'information et des m\'ethodes
formelles.

\subsubsection{La formalisation renforce les th\'eor\`emes interdisciplinaires}
\label{subsubsec:formalisation_renforce}

La v\'erification formelle expose les hypoth\`eses cach\'ees et clarifie
les d\'ependances logiques, l\`a o\`u le raisonnement informel les
absorbe tacitement~\cite{Wiedijk2008,Avigad2014}. Ce r\'esultat est
coh\'erent avec Scerri~\cite{Scerri14}, pour qui revisiter les mod\`eles de
preuve r\'ev\`ele des probl\`emes fondationnels que ni la th\'eorie pure
ni l'exp\'erimentation n'identifient.

\subsubsection{L'axiomatisation est une n\'ecessit\'e pragmatique}
\label{subsubsec:axiomatisation_pragmatique}

Tant que les math\'ematiques sous-jacentes ne sont pas enti\`erement
formalis\'ees dans une biblioth\`eque partag\'ee, int\'egrer les
r\'esultats profonds comme propri\'et\'es d'instance est une approche
valide et transparente. La cl\'e est de s\'eparer nettement le noyau
certifi\'e des fondations accept\'ees comme
vraies~\cite{BartheEA2013,Doussot23}.

\subsubsection{Les constantes explicites sont n\'ecessaires}
\label{subsubsec:constantes_explicites}

Masquer les constantes dans $O(\cdot)$ est un raccourci acceptable en
s\'ecurit\'e classique~; en cryptographie institutionnelle, o\`u
$\Delta$, $\alpha$, $\beta$, $H(C)$ doivent \^etre mesur\'es pour
permettre la d\'ecision juridique et administrative, cette abstraction
devient un obstacle op\'erationnel. Notre borne explicite
$\gamma^{-1}(\Delta/\lambda + \kappa(\Delta+1)/\lambda^2 +
\kappa^2/\lambda^3)$ rend calculable la borne sur le produit CRO pour
tout ensemble de param\`etres institutionnels mesur\'es.

\subsubsection{La transparence scientifique comme mod\`ele}
\label{subsubsec:transparence_scientifique}

Judson~\cite{Judson22} note un sch\'ema pr\'eoccupant, \`a savoir que les mises en garde
et les d\'ecisions de conception difficiles sont souvent renvoy\'ees aux
praticiens \`a travers des heuristiques mal comprises. Notre
documentation explicite des bornes non prouv\'ees, des hypoth\`eses
n\'ecessaires et des limitations empiriques r\'epond \`a cette
pr\'eoccupation, car la rigueur consiste \`a d\'elimiter pr\'ecis\'ement les probl\`emes
ouverts au lieu de les \'eviter.

\subsubsection{La v\'erification formelle comme n\'ecessit\'e \'epist\'emologique}
\label{subsubsec:verification_necessite}

La formalisation Lean4 constitue une n\'ecessit\'e
\'epist\'emologique pour les r\'esultats interdisciplinaires combinant
th\'eorie de l'information, sciences cognitives et cryptographie, dont
aucun des trois champs n'a la ma\^itrise compl\`ete. En for\c{c}ant chaque
hypoth\`ese, quantification et approximation \`a \^etre explicite, la
v\'erification m\'ecanique approfondit la compr\'ehension du
th\'eor\`eme lui-m\^eme. Les formulations alternatives et le cas limite sont
document\'es \`a la sous-section~\ref{subsec:lean4_erreurs}.
\end{sloppypar}

\begin{boiteArbitre}
L'arbitre retient quatre certitudes. La cryptographie institutionnelle
constitue un domaine propre, car sa question centrale n'\'etait pos\'ee nulle part.
Une impossibilit\'e prouv\'ee repr\'esente la forme la plus nette de
connaissance. Ce qu'il ne peut pas obtenir simultan\'ement a
d\'esormais un nom, une mesure et une borne. Enfin, la rigueur formelle a
mis au jour des propri\'et\'es
informationnelles transform\'ees en conditions v\'erifiables, une
structure de preuve pr\'eservant la borne dominante et des constantes que
la notation $O(\cdot)$ dissimulait. Il sait maintenant qu'il attend
de la section suivante les perspectives qui transforment ces contraintes
en programme.
\end{boiteArbitre}

\begin{table}[htbp]
\centering
\caption{Synth\`ese~: r\'ev\'elations, implications, limitations et
le\c{c}ons par contribution}
\label{tab:synthese_discussions}
\begin{tabular}{@{}p{1.8cm}p{3cm}p{3cm}p{3.5cm}@{}}
\toprule
\textbf{Contrib.} & \textbf{R\'ev\'elation / R\'esultat}
  & \textbf{Implication pratique} & \textbf{Travail futur} \\
\midrule
C1~: CLO & M\'etrique $\Opp$ in\'edite~\cite{MINKA2026CLO}
  & Quantification des compromis
  & \'Etudes contr\^ol\'ees \\
C2~: Trilemme & Borne $\gamma^{-1}(\Delta/\lambda + \kappa(\Delta+1)/\lambda^2 + \kappa^2/\lambda^3)$~\cite{MINKA2026CRO}
  & Guide conception classique
  & Extension Lean4 \\
C3~: PQ & Borne $\gamma_q^{-1}(\Delta_q/\lambda + \kappa(\Delta_q+1)/\lambda^2 + \kappa^2/\lambda^3)$~\cite{MINKA2026PQCRO};
  d\'ecomp. $\delta_{\text{enh}},\delta_{\text{degr}}$
  & Standards NIST~PQC
  & Canaux quantiques r\'eels \\
C4~: SOTP & $\Conf$, $\Relver$ prouv\'ees~; opposabilit\'e heuristique~(section~\ref{sec:sotp_construction})
  & Protocoles de r\'ef\'erence
  & Preuves interactives \\
C5~: Lean4 & Propri\'et\'es d'instance explicites~;
  forme raffin\'ee certifi\'ee~(section~\ref{sec:lean4})~;
  t\'emoins (mod\`ele abstrait)
  & Certification m\'ecanique
  & Extension SOTP et PQ \\
C6~: Protocoles exp\'erimentaux & Programme de validation falsifiable
  & Feuille de route empirique
  & D\'eploiement effectif \\
\midrule
\multicolumn{4}{l}{\textit{Limitations transversales~:}} \\
L1--L6 & Estimations analytiques~; $H(\mathcal{C})$ non standardis\'e~;
  r\'ealisme computationnel~; bruit quantique simplifi\'e~;
  preuves non interactives
  & \multicolumn{2}{l}{Programme exp\'erimental
    futur} \\
\bottomrule
\end{tabular}
\end{table}
% ============================================================================
% CHAPITRE 4 — SECTION 4.4 : PERSPECTIVES DE RECHERCHE
% Fichier : partie4/perspectives.tex
% Version : 2.2 — alignée sur la formalisation Lean4 révisée
% ============================================================================

\section{PERSPECTIVES DE RECHERCHE}
\label{sec:perspectives}
% ============================================================================
\begin{sloppypar}
Le corpus th\'eorique de la cryptographie institutionnelle est maintenant
constitu\'e. Il r\'eunit des m\'etriques formalis\'ees (C1)~\cite{MINKA2026CLO}, une borne
d'impossibilit\'e \'etablie (C2)~\cite{MINKA2026CRO}, une extension quantique
d\'emontr\'ee (C3)~\cite{MINKA2026PQCRO}, des constructions de r\'ef\'erence aux propri\'et\'es
cryptographiques prouv\'ees (C4) (section~\ref{sec:sotp_construction}), un noyau Lean4 certifi\'e (C5)
(section~\ref{sec:lean4}) et un programme de validation sp\'ecifi\'e (C6)
(chapitre~\ref{chap:experimentation}). Que faire de ces fondations~?
Cette section y r\'epond en
tra\c{c}ant la feuille de route th\'eorique d'un sous-domaine dont les
premiers th\'eor\`emes viennent d'\^etre pos\'es, sans projeter de d\'eploiements.

La progression F$\to$B$\to$E$\to$S$\to$C$\to$X
(formaliser, borner, \'etendre, saturer, certifier, eXp\'erimenter)
qui a guid\'e cette th\`ese se prolonge au-del\`a de ces r\'esultats, comme le montrent les quatre axes qui structurent cette
prospective.

Chacun de ces quatre axes prolonge un aspect
distinct du cadre CLO-CRO-PQ-SOTP. Il s'agit des fondations computationnelles
post-quantiques port\'ees par le probl\`eme AIIP~\cite{MINKA2025AIIP} et le
CASH Toolkit (section~\ref{subsec:axe1_aiip_cash})~; la voie protocolaire
ouverte par Q2CSI et la s\'erie
ZK-NR~\cite{MINKA2025Q2CSI,MINKA2025ZKNR_Exploration,MINKA2025ZKNR_Design,MINKA2025ZKNR_UC}
(section~\ref{subsec:axe2_q2csi_zknr})~; les extensions th\'eoriques du
Trilemme CRO vers le calcul homomorphe, le multi-parties et le
post-quantique quantitatif (section~\ref{subsec:axe3_extensions})~; et
l'interaction encore largement inexploit\'ee entre cryptographie
institutionnelle et apprentissage automatique (section~\ref{subsec:axe4_ia}).
Un programme \`a horizons structur\'es conclut la section
(section~\ref{subsec:programme_recherche}).
\end{sloppypar}

\begin{boiteArbitre}
Le voyage de l'arbitre ne s'arr\^ete pas avec cette th\`ese. Les fondations
sont pos\'ees (m\'etriques, bornes, constructions, certification,
protocoles) mais le programme ouvert est plus vaste que ce qui vient
d'\^etre r\'esolu. Il reste \`a obtenir une borne dont le terme dominant serait en $O(\lambda^{-2})$,
\`a \'etablir l'optimalit\'e des constructions SOTP, \`a calibrer empiriquement les param\`etres
quantiques~$\Delta_q$ et~$\eta_q$ ainsi qu'\`a \'etendre les constructions
au calcul homomorphe et multi-parties. L'arbitre de demain
aura besoin de protocoles que la cryptographie institutionnelle ne sait pas
encore concevoir. Ce chapitre en trace la feuille de route.
\end{boiteArbitre}

% -----------------------------------------------------------------------
\subsection{Maturit\'e du sous-domaine}
\label{subsec:maturite_sous_domaine}
% -----------------------------------------------------------------------
\begin{sloppypar}
Un sous-domaine scientifique est constitu\'e lorsqu'il dispose d'un formalisme
propre, d'une impossibilit\'e fondamentale, de constructions et d'une
certification m\'ecanique~\cite{Wiedijk2007,Gonthier08}. Le cadre
CLO-CRO-CRO$^q$-SOTP-Lean4 r\'eunit ces \'el\'ements, sous deux r\'eserves. D'une part,
l'optimalit\'e des constructions est conjecturale~; d'autre part, les param\`etres
institutionnels restent \`a mesurer. Il s'agit d'un programme de recherche au sens
de Lakatos~\cite{Lakatos1978}, dont le noyau est fix\'e et la ceinture
protectrice explicite.

La question devient \guillemotleft~quelles directions ouvertes en
d\'efiniront la croissance et quelles mesures en confirmeront ou en
infirmeront les hypoth\`eses~?~\guillemotright.

L'\'etat des contributions \`a la date de soutenance figure au
tableau~\ref{tab:contributions_these}.

Trois classes de probl\`emes ouverts structurent la fronti\`ere du
sous-domaine. La premi\`ere concerne les fondations computationnelles, c'est-\`a-dire la n\'ecessit\'e
d'asseoir les instanciations pratiques sur des hypoth\`eses de duret\'e
nouvelles, r\'esistantes aux adversaires quantiques~\cite{Bernstein2009,Chen2016}
et analys\'ees par la communaut\'e IACR. La deuxi\`eme concerne la
g\'en\'eralisation des bornes. Dans ses formes classique et
quantique, le Trilemme CRO couvre les preuves non interactives~\cite{Groth16}~; le passage
aux protocoles \`a rounds multiples, au calcul multi-parties~\cite{Evans2018}
et aux primitives homomorphes~\cite{Gentry09} reste entier. La troisi\`eme
est l'ouverture d'un territoire in\'edit, car l'interaction entre interpr\`etes
institutionnels artificiels et protocoles SOTP modifie structurellement les
hypoth\`eses du mod\`ele.
\end{sloppypar}

% -----------------------------------------------------------------------
\subsection{Axe I. Fondations computationnelles~: AIIP et le CASH Toolkit}
\label{subsec:axe1_aiip_cash}
% -----------------------------------------------------------------------
\begin{sloppypar}

\subsubsection{Le probl\`eme AIIP~: hypoth\`ese de duret\'e post-quantique}
\label{subsubsec:probleme_aiip}

Toute discipline cryptographique repose sur des hypoth\`eses de duret\'e dont
le statut est \'etabli par cryptanalyse
communautaire~\cite{Goldreich01,KatzLindell14}. La cryptographie
institutionnelle ne fait pas exception. En effet, si les Trilemmes CRO classiques et
quantiques sont prouv\'es au sens de la th\'eorie de l'information, leurs
instanciations pratiques requi\`erent des primitives r\'esistant aux
adversaires quantiques \`a temps polynomial~(QPT)~\cite{Boneh2011,Watrous2009}.

Le \textbf{probl\`eme AIIP} (\textit{Affine Iterated Inversion
Problem})~\cite{MINKA2025AIIP} propose une telle hypoth\`ese. \'Etant
donn\'e $f \in \mathbb{F}_q[x]$ de degr\'e $d \geq 2$, un param\`etre
d'it\'eration $n$ et une cible $y \in \mathbb{F}_q$, AIIP requiert de trouver
$x$ tel que $f^{(n)}(x) = y$, o\`u $f^{(n)}$ d\'esigne la $n$-i\`eme
composition it\'er\'ee. Sa duret\'e s'\'etablit par deux r\'eductions
ind\'ependantes~\cite{MINKA2025AIIP}, \`a savoir une r\'eduction au logarithme discret dans
le Jacobien d'une courbe hyperelliptique de genre exponentiellement
grand~(HCDLP)~\cite{Menezes2001,Enge2002} et une r\'eduction polynomiale \`a la
r\'esolution de syst\`emes MQ structur\'es~(MQ-NP)~\cite{Courtois2000,Ding2004}.
Ces r\'eductions et leur port\'ee restent \`a \'etablir de fa\c{c}on
ind\'ependante. En particulier, une objection doit \^etre lev\'ee avant de
consid\'erer AIIP comme une hypoth\`ese candidate. En effet, sauf argument contraire,
r\'esoudre $f^{(n)}(x) = y$ se ram\`ene \`a $n$ extractions successives de racines
d'un polyn\^ome de degr\'e $d$ sur $\mathbb{F}_q$, ce qui se fait en temps
polynomial en $n$, $d$ et $\log q$ (algorithme de Cantor--Zassenhaus), en
suivant les ant\'ec\'edents successifs. Un argument qui exclut cette attaque
(par exemple un contr\^ole de la ramification de l'arbre des ant\'ec\'edents)
est donc un pr\'ealable. La cryptanalyse (attaques de
Gr\"obner~\cite{Faugere02,Faugere99}, algorithmes de
point-count~\cite{Schoof85,Kedlaya2001}, variantes
quantiques~(Grover~\cite{Grover96}, MNFS quantique~\cite{Jeppesen2022}))
reste \`a conduire.

\subsubsection{Le CASH Toolkit~: primitives align\'ees sur les propri\'et\'es CLO}
\label{subsubsec:cash_toolkit}

Le \textbf{CASH Toolkit} (\textit{Chaotic Affine Secure Hash}) regroupe trois
primitives post-quantiques construites en correspondance directe avec les
dimensions du tripl\'e~$(\Conf, \Rel, \Opp)$. Leur conception est guid\'ee par
un principe inhabituel dans la cryptographie classique, selon lequel l'opposabilit\'e
future de la preuve conditionne d\`es la conception les choix de transparence
alg\'ebrique~\cite{MINKA2026CLO,Casey11}.

\paragraph{CEE. Couche confidentialit\'e ($\Conf$).}
La primitive \textbf{CEE} (\textit{Chaotic Entropic
Expansion})~\cite{MINKA2025CEE} construit une fonction \`a sens unique sur des
applications polynomiales chaotiques it\'er\'ees sur
$\mathbb{F}_q$~\cite{Menezes1996,Stinson2005}. Pour $N$ it\'erations d'un
polyn\^ome $f \in \mathcal{F}_d$, la min-entropie r\'esiduelle est au moins
$\log_2 q - N\log_2 d$ bits. La borne de distance statistique \`a l'uniforme
annonc\'ee dans~\cite{MINKA2025CEE}, $(q-1)(d^N-1)/(2\sqrt{q})$, d\'epasse $1$ d\`es
que $d^N > 2\sqrt{q}/(q-1) + 1$, c'est-\`a-dire pour tout $N$ utile. Elle est
alors vide, si bien qu'une borne non triviale reste \`a \'etablir. La s\'ecurit\'e
quantique est de l'ordre $O(2^{\lambda/2})$ pour un niveau classique
$\lambda$~\cite{NielsenChuang2010}. La construction a pour trait caract\'eristique que sa transparence
alg\'ebrique est d\'elib\'er\'ement pr\'eserv\'ee, permettant
\`a $\Conf$ de coexister avec $\Opp$ sans conflit structurel~\cite{MINKA2026CLO}.

\paragraph{AOW. Couche fiabilit\'e ($\Rel$).}
La primitive \textbf{AOW} (\textit{Affine One-Wayness})~\cite{MINKA2025AOW}
fournit une v\'erification temporelle par \'evaluation polynomiale it\'erative
sur corps finis. Sa s\'ecurit\'e est r\'eduite, dans~\cite{MINKA2025AOW}, \`a HCDLP
et \`a AIIP, dont le statut est discut\'e ci-dessus~; elle s'accompagne d'une int\'egration
STARK~\cite{STARKs18,BenSasson2018} \`a mise \`a l'\'echelle logarithmique
pour la v\'erification zero-knowledge. La tra\c{c}abilit\'e temporelle qu'AOW
produit, ancrage de s\'equences v\'erifiables sans autorit\'e de
confiance~\cite{Haber1991,Bayer1993}, alimente directement $\Opp$, puisque
l'institution peut \'etablir \emph{quand} une preuve a \'et\'e g\'en\'er\'ee.

\paragraph{SH. Couche opposabilit\'e ($\Opp$).}
La primitive \textbf{SH} (\textit{Semantic Holder})~\cite{MINKA2025SH}
op\'erationnalise l'extraction alg\'ebrique depuis les traces de polyn\^omes
it\'er\'es pour produire des interpr\'etations l\'egalement
opposables~\cite{Kerr04,Casey11}. Elle vise trois propri\'et\'es, \`a savoir la
pr\'eservation entropique, l'explicabilit\'e institutionnelle et la contestabilit\'e
l\'egale v\'erifiable par STARK. Le score agr\'eg\'e $\Omega$ et la valeur
$\Omega \geq 0{,}60$ avanc\'es dans~\cite{MINKA2025SH} ne sont pas d\'efinis dans le
cadre de cette th\`ese~; ils ne sont pas repris ici comme r\'esultats.

Plusieurs directions restent ouvertes. Peut-on prouver que l'assemblage
$(CEE \;\|\; AOW \;\|\; SH)$ en architecture Gold-Iron-Clay \emph{pr\'eserve}
les bornes CRO \`a un facteur constant pr\`es~? Peut-on d\'efinir un score
d'opposabilit\'e pour SH qui soit reli\'e \`a $\Opp$~? Enfin, les preuves actuelles
op\'erant dans le mod\`ele de l'oracle al\'eatoire~\cite{Bellare1993,Canetti2004},
des preuves en mod\`ele standard renforceraient substantiellement la fondation
th\'eorique du toolkit.
\end{sloppypar}

% -----------------------------------------------------------------------
\subsection{Axe II. Voie protocolaire~: Q2CSI, ZK-NR et l'instanciation de l'esprit SOTP}
\label{subsec:axe2_q2csi_zknr}
% -----------------------------------------------------------------------
\begin{sloppypar}

\subsubsection{Q2CSI comme traduction architecturale de l'esprit SOTP}
\label{subsubsec:q2csi}

La famille SOTP propose des constructions de r\'ef\'erence par r\'egime,
param\'etr\'ees par $H(\mathcal{C})$ (section~\ref{sec:sotp_construction}), dont
l'optimalit\'e reste conjecturale. La question qui
s'ouvre naturellement est celle de l'\emph{instanciation concr\`ete}. Comment
traduire cette philosophie en architectures d\'eployables dans des
environnements cryptographiques composables~\cite{Canetti01,Canetti2000}~?

Le m\'eta-protocole \textbf{Q2CSI} (\textit{Quantum Composable and Contextual
Security Infrastructure})~\cite{MINKA2025Q2CSI} apporte une r\'eponse
pr\'ecise. En d\'ecomposant les garanties de s\'ecurit\'e en trois couches
fonctionnellement isol\'ees (Iron~($\Rel$), Gold~($\Conf$),
Clay~($\Opp$)) Q2CSI traduit architecturalement le paradigme
Gold-Iron-Clay dans le mod\`ele de Canetti~\cite{Canetti01,Canetti2000}.
L'article~\cite{MINKA2025Q2CSI} r\'esume ses performances par un indice
$\Gamma_\text{CRO}$ qui n'est pas d\'efini dans le cadre de cette th\`ese~; les
valeurs chiffr\'ees qu'il en donne ne sont donc pas reprises ici. Relier
cet indice au produit $\Conf\cdot\Relrec\cdot\Opp$ est une question ouverte. Q2CSI incarne
l'esprit SOTP sous la forme d'un sch\'ema d'architecture, distinct d'une
construction abstraite. Ce sch\'ema repose sur la s\'eparation dialectique des couches, l'isolation stricte des
fuites contextuelles et la compatibilit\'e avec les adversaires
quantiques~\cite{Unruh2010,Unruh2015}.

\subsubsection{La s\'erie ZK-NR~: instanciation composable par zero-knowledge}
\label{subsubsec:zk_nr}

La s\'erie \textbf{ZK-NR} (\textit{Zero-Knowledge Non-Repudiation}) protocol
mat\'erialise la maturation progressive d'un protocole institutionnel depuis
l'exploration architecturale jusqu'\`a la s\'ecurit\'e UC
compl\`ete~\cite{Canetti01}. Elle constitue en elle-m\^eme un paradigme
m\'ethodologique, qui consiste \`a s\'eparer l'exploration, la formalisation du design et la
certification de s\'ecurit\'e en trois stades distincts, chacun s'appuyant
rigoureusement sur le pr\'ec\'edent.

ZK-NR~v1~\cite{MINKA2025ZKNR_Exploration} compose commitments
Merkle~\cite{Merkle1987}, preuves STARKs~\cite{STARKs18}, signatures BLS \`a
seuil~\cite{Boldyreva2003,Boneh2003} et authentification
Dilithium~\cite{Dilithium2022} en un protocole de non-r\'epudiation sans
divulgation~; il \'etablit la faisabilit\'e conceptuelle, les preuves formelles
\'etant r\'eserv\'ees aux versions ult\'erieures. ZK-NR~v2~\cite{MINKA2025ZKNR_Design}
construit l'architecture compl\`ete satisfaisant les propri\'et\'es Q2CSI et le
mod\`ele NIZK-E~\cite{Groth16}, avec s\'eparation structurelle entre preuves
contextuelles et explications born\'ees~; la compatibilit\'e UC est prouv\'ee
symboliquement. ZK-NR~v3~\cite{MINKA2025ZKNR_UC} \'etablit la preuve
compl\`ete de s\'ecurit\'e UC pour les trois couches, mod\'elis\'ees comme
fonctionnalit\'es id\'eales avec s\'emantique d'effacement~\cite{Canetti2016}.
Il prouve aussi l'indistinguabilit\'e r\'eel/id\'eal sous hypoth\`eses
post-quantiques~\cite{Unruh2015}.

Trois questions ouvertes d\'elimitent la suite. Le sch\'ema Q2CSI est-il
\emph{complet}, c'est-\`a-dire, tout protocole institutionnel qui s\'epare
les trois couches est-il UC-\'equivalent \`a une instanciation Q2CSI~? \`a l'image du th\'eor\`eme de compl\'etude de
Canetti-Cohen-Lindell~\cite{CCL15}. La r\'eduction formelle
SOTP-High~$\to$~ZK-NR~v3 est-elle exhibable, comblant l'\'ecart entre la
th\'eorie abstraite de SOTP~(section~\ref{sec:sotp_construction}) et les
protocoles composables d\'eployables~? La structure \`a couches de ZK-NR
s'\'etend-elle aux r\'eseaux de preuves institutionnelles multi-parties
impliquant $n > 2$ institutions, via des sch\'emas de partage de
secret~\cite{Shamir1979,Blakley1979} et des preuves zero-knowledge
agr\'eg\'ees~\cite{Bunz2018}~?
\end{sloppypar}

% -----------------------------------------------------------------------
\subsection{Axe III. Extensions th\'eoriques du Trilemme CRO}
\label{subsec:axe3_extensions}
% -----------------------------------------------------------------------
\begin{sloppypar}

\subsubsection{Borne forte sur le produit CRO~: vers la caract\'erisation compl\`ete}
\label{subsubsec:borne_forte}

Le Trilemme CRO classique~(Th\'eor\`eme~\ref{thm:trilemme_cro_conditionnel})
\'etablit la borne explicite
\begin{equation}
\Conf \cdot \Relrec \cdot \Opp \leq \frac{1}{\gamma}\left(\frac{\Delta}{\lambda} + \frac{\kappa(\Delta+1)}{\lambda^2} + \frac{\kappa^2}{\lambda^3}\right)
\label{eq:borne_forte_classique}
\end{equation}
dont le terme dominant $\Delta/(\gamma\lambda)$ est en $O(\lambda^{-1})$ et les
termes correctifs en $O(\lambda^{-2})$, sous l'Hypoth\`ese~\ref{ass:regime}
($I(M;\Sigma|C) \leq \kappa/\lambda$). Le Trilemme CRO post-quantique
(Contribution C3, section~\ref{sec:cro_pq}) \'etablit la borne correspondante
\begin{equation}
\Conf^q \cdot \Relrec^q \cdot \Opp^q \leq \frac{1}{\gamma_q}\left(\frac{\Delta_q}{\lambda} + \frac{\kappa(\Delta_q+1)}{\lambda^2} + \frac{\kappa^2}{\lambda^3}\right)
\label{eq:borne_forte_quantique}
\end{equation}
o\`u $\Delta_q$ repr\'esente la capacit\'e d'extraction institutionnelle
adapt\'ee au quantum et $\eta_q$ le coefficient de perte d'interpr\'etabilit\'e
quantique.

\begin{remark}[Statut d'une borne forte en $O(\lambda^{-2})$]
\label{rem:borne_forte_statut}
La borne~\eqref{eq:borne_forte_classique}, v\'erifi\'ee dans
l'annexe~\ref{chap:annexe_lean4}, a un terme dominant $\Delta/(\gamma\lambda)$ en
$O(\lambda^{-1})$~; seuls ses termes correctifs, $\kappa(\Delta+1)/\lambda^2$ et
$\kappa^2/\lambda^3$, sont d'ordre $\lambda^{-2}$ et $\lambda^{-3}$. Une borne dont
le terme dominant serait en $O(\lambda^{-2})$ reste une question ouverte.
La forme raffin\'ee~\eqref{eq:trilemme_cro_raffine} montre d'ailleurs qu'en
r\'egime \og rare\fg{} la borne est gouvern\'ee par Fano, en $1/\lambda$, ind\'ependamment
de $\Delta$. Les autres probl\`emes ouverts concernent l'extension de la
formalisation \`a SOTP et au cadre post-quantique ainsi que la d\'erivation dans
Mathlib des in\'egalit\'es de Fano et Pinsker, aujourd'hui pos\'ees comme
propri\'et\'es des instances.
\end{remark}

La formalisation de l'optimisation analytique sur l'intervalle complet
$[0, \lambda]$ reste un probl\`eme ouvert, car elle n\'ecessite des outils
d'analyse r\'eelle (d\'eriv\'ees de fonctions avec racine carr\'ee) non encore
disponibles dans Mathlib. Elle permettrait d'\'etudier le comportement de la
borne hors du r\'egime de l'Hypoth\`ese~\ref{ass:regime}.

\subsubsection{G\'en\'eralisation \`a la r\'egion r\'ealisable compl\`ete}
\label{subsubsec:region_realisable}

La borne sur le produit scalaire ne r\'epond pas \`a la question
de savoir quelle est la r\'egion r\'ealisable compl\`ete
$\mathcal{R} \subset [0,1]^3$ des triplets $(\Conf, \Rel, \Opp)$
atteignables~\cite{MINKA2026CLO}. Deux syst\`emes de produit CRO
identique~($= 0{,}4$) peuvent pr\'esenter des profils radicalement
diff\'erents~; seule la caract\'erisation de $\mathcal{R}$ permet de
d\'eterminer quels profils sont r\'ealisables et quels sont les arbitrages
Pareto-optimaux~\cite{Boyd2004}.

Deux conjectures guident l'exploration. La premi\`ere, dite de \emph{convexit\'e}, pose que $\mathcal{R}$
est convexe dans $[0,1]^3$, car la composition probabiliste de deux protocoles
g\'en\`ere un troisi\`eme avec m\'etriques interpol\'ees~\cite{MINKA2026CLO}.
La seconde, dite de \emph{fronti\`ere poly\'edrale}, pose que la fronti\`ere de Pareto de $\mathcal{R}$
est d\'efinie par un syst\`eme fini d'in\'egalit\'es lin\'eaires issues de
Fano~\cite{Fano61}, Pinsker~\cite{Pinsker1964} et du principe de traitement
des donn\'ees~\cite{CoverThomas06}. Des techniques issues de la programmation
convexe~\cite{Boyd2004,Nesterov2018} et de la th\'eorie taux-distortion
\'etendue~\cite{Berger71,Gray2011} semblent prometteuses pour attaquer ces
questions.

\subsubsection{G\'en\'eralisation au chiffrement homomorphe (FHE)}
\label{subsubsec:generalisation_fhe}

Le chiffrement totalement homomorphe~\cite{Gentry09,Vaikuntanathan2012,Brakerski2014}
maximise $\Conf$, les donn\'ees restent chiffr\'ees tout au long du calcul,
mais d\'egrade conjecturalement $\Opp$, car l'interpr\`ete institutionnel ne peut
extraire d'information verdictuelle d'un chiffr\'e homomorphe sans cl\'e de
d\'echiffrement~\cite{Gentry2010,Halevi2012}. Cette tension n'est pas encore
formalis\'ee dans le cadre CLO. La r\'esolution pourrait passer par des
constructions autorisant un d\'ecryptage partiel
contr\^ol\'e~\cite{Brakerski2013,Boneh2013} vers la couche Clay sans
compromettre la couche Gold~; il s'agit d'une conjecture.

\subsubsection{G\'en\'eralisation au calcul multi-parties s\'ecuris\'e (MPC)}
\label{subsubsec:generalisation_mpc}

Dans un protocole MPC \`a $n$
parties~\cite{Yao82,GMW87,Goldreich1987,Evans2018,Garay15}, l'opposabilit\'e
agr\'eg\'ee
\begin{equation}
  \Opp_\text{MPC} = \frac{I(V_1,\ldots,V_n;\,\Sigma \mid C)}{H(V_1,\ldots,V_n \mid C)}
\end{equation}
peut exhiber des comportements radicalement diff\'erents de l'opposabilit\'e
centralis\'ee~\cite{MINKA2026CLO}. Des effets d'amplification~(chaque partie
apportant une vue partielle enrichissant l'information de verdict) ou de
cancellation~(vues partielles contradictoires r\'eduisant
$I(V_1,\ldots,V_n;\,\Sigma \mid C)$) sont th\'eoriquement envisageables, sous
r\'eserve de formalisation. La structure du Trilemme CRO se pr\'eserve-t-elle
dans ce cadre~? Des techniques issues du mod\`ele
UC~\cite{Canetti01,Canetti2000} et de la th\'eorie des jeux
coop\'eratifs~\cite{Peleg2007,Owen2013} semblent indiqu\'ees.

\subsubsection{Cadre quantique quantitatif}
\label{subsubsec:cadre_quantique_quantitatif}

La migration vers des constructions quantiques, commitments bas\'es sur les
r\'eseaux~\cite{Ajtai96,Regev05,Peikert2016},
STARKs~\cite{STARKs18,BenSasson2018}, pr\'eserve-t-elle les propri\'et\'es
de confidentialit\'e et de fiabilit\'e de v\'erification \'etablies \`a la
section~\ref{sec:sotp_construction}, sachant que les engagements de Pedersen
reposent sur le logarithme discret et ne sont pas post-quantiques~? Quelles bornes CLO quantitatives
s'obtiennent pour des adversaires Grover~\cite{Grover96,Zhandry2012}~?
L'opposabilit\'e naturelle pour des protocoles \`a $r$ rounds
\begin{equation}
  \Opp_r = \frac{I(V;\,\Sigma_1,\ldots,\Sigma_r \mid C)}{H(V \mid C)}
\end{equation}
satisfait-elle une variante du Trilemme CRO~? Ces questions prolongent l'\'etape B
de la progression, par l'\'etablissement de bornes g\'en\'eralis\'ees.

L'extension quantique (Contribution C3, section~\ref{sec:cro_pq}) a ouvert
cette voie en introduisant les param\`etres $\Delta_q$ et $\eta_q$. Le
programme de recherche vise maintenant \`a calibrer ces param\`etres pour les
signatures NIST PQC (ML-DSA~\cite{Dilithium2022}, SLH-DSA~\cite{SPHINCS2022}),
pour lesquelles $\eta_q = 0$, puis \`a d\'evelopper une th\'eorie quantitative des
pertes d'interpr\'etabilit\'e pour le sc\'enario d'un support quantique~\cite{NielsenChuang2010,Preskill1998}.

Trois pistes th\'eoriques prolongent cette extension. \textbf{(D) Contraction.}
Pour un canal $\Lambda$ de coefficient de contraction $\eta_\Lambda$
(in\'egalit\'e de traitement des donn\'ees forte), on aurait
$\Opp^q \leq \eta_\Lambda\,\Opp$, ce qui remplacerait le terme additif
$\eta_q H(C)$ par un facteur multiplicatif. \textbf{(E) Granularit\'e et
capacit\'e.} Lier $\alpha$ et $\Delta$ (un interpr\`ete plus capable distingue
des verdicts plus fins) ferait d\'ependre $\gamma$ de $\Delta$ et modifierait le
r\'egime de la forme raffin\'ee. \textbf{(B2) Contexte priv\'e.} Pour SOTP, un
mod\`ele o\`u une partie $\tilde C$ du contexte est transmise \`a l'institution
seule (et non publique) permettrait de donner un statut prouv\'e aux
minorations d'opposabilit\'e, aujourd'hui heuristiques.

\subsubsection{Enseignements m\'ethodologiques de la formalisation Lean4}
\label{subsubsec:enseignements_lean4}

La formalisation Lean4 du noyau alg\'ebrique du Trilemme CRO
(section~\ref{sec:lean4}) a r\'ev\'el\'e trois le\c{c}ons structurelles pour
le d\'eveloppement du sous-domaine.

\paragraph{Propri\'et\'es informationnelles comme conditions v\'erifiables.}
La formalisation int\`egre Fano et Pinsker comme champs explicites de la
structure \texttt{ProtocolInstance} tandis qu'elle traite l'extraction born\'ee et le r\'egime de
confidentialit\'e comme hypoth\`eses du th\'eor\`eme. Ce choix transforme des
propri\'et\'es profondes en conditions v\'erifiables \`a chaque instance,
clarifiant le statut \'epist\'emique du th\'eor\`eme.

\paragraph{Coh\'erence et port\'ee du th\'eor\`eme.}
La formalisation produit deux t\'emoins qui d\'emontrent la non-vacuit\'e et
la n\'ecessit\'e des hypoth\`eses. Le \emph{t\'emoin de coh\'erence} atteint
exactement la borne th\'eorique pour une instance du mod\`ele abstrait,
confirmant que les hypoth\`eses formelles sont compatibles (ce point n'est
pas r\'ealisable par un protocole, Remarque~\ref{rem:portee_temoin}). Le \emph{t\'emoin de port\'ee} v\'erifie
les Hypoth\`eses 1, 2$'$ et 3 mais s'\'ecarte du r\'egime de confidentialit\'e,
atteignant $1/4$ constant, ce qui \'etablit que ce r\'egime est n\'ecessaire
\`a la borne. Le cas limite complet est document\'e \`a la
sous-section~\ref{subsec:lean4_erreurs}.

\paragraph{Couverture du corpus et scalabilit\'e.}
La formalisation couvre le noyau alg\'ebrique du th\'eor\`eme.
L'extension aux constructions SOTP compl\`etes et aux r\'esultats
post-quantiques constitue une direction de recherche future, estim\'ee \`a
un investissement substantiel mais r\'ealiste pour une \'equipe d\'edi\'ee.

\subsubsection{Extension de la formalisation Lean4~: SOTP et PQ}
\label{subsubsec:extension_lean4_sotp_pq}

La formalisation Lean4 actuelle (annexe~\ref{chap:annexe_lean4}) couvre
int\'egralement le noyau alg\'ebrique du Trilemme CRO
classique, y compris les t\'emoins de coh\'erence et de port\'ee. Deux
extensions restent \`a d\'evelopper.

\begin{itemize}
    \item \textbf{Extension SOTP.} La formalisation des constructions
          SOTP-Low, SOTP-Base et SOTP-High n\'ecessiterait la
          formalisation des r\'eductions cryptographiques (PRG, DLog,
          q-SDH) et des algorithmes de compression
          Johnson-Lindenstrauss. L'effort est estim\'e \`a 1000--1500 lignes
          additionnelles.
    \item \textbf{Extension post-quantique.} La formalisation du
          Trilemme CRO PQ n\'ecessiterait l'int\'egration des
          in\'egalit\'es quantiques (Pinsker quantique, Gentle
          Measurement, Fano quantique) dans une biblioth\`eque de
          th\'eorie de l'information quantique. L'effort est estim\'e \`a
          1000--1500 lignes additionnelles.
\end{itemize}

Cette extension, estim\'ee \`a 2000--3000 lignes additionnelles au total,
constituerait la premi\`ere certification m\'ecanique compl\`ete d'un
cadre cryptographique institutionnel.
\end{sloppypar}

% -----------------------------------------------------------------------
\subsection{Axe IV. Cryptographie institutionnelle et intelligence artificielle}
\label{subsec:axe4_ia}
% -----------------------------------------------------------------------
\begin{sloppypar}

\subsubsection{Interpr\`etes institutionnels artificiels}
\label{subsubsec:interpretes_artificiels}

Le cadre CLO mod\'elise l'interpr\`ete institutionnel par sa capacit\'e
cognitive $\Delta$ et un espace de verdicts $\mathcal{V}$~\cite{MINKA2026CLO,Simon82}.
Un mod\`ele de langage de grande taille~(LLM) utilis\'e comme interpr\`ete
pr\'esente des propri\'et\'es qualitativement diff\'erentes, puisque $\Delta_\text{LLM}$
peut d\'epasser $H(\mathcal{C})$ gr\^ace \`a des connaissances
hors-contexte~\cite{Brown2020,OpenAI2023}, mais est \'egalement stochastique
et potentiellement adversarial~\cite{Carlini2021,Zou2023}. Si $\Delta \gtrsim
H(\mathcal{C})$, la borne CRO (forme raffin\'ee~\eqref{eq:trilemme_cro_raffine}) peut-elle
\^etre approch\'ee arbitrairement~? Autrement dit, des interpr\`etes IA
rel\`event-ils fondamentalement la fronti\`ere de Pareto dans l'espace
CLO~? Cette question est ouverte~\cite{Bommasani2021}. Un adversaire IA, LLM
finement ajust\'e pour l'extraction d'information institutionnelle,
constitue-t-il une menace qualitativement nouvelle pour la couche Clay, ou
reste-t-il born\'e par les limites issues de la th\'eorie de
l'information~\cite{CoverThomas06}~? Les deux questions rel\`event de l'\'etape F
du programme, car elles requi\`erent d'abord une d\'efinition formelle de
$\Delta_\text{LLM}$ avant toute borne~\cite{Weidinger2023}.

\subsubsection{Apprentissage automatique comme outil d'estimation CLO}
\label{subsubsec:apprentissage_estimation}

Les m\'etriques CLO sont d\'efinies au sens de la th\'eorie de
l'information~\cite{Shannon48,CoverThomas06}; leur estimation sur donn\'ees
r\'eelles est computationnellement difficile~\cite{Paninski2003}. Les
techniques d'estimation d'information mutuelle par apprentissage~(MINE~\cite{Belghazi18},
m\'ethodes par diff\'erence de log-vraisemblance~\cite{Poole2019,Oord2018})
ouvrent une voie vers des \'evaluations empiriques automatis\'ees des
propri\'et\'es CLO. L'estimation de $H(\mathcal{C})$ pour des contextes
institutionnels r\'eels via des mod\`eles de compression
neuronaux~\cite{Hinton2007,Kingma2013}, l'estimation de $\Delta$ par des tests
de compr\'ehension automatis\'es sur corpus de d\'ecisions institutionnelles,
ces directions permettraient de remplacer les estimations analytiques
actuelles~(ECDSA, Helios, zk-SNARKs) par des mesures directes. Leurs
r\'esultats resteraient soumis \`a calibration
exp\'erimentale~\cite{Campbell1976,Saltelli2008}; ils constitueraient
n\'eanmoins la premi\`ere validation empirique directe du cadre.

\subsubsection{Audit cryptographique de l'intelligence artificielle}
\label{subsubsec:audit_ia}

La direction inverse est \'egalement possible. Des protocoles SOTP permettant
de prouver des propri\'et\'es de mod\`eles IA,
robustesse~\cite{Goodfellow2015,Madry2018},
\'equit\'e~\cite{Barocas2019,Mitchell2021}, conformit\'e
r\'eglementaire~\cite{EU2024AI,EU2021AI}, \`a des r\'egulateurs
institutionnels sans r\'ev\'eler les poids du mod\`ele constituent une
application naturelle du paradigme CLO. L'audit cryptographique d'un mod\`ele
de d\'ecision consisterait ainsi \`a prouver sa conformit\'e via un sch\'ema SOTP-High sans exposer
son architecture. Ce programme croise la v\'erification du
calcul~(\textit{verifiable computation})~\cite{Gennaro2010,Walfish2015} et la
gouvernance de l'IA~\cite{Floridi2018,Jobin2019}, territoire d'application
neuf pour lequel le cadre CLO-CRO-SOTP est th\'eoriquement arm\'e.
\end{sloppypar}

% -----------------------------------------------------------------------
\subsection{Programme de recherche \`a horizon structur\'e}
\label{subsec:programme_recherche}
% -----------------------------------------------------------------------
\begin{sloppypar}

\begin{boiteArbitre}
L'arbitre quitte l'exploration. Quatre axes ont \'et\'e ouverts, qui portent sur les fondations
computationnelles, la voie protocolaire, les extensions th\'eoriques et l'intelligence
artificielle. Chacun ouvre des territoires prometteurs, mais en partie encore sp\'eculatifs.
Il faut d\'esormais planifier, en distinguant ce qui peut \^etre
fait en un an, en trois ans ou en dix ans selon que cela d\'epend d'une preuve, d'un
financement ou d'une communaut\'e. L'arbitre note que le propos rel\`eve d\'esormais moins d'un
chercheur isol\'e que d'un programme qui se donne des
\'ech\'eances.
\end{boiteArbitre}

\subsubsection{Horizon court terme (1--3 ans)~: consolidation des fondations}
\label{subsubsec:horizon_court_terme}

Cinq t\^aches conditionnent les horizons ult\'erieurs. D'abord,
\'ecrire l'argument qui exclut l'inversion successive par extraction de
racines, puis soumettre AIIP \`a la cryptanalyse~; la robustesse de cette
hypoth\`ese est pr\'ealable \`a tout usage applicatif du CASH
Toolkit~\cite{MINKA2025AIIP}. Ensuite,
\'etendre la formalisation Lean4 \`a SOTP et au cadre post-quantique
(cf.~sous-section~\ref{subsubsec:extension_lean4_sotp_pq}), avec pour objectif
d'\'elargir la couverture de l'annexe~\ref{chap:annexe_lean4} sans modifier la
borne. Parall\`element,
exhiber la r\'eduction UC entre SOTP-High et ZK-NR~v3, ce qui fermerait
formellement la boucle entre la th\'eorie abstraite de
SOTP~(section~\ref{sec:sotp_construction}) et son instanciation
composable~\cite{MINKA2025ZKNR_UC,Canetti01}. Puis, prouver un
th\'eor\`eme de composition garantissant que l'assemblage
$(CEE \;\|\; AOW \;\|\; SH)$ en Gold-Iron-Clay pr\'eserve les bornes CRO
\`a un facteur constant. Enfin, ex\'ecuter le
programme exp\'erimental sp\'ecifi\'e au
chapitre~\ref{chap:experimentation}.

Les six protocoles y sont d\'ecrits jusqu'au niveau op\'erationnel (stratification des \'echantillons, estimateurs, crit\`eres de fiabilit\'e
inter-\'evaluateurs, seuils de falsification) mais aucun n'a \'et\'e
d\'eploy\'e (sous-section~\ref{subsubsec:statut_deploiement}). Tant que
$\Delta$, $H(\mathcal{C})$, $\Delta_q$ et $\eta_q$ demeurent estim\'es
par triangulation documentaire, sans mesure aupr\`es d'acteurs
institutionnels r\'eels, les bornes CRO conservent le statut de
r\'esultats conditionnels, puisqu'elles sont math\'ematiquement \'etablies sous
hypoth\`eses, sans calibration empirique. Le d\'eploiement suppose trois pr\'ealables non
techniques, \`a savoir l'accord institutionnel des juridictions et r\'egulateurs
approch\'es, l'approbation d'un comit\'e d'\'ethique de la recherche et un
financement sur un horizon de vingt-six mois. L'audit post-\'electoral
Helios de la section~\ref{subsec:protocole_opp}, le plus l\'eger en
ressources et le seul enti\`erement internalisable \`a l'ENSP, constitue
le point d'entr\'ee naturel, puisqu'il valide $\Opp$ sur un syst\`eme
r\'eellement d\'eploy\'e et fournit le premier point de calibration du
cadre.

\subsubsection{Horizon moyen terme (3--5 ans)~: \'elargissement du domaine}
\label{subsubsec:horizon_moyen_terme}

La th\'eorie CLO computationnelle, d\'efinitions de
$(\Conf_\text{comp}, \Rel_\text{comp}, \Opp_\text{comp})$ via adversaires
PPT~\cite{Goldreich01,Bellare97} et \'etablissement du trilemme dans ce cadre
renforc\'e, est la priorit\'e th\'eorique centrale de cet horizon. En
parall\`ele, les variantes du Trilemme CRO pour FHE et
MPC~\cite{Gentry09,Evans2018} doivent \^etre d\'emontr\'ees ou r\'efut\'ees~;
les bornes d'opposabilit\'e atteignables dans ces extensions restent
inconnues. Des outils d'estimation de $H(\mathcal{C})$ et $\Delta$ par
apprentissage automatique peuvent \^etre d\'evelopp\'es et calibr\'es sur des
preuves existantes dont les m\'etriques CLO ont \'et\'e estim\'ees
analytiquement ($0{,}437$ pour Helios, $0{,}572$ pour zk-SNARKs et $0{,}273$ pour ECDSA), sous
r\'eserve que les estimations analytiques constituent une r\'ef\'erence
raisonnablement proche des mesures directes. La contribution \`a la
standardisation, int\'egration de l'opposabilit\'e dans les standards NIST
post-quantiques, cadrage d'un futur standard ISO/IEC, peut d\'ebuter sur cet
horizon.

\subsubsection{Horizon long terme (5--10 ans)~: sous-domaine \'etabli}
\label{subsubsec:horizon_long_terme}

La vision d\'ecennale est celle d'un sous-domaine reconnu dans les
conf\'erences de cryptologie~(CRYPTO, EUROCRYPT, ASIACRYPT), avec sa propre
journ\'ee de travail sp\'ecialis\'ee, des programmes doctoraux d\'edi\'es dans
les universit\'es de recherche ainsi que des venues
propres~(\textit{International Symposium on Institutional Cryptography},
\textit{Journal of Institutional Cryptography}). Scientifiquement, l'horizon
d\'ecennal vise la caract\'erisation compl\`ete de $\mathcal{R} \subset [0,1]^3$,
fronti\`ere de Pareto exacte comprise, une th\'eorie unifi\'ee couvrant
protocoles non interactifs, interactifs et multi-parties sous un cadre
CLO-CRO \'etendu. Si AIIP r\'esiste \`a la cryptanalyse, cet horizon inclut aussi des primitives
CASH soumises \`a standardisation.

\begin{table}[htbp]
\centering
\caption{Synth\`ese des perspectives par contribution et horizon}
\label{tab:perspectives_horizons}
\small
\begin{tabular}{@{}p{2.5cm}p{3cm}p{3cm}p{3cm}@{}}
\toprule
\textbf{Contribution} & \textbf{Court terme} & \textbf{Moyen terme} & \textbf{Long terme} \\
\midrule
C1 (CLO) & D\'eploiement du protocole (chap.~\ref{chap:experimentation}, non ex\'ecut\'e) & Outils ML d'estimation & Standards ISO/IEC \\
C2 (CRO classique) & Extension Lean4 \`a SOTP et PQ & Extensions FHE/MPC & Caract\'erisation $\mathcal{R}$ \\
C3 (CRO post-quantique) & Mesure de $\Delta_q$ (ML-DSA, SLH-DSA) & Th\'eorie $\eta_q$ (options D, E) & PQC certifi\'e CLO \\
C4 (SOTP) & Mod\`ele \`a contexte priv\'e (B2), R\'eduction SOTP$\to$ZK-NR & Preuve ou r\'efutation de la conjecture d'optimalit\'e & Th\'eorie unifi\'ee, impl\'ementations de r\'ef\'erence \\
C5 (Lean4) & Extension couverture SOTP et PQ & Certification UC & V\'erificateur automatique, int\'egration Mathlib et Coq/Isabelle \\
C6 (Protocoles) & D\'eploiement Helios audit post-\'electoral & Validation multi-institutionnelle & Standards empiriques \\
\bottomrule
\end{tabular}
\end{table}

\subsubsection{Financements envisag\'es}
\label{subsubsec:financements_envisages}

\noindent\textbf{Agences nationales.} MINRESI (Cameroun), ANR (France).

\noindent\textbf{Agences internationales.} NSF SaTC (\'Etats-Unis), IDRC
(Canada), Horizon Europe (EU).

\noindent\textbf{Partenariats institutionnels.} LIMSI (Cameroun), NIST (PQC
validation), IBM Quantum (acc\`es simulateur).

\subsubsection{Suivi longitudinal}
\label{subsubsec:suivi_longitudinal}

Les travaux ant\'erieurs~\cite{MINKA2026CRO} et la
section~\ref{sec:cro_pq} invitent \`a
mesurer l'\'evolution des param\`etres \`a mesure que les institutions
acqui\`erent de l'exp\'erience, selon les lois
\begin{align}
\Delta(t) &= \Delta_\infty(1 - e^{-t/\tau})
\label{eq:suivi_delta} \\
\Delta_q(t) &= \Delta_{q,\infty}(1 - e^{-t/\tau_q})
\label{eq:suivi_delta_q} \\
\eta_q(t) &= \eta_{q,0} \cdot e^{-t/\tau_{\mathrm{hw}}}
\label{eq:suivi_eta_q}
\end{align}
o\`u $\tau$ est la constante de temps d'apprentissage institutionnel et
$\tau_{\mathrm{hw}}$ refl\`ete l'am\'elioration du mat\'eriel quantique. Ce
programme longitudinal repr\'esente une direction de recherche \`a 3--5 ans.
\end{sloppypar}

\begin{boiteArbitre}
L'arbitre retient que le voyage ne s'arr\^ete pas ici. Les bornes existent
sous hypoth\`eses, tandis que les constructions ont des propri\'et\'es prouv\'ees et une
opposabilit\'e encore heuristique. Une borne de terme dominant
$O(\lambda^{-2})$ reste ouverte, de m\^eme que les extensions au calcul homomorphe
et multi-parties~; les param\`etres quantiques
$\Delta_q$ et $\eta_q$ restent quant \`a eux \`a calibrer sur des acteurs r\'eels. Ce qui a
\'et\'e construit est prouv\'e l\`a o\`u il est dit prouv\'e et conjectur\'e l\`a o\`u il est dit conjectur\'e.
Ce qui reste \`a construire est plus vaste encore, mais une carte existe d\'esormais pour s'y orienter.
Un sous-domaine est en marche avec ses questions ouvertes, ses hypoth\`eses
\`a valider et sa communaut\'e qui commence \`a se constituer autour d'elles.
\end{boiteArbitre}
% ============================================================================
% CONCLUSION DU CHAPITRE 4 
% source    : partie4/conclusion_partie3.tex
% ============================================================================

\section{CONCLUSION DU CHAPITRE~\ref{chap:resultats_discussions}}
\label{sec:conclusion_resultats_discussions}

\begin{sloppypar}
Ce chapitre a soumis le cadre th\'eorique
construit aux chapitres~\ref{chap:cadre_formel} et~\ref{chap:constructions} \`a l'\'epreuve de la coh\'erence
interne et de l'interpr\'etation. Quatre r\'esultats structurent le bilan.

Le premier est une calibration. Les ordres de grandeur retenus pour
les syst\`emes \'etudi\'es sont mutuellement coh\'erents et plausibles au regard
de la litt\'erature cognitive. Cette coh\'erence ne vaut pas validation, car les
valeurs ont \'et\'e fix\'ees par calibration r\'etrospective et les syst\`emes
d\'eploy\'es ne sont pas des instances du Trilemme
(Remarque~\ref{rem:portee_cas_etude}). De plus, les produits post-quantiques
illustratifs d\'epassent la majoration.

Le deuxi\`eme est \'epist\'emologique. La formalisation Lean4 et la relecture
des d\'efinitions ont conduit \`a pr\'eciser quatre points, \`a savoir la n\'ecessit\'e du r\'egime de
confidentialit\'e, le rapport entre la formulation initiale de la borne et sa
formulation certifi\'ee, la distinction entre fiabilit\'e de reconstruction et
fiabilit\'e de v\'erification ainsi que le caract\`ere non r\'ealisable du t\'emoin de
coh\'erence. Ces pr\'ecisions d\'elimitent la port\'ee des r\'esultats et montrent
l'int\'er\^et de combiner certification m\'ecanique et examen des d\'efinitions.

Le troisi\`eme est structurel. Sous ses hypoth\`eses, le Trilemme CRO
formalise la tension que le paradoxe de l'authenticit\'e invisible
d\'ecrivait. Une opposabilit\'e \'elev\'ee exige que la preuve informe le verdict,
ce qui, pour un interpr\`ete admissible, passe par de l'information sur le
message ou par la capacit\'e $\Delta$, alors que l'information sur le message d\'egrade la
confidentialit\'e. La forme raffin\'ee pr\'ecise quel facteur limite
l'opposabilit\'e selon le r\'egime. L'extension post-quantique montre que la
structure de la borne est la m\^eme pour des preuves quantiques. Elle montre aussi que le
terme de perte $\eta_q$ est nul pour les signatures post-quantiques
actuelles, qui sont des cha\^{\i}nes de bits.

Le quatri\`eme est prospectif. La cryptographie institutionnelle dispose d'un
formalisme, d'une borne prouv\'ee sous hypoth\`eses, de constructions et d'une
certification m\'ecanique, mais il lui manque des mesures. Quatre directions
restent ouvertes, en l'occurrence la calibration empirique de $\Delta$ et $\Delta_q$ par des
\'etudes contr\^ol\'ees, avec les crit\`eres de r\'efutation fix\'es au
chapitre~\ref{chap:experimentation}~; l'extension du cadre aux protocoles
interactifs et multi-parties~; la preuve ou la r\'efutation de la conjecture
d'optimalit\'e des constructions SOTP~; enfin la compl\'etion de la couverture
formelle Lean4. Ces ouvertures d\'elimitent le p\'erim\`etre de ce travail et
tracent le programme d'un sous-domaine dont le d\'eveloppement commence.

La conclusion g\'en\'erale dresse le bilan synth\'etique de l'ensemble.
\end{sloppypar}



% -----------------------------------------------------------------------------
% CONCLUSION GÉNÉRALE
% -----------------------------------------------------------------------------
\chapitreNonNumerote{CONCLUSION G\'EN\'ERALE}{CG}
% ============================================================================
% CONCLUSION G\'eN\'eRALE
% source : intro_conclu_gen/conclusion_generale.tex
% Version : 2.2 — alignée sur la formalisation Lean4 révisée
% ============================================================================

%\chapter*{CONCLUSION G\'eN\'eRALE}
%\addcontentsline{toc}{chapter}{Conclusion g\'en\'erale}
\label{chap:conclusion_generale}

% ----------------------------------------------------------------------------
\section{RETOUR SUR LES OBJECTIFS DE RECHERCHE}
\label{sec:cg_questions}
% ----------------------------------------------------------------------------
\begin{sloppypar}

Le \textbf{paradoxe de l'authenticit\'e invisible}~\cite{MINKA2026CLO}, selon lequel une preuve
cryptographique math\'ematiquement irr\'eprochable peut demeurer institutionnellement
opaque, formalise une tension latente entre les exigences cryptographiques et
l'utilit\'e institutionnelle qui passe par l'interpr\'etabilit\'e des preuves par
les arbitres (juge, r\'egulateur, auditeur) non-experts. Nous nous sommes servi
d'exemples illustratifs~\cite{IACR2025incident,Adida09UCL} d\'ecrits \`a la
section~\ref{subsec:cas_etude} pour mat\'erialiser la r\'ealit\'e de cette
situation et en faciliter la compr\'ehension. Ce paradoxe g\'en\`ere la
probl\'ematique centrale formul\'ee ci-dessous.

\begin{boiteProblematique}{Probl\'ematique centrale}
\textit{Comment, d'une part, concilier la s\'ecurit\'e cryptographique qui
requiert l'opacit\'e computationnelle avec l'interpr\'etabilit\'e
institutionnelle qui requiert la transparence s\'emantique~? D'autre part,
existe-t-il des limites fondamentales aux compromis r\'ealisables entre
confidentialit\'e, fiabilit\'e et opposabilit\'e~?}
\end{boiteProblematique}

Cette th\`ese y r\'epond en traitant les six objectifs qui en
d\'ecoulent, avec le statut propre \`a chacun.

\medskip

\textbf{Obj~\ref{obj:obj1_unification}. Formalisation.} Peut-on formaliser
l'utilit\'e institutionnelle d'un protocole cryptographique~? Nous avons
d\'efini trois mesures issues de la th\'eorie de
l'information~\cite{Shannon48,CoverThomas06}, \`a savoir $\Conf$ (statistique), la
fiabilit\'e distingu\'ee en $\Relrec$ (reconstruction) et $\Relver$
(v\'erification), enfin $\Opp = I(V;\Sigma|C)/H(V|C)$, dont l'in\'egalit\'e
$I(V;\Sigma|C) \leq \min\{H(\Sigma|C),H(V|C)\} \leq \alpha H(\mathcal{C})$ formalise
le goulot contextuel.
L'opposabilit\'e $\Opp$ \'etait absente de tout cadre cryptographique
ant\'erieur.

\textbf{Obj~\ref{obj:obj2_impossibilite}. Bornes.} Existent-ils des limites
fondamentales sur les compromis r\'ealisables~? Sous quatre hypoth\`eses
institutionnelles explicites, nous avons \'etabli la borne
\begin{equation}
  \Conf \cdot \Relrec \cdot \Opp \;\leq\;
  \frac{1}{\gamma}\left(\frac{\Delta}{\lambda} + \frac{\kappa(\Delta+1)}{\lambda^2} + \frac{\kappa^2}{\lambda^3}\right)
\end{equation}
o\`u $\gamma = \alpha H(C) - 1 \geq 1$ et $\kappa$ est la constante du
r\'egime de pr\'eservation de la confidentialit\'e
(Hypoth\`ese~\ref{ass:regime}). Le terme dominant de cette borne est
$\Delta/(\gamma\lambda)$~; les autres sont en $O(\lambda^{-2})$. La forme raffin\'ee
\eqref{eq:trilemme_cro_raffine} identifie le r\'egime ($H(\mathcal{C}) > H^\sharp$) o\`u la
capacit\'e d'extraction institutionnelle est la contrainte liante, la puissance
de calcul de l'adversaire passant au second plan. La formulation
publi\'ee~\cite{MINKA2026CRO} s'obtient comme cas particulier.

\textbf{Obj~\ref{obj:obj3_postquantique}. Extension quantique.} Les bornes CRO
se pr\'eservent-elles face aux adversaires quantiques~? Nous avons \'etabli
le Trilemme CRO post-quantique
\begin{equation}
  \Conf^q \cdot \Relrec^q \cdot \Opp^q \;\leq\;
  \frac{1}{\gamma_q}\left(\frac{\Delta_q}{\lambda} + \frac{\kappa(\Delta_q+1)}{\lambda^2} + \frac{\kappa^2}{\lambda^3}\right)
\end{equation}
o\`u $\Delta_q$ est la capacit\'e d'extraction quantique et $\eta_q$ quantifie
la perte d'interpr\'etabilit\'e, nulle pour les preuves qui sont des cha\^{\i}nes
de bits classiques, dont les signatures ML-DSA et SLH-DSA.

\textbf{Obj~\ref{obj:obj4_saturation}. Optimalit\'e.} Peut-on approcher ces
bornes~? La famille SOTP propose une construction par r\'egime, SOTP-Low,
SOTP-Base (projection Johnson-Lindenstrauss) et SOTP-High (Groth16). Leur
confidentialit\'e et leur fiabilit\'e de v\'erification sont prouv\'ees sous
des hypoth\`eses standard~; leur opposabilit\'e est guid\'ee par une heuristique
de conception, tandis que leur optimalit\'e demeure conjecturale
(Conjecture~\ref{conj:optimalite_universelle}). Cet objectif est donc
partiellement atteint.

\textbf{Obj~\ref{obj:obj5_verification}. V\'erification.} Comment certifier la
correction d'un th\'eor\`eme complexe~? La formalisation Lean4 a
certifi\'e m\'ecaniquement le noyau alg\'ebrique du th\'eor\`eme
(section~\ref{sec:lean4}), ainsi que la forme raffin\'ee de la borne. Fano et
Pinsker y sont pos\'es comme propri\'et\'es des instances de protocole~; les
hypoth\`eses institutionnelles y sont aussi explicit\'ees. La formalisation a produit
deux t\'emoins, l'un de coh\'erence (un point du mod\`ele abstrait o\`u la
borne est atteinte, non r\'ealisable par un protocole), l'autre de port\'ee
(une instance hors r\'egime atteint $1/4$). Le d\'etail de ce que la certification a rendu explicite, hypoth\`eses,
constantes et r\'egime de validit\'e, est pr\'esent\'e \`a la
sous-section~\ref{subsec:lean4_erreurs}.

\textbf{Obj~\ref{obj:obj6_experimentation}. Exp\'erimentation.} Comment rendre
ces bornes testables~? Six protocoles op\'erationnels dotent chacun des
param\`etres institutionnels d'une d\'efinition mesurable, d'un estimateur et
d'un plan d'\'echantillonnage. Ils fixent aussi les seuils \`a partir desquels le
mod\`ele serait r\'efut\'e (tableau~\ref{tab:criteres_refutation}). Une
calibration illustrative sur des syst\`emes d\'eploy\'es v\'erifie que les ordres de
grandeur retenus sont mutuellement coh\'erents~; ces syst\`emes ne sont pas des
instances du th\'eor\`eme. Le programme exp\'erimental lui-m\^eme n'a d'ailleurs pas \'et\'e
d\'eploy\'e dans le cadre de cette th\`ese.

\medskip
\`A la question de savoir si confidentialit\'e, fiabilit\'e et opposabilit\'e
peuvent \^etre simultan\'ement maximis\'ees, la th\`ese r\'epond par la n\'egative
sous quatre hypoth\`eses institutionnelles explicites. Le produit y est born\'e
par une quantit\'e dont le terme dominant est $\Delta/(\gamma\lambda)$, tandis que la
forme raffin\'ee identifie exactement le r\'egime ($H(\mathcal{C}) > H^\sharp$) o\`u la
capacit\'e d'extraction institutionnelle est la contrainte liante. Le noyau
alg\'ebrique de ce r\'esultat est certifi\'e en Lean4 et sa port\'ee est
d\'elimit\'ee, puisque les syst\`emes d\'eploy\'es \`a haute fiabilit\'e de
v\'erification n'en sont pas des instances et servent \`a calibrer les
ordres de grandeur. L'existence de constructions atteignant la borne est
formul\'ee comme conjecture et constitue, avec la mesure exp\'erimentale de
$\Delta$ par l'estimateur $\hat{I}(V;\Sigma|M,C)$, le programme ouvert par ce travail.

\end{sloppypar}

% ----------------------------------------------------------------------------
\section{PRINCIPALES CONTRIBUTIONS}
\label{sec:cg_contributions}
% ----------------------------------------------------------------------------
\begin{sloppypar}

Le cycle Obj~\ref{obj:obj1_unification}~$\to$~Obj~\ref{obj:obj2_impossibilite}~$\to$~Obj~\ref{obj:obj3_postquantique}~$\to$~Obj~\ref{obj:obj4_saturation}~$\to$~Obj~\ref{obj:obj5_verification}~$\to$~Obj~\ref{obj:obj6_experimentation} suivant la progression
F$\to$B$\to$E$\to$S$\to$C$\to$X
constitue le squelette logique de la th\`ese. Six contributions
techniques (C1--C6) et trois apports transversaux en \'emergent.

\subsection*{Contributions techniques (C1--C6)}

\begin{table}[htbp]
\centering
\caption{Synth\`ese des six contributions}
\label{tab:synthese_contributions}
\small
\begin{tabular}{@{}llp{6cm}@{}}
\toprule
\textbf{\'Etape} & \textbf{Contribution} & \textbf{R\'esultat cl\'e} \\
\midrule
F (Formaliser) & C1~: Cadre CLO & M\'etriques $\Conf$, $\Relrec$/$\Relver$, $\Opp$~; goulot contextuel \\
B (Borner) & C2~: Trilemme CRO classique & Borne $\gamma^{-1}(\Delta/\lambda + \kappa(\Delta+1)/\lambda^2 + \kappa^2/\lambda^3)$ et forme raffin\'ee \\
E (\'Etendre) & C3~: Trilemme CRO post-quantique & Borne avec $\Delta_q$, $\eta_q$~; calibration NIST PQC \\
S (Saturer) & C4~: Famille SOTP & Trois constructions ($\Conf$, $\Relver$ prouv\'ees)~; optimalit\'e conjectur\'ee \\
C (Certifier) & C5~: Certification Lean4 & Noyau v\'erifi\'e, t\'emoins de coh\'erence et port\'ee \\
X (eXp\'erimenter) & C6~: Programme de validation & Protocoles exp\'erimentaux sp\'ecifi\'es, falsifiables \\
\bottomrule
\end{tabular}
\end{table}

\textbf{C1. Cadre CLO}~\cite{MINKA2026CLO}. Premier cadre unifi\'e
formalisant $\Conf$, $\Rel$ et $\Opp$ via la th\'eorie de l'information, avec
la distinction entre fiabilit\'e de reconstruction et de v\'erification.
Le th\'eor\`eme du goulot d'\'etranglement contextuel \'etablit que
$I(V;\Sigma|C) \leq \alpha H(\mathcal{C})$~; un verdict d\'eterministe donne $\Opp = 1$.

\textbf{C2. Trilemme CRO conditionnel}~\cite{MINKA2026CRO}. Premi\`ere
borne prouv\'ee sur les compromis $(\Conf, \Rel, \Opp)$, conditionnelle
aux param\`etres $(\Delta, \alpha, \beta)$ et au r\'egime de
confidentialit\'e $I(M;\Sigma|C) \leq \kappa/\lambda$. La forme explicite
est $\gamma^{-1}(\Delta/\lambda + \kappa(\Delta+1)/\lambda^2 + \kappa^2/\lambda^3)$~;
la forme raffin\'ee, fond\'ee sur l'identit\'e
$I(V;\Sigma|C) = I(V;M|C) + I(V;\Sigma|M,C)$ pour un interpr\`ete admissible, distingue
un r\'egime o\`u Fano limite l'opposabilit\'e et un r\'egime o\`u $\Delta$ la limite.
La finesse de la borne reste une question ouverte.

\textbf{C3. Trilemme CRO post-quantique}~\cite{MINKA2026PQCRO}. Extension
aux adversaires quantiques avec param\`etres $\Delta_q$ et $\eta_q$.
Calibration illustrative sur ML-DSA et SLH-DSA, pour lesquelles $\eta_q = 0$.

\textbf{C4. Famille SOTP}. Trois constructions, une par r\'egime~;
confidentialit\'e et fiabilit\'e de v\'erification prouv\'ees sous PRG, DLog et
q-PKE~; opposabilit\'e guid\'ee par une heuristique de conception~; optimalit\'e
conjectur\'ee. Introduction du paradigme Gold-Iron-Clay.

\textbf{C5. Formalisation Lean4}. Certification m\'ecanis\'ee
du noyau alg\'ebrique et de la forme raffin\'ee. Fano et Pinsker pos\'es comme
propri\'et\'es d'instance, explicitation des hypoth\`eses institutionnelles
et production de deux t\'emoins, l'un de coh\'erence (mod\`ele abstrait), l'autre de
port\'ee. Les formulations alternatives issues de
la certification sont pr\'esent\'ees \`a la sous-section~\ref{subsec:lean4_erreurs}.

\textbf{C6. Programme de validation falsifiable}. Les protocoles
exp\'erimentaux permettant de mesurer $\Delta$, $\alpha$, $\beta$,
$\Delta_q$ et $\eta_q$ sont sp\'ecifi\'es jusqu'au niveau
op\'erationnel~(chapitre~\ref{chap:experimentation}).

\subsection*{Apports transversaux}

\textbf{Apport constructif (de la borne au guide).} La famille SOTP
transforme le Trilemme CRO en guide de conception, sans \'etablir que la
borne est atteinte, question qui fait l'objet de la
Conjecture~\ref{conj:optimalite_universelle}. Comme CAP trouve son sens dans
les syst\`emes distribu\'es r\'eels et FLP dans les protocoles partiels, le
Trilemme CRO trouve un prolongement dans les trois variantes
SOTP~(section~\ref{subsec:portee_constructive}).

\textbf{Apport m\'ethodologique (la le\c{c}on de la formalisation).}
La formalisation Lean4 a conduit \`a des formulations plus pr\'ecises de trois
\'etapes des preuves publi\'ees (sens d'une in\'egalit\'e, statut de
l'Hypoth\`ese~2, r\'esolution d'une optimisation). Elle a \'etabli par un cas
limite la n\'ecessit\'e du r\'egime de confidentialit\'e. La le\c{c}on, selon
laquelle la formalisation gagne \`a pr\'ec\'eder la publication, est d\'evelopp\'ee
\`a la sous-section~\ref{subsec:verification_corrige_lecon}. Les formulations
alternatives sont pr\'esent\'ees \`a la
sous-section~\ref{subsec:lean4_erreurs}.

\textbf{Apport \'epist\'emique (un programme falsifiable).} Le
chapitre~\ref{chap:experimentation} sp\'ecifie comment tester le
cadre, transformant un r\'esultat conditionnel en programme scientifique
falsifiable au sens de Popper~\cite{Popper1934}. Sans cette
sp\'ecification, le Trilemme CRO serait non testable~; avec elle, il est
testable et critiquable~(section~\ref{subsec:statut_protocoles_experimentaux}).

\end{sloppypar}

% ----------------------------------------------------------------------------
\section{POSITIONNEMENT DANS LA LITT\'ERATURE}
\label{sec:cg_positionnement}
% ----------------------------------------------------------------------------
\begin{sloppypar}

Quatre corpus ont abord\'e des aspects partiels du probl\`eme, \`a savoir la
QIF~\cite{Smith09}, la s\'ecurit\'e utilisable~\cite{Bonneau12}, les \'etudes
juridiques~\cite{Kerr04} et les trilemmes d'impossibilit\'e
(CAP~\cite{Brewer00,GilbertLynch02}, FLP~\cite{FLP85},
Blockchain~\cite{Nakamoto08,Buterin17}). Le cadre CLO-CRO s'en distingue
en ce qu'il formalise simultan\'ement les trois dimensions
$(\Conf, \Rel, \Opp)$, prouve une borne quantitative conditionnelle
classique et quantique, propose des constructions dont l'optimalit\'e est
conjectur\'ee et certifie m\'ecaniquement une partie du r\'esultat.

\begin{table}[htbp]
\centering
\caption{Positionnement par rapport aux travaux existants.}
\label{tab:positionnement_travaux}
\resizebox{\textwidth}{!}{%
\begin{tabular}{@{}lllllll@{}}
\toprule
\textbf{Cadre} & \textbf{Conf} & \textbf{Rel} & \textbf{Opp} & \textbf{Borne} & \textbf{Optimalit\'e} & \textbf{Validation} \\
\midrule
QIF                    & Total     & Total     & Absent      & Partielle & --- & --- \\
S\'ecurit\'e utilisable    & Partielle & Absent    & Qualitative & Absent    & --- & Empirique \\
\'Etudes juridiques      & Absent    & Absent    & Qualitative & Absent    & --- & --- \\
CAP/FLP/Blockchain     & Partielle & Partielle & Absent      & Total     & Oui & --- \\
\textbf{CLO-CRO (ce travail)} & \textbf{Total} & \textbf{Total} & \textbf{Total} & \textbf{Oui (conditionnel)} & \textbf{Partiel (conjecture)} & \textbf{Sp\'ecifi\'e, non d\'eploy\'e} \\
\bottomrule
\end{tabular}}

\vspace{0.4em}
\footnotesize
\textit{\'Echelle de couverture.} Total (traitement formel complet)~;
Partielle (traitement formel restreint \`a certains cas)~; Qualitative
(traitement non formalis\'e)~; Absent (dimension non trait\'ee).
\end{table}

\end{sloppypar}

% ----------------------------------------------------------------------------
\section{PORT\'EE DE LA TH\`ESE}
\label{sec:cg_portee}
% ----------------------------------------------------------------------------
\begin{sloppypar}

La port\'ee de cette th\`ese est quadruple.

\textbf{Th\'eorique.} Les bornes C1--C3 \'etablissent des limites
fondamentales conditionnelles \`a des param\`etres institutionnels
mesurables.

\textbf{Constructive.} La famille SOTP (C4) propose, pour chaque r\'egime, une
construction aux propri\'et\'es cryptographiques prouv\'ees~; qu'elle approche
les bornes \`a facteur constant est conjectur\'e.

\textbf{Certifi\'ee.} La formalisation Lean4 (C5) certifie
m\'ecaniquement le noyau alg\'ebrique du Trilemme CRO classique.

\textbf{M\'ethodologique.} Le programme de validation (C6) et la
le\c{c}on de la formalisation
(section~\ref{subsec:verification_corrige_lecon}) fournissent \`a la
communaut\'e des outils et des m\'ethodes transf\'erables.

\end{sloppypar}

% ----------------------------------------------------------------------------
\section{LIMITATIONS ET TRAVAUX FUTURS}
\label{sec:cg_limitations}
% ----------------------------------------------------------------------------
\begin{sloppypar}

Quatre limitations d\'elimitent le p\'erim\`etre des contributions. Les
hypoth\`eses institutionnelles $(\Delta, \alpha, \beta)$ ne sont pas
encore valid\'ees par des \'etudes exp\'erimentales~; les valeurs retenues,
de quelques dizaines de bits, sont une calibration illustrative rendue
plausible par la litt\'erature sur la rationalit\'e born\'ee. L'opposabilit\'e
des constructions SOTP est heuristique et leur optimalit\'e conjectur\'ee. La formalisation Lean4 couvre le noyau
alg\'ebrique~; les constructions SOTP et l'extension post-quantique restent
hors de son p\'erim\`etre. Les constructions SOTP disposent d'impl\'ementations
de r\'ef\'erence dans le d\'ep\^ot compagnon, sans d\'eploiement en conditions
r\'eelles. Une
analyse compl\`ete des limitations est pr\'esent\'ee \`a la
section~\ref{subsec:limitations_problemes_ouverts}.

Les travaux futurs par contribution et par horizon sont d\'etaill\'es au
tableau~\ref{tab:perspectives_horizons}.

Le programme futur s'organise selon trois directions structurantes.

\textbf{Validation empirique.} Les param\`etres $\Delta$, $\alpha$,
$\beta$, $\Delta_q$, $\eta_q$ doivent \^etre mesur\'es sur des acteurs
institutionnels r\'eels. Le programme du
chapitre~\ref{chap:experimentation} fournit la feuille de route.

\textbf{Extension de la formalisation.} La formalisation Lean4 doit
\^etre \'etendue \`a SOTP et au cadre post-quantique, une fois les
biblioth\`eques de th\'eorie de l'information quantique disponibles dans
Mathlib~(section~\ref{subsubsec:extension_lean4_sotp_pq}).

\textbf{Extension th\'eorique.} Les variantes du Trilemme CRO pour le
chiffrement homomorphe, le calcul multi-parties et les contextes
dynamiques restent \`a d\'evelopper~(section~\ref{subsec:axe3_extensions}).

\end{sloppypar}

% ----------------------------------------------------------------------------
\section{SIGNIFICATION}
\label{sec:cg_signification}
% ----------------------------------------------------------------------------
\begin{sloppypar}

La cryptographie classique s'attache \`a la probabilit\'e qu'un
adversaire casse un protocole, alors que la cryptographie institutionnelle
s'interroge sur la probabilit\'e qu'un juge comprenne une
preuve. Loin d'\^etre une m\'etrique de confort
utilisateur, l'opposabilit\'e $\Opp$ traduit une contrainte propre \`a la nature de l'information,
de m\^eme rang que la confidentialit\'e. Le Trilemme CRO dit que,
sous ses hypoth\`eses, certains compromis sont impossibles, ind\'ependamment
du talent des ing\'enieurs.

La certification Lean4 illustre un pari \'epist\'emologique selon lequel les
th\'eor\`emes importants m\'eritent une v\'erification m\'ecanique.

La cryptographie institutionnelle dispose d\'esormais d'un cadre
compos\'e de m\'etriques (CLO), d'une borne d'impossibilit\'e conditionnelle classique et
quantique (CRO), de constructions (SOTP), d'une v\'erification formelle (Lean4)
et d'un programme de validation (C6). Le travail restant (finesse de la borne,
optimalit\'e des constructions, calibration des param\`etres,
extensions FHE/MPC, validation empirique, standardisation) constitue un
programme de recherche ouvert.

La cryptographie n'existe pas dans le vide. Elle sert
des institutions humaines soumises \`a des contraintes cognitives,
organisationnelles et proc\'edurales. L'ignorance de ces
contraintes engendre le paradoxe de l'authenticit\'e invisible.

\end{sloppypar}

% ----------------------------------------------------------------------------
\section*{MOT DE LA FIN}
\label{sec:cg_mot_fin}
% ----------------------------------------------------------------------------
\begin{sloppypar}

Cette th\`ese a transform\'e en cadre quantitatif l'intuition qualitative
selon laquelle les preuves cryptographiques sont opaques pour les interpr\`etes
institutionnels. Ce cadre r\'eunit des m\'etriques d\'efinies, des bornes
prouv\'ees sous hypoth\`eses, des constructions, une certification m\'ecanique
et un programme de validation falsifiable. Dans le r\'egime que la forme
raffin\'ee identifie, la contrainte dominante devient la capacit\'e d'extraction de l'interpr\`ete, la
puissance de calcul de l'adversaire passant au second plan.
Cette inversion de perspective est l'id\'ee centrale de cette th\`ese~; sa
confirmation empirique reste \`a apporter.

L'arbitre le rappelait, \emph{«~La rigueur ne se suppose pas. Elle se
v\'erifie.~»} Cette th\`ese a v\'erifi\'e et pr\'ecis\'e ses r\'esultats, document\'e
ces pr\'ecisions et sp\'ecifi\'e comment tester ce qui reste \`a
\'etablir. En cela, la science se distingue de la sp\'eculation.

\end{sloppypar}

% ----------------------------------------------------------------------------
% Bo\^ite finale
% ----------------------------------------------------------------------------
\begin{boiteArbitre}
L'arbitre referme le dossier sans avoir r\'esolu le paradoxe, puisque le Trilemme
CRO \'etablit qu'aucun protocole ne peut le r\'esoudre enti\`erement, que
ce soit en environnement classique~\cite{MINKA2026CRO} ou
quantique~\cite{MINKA2026PQCRO}. Il dispose toutefois d\'esormais de ce
que la cryptographie classique ne lui avait jamais offert, \`a savoir une
m\'etrique pour mesurer \`a quel point un protocole s'approche de ce qu'il
peut interpr\'eter~\cite{MINKA2026CLO} ainsi qu'une famille de constructions pour
commencer. Le paradoxe de l'authenticit\'e invisible n'est plus invisible
puisqu'il est formalis\'e et born\'e, sous des hypoth\`eses qu'il reste \`a mesurer. Ce que l'arbitre ne comprend pas
encore enti\`erement, la th\'eorie lui dit maintenant combien cela lui
co\^ute en bits.
\end{boiteArbitre}

% -----------------------------------------------------------------------------
% BIBLIOGRAPHIE
% -----------------------------------------------------------------------------
\clearpage
\markboth{R\'EF\'ERENCES BIBLIOGRAPHIQUES}{}
\printbibliography[
  title={R\'EF\'ERENCES BIBLIOGRAPHIQUES},
  heading=bibintoc
]
% =============================================================================
% ANNEXES
% =============================================================================
\appendix
\addtocontents{toc}{\protect\etocsettocdepth{chapter}}
\entreePartie{ANNEXES}     
\debutAnnexes

% ============================================================================
% S\'eCURIT\'e GAME-BASED ET R\'eDUCTIONS FORMELLES
% source : annexes/Gamebased.tex
% ============================================================================

\chapter{S\'ecurit\'e game-based et r\'eductions formelles pour la famille SOTP}
\label{chap:annexe_gamebased}

\begin{sloppypar}
Le cadre CLO d\'efinit $\Conf$, $\Rel$ et $\Opp$ avec des termes issus de la théorie de l'information, \`a savoir les entropies conditionnelles, l'information
mutuelle et la distance de variation totale~\cite{MINKA2026CLO}. Cette approche
permet des r\'esultats g\'en\'eraux (th\'eor\`eme du goulot d'\'etranglement
contextuel, Trilemme CRO conditionnel), mais ne pr\'ecise pas directement
\emph{ce qu'un adversaire calculatoirement born\'e peut ou ne peut pas
faire}. En cryptographie computationnelle, la s\'ecurit\'e s'exprime
classiquement via des \emph{jeux} (games) entre un challenger et un
adversaire probabiliste polynomial-time (PPT).

\begin{cmtEtape}{Pourquoi cette annexe.}
Le cadre CLO raisonne en bits et en entropies, alors que la cryptographie classique
raisonne en jeux et en avantages d'adversaires. Cette annexe construit
le pont entre les deux en capturant chaque m\'etrique du cadre CLO
par un jeu cryptographique, puis en reliant chaque jeu \`a une hypoth\`ese
de duret\'e standard (PRG, DLOG, JL, Groth16).
\end{cmtEtape}

Cette annexe \'etablit le pont formel entre les deux mondes en trois \'etapes.
\begin{enumerate}
\item D\'efinir des jeux cryptographiques (\emph{IND-CONF},
      \emph{BIND-REL}, \emph{EXT-OPP}) capturant op\'erationnellement
      $\Conf$ (dans sa version calculatoire), $\Relver$ et $\Opp$
\item Prouver des \emph{r\'eductions} formelles \'etablissant que, si un adversaire PPT
      gagne l'un de ces jeux, il r\'esout un probl\`eme suppos\'e difficile
      (PRG, logarithme discret, JL, Groth16)
\item D\'eriver une minoration game-based du produit CLO
      (Corollaire~\ref{cor:ann_cro_from_games}), \`a l'\'echelle op\'erationnelle
      (confidentialit\'e calculatoire, fiabilit\'e de v\'erification), qui ne se
      compare pas directement au Trilemme
      (Remarque~\ref{rem:ann_pas_de_contradiction})
\end{enumerate}

\begin{cmtEtape}{Trois jeux, trois m\'etriques.}
\textbf{IND-CONF} capte la confidentialit\'e en exigeant que l'adversaire ne puisse
distinguer les preuves de deux messages. \textbf{BIND-REL} capte la
fiabilit\'e, puisque l'adversaire ne doit pas pouvoir ouvrir une preuve sur
deux messages distincts. Enfin, \textbf{EXT-OPP} capte l'opposabilit\'e,
l'interpr\`ete institutionnel devant extraire le bon verdict depuis la
preuve.
\end{cmtEtape}

\noindent\textbf{Convention.} Tout adversaire $\mathcal{A}$ est
probabiliste polynomial-time (PPT) en $\lambda$ sauf mention contraire.
Le challenger est not\'e $\mathcal{C}$. Une r\'eduction construit un
algorithme $\mathcal{B}$ (r\'educteur) qui exploite $\mathcal{A}$ pour
r\'esoudre un probl\`eme difficile.
\cmt{Selon la convention standard en cryptographie, PPT signifie temps polynomial probabiliste.}
\cmt{Le r\'educteur $\mathcal{B}$ est l'algorithme qui transforme une attaque en solution d'un probl\`eme difficile.}

\begin{table}[htbp]
\centering
\caption{Structure de l'Annexe \ref{chap:annexe_gamebased}}
\label{tab:ann_g_structure}
\small
\begin{tabular}{@{}llll@{}}
\toprule
\textbf{Section} & \textbf{Contenu} & \textbf{Propri\'et\'e}
& \textbf{Hypoth\`ese} \\
\midrule
\ref{sec:gamebased_primitives} & Primitives et d\'efinitions pr\'ealables & -- & -- \\
\ref{sec:gamebased_confidentialite} & Jeu IND-CONF + lien $\Conf$ & $\Conf$ & PRG \\
\ref{sec:gamebased_fiabilite} & Jeu BIND-REL + lien $\Relver$ & $\Relver$ & DLOG \\
\ref{sec:gamebased_opposabilite} & Jeu EXT-OPP + lien $\Opp$ & $\Opp$ & Fano \\
\ref{sec:gamebased_SOTP-Low} & R\'eductions SOTP-Low & $\Conf_{\mathrm{calc}}$, $\Relver$ & PRG, DLOG \\
\ref{sec:gamebased_SOTP-base} & Dimensionnement SOTP-Base & -- & JL \\
\ref{sec:gamebased_SOTP-High} & R\'eductions SOTP-High & $\Conf$, $\Relver$ & Groth16 \\
\ref{sec:gamebased_produitCRO} & \'enonc\'es complets + produit CLO & $\Conf_{\mathrm{calc}}\cdot\Relver\cdot\Opp$ & Tous \\
\ref{sec:gamebased_limitations} & Limitations et port\'ee & -- & -- \\
\bottomrule
\end{tabular}
\end{table}
\end{sloppypar}

% ============================================================================
\section{PRIMITIVES CRYPTOGRAPHIQUES ET D\'EFINITIONS PR\'EALABLES}
\addcontentsline{toc}{section}{Primitives et d\'efinitions pr\'ealables}
\label{sec:gamebased_primitives}
% ============================================================================
\begin{sloppypar}

\begin{cmtEtape}{Objectif de cette section.}
Rappeler les quatre primitives cryptographiques utilis\'ees dans la
famille SOTP, \`a savoir le PRG (couche Gold), l'engagement de Pedersen (couche Iron),
le zk-SNARK Groth16 (SOTP-High) et la transformation Johnson-Lindenstrauss
(SOTP-Base), ainsi que deux outils th\'eoriques, l'in\'egalit\'e de Fano et le DPI.
\end{cmtEtape}

\subsection{G\'en\'erateur pseudo-al\'eatoire (PRG)}
\noindent\textbf{Intuition.} Un PRG prend une graine courte
$s \in \{0,1\}^\lambda$ et produit une longue s\'equence
$G(s) \in \{0,1\}^{L(\lambda)}$, $L > \lambda$, indiscernable
d'une s\'equence vraiment al\'eatoire. Il fonde la couche
Gold de SOTP-Low, dans laquelle le pad $G(s)$ masque le message de sorte que, sans $s$,
aucun adversaire PPT ne peut r\'ecup\'erer $m$ depuis $m \oplus G(s)$.
\cmt{Le PRG \'etire une petite graine en un long pad pseudo-al\'eatoire.}
\cmt{Sans la graine, le pad est indiscernable d'un al\'ea v\'eritable.}

\begin{definition}[Sécurité PRG~\cite{GoldreichKrawczyk99}]\label{def:ann_prg}
Une fonction calculable en temps polynomial
$G\colon \{0,1\}^\lambda \to \{0,1\}^{L(\lambda)}$ avec $L(\lambda) > \lambda$
est un PRG sûr si, pour tout adversaire PPT $\mathcal{A}$, 
\begin{equation}
\mathsf{Adv}^{\mathsf{PRG}}_{\mathcal{A}}(\lambda) =
\left|
\Pr_{s \leftarrow \{0,1\}^\lambda}\!\big[\mathcal{A}(G(s)) = 1\big]
- \Pr_{y \leftarrow \{0,1\}^{L(\lambda)}}\!\big[\mathcal{A}(y) = 1\big]
\right| \leq \negl(\lambda)
\end{equation}
\cmtFormule{L'avantage mesure la capacit\'e \`a distinguer $G(s)$ d'un al\'ea pur.}
\cmt{L'avantage doit \^etre n\'egligeable pour tout adversaire PPT.}
\end{definition}

\noindent\textbf{Impl\'ementations pratiques.} SOTP-Low utilise
ChaCha20/AES-CTR, deux PRG standardis\'es atteignant cette d\'efinition.
\cmt{ChaCha20 et AES-CTR sont des PRG concrets, utilis\'es en pratique.}

\subsection{Engagement de Pedersen}

\noindent\textbf{Intuition.} L'engagement de Pedersen permet \`a un
prouveur de \guillemotleft~s'engager~\guillemotright sur $m$ sans le r\'ev\'eler (\emph{perfect hiding}),
tout en \'etant incapable de changer d'avis (\emph{computational binding}).
La construction $\mathsf{Com}(m, r) = g^m h^r$ exploite l'alg\`ebre des
groupes cycliques, de sorte que toute collision d'engagements permettrait de calculer
$\log_g h$, r\'esolvant le logarithme discret.
\cmt{Le hiding parfait signifie que le message est cach\'e inconditionnellement.}
\cmt{Par le binding computationnel, ouvrir sur deux messages casserait DLOG.}

\begin{definition}[Engagement~de~Pedersen~\cite{Pedersen91}]\label{def:ann_pedersen}
Soit $\mathbb{G}$ un groupe cyclique d'ordre premier $q$,
$g, h \in \mathbb{G}$ deux générateurs tels que $\log_g h$ est inconnu.
\begin{itemize}
\item $\mathsf{Com}(m, r) = g^m h^r \in \mathbb{G}$, $\quad m,r \in \mathbb{Z}_q$ (op\'eration de groupe)
\item $\mathsf{Open}(\sigma_2, m, r)$ accepte si $\sigma_2 = g^m h^r$
\end{itemize}
\textbf{Perfect Hiding.} Pour tout $m_0, m_1$, les engagements
$\mathsf{Com}(m_0, r_0)$ et $\mathsf{Com}(m_1, r_1)$ sont identiquement
distribués pour $r_0, r_1 \xleftarrow{\$} \mathbb{Z}_q$.
\cmt{Le masque $h^r$ est uniforme, ind\'ependant du message.}
\textbf{Computational Binding.} Aucun adversaire PPT ne peut trouver
$(m_0, r_0, m_1, r_1)$ avec $m_0 \neq m_1$ tel que
$g^{m_0} h^{r_0} = g^{m_1} h^{r_1}$ dans $\mathbb{G}$.
\cmt{Deux ouvertures distinctes r\'ev\`elent $\log_g h$.}
\end{definition}

\noindent\textbf{R\^ole dans SOTP.} La couche Iron de toute variante SOTP
est un engagement de Pedersen (sur la graine $s \in \mathbb{Z}_q$ pour SOTP-Low et
SOTP-Base, sur $m$ pour SOTP-High). Il garantit la fiabilit\'e de v\'erification
$\Relver$, puisqu'une preuve ne peut pas \^etre \guillemotleft~ouverte~\guillemotright sur deux messages
distincts.
\cmt{C'est la source de la fiabilit\'e de v\'erification dans toutes les variantes SOTP.}

\subsection{zk-SNARK Groth16}

\noindent\textbf{Intuition.} Un zk-SNARK pour une relation $\mathcal{R}$
permet de prouver qu'on conna\^it un t\'emoin $w$ tel que $(x,w) \in
\mathcal{R}$, sans r\'ev\'eler $w$, en une preuve de taille constante.
Une preuve Groth16 compte trois \'el\'ements de groupe, soit 128 octets (points compress\'es) sur BN254
ou 192 octets sur BLS12-381~; elle se v\'erifie en $O(1)$ couplages.
\cmt{Preuve de taille constante quel que soit le circuit.}
\cmt{La v\'erification se fait en temps constant via des couplages.}

\begin{definition}[zk-SNARK~Groth16~\cite{Groth16}]\label{def:ann_groth16}
$(\mathsf{Setup}, \mathsf{Prove}, \mathsf{Verify})$ satisfait les propri\'et\'es suivantes.

\noindent\textbf{Complétude parfaite.}
$\Pr[\mathsf{Verify}(\sigma_v,x,\pi)=1 \;|\; (x,w)\in\mathcal{R}] = 1$
\cmt{Une preuve honn\^ete est toujours accept\'ee.}

\noindent\textbf{Solidité connaissance.} Pour tout PPT $\mathcal{A}$,
il existe un extracteur $\mathcal{E}$ tel que
\begin{equation}
\Pr\!\left[
\mathsf{Verify}(\sigma_v,x,\pi)=1 \land (x,w)\notin\mathcal{R} :
(x,\pi)\leftarrow\mathcal{A}(\sigma_p),\; w\leftarrow\mathcal{E}(\sigma_p)
\right] \leq \negl(\lambda)
\end{equation}
\cmtFormule{L'extracteur r\'ecup\`ere un t\'emoin valide si la preuve passe.}
\cmt{Cette propri\'et\'e est plus forte que la soundness standard puisqu'on peut extraire le t\'emoin.}

\noindent\textbf{Zéro-connaissance parfaite.}
$\exists$ simulateur $\mathcal{S}$ tel que pour tout $(x,w)\in\mathcal{R}$,
\begin{equation}
\{\pi : \pi \leftarrow \mathsf{Prove}(\sigma_p,x,w)\}
\equiv
\{\pi : \pi \leftarrow \mathcal{S}(\sigma_p,x)\}
\quad\text{(distributions identiques)}
\end{equation}
\cmtFormule{Les vraies preuves et les preuves simul\'ees sont statistiquement identiques.}
\cmt{Cette \'egalit\'e exacte explique $\Conf \geq 1 - \negl$.}
\end{definition}

\noindent\textbf{Lecture pour SOTP-High.} La z\'ero-connaissance
\emph{parfaite}, jointe au masquage parfait de l'engagement $c_m$, implique que
$(\pi, c_m)$ est ind\'ependant de $m$ sachant $c$, d'o\`u $\Conf = 1$ au sens statistique.
La solidit\'e connaissance implique la fiabilit\'e de v\'erification
$\Relver \geq 1 - \negl(\lambda)$~; elle ne dit rien de la fiabilit\'e de
reconstruction, qui vaut $2^{-\lambda}$ lorsque $\Conf = 1$
(Remarque~\ref{rem:rel_ver}).
\cmt{La ZK parfaite donne la confidentialit\'e~; la solidit\'e donne la fiabilit\'e de v\'erification.}

\subsection{Transformation Johnson-Lindenstrauss (JL)}

\noindent\textbf{Intuition.} La JL-transform projette $N$ points de
$\mathbb{R}^n$ dans $\mathbb{R}^d$ ($d \ll n$) en pr\'eservant
approximativement toutes les distances par paires. Dans SOTP-Base,
les contextes $c \in \mathbb{R}^n$ sont compress\'es en $\Phi c \in
\mathbb{R}^d$, conservant la structure permettant l'extraction
institutionnelle.
\cmt{On compresse un espace de grande dimension en pr\'eservant les distances.}
\cmt{Cela permet de r\'eduire la taille de la preuve sans perdre l'information utile.}

\begin{definition}[Transformation~JL~\cite{JohnsonLindenstrauss84,DasguptaGupta03}]\label{def:ann_jl}
Pour $0 < \varepsilon < 1$, $N > 0$, $S \subset \mathbb{R}^n$ avec
$|S| = N$, il existe $\Phi\colon \mathbb{R}^n \to \mathbb{R}^d$ avec
$d = O(\varepsilon^{-2} \log N)$ telle que, pour tout $x,y \in S$,
\begin{equation}
(1-\varepsilon)\|x-y\|^2 \leq \|\Phi x-\Phi y\|^2 \leq (1+\varepsilon)\|x-y\|^2
\end{equation}
\cmtFormule{Les distances sont pr\'eserv\'ees \`a $\varepsilon$ pr\`es.}
\cmt{$d$ d\'epend logarithmiquement de $N$, pas lin\'eairement.}
\end{definition}

\noindent\textbf{Lemme de dimension pour SOTP-Base.}
La dimension minimale $d$ pour pr\'eserver $\Delta$ bits d'extraction
avec tol\'erance $\varepsilon^2$ est donn\'ee par la condition
\begin{equation}
d \geq \frac{4\Delta H(\mathcal{C})}{\varepsilon^2}
\label{eq:ann_jl_dim}
\end{equation}
\cmtFormule{La dimension cro\^it lin\'eairement en $\Delta H(C)$.}
\cmt{C'est le co\^ut de la compression pr\'eservant l'extraction.}

\subsection{Outils th\'eoriques fondamentaux}

\begin{lemma}[Inégalité~de~Fano~\cite{CoverThomas06}]\label{lem:ann_g_fano}
Soit $X$ une variable à $k$ valeurs, $\hat{X}$ un estimateur,
$p_e = \Pr[X \neq \hat{X}]$. Alors
\begin{equation}
H(X|\hat{X}) \leq h(p_e) + p_e \log_2(k-1)
\end{equation}
\cmtFormule{L'incertitude r\'esiduelle est born\'ee par la probabilit\'e d'erreur.}
\end{lemma}

\noindent\textbf{Lecture.} Faible erreur d'extraction $p_e$
$\Rightarrow$ faible incertitude r\'esiduelle $H(X|\hat{X})$
$\Rightarrow$ grande information mutuelle extraite. Ce lemme est le
pont direct entre le jeu EXT-OPP et la m\'etrique $\Opp$.
\cmt{Plus l'interpr\`ete se trompe peu, plus il extrait d'information.}
\cmt{C'est le lien formel entre le jeu EXT-OPP et la m\'etrique $\Opp$.}

\begin{lemma}[Traitement~de~l'Information~(DPI)~\cite{CoverThomas06}]\label{lem:ann_dpi}
Si $X \to Y \to Z$ forme une chaîne de Markov, alors
$I(X;Z) \leq I(X;Y)$.
\cmtFormule{Traiter une donn\'ee ne peut qu'\'eliminer de l'information.}
\end{lemma}

\noindent\textbf{Lecture.} Toute transformation de $Y$ ne peut
qu'augmenter l'incertitude sur $X$. Un adversaire observant $\sigma$
ne peut extraire au plus que $I(M;\Sigma|C)$ bits sur $m$, fondation
du lien entre $\mathsf{Adv}^{\mathsf{IND\text{-}CONF}}$ et $I(M;\Sigma|C)$.
\cmt{Le DPI justifie que la meilleure attaque est born\'ee par $I(M;\Sigma|C)$.}
\end{sloppypar}

% ============================================================================
\section{JEU IND-CONF~: CONFIDENTIALIT\'E GAME-BASED}
\addcontentsline{toc}{section}{Jeu IND-CONF}
\label{sec:gamebased_confidentialite}
% ============================================================================
\begin{sloppypar}

\begin{cmtEtape}{Objectif de cette section.}
D\'efinir le jeu IND-CONF qui capture la confidentialit\'e, puis
\'etablir son lien avec la m\'etrique $\Conf$ du cadre CLO.
\end{cmtEtape}

\subsection{D\'efinition du jeu}

\noindent\textbf{Intuition.} L'adversaire choisit deux messages
$m_0, m_1$, re\c{c}oit la preuve de l'un des deux (au choix secret du
challenger) puis doit deviner lequel. Si le protocole est confidentiel,
l'adversaire ne fait pas mieux qu'une devinette al\'eatoire
(avantage $\approx 0$).
\cmt{L'adversaire ne doit pas pouvoir distinguer les preuves des deux messages.}

\begin{definition}[Jeu~IND-CONF (section~\ref{sec:sotp_construction})]\label{def:ann_ind_conf}
Pour un protocole $\mathcal{P} = (\mathsf{Setup},\mathsf{Prove},
\mathsf{Verify})$ et un adversaire $\mathcal{A}$, le jeu se d\'eroule comme suit.
\begin{enumerate}
\item $\mathcal{A}$ reçoit $pp \leftarrow \mathsf{Setup}(1^\lambda)$
\item $\mathcal{A}$ choisit $m_0,m_1 \in \mathcal{M}_\lambda$
      ($|m_0|=|m_1|$) et $c \in \mathcal{C}_\lambda$
      \hfill\textit{(phase de défi)}
      \cmt{L'adversaire choisit les deux messages \`a distinguer.}
\item Le challenger tire $b \xleftarrow{\$} \{0,1\}$ et calcule
      $\sigma \leftarrow \mathsf{Prove}(pp, m_b, c)$
      \cmt{Le challenger chiffre l'un des deux messages au hasard.}
\item $\mathcal{A}$ reçoit $\sigma$ et émet $b' \in \{0,1\}$
      \hfill\textit{(phase de réponse)}
      \cmt{L'adversaire devine lequel des deux messages a \'et\'e chiffr\'e.}
\item L'expérience produit 1 si $b' = b$, sinon 0
\end{enumerate}
\begin{equation}
\mathsf{Adv}^{\mathsf{IND\text{-}CONF}}_{\mathcal{A},\mathcal{P}}(\lambda)
= \left|\Pr\!\left[
\mathsf{Exp}^{\mathsf{IND\text{-}CONF}}_{\mathcal{A},\mathcal{P}}(\lambda)
= 1\right] - \frac{1}{2}\right|
\end{equation}
\cmtFormule{L'avantage est l'\'ecart \`a la devinette al\'eatoire (probabilit\'e 1/2).}
\cmt{Un avantage n\'egligeable signifie confidentialit\'e.}
\end{definition}

\subsection{Lien avec $\Conf$ au sens de la th\'eorie de l'information}

\begin{remark}[Pont IND-CONF $\leftrightarrow$ $\Conf$]
\label{rem:ann_ind_conf_link}
Pour tout distingueur $\mathcal{A}$, m\^eme non born\'e en calcul, on a
\begin{equation}
\mathsf{Adv}^{\mathsf{IND\text{-}CONF}}_{\mathcal{A},\mathcal{P}}(\lambda)
\;\leq\; \frac{1 - \Conf(\mathcal{P})}{2},
\end{equation}
\cmtFormule{L'avantage IND-CONF est born\'e par la distinguabilit\'e statistique.}
avec \'egalit\'e pour le test de Neyman--Pearson sur la paire la plus
distinguable (Remarque~\ref{rem:jeu_ind_conf_link}). Un protocole avec
$\Conf \geq 1 - \negl(\lambda)$ r\'esiste donc au jeu IND-CONF. La r\'eciproque
est fausse, car la r\'esistance \`a IND-CONF contre les adversaires PPT est une
propri\'et\'e calculatoire, compatible avec une confidentialit\'e statistique faible.
Les r\'eductions de cette annexe \'etablissent des garanties calculatoires, que
nous notons $\Conf_{\mathrm{calc}}$ lorsqu'il faut les distinguer de $\Conf$.
\cmt{La confidentialit\'e statistique implique la confidentialit\'e calculatoire, pas l'inverse.}
\end{remark}

\end{sloppypar}

% ============================================================================
\section{JEU BIND-REL~: FIABILIT\'E GAME-BASED}
\addcontentsline{toc}{section}{Jeu BIND-REL}
\label{sec:gamebased_fiabilite}
% ============================================================================
\begin{sloppypar}

\begin{cmtEtape}{Objectif de cette section.}
D\'efinir le jeu BIND-REL qui capture la fiabilit\'e (non-\'equivocit\'e),
puis \'etablir son lien avec la m\'etrique $\Relver$.
\end{cmtEtape}

\subsection{D\'efinition du jeu}

\noindent\textbf{Intuition.} Le jeu capture la \emph{non-\'equivocit\'e},
selon laquelle un adversaire ne doit pas pouvoir produire une preuve $\sigma$ ouvrable
sur deux messages distincts $m_0 \neq m_1$. Contrairement \`a IND-CONF,
l'adversaire choisit librement $(c, \sigma, s_0, s_1)$, mod\'elisant
un agent qui forge d\'elib\'er\'ement une preuve ambiguë.
\cmt{L'adversaire essaie d'ouvrir une m\^eme preuve sur deux messages diff\'erents.}

\begin{definition}[Jeu~BIND-REL (section~\ref{sec:sotp_construction})]\label{def:ann_bind_rel}
Pour un protocole $\mathcal{P}$ avec $\mathsf{Open}$ et un adversaire
$\mathcal{A}$, le jeu se d\'eroule comme suit.
\begin{enumerate}
\item $\mathcal{A}$ reçoit $pp \leftarrow \mathsf{Setup}(1^\lambda)$
\item $\mathcal{A}$ produit $(c, \sigma, s_0, s_1)$
      \cmt{L'adversaire forge une preuve et deux t\'emoins.}
\item Calculer $m_i = \mathsf{Open}(pp, \sigma, s_i)$ pour $i \in \{0,1\}$
\item $\mathcal{A}$ \emph{gagne} si et seulement si
  $\mathsf{Verify}(pp,\sigma,c)=1$,
  $m_0 \neq m_1$,
  $m_0 \neq \bot$ et $m_1 \neq \bot$
  \cmt{Il gagne si la preuve passe et les deux messages sont distincts et valides.}
\end{enumerate}
\begin{equation}
\mathsf{Win}^{\mathsf{BIND\text{-}REL}}_{\mathcal{A},\mathcal{P}}(\lambda)
= \Pr\!\left[\mathcal{A} \text{ gagne dans }
\mathsf{Exp}^{\mathsf{BIND\text{-}REL}}_{\mathcal{A},\mathcal{P}}(\lambda)
\right]
\end{equation}
\cmtFormule{Probabilit\'e que l'adversaire produise une collision d'engagement.}
\end{definition}

\subsection{Lien avec $\Relver$ (fiabilit\'e de v\'erification)}

\begin{remark}[Pont BIND-REL $\leftrightarrow$ $\Relver$]
\label{rem:ann_bind_rel_link}
Pour un protocole complet dont la probabilit\'e de gain BIND-REL est au plus
$\varepsilon_{\mathsf{bind}}$, on a
\begin{equation}
\Relver \geq 1 - \varepsilon_{\mathsf{bind}} - \negl(\lambda)
\end{equation}
\cmtFormule{Fiabilit\'e de v\'erification born\'ee inf\'erieurement par $1 - \varepsilon_{\mathsf{bind}}$.}
(Remarque~\ref{rem:jeu_bind_rel_link}). BIND-REL ne dit rien de la fiabilit\'e de
reconstruction $\Relrec$, puisqu'un engagement parfaitement masquant et liant a
$\Relrec = 2^{-\lambda}$.
\cmt{Un engagement liant ne r\'ev\`ele pas pour autant le message.}
\end{remark}
\end{sloppypar}

% ============================================================================
\section{JEU EXT-OPP~: OPPOSABILIT\'E GAME-BASED}
\addcontentsline{toc}{section}{Jeu EXT-OPP}
\label{sec:gamebased_opposabilite}
% ============================================================================
\begin{sloppypar}

\begin{cmtEtape}{Objectif de cette section.}
D\'efinir le jeu EXT-OPP qui capture l'opposabilit\'e institutionnelle,
puis \'etablir son lien avec la m\'etrique $\Opp$ via l'in\'egalit\'e de Fano.
\end{cmtEtape}

\subsection{D\'efinition du jeu}

\noindent\textbf{Intuition.} Le jeu mod\'elise la capacit\'e d'un interpr\`ete
institutionnel \`a extraire le bon verdict depuis une preuve. L'\guillemotleft~adversaire~\guillemotright
ici est l'interpr\`ete lui-m\^eme (juge, r\'egulateur), dont on mesure la probabilit\'e
d'\emph{erreur}. Le verdict de r\'ef\'erence $v^*$ est celui que l'institution rendrait en
pleine connaissance de cause (D\'efinition~\ref{def:interp_admissible})~;
l'institution doit rendre une d\'ecision align\'ee sur le contexte l\'egal, au-del\`a
de la seule v\'erification de la preuve.
\cmt{Contrairement aux autres jeux, l'adversaire est ici l'interpr\`ete, qui n'est pas malveillant.}
\cmt{On mesure sa probabilit\'e de se tromper sur le verdict.}

\begin{definition}[Jeu~EXT-OPP (section~\ref{sec:sotp_construction})]\label{def:ann_ext_opp}
Pour un protocole $\mathcal{P}$, un interprète $\mathsf{Interpret}$
de capacité $\Delta$ et une distribution $\mathcal{D}_C$ fixée à l'avance, le jeu se d\'eroule comme suit.
\begin{enumerate}
\item Le challenger tire $c \leftarrow \mathcal{D}_C$,
      $m \leftarrow \mathcal{M}_\lambda$
      \cmt{Le contexte est tir\'e selon la distribution du domaine.}
\item Le challenger calcule $\sigma \leftarrow \mathsf{Prove}(pp,m,c)$
\item L'interprète reçoit $(\sigma,c)$ et produit
      $v \leftarrow \mathsf{Interpret}(pp,\sigma,c,\Delta)$
      \cmt{L'interpr\`ete rend un verdict \`a partir de la preuve et du contexte.}
\item Soit $v^*$ le verdict de r\'ef\'erence, r\'ealisation de la variable $V$
      (qui peut d\'ependre du message, du contexte et des ressources
      institutionnelles, D\'efinition~\ref{def:interp_admissible})
      \cmt{Le verdict de r\'ef\'erence est celui que l'institution rendrait en pleine connaissance de cause.}
\item L'interprète \emph{gagne} si $v = v^*$
\end{enumerate}
\begin{equation}
\mathsf{Err}^{\mathsf{EXT\text{-}OPP}}_{\mathsf{Interpret},\mathcal{P},\mathcal{D}_C}(\lambda)
= \Pr\!\left[
\mathsf{Interpret} \text{ perd dans }
\mathsf{Exp}^{\mathsf{EXT\text{-}OPP}}\right]
\end{equation}
\cmtFormule{Probabilit\'e que l'interpr\`ete rende un verdict incorrect.}
\end{definition}

\noindent\textbf{Pourquoi $\mathcal{D}_C$ est fix\'ee \`a l'avance ?}
L'opposabilit\'e d\'epend du domaine institutionnel. Un protocole tr\`es
opposable dans un contexte contractuel simple peut l'\^etre peu dans un
contexte r\'eglementaire complexe. Fixer $\mathcal{D}_C$ param\'etrise le
jeu par le domaine.
\cmt{Le domaine institutionnel conditionne la difficult\'e de l'extraction.}

\subsection{Lien avec $\Opp$ via Fano}

\begin{remark}[Pont EXT-OPP $\leftrightarrow$ $\Opp$]
\label{rem:ann_ext_opp_link}
Par Fano (Lemme~\ref{lem:ann_g_fano}) appliqu\'e \`a l'estimation de $V$ \`a
partir de $(\Sigma,C)$, on obtient
\begin{equation}
\Opp = 1 - \frac{H(V|\Sigma,C)}{H(V|C)}
\geq 1 - \frac{h(\varepsilon_{\mathsf{ext}})
  + \varepsilon_{\mathsf{ext}}\log_2(|\mathcal{V}_\lambda|-1)}
  {H(V|C)}
\label{eq:ann_fano_opp}
\end{equation}
\cmtFormule{L'opposabilit\'e est born\'ee inf\'erieurement par $1 - \varepsilon_{\mathsf{ext}}$.}
où $\varepsilon_{\mathsf{ext}} = \mathsf{Err}^{\mathsf{EXT\text{-}OPP}}$
et $h(\cdot)$ est l'entropie binaire.
Un interprète avec faible erreur d'extraction implique une forte
opposabilité.
\cmt{Moins l'interpr\`ete se trompe, plus l'opposabilit\'e est grande.}
\end{remark}
\end{sloppypar}

% ============================================================================
\section{R\'EDUCTIONS FORMELLES POUR SOTP-LOW}
\addcontentsline{toc}{section}{R\'eductions SOTP-Low}
\label{sec:gamebased_SOTP-Low}
% ============================================================================
\begin{sloppypar}
SOTP-Low op\`ere dans le r\'egime \`a contexte rare ($H(\mathcal{C}) < H^*(\lambda)$).
Ses deux r\'eductions principales ciblent PRG (confidentialit\'e) et
logarithme discret (fiabilit\'e).

\begin{cmtEtape}{Objectif de cette section.}
\'Etablir les deux r\'eductions formelles qui relient SOTP-Low aux
hypoth\`eses standard PRG (couche Gold) et DLOG (couche Iron).
\end{cmtEtape}

\subsection{R\'eduction PRG pour la confidentialit\'e}

\noindent\textbf{Intuition.} La couche Gold chiffre $m$ par le pad $G(s)$.
Si $\mathcal{A}$ distingue les preuves de $m_0$ vs $m_1$, il distingue
n\'ecessairement $G(s)$ d'une cha\^ine al\'eatoire, contredisant la s\'ecurit\'e PRG.
\cmt{Toute attaque contre la confidentialit\'e devient une attaque contre le PRG.}

\begin{theorem*}[SOTP-Low~:~Réduction~PRG (section~\ref{sec:sotp_construction})]\label{thm:ann_sotp_low_conf}
Soit $\mathcal{A}$ adversaire PPT contre IND-CONF de SOTP-Low avec
avantage $\varepsilon$. Alors $\exists$ adversaire PPT $\mathcal{B}$
contre le PRG $G$ avec avantage $\geq \varepsilon/2 - \negl(\lambda)$.
\cmt{L'avantage contre le PRG est au moins la moiti\'e de l'avantage contre IND-CONF.}
\end{theorem*}

\begin{Preuve}

\begin{cmtEtape}{Construction du r\'educteur $\mathcal{B}$.}
$\mathcal{B}$ re\c{c}oit un d\'efi PRG (soit $G(s)$, soit uniforme) et
simule l'environnement IND-CONF pour $\mathcal{A}$. Il utilise son
d\'efi comme pad masquant, puis observe si $\mathcal{A}$ devine
correctement $b$.
\end{cmtEtape}

\textbf{Construction de $\mathcal{B}$.} $\mathcal{B}$ re\c{c}oit un d\'efi
PRG $y \in \{0,1\}^L$ (soit $G(s)$ soit uniforme) et simule le jeu
IND-CONF pour $\mathcal{A}$ de la mani\`ere suivante.
\begin{enumerate}
\item G\'en\'erer $(\mathbb{G},q,g,h)$ comme dans le Setup SOTP-Low
      \cmt{$\mathcal{B}$ peut g\'en\'erer les param\`etres publics car ils ne d\'ependent pas du PRG.}
\item Quand $\mathcal{A}$ soumet $(m_0,m_1,c)$, choisir
      $b \xleftarrow{\$} \{0,1\}$
      \cmt{$\mathcal{B}$ choisit un bit secret comme dans le vrai jeu.}
\item Poser $\sigma_1 = m_b \oplus y$ \hfill (d\'efi PRG comme pad)
      \cmt{Le d\'efi PRG joue le r\^ole du pad.}
\item Poser $\sigma_2 = g^{s'}h^\rho$ pour
      $s',\rho \xleftarrow{\$} \mathbb{Z}_q$ \hfill (Iron simul\'ee)
      \cmt{$\mathcal{B}$ simule la couche Iron sans conna\^itre $s$.}
\item Poser $\sigma_3 = (\phi(c), \mathsf{Hash}(c), V(c))$
      \hfill (Clay d\'eterministe)
      \cmt{La couche Clay est simul\'ee exactement comme dans SOTP-Low.}
\item Retourner $\sigma = (\sigma_1,\sigma_2,\sigma_3)$ \`a $\mathcal{A}$
\end{enumerate}
Si $\mathcal{A}$ \'emet $b'=b$, $\mathcal{B}$ retourne 1 ; sinon 0.
\cmt{$\mathcal{B}$ retourne 1 si $\mathcal{A}$ a devin\'e correctement.}

\textbf{Analyse.}
\begin{itemize}
\item \emph{Cas $y = G(s)$.} $\mathcal{A}$ voit une vraie preuve SOTP-Low,
      gagne avec probabilit\'e $1/2 + \varepsilon$.
      $\mathcal{B}$ retourne 1 avec probabilit\'e $1/2+\varepsilon$.
      \cmt{Dans ce cas, $\mathcal{B}$ a bien simul\'e le vrai jeu.}
\item \emph{Cas $y$ uniforme.} $\sigma_1 = m_b \oplus U_{|m|}$ est un
      one-time pad parfait, $\mathcal{A}$ n'a aucune information sur $b$.
      $\mathcal{B}$ retourne 1 avec probabilit\'e $1/2$.
      \cmt{Dans ce cas, $\mathcal{A}$ ne peut pas mieux faire que 1/2.}
\end{itemize}
Donc $\mathsf{Adv}^{\mathsf{PRG}}_{\mathcal{B}}(\lambda) \geq
\varepsilon/2 - \negl(\lambda)$.
\cmtFormule{Conclusion de la r\'eduction.}
\end{Preuve}

\subsection{R\'eduction logarithme discret pour la fiabilit\'e}

\noindent\textbf{Intuition.} La couche Iron est
$\sigma_2 = g^s h^\rho$. Pour gagner BIND-REL, l'adversaire doit trouver
deux ouvertures $s_0 \neq s_1$ pour le m\^eme $\sigma_2$, ce qui implique
$g^{s_0}h^{\rho_0} = g^{s_1}h^{\rho_1}$, permettant de calculer
$\log_g h$, r\'esolvant DLOG.
\cmt{Deux ouvertures distinctes donnent directement le logarithme discret.}

\begin{theorem}[SOTP-Low~:~Réduction~DLOG (section~\ref{sec:sotp_construction})]\label{thm:ann_sotp_low_rel}
On suppose la graine tir\'ee dans $\mathbb{Z}_q$ et l'engagement de la forme
$\sigma_2 = g^{s}h^{\rho}$ (Algorithme~\ref{alg:sotp_low}). Soit $\mathcal{A}$ PPT
gagnant BIND-REL pour SOTP-Low (proc\'edure d'ouverture $\mathsf{Open}$) avec
probabilit\'e $\varepsilon$. Alors $\exists$ adversaire PPT $\mathcal{B}$ r\'esolvant
DLOG dans $\mathbb{G}$ avec probabilit\'e $\geq \varepsilon$.
\cmt{L'avantage contre DLOG est le m\^eme que contre BIND-REL.}
\end{theorem}

\begin{Preuve}

\begin{cmtEtape}{Construction du r\'educteur DLOG.}
$\mathcal{B}$ re\c{c}oit une instance DLOG $(g, h)$ et utilise l'attaque
de $\mathcal{A}$ contre BIND-REL pour extraire $\log_g h$. La cl\'e est
que deux ouvertures valides sur la m\^eme preuve donnent une \'equation
alg\'ebrique r\'esoluble.
\end{cmtEtape}

Soit une instance DLOG $(g,h)$ avec $x = \log_g h$ inconnu.
\begin{enumerate}
\item $\mathcal{B}$ configure SOTP-Low avec $(g,h)$ donn\'es
      \cmt{$\mathcal{B}$ utilise l'instance DLOG comme param\`etres publics~; leur loi est celle du Setup.}
\item $\mathcal{A}$ produit $(c,\,\sigma{=}(\sigma_1,\sigma_2,\sigma_3))$ et deux
      ouvertures accept\'ees $(m_0, s_0, \rho_0)$, $(m_1, s_1, \rho_1)$ avec
      $m_0 \neq m_1$
      \cmt{$\mathcal{A}$ fournit sa collision d'ouverture.}
\item Comme $\sigma_1 = m_0 \oplus \mathsf{PRG}(s_0) = m_1 \oplus \mathsf{PRG}(s_1)$ et
      $m_0 \neq m_1$, on a $\mathsf{PRG}(s_0) \neq \mathsf{PRG}(s_1)$, donc
      $s_0 \neq s_1$ dans $\mathbb{Z}_q$
      \cmt{C'est ici que le tirage de la graine dans $\mathbb{Z}_q$ est indispensable, car avec une graine de $\{0,1\}^\lambda$ engag\'ee par $\mathsf{int}(s) \bmod q$, $s$ et $s+q$ donneraient le m\^eme engagement.}
\item Les deux ouvertures donnent $g^{s_0}h^{\rho_0} = g^{s_1}h^{\rho_1}$, donc
      $g^{s_0-s_1} = h^{\rho_1-\rho_0}$~; comme $s_0 \neq s_1$ et $g$ engendre
      $\mathbb{G}$, on a $\rho_1 \neq \rho_0$
      \cmtFormule{On isole $\log_g h = x$ dans l'\'equation.}
\item Calculer
\begin{equation}
x = \bigl(s_0-s_1\bigr)\cdot\bigl(\rho_1-\rho_0\bigr)^{-1} \bmod q
\end{equation}
\cmtFormule{On r\'esout l'\'equation lin\'eaire pour obtenir $x = \log_g h$.}
\end{enumerate}
$\mathcal{B}$ r\'eussit chaque fois que $\mathcal{A}$ produit une
attaque valide, avec probabilit\'e $\varepsilon$.
\cmt{La r\'eduction est exacte, sans perte d'avantage.}
\end{Preuve}

\noindent\textbf{Signification pratique.} Avec Curve25519 ($q \approx 2^{252}$,
s\'ecurit\'e 128 bits), casser la fiabilit\'e de SOTP-Low est aussi difficile
que r\'esoudre le DLP sur les courbes elliptiques.
\cmt{Le niveau de s\'ecurit\'e de SOTP-Low h\'erite de celui de DLOG sur la courbe.}
\end{sloppypar}

% ============================================================================
\section{R\'EDUCTION JL POUR SOTP-BASE}
\addcontentsline{toc}{section}{R\'eduction SOTP-Base}
\label{sec:gamebased_SOTP-base}
% ============================================================================
\begin{sloppypar}
SOTP-Base op\`ere dans le r\'egime de transition
($H^*(\lambda) \leq H(\mathcal{C}) \leq \lambda/\log_2\lambda$).
Sa couche Clay applique une projection JL $\Phi\colon \mathbb{R}^n \to
\mathbb{R}^d$ avec $d = \lceil 4\Delta H(\mathcal{C})/\varepsilon^2
\rceil$.

\begin{cmtEtape}{Objectif de cette section.}
Montrer que l'erreur d'extraction de l'interpr\`ete est directement
reli\'ee \`a la qualit\'e de la pr\'eservation JL. Si la projection
d\'eforme trop les distances, l'extraction devient erron\'ee.
\end{cmtEtape}

\noindent\textbf{Intuition.} Si $\Phi$ d\'eforme trop certaines distances,
les contextes voisins deviennent indiscernables dans l'espace compress\'e,
causant des erreurs d'extraction. La r\'eduction ci-dessous quantifie
pr\'ecis\'ement cette relation.
\cmt{La compression JL est le seul point de perte en SOTP-Base.}

\begin{theorem}[SOTP-Base~:~Dimensionnement~JL (section~\ref{sec:sotp_construction})]\label{thm:ann_sotp_base_opp}
Soit $\mathcal{C} = \{c_1,\ldots,c_N\}$ l'ensemble des contextes plong\'es, $N =
2^{H(\mathcal{C})}$ et $\delta, \varepsilon' \in (0,1)$. Si la dimension de
projection v\'erifie
\begin{equation}
d \;\geq\; \frac{8}{\delta^2}\,\ln\!\left(\frac{N^2}{\varepsilon'}\right)
= O\!\left(\delta^{-2}\bigl(H(\mathcal{C}) + \log(1/\varepsilon')\bigr)\right),
\end{equation}
alors, avec probabilit\'e au moins $1 - \varepsilon'$ sur le tirage de $\Phi$,
$\Phi$ pr\'eserve toutes les distances par paires de $\mathcal{C}$ \`a un facteur
$(1 \pm \delta)$ pr\`es.
\cmt{La pr\'eservation JL conditionne l'extraction~; la r\'eciproque est fausse.}
\end{theorem}

\begin{Preuve}
Par le lemme JL, pour chaque paire $(i,j)$,
$\Pr\bigl[\,\bigl|\|\Phi c_i-\Phi c_j\|^2/\|c_i-c_j\|^2 - 1\bigr| > \delta\,\bigr]
\leq 2e^{-d\delta^2/8}$~; une borne de l'union sur les $N^2/2$ paires donne une
probabilit\'e d'\'echec au plus $N^2 e^{-d\delta^2/8}$, inf\'erieure \`a $\varepsilon'$
d\`es que $d \geq (8/\delta^2)\ln(N^2/\varepsilon')$.
\cmtFormule{Borne de l'union sur les paires de contextes.}
La pr\'eservation des distances est une condition \emph{n\'ecessaire} pour qu'une
fonction d'extraction fond\'ee sur les distances distingue les contextes~; elle
n'est pas suffisante pour garantir l'information extractible sur le verdict.
Elle ne fournit donc aucune minoration de $\Opp$. Cette dimension, en
$O(H(\mathcal{C}))$, est distincte de la r\`egle de dimensionnement
$\lceil 4\Delta H(\mathcal{C})/\varepsilon^2 \rceil$ du texte principal, qui est
un choix de conception.
\end{Preuve}

\noindent\textbf{Interpr\'etation.} SOTP-Base choisit la dimension
$d$ garantissant la pr\'eservation des distances avec la probabilit\'e voulue~;
le lien avec l'extraction du verdict reste une hypoth\`ese de conception.
\cmt{Le choix de $d$ est un compromis entre taille de preuve et qualit\'e d'extraction.}
\end{sloppypar}

% ============================================================================
\section{R\'EDUCTIONS POUR SOTP-HIGH}
\addcontentsline{toc}{section}{R\'eductions SOTP-High}
\label{sec:gamebased_SOTP-High}
% ============================================================================
\begin{sloppypar}
SOTP-High op\`ere dans le r\'egime \`a contexte riche
($H(\mathcal{C}) > \lambda/\log_2\lambda$) et s'appuie sur Groth16
pour sa couche Gold. Sa couche Clay encode $h_c = \mathsf{Poseidon}(c)$.

\begin{cmtEtape}{Objectif de cette section.}
Montrer que $\Conf$ et $\Relver$ de SOTP-High h\'eritent des propri\'et\'es de
Groth16 (z\'ero-connaissance parfaite, solidit\'e connaissance)~; l'opposabilit\'e
n'est pas couverte par une r\'eduction.
\end{cmtEtape}

\begin{theorem}[SOTP-High~:~Sécurité~héritée~de~Groth16 (section~\ref{sec:sotp_construction})]\label{thm:ann_sotp_high}
Pour la relation $\mathcal{R} = \{((h_c, c_m), (m, c, \rho)) : h_c = \mathsf{Poseidon}(c)
\wedge c_m = g^m h^\rho \wedge \mathsf{Conforme}(m,c) = 1\}$, les deux propri\'et\'es suivantes sont v\'erifi\'ees.
\begin{enumerate}
\item \textbf{Confidentialité.} La z\'ero-connaissance parfaite
      de Groth16 et le masquage parfait de $c_m$ donnent $\Conf = 1$.
      \cmt{La preuve et l'engagement sont ind\'ependants de $m$ sachant $c$.}
\item \textbf{Fiabilité de v\'erification.} La solidité connaissance de
      Groth16 donne $\Relver \geq 1 - \negl(\lambda)$.
      \cmt{Un prouveur qui fait accepter conna\^it un t\'emoin valide.}
\end{enumerate}
\end{theorem}

\begin{Preuve}

\textbf{Confidentialit\'e.} Par la ZK parfaite, les preuves honn\^etes et les
preuves simul\'ees $\mathcal{S}(\sigma_p, (h_c, c_m))$ ont la m\^eme loi. Le
simulateur n'utilise que les entr\'ees publiques~; $h_c$ ne d\'epend que de $c$ tandis que
$c_m$ est uniforme dans $\mathbb{G}$ quel que soit $m$. La loi de
$\Sigma = (h_c, \pi, c_m, \mathsf{metadata})$ sachant $(M = m, C = c)$ ne d\'epend donc
pas de $m$, d'o\`u $\Conf = 1$.
\cmtFormule{La ZK parfaite implique l'indistinguabilit\'e inconditionnelle.}

\textbf{Fiabilit\'e de v\'erification.} Pour tout $\mathcal{A}$ PPT produisant une
preuve acceptante $\pi$, l'extracteur $\mathcal{E}$ de Groth16 (qui dispose du
code de $\mathcal{A}$ et exploite l'hypoth\`ese q-PKE~; ce n'est pas le
v\'erificateur) produit un t\'emoin $(m,c,\rho)$ avec $((h_c,c_m),(m,c,\rho)) \in
\mathcal{R}$, sauf avec probabilit\'e n\'egligeable. Ainsi
\begin{equation}
\Pr\!\left[\mathsf{Verify}=1 \land (m,c,\rho)\notin\mathcal{R}\right]
\leq \mathsf{Adv}^{\mathsf{sound}}(\lambda) \leq \negl(\lambda)
\;\Rightarrow\; \Relver \geq 1-\negl(\lambda).
\end{equation}
\cmtFormule{La solidit\'e connaissance implique l'absence de forge.}

\textbf{Opposabilit\'e.} Aucune borne n'est d\'eriv\'ee ici, car $h_c$ et les
m\'etadonn\'ees sont des fonctions du contexte public, donc $I(V;h_c|C) = 0$~; ce que
la preuve apporte au verdict tient \`a l'assertion certifi\'ee
$\mathsf{Conforme}(m,c) = 1$ (Remarque~\ref{rem:heuristique_sotp_high}).
\end{Preuve}

\begin{remark}[La ZK parfaite] La ZK \emph{parfaite}
de Groth16 signifie que m\^eme un adversaire \`a calcul infini ne peut
distinguer vraies preuves de simul\'ees. C'est une garantie inconditionnelle
une fois le CRS fix\'e, propri\'et\'e forte, qui fait de SOTP-High
la variante la plus protectrice pour les contextes \`a haute entropie.
\cmt{La ZK parfaite est inconditionnelle.}
\end{remark}
\end{sloppypar}

% ============================================================================
\section{\'ENONC\'ES COMPLETS ET PRODUIT CRO GAME-BASED}
\addcontentsline{toc}{section}{\'Enonc\'es complets et produit CLO}
\label{sec:gamebased_produitCRO}
% ============================================================================
\begin{sloppypar}

\begin{cmtEtape}{Objectif de cette section.}
Rassembler les bornes game-based des trois jeux et en d\'eduire une
minoration du produit CLO \`a l'\'echelle op\'erationnelle. Elle ne se compare pas
directement au Trilemme CRO (Remarque~\ref{rem:ann_pas_de_contradiction}).
\end{cmtEtape}

\subsection{S\'ecurit\'e compl\`ete de SOTP-Low}

\begin{theorem}[SOTP-Low~:~Bornes~game-based~compl\`etes (section~\ref{sec:sotp_construction})]\label{thm:ann_sotp_low_games}
Sous PRG et DLOG dans $\mathbb{G}$ d'ordre $q$, on a
\begin{align}
\mathsf{Adv}^{\mathsf{IND\text{-}CONF}}_{\mathcal{A},\mathsf{SOTP\text{-}Low}}(\lambda)
&\leq 2\cdot\mathsf{Adv}^{\mathsf{PRG}}_{\mathcal{B}}(\lambda)
    + \frac{1}{q}
\label{eq:ann_sotp_low_conf} \\[4pt]
\mathsf{Win}^{\mathsf{BIND\text{-}REL}}_{\mathcal{A},\mathsf{SOTP\text{-}Low}}(\lambda)
&\leq \mathsf{Adv}^{\mathsf{DLOG}}_{\mathcal{B}'}(\lambda)
    + \frac{1}{q}
\label{eq:ann_sotp_low_rel}
\end{align}
\cmtFormule{Confidentialit\'e calculatoire born\'ee par l'avantage PRG.}
\cmtFormule{Fiabilit\'e de v\'erification born\'ee par l'avantage DLOG.}
pour tout PPT $\mathcal{A}$.
\end{theorem}

\begin{remark}[Opposabilit\'e de SOTP-Low~: pas de borne game-based]
\label{rem:ann_sotp_low_opp_statut}
La formulation initiale de ce th\'eor\`eme contenait une troisi\`eme in\'egalit\'e,
$\mathsf{Err}^{\mathsf{EXT\text{-}OPP}} \leq 1 - \Delta/(\alpha H(\mathcal{C}))$,
tir\'ee du fait que la couche Clay $\sigma_3 = (\phi(c), \mathsf{Hash}(c), V(c))$ d\'etermine
le verdict. La pr\'esente formulation ne la reprend pas, car $\sigma_3$ est une fonction du
contexte public $C$, donc $I(V;\sigma_3|C) = 0$~; elle ne contribue ainsi en rien \`a
l'opposabilit\'e au sens de la D\'efinition~\ref{def:interp_admissible}. Ce que
SOTP-Low apporte au verdict passe par l'ouverture $(m,s,\rho)$ fournie \`a
l'institution, ce qui rel\`eve de l'heuristique de conception
(Remarque~\ref{rem:facteur_limitant_low}), sans r\'eduction associ\'ee.
\end{remark}

\subsection{Tableau r\'ecapitulatif des r\'eductions}

\begin{table}[htbp]
\centering
\caption{Synth\`ese des r\'eductions game-based pour la famille SOTP}
\label{tab:ann_g_reductions}
\begin{tabular}{@{}lllll@{}}
\toprule
\textbf{Variante} & \textbf{Propri\'et\'e} & \textbf{Jeu} &
\textbf{Hypoth\`ese} & \textbf{Borne} \\
\midrule
SOTP-Low & $\Conf_{\mathrm{calc}}$ & IND-CONF & PRG
  & $\mathsf{Adv} \leq 2\varepsilon_{\mathsf{PRG}} + 1/q$ \\
SOTP-Low & $\Relver$ & BIND-REL & DLOG
  & $\mathsf{Win} \leq \varepsilon_{\mathsf{DLOG}} + 1/q$ \\
SOTP-Low & $\Opp$ & -- & --
  & aucune (Rem.~\ref{rem:ann_sotp_low_opp_statut}) \\
\midrule
SOTP-Base & $\Conf_{\mathrm{calc}}$ & IND-CONF & PRG
  & $\mathsf{Adv} \leq 2\varepsilon_{\mathsf{PRG}} + \negl$ \\
SOTP-Base & $\Relver$ & BIND-REL & DLOG
  & $\mathsf{Win} \leq \varepsilon_{\mathsf{DLOG}} + \negl$ \\
SOTP-Base & -- & -- & Lemme JL
  & dimensionnement de $d$ (Th.~\ref{thm:ann_sotp_base_opp}) \\
\midrule
SOTP-High & $\Conf$ & IND-CONF & Groth16 ZK parfaite
  & $\Conf = 1$ \\
SOTP-High & $\Relver$ & BIND-REL & Groth16 (q-PKE)
  & $\geq 1-\negl(\lambda)$ \\
SOTP-High & $\Opp$ & -- & --
  & aucune (Rem.~\ref{rem:heuristique_sotp_high}) \\
\bottomrule
\end{tabular}
\end{table}
\cmtFormule{Tableau r\'ecapitulatif des neuf r\'eductions.}

\subsection{Corollaire central~: produit CLO game-based}

\begin{corollary}[Produit CLO depuis les jeux (section~\ref{sec:sotp_construction})]\label{cor:ann_cro_from_games}
Si IND-CONF, BIND-REL et EXT-OPP sont gagn\'es avec avantages respectifs au plus
$\varepsilon_{\mathsf{ind}}$, $\varepsilon_{\mathsf{bind}}$, $\varepsilon_{\mathsf{ext}}$, alors
\begin{equation}
\Conf_{\mathrm{calc}} \cdot \Relver \cdot \Opp
\geq \!\left(1-2\varepsilon_{\mathsf{ind}}\right)\!
\left(1-\varepsilon_{\mathsf{bind}}\right)\!
\left(1 - \frac{h(\varepsilon_{\mathsf{ext}})
  + \varepsilon_{\mathsf{ext}}\log_2(|\mathcal{V}_\lambda|-1)}
  {H(V|C)}\right)
\label{eq:ann_cro_product_games}
\end{equation}
\cmtFormule{Produit des trois bornes inf\'erieures.}
o\`u $\Conf_{\mathrm{calc}} = 1 - 2\varepsilon_{\mathsf{ind}}$ d\'esigne la confidentialit\'e
calculatoire. Pour $\varepsilon_{\mathsf{ind}},\varepsilon_{\mathsf{bind}},
\varepsilon_{\mathsf{ext}} \ll 1$, il vient
\begin{equation}
\Conf_{\mathrm{calc}} \cdot \Relver \cdot \Opp
\geq 1 - O\!\left(\varepsilon_{\mathsf{ind}}
  + \varepsilon_{\mathsf{bind}}
  + \varepsilon_{\mathsf{ext}}\log_2|\mathcal{V}_\lambda| + h(\varepsilon_{\mathsf{ext}})\right)
\label{eq:ann_cro_approx}
\end{equation}
\cmtFormule{Approximation au premier ordre pour petits $\varepsilon$.}
\end{corollary}

\begin{Preuve}
\textbf{Terme $\Conf_{\mathrm{calc}}$.} D\'efinition, avec le jeu IND-CONF
(Remarque~\ref{rem:ann_ind_conf_link}).
\textbf{Terme $\Relver$.} La Remarque~\ref{rem:ann_bind_rel_link} donne
$\Relver \geq 1-\varepsilon_{\mathsf{bind}} - \negl(\lambda)$.
\textbf{Terme $\Opp$.} Remarque~\ref{rem:ann_ext_opp_link}
(in\'egalit\'e de Fano, \eqref{eq:ann_fano_opp}).
Le produit $(1-a)(1-b)(1-c) \geq 1-a-b-c$ pour $a,b,c \in [0,1]$
donne~\eqref{eq:ann_cro_approx}.
\end{Preuve}

\begin{remark}[Pourquoi ce corollaire ne contredit pas le Trilemme]
\label{rem:ann_pas_de_contradiction}
Le Trilemme CRO (Th\'eor\`eme~\ref{thm:trilemme_cro_conditionnel}) majore le produit
$\Conf\cdot\Relrec\cdot\Opp$, o\`u $\Conf$ est statistique et $\Relrec$ mesure la
reconstruction du message. Le corollaire minore un produit d'une autre nature, qui fait intervenir la
confidentialit\'e \emph{calculatoire} et la fiabilit\'e de \emph{v\'erification}.
Les deux \'echelles ne se comparent pas. Un protocole peut avoir
$\Conf_{\mathrm{calc}}\cdot\Relver\cdot\Opp$ proche de~1 tout en ob\'eissant au
Trilemme, puisque $\Relrec$ reste alors de l'ordre de $2^{-\lambda}$
(Remarque~\ref{rem:rel_ver}). L'hypoth\`ese que
$\varepsilon_{\mathsf{ext}}$ est petit est pr\'ecis\'ement ce que le Trilemme rend
co\^uteux lorsque $\Conf$ statistique est proche de~1~; la
Remarque~\ref{rem:activite_delta} indique quelle quantit\'e ($\Delta$ ou Fano) limite
alors $\Opp$. La formulation initiale de ce paragraphe annon\c{c}ait un
\guillemotleft~encadrement~\guillemotright{} du produit par $\Delta/(\alpha\lambda)$ et une
constante $c \in [0{,}45;\,0{,}85]$~; la pr\'esente formulation distingue les deux \'echelles.
\cmt{Les deux produits rel\`event d'\'echelles distinctes et n'admettent pas d'encadrement commun.}
\end{remark}
\end{sloppypar}

% ============================================================================
\section{LIMITATIONS ET PORT\'EE DES R\'EDUCTIONS}
\addcontentsline{toc}{section}{Limitations et port\'ee}
\label{sec:gamebased_limitations}
% ============================================================================
\begin{sloppypar}

\begin{cmtEtape}{Objectif de cette section.}
\'Enum\'erer les cinq limitations des r\'eductions game-based, \`a savoir la
conditionnalit\'e, l'interpr\`ete honn\^ete, la r\'eduction one-way, le terme $1/q$
et la relation avec Lean4.
\end{cmtEtape}

\begin{enumerate}
\item \textbf{Conditionnalit\'e aux hypoth\`eses calculatoires.}
      La s\'ecurit\'e de SOTP-Low est prouv\'ee sous PRG et DLOG.
      SOTP-High via Groth16 repose sur la difficult\'e des
      couplages sur BN254, non post-quantique. Une extension aux
      primitives r\'esistantes aux quantiques (STARKs, r\'eseaux
      euclidiens) est une direction de recherche ouverte (section~\ref{sec:sotp_construction}).
      \cmt{La s\'ecurit\'e d\'epend d'hypoth\`eses non prouv\'ees (PRG, DLOG, couplages).}

\item \textbf{Mod\`ele d'interpr\`ete honn\^ete dans EXT-OPP.}
      Le jeu suppose que l'interpr\`ete tente sinc\`erement d'extraire
      le bon verdict. Un interpr\`ete \emph{adversarial} (juge
      corrompu) n'est pas mod\'elis\'e. Des extensions game-th\'eoriques
      (r\'egulateur vs entit\'e audit\'ee) sont identifi\'ees comme
      direction future (section~\ref{sec:sotp_construction}).
      \cmt{Un juge corrompu n'est pas couvert par le mod\`ele actuel.}

\item \textbf{R\'eduction SOTP-Base one-way.}
      Le Th\'eor\`eme~\ref{thm:ann_sotp_base_opp} garantit seulement l'implication
      \textit{dimension suffisante $\Rightarrow$ pr\'eservation JL}. Il ne garantit
      pas que la pr\'eservation des distances permette d'extraire le verdict,
      car cela d\'epend de la fonction $f_\Delta$ choisie par l'institution, dont ni
      l'existence ni la calculabilit\'e efficace ne sont \'etablies.
      \cmt{La pr\'eservation JL est n\'ecessaire, pas suffisante, pour l'opposabilit\'e.}

\item \textbf{Terme $1/q$ dans les bornes SOTP-Low.}
      Les \'equations~\eqref{eq:ann_sotp_low_conf}
      et~\eqref{eq:ann_sotp_low_rel} contiennent $1/q$
      (probabilit\'e de deviner l'exposant). Il est n\'egligeable en $\lambda$ au
      sens strict d\`es que l'ordre du groupe v\'erifie $\log_2 q \geq \lambda$, ce
      qui est le cas des param\`etres utilis\'es ($q \approx 2^{252}$, soit
      $1/q \approx 10^{-76}$). Le terme est conserv\'e pour que les bornes restent
      exactes \`a $\lambda$ fix\'e.
      \cmt{Le terme $1/q$ est n\'egligeable d\`es que $\log_2 q \geq \lambda$.}

\item \textbf{Relation avec la formalisation Lean4 (Annexe~\ref{chap:annexe_lean4}).}
      Les deux annexes se compl\`etent de la mani\`ere suivante.
      \begin{itemize}
      \item \textbf{Lean4 (Annexe~\ref{chap:annexe_lean4})} certifie m\'ecaniquement la
            d\'erivation alg\'ebrique de la borne
            $\Conf\cdot\Relrec\cdot\Opp \leq \Delta/\lambda + K/\lambda^2 + \varepsilon(\lambda)$
            (Th\'eor\`eme~\ref{thm:conditional_CRO_trilemma}), en posant Fano et
            Pinsker comme propri\'et\'es de chaque instance et l'extraction born\'ee
            comme hypoth\`ese
            \cmt{Certification m\'ecanique de la logique interne.}
      \item \textbf{Annexe~\ref{chap:annexe_gamebased} (ici)} \'etablit les
            propri\'et\'es cryptographiques de la famille SOTP (confidentialit\'e
            calculatoire, fiabilit\'e de v\'erification) sous PRG, DLOG et Groth16~;
            elle n'instancie pas les hypoth\`eses du Trilemme, qui portent sur
            $\Conf$ statistique et $\Relrec$
            \cmt{Instanciation dans des primitives calculatoires concr\`etes.}
      \end{itemize}
      Les deux niveaux sont donc compl\'ementaires sans se composer en une preuve
      unique, puisque le niveau \emph{formel} certifie la logique de la d\'erivation du
      Trilemme tandis que le niveau \emph{game-based} justifie les garanties
      cryptographiques des constructions.
      \cmt{La preuve se d\'eploie sur deux niveaux, formel (Lean4) et game-based (ici).}
\end{enumerate}
\end{sloppypar}
% ============================================================================
% CODE LEAN4 COMPLET ET ANNOTÉ — Version 2.2
% Annexe de thèse — Formalisation du Trilemme CRO conditionnel
% ============================================================================

\chapter{Code Lean4 complet et annot\'e}
\label{chap:annexe_lean4}

\begin{sloppypar}
Cette annexe pr\'esente le code Lean4 complet de la formalisation du
Trilemme CRO conditionnel, dans sa version 2.2 (octobre 2026), qui reprend
\`a l'identique les \'enonc\'es de la version~2.1 et ajoute les extensions de la
sous-section~\ref{subsec:lean4_extensions}.

\medskip
\noindent\textbf{Caract\'eristiques du code.} Le fichier compile sans erreur
ni \texttt{sorry} sur \url{https://live.lean-lang.org} (projet Mathlib, version
courante, v\'erification d'octobre 2026). Le d\'ep\^ot
compagnon~\cite{CROTrilemmaDepot} le fournit aussi dans un projet Lake, avec des
versions de Lean et de Mathlib fig\'ees (r\'epertoire
\texttt{05\_formalisation\_lean/}). Il ne contient aucun axiome ad
hoc, car les bornes informationnelles (Pinsker, Fano) sont des propri\'et\'es des
instances de protocole tandis que les hypoth\`eses institutionnelles
(Hypoth\`eses~1, 2$'$, 3, 4) sont des hypoth\`eses explicites des
th\'eor\`emes. Toutes les preuves ne d\'ependent que des trois axiomes
standards de Lean/Mathlib, \`a savoir \texttt{propext}, \texttt{Classical.choice}
et \texttt{Quot.sound}.

\medskip
\noindent\textbf{Mode d'emploi.} Pour v\'erifier le code (d\'etails \`a la
section~\ref{sec:lean4_reproduction}), on proc\`ede comme suit.
\begin{enumerate}
\item Dans le navigateur, coller le Listing~\ref{lst:lean4_complet} dans
      \url{https://live.lean-lang.org} (projet Mathlib), qui utilise la version
      courante de Mathlib.
\item Localement, ex\'ecuter \texttt{lake exe cache get \&\& lake build} dans le
      r\'epertoire \texttt{05\_formalisation\_lean/} du d\'ep\^ot.
\item Dans les deux cas, le r\'esultat attendu ne comporte ni erreur ni \texttt{sorry}. Chacun des dix
      \texttt{\#print axioms} du bloc final affiche uniquement
      \texttt{[propext, Classical.choice, Quot.sound]}.
\end{enumerate}
\end{sloppypar}

% ----------------------------------------------------------------------------
\section{ARCHITECTURE DU CODE}
\label{sec:lean4_architecture}
% ----------------------------------------------------------------------------

\begin{sloppypar}
Le code est organis\'e en dix sections qui suivent la structure
math\'ematique du Trilemme CRO.
\begin{enumerate}
\item \textbf{Constante de Pinsker} $c_0 = 1/(2\ln 2)$ et ses propri\'et\'es.
\item \textbf{N\'egligeabilit\'e} de $\varepsilon_0(\lambda) = 2^{-\lambda}$,
      prouv\'ee via les outils d'analyse r\'eelle de Mathlib.
\item \textbf{Param\`etres institutionnels} $(\Delta, \alpha, \beta, H(C))$
      et le coefficient $\gamma = \alpha H(C) - 1$.
\item \textbf{Instance de protocole} avec les d\'efinitions de Conf, Rel,
      Opp ainsi que les lemmes de Pinsker et Fano int\'egr\'es comme propri\'et\'es.
\item \textbf{Hypoth\`eses institutionnelles} 1, 2$'$ et 3.
\item \textbf{Lemme 4} (garantie de confidentialit\'e) et
      \textbf{Th\'eor\`eme 1} (forme explicite et instance).
\item \textbf{Th\'eor\`eme principal} (forme asymptotique) et t\'emoins
      de coh\'erence et de port\'ee.
\item \textbf{R\'egimes de capacit\'e institutionnelle} (classification indicative).
\item \textbf{Extensions v2.2}, soit le verdict d\'eterministe et la forme raffin\'ee.
\item \textbf{Bloc de v\'erification} (\texttt{\#check}, \texttt{\#print axioms}).
\end{enumerate}

\medskip
\noindent\textbf{T\'emoins de coh\'erence et de port\'ee.} Deux instances
concr\`etes accompagnent le code pour en d\'emontrer la non-vacuit\'e et
la n\'ecessit\'e des hypoth\`eses.
\begin{itemize}
\item \texttt{witnessInstance} est une instance du mod\`ele abstrait v\'erifiant toutes les
      hypoth\`eses, avec $\Conf \cdot \Rel \cdot \Opp = 1/\lambda$
      exactement (\'egalit\'e dans la borne rel\^ach\'ee~; voir la
      Remarque~\ref{rem:portee_temoin}).
\item \texttt{scopeInstance} est une instance v\'erifiant les Hypoth\`eses 1,
      2$'$ et 3 mais hors du r\'egime de l'Hypoth\`ese 4, avec
      $\Conf \cdot \Rel \cdot \Opp = 1/4$ constant.
\end{itemize}
\end{sloppypar}

% ----------------------------------------------------------------------------
\section{GUIDE DE LECTURE}
\label{sec:lean4_guide}
% ----------------------------------------------------------------------------

\begin{sloppypar}
Pour les lecteurs non familiers avec Lean4, rappelons les conventions
syntaxiques utilis\'ees dans ce code.

\texttt{structure} d\'efinit un type de donn\'ees structur\'e (analogue
\`a un enregistrement). \texttt{def} introduit une d\'efinition nomm\'ee.
\texttt{lemma} et \texttt{theorem} introduisent un \'enonc\'e \`a prouver.
\texttt{by} ouvre un bloc de tactiques de preuve. \texttt{norm\_num} est
une tactique d'arithm\'etique num\'erique automatique. \texttt{linarith}
r\'esout des in\'egalit\'es lin\'eaires automatiquement. \texttt{nlinarith}
\'etend \texttt{linarith} aux in\'egalit\'es non lin\'eaires.
\texttt{field\_simp} simplifie les fractions rationnelles. \texttt{ring}
normalise les expressions polynomiales. L'expression
$\forall^{\mathtt{f}}\,n\;\texttt{in atTop},\;P\,n$ signifie \og pour tout
$n$ suffisamment grand, $P(n)$ est vrai \fg. \texttt{sorry} marque une
\'etape admise (aucune dans ce fichier).

L'instruction \texttt{import Mathlib} en t\^ete de fichier charge la biblioth\`eque
Mathlib, ce qui donne acc\`es aux lemmes d'analyse r\'eelle et d'alg\`ebre
utilis\'es.
\end{sloppypar}

% ----------------------------------------------------------------------------
\section{CODE SOURCE COMPLET}
\label{sec:lean4_code}
% ----------------------------------------------------------------------------

\begin{sloppypar}
Le code ci-dessous est la version compl\`ete et autonome de la
formalisation, identique au fichier \texttt{05\_formalisation\_lean/ConditionalCRO.lean} du
d\'ep\^ot~\cite{CROTrilemmaDepot}. Il peut \^etre copi\'e int\'egralement dans
\url{https://live.lean-lang.org} pour v\'erification.
\end{sloppypar}

\begin{lstlisting}[caption={Formalisation Lean4 du Trilemme CRO conditionnel (version 2.2)}, label={lst:lean4_complet}, language=Lean4]
/-
  Formalisation du Trilemme CRO Conditionnel — Annexe de thèse
  =============================================================
  Version 2.2 (octobre 2026) — fichier unique, autonome.

  Auteurs : Thierry Emmanuel MINKA Mi Nguidjoï,
            Flavien Serge Mani Onana,
            Thomas Djotio Ndié
  LIMSI — École Nationale Supérieure Polytechnique, Université de Yaoundé I

  MODE D'EMPLOI POUR L'ÉVALUATEUR
  1. Ouvrir https://live.lean-lang.org (projet Mathlib, choix par défaut).
  2. Coller l'intégralité de ce fichier et attendre la fin de la vérification.
  3. Résultat attendu : aucune erreur, aucun `sorry`. Le bloc final
     (`#print axioms`) n'affiche que propext, Classical.choice, Quot.sound.

  ARCHITECTURE
  • Couche « théorie de l'information » : Pinsker (Lemme 1) et Fano
    (Lemme 2, étape 1) sont des résultats classiques [Cover & Thomas] ; ils
    sont intégrés comme propriétés de toute instance de protocole.
  • Couche « cryptographie institutionnelle » : les Hypothèses 1 (extraction
    bornée par Δ), 2′ (verdict efficace, γ = αH(C) − 1), 3 (β-régularité)
    et 4 (régime de préservation de la confidentialité, I ≤ κ/λ) sont des
    hypothèses explicites des théorèmes.
  • Aucun `axiom` n'est déclaré. Le paramètre λ est noté `lam` (en Lean 4,
    `λ` est un mot réservé).
-/

import Mathlib

open Real Filter

namespace ConditionalCRO

/-! ## 1. Constante de Pinsker c₀ = 1/(2 ln 2) -/

noncomputable def c0 : ℝ := 1 / (2 * Real.log 2)

lemma c0_pos : 0 < c0 := by
  unfold c0
  have h : 0 < Real.log 2 := Real.log_pos (by norm_num)
  positivity

/-- c₀ > 1/2, car ln 2 < 1. -/
lemma c0_gt_half : 1 / 2 < c0 := by
  unfold c0
  have h1 := Real.log_two_lt_d9
  have h2 : 0 < Real.log 2 := Real.log_pos (by norm_num)
  rw [lt_div_iff₀ (by positivity)]
  nlinarith

/-! ## 2. Négligeabilité de ε₀(λ) = 2^(−λ) -/

/-- f est négligeable : ∀ c > 0, |f(n)| ≤ n^(−c) à partir d'un certain rang. -/
def Negligible (f : ℕ → ℝ) : Prop :=
  ∀ c : ℝ, 0 < c → ∀ᶠ n : ℕ in atTop, |f n| ≤ (n : ℝ) ^ (-c)

noncomputable def ε₀ (lam : ℕ) : ℝ := (1 / 2 : ℝ) ^ lam

lemma ε₀_nonneg (lam : ℕ) : 0 ≤ ε₀ lam := by
  unfold ε₀; positivity

theorem negligible_eps0 : Negligible ε₀ := by
  intro c hc
  have hlim : Tendsto (fun n : ℕ => (n : ℝ) ^ (⌈c⌉₊) / (2 : ℝ) ^ n) atTop (nhds 0) :=
    tendsto_pow_const_div_const_pow_of_one_lt ⌈c⌉₊ (by norm_num)
  have hev : ∀ᶠ n : ℕ in atTop, (n : ℝ) ^ (⌈c⌉₊) / (2 : ℝ) ^ n < 1 :=
    hlim.eventually (gt_mem_nhds one_pos)
  filter_upwards [hev, eventually_ge_atTop 1] with n hn hn1
  have hn1' : (1 : ℝ) ≤ n := by exact_mod_cast hn1
  have hnpos : (0 : ℝ) < n := by linarith
  have h2pos : (0 : ℝ) < 2 ^ n := by positivity
  have hck : (n : ℝ) ^ c ≤ (n : ℝ) ^ (⌈c⌉₊) := by
    rw [← Real.rpow_natCast]
    exact Real.rpow_le_rpow_of_exponent_le hn1' (Nat.le_ceil c)
  have hk : (n : ℝ) ^ (⌈c⌉₊) ≤ 2 ^ n := by
    rw [div_lt_one h2pos] at hn
    exact hn.le
  have hcpos : (0 : ℝ) < (n : ℝ) ^ c := Real.rpow_pos_of_pos hnpos c
  unfold ε₀
  rw [abs_of_nonneg (by positivity), Real.rpow_neg hnpos.le, one_div, inv_pow]
  exact inv_anti₀ hcpos (hck.trans hk)

/-! ## 3. Paramètres institutionnels -/

structure InstitutionalParams where
  H : ℝ      -- H(C) : entropie du contexte institutionnel (bits)
  Δ : ℝ      -- capacité d'extraction institutionnelle (bits)
  α : ℝ      -- granularité des verdicts
  β : ℝ      -- constante de régularité du protocole
  hH_nonneg : 0 ≤ H
  hΔ_nonneg : 0 ≤ Δ
  hα_pos    : 0 < α
  hβ_ge_one : 1 ≤ β

/-- γ = αH(C) − 1 : borne inférieure de H(V|C) (Hypothèse 2′). -/
def gamma (p : InstitutionalParams) : ℝ := p.α * p.H - 1

-- Familles canoniques (valeurs empiriques de l'article CLO)
noncomputable def HeliosParams : InstitutionalParams :=
  { H := 50, Δ := 56, α := 1.2, β := 1.4,
    hH_nonneg := by norm_num, hΔ_nonneg := by norm_num,
    hα_pos    := by norm_num, hβ_ge_one := by norm_num }

noncomputable def zkSNARKParams : InstitutionalParams :=
  { H := 70, Δ := 73, α := 1.5, β := 1.5,
    hH_nonneg := by norm_num, hΔ_nonneg := by norm_num,
    hα_pos    := by norm_num, hβ_ge_one := by norm_num }

noncomputable def TraditionalAuditParams : InstitutionalParams :=
  { H := 20, Δ := 35, α := 1.0, β := 1.0,
    hH_nonneg := by norm_num, hΔ_nonneg := by norm_num,
    hα_pos    := by norm_num, hβ_ge_one := by norm_num }

/-! ## 4. Instance de protocole

Grandeurs d'une instance au paramètre de sécurité `lam` (M uniforme,
H(M) = λ, M ⊥ C) :
  Conf = 1 − max TV (Déf. 7), Rel (Déf. 8), Opp = I(V;Σ|C)/H(V|C) (Déf. 9).
Les champs `pinsker` et `fano` encodent les Lemmes 1 et 2 de l'article
(étape 1 de la preuve : I(M;Σ|C) ≥ Rel·λ − 1, valable pour tout Rel). -/
structure ProtocolInstance (lam : ℕ) where
  params : InstitutionalParams
  Conf   : ℝ   -- confidentialité
  Rel    : ℝ   -- fiabilité
  Opp    : ℝ   -- opposabilité
  I      : ℝ   -- I(M;Σ|C)
  IV     : ℝ   -- I(V;Σ|C)
  HV     : ℝ   -- H(V|C)
  TVmax  : ℝ   -- max_{m₀,m₁,c} TV(P_{Σ|m₀,c}, P_{Σ|m₁,c})
  KLmax  : ℝ   -- max_{m₀,m₁,c} D_KL(P_{Σ|m₀,c} ‖ P_{Σ|m₁,c})
  hConf  : Conf = 1 - TVmax
  hOpp   : Opp = IV / HV
  hTV    : 0 ≤ TVmax ∧ TVmax ≤ 1
  hRel   : 0 ≤ Rel ∧ Rel ≤ 1
  hI     : 0 ≤ I
  hIV    : 0 ≤ IV
  pinsker : TVmax ^ 2 ≤ c0 * KLmax             -- Lemme 1 (Pinsker)
  fano    : I ≥ Rel * lam - 1                   -- Lemme 2 (Fano), étape 1

/-- Hypothèses institutionnelles 1, 2′ et 3. -/
structure InstitutionalAssumptions {lam : ℕ} (P : ProtocolInstance lam) : Prop where
  /-- Hypothèse 1 (extraction bornée) : I(V;Σ|C) ≤ I(M;Σ|C) + Δ. -/
  A1 : P.IV ≤ P.I + P.params.Δ
  /-- Hypothèse 2′ (verdict efficace) : H(V|C) ≥ γ = αH(C) − 1 … -/
  A2 : gamma P.params ≤ P.HV
  /-- … avec γ ≥ 1. -/
  A2γ : 1 ≤ gamma P.params
  /-- Hypothèse 3 (β-régularité) : D_KL ≤ β·I(M;Σ|C). -/
  A3 : P.KLmax ≤ P.params.β * P.I

/-! ## 5. Lemme 4 : garantie de confidentialité -/

/-- **Lemme 4.** Pour tout protocole β-régulier : Conf ≥ 1 − √(c₀·β·I(M;Σ|C)). -/
theorem lemma4_confidentiality_guarantee {lam : ℕ} (P : ProtocolInstance lam)
    (hA3 : P.KLmax ≤ P.params.β * P.I) :
    1 - Real.sqrt (c0 * P.params.β * P.I) ≤ P.Conf := by
  have hTV0 := P.hTV.1
  -- TV ≤ √(c₀·KL) (Pinsker)
  have h1 : P.TVmax ≤ Real.sqrt (c0 * P.KLmax) := by
    have := Real.abs_le_sqrt P.pinsker
    rwa [abs_of_nonneg hTV0] at this
  -- √(c₀·KL) ≤ √(c₀·β·I) (régularité)
  have h2 : Real.sqrt (c0 * P.KLmax) ≤ Real.sqrt (c0 * P.params.β * P.I) := by
    apply Real.sqrt_le_sqrt
    rw [mul_assoc]
    exact mul_le_mul_of_nonneg_left hA3 c0_pos.le
  rw [P.hConf]
  linarith

/-! ## 6. Théorème 1 : Trilemme CRO conditionnel -/

/-- **Théorème 1, forme explicite.** Sous les Hypothèses 1, 2′, 3 et 4
    (I(M;Σ|C) ≤ κ/λ), pour λ ≥ 1 :
    Conf·Rel·Opp ≤ (1/γ)·(Δ/λ + κ(Δ+1)/λ² + κ²/λ³). -/
theorem trilemma_explicit {lam : ℕ} (P : ProtocolInstance lam)
    (hA : InstitutionalAssumptions P) (κ : ℝ) (hκ : 0 ≤ κ)
    (hlam : 1 ≤ lam) (hA4 : P.I ≤ κ / lam) :
    P.Conf * P.Rel * P.Opp ≤
      (1 / gamma P.params) * (P.params.Δ / lam + κ * (P.params.Δ + 1) / (lam : ℝ) ^ 2
        + κ ^ 2 / (lam : ℝ) ^ 3) := by
  have hl : (1 : ℝ) ≤ lam := by exact_mod_cast hlam
  have hl0 : (0 : ℝ) < lam := by linarith
  have hγ1 := hA.A2γ
  have hγ0 : 0 < gamma P.params := by linarith
  have hHV : 0 < P.HV := by linarith [hA.A2]
  have hx0 := P.hI
  have hΔ := P.params.hΔ_nonneg
  have hκl : 0 ≤ κ / lam := div_nonneg hκ hl0.le
  -- Conf ∈ [0, 1]
  have hC1 : P.Conf ≤ 1 := by rw [P.hConf]; linarith [P.hTV.1]
  -- Étape 1 : Rel ≤ (I + 1)/λ (Fano)
  have hR0 := P.hRel.1
  have hR : P.Rel ≤ (P.I + 1) / lam := by
    rw [le_div_iff₀ hl0]; linarith [P.fano]
  -- Étape 2 : Opp ≤ (I + Δ)/γ (Hypothèses 1 et 2′)
  have hO0 : 0 ≤ P.Opp := by rw [P.hOpp]; exact div_nonneg P.hIV hHV.le
  have hO : P.Opp ≤ (P.I + P.params.Δ) / gamma P.params := by
    rw [P.hOpp, div_le_div_iff₀ hHV hγ0]
    have e1 := mul_le_mul_of_nonneg_right hA.A1 hγ0.le
    have e2 := mul_le_mul_of_nonneg_left hA.A2 (by linarith : (0 : ℝ) ≤ P.I + P.params.Δ)
    linarith
  -- Étape 3 : monotonie en I et Hypothèse 4
  have hRx : (P.I + 1) / lam ≤ (κ / lam + 1) / lam := by
    rw [div_le_div_iff_of_pos_right hl0]; linarith
  have hOx : (P.I + P.params.Δ) / gamma P.params ≤ (κ / lam + P.params.Δ) / gamma P.params := by
    rw [div_le_div_iff_of_pos_right hγ0]; linarith
  have hb0 : 0 ≤ (P.I + 1) / lam := div_nonneg (by linarith) hl0.le
  have hb0' : 0 ≤ (κ / lam + 1) / lam := div_nonneg (by linarith) hl0.le
  have hl0' : (lam : ℝ) ≠ 0 := hl0.ne'
  have hγ' : gamma P.params ≠ 0 := hγ0.ne'
  calc P.Conf * P.Rel * P.Opp
      ≤ 1 * ((P.I + 1) / lam) * ((P.I + P.params.Δ) / gamma P.params) :=
        mul_le_mul (mul_le_mul hC1 hR hR0 zero_le_one) hO hO0 (by rw [one_mul]; exact hb0)
    _ ≤ 1 * ((κ / lam + 1) / lam) * ((κ / lam + P.params.Δ) / gamma P.params) :=
        mul_le_mul (by rw [one_mul, one_mul]; exact hRx) hOx
          (div_nonneg (by linarith) hγ0.le) (by rw [one_mul]; exact hb0')
    _ = (1 / gamma P.params) * (P.params.Δ / lam + κ * (P.params.Δ + 1) / (lam : ℝ) ^ 2
          + κ ^ 2 / (lam : ℝ) ^ 3) := by
        field_simp
        ring

/-- **Théorème 1, forme publiée.** Comme γ ≥ 1 :
    Conf·Rel·Opp ≤ Δ/λ + κ(Δ+1)/λ² + κ²/λ³. -/
theorem trilemma_instance {lam : ℕ} (P : ProtocolInstance lam)
    (hA : InstitutionalAssumptions P) (κ : ℝ) (hκ : 0 ≤ κ)
    (hlam : 1 ≤ lam) (hA4 : P.I ≤ κ / lam) :
    P.Conf * P.Rel * P.Opp ≤
      P.params.Δ / lam + κ * (P.params.Δ + 1) / (lam : ℝ) ^ 2 + κ ^ 2 / (lam : ℝ) ^ 3 := by
  have h := trilemma_explicit P hA κ hκ hlam hA4
  have hl0 : (0 : ℝ) < lam := by exact_mod_cast hlam
  have hB : 0 ≤ P.params.Δ / lam + κ * (P.params.Δ + 1) / (lam : ℝ) ^ 2
      + κ ^ 2 / (lam : ℝ) ^ 3 := by
    have hΔ := P.params.hΔ_nonneg
    have hl2 : (0 : ℝ) < (lam : ℝ) ^ 2 := by positivity
    have hl3 : (0 : ℝ) < (lam : ℝ) ^ 3 := by positivity
    exact add_nonneg (add_nonneg (div_nonneg hΔ hl0.le)
      (div_nonneg (mul_nonneg hκ (by linarith)) hl2.le)) (div_nonneg (sq_nonneg κ) hl3.le)
  have hγ : 1 / gamma P.params ≤ 1 := by
    rw [div_le_one (by linarith [hA.A2γ])]; exact hA.A2γ
  exact h.trans (mul_le_of_le_one_left hB hγ)

/-- Famille de protocoles : paramètres institutionnels communs à toutes les instances. -/
structure ProtocolFamily where
  params  : InstitutionalParams
  inst    : (lam : ℕ) → ProtocolInstance lam
  hparams : ∀ lam, (inst lam).params = params

/-- **Théorème 1 (Trilemme CRO conditionnel), forme asymptotique.**
    Sous les Hypothèses 1, 2′, 3 et 4, il existe K et ε négligeable tels que,
    pour tout λ ≥ 64 : Conf·Rel·Opp ≤ Δ/λ + K/λ² + ε(λ). -/
theorem conditional_CRO_trilemma (F : ProtocolFamily) (κ : ℝ) (hκ : 0 ≤ κ)
    (hA : ∀ lam, InstitutionalAssumptions (F.inst lam))
    (hA4 : ∀ lam, (F.inst lam).I ≤ κ / lam) :
    ∃ K : ℝ, ∃ ε : ℕ → ℝ, Negligible ε ∧ ∀ lam : ℕ, 64 ≤ lam →
      (F.inst lam).Conf * (F.inst lam).Rel * (F.inst lam).Opp ≤
        F.params.Δ / lam + K / (lam : ℝ) ^ 2 + ε lam := by
  refine ⟨κ * (F.params.Δ + 1) + κ ^ 2, ε₀, negligible_eps0, fun lam hlam => ?_⟩
  have h := trilemma_instance (F.inst lam) (hA lam) κ hκ (by omega) (hA4 lam)
  rw [F.hparams lam] at h
  have hl : (1 : ℝ) ≤ lam := by exact_mod_cast (by omega : 1 ≤ lam)
  have hl0 : (0 : ℝ) < lam := by linarith
  have hε := ε₀_nonneg lam
  have h3 : κ ^ 2 / (lam : ℝ) ^ 3 ≤ κ ^ 2 / (lam : ℝ) ^ 2 := by
    have e : κ ^ 2 / (lam : ℝ) ^ 3 = κ ^ 2 / (lam : ℝ) ^ 2 * (1 / lam) := by
      field_simp
    rw [e]
    exact mul_le_of_le_one_right (by positivity) (by rw [div_le_one hl0]; exact hl)
  have hK : (κ * (F.params.Δ + 1) + κ ^ 2) / (lam : ℝ) ^ 2 =
      κ * (F.params.Δ + 1) / (lam : ℝ) ^ 2 + κ ^ 2 / (lam : ℝ) ^ 2 := add_div _ _ _
  linarith

/-! ## 7. Cohérence et portée des hypothèses -/

/-- Paramètres du témoin de cohérence : H(C) = 2, Δ = 1, α = 1, β = 1 (γ = 1). -/
noncomputable def WitnessParams : InstitutionalParams :=
  { H := 2, Δ := 1, α := 1, β := 1,
    hH_nonneg := by norm_num, hΔ_nonneg := by norm_num,
    hα_pos    := by norm_num, hβ_ge_one := by norm_num }

/-- Instance vérifiant toutes les hypothèses : I = 0, Rel = 1/λ, Conf = Opp = 1. -/
noncomputable def witnessInstance (lam : ℕ) (hlam : 1 ≤ lam) : ProtocolInstance lam :=
  { params := WitnessParams, Conf := 1, Rel := 1 / lam, Opp := 1, I := 0,
    IV := 1, HV := 1, TVmax := 0, KLmax := 0,
    hConf := by norm_num, hOpp := by norm_num, hTV := by norm_num,
    hRel := by
      have hl : (1 : ℝ) ≤ lam := by exact_mod_cast hlam
      exact ⟨by positivity, by rw [div_le_one (by linarith)]; exact hl⟩,
    hI := le_refl _, hIV := by norm_num,
    pinsker := by norm_num,
    fano := by
      have hl : (1 : ℝ) ≤ lam := by exact_mod_cast hlam
      rw [one_div_mul_cancel (ne_of_gt (by linarith))]
      norm_num }

/-- **Cohérence (non-vacuité).** Les Hypothèses 1, 2′, 3 et 4 sont
    simultanément satisfiables dans le modèle abstrait, et la borne du
    Théorème 1 y est atteinte : Conf·Rel·Opp = Δ/(γλ) = 1/λ. Le témoin est
    un point du modèle relâché ; il ne construit pas une loi de probabilité. -/
theorem assumptions_consistent_and_bound_attained (lam : ℕ) (hlam : 1 ≤ lam) :
    InstitutionalAssumptions (witnessInstance lam hlam) ∧
    (witnessInstance lam hlam).I ≤ 0 / lam ∧
    (witnessInstance lam hlam).Conf * (witnessInstance lam hlam).Rel *
      (witnessInstance lam hlam).Opp =
      (witnessInstance lam hlam).params.Δ / (gamma (witnessInstance lam hlam).params * lam) := by
  refine ⟨⟨?_, ?_, ?_, ?_⟩, ?_, ?_⟩ <;>
    simp [witnessInstance, WitnessParams, gamma] <;> norm_num

/-- Paramètres du témoin de portée : H(C) = 2, Δ = 0, α = 1, β = 2 (γ = 1). -/
noncomputable def ScopeParams : InstitutionalParams :=
  { H := 2, Δ := 0, α := 1, β := 2,
    hH_nonneg := by norm_num, hΔ_nonneg := le_refl _,
    hα_pos    := by norm_num, hβ_ge_one := by norm_num }

/-- Instance hors régime : I = λ/2 − 1 (révélation partielle du message). -/
noncomputable def scopeInstance (lam : ℕ) (hlam : 64 ≤ lam) : ProtocolInstance lam :=
  { params := ScopeParams, Conf := 1 / 2, Rel := 1 / 2, Opp := 1, I := lam / 2 - 1,
    IV := 1, HV := 1, TVmax := 1 / 2, KLmax := lam / 2,
    hConf := by norm_num, hOpp := by norm_num, hTV := by norm_num,
    hRel := by norm_num,
    hI := by
      have hl : (64 : ℝ) ≤ lam := by exact_mod_cast hlam
      linarith,
    hIV := by norm_num,
    pinsker := by
      have hl : (64 : ℝ) ≤ lam := by exact_mod_cast hlam
      have hc := c0_gt_half
      nlinarith,
    fano := by linarith }

/-- **Portée : l'Hypothèse 4 est nécessaire.** Pour tout λ ≥ 64, une instance
    satisfait les Hypothèses 1, 2′ et 3 (même avec Δ = 0) avec
    Conf·Rel·Opp = 1/4 ; elle viole l'Hypothèse 4 (I = λ/2 − 1). -/
theorem regime_assumption_necessary (lam : ℕ) (hlam : 64 ≤ lam) :
    InstitutionalAssumptions (scopeInstance lam hlam) ∧
    (scopeInstance lam hlam).params.Δ = 0 ∧
    (scopeInstance lam hlam).Conf * (scopeInstance lam hlam).Rel *
      (scopeInstance lam hlam).Opp = 1 / 4 ∧
    (scopeInstance lam hlam).I = lam / 2 - 1 := by
  have hl : (64 : ℝ) ≤ lam := by exact_mod_cast hlam
  refine ⟨⟨?_, ?_, ?_, ?_⟩, ?_, ?_, ?_⟩ <;>
    simp [scopeInstance, ScopeParams, gamma] <;> linarith

/-! ## 8. Régimes de capacité institutionnelle (Définition 10) -/

inductive InstitutionalRegime where
  | LowCapacity | MidCapacity | HighCapacity
  deriving DecidableEq, Repr

/-- Low si Δ < 0,1·λ ; Mid si 0,1·λ ≤ Δ ≤ 0,5·λ ; High sinon. -/
noncomputable def capacity_regime (Δ lam : ℝ) : InstitutionalRegime :=
  if Δ < 0.1 * lam then .LowCapacity
  else if Δ ≤ 0.5 * lam then .MidCapacity
  else .HighCapacity

-- Helios (Δ = 56) à λ = 128 : 12,8 ≤ 56 ≤ 64  →  Mid
theorem Helios_regime_at_128 :
    capacity_regime HeliosParams.Δ 128 = .MidCapacity := by
  show capacity_regime 56 128 = _
  unfold capacity_regime
  norm_num

-- zk-SNARKs (Δ = 73) à λ = 256 : 25,6 ≤ 73 ≤ 128  →  Mid
theorem zkSNARK_regime_at_256 :
    capacity_regime zkSNARKParams.Δ 256 = .MidCapacity := by
  show capacity_regime 73 256 = _
  unfold capacity_regime
  norm_num

-- Audit traditionnel ECDSA (Δ = 35) à λ = 128 : 12,8 ≤ 35 ≤ 64  →  Mid
theorem TraditionalAudit_regime_at_128 :
    capacity_regime TraditionalAuditParams.Δ 128 = .MidCapacity := by
  show capacity_regime 35 128 = _
  unfold capacity_regime
  norm_num

/-! ## 9. Verdict déterministe et forme raffinée -/

/-- **T1.** Verdict déterministe : I(V;Σ|C) = H(V|C) ⇒ Opp = 1. -/
theorem opp_eq_one_of_deterministic {lam : ℕ} (P : ProtocolInstance lam)
    (hdet : P.IV = P.HV) (hHV : 0 < P.HV) : P.Opp = 1 := by
  rw [P.hOpp, hdet]; exact div_self hHV.ne'

/-- **T1 bis.** Verdict déterministe + Hyp. 1 et 2′ ⇒ γ − I(M;Σ|C) ≤ Δ :
    la contrainte institutionnelle ne peut pas être active. -/
theorem deterministic_forces_capacity {lam : ℕ} (P : ProtocolInstance lam)
    (hA : InstitutionalAssumptions P) (hdet : P.IV = P.HV) :
    gamma P.params - P.I ≤ P.params.Δ := by
  have := hA.A1; have := hA.A2; linarith

/-- **T2 (forme raffinée).** Avec I(V;Σ|C) ≤ H(V|C) (toujours vrai) :
    Conf·Rel·Opp ≤ ((κ/λ+1)/λ) · min 1 ((κ/λ+Δ)/γ). -/
theorem trilemma_refined {lam : ℕ} (P : ProtocolInstance lam)
    (hA : InstitutionalAssumptions P) (κ : ℝ) (hκ : 0 ≤ κ)
    (hlam : 1 ≤ lam) (hA4 : P.I ≤ κ / lam) (hIVHV : P.IV ≤ P.HV) :
    P.Conf * P.Rel * P.Opp ≤
      ((κ / lam + 1) / lam) * min 1 ((κ / lam + P.params.Δ) / gamma P.params) := by
  have hl0 : (0 : ℝ) < lam := by exact_mod_cast hlam
  have hγ0 : 0 < gamma P.params := by linarith [hA.A2γ]
  have hHV : 0 < P.HV := by linarith [hA.A2, hA.A2γ]
  have hκl : 0 ≤ κ / lam := div_nonneg hκ hl0.le
  have hC1 : P.Conf ≤ 1 := by rw [P.hConf]; linarith [P.hTV.1]
  have hC0 : 0 ≤ P.Conf := by rw [P.hConf]; linarith [P.hTV.2]
  have hR0 := P.hRel.1
  have hR : P.Rel ≤ (κ / lam + 1) / lam := by
    rw [le_div_iff₀ hl0]; linarith [P.fano]
  have hO0 : 0 ≤ P.Opp := by rw [P.hOpp]; exact div_nonneg P.hIV hHV.le
  have hO1 : P.Opp ≤ 1 := by
    rw [P.hOpp, div_le_one hHV]; exact hIVHV
  have hO2 : P.Opp ≤ (κ / lam + P.params.Δ) / gamma P.params := by
    rw [P.hOpp, div_le_div_iff₀ hHV hγ0]
    have e1 := mul_le_mul_of_nonneg_right hA.A1 hγ0.le
    have e2 := mul_le_mul_of_nonneg_left hA.A2
      (by linarith [P.params.hΔ_nonneg] : (0 : ℝ) ≤ κ / lam + P.params.Δ)
    have e3 : P.I * gamma P.params ≤ κ / lam * gamma P.params :=
      mul_le_mul_of_nonneg_right hA4 hγ0.le
    nlinarith
  have hOm : P.Opp ≤ min 1 ((κ / lam + P.params.Δ) / gamma P.params) := le_min hO1 hO2
  calc P.Conf * P.Rel * P.Opp ≤ 1 * ((κ / lam + 1) / lam) * min 1 ((κ / lam + P.params.Δ) / gamma P.params) :=
        mul_le_mul (mul_le_mul hC1 hR hR0 zero_le_one) hOm hO0
          (by rw [one_mul]; exact div_nonneg (by linarith) hl0.le)
    _ = _ := by ring

/-- **T2 bis.** Si Δ ≥ γ, la borne raffinée coïncide avec la seule borne de Fano. -/
theorem refined_bound_is_fano_when_delta_ge_gamma {lam : ℕ} (P : ProtocolInstance lam)
    (hA : InstitutionalAssumptions P) (κ : ℝ) (hκ : 0 ≤ κ) (hlam : 1 ≤ lam)
    (hΔγ : gamma P.params ≤ P.params.Δ) :
    min 1 ((κ / lam + P.params.Δ) / gamma P.params) = 1 := by
  have hl0 : (0 : ℝ) < lam := by exact_mod_cast hlam
  have hγ0 : 0 < gamma P.params := by linarith [hA.A2γ]
  have hκl : 0 ≤ κ / lam := div_nonneg hκ hl0.le
  apply min_eq_left
  rw [le_div_iff₀ hγ0]; linarith


end ConditionalCRO

/-! ## 10. Bloc de vérification pour l'évaluateur

Chaque `#print axioms` doit afficher uniquement
`[propext, Classical.choice, Quot.sound]` : ni `sorryAx`, ni axiome ad hoc. -/

#check @ConditionalCRO.lemma4_confidentiality_guarantee
#check @ConditionalCRO.conditional_CRO_trilemma
#print axioms ConditionalCRO.lemma4_confidentiality_guarantee
#print axioms ConditionalCRO.trilemma_explicit
#print axioms ConditionalCRO.conditional_CRO_trilemma
#print axioms ConditionalCRO.assumptions_consistent_and_bound_attained
#print axioms ConditionalCRO.regime_assumption_necessary
#print axioms ConditionalCRO.zkSNARK_regime_at_256
#print axioms ConditionalCRO.opp_eq_one_of_deterministic
#print axioms ConditionalCRO.deterministic_forces_capacity
#print axioms ConditionalCRO.trilemma_refined
#print axioms ConditionalCRO.refined_bound_is_fano_when_delta_ge_gamma
\end{lstlisting}

% ----------------------------------------------------------------------------
\section{STATUT DE V\'ERIFICATION}
\label{sec:lean4_statut}
% ----------------------------------------------------------------------------

\begin{sloppypar}
Le tableau~\ref{tab:lean4_status} r\'ecapitule le statut de v\'erification
de chaque composant de la formalisation, tel qu'il ressort de la
version~2.2 du fichier \texttt{ConditionalCRO.lean}, v\'erifi\'ee sur
\url{https://live.lean-lang.org}~; la sortie de r\'ef\'erence du projet local est
\texttt{05\_formalisation\_lean/results/lean\_build.log} du
d\'ep\^ot~\cite{CROTrilemmaDepot}.
\end{sloppypar}

\begin{table}[htbp]
\centering
\caption{Statut de v\'erification des composants Lean4 (version~2.2)}
\label{tab:lean4_status}
\begin{tabular}{@{}p{7.2cm}p{3cm}p{3cm}@{}}
\toprule
\textbf{Composant} & \textbf{Statut} & \textbf{Note} \\
\midrule
Structures de base (\texttt{InstitutionalParams}, etc.)
  & $\checkmark$ V\'erifi\'e & Entier \\
Constante $c_0$ et propri\'et\'es
  & $\checkmark$ V\'erifi\'e & Entier \\
N\'egligeabilit\'e de $2^{-\lambda}$
  & $\checkmark$ V\'erifi\'e & Via Mathlib \\
Structure \texttt{ProtocolInstance}
  & $\checkmark$ V\'erifi\'e & Entier \\
Structure \texttt{InstitutionalAssumptions}
  & $\checkmark$ V\'erifi\'e & Entier \\
Lemme 4 (garantie de confidentialit\'e)
  & $\checkmark$ V\'erifi\'e & Entier \\
Th\'eor\`eme 1 (forme explicite)
  & $\checkmark$ V\'erifi\'e & Entier \\
Th\'eor\`eme 1 (forme publi\'ee)
  & $\checkmark$ V\'erifi\'e & Entier \\
Th\'eor\`eme 1 (forme asymptotique)
  & $\checkmark$ V\'erifi\'e & Entier \\
T\'emoin de coh\'erence (mod\`ele abstrait)
  & $\checkmark$ V\'erifi\'e & Entier \\
T\'emoin de port\'ee
  & $\checkmark$ V\'erifi\'e & Entier \\
Classification des r\'egimes (3 th\'eor\`emes)
  & $\checkmark$ V\'erifi\'e & \texttt{norm\_num} \\
Extensions v2.2 (4 th\'eor\`emes)
  & $\checkmark$ V\'erifi\'e & Entier \\
\midrule
\textbf{Total} & \textbf{0 sorry} & \textbf{0 axiome ad hoc} \\
\bottomrule
\end{tabular}
\end{table}

\begin{sloppypar}
\noindent\textbf{Axiomes utilis\'es.} Le bloc final \texttt{\#print axioms}
confirme que les th\'eor\`emes ne d\'ependent que des trois axiomes
standards de Lean/Mathlib, rappel\'es ci-dessous.
\begin{itemize}
\item \texttt{propext} (extensionnalit\'e propositionnelle),
\item \texttt{Classical.choice} (axiome du choix),
\item \texttt{Quot.sound} (soundness des quotients).
\end{itemize}
Ces axiomes sont utilis\'es par la quasi-totalit\'e des th\'eor\`emes de
Mathlib, y compris les plus \'el\'ementaires. Leur pr\'esence est
standard et ne remet pas en cause la validit\'e des preuves.
Aucun axiome \emph{ad hoc} (num\'erique ou informationnel) n'est d\'eclar\'e
dans ce fichier.
\end{sloppypar}

% ----------------------------------------------------------------------------
\section{INSTRUCTIONS DE REPRODUCTION}
\label{sec:lean4_reproduction}
% ----------------------------------------------------------------------------

\begin{sloppypar}
\textbf{Version de r\'ef\'erence.} Lean~v4.34.0 et Mathlib~v4.34.0, fig\'es dans les
fichiers \texttt{lean-toolchain}, \texttt{lakefile.toml} et
\texttt{lake-manifest.json} du d\'ep\^ot
\url{https://github.com/MINKAThierry/PhDThesis}~\cite{CROTrilemmaDepot}
(licence MIT).

\medskip
\noindent\textbf{Proc\'edure locale.}
\begin{enumerate}
\item Installer \texttt{elan} (gestionnaire de versions de Lean).
\item Cloner le d\'ep\^ot et se placer dans son r\'epertoire \texttt{05\_formalisation\_lean/}.
\item Ex\'ecuter \texttt{lake exe cache get} (t\'el\'echarge Mathlib pr\'ecompil\'ee),
      puis \texttt{lake build} (environ une minute et demie).
\item Ex\'ecuter \texttt{lake env lean ConditionalCRO.lean}, qui affiche les sorties des
      commandes \texttt{\#check} et \texttt{\#print axioms}.
\end{enumerate}

\medskip
\noindent\textbf{Proc\'edure dans le navigateur.}
\begin{enumerate}
\item Ouvrir \url{https://live.lean-lang.org} et choisir le projet \og Mathlib\fg.
\item Coller l'int\'egralit\'e du Listing~\ref{lst:lean4_complet} et attendre la
      fin de la v\'erification. Le site utilise la version courante de Mathlib~;
      le fichier y a \'et\'e v\'erifi\'e en octobre 2026.
\end{enumerate}

\medskip
\noindent\textbf{R\'esultat attendu.} La compilation doit satisfaire les conditions suivantes.
\begin{itemize}
\item Aucune erreur de compilation.
\item Aucun \texttt{sorry} (pas de \texttt{sorryAx} dans les listes d'axiomes).
\item Aucun axiome \emph{ad hoc} d\'eclar\'e.
\item Chacun des dix \texttt{\#print axioms} affiche exactement
      \texttt{[propext, Classical.choice, Quot.sound]}.
\end{itemize}

\medskip
\noindent\textbf{D\'ep\^ot.} Le d\'ep\^ot contient aussi les scripts et donn\'ees
du chapitre~\ref{chap:experimentation}, une batterie de tests qui recalcule
les valeurs num\'eriques cit\'ees dans cette th\`ese ainsi qu'une int\'egration continue
qui recompile la formalisation \`a chaque modification.
\end{sloppypar}
% ============================================================================
% PUBLICATIONS ISSUES DE LA TH\`eSE
% ============================================================================
%\markboth{PUBLICATIONS ISSUES DE LA TH\`ESE}{}
\chapter{Publications issues de la th\`ese}
%\addcontentsline{toc}{chapter}{Publications Issues de la Th\`ese}
\label{chap:annexe_publications}
\begin{sloppypar}


Cette annexe recense les publications scientifiques issues des travaux de
cette th\`ese. Seules sont list\'ees les publications effectivement fournies
dans le dossier de th\`ese, publi\'ees dans des revues \`a comit\'e de lecture.
Trois articles r\'epondent \`a ce crit\`ere \`a la date de soutenance~; ils
couvrent les contributions C1, C2 et C3, soit les \'etapes F, B et E de la
progression F$\to$B$\to$E$\to$S$\to$C$\to$X. Le statut et la r\'ef\'erence de
l'ensemble des contributions figurent au tableau~\ref{tab:contributions_these}.
\end{sloppypar}



% ============================================================================
\section{ARTICLES PUBLI\'ES}
\addcontentsline{toc}{section}{Articles publi\'es}
\label{sec:publications_publies}
% ============================================================================
\begin{sloppypar}
\subsection{[P1] The CLO Framework}

\noindent\textbf{Auteurs} : Thierry Emmanuel MINKA MI NGUIDJO\"{I},
Flavien Serge MANI ONANA, Thomas DJOTIO NDI\'E.

\noindent\textbf{R\'ef\'erence} : \textit{Cryptographic Legal Opposability (CLO) 
Framework: Confidentiality-Reliability-Opposability Trade-offs in Institutional 
Cryptography}. \textit{IEEE Access}, 14, 20427--20445, 2026.
\href{https://doi.org/10.1109/ACCESS.2026.3660105}{\texttt{doi:10.1109/ACCESS.2026.3660105}}

\noindent\textbf{Statut} : Publi\'e (2 f\'evrier 2026).

\noindent\textbf{Contribution} : Cet article \'etablit les fondations de la 
cryptographie institutionnelle (Contribution C1, \'etape \textbf{F} de la progression
\textbf{F}$\to$B$\to$E$\to$S$\to$C$\to$X). Il introduit les m\'etriques CLO ($\Conf$, $\Rel$, $\Opp$), 
le th\'eor\`eme du goulot contextuel (dans la formulation plus pr\'ecise
$I(V;\Sigma|C) \leq \alpha H(\mathcal{C})$ retenue par la th\`ese)
et le paradigme architectural Gold-Iron-Clay. Il constitue la section~\ref{sec:clo}
de cette th\`ese et ouvre la voie aux perspectives de l'Axe~I 
(fondations computationnelles AIIP/CASH Toolkit) identifi\'ees \`a la section~\ref{sec:perspectives}.

\vspace{0.5cm}

\subsection{[P2] The Conditional CRO Trilemma}

\noindent\textbf{Auteurs} : Thierry Emmanuel MINKA MI NGUIDJO\"{I},
Flavien Serge MANI ONANA, Thomas DJOTIO NDI\'E.

\noindent\textbf{R\'ef\'erence} : \textit{The Conditional CRO Trilemma: 
Impossibility Bounds on Confidentiality, Reliability, and Opposability 
under Institutional Extraction Limits}. \textit{IEEE Access}, 14, 33263--33274, 2026.
\href{https://doi.org/10.1109/ACCESS.2026.3669065}{\texttt{doi:10.1109/ACCESS.2026.3669065}}

\noindent\textbf{Statut} : Publi\'e (27 f\'evrier 2026).

\noindent\textbf{Contribution} : Cet article \'etablit la borne 
d'impossibilit\'e fondamentale (Contribution C2, \'etape \textbf{B} de la progression 
F$\to$\textbf{B}$\to$E$\to$S$\to$C$\to$X)~:
\begin{equation}
\Conf \cdot \Rel \cdot \Opp \leq \frac{\Delta}{\lambda} + O(\lambda^{-2}) + \negl(\lambda)
\end{equation}
Il introduit les param\`etres institutionnels $(\Delta, \alpha, \beta)$ et
propose une calibration sur trois syst\`emes d\'eploy\'es. La th\`ese en donne la
forme explicite $\gamma^{-1}(\Delta/\lambda + \dots)$ et la forme raffin\'ee
(section~\ref{sec:cro_classique}) tout en pr\'ecisant que la fiabilit\'e de la borne est
la fiabilit\'e de reconstruction $\Relrec$. Il constitue la section~\ref{sec:cro_classique} de cette th\`ese et 
fonde les perspectives de l'Axe III (extensions th\'eoriques vers la borne forte, 
FHE et MPC).

\vspace{0.5cm}

\subsection{[P3] The Post-Quantum Conditional CRO Trilemma}

\noindent\textbf{Auteurs} : Thierry Emmanuel MINKA MI NGUIDJO\"{I},
Flavien Serge MANI ONANA, Thomas DJOTIO NDI\'E.

\noindent\textbf{R\'ef\'erence} : \textit{Post-Quantum Extension of the
Conditional CRO Trilemma: Impossibility Bounds on Confidentiality,
Reliability, and Opposability under Quantum Institutional Extraction Limits}.
\textit{IEEE Access}, 14, 125130--125155,  2026.
\href{https://doi.org/10.1109/ACCESS.2026.3722742}{\texttt{doi:10.1109/ACCESS.2026.3722742}}

\noindent\textbf{Statut} : Publi\'e (12 ao\^ut 2026).

\noindent\textbf{Contribution} : Cet article \'etend la borne d'impossibilit\'e
au r\'egime quantique (Contribution C3, \'etape \textbf{E} de la progression
F$\to$B$\to$\textbf{E}$\to$S$\to$C$\to$X)~:
\begin{equation}
\Conf^q \cdot \Rel^q \cdot \Opp^q \leq \frac{\Delta_q + \eta_q \cdot H(\mathcal{C})}{\lambda}
  + O(\lambda^{-2})
\end{equation}
Il introduit la capacit\'e d'extraction institutionnelle quantique $\Delta_q$
et le coefficient de perte d'interpr\'etabilit\'e quantique $\eta_q$. La th\`ese
en retient une d\'efinition unique (une perte, nulle pour les preuves qui sont
des cha\^{\i}nes de bits classiques) et pr\'ecise que le terme $\eta_q H(\mathcal{C})$
n'appara\^it que dans la majoration affaiblie. L'analyse de calibration
porte sur les signatures ML-DSA et SLH-DSA~; ML-KEM, m\'ecanisme d'encapsulation
de cl\'e, est hors du champ du Trilemme.
Il constitue la section~\ref{sec:cro_pq} de cette th\`ese~; les protocoles de
mesure de $\Delta_q$ et $\eta_q$ qu'il appelle sont sp\'ecifi\'es au
chapitre~\ref{chap:experimentation}, tandis que ses prolongements rel\`event de l'Axe~III
(cadre quantique quantitatif) identifi\'e \`a la
section~\ref{sec:perspectives}.

\end{sloppypar}
\cloturerAnnexes



\end{document}
